\documentclass[11pt]{article}

\usepackage[T1]{fontenc}
\usepackage{lmodern}
\usepackage[margin=1in]{geometry}
\usepackage{amsmath,amssymb,amsthm,mathtools}
\usepackage{enumitem}
\usepackage{booktabs,longtable,array}
\usepackage{xcolor}
\usepackage[colorlinks=true,linkcolor=blue!55!black,urlcolor=blue!55!black]{hyperref}
\usepackage[nameinlink,noabbrev]{cleveref}

\newtheorem{theorem}{Theorem}[section]
\newtheorem{lemma}[theorem]{Lemma}
\newtheorem{proposition}[theorem]{Proposition}
\newtheorem{corollary}[theorem]{Corollary}
\newtheorem{counterexample}[theorem]{Exact separator}
\newtheorem{openproblem}[theorem]{Open problem}
\theoremstyle{definition}
\newtheorem{definition}[theorem]{Definition}
\newtheorem{assessment}[theorem]{Assessment}
\newtheorem{firewall}[theorem]{Firewall}
\newtheorem{programme}[theorem]{Programme}
\theoremstyle{remark}
\newtheorem{remark}[theorem]{Remark}

\newcommand{\C}{\mathbb C}
\newcommand{\R}{\mathbb R}
\newcommand{\Ran}{\operatorname{Ran}}
\newcommand{\Ker}{\operatorname{Ker}}
\newcommand{\Tr}{\operatorname{Tr}}
\newcommand{\dd}{\mathrm d}
\newcommand{\id}{\mathrm{id}}

\title{Renewal Standard--Model--Einstein Continuum Research Notes\\
Fourth Volume, Version V1: Result Map and the Signed Gaussian Parent Measure}
\author{Append-only research continuation}
\date{September 2026}

\begin{document}
\maketitle

\begin{abstract}
This note opens the fourth volume of the Renewal Standard--Model--Einstein
continuum programme after freezing the third volume at V100.  It begins with
a compact map of the rigorous results now available and the hypotheses which
remain open.  The new argument then implements the proof order forced by the
third-volume Gaussian counterexample.  The bosonic quadratic action and the
complete signed polarization are normalized as one parent Gaussian measure.
Exact shell Schur complements preserve positivity and make the determinant
normalizations telescope.  Uniform ellipticity gives cutoff-independent
tightness of gauge potentials in every \(H^{-1-\varepsilon}\) and of
curvatures in every \(H^{-2-\varepsilon}\).  A relative-entropy lemma then
transfers this tightness to nonlinear cutoff measures if their entropy with
respect to the combined parent is uniformly controlled.  Reflection
positivity requires an additional positive spectral representation, and the
hypercharge ultraviolet direction remains open.  No full interacting
SM--Einstein continuum is claimed.
\end{abstract}

\tableofcontents

\section{Context, archive lineage, and objective}
\label{sec:v4v1-context}

The target remains a common continuum construction for the gauge, Higgs,
fermion, ghost, and metric sectors in which the normalized observables,
physical state space, determinant line, composite stress tensor, Ward
identities, and sourced Einstein evolution arise from compatible regulators.
Anomaly cancellation, existence of a measure, reflection positivity,
composite-field convergence, and nonlinear Einstein stability are separate
claims.  They are not substitutes for one another.

The frozen proof archive is now
\begin{center}
\begin{tabular}{lll}
\toprule
volume & frozen file & SHA--256\\
\midrule
first &
\texttt{Renewal\_SM\_Einstein\_Continuum\_Research\_Notes\_V157\_f1.tex}
&\texttt{EA08CE75\ldots563AAD94}\\
second &
\texttt{Renewal\_SM\_Einstein\_Continuum\_Research\_Notes\_V100\_f2.tex}
&\texttt{497249F6\ldotsCB6E673E}\\
third &
\texttt{Renewal\_SM\_Einstein\_Continuum\_Research\_Notes\_V100\_f3.tex}
&\texttt{E93FC354\ldotsF137D5BB}\\
\bottomrule
\end{tabular}
\end{center}
The complete third-volume digest is
\begin{equation}
 \texttt{E93FC354593735071C6742E06ABFC240C617B3CB8120A5A9EFBA4970F137D5BB}.
\label{eq:v4v1-f3-hash}
\end{equation}
The frozen files retain the full hypotheses and proofs.  The map below is a
guide and does not replace them.  The supplied monograph and programme files
remain mathematical inputs; their proposed steps are not promoted to
theorems without proof.

\section{Major mathematical results obtained so far}
\label{sec:v4v1-results}

\subsection{Finite-cutoff calculus, Ward structure, and stress response}
\label{sec:v4v1-frozen-early}

The first two frozen volumes construct a finite-cutoff BV pushforward which
preserves the quantum master equation and composes under shell integration.
Connected-cumulant and coordinate-resampling identities retain one physical
shell occurrence and prevent false products of covariance marks.  Local
Taylor jets are separated from irrelevant remainders; the four-dimensional
degree-four shell is marginal until the complete local jet is subtracted.

The metric derivative is decomposed into Ward-completed contact terms and a
causal stress--stress response.  Fourth-order spectral growth produces a
nonlocal logarithmic polarization and disproves a naive same-order
second-order source estimate.  Conditional graph-space and Volterra
theorems control the response after its linear polarization is inverted
together with the Einstein operator.  Compactness results isolate the
microlocal, diagonal-remainder, uniform-integrability, and metric-parametric
hypotheses needed for convergence of the composite stress.  Explicit
concentrating sequences show that ordinary energy bounds can leave a stress
defect.

\subsection{Boundary complexes and causal Einstein evolution}
\label{sec:v4v1-boundary}

The third volume constructs BRST-compatible boundary lifts by homological
repair and identifies the connecting cohomology class as the exact
obstruction.  Hodge--Poisson collar extensions, relative cokernel lifts,
one-port energy identities, moving-boundary Kato transport, Feshbach memory,
and maximal dissipativity are proved under uniform gaps.  The general
symmetric-hyperbolic boundary symbol is analyzed through glancing, and local
physical reservoirs repair its explicitly identified scalar and elastic
roots.

The resulting TT--material--bath junction yields a common boundary domain,
linear solution operator, invariant Picard ball, constraint propagation,
and a conditional classical local Einstein--matter theorem.  These results
retain a low-energy causal margin and a bounded-geometry time interval; they
do not prove all-frequency poles absent or global nonlinear stability.

\subsection{Chiral regulator, anomaly, and determinant line}
\label{sec:v4v1-fermion}

One finite regulator is built for the gauge, constraint, carrier, and chiral
sectors.  The local Standard Model anomaly coefficients cancel in the
physical left-handed convention, including the mixed
gauge--gravitational coefficient.  Domain-wall analysis proves the exact
one-species Wilson degree, kink--antikink mixing, heavy resolvent bounds,
local correlator convergence, and anomaly inflow.  The finite overlap Schur
complement satisfies the Ginsparg--Wilson algebra.

The later third-volume chain proves a volume-uniform Wilson gap for
quantitatively admissible almost-commuting shifts, an overlap local-index
theorem on smooth flat bundle sectors, and the corresponding family
determinant curvature.  The one-species Wilson symbol is identified with the
Spin(4)-equivariant Bott--Thom class.  Its curved pairing is
\begin{equation}
 \operatorname{ind}D^+_{E_R}
 =\int_X[\widehat A(TX)\operatorname{ch}(E_R)]_4.
\label{eq:v4v1-curved-index}
\end{equation}
A two-scale finite-element compression and geodesic tail cut give finite
curved matrices with this exact index.  A diagonal schedule makes their
first two family derivatives summably close to the global regulator, though
the mesh-unit propagation range grows.

Spectral sections and generalized norm-resolvent comparison control
physical zero modes, crossing orientations, and determinant transition
lines without identifying physical and Wilson eigenvalues.  The actual
global gauge group is
\(G_{\rm SM}=S(U(3)\times U(2))\), and the calculation
\begin{equation}
 \Omega_5^{\rm Spin}(BG_{\rm SM})=0
\label{eq:v4v1-bordism-zero}
\end{equation}
removes the residual Standard Model global gauge-anomaly character on spin
manifolds.  Kato transport along the compression homotopy preserves torsion
and gives an exact determinant-line isomorphism.

\subsection{Renormalized determinant and rough geometric sectors}
\label{sec:v4v1-rough}

The regulator eta form has one exact heat-time partition.  Its UV part is a
finite local jet plus an \(O(a)\) remainder; its complementary IR part is the
physical Bismut--Freed determinant connection with the correct crossing
transitions.  An admissible rough link field reconstructs a genuine smooth
\(G_{\rm SM}\)-bundle with exact edge holonomies and satisfies the one-use
bound
\begin{equation}
 \|F_{A_U}\|_{L^2}^2\le C\sum_p\|1-U_p\|^2.
\label{eq:v4v1-reconstruction-bound}
\end{equation}
The derivative bounds deteriorate at the critical scale, so this does not
put typical quantum fields into a smooth bounded-geometry set.

On concentration balls, Uhlenbeck gauge gives \(A\in L^4\).  The Dirac
perturbation is in weak \(\mathcal S_4\), and the fifth regularized
determinant has an action-absorbable upper bound.  The four lower physical
Born terms assemble by Ward identities and require the logarithmic curvature
moment
\begin{equation}
 \langle F,\log(1+\Delta_A/\mu^2)F\rangle.
\label{eq:v4v1-log-moment}
\end{equation}
A transverse plane wave proves that \(L^2\)-curvature and plaquette
admissibility do not bound this quantity.

The V99 Gaussian test proves that a separate cutoff-uniform exponential
moment of \eqref{eq:v4v1-log-moment} fails even in the free field.  The
correct parent object is the signed sum of the bosonic inverse covariance
and the complete polarization.  The nonabelian Standard Model one-loop signs
are asymptotically free; the hypercharge sign is Landau-like and remains an
ultraviolet obstruction.  Maximum plaquette admissibility is replaced by a
conditional punctured bad-cluster reconstruction.

\subsection{Current theorem boundary}
\label{sec:v4v1-boundary-status}

\begin{firewall}[The full continuum has not been reached]
\label{fw:v4v1-not-complete}
No frozen volume proves the interacting four-dimensional Yang--Mills measure
and mass gap, a nontrivial hypercharge ultraviolet completion, an Einstein
metric measure, the punctured bad-cluster estimate, reflection positivity of
the curved compressed fermion action, composite-stress uniform
integrability, or global nonlinear Einstein stability.  The third-volume
V100 theorem proves that these hypotheses would close the assembled chain;
it does not prove them.
\end{firewall}

\section{The signed Gaussian parent at finite cutoff}
\label{sec:v4v1-parent}

We now continue the programme.  Let \(\mathcal H_\Lambda\) be the real
finite-dimensional transverse gauge-field space below momentum cutoff
\(\Lambda\), after the finite ghost and longitudinal cancellations have
been performed.  Let
\begin{equation}
 K_\Lambda=K_{\Lambda}^{\rm bos}+\Pi_\Lambda
\label{eq:v4v1-total-kernel}
\end{equation}
be the complete renormalized quadratic kernel, including gauge, ghost,
Higgs, fermion, and any proved metric polarization.  Assume for the moment
that
\begin{equation}
 c_-\langle A,(-\Delta)A\rangle
 \le\langle A,K_\Lambda A\rangle
 \le c_+\langle A,(1-\Delta)A\rangle
\label{eq:v4v1-uniform-ellipticity}
\end{equation}
on the transverse space, with constants independent of \(\Lambda\).

Define the normalized parent measure
\begin{equation}
 d\nu_\Lambda(A)=
 Z_\Lambda^{-1}
 \exp\left(-\frac12\langle A,K_\Lambda A\rangle\right)dA,
 \qquad
 Z_\Lambda=(2\pi)^{N_\Lambda/2}\det(K_\Lambda)^{-1/2}.
\label{eq:v4v1-parent-measure}
\end{equation}

\subsection{Exact shell normalization}
\label{sec:v4v1-shells}

Let \(\mathcal H_\Lambda=H_<\oplus H_>\) be one orthogonal shell split and
write
\begin{equation}
 K_\Lambda=
 \begin{pmatrix}A&B^*\\B&D\end{pmatrix}.
\label{eq:v4v1-block-kernel}
\end{equation}

\begin{theorem}[Quadratic shell Schur identity]
\label{thm:v4v1-Schur-shell}
If \(K_\Lambda>0\), then \(D>0\), the Schur complement
\begin{equation}
 S=A-B^*D^{-1}B
\label{eq:v4v1-Schur-complement}
\end{equation}
is positive, and integrating the high shell gives exactly
\begin{equation}
 \int_{H_>}
 e^{-\frac12\langle(A_<,A_>),K_\Lambda(A_<,A_>)\rangle}
 dA_>
 =(2\pi)^{\dim H_>/2}\det(D)^{-1/2}
 e^{-\frac12\langle A_<,SA_<\rangle}.
\label{eq:v4v1-shell-integral}
\end{equation}
Moreover
\begin{equation}
 \det K_\Lambda=\det D\,\det S.
\label{eq:v4v1-Schur-determinant}
\end{equation}
Thus iterated shell normalizations telescope to the one full determinant in
\eqref{eq:v4v1-parent-measure} even when the polarization has cross-shell
blocks.
\end{theorem}

\begin{proof}
Complete the square:
\begin{align}
 \langle(A_<,A_>),K_\Lambda(A_<,A_>)\rangle
 &=\langle A_>,D A_>\rangle
 +2\langle BA_<,A_>\rangle+\langle A_<,AA_<\rangle\nonumber\\
 &=\langle A_>+D^{-1}BA_<,
 D(A_>+D^{-1}BA_<)\rangle
 +\langle A_<,SA_<\rangle.
\label{eq:v4v1-complete-square}
\end{align}
Positivity of \(D\) follows by restriction.  Positivity of \(S\) follows by
evaluating the positive form on
\((A_<,-D^{-1}BA_<)\).  Translation of the high variable and the ordinary
Gaussian integral prove \eqref{eq:v4v1-shell-integral}.  Block Gaussian
elimination proves \eqref{eq:v4v1-Schur-determinant}.  Repeating the identity
over a finite shell tower makes adjacent determinant factors cancel against
the normalization inherited from the next marginal, leaving the full
determinant exactly once.
\end{proof}

\begin{firewall}[A bounded cross-shell block is not removable]
\label{fw:v4v1-cross-shell}
The term \(B^*D^{-1}B\) is the exact response of the eliminated shell.  An
operator-norm bound on \(B\) does not permit replacing \(S\) by \(A\).  Such
a replacement changes both the low covariance and the determinant
normalization and repeats the V100 ancestry error.
\end{firewall}

\subsection{Distributional tightness of the parent}
\label{sec:v4v1-tightness}

Let \(H^{-s}\) denote the transverse Sobolev space on a fixed compact
four-manifold, defined by a reference Laplacian.  Extend each cutoff field by
zero on modes above \(\Lambda\).

\begin{theorem}[Gaussian parent tightness]
\label{thm:v4v1-Gaussian-tightness}
Under \eqref{eq:v4v1-uniform-ellipticity}, the family
\(\{\nu_\Lambda\}\) is tight in \(H^{-1-\varepsilon}\) for every
\(\varepsilon>0\).  The induced curvature distributions are tight in
\(H^{-2-\varepsilon}\).  Every subsequential limit is a centred Gaussian
distribution with covariance \(C\) satisfying, as quadratic forms,
\begin{equation}
 c_+^{-1}(1-\Delta)^{-1}
 \le C\le c_-^{-1}(-\Delta)^{-1}
\label{eq:v4v1-covariance-bounds}
\end{equation}
on the transverse nonzero modes.
\end{theorem}

\begin{proof}
Let \(\lambda_j\asymp j^{1/2}\) be the four-dimensional Laplacian
eigenvalues.  The covariance bound gives
\begin{equation}
 \mathbb E_{\nu_\Lambda}
 \|A\|_{H^{-1-\varepsilon/2}}^2
 \le C\sum_j(1+\lambda_j)^{-1-\varepsilon/2}\lambda_j^{-1}.
\label{eq:v4v1-Sobolev-second-moment}
\end{equation}
In momentum notation the summand density is
\(r^3r^{-2-\varepsilon}r^{-2}dr=r^{-1-\varepsilon}dr\), which is
integrable at infinity.  The bound is uniform in \(\Lambda\).  Closed balls
of \(H^{-1-\varepsilon/2}\) are compact in
\(H^{-1-\varepsilon}\).  Markov's inequality therefore gives tightness.

One derivative maps \(H^{-1-\varepsilon}\) continuously into
\(H^{-2-\varepsilon}\), proving curvature tightness.  Prokhorov gives
subsequential limits.  Their characteristic functions are limits of finite
Gaussian characteristic functions and hence Gaussian.  Passing the uniform
quadratic-form inequalities for the covariances to cylinder test functions
gives \eqref{eq:v4v1-covariance-bounds}.
\end{proof}

This theorem puts the free parent on its correct support.  Typical curvature
is a distribution in \(H^{-2-\varepsilon}\), not a bounded classical
\(L^2\) field.  The V96 reconstruction and V97 stopping balls must therefore
be applied to controlled events or renormalized block fields, not declared
typical pointwise configurations.

\begin{corollary}[Uniqueness from covariance convergence]
\label{cor:v4v1-Gaussian-uniqueness}
If
\begin{equation}
 \langle f,K_\Lambda^{-1}g\rangle
 \longrightarrow\langle f,Cg\rangle
\label{eq:v4v1-cylinder-covariance}
\end{equation}
for every pair of smooth transverse test forms, then the entire parent
sequence converges weakly in \(H^{-1-\varepsilon}\) to the Gaussian measure
of covariance \(C\).
\end{corollary}

\begin{proof}
Equation \eqref{eq:v4v1-cylinder-covariance} gives convergence of every
finite-dimensional characteristic function.  Tightness from
Theorem~\ref{thm:v4v1-Gaussian-tightness} makes every subsequence have a
limit, and the characteristic functions identify every limit with the same
Gaussian measure.  Hence the full sequence converges.
\end{proof}

\section{Entropy transfer to the nonlinear measure}
\label{sec:v4v1-entropy}

Let \(\mu_\Lambda\) be the normalized measure after multiplying
\(\nu_\Lambda\) by the remaining factors in the V100 ancestry ledger.
Whenever \(\mu\ll\nu\), write
\begin{equation}
 H(\mu\mid\nu)=
 \int\log\left(\frac{d\mu}{d\nu}\right)d\mu.
\label{eq:v4v1-relative-entropy}
\end{equation}

\begin{lemma}[Entropy bound for an event]
\label{lem:v4v1-entropy-event}
For every measurable event \(E\) with \(0<\nu(E)<1\),
\begin{equation}
 \mu(E)\le
 \frac{H(\mu\mid\nu)+\log2}{\log(1/\nu(E))}.
\label{eq:v4v1-entropy-event}
\end{equation}
\end{lemma}

\begin{proof}
Map the underlying space to the two-point partition \(\{E,E^c\}\).  Relative
entropy decreases under measurable maps, so
\begin{align}
 H(\mu\mid\nu)
 &\ge \mu(E)\log\frac{\mu(E)}{\nu(E)}
 +(1-\mu(E))\log\frac{1-\mu(E)}{1-\nu(E)}\nonumber\\
 &\ge \mu(E)\log\frac1{\nu(E)}-log2.
\end{align}
Rearrange to obtain \eqref{eq:v4v1-entropy-event}.
\end{proof}

\begin{theorem}[Nonlinear tightness from parent-relative entropy]
\label{thm:v4v1-entropy-tightness}
Suppose \(\{\nu_\Lambda\}\) is the Gaussian parent family of
Theorem~\ref{thm:v4v1-Gaussian-tightness} and
\begin{equation}
 \sup_\Lambda H(\mu_\Lambda\mid\nu_\Lambda)<\infty.
\label{eq:v4v1-uniform-entropy}
\end{equation}
Then \(\{\mu_\Lambda\}\) is tight in every
\(H^{-1-\varepsilon}\).  The same conclusion holds locally in infinite
volume if the entropy bound is uniform on every fixed compact region and
the parent has uniform local exponential Sobolev tails.
\end{theorem}

\begin{proof}
The Gaussian second-moment proof can be strengthened by the Gaussian
exponential-tail estimate: for every \(M>0\), choose a compact Sobolev ball
\(K_M\) such that
\begin{equation}
 \sup_\Lambda\nu_\Lambda(K_M^c)\le e^{-M}.
\label{eq:v4v1-parent-exp-tight}
\end{equation}
Apply Lemma~\ref{lem:v4v1-entropy-event} with \(E=K_M^c\).  If the entropy
bound is \(H_*\), then
\begin{equation}
 \sup_\Lambda\mu_\Lambda(K_M^c)
 \le\frac{H_*+\log2}{M}.
\end{equation}
Let \(M\to\infty\) to obtain tightness.  The local assertion applies the
same proof to restrictions and then uses a diagonal compactness argument.
\end{proof}

\begin{firewall}[Extensive entropy is not the required bound]
\label{fw:v4v1-extensive-entropy}
A relative entropy growing like the number of cutoff modes does not satisfy
\eqref{eq:v4v1-uniform-entropy} and gives no continuum tightness through this
theorem.  Vacuum-energy and quadratic determinant normalizations must be
removed before the entropy is measured.  In infinite volume one needs local
entropy or entropy density together with a separate thermodynamic-limit
argument.
\end{firewall}

\begin{corollary}[A concrete next bosonic interface]
\label{cor:v4v1-next-interface}
After the exact quadratic normalization, it is sufficient for
distributional tightness to prove a cutoff-uniform relative-entropy bound for
the remaining nonlinear gauge--ghost--Higgs--metric density with respect to
\(\nu_\Lambda\).  This criterion does not require maximum plaquette
admissibility and is compatible with fields supported only in negative
Sobolev spaces.
\end{corollary}

\begin{proof}
This is Theorem~\ref{thm:v4v1-entropy-tightness}.  The V100 ledger ensures
that the quadratic determinant normalization is not counted again inside
the relative density.
\end{proof}

\section{Reflection positivity of the Gaussian parent}
\label{sec:v4v1-RP}

Uniform positivity of \(K_\Lambda\) is not Osterwalder--Schrader reflection
positivity.  A sufficient translation-invariant condition is a positive
Kallen--Lehmann representation in the Euclidean time frequency:
\begin{equation}
 C_\Lambda(p_0,\mathbf p)=
 \int_0^\infty
 \frac{d\rho_{\Lambda,\mathbf p}(m^2)}
 {p_0^2+\mathbf p^2+m^2},
 \qquad d\rho_{\Lambda,\mathbf p}\ge0.
\label{eq:v4v1-KL}
\end{equation}

\begin{theorem}[Gaussian reflection-positive closure]
\label{thm:v4v1-Gaussian-RP}
If every transverse covariance has the representation
\eqref{eq:v4v1-KL}, then the Gaussian parent is reflection positive on
gauge-invariant positive-time linear observables.  If its cylinder
Schwinger functions converge, the limiting Gaussian parent is reflection
positive.
\end{theorem}

\begin{proof}
For a positive-time test function \(f\), Fourier inversion in \(p_0\) gives
\begin{equation}
 \int_{t,s>0}\overline{f(t)}
 C_\Lambda(-t-s,\mathbf p)f(s)\,dt,ds
 =\int_0^\infty\frac{d\rho(m^2)}{2\omega_m}
 \left|\int_0^\infty e^{-\omega_mt}f(t)\,dt\right|^2\ge0,
\end{equation}
where \(\omega_m=(\mathbf p^2+m^2)^{1/2}\).  Gaussian Wick expansion
extends this positivity from linear to polynomial observables.  The limit
statement is the reflection-positivity closure theorem proved in frozen
V100.
\end{proof}

\begin{firewall}[Ellipticity does not imply a positive spectral measure]
\label{fw:v4v1-RP-open}
The bounds \eqref{eq:v4v1-uniform-ellipticity} control the size of the
covariance but do not imply \eqref{eq:v4v1-KL}.  Logarithmic polarization,
gauge fixing, curved finite compression, and the conformal gravitational
sector must be assembled before the spectral measure can be shown positive.
No such proof is supplied here.
\end{firewall}

\section{Standard Model diagnosis and next acquisition}
\label{sec:v4v1-next}

For \(SU(3)\) and \(SU(2)\), perturbative asymptotic freedom is consistent
with the lower ellipticity bound in
\eqref{eq:v4v1-uniform-ellipticity}, but it does not prove it for the full
interacting cutoff measure.  For \(U(1)_Y\), the ordinary Standard Model
one-loop coefficient has the Landau sign.  Neither anomaly cancellation nor
the determinant-line trivialization changes this quadratic sign.  A
nontrivial ultraviolet completion, rigorous embedding, or complete quantum-
gravitational fixed point remains required before the parent measure can be
uniformly positive in that direction.

\begin{programme}[Fourth-volume V2 acquisition order]
\label{prog:v4v1-V2}
\begin{enumerate}
 \item Decompose the nonlinear cutoff density relative to the exact combined
 Gaussian parent, retaining the shell Schur ancestry of every term.
 \item Derive a finite-volume relative-entropy identity from the connected
 shell free energy.  Subtract vacuum and already-normalized quadratic terms
 before estimating it.
 \item Test the entropy bound first in the superrenormalizable Higgs--Yukawa
 and finite BRST quartet sectors, where the covariance is fixed.
 \item Isolate the nonabelian gauge interaction and determine whether its
 signed ghost-completed shell increment has a uniform entropy upper bound.
 \item Keep the hypercharge direction separate; do not infer its ultraviolet
 control from the nonabelian factors.
 \item If uniform entropy fails by an extensive shell term, identify the
 exact local counterterm or thermodynamic entropy density before attempting
 another norm estimate.
\end{enumerate}
\end{programme}

\section{V1 conclusion}
\label{sec:v4v1-conclusion}

The fourth volume starts from the measure object which survived the third-
volume falsification test.  Exact shell Schur complements show how the
signed bosonic--polarization quadratic form and its determinant normalization
flow together without duplicated factors.  Uniform ellipticity constructs a
tight Gaussian parent on the correct negative-Sobolev support.  A new
relative-entropy theorem reduces nonlinear distributional tightness to a
specific cutoff-uniform bound after vacuum and quadratic normalizations are
removed.  Reflection positivity still requires a positive spectral
representation, and the hypercharge ultraviolet direction remains open.

V2 will attack the connected-shell relative entropy of the nonlinear density
with respect to this parent measure.  The full SM--Einstein continuum remains
the objective and has not yet been proved.

\clearpage
\section{V2: connected-shell relative entropy after quadratic normalization}
\label{sec:v4v2}

Fourth-volume V1 constructs the signed Gaussian parent and reduces nonlinear
tightness to a relative-entropy bound.  This note derives the exact entropy
identity which must be estimated.  Interpolating the residual interaction
shows that its entropy is the integral of one connected variance.  A second
chain rule decomposes that quantity into conditional shell entropies without
copying an interaction occurrence across scales.  Summability follows if the
complete nondecaying local jet has already been placed in running parameters
and the conditional variance of the remaining shell increment gains one
positive power.  A Brascamp--Lieb criterion supplies a checkable sufficient
estimate on conditionally convex sectors.  Gauge, ghost, conformal, and
hypercharge directions which do not satisfy that convexity remain explicit.

\subsection{The interpolating free energy}
\label{sec:v4v2-free-energy}

Let \(\nu\) be one finite-dimensional combined Gaussian parent of V1.  Let
\(V\) be a real residual interaction for which
\begin{equation}
 Z(t)=\int e^{-tV}\,d\nu
\label{eq:v4v2-Zt}
\end{equation}
is finite and nonzero for \(0\le t\le1\).  Define
\begin{equation}
 d\mu_t=Z(t)^{-1}e^{-tV}d\nu,
 \qquad
 F(t)=-\log Z(t).
\label{eq:v4v2-mut-F}
\end{equation}
Constants, the full signed quadratic polarization, and its Gaussian
determinant normalization are absent from \(V\); they already belong to
\(\nu\).

\begin{theorem}[Exact connected-variance entropy identity]
\label{thm:v4v2-entropy-variance}
If \(V\in L^2(\mu_t)\) uniformly on compact subintervals of \([0,1]\), then
\begin{align}
 F'(t)&=\mathbb E_{\mu_t}V,
 \label{eq:v4v2-Fprime}\\
 F''(t)&=-\operatorname{Var}_{\mu_t}(V),
 \label{eq:v4v2-Fsecond}\\
 H(\mu_t\mid\nu)
 &=F(t)-tF'(t)
 =\int_0^t s\operatorname{Var}_{\mu_s}(V)\,ds.
 \label{eq:v4v2-entropy-identity}
\end{align}
In particular,
\begin{equation}
 H(\mu_1\mid\nu)
 =\int_0^1s\,langle V;V\rangle_{\mu_s}^{\rm conn}\,ds.
\label{eq:v4v2-one-connected-occurrence}
\end{equation}
\end{theorem}

\begin{proof}
Differentiate \(Z(t)\) under the integral.  Dominated truncation followed by
the assumed \(L^2\) bound justifies the derivative and gives
\begin{equation}
 F'(t)=\frac{\int Ve^{-tV}d\nu}{Z(t)}.
\end{equation}
Differentiating once more gives minus the connected second moment,
\eqref{eq:v4v2-Fsecond}.  Since
\begin{equation}
 \log\frac{d\mu_t}{d\nu}=-tV+F(t),
\end{equation}
integration against \(\mu_t\) gives
\(H=F-tF'\).  Its derivative is
\(-tF''(t)=t\operatorname{Var}_{\mu_t}(V)\), and it vanishes at \(t=0\).
Integration proves \eqref{eq:v4v2-entropy-identity}.
\end{proof}

The identity uses one complete residual interaction on both legs of a
connected covariance.  Expanding \(V\) into local monomials is allowed only
inside that covariance; estimating every monomial pair separately can
reintroduce volume and shell multiplicities which the connected expression
does not possess.

\begin{corollary}[A direct entropy criterion]
\label{cor:v4v2-direct-criterion}
If
\begin{equation}
 \sup_a\int_0^1s\operatorname{Var}_{\mu_{a,s}}(V_a)\,ds<\infty,
\label{eq:v4v2-direct-bound}
\end{equation}
then the nonlinear measures have uniformly bounded relative entropy with
respect to the V1 Gaussian parents and are tight in every
\(H^{-1-\varepsilon}\).
\end{corollary}

\begin{proof}
Apply Theorem~\ref{thm:v4v2-entropy-variance}, then
Theorem~\ref{thm:v4v1-entropy-tightness}.
\end{proof}

\subsection{Conditional entropy and the exact shell chain rule}
\label{sec:v4v2-chain}

Let \(X=(X_0,\ldots,X_N)\) be the exact orthogonal shell coordinates of the
combined parent after the Schur updates of
Theorem~\ref{thm:v4v1-Schur-shell}.  For \(j\le N\), write
\(\mathcal F_j=\sigma(X_0,ldots,X_j)\).  Let \(\mu^{(j)}\) and
\(\nu^{(j)}\) denote the restrictions to \(\mathcal F_j\), and let
\(\mu(dx_j\mid\mathcal F_{j-1})\) and
\(\nu(dx_j\mid\mathcal F_{j-1})\) be regular conditional laws.

\begin{theorem}[Relative-entropy shell chain rule]
\label{thm:v4v2-chain-rule}
If \(\mu\ll\nu\), then
\begin{equation}
 H(\mu\mid\nu)=H(\mu^{(0)}\mid\nu^{(0)})
 +\sum_{j=1}^N
 \mathbb E_\mu
 H\left(
 \mu(dx_j\mid\mathcal F_{j-1})
 \,middle|\,
 \nu(dx_j\mid\mathcal F_{j-1})
 \right).
\label{eq:v4v2-chain-rule}
\end{equation}
Every shell coordinate occurs in exactly one conditional entropy.
Cross-shell polarization is already present in the conditional Gaussian
means and Schur covariances and creates no extra summand.
\end{theorem}

\begin{proof}
For two blocks \((Y,Z)\), disintegrate both measures and factor their density:
\begin{equation}
 \frac{d\mu}{d\nu}(Y,Z)=
 \frac{d\mu_Y}{d\nu_Y}(Y)
 \frac{d\mu(dZ\mid Y)}{d\nu(dZ\mid Y)}(Z).
\label{eq:v4v2-density-chain}
\end{equation}
Take the logarithm and integrate against \(\mu\).  This gives the marginal
entropy plus the expected conditional entropy.  Induct over the finite shell
tower to obtain \eqref{eq:v4v2-chain-rule}.  The conditional law of a block
Gaussian is exactly the completed-square law in
\eqref{eq:v4v1-complete-square}, so cross-shell blocks modify that one
conditional factor rather than adding another occurrence.
\end{proof}

\begin{firewall}[A marginal interaction is not paid on every descendant shell]
\label{fw:v4v2-no-repeat}
An interaction assigned to the conditional density of shell \(j\) remains
inside its conditional entropy in \eqref{eq:v4v2-chain-rule}.  Reusing the
same unconditioned interaction in all later variances would replace one
density factor by \(N-j\) copies and create a false logarithmic divergence.
\end{firewall}

\subsection{Connected shell free energies}
\label{sec:v4v2-shell-free-energy}

Integrate shells from the ultraviolet downward.  Conditioned on the lower
field \(\phi_{<j}\), let \(\nu_j(\cdot\mid\phi_{<j})\) be the exact Gaussian
shell law produced by the current Schur complement, and let
\(V_j(\phi_j;\phi_{<j})\) be the residual interaction first assigned to
that shell.  Define
\begin{equation}
 Z_j(t;\phi_{<j})=
 \int e^{-tV_j(\phi_j;\phi_{<j})}
 d\nu_j(\phi_j\mid\phi_{<j}).
\label{eq:v4v2-shell-Z}
\end{equation}

\begin{proposition}[Conditional connected-variance identity]
\label{prop:v4v2-conditional-variance}
For almost every lower field,
\begin{equation}
 H(\mu_j(\cdot\mid\phi_{<j})mid
 \nu_j(\cdot\mid\phi_{<j}))
 =\int_0^1t\,operatorname{Var}_{\mu_{j,t}}
 (V_j\mid\phi_{<j})\,dt.
\label{eq:v4v2-conditional-variance}
\end{equation}
Consequently, if deterministic numbers \(h_j\) satisfy
\begin{equation}
 \mathbb E_\mu\int_0^1t\,operatorname{Var}_{\mu_{j,t}}
 (V_j\mid\mathcal F_{j-1})\,dt\le h_j,
\label{eq:v4v2-shell-hj}
\end{equation}
then
\begin{equation}
 H(\mu\mid\nu)\le H_0+\sum_{j=1}^Nh_j.
\label{eq:v4v2-entropy-sum}
\end{equation}
\end{proposition}

\begin{proof}
Apply Theorem~\ref{thm:v4v2-entropy-variance} to the conditional probability
space for each fixed lower field.  Measurability follows from the finite
dimensional conditional integral.  Take the lower-field expectation and
insert the result into Theorem~\ref{thm:v4v2-chain-rule}.
\end{proof}

\begin{theorem}[Positive-power entropy summability]
\label{thm:v4v2-positive-power}
Let the shell scale be \(a_j=q^{-j}a_0\).  Suppose the complete nondecaying
local jet of every \(V_j\) has been transported into the running quadratic
parent and the prescribed local couplings, and suppose the normalized
remainder obeys
\begin{equation}
 h_j\le Ca_j^\omega,
 \qquad \omega>0.
\label{eq:v4v2-positive-gain}
\end{equation}
Then the cutoff relative entropies are uniformly bounded and the nonlinear
measures are tight in \(H^{-1-\varepsilon}\).
\end{theorem}

\begin{proof}
The geometric series
\(\sum_ja_j^\omega=a_0^\omega/(1-q^{-\omega})\) converges.  Proposition
\ref{prop:v4v2-conditional-variance} gives a uniform entropy bound, and
Theorem~\ref{thm:v4v1-entropy-tightness} gives tightness.
\end{proof}

\begin{counterexample}[Marginal shell entropy accumulates]
\label{sep:v4v2-marginal}
Let every shell contain one independent standard Gaussian \(X_j\), and tilt
it by the same nonconstant bounded potential \(v(X_j)\).  If the one-shell
relative entropy is \(h_*>0\), then
\begin{equation}
 H(\mu^{(N)}\mid\nu^{(N)})=(N+1)h_*.
\label{eq:v4v2-marginal-growth}
\end{equation}
Thus a uniform order-one shell estimate does not imply continuum tightness
through relative entropy.
\end{counterexample}

\begin{proof}
Both measures are products.  Apply the chain rule; every conditional entropy
equals \(h_*\).  Their sum is \eqref{eq:v4v2-marginal-growth}.
\end{proof}

This is the entropy analogue of the V85 marginal determinant obstruction.
It also shows why vacuum energy, wave-function polarization, and marginal
couplings must be renormalized before the residual entropy is estimated.

\subsection{A conditionally convex sufficient estimate}
\label{sec:v4v2-convex}

Let a conditional shell Gaussian have inverse covariance \(K_j\), and use
the Cameron--Martin norm
\begin{equation}
 \|f\|_{K_j^{-1}}^2=
 \langle f,K_j^{-1}f\rangle.
\label{eq:v4v2-CM-norm}
\end{equation}

\begin{theorem}[Brascamp--Lieb shell entropy criterion]
\label{thm:v4v2-BL}
Suppose for every lower field and \(0\le t\le1\),
\begin{equation}
 K_j+t\nabla_j^2V_j\ge\rho K_j,
 \qquad \rho>0,
\label{eq:v4v2-uniform-convexity}
\end{equation}
as quadratic forms.  Then
\begin{equation}
 H(\mu_j\mid\nu_j)
 \le\frac1{2\rho}
 \sup_{0\le t\le1}
 \mathbb E_{\mu_{j,t}}
 \|\nabla_jV_j\|_{K_j^{-1}}^2.
\label{eq:v4v2-BL-entropy}
\end{equation}
If the right side is at most \(Ca_j^\omega\), \(\omega>0\), the entropy
contribution is summable.
\end{theorem}

\begin{proof}
The Brascamp--Lieb variance inequality for the measure with potential
\(\frac12\langle x,K_jx\rangle+tV_j(x)\) gives
\begin{equation}
 \operatorname{Var}_{\mu_{j,t}}(V_j)
 \le\mathbb E_{\mu_{j,t}}
 \left\langle\nabla_jV_j,
 (K_j+t\nabla_j^2V_j)^{-1}\nabla_jV_j\right\rangle.
\end{equation}
The form bound \eqref{eq:v4v2-uniform-convexity} makes this at most
\(\rho^{-1}\mathbb E\|\nabla_jV_j\|_{K_j^{-1}}^2\).  Insert it in
\eqref{eq:v4v2-conditional-variance} and use
\(\int_0^1t\,dt=1/2\).
\end{proof}

\begin{corollary}[Stable polynomial sectors]
\label{cor:v4v2-stable-polynomial}
On a finite shell, a positive Higgs quartic interaction and any bounded
Yukawa/Higgs remainder whose negative Hessian is strictly smaller than the
parent quadratic form satisfy \eqref{eq:v4v2-uniform-convexity}.  Their
entropy is therefore reduced to the Cameron--Martin gradient estimate in
\eqref{eq:v4v2-BL-entropy}.
\end{corollary}

\begin{proof}
A positive quartic has nonnegative radial Hessian outside its finite
mixed-component correction.  Bound the negative part of the full finite
matrix Hessian relative to \(K_j\).  If its relative norm is below
\(1-\rho\), then \eqref{eq:v4v2-uniform-convexity} follows.  Apply
Theorem~\ref{thm:v4v2-BL}.
\end{proof}

\begin{firewall}[Gauge and ghost sectors are not globally convex]
\label{fw:v4v2-not-convex}
Nonabelian gauge interactions in linear coordinates, the Faddeev--Popov
determinant across Gribov horizons, the gravitational conformal direction,
and a hypercharge kernel without ultraviolet positivity need not satisfy
\eqref{eq:v4v2-uniform-convexity}.  Theorem~\ref{thm:v4v2-BL} is a sufficient
criterion for controlled sectors, not a proof of the full SM--Einstein
entropy bound.
\end{firewall}

\subsection{Complex phases and determinant lines}
\label{sec:v4v2-phase}

Relative entropy applies to positive measures.  Write the globally
trivialized anomaly-free determinant section as
\begin{equation}
 W_a=|W_a|e^{i\vartheta_a}.
\label{eq:v4v2-phase-split}
\end{equation}
V87--V95 remove the gauge-anomalous and torsion holonomy, but a physical
CP/Yukawa phase can remain.

\begin{proposition}[Modulus tightness and phase convergence]
\label{prop:v4v2-modulus-phase}
If the positive measures proportional to \(|W_a|d\nu_a\) have uniformly
bounded parent-relative entropy, and if
\(e^{i\vartheta_a}\to e^{i\vartheta}\) in probability with a nonzero
limiting normalized expectation, then the complex normalized Schwinger
functionals converge on bounded cylinder observables along every convergent
subsequence of the positive measures.
\end{proposition}

\begin{proof}
Entropy tightness gives a convergent subsequence of the positive measures.
The phases have modulus one and are therefore uniformly integrable.
Convergence in probability implies \(L^1\) convergence along a Skorokhod
realization.  Apply the normalized-weight Vitali theorem from frozen V100 to
the phase weight.
\end{proof}

\begin{firewall}[A trivial anomaly line is not a positive density]
\label{fw:v4v2-line-not-positive}
Trivial determinant-line holonomy permits a global phase convention.  It
does not make the determinant pointwise nonnegative and does not prove
reflection positivity.  Entropy controls the modulus; the phase and the
cross-reflection cone require their own estimates.
\end{firewall}

\subsection{V2 conclusion and V3 target}
\label{sec:v4v2-conclusion}

The nonlinear tightness gate now has an exact connected formulation.
Relative entropy is the interpolation integral of one connected variance,
and its shell chain rule assigns every coordinate and cross-shell response
to one conditional law.  A positive power after complete local
normalization yields summability; an order-one shell contribution diverges
linearly.  Conditional Brascamp--Lieb convexity reduces controlled Higgs and
bounded Yukawa sectors to a Cameron--Martin gradient estimate, while the
gauge, ghost, conformal, and hypercharge sectors remain open.

V3 should compute the shell Cameron--Martin gradient for the normalized
Higgs--Yukawa sector in four dimensions, including its gauge-covariant
couplings, and determine whether its complete nondecaying jet leaves a
positive entropy power.  The calculation must retain the common quadratic
parent and must not use a fermion polarization term a second time.

\clearpage
\section{V3: the marginal Higgs shell falsifies uniform Gaussian-relative entropy}
\label{sec:v4v3}

V2 gives a sufficient nonlinear-tightness criterion in terms of entropy
relative to the signed Gaussian parent.  The required marginal test is
negative.  A Wick-ordered quartic interaction on one four-dimensional
Gaussian shell has connected variance proportional to the number of shell
cells.  Its Cameron--Martin gradient has the same scaling, so the
Brascamp--Lieb estimate cannot produce a positive cutoff power.  The same
dimensional obstruction occurs for a marginal Yukawa vertex.  Vacuum and
quadratic subtractions do not remove it.  Thus the Gaussian-relative entropy
criterion is useful for irrelevant residuals but is too strong for the full
marginal theory.  The correct upstream parent must include the complete
running marginal shell law; entropy is then applied only to its irrelevant
remainder.

\subsection{A single band-limited scalar shell}
\label{sec:v4v3-shell}

Work first with one real scalar on a flat unit four-torus.  Let \(h_K\) be a
centred Gaussian field with covariance
\begin{equation}
 C_K(x-y)=
 \sum_{K\le|p|<2K}\frac{\chi(p/K)}{p^2+m_K^2}
 e^{ip\cdot(x-y)},
\label{eq:v4v3-shell-covariance}
\end{equation}
where \(\chi\) is a fixed nonzero smooth band cutoff and
\(0\le m_K\le cK\).  Fourier normalization constants are absorbed into
\(\chi\).  Scaling gives
\begin{equation}
 C_K(x)=K^2c_K(Kx),
\label{eq:v4v3-covariance-scaling}
\end{equation}
where \(c_K\) stays in a bounded set of Schwartz functions and converges
along every convergent ratio \(m_K/K\).

Let Wick ordering refer to \(C_K(0)\), and define
\begin{equation}
 V_K(h)=\frac{\lambda_K}{4!}
 \int_{T^4}:h_K(x)^4:\,dx.
\label{eq:v4v3-quartic-shell}
\end{equation}

\begin{lemma}[Exact quartic connected variance]
\label{lem:v4v3-quartic-variance}
One has
\begin{equation}
 \operatorname{Var}_{\nu_K}(V_K)
 =\frac{\lambda_K^2}{4!}
 \int_{T^4}\int_{T^4}C_K(x-y)^4\,dx,dy.
\label{eq:v4v3-quartic-variance}
\end{equation}
There are constants \(0<c<C<\infty\), independent of large \(K\), such
that
\begin{equation}
 c\lambda_K^2K^4
 \le\operatorname{Var}_{\nu_K}(V_K)
 \le C\lambda_K^2K^4.
\label{eq:v4v3-volume-scaling}
\end{equation}
\end{lemma}

\begin{proof}
Wick's theorem gives
\begin{equation}
 \mathbb E[:h(x)^4::h(y)^4:]=4!\,C_K(x-y)^4.
\end{equation}
The Wick polynomial has mean zero, so inserting this identity into the
square of \eqref{eq:v4v3-quartic-shell} proves
\eqref{eq:v4v3-quartic-variance}.

Translation invariance removes one integral.  Using
\eqref{eq:v4v3-covariance-scaling} and changing variables \(z=K(x-y)\)
gives
\begin{equation}
 \int_{T^4}C_K(z)^4dz
 =K^4\int_{KT^4}c_K(w)^4dw.
\end{equation}
The upper bound follows from the uniform Schwartz bounds.  The cutoff is
nonzero, so a fixed compact set has a positive limiting fourth-power
integral; this gives the lower bound.
\end{proof}

The factor \(K^4\) is the number of correlation cells of side \(K^{-1}\)
in fixed physical volume.  It is not a vacuum expectation removed by Wick
ordering: it is the connected fluctuation of the interaction itself.

\subsection{Small-coupling entropy lower asymptotics}
\label{sec:v4v3-entropy}

Let
\begin{equation}
 d\mu_{K,t}=Z_K(t)^{-1}e^{-tV_K}d\nu_K.
\label{eq:v4v3-tilted-shell}
\end{equation}
For a positive quartic and bounded \(|\lambda_K|\), this is well defined for
\(0\le t\le1\).

\begin{theorem}[Marginal shell entropy]
\label{thm:v4v3-marginal-entropy}
There are \(\lambda_0>0\) and \(c_0>0\) such that, for
\(0<\lambda_K\le\lambda_0\),
\begin{equation}
 H(\mu_{K,1}\mid\nu_K)
 \ge c_0\lambda_K^2K^4
\label{eq:v4v3-entropy-lower}
\end{equation}
for all sufficiently large \(K\), provided
\(\lambda_KK^0\) remains in the perturbative interval.  In particular, a
nonzero marginal coupling cannot satisfy the uniform entropy hypothesis
\eqref{eq:v4v1-uniform-entropy} relative to the Gaussian shell parent.
\end{theorem}

\begin{proof}
The exact V2 identity gives
\begin{equation}
 H(\mu_{K,1}\mid\nu_K)
 =\int_0^1t\operatorname{Var}_{\mu_{K,t}}(V_K)dt.
\label{eq:v4v3-entropy-integral}
\end{equation}
At \(t=0\), Lemma~\ref{lem:v4v3-quartic-variance} gives the lower bound
\(c\lambda_K^2K^4\).  Rescale the torus into \(K^4\) unit correlation
cells.  The connected moments of the normalized Wick polynomial per cell
are bounded uniformly for \(|t\lambda_K|\le\lambda_0\), by Gaussian
hypercontractivity and the finite-range decay of \(c_K\).  Hence
\begin{equation}
 \left|operatorname{Var}_{\mu_{K,t}}(V_K)
 -\operatorname{Var}_{\nu_K}(V_K)\right|
 \le C t\lambda_K^3K^4.
\label{eq:v4v3-variance-perturbation}
\end{equation}
Choose \(\lambda_0\) so \(C\lambda_0<c/2\).  The variance in
\eqref{eq:v4v3-entropy-integral} is then at least
\((c/2)\lambda_K^2K^4\) for \(0\le t\le1\).  Integration of \(t\) proves
\eqref{eq:v4v3-entropy-lower}.
\end{proof}

\begin{remark}[Meaning of the perturbative interval]
The proof is a one-shell cluster estimate with constants uniform after
rescaling.  It does not assume that perturbation theory controls an
arbitrarily long RG trajectory.  It shows that even one sufficiently high
shell has extensive Gaussian-relative entropy when its dimensionless
quartic coupling is bounded away from zero.
\end{remark}

\subsection{The Cameron--Martin gradient has the same obstruction}
\label{sec:v4v3-gradient}

The Gaussian gradient of \eqref{eq:v4v3-quartic-shell} is
\begin{equation}
 \nabla V_K(x)=\frac{\lambda_K}{3!}:h_K(x)^3:.
\label{eq:v4v3-gradient}
\end{equation}

\begin{proposition}[Exact expected gradient energy]
\label{prop:v4v3-gradient-energy}
At the Gaussian point,
\begin{equation}
 \mathbb E_{\nu_K}
 \|\nabla V_K\|_{C_K}^2
 =\frac{\lambda_K^2}{3!}
 \int_{T^4}\int_{T^4}C_K(x-y)^4dxdy
 \asymp\lambda_K^2K^4.
\label{eq:v4v3-gradient-energy}
\end{equation}
Thus the V2 Brascamp--Lieb criterion gives an extensive upper bound and no
positive irrelevant power.
\end{proposition}

\begin{proof}
By definition of the Cameron--Martin metric,
\begin{equation}
 \|\nabla V_K\|_{C_K}^2
 =\int\int \nabla V_K(x)C_K(x-y)\nabla V_K(y)dxdy.
\end{equation}
Wick's theorem gives
\(\mathbb E[:h(x)^3::h(y)^3:]=3!C_K(x-y)^3\).  Insert this and
\eqref{eq:v4v3-gradient}; the result is
\eqref{eq:v4v3-gradient-energy}.  Lemma
\ref{lem:v4v3-quartic-variance} supplies the scaling.
\end{proof}

\begin{firewall}[Convexity is not a positive cutoff power]
\label{fw:v4v3-convexity}
The positive quartic makes the scalar density stable and can satisfy the V2
conditional-convexity hypothesis.  Its gradient energy nevertheless grows
like \(K^4\).  Stability and ultraviolet entropy summability are different
properties.
\end{firewall}

\subsection{Four Higgs components and gauge-covariant backgrounds}
\label{sec:v4v3-Higgs}

Let \(H\in\C^2\cong\R^4\) and
\begin{equation}
 V_{H,K}=\lambda_K\int:(H^*H)^2:dx.
\label{eq:v4v3-Higgs-quartic}
\end{equation}

\begin{corollary}[Standard Model Higgs marginality]
\label{cor:v4v3-Higgs-marginality}
For a nondegenerate four-component shell covariance, both
\(\operatorname{Var}(V_{H,K})\) and its expected Cameron--Martin gradient
energy are bounded above and below by positive constants times
\(\lambda_K^2K^4\).  A smooth gauge background and bounded Yukawa mass
matrix change only the constants at leading shell order.
\end{corollary}

\begin{proof}
Expand \((H^*H)^2\) into finitely many fourth-order real Wick monomials.
Their covariance matrix is positive, and at least one diagonal monomial has
the scalar lower bound of Lemma~\ref{lem:v4v3-quartic-variance}.  Finite
component count gives the upper bound.  A smooth connection and bounded
endomorphism are lower-order perturbations of the shell covariance; after
rescaling their symbol errors are \(O(K^{-1})\) and do not change the
leading positive constant.
\end{proof}

Relevant Higgs mass and vacuum terms must be included in the running parent.
Doing so changes the ratio \(m_K/K\) in
\eqref{eq:v4v3-shell-covariance} but does not remove the quartic lower bound
while that ratio stays bounded.

\subsection{The Yukawa shell has marginal entropy scale}
\label{sec:v4v3-Yukawa}

Let \(\psi_K,\bar\psi_K\) be a fermion shell and \(H_K\) a Higgs shell at
comparable momentum.  For one Yukawa matrix \(Y_K\), set
\begin{equation}
 V_{Y,K}=\int \bar\psi_KY_KH_K\psi_K\,dx+\text{Hermitian conjugate}.
\label{eq:v4v3-Yukawa-vertex}
\end{equation}
The Grassmann variables are integrated before positive entropy is applied,
so the statement below concerns the induced determinant modulus.

\begin{proposition}[Yukawa marginal shell count]
\label{prop:v4v3-Yukawa-scaling}
The connected second-order Yukawa contribution per fixed physical volume has
the scaling
\begin{equation}
 \langle V_{Y,K};V_{Y,K}\rangle^{\rm conn}
 =c_Y\operatorname{tr}(Y_K^*Y_K)K^4+o(K^4),
 \qquad c_Y>0,
\label{eq:v4v3-Yukawa-scaling}
\end{equation}
after the local vacuum and quadratic mass contractions are assigned to their
running parameters.  Hence a Yukawa coupling bounded away from zero also
fails the uniform Gaussian-relative entropy test.
\end{proposition}

\begin{proof}
Contract the two Yukawa vertices with two shell Dirac propagators and one
shell Higgs covariance.  Under \(x-y=K^{-1}z\), a four-dimensional Dirac
kernel has scale \(K^3\), the scalar kernel has scale \(K^2\), and the
relative-volume element has scale \(K^{-4}\).  The product therefore has
scale \(K^{3+3+2-4}=K^4\).  The spin and internal trace of the modulus
packet is a positive multiple of \(\operatorname{tr}(Y^*Y)\), nonzero for a
nonzero Yukawa matrix.  Wick/local contractions at one vertex are precisely
the already-assigned mass and vacuum terms; the separated connected
two-vertex integral retains the displayed leading coefficient.
\end{proof}

Anomaly cancellation does not cancel
\(\operatorname{tr}(Y^*Y)\), and Higgs stability does not supply an
irrelevant power.  The marginal gauge, Yukawa, and quartic trajectories must
therefore be treated jointly.

\subsection{The corrected non-Gaussian parent}
\label{sec:v4v3-corrected-parent}

Let \(\nu_j^{\rm marg}(d\phi_j\mid\phi_{<j})\) be a normalized conditional
shell law which contains:
\begin{enumerate}
 \item the exact Schur-updated signed quadratic form;
 \item the running Higgs quartic and Yukawa determinant packet;
 \item all other marginal gauge-invariant interactions at that scale; and
 \item the local counterterms fixed by the common normalization conditions.
\end{enumerate}
Let \(R_j\) denote only the remaining irrelevant interaction.

\begin{theorem}[Marginal-parent entropy criterion]
\label{thm:v4v3-marginal-parent}
Suppose the iterated marginal parent measures \(\nu_N^{\rm marg}\) are tight
on a common distribution space.  Let the full measures have conditional
density proportional to \(e^{-R_j}\) relative to the marginal shell laws.
If
\begin{equation}
 \sum_{j\ge0}
 \mathbb E_{\mu}
 H\left(
 \mu(d\phi_j\mid\mathcal F_{j-1})
 \,middle|\,
 \nu_j^{\rm marg}(d\phi_j\mid\mathcal F_{j-1})
 \right)<\infty,
\label{eq:v4v3-relative-marginal-entropy}
\end{equation}
then the full measures are tight.  A sufficient bound is
\(h_j\le Ca_j^\omega\), \(\omega>0\), for the connected interpolation
variance of \(R_j\).
\end{theorem}

\begin{proof}
Apply the relative-entropy shell chain rule of V2 with
\(\nu^{\rm marg}\), rather than the Gaussian, as the reference measure.
The summability hypothesis gives a uniform relative entropy.  The event
entropy lemma and tightness of the reference parents then transfer tightness
to the full measures.  The positive-power condition implies the summability
by the same geometric-series calculation as V2.
\end{proof}

\begin{firewall}[Tightness of the marginal parent is the fixed-point problem]
\label{fw:v4v3-fixed-point}
Theorem~\ref{thm:v4v3-marginal-parent} moves the quartic and Yukawa marginal
entropy into the reference law; it does not construct that law.  Its
tightness is the joint ultraviolet RG problem.  In four dimensions a
non-Gaussian Higgs--Yukawa fixed trajectory is not established here, and the
hypercharge Landau direction remains coupled to it.
\end{firewall}

\begin{firewall}[Uniform Gaussian-relative entropy is withdrawn as a full gate]
\label{fw:v4v3-withdraw-Gaussian-entropy}
The fourth-volume V1 theorem remains a correct sufficient criterion for
irrelevant tilts, but
Theorem~\ref{thm:v4v3-marginal-entropy} proves its uniform entropy hypothesis
fails for a nonzero marginal Higgs quartic relative to a purely Gaussian
parent.  It must not be used as the full interacting completion hypothesis.
The corrected criterion is Theorem~\ref{thm:v4v3-marginal-parent}.
\end{firewall}

\subsection{V3 conclusion and V4 target}
\label{sec:v4v3-conclusion}

The first nonlinear entropy test has done its job by failing sharply.  A
Wick-ordered four-dimensional Higgs quartic has variance and
Cameron--Martin gradient energy of order \(\lambda_K^2K^4\); the Yukawa
packet has the same cell count.  Neither Wick ordering, anomaly cancellation,
nor convexity produces a positive cutoff power.  Gaussian-relative entropy
can control only the irrelevant remainder after the complete marginal law
has been promoted into the parent measure.

V4 should formulate the exact one-step RG map for that marginal parent.  It
should use the shell Schur identity for the quadratic block, connected
cumulants for the gauge--Higgs--Yukawa interactions, and a finite projection
onto the running marginal basis.  The target is a contraction criterion for
the irrelevant complement and a precise fixed-point equation for the
nonabelian, hypercharge, Yukawa, quartic, and gravitational couplings.  No
perturbative beta function should be treated as a nonperturbative fixed
point.

\clearpage
\section{V4: exact marginal-parent shell map and the irrelevant contraction criterion}
\label{sec:v4v4}

V3 proves that a Gaussian parent cannot absorb the entropy of a nonzero
four-dimensional quartic or Yukawa shell.  The correct reference law must
contain the complete running marginal interaction.  This note defines that
law through the exact connected shell map.  A finite local projection gives
the running coupling vector; the complementary operator is the irrelevant
remainder.  First and second derivatives of the shell map are conditional
expectations and connected covariances, so the contraction hypothesis is
stated on the complete packet before norms.  A backward recursion theorem
then proves uniform decay of the irrelevant remainder along any prescribed
bounded marginal trajectory.  Existence of such a trajectory, especially in
the hypercharge and gravitational directions, remains the fixed-point gate.

\subsection{Local marginal coordinates and the normalized complement}
\label{sec:v4v4-coordinates}

At scale \(a_j=q^{-j}a_0\), let \(\mathfrak X_j\) be a Banach space of
gauge- and BRST-compatible local and quasilocal interactions with analytic
norm \(\|\cdot\|_j\).  Choose a finite-dimensional projection
\begin{equation}
 \mathfrak M_j:\mathfrak X_j\longrightarrow\mathfrak G_j
\label{eq:v4v4-marginal-projection}
\end{equation}
onto the complete prescribed basis of local operators of engineering
dimension at most four.  It includes vacuum energy, gauge and Higgs kinetic
forms, Higgs mass and quartic, Yukawa matrices, gauge fixing and BRST-exact
normalizations, the cosmological and Einstein terms, and every required
curvature-squared coupling.  Write
\begin{equation}
 \mathfrak I_j=1-\mathfrak M_j.
\label{eq:v4v4-irrelevant-projection}
\end{equation}
The coefficient vector \(g_j\) parametrizes
\(M_j(g_j)\in\mathfrak G_j\), while
\(R_j\in\Ran\mathfrak I_j\) is normalized by
\(\mathfrak M_jR_j=0\).

The finite list must contain all symmetry-allowed nondecaying operators.  If
one is omitted, V3 shows that it reappears as an extensive residual rather
than as an irrelevant error.

\subsection{The exact one-shell map}
\label{sec:v4v4-shell-map}

Let \(\zeta_j\) be the shell field with normalized conditional carrier
\(\nu_j(d\zeta_j\mid\phi)\).  Its quadratic block is the exact Schur
covariance of V1 and its gauge-fixed realization is understood only after
the BRST packet has been assembled.  It contains no second copy of the
interaction \(M_j(g_j)+R_j\).  For a full interaction
\(V\in\mathfrak X_j\), define
\begin{equation}
 \mathcal R_j(V)(\phi)=
 -\log\int
 \exp[-V(\phi+\zeta_j)]
 d\nu_j(\zeta_j\mid\phi).
\label{eq:v4v4-exact-RG}
\end{equation}
An additive constant is retained in the vacuum-energy coordinate of
\(\mathfrak G_{j-1}\), rather than discarded.

\begin{theorem}[Derivatives of the connected shell map]
\label{thm:v4v4-RG-derivatives}
On every ball on which the exponential moment in
\eqref{eq:v4v4-exact-RG} is analytic, let
\(\mu_{j,V}^{\phi}\) denote the tilted conditional shell law.  Then
\begin{align}
 D\mathcal R_j(V)[W]
 &=\mathbb E_{\mu_{j,V}^{\phi}}W(\phi+\zeta_j),
 \label{eq:v4v4-first-derivative}\\
 D^2\mathcal R_j(V)[W_1,W_2]
 &=-\langle W_1;W_2\rangle_{\mu_{j,V}^{\phi}}^{\rm conn}.
 \label{eq:v4v4-second-derivative}
\end{align}
Higher derivatives are the corresponding signed connected cumulants.  The
formulas include every shell Schur response and determinant normalization
exactly once.
\end{theorem}

\begin{proof}
Differentiate the logarithm of the normalized conditional integral.
The first derivative is the tilted expectation.  Differentiating that
expectation produces the negative covariance, proving
\eqref{eq:v4v4-second-derivative}.  Induction gives connected cumulants.
The conditional carrier was obtained by completing the quadratic square,
so its mean, covariance, and determinant normalization are already inside
the single expectation and are not additional derivative terms.
\end{proof}

The exact running and residual maps are
\begin{align}
 g_{j-1}
 &=\mathfrak M_{j-1}
 \mathcal S_j\mathcal R_j(M_j(g_j)+R_j),
 \label{eq:v4v4-beta-map}\\
 R_{j-1}
 &=\mathfrak I_{j-1}
 \mathcal S_j\mathcal R_j(M_j(g_j)+R_j),
 \label{eq:v4v4-remainder-map}
\end{align}
where \(\mathcal S_j\) rescales the remaining field and coordinates from
scale \(j\) to \(j-1\).  Equation \eqref{eq:v4v4-beta-map}, not a truncated
perturbative polynomial, is the fixed-point equation for the marginal
parent.

\begin{proposition}[Exact Taylor allocation]
\label{prop:v4v4-Taylor-allocation}
For \(V=M_j(g)+R\),
\begin{align}
 \mathcal S_j\mathcal R_j(V)
 &=\mathcal S_j\mathcal R_j(M_j(g))
 +\mathcal S_jD\mathcal R_j(M_j(g))[R]
 +\mathcal N_j(g,R),
 \label{eq:v4v4-Taylor}\
 \mathcal N_j(g,R)
 &=\int_0^1(1-t)
 \mathcal S_jD^2\mathcal R_j(M_j(g)+tR)[R,R]dt.
 \label{eq:v4v4-Taylor-remainder}
\end{align}
Projecting this single identity by \(\mathfrak M_{j-1}\) and
\(\mathfrak I_{j-1}\) assigns every connected occurrence to exactly one
running or irrelevant coordinate.
\end{proposition}

\begin{proof}
This is Taylor's formula with integral remainder in the Banach space
\(\mathfrak X_j\).  The two projections sum to the identity, so their
outputs are disjoint and exhaustive.
\end{proof}

\subsection{The irrelevant contraction theorem}
\label{sec:v4v4-contraction}

Fix a prescribed marginal trajectory \(g_j\in\mathcal K\), where
\(\mathcal K\) is compact.  Define
\begin{align}
 E_j(g)&=\mathfrak I_{j-1}
 \mathcal S_j\mathcal R_j(M_j(g)),
 \label{eq:v4v4-inhomogeneity}\\
 L_j(g)&=\mathfrak I_{j-1}\mathcal S_j
 D\mathcal R_j(M_j(g))\big|_{\Ran\mathfrak I_j}.
 \label{eq:v4v4-linear-irrelevant}
\end{align}

\begin{theorem}[Backward irrelevant stability]
\label{thm:v4v4-backward-stability}
Suppose uniformly for \(g\in\mathcal K\):
\begin{align}
 \|L_j(g)\|_{j\to j-1}&\le\kappa<1,
 \label{eq:v4v4-linear-contraction}\\
 \|E_j(g)\|_{j-1}&\le C_Ea_j^\omega,
 \qquad\omega>0,
 \label{eq:v4v4-source-gain}\\
 \|\mathfrak I_{j-1}\mathcal N_j(g,R)\|_{j-1}
 &\le C_N\|R\|_j^2
 \quad\text{for }\|R\|_j\le r_*.
 \label{eq:v4v4-quadratic-remainder}
\end{align}
For sufficiently small \(a_0\) and compatible constants, the terminal
problems \(R_N=0\) have solutions satisfying, uniformly in \(N\),
\begin{equation}
 \|R_j^{(N)}\|_j
 \le C\sum_{k=j+1}^{N}
 \kappa^{k-j-1}a_k^\omega
 \le C' a_{j+1}^\omega.
\label{eq:v4v4-remainder-bound}
\end{equation}
For each fixed \(j\), \(R_j^{(N)}\) converges as \(N\to\infty\), and the
limit is the unique small irrelevant trajectory subordinate to the given
\(g_j\).
\end{theorem}

\begin{proof}
Equation \eqref{eq:v4v4-remainder-map} and
Proposition~\ref{prop:v4v4-Taylor-allocation} give
\begin{equation}
 \|R_{j-1}\|_{j-1}
 \le C_Ea_j^\omega+\kappa\|R_j\|_j+C_N\|R_j\|_j^2.
\label{eq:v4v4-recursion-bound}
\end{equation}
Start from \(R_N=0\) and iterate toward smaller \(j\).  Choose the small
ball so \(C_Nr_*\le(1-\kappa)/2\).  The quadratic term can then be absorbed
into an effective contraction
\(\kappa'=(1+\kappa)/2<1\).  Induction gives the convolution bound in
\eqref{eq:v4v4-remainder-bound}; since
\(a_k^\omega=a_{j+1}^\omega q^{-\omega(k-j-1)}\), the final series is
geometric.

For two terminal cutoffs \(N'>N\), their difference at level \(j\) obeys the
same contraction, with a tail source beginning above \(N\).  Its norm is
bounded by the tail of the convergent series in
\eqref{eq:v4v4-remainder-bound}, so \(R_j^{(N)}\) is Cauchy.  The same
contraction proves uniqueness in the small ball.
\end{proof}

\begin{corollary}[Summable residual entropy]
\label{cor:v4v4-residual-entropy}
If the conditional connected variance of the normalized remainder is
bounded by
\begin{equation}
 h_j\le C\|R_j\|_j^2,
\label{eq:v4v4-entropy-norm}
\end{equation}
then the residual entropy relative to the marginal parent is uniformly
bounded and the V3 marginal-parent tightness theorem applies.
\end{corollary}

\begin{proof}
Theorem~\ref{thm:v4v4-backward-stability} gives
\(h_j\le C a_j^{2\omega}\).  Sum the geometric series and apply
Theorem~\ref{thm:v4v3-marginal-parent}.
\end{proof}

\subsection{What the contraction hypothesis means}
\label{sec:v4v4-meaning}

By Theorem~\ref{thm:v4v4-RG-derivatives}, the linear map in
\eqref{eq:v4v4-linear-irrelevant} is
\begin{equation}
 L_j(g)R=\mathfrak I_{j-1}\mathcal S_j
 \mathbb E_{\mu_{j,M(g)}^\phi}R(\phi+\zeta_j).
\label{eq:v4v4-L-expectation}
\end{equation}
Its norm must be taken after the complete conditional expectation and after
the marginal projection removes every nondecaying local jet.  Estimating
individual Wick contractions first can produce an order-one row and hide
the cancellation which yields \(q^{-\omega}\).

\begin{proposition}[Scaling certificate for the linear contraction]
\label{prop:v4v4-scaling-certificate}
Suppose every monomial in \(R\) has scaling dimension at least
\(4+\omega\), its conditional expectation preserves the Ward-completed
operator basis, and the analytic norm controls all displacement moments
through order \(\lceil\omega\rceil\).  Then its rescaled marginally projected
linear image satisfies
\begin{equation}
 \|L_j(g)R\|_{j-1}\le Cq^{-\omega}\|R\|_j.
\label{eq:v4v4-scaling-contraction}
\end{equation}
If \(Cq^{-\omega}<1\), this verifies
\eqref{eq:v4v4-linear-contraction}.
\end{proposition}

\begin{proof}
Rescaling a dimension-\(4+\omega\) integrated operator from cell size
\(a_j\) to \(a_{j-1}=qa_j\) contributes \(q^{-\omega}\).  Conditional
expectation contracts only shell fields and cannot lower the dimension
except through local Taylor jets.  Those jets lie in
\(\Ran\mathfrak M_{j-1}\) and are removed before the norm.  The displacement
moments bound the Taylor remainder with the same factor.  The finite operator
basis and compact coupling set give the constant \(C\).
\end{proof}

\begin{firewall}[Power counting is conditional on the complete local projection]
\label{fw:v4v4-power-counting}
If a symmetry-allowed marginal term is absent from \(\mathfrak G_j\), the
conditional expectation can generate it from an apparently irrelevant
operator.  Then \(\mathfrak I_{j-1}\) retains an order-one component and
\eqref{eq:v4v4-scaling-contraction} is false.  The projection must include
the full gauge--Higgs--Yukawa--gravity basis before power counting.
\end{firewall}

\subsection{The marginal fixed-trajectory equation}
\label{sec:v4v4-fixed-trajectory}

The irrelevant theorem is subordinate to the exact finite-dimensional
recursion
\begin{equation}
 g_{j-1}=\mathcal B_j(g_j,R_j),
 \qquad
 \mathcal B_j:=\mathfrak M_{j-1}\mathcal S_j
 \mathcal R_j(M_j(g_j)+R_j).
\label{eq:v4v4-exact-beta}
\end{equation}

\begin{theorem}[Conditional assembly of the marginal parent]
\label{thm:v4v4-parent-assembly}
Suppose there is a trajectory \(g_j\in\mathcal K\) solving
\eqref{eq:v4v4-exact-beta} for all \(j\), the corresponding normalized
marginal shell laws are positive and tight, and the irrelevant estimates of
Theorem~\ref{thm:v4v4-backward-stability} hold.  Then the iterated marginal
parents exist as a tight projective family, their residual tilts have
uniformly bounded relative entropy, and every subsequential full measure has
the prescribed running couplings.
\end{theorem}

\begin{proof}
The exact shell normalization theorem of V1 makes the finite parent laws
projectively consistent along the solved coupling trajectory.  Their
assumed tightness gives a projective subsequential limit.  The irrelevant
trajectory and Corollary~\ref{cor:v4v4-residual-entropy} give a uniformly
bounded relative entropy, so V3 transfers tightness to the full measures.
The marginal projection of every one-step effective interaction is exactly
\eqref{eq:v4v4-exact-beta}; hence the limiting cylinder interactions retain
the prescribed coupling coordinates.
\end{proof}

\begin{firewall}[The exact beta map has not been solved]
\label{fw:v4v4-beta-open}
No theorem here constructs a global solution of
\eqref{eq:v4v4-exact-beta}.  Perturbative coefficients diagnose its tangent
directions but do not prove a nonperturbative trajectory, positivity, or
tightness.  The hypercharge Landau sign, Higgs--Yukawa marginality,
nonabelian gauge measure, and gravitational conformal sector remain coupled
fixed-point gates.
\end{firewall}

\begin{firewall}[Ghost completion precedes the Banach norm]
\label{fw:v4v4-ghost}
The marginal projection and connected cumulants must be formed in the
BRST-complete gauge--ghost packet.  A positive probability norm is applied
only after the quartet cancellations and determinant phases have been
assigned.  Treating a Grassmann ghost factor as an independent positive
measure would invalidate both \eqref{eq:v4v4-exact-RG} and the Ward
projection.
\end{firewall}

\subsection{V4 conclusion and V5 audit target}
\label{sec:v4v4-conclusion}

The corrected marginal parent now has an exact one-step definition.  Its
derivatives are conditional connected correlators, its running couplings are
the finite local projection, and its irrelevant remainder is the
complementary projection of the same shell occurrence.  Along any bounded
solution of the exact marginal recursion, a uniform irrelevant contraction
and an \(O(a_j^\omega)\) source give a unique summable remainder and bounded
residual entropy.  The fixed marginal trajectory itself remains open.

V5 is the next compulsory NS/YM audit.  It should then select one marginal
direction whose exact shell map can be controlled without assuming the full
fixed point.  The likely candidates are the BRST-completed nonabelian
quadratic/marginal block or the positive Higgs radial sector.  Hypercharge
must remain separate unless a genuinely new ultraviolet mechanism is
proved.


\clearpage
\section{V5: companion audit and the exact nonabelian quadratic shell coordinate}
\label{sec:v4v5}

V4 reduced the measure problem to two distinct tasks: solve the finite
marginal recursion and contract its irrelevant complement.  The latter has
a complete conditional theorem; the former is still open.  This compulsory
fifth-version audit first records what can be imported from the current
Yang--Mills and Navier--Stokes programmes.  It then applies those lessons to
one part of the marginal recursion which admits an exact treatment: the
background-gauge quadratic coordinate of each compact nonabelian factor.

The result is an exact finite-cutoff Ward identity, an exact connected
formula for the shell polarization, and a compactification of the
asymptotically free coordinate.  These facts remove a false boundedness
obstruction in V4.  They do not establish the sign or uniform size of the
nonabelian beta increment.

\subsection{Compulsory companion audit}
\label{sec:v4v5-audit}

The files read for this audit were
\begin{enumerate}
\item \texttt{Renewal\_Yang\_Mills\_Mass\_Gap\_Research\_Notes\_V24.tex},
\(459461\) bytes, with SHA--256
\begin{center}\scriptsize\ttfamily
7ACF32A855B728FB2C3B57C257B38FCDDBBBA88844FE3455A6F056600D617D5F.
\end{center}
\item \texttt{Renewal\_NS\_Stability\_Research\_Notes\_V31.tex},
\(887537\) bytes, with SHA--256
\begin{center}\scriptsize\ttfamily
F1121B62238AD5491BB7CF03635EA9934942EDF573ED8FAAA9EE79E1D38A6696.
\end{center}
\end{enumerate}

The Yang--Mills V24 note proves its clean one-occurrence triple-response
bound by estimating the complete multiplier once and attaching the scalar
response product afterwards.  Its seven signature atoms and five equality
strata are exact partitions of one coordinate bank; they are not additional
physical occurrences.  The lesson for the present RG map is that several
background derivatives of one logarithmic shell integral must first remain
inside their complete connected cumulant.  A later tensor or momentum
partition may classify that cumulant, but cannot turn it into several
independent entropy or covariance payers.

The Navier--Stokes V31 note proves that a normalized dyadic mark does not
control the critical first moment: removing the mark restores the exact
scale factor.  Heat lifting merely moves that factor into the critical
source norm unless the signed nonlinearity is retained before taking norms.
The corresponding SM lesson is that an order-one gauge polarization on each
logarithmic shell is not an irrelevant error.  Summing it produces the
running inverse coupling.  Probability normalization cannot make this sum
convergent; the complete marginal coordinate must carry it.

\begin{proposition}[Audit transfer rule]
\label{prop:v4v5-audit-transfer}
In the V4 connected shell map, the following allocation is forced.
\begin{enumerate}
\item Every mixed derivative belonging to one shell integral is assembled
as one connected cumulant before any absolute estimate.
\item Local momentum and internal-index partitions may be applied after
that assembly, provided they form a disjoint identity decomposition.
\item The curvature-quadratic projection is retained in the running
coupling.  It is not charged to the summable irrelevant entropy.
\end{enumerate}
\end{proposition}

\begin{proof}
For a normalized shell law and observables \(X_1,\ldots,X_m\), put
\[
 L(t)=-\log\mathbb E
 \exp\!\left[-\sum_{r=1}^m t_rX_r\right].
\]
The mixed derivative of \(L\) is the signed joint cumulant, obtained by the
single moment--cumulant partition identity.  Thus there is one parent
expectation, even when several insertion labels coincide.  A disjoint
projection of the resulting tensor sums to the original tensor and creates
no second expectation.  Finally, a four-dimensional curvature square has
engineering dimension four.  Its rescaled coefficient has no positive
power of the lattice spacing, so the scaling certificate of
Proposition~\ref{prop:v4v4-scaling-certificate} cannot place it in the
irrelevant complement.  This proves all three assertions.
\end{proof}

\subsection{An equivariant compact-lattice block transform}
\label{sec:v4v5-equivariant-block}

Let \(G_r\) be either \(SU(3)\) or \(SU(2)\).  On two finite periodic
lattices, let \(U\in G_r^{E_j}\) and \(V\in G_r^{E_{j-1}}\) denote fine and
coarse link fields.  Write \(dU\) for product Haar measure.  Suppose
\(S_j(U)\) is fine-lattice gauge invariant and the nonnegative block kernel
\(K_j(V,U)\) satisfies
\begin{equation}
 K_j(V^h,U^{\iota_jh})=K_j(V,U)
\label{eq:v4v5-kernel-equivariance}
\end{equation}
for every coarse gauge transformation \(h\), where \(\iota_jh\) is its
fixed fine-lattice lift.  Define
\begin{align}
 W_{j-1}(V)&=\int K_j(V,U)e^{-S_j(U)}\,dU,
\label{eq:v4v5-effective-weight}\\
 S_{j-1}(V)&=-\log W_{j-1}(V)
\label{eq:v4v5-effective-action}
\end{align}
where the weight is positive.

\begin{theorem}[Exact preservation of background gauge invariance]
\label{thm:v4v5-gauge-invariance}
Under \eqref{eq:v4v5-kernel-equivariance}, the effective action is exactly
coarse-lattice gauge invariant:
\begin{equation}
 S_{j-1}(V^h)=S_{j-1}(V).
\label{eq:v4v5-effective-invariance}
\end{equation}
No expansion in the gauge coupling is used.
\end{theorem}

\begin{proof}
Use \eqref{eq:v4v5-kernel-equivariance} and make the change of variables
\(U'=U^{\iota_jh^{-1}}\) in \eqref{eq:v4v5-effective-weight}.  Product Haar
measure is invariant and \(S_j(U')=S_j(U)\).  Hence
\(W_{j-1}(V^h)=W_{j-1}(V)\).  Taking the negative logarithm proves the
claim.
\end{proof}

This finite-dimensional construction supplies the positive compact gauge
carrier required by V4 before gauge fixing.  Ghosts may be introduced to
compute local coordinates, but the invariant integral is the defining
object and prevents a ghost determinant from being counted as an
independent probability law.

\subsection{The exact quadratic Ward identity}
\label{sec:v4v5-ward}

Use the exponential chart
\[
 V_e(A)=\exp(a_{j-1}A_e),
 \qquad A_e\in\mathfrak g_r,
\]
near the trivial coarse field, and put
\(\Gamma_{j-1}(A)=S_{j-1}(V(A))\).  Let
\(d_{j-1}\alpha\) be the infinitesimal lattice gradient of a site field
\(\alpha\).  Assume \(\Gamma_{j-1}\) is twice differentiable at zero and
stationary there.

\begin{theorem}[Finite-cutoff Hessian Ward identity]
\label{thm:v4v5-hessian-ward}
The Hessian \(H_{j-1}=D^2\Gamma_{j-1}(0)\) satisfies
\begin{equation}
 H_{j-1}(d_{j-1}\alpha,B)=0
\label{eq:v4v5-hessian-ward}
\end{equation}
for every site field \(\alpha\) and link field \(B\).
If the action is translation invariant, its finite-volume Fourier kernel
obeys
\begin{equation}
 \sum_\mu
 \overline{\widehat d_\mu(p)}
 H_{\mu\nu}^{ab}(p)=0,
 \qquad
 \widehat d_\mu(p)=a_{j-1}^{-1}(e^{ia_{j-1}p_\mu}-1).
\label{eq:v4v5-fourier-ward}
\end{equation}
If, along an increasing-volume family, these kernels have a common
quasilocal continuous symbol near \(p=0\) for which
\eqref{eq:v4v5-fourier-ward} holds at momenta accumulating at zero, that
symbol has no local gauge-boson mass term:
\begin{equation}
 H_{\mu\nu}^{ab}(0)=0.
\label{eq:v4v5-no-mass}
\end{equation}
\end{theorem}

\begin{proof}
Theorem~\ref{thm:v4v5-gauge-invariance} gives the infinitesimal identity
\[
 D\Gamma_{j-1}(A)[X_\alpha(A)]=0,
\]
where \(X_\alpha(0)=d_{j-1}\alpha\).  Differentiate in the direction
\(B\) at \(A=0\).  The product rule gives
\[
 D^2\Gamma_{j-1}(0)[B,d_{j-1}\alpha]
 +D\Gamma_{j-1}(0)[DX_\alpha(0)B]=0.
\]
The second term vanishes by stationarity, proving
\eqref{eq:v4v5-hessian-ward}.  Fourier transformation converts the lattice
gradient into \(\widehat d_\mu(p)\), proving
\eqref{eq:v4v5-fourier-ward}.

For the asserted common symbol, write
\(M_{\mu\nu}^{ab}=H_{\mu\nu}^{ab}(0)\).  Choose allowed momenta tending to
zero along the \(\rho\)-axis, divide
\eqref{eq:v4v5-fourier-ward} by their magnitude, and pass to the limit.
Since \(t^{-1}\widehat d_\mu(te_\rho)\to i\delta_{\mu\rho}\), one gets
\(M_{\rho\nu}^{ab}=0\).  This holds for every \(\rho\), proving
\eqref{eq:v4v5-no-mass}.  The additional limit hypothesis is essential:
finite-volume gauge invariance alone does not exclude nonlocal holonomy
dependence of the constant mode.
\end{proof}

If the kernel is twice differentiable in momentum and has the usual
hypercubic, parity, and global adjoint symmetries, its leading quadratic term
on a simple factor is
\begin{equation}
 H_{\mu\nu}^{ab}(p)
 =z_{r,j-1}\delta^{ab}
 (|p|^2\delta_{\mu\nu}-p_\mu p_\nu)+O(|p|^4).
\label{eq:v4v5-kinetic-coordinate}
\end{equation}
Thus the Ward-completed Hessian defines a single inverse-coupling coordinate
\(z_{r,j-1}=g_{r,j-1}^{-2}\) for each simple factor.  The order-
\(|p|^0\) coordinate is absent exactly, rather than tuned shell by shell.

\subsection{One complete shell-polarization occurrence}
\label{sec:v4v5-polarization}

For the same block integral, suppose locally that
\begin{equation}
 K_j(V(A),U)e^{-S_j(U)}=e^{-F_j(A,U)}
\label{eq:v4v5-F-definition}
\end{equation}
with an \(A\)-independent reference measure and sufficient differentiable
domination.  Let \(\mathbb E_A\) denote expectation under the normalized
density proportional to \(e^{-F_j(A,U)}\).

\begin{theorem}[Exact connected shell polarization]
\label{thm:v4v5-exact-polarization}
For link variations \(B,C\),
\begin{align}
 D\Gamma_{j-1}(A)[B]
 &=\mathbb E_A D_AF_j(A,U)[B],
\label{eq:v4v5-first-response}\\
 D^2\Gamma_{j-1}(A)[B,C]
 &=\mathbb E_A D_A^2F_j(A,U)[B,C]
 -\left\langle D_AF_j[B];D_AF_j[C]\right\rangle_A^{\rm conn}.
\label{eq:v4v5-second-response}
\end{align}
Both terms on the right of \eqref{eq:v4v5-second-response} belong to one
logarithmic shell occurrence.
\end{theorem}

\begin{proof}
Differentiate the negative logarithm of the integral in
\eqref{eq:v4v5-effective-weight}.  The first derivative is the normalized
expectation of \(D_AF_j\).  Differentiating this expectation gives the
expectation of \(D_A^2F_j\) and minus the connected covariance of the two
first derivatives.  These are the two terms generated by one application of
the product rule to one normalized integral, so no separate occurrence is
present.
\end{proof}

Let \(\mathfrak P_r\) denote the normalized projection of a transverse
Hessian onto the tensor in \eqref{eq:v4v5-kinetic-coordinate}.  The exact
one-shell relation is
\begin{align}
 z_{r,j-1}
 &=\mathfrak P_r\!\left(
 \mathbb E_0D_A^2F_j
 -\langle D_AF_j;D_AF_j\rangle_0^{\rm conn}
 \right),
\label{eq:v4v5-z-output}\\
 \beta^{\rm ex}_{r,j}
 &:=z_{r,j}-z_{r,j-1}.
\label{eq:v4v5-exact-beta-increment}
\end{align}
Equations \eqref{eq:v4v5-z-output}--%
\eqref{eq:v4v5-exact-beta-increment} define the exact nonabelian beta
increment without a perturbative truncation.

\begin{firewall}[The two response terms have no separate signs]
\label{fw:v4v5-response-sign}
The covariance in \eqref{eq:v4v5-second-response} is nonnegative only when
the two variations coincide as real observables.  It is subtracted from the
second-derivative term.  Neither positivity of the compact measure nor the
Ward identity determines the sign of their projected difference.  Bounding
the two terms independently and assigning each a new shell payer would also
violate Proposition~\ref{prop:v4v5-audit-transfer}.
\end{firewall}

\subsection{Compactifying the asymptotically free coordinate}
\label{sec:v4v5-compactification}

V4 placed its marginal trajectory in a compact coordinate set.  The inverse
nonabelian coupling \(z_{r,j}\) should instead diverge along an
asymptotically free ultraviolet trajectory.  Put
\begin{equation}
 u_{r,j}=z_{r,j}^{-1}=g_{r,j}^2.
\label{eq:v4v5-compact-coordinate}
\end{equation}
The change from \(z\) to \(u\) is not by itself sufficient: the coefficient
of \(F^2\) would still diverge in an unrescaled gauge field.  One must also
use the canonically normalized local field.

\begin{lemma}[Canonical Gaussian-boundary chart]
\label{lem:v4v5-canonical-chart}
In a local continuum chart, set \(A=gB\).  Then
\begin{equation}
 \frac1{4g^2}\int
 \left|dA+[A,A]\right|^2
 =\frac14\int
 \left|dB+g[B,B]\right|^2.
\label{eq:v4v5-canonical-action}
\end{equation}
Equivalently, use \(U_e=\exp(a gB_e)\) in the lattice identity chart.
The quadratic form in \(B\) is independent of \(g\), while the cubic and
quartic vertices carry \(g\) and \(g^2\).  Hence \(g=0\) is a regular
Gaussian boundary for local perturbative coordinates.
\end{lemma}

\begin{proof}
Substitute \(A=gB\) into the curvature:
\[
 dA+[A,A]=g\bigl(dB+g[B,B]\bigr).
\]
The factor \(g^2\) cancels the inverse coupling in front of the action,
giving \eqref{eq:v4v5-canonical-action}.  The lattice statement is the same
change of variables in the exponential chart.  Expansion at the identity
then gives a coupling-independent quadratic term and vertices with the
displayed powers of \(g\).
\end{proof}

\begin{theorem}[Exact telescope under a signed shell bound]
\label{thm:v4v5-beta-telescope}
Suppose for one simple factor that the exact trajectory exists and
\begin{equation}
 0<b_-\le\beta^{\rm ex}_{r,j}\le b_+<\infty
\label{eq:v4v5-beta-bounds}
\end{equation}
for every ultraviolet shell \(j\ge1\).  Then
\begin{align}
 z_{r,j}&=z_{r,0}+\sum_{k=1}^j\beta^{\rm ex}_{r,k},
\label{eq:v4v5-z-telescope}\\
 \frac1{z_{r,0}+b_+j}
 &\le u_{r,j}\le
 \frac1{z_{r,0}+b_-j},
\label{eq:v4v5-u-bounds}\\
 u_{r,j}
 &=\frac{u_{r,j-1}}
 {1+\beta^{\rm ex}_{r,j}u_{r,j-1}}.
\label{eq:v4v5-u-recursion}
\end{align}
In particular \(u_{r,j}\to0\), and the trajectory is compact in the
\(u\)-coordinate.
\end{theorem}

\begin{proof}
The definition \eqref{eq:v4v5-exact-beta-increment} telescopes to
\eqref{eq:v4v5-z-telescope}.  Insert
\eqref{eq:v4v5-beta-bounds} and take reciprocals of positive numbers to get
\eqref{eq:v4v5-u-bounds}.  Finally,
\(z_{r,j}=z_{r,j-1}+\beta^{\rm ex}_{r,j}\); inversion gives
\eqref{eq:v4v5-u-recursion}.
\end{proof}

\begin{corollary}[Compatibility with the V4 irrelevant theorem]
\label{cor:v4v5-v4-compatible}
Rewrite the local gauge interactions in the canonical fields of
Lemma~\ref{lem:v4v5-canonical-chart}.  Assume the V4 analytic norms, local
projections, large-field regulators, source gain, and irrelevant contraction
are uniform for \(g_r\in[0,g_*]\), including the Gaussian boundary.  If
\eqref{eq:v4v5-beta-bounds} holds for both \(SU(3)\) and \(SU(2)\), then
their unbounded inverse couplings do not violate the compact marginal-set
hypothesis of Theorem~\ref{thm:v4v4-backward-stability}.  In the coordinates
\((g_3,g_2)\), the unique irrelevant trajectory and its summable residual
entropy extend to the ultraviolet endpoint.
\end{corollary}

\begin{proof}
Theorem~\ref{thm:v4v5-beta-telescope} keeps each
\(g_r=\sqrt{u_r}\) in a compact interval and sends it to zero.  By
Lemma~\ref{lem:v4v5-canonical-chart}, the divergent coefficient \(z_r\) has
been removed from the quadratic reference form and the local vertices are
regular at that endpoint.  The assumed V4 constants are uniform in this
canonical chart, so Theorem~\ref{thm:v4v4-backward-stability} and
Corollary~\ref{cor:v4v4-residual-entropy} apply without using a bound on
\(z_r\).
\end{proof}

\begin{firewall}[The signed shell bound is the new nonabelian gate]
\label{fw:v4v5-beta-open}
No proof of \eqref{eq:v4v5-beta-bounds} has been given.  The one-loop
asymptotic-freedom coefficients predict it near \(u=0\), but an asymptotic
series does not provide the uniform connected-cumulant remainder estimate
needed here.  Establishing that estimate while preserving the complete
BRST/Ward packet is the next nonabelian task.
\end{firewall}

\begin{firewall}[Hypercharge remains a separate ultraviolet problem]
\label{fw:v4v5-hypercharge}
The compact \(U(1)_Y\) block transform satisfies the same formal Ward and
connected-response identities, but the SM matter polarization has the
opposite perturbative ultraviolet sign.  Theorem~\ref{thm:v4v5-beta-telescope}
therefore cannot be applied to hypercharge by analogy.  A UV completion,
an embedding, or a genuinely nonperturbative alternative remains necessary.
\end{firewall}

\subsection{V5 result ledger and V6 target}
\label{sec:v4v5-conclusion}

Relative to V1--V4 of this volume, V5 proves:
\begin{enumerate}
\item the compulsory companion audit with exact source fingerprints;
\item an occurrence-preserving transfer rule for connected shell responses;
\item exact preservation of coarse gauge invariance by an equivariant
compact-lattice block transform;
\item the finite-cutoff Hessian Ward identity and, under a common quasilocal
infinite-volume symbol, exact exclusion of a local quadratic gauge-boson
mass term;
\item the complete expectation-minus-connected-covariance formula for the
shell polarization;
\item an exact definition of the nonabelian beta increment; and
\item compactification of an asymptotically free inverse-coupling trajectory
by the canonical field and the coordinate \(g=\sqrt u\), conditional on a
uniform signed shell bound.
\end{enumerate}

The V4 compactness hypothesis is therefore compatible with asymptotic
freedom after canonical field normalization and coupling compactification.  The first remaining
nonabelian gate is no longer the Ward identity or a spurious mass tuning; it
is a uniform remainder theorem which turns the perturbative sign of
\(\beta^{\rm ex}_{r,j}\) into the exact bound
\eqref{eq:v4v5-beta-bounds}.  V6 should formulate that remainder theorem on
one finite shell, retain the diamagnetic and covariance pieces as one Ward
packet, and identify the weakest cluster or spectral hypotheses under which
the one-loop coefficient dominates the exact remainder.  The hypercharge,
Higgs--Yukawa, and gravitational marginal directions remain open and
separate.


\clearpage
\section{V6: one-shell dominance, small/large fields, and the hypercharge obstruction}
\label{sec:v4v6}

V5 identified the exact nonabelian beta increment as the transverse
projection of one complete expectation-minus-covariance shell packet.  Its
perturbative sign becomes a theorem only if the exact remainder is uniform
in shell number and volume.  This note states and proves a sufficient
small-field/large-field criterion for that conclusion.  The small-field
piece is controlled by a Ward-complete polymer expansion; the large-field
piece is controlled by an exponentially rare-event estimate before the
connected packet is split.

The same criterion has a sharp consequence for \(U(1)_Y\).  In the Standard
Model field content its Gaussian-boundary coefficient has the screening
sign.  Therefore an exact trajectory which remains in the small-coupling
wedge cannot extend through infinitely many ultraviolet shells with positive
inverse hypercharge coupling.  This is a conditional no-go theorem for the
unmodified hypercharge direction, not a construction of a replacement UV
sector.

\subsection{Flat-shell normalization and the one-loop constants}
\label{sec:v4v6-one-loop}

Work first at a fixed flat metric and omit dynamical graviton shells.  Use
the canonical gauge field of Lemma~\ref{lem:v4v5-canonical-chart}.  For a
simple factor \(G_r\), normalize the beta function by
\begin{equation}
 \mu\frac{dg_r}{d\mu}
 =-\frac{b_r}{16\pi^2}g_r^3+O(g_r^5).
\label{eq:v4v6-beta-convention}
\end{equation}
One logarithmic shell of ratio \(q>1\) then has inverse-coupling increment
\begin{equation}
 c_r=\frac{b_r\log q}{8\pi^2}.
\label{eq:v4v6-shell-constant}
\end{equation}

\begin{proposition}[Ward-complete one-loop coefficient]
\label{prop:v4v6-one-loop}
For left-handed Weyl fermions and complex scalars,
\begin{equation}
 b_r=\frac{11}{3}C_A(G_r)
 -\frac23\sum_f T_r(R_f)
 -\frac13\sum_s T_r(R_s).
\label{eq:v4v6-general-b}
\end{equation}
For three Standard Model generations and one complex Higgs doublet,
\begin{equation}
 b_3=7,
 \qquad b_2=\frac{19}{6},
 \qquad b_Y=-\frac{41}{6},
\label{eq:v4v6-SM-b}
\end{equation}
where the hypercharge convention is \(Q=T_3+Y\) and the sign convention is
\eqref{eq:v4v6-beta-convention}.
\end{proposition}

\begin{proof}
The background-field Hessian combines the gauge and ghost contractions into
\(11C_A/3\).  A Weyl loop contributes \(-2T(R)/3\), and a complex scalar
loop, including its seagull term, contributes \(-T(R)/3\).  These constants
are the transverse coefficients after the longitudinal pieces cancel in the
Ward packet.  The only radial integral is
\[
 \int_{K/q<|p|<K}\frac{d^4p}{(2\pi)^4|p|^4}
 =\frac{log q}{8\pi^2},
\]
which gives \eqref{eq:v4v6-shell-constant} after differentiating
\(g^{-2}\).

For \(SU(3)\), \(C_A=3\), the twelve Weyl triplet occurrences have
\(\sum_fT(R_f)=6\), and there is no colored scalar.  Hence
\(b_3=11-4=7\).  For \(SU(2)\), \(C_A=2\), the twelve Weyl doublets have
\(\sum_fT(R_f)=6\), and the Higgs has \(T(\mathbf2)=1/2\).  Thus
\[
 b_2=\frac{22}{3}-4-\frac16=\frac{19}{6}.
\]
For hypercharge, one generation has
\[
 \sum_fY^2
 =6\left(\frac16\right)^2
 +3\left(\frac23\right)^2
 +3\left(\frac13\right)^2
 +2\left(\frac12\right)^2+1
 =\frac{10}{3}.
\]
Three generations give \(10\), while the two complex Higgs components give
\(\sum_sY^2=1/2\).  The abelian gauge term is zero and matter has the
screening sign, so
\[
 b_Y=-\frac23(10)-\frac13\left(\frac12\right)
 =-\frac{41}{6}.
\]
\end{proof}

\begin{firewall}[Scope of the coefficient calculation]
\label{fw:v4v6-one-loop-scope}
Proposition~\ref{prop:v4v6-one-loop} is an exact calculation of the first
Gaussian-boundary coefficient in the fixed-flat-metric shell expansion.  It
does not prove that the full shell integral is analytic at zero coupling,
that its remainder is uniform, or that dynamical gravity leaves the
coefficient unchanged.  Those are separate estimates below.
\end{firewall}

\subsection{Canonical factorization of the two-point response}
\label{sec:v4v6-factorization}

Let \(\boldsymbol\xi\) collect the canonical dimensionless couplings on one
shell, including \(g_3,g_2,g_Y\), Yukawa entries, and the Higgs quartic.  For
this flat-shell argument, gravitational fluctuation couplings are held at
zero.  Let \(\delta Z_{r,j}(\boldsymbol\xi)\) be the transverse correction
to the canonically normalized background two-point form after the complete
V5 Ward packet has been assembled.

\begin{lemma}[Quadratic external-coupling factor]
\label{lem:v4v6-external-factor}
Suppose the localized small-field shell integral is twice differentiable in
the background and holomorphic in the canonical vertices near
\(\boldsymbol\xi=0\).  Then
\begin{equation}
 \delta Z^{\rm sf}_{r,j}(\boldsymbol\xi)
 =g_r^2\Psi^{\rm sf}_{r,j}(\boldsymbol\xi)
\label{eq:v4v6-g2-factor}
\end{equation}
with \(\Psi^{\rm sf}_{r,j}\) holomorphic near zero.
\end{lemma}

\begin{proof}
At \(g_r=0\), the canonical background of factor \(r\) decouples from all
charged shell fields, so the correction vanishes.  A term linear in \(g_r\)
would contain one external gauge vertex.  Its internal Lie-algebra trace is
zero, equivalently it is odd under simultaneous reversal of that background
in the centered Gaussian coefficient.  Every nonzero two-point contraction
therefore contains either two linear gauge vertices or one seagull vertex,
and carries \(g_r^2\).  Holomorphy makes the vanishing of the constant and
linear Taylor coefficients an exact factorization by \(g_r^2\).
\end{proof}

With the scale orientation of V5, recanonicalization gives
\begin{equation}
 \beta^{\rm sf}_{r,j}
 =-\frac{\delta Z^{\rm sf}_{r,j}}{g_r^2}
 =-\Psi^{\rm sf}_{r,j}.
\label{eq:v4v6-beta-quotient}
\end{equation}
The sign in this equation is fixed by
\(\beta^{\rm ex}_{r,j}=z_{r,j}-z_{r,j-1}\).  Proposition~\ref{prop:v4v6-one-loop}
says \(-\Psi^{\rm sf}_{r,j}(0)=c_r\).

\subsection{Uniform small-field control}
\label{sec:v4v6-small-field}

For \(d\) canonical coordinates, let
\[
 \mathbb D_\rho^d
 =\{\boldsymbol\xi\in\mathbb C^d:
 |\xi_k|<\rho\text{ for every }k\}.
\]

\begin{lemma}[Polydisc Cauchy remainder]
\label{lem:v4v6-cauchy}
Let \(f\) be holomorphic on \(\mathbb D_\rho^d\) and satisfy
\(\sup_{\mathbb D_\rho^d}|f|\le B\).  If
\(\|\boldsymbol\xi\|_\infty\le\rho/2\), then
\begin{equation}
 |f(\boldsymbol\xi)-f(0)|
 \le\frac{2B}{\rho}\|\boldsymbol\xi\|_1.
\label{eq:v4v6-cauchy-bound}
\end{equation}
\end{lemma}

\begin{proof}
Join zero to \(\boldsymbol\xi\) by changing one coordinate at a time.  On
each segment, the circle of radius \(\rho/2\) about the current coordinate
remains inside its radius-\(\rho\) disc.  The one-variable Cauchy estimate
therefore gives \(|\partial_kf|\le2B/\rho\).  Integrating along the coordinate
segments and summing proves \eqref{eq:v4v6-cauchy-bound}.
\end{proof}

For background derivatives in factor \(r\), write
\[
 \mathcal D_{r,A}^{(m)}=g_r^{-m}D_A^m,
 \qquad 0\le m\le2,
\]
with the value at \(g_r=0\) supplied by the removable factor in
Lemma~\ref{lem:v4v6-external-factor}.

\begin{theorem}[A uniform polymer hypothesis supplies the Cauchy bound]
\label{thm:v4v6-polymer}
Suppose the localized shell logarithm has polymer activities
\(w_j(X;\boldsymbol\xi,A)\), holomorphic on \(\mathbb D_\rho^d\), and for
some \(\theta>0\) and \(\eta<1\) satisfies, uniformly in shell, volume, and
small background,
\begin{equation}
 \sup_x\sum_{X\ni x}
 e^{\theta|X|}
 \max_{0\le m\le2}
 \left\|\mathcal D_{r,A}^{(m)}
 w_j(X;\boldsymbol\xi,A)\right\|
 \le\eta.
\label{eq:v4v6-KP}
\end{equation}
Assume the transverse local projection is bounded in the same norm.  Then
the projected connected shell logarithm and its first two background
derivatives converge absolutely, are holomorphic on the polydisc, and obey
a shell- and volume-uniform bound
\begin{equation}
 \sup_{\mathbb D_\rho^d}
 |\Psi^{\rm sf}_{r,j}|\le B_r(\eta,\theta).
\label{eq:v4v6-Psi-bound}
\end{equation}
\end{theorem}

\begin{proof}
Expand the logarithm as the sum over connected families of polymers.  Give
each incompatibility graph a rooted spanning tree and sum leaves toward the
root.  At every leaf, \eqref{eq:v4v6-KP} pays its incident polymer together
with the factor \(e^{\theta|X|}\).  The remaining tree series is dominated by
a geometric series because \(\eta<1\).  Marking at most two coupling-normalized background derivatives chooses
one or two activities; the maximum over \(0\le m\le2\) in
\eqref{eq:v4v6-KP} pays those marks without changing convergence.  Absolute convergence permits termwise differentiation and
preserves holomorphy.  The bounded local transverse projection then gives
\eqref{eq:v4v6-Psi-bound}.  All estimates are independent of the number of
available polymer roots after taking the local coefficient per root.
\end{proof}

Theorem~\ref{thm:v4v6-polymer} is a sufficient hypothesis, not an assertion
that the four-dimensional gauge shell already satisfies
\eqref{eq:v4v6-KP}.  Its value is that it identifies a finite local estimate
which is strong enough and which keeps the Ward packet intact.

\subsection{Putting back the large-field sector}
\label{sec:v4v6-large-field}

Let \(\mathsf S_j\) be the small-field event used in the polymer expansion
and \(\mathsf L_j=\mathsf S_j^c\).  The next lemma records the precise cost
of replacing a complete expectation by its small-field conditioning.

\begin{lemma}[Rare-event comparison of a connected pair]
\label{lem:v4v6-rare-event}
Let \(P(\mathsf L)=\varepsilon<1/2\).  For real random variables \(X,Y,Z\)
with
\begin{equation}
 \max(\|X\|_4,\|Y\|_4,\|Z\|_2)\le M,
\label{eq:v4v6-moment-bound}
\end{equation}
there is a numerical constant \(C\) such that
\begin{align}
 |\mathbb E Z-\mathbb E[Z\mid\mathsf S]|
 &\le CM\varepsilon^{1/2},
\label{eq:v4v6-mean-rare}\\
 |\langle X;Y\rangle^{\rm conn}
 -\langle X;Y\rangle_{\mathsf S}^{\rm conn}|
 &\le CM^2\varepsilon^{1/2}.
\label{eq:v4v6-cov-rare}
\end{align}
\end{lemma}

\begin{proof}
For any \(W\in L^2\), Cauchy--Schwarz gives
\(|\mathbb E(W1_{\mathsf L})|\le\|W\|_2\varepsilon^{1/2}\).
Also \(|(1-\varepsilon)^{-1}-1|\le2\varepsilon\).  Apply these two facts to
the conditional means.  For the covariance, first apply them to \(XY\),
using \(\|XY\|_2\le\|X\|_4\|Y\|_4\), and then to the two products of
means.  Enlarging a numerical constant proves both inequalities.
\end{proof}

\begin{proposition}[Large-field shell error]
\label{prop:v4v6-large-field}
Suppose, uniformly in shell and volume,
\begin{equation}
 P(\mathsf L_j)\le C_0e^{-\kappa/u_*},
 \qquad
 M_j\le C_1u_*^{-m},
\label{eq:v4v6-large-field-hyp}
\end{equation}
where \(u_*=\max_r g_r^2\) and \(M_j\) bounds the moments of the
coupling-normalized insertions
\(g_r^{-1}D_AF_j\) and \(g_r^{-2}D_A^2F_j\) in
Lemma~\ref{lem:v4v6-rare-event}.  Then the difference between the complete
V5 beta packet and its small-field-conditioned value satisfies
\begin{equation}
 |R^{\rm lf}_{r,j}|
 \le C_2u_*^{-2m}e^{-\kappa/(2u_*)}.
\label{eq:v4v6-large-field-error}
\end{equation}
\end{proposition}

\begin{proof}
Divide each first background insertion by \(g_r\) and the second
background insertion by \(g_r^2\) before applying
Lemma~\ref{lem:v4v6-rare-event} to the complete packet
\eqref{eq:v4v5-second-response}.  This is an algebraic extraction of the
two external coupling factors, not a division of an estimated remainder.
The largest resulting power is \(M_j^2\), and conditioning costs
\(e^{-\kappa/(2u_*)}\).  Absorb smaller powers and the bounded transverse
projection into \(C_2\).
\end{proof}

\subsection{The exact one-shell dominance theorem}
\label{sec:v4v6-dominance}

\begin{theorem}[One-loop sign from uniform small/large-field control]
\label{thm:v4v6-dominance}
Fix a gauge factor with \(c_r\ne0\).  Suppose
Theorem~\ref{thm:v4v6-polymer} and
Proposition~\ref{prop:v4v6-large-field} hold and the small-field value at the
Gaussian boundary is \(c_r\).  If
\begin{align}
 \|\boldsymbol\xi\|_\infty&\le\rho/2,
\label{eq:v4v6-small-wedge-a}\\
 \frac{2B_r}{\rho}\|\boldsymbol\xi\|_1
 &\le\frac{|c_r|}{4},
\label{eq:v4v6-small-wedge-b}\\
 C_2u_*^{-2m}e^{-\kappa/(2u_*)}
 &\le\frac{|c_r|}{4},
\label{eq:v4v6-small-wedge-c}
\end{align}
then the exact shell increment has the sign of \(c_r\) and
\begin{equation}
 \frac{|c_r|}{2}
 \le |\beta^{\rm ex}_{r,j}|
 \le\frac{3|c_r|}{2}.
\label{eq:v4v6-dominance-bound}
\end{equation}
The constants are uniform in shell and volume.
\end{theorem}

\begin{proof}
Lemma~\ref{lem:v4v6-cauchy} and
\eqref{eq:v4v6-Psi-bound} show that the small-field beta increment differs
from \(c_r\) by at most \(|c_r|/4\).  Proposition~\ref{prop:v4v6-large-field}
and \eqref{eq:v4v6-small-wedge-c} bound the restoration of the omitted
sector by another \(|c_r|/4\).  Therefore
\(|\beta^{\rm ex}_{r,j}-c_r|\le|c_r|/2\).  The claimed sign and two-sided
bound follow from the triangle inequality.
\end{proof}

\begin{corollary}[Conditional nonabelian signed shell bound]
\label{cor:v4v6-nonabelian}
If an ultraviolet SM trajectory stays in the wedge
\eqref{eq:v4v6-small-wedge-a}--%
\eqref{eq:v4v6-small-wedge-c}, then the \(SU(3)\) and \(SU(2)\) increments
satisfy \eqref{eq:v4v5-beta-bounds} with
\begin{equation}
 b_-^{(r)}=\frac{b_r\log q}{16\pi^2},
 \qquad
 b_+^{(r)}=\frac{3b_r\log q}{16\pi^2}.
\label{eq:v4v6-nonabelian-bounds}
\end{equation}
Consequently the V5 inverse-coupling telescope and the V4 irrelevant
theorem apply, subject to their remaining hypotheses.
\end{corollary}

\begin{proof}
For \(r=3,2\), Proposition~\ref{prop:v4v6-one-loop} gives \(c_r>0\).
Insert \eqref{eq:v4v6-shell-constant} into
\eqref{eq:v4v6-dominance-bound} and apply
Theorem~\ref{thm:v4v5-beta-telescope} and
Corollary~\ref{cor:v4v5-v4-compatible}.
\end{proof}

\begin{theorem}[Conditional hypercharge Gaussian-boundary obstruction]
\label{thm:v4v6-hypercharge-obstruction}
Assume an ultraviolet hypercharge trajectory remains in the same uniform
dominance wedge for every shell and has finite positive
\(z_{Y,0}=g_{Y,0}^{-2}\).  Then it cannot remain a positive-coupling
trajectory for infinitely many ultraviolet shells.
\end{theorem}

\begin{proof}
Proposition~\ref{prop:v4v6-one-loop} gives
\(c_Y=-41\log q/(48\pi^2)<0\).  Theorem~\ref{thm:v4v6-dominance} therefore
gives
\[
 \beta^{\rm ex}_{Y,j}\le-\frac{|c_Y|}{2}.
\]
The exact telescope yields
\[
 z_{Y,j}
 =z_{Y,0}+\sum_{k=1}^j\beta^{\rm ex}_{Y,k}
 \le z_{Y,0}-\frac{j|c_Y|}{2}.
\]
The right side is negative once \(j>2z_{Y,0}/|c_Y|\), contradicting
\(z_{Y,j}=g_{Y,j}^{-2}>0\).
\end{proof}

\begin{firewall}[Meaning of the hypercharge obstruction]
\label{fw:v4v6-hypercharge-scope}
Theorem~\ref{thm:v4v6-hypercharge-obstruction} does not prove a Landau pole
for every nonperturbative completion.  It proves that the unmodified SM
hypercharge factor cannot supply an infinite positive Gaussian-boundary
trajectory while all hypotheses of the uniform small-coupling dominance
theorem remain valid.  Escaping the wedge, embedding \(U(1)_Y\) in a
nonabelian group, or adding new ultraviolet degrees of freedom lies outside
the theorem.
\end{firewall}

\subsection{V6 result ledger and V7 target}
\label{sec:v4v6-conclusion}

V6 has converted the V5 sign question into two explicit local estimates.  A
uniform Ward-complete polymer bound controls the small-field connected
remainder by Cauchy's estimate.  An exponential large-field probability,
together with insertion moments, restores the omitted sector at cost
\(u_*^{-M}e^{-\kappa/(2u_*)}\).  Under these hypotheses, the exact beta
increments inherit their one-loop signs with quantitative two-sided bounds.

For \(SU(3)\) and \(SU(2)\), this supplies precisely the signed increment
assumed in V5.  For hypercharge, the same argument gives a conditional
obstruction to an infinite Gaussian-boundary trajectory.  The full
SM--Einstein continuum has therefore not been constructed: the polymer and
large-field hypotheses remain unproved, the coupled Higgs--Yukawa trajectory
has not been kept inside the wedge, dynamical gravity has been excluded from
this shell calculation, and hypercharge requires a new ultraviolet
mechanism.

V7 should go upstream to the first unproved estimate rather than repeat the
one-loop calculation.  It should build a gauge-invariant small-field/large-
field partition for one compact nonabelian shell, prove the exact plaquette
energy suppression available from the Wilson action, and determine whether
that suppression is strong enough to pay the derivative-marked polymer norm
\eqref{eq:v4v6-KP} without treating gauge orbit directions as massive.


\clearpage
\section{V7: gauge-invariant plaquette sectors and reduction to a local repair ratio}
\label{sec:v4v7}

V6 reduced exact one-loop dominance to a coupling-normalized polymer bound
and a negligible large-field contribution.  This note constructs the
small-field/large-field partition in gauge-invariant variables.  A spanning-
tree gauge then converts small plaquette curvature into a local link chart
without penalizing pure-gauge directions.  Wilson energy gives the correct
stretched-exponential cost of every bad plaquette.  Connected-set counting
shows that this cost is more than sufficient for the derivative-marked
polymer norm, provided one additional local comparison holds: a uniform
repair ratio between a prescribed bad component and its repaired interior.

That repair ratio, rather than plaquette energy or gauge fixing, is the first
unproved large-field gate.  It is stated below in a form stable under
arbitrary exterior links.  The note also corrects an unnecessarily strong
V6 hypothesis: \(e^{-\kappa/u}\) is not required.  Any
\(e^{-\kappa u^{-\alpha}}\), \(\alpha>0\), dominates the polynomial losses
arising from response insertions.

\subsection{Wilson energy and invariant bad plaquettes}
\label{sec:v4v7-wilson}

Let \(G=SU(N)\) with its bi-invariant Riemannian distance \(d_G\), and put
\begin{equation}
 \ell(U)=1-\frac1N\operatorname{Re}\operatorname{Tr}U.
\label{eq:v4v7-wilson-potential}
\end{equation}
On a finite four-dimensional cubical lattice, use
\begin{equation}
 S_g(U)=\frac1{g^2}\sum_p\ell(U_p),
\label{eq:v4v7-wilson-action}
\end{equation}
where \(U_p\) is the oriented plaquette holonomy.

\begin{lemma}[Quadratic comparison on the compact group]
\label{lem:v4v7-group-comparison}
There are \(r_0,c_-,c_+>0\), depending only on \(N\) and the chosen metric,
such that
\begin{equation}
 c_-d_G(U,1)^2\le\ell(U)\le c_+d_G(U,1)^2
\label{eq:v4v7-local-comparison}
\end{equation}
whenever \(d_G(U,1)\le r_0\).  Moreover, for every \(0<r\le r_0\),
\begin{equation}
 d_G(U,1)\ge r
 \quad\Longrightarrow\quad
 \ell(U)\ge c_-r^2.
\label{eq:v4v7-global-lower}
\end{equation}
\end{lemma}

\begin{proof}
In exponential coordinates \(U=e^X\), the Hessian of \(\ell(e^X)\) at
zero is a positive multiple of the invariant Hilbert--Schmidt form on
\(\mathfrak{su}(N)\).  Taylor's theorem gives
\eqref{eq:v4v7-local-comparison} after reducing \(r_0\).  On the compact set
\(\{U:d_G(U,1)\ge r_0\}\), \(\ell\) has a strictly positive minimum because
its unique zero is the identity.  Reduce \(c_-\), if needed, to combine that
minimum with the local lower bound for every \(r\le r_0\).  This proves
\eqref{eq:v4v7-global-lower}.
\end{proof}

Fix \(0<\alpha<1\) and, for small \(g\), define the invariant threshold
\begin{equation}
 r_g=g^{1-\alpha}.
\label{eq:v4v7-threshold}
\end{equation}
A plaquette is good when \(d_G(U_p,1)<r_g\) and bad otherwise.  This
classification depends only on conjugacy-invariant plaquette data and hence
does not constrain a pure-gauge link field.

\begin{proposition}[Deterministic bad-plaquette energy]
\label{prop:v4v7-bad-energy}
If \(B(U)\) is the set of bad plaquettes, then
\begin{equation}
 S_g(U)\ge c_-g^{-2\alpha}|B(U)|.
\label{eq:v4v7-bad-energy}
\end{equation}
\end{proposition}

\begin{proof}
For each bad plaquette, Lemma~\ref{lem:v4v7-group-comparison} and
\eqref{eq:v4v7-threshold} give
\(g^{-2}\ell(U_p)\ge c_-g^{-2}r_g^2=c_-g^{-2\alpha}\).
Sum these nonnegative contributions.
\end{proof}

\subsection{Tree gauge does not create a gauge mass}
\label{sec:v4v7-tree-gauge}

Let \(Q\) be a fixed simply connected block of side \(q\) lattice units.
Choose a root vertex and a spanning tree \(T\) in its one-skeleton.

\begin{lemma}[Exact spanning-tree gauge factorization]
\label{lem:v4v7-tree-factorization}
Every link field on \(Q\) has a unique gauge transform, with the root gauge
fixed to the identity, for which every tree link equals the identity.
Product Haar measure factors as the product of Haar measure on the nonroot
site gauges and product Haar measure on the remaining links.  In particular,
tree gauge introduces no field-dependent Faddeev--Popov determinant.
\end{lemma}

\begin{proof}
Starting at the root, define the gauge value at each new tree vertex by the
ordered product of tree links along its unique root path.  This makes every
tree link equal to one and proves existence and uniqueness.  The change of
variables is a finite sequence of left and right translations of group
variables.  Haar invariance makes its Jacobian one, giving the stated
factorization.
\end{proof}

\begin{lemma}[Small curvature gives a local link chart]
\label{lem:v4v7-small-curvature-chart}
There is a constant \(C_q\), depending only on the fixed block, such that if
all plaquettes of \(Q\) satisfy
\(d_G(U_p,1)<r\), then in a suitable tree gauge every remaining link obeys
\begin{equation}
 d_G(U_e,1)\le C_qr.
\label{eq:v4v7-link-chart}
\end{equation}
Consequently, on the good event \(r=r_g\), the canonical coordinate in
\(U_e=\exp(agB_e)\) has size
\begin{equation}
 \|aB_e\|\le C'_qg^{-\alpha}
\label{eq:v4v7-canonical-size}
\end{equation}
provided \(g\) is small enough for the exponential chart.
\end{lemma}

\begin{proof}
Choose the tree so that each non-tree edge closes a loop bounding a fixed
cubical surface in \(Q\).  In tree gauge, its link holonomy is an ordered
product of conjugates of the plaquette holonomies on that surface.  The
number of plaquettes is bounded by a constant depending only on \(q\).
Bi-invariance and the triangle inequality give
\[
 d_G(U_e,1)
 \le\sum_{p\text{ in the surface}}d_G(U_p,1)
 \le C_qr.
\]
For \(r=r_g\), local equivalence of exponential-coordinate norm and group
distance gives
\(ag\|B_e\|\le C''_qr_g\), which is
\eqref{eq:v4v7-canonical-size}.
\end{proof}

\begin{firewall}[The shell gap is not a gauge-boson mass]
\label{fw:v4v7-no-false-mass}
Lemma~\ref{lem:v4v7-small-curvature-chart} fixes only a redundant tree
coordinate and bounds links by curvature on a fixed block.  The quadratic
shell covariance may have a lower momentum \(K/q\) because it is a band
covariance.  Neither fact inserts a term \(m^2B^2\) into the gauge-invariant
effective action.  The Ward identity of Theorem~\ref{thm:v4v5-hessian-ward}
remains exact.
\end{firewall}

\subsection{Bad components and their entropy}
\label{sec:v4v7-components}

Declare two plaquettes adjacent when they share an edge.  Let \(D_4\) be the
maximum degree of this adjacency graph; one may take \(D_4\le20\).  The bad
set is the disjoint union of its connected components.

\begin{lemma}[Connected plaquette count]
\label{lem:v4v7-connected-count}
In a graph of maximum degree \(D\), the number of connected vertex sets of
cardinality \(m\) containing a fixed vertex is at most
\begin{equation}
 (eD)^{m-1}.
\label{eq:v4v7-animal-count}
\end{equation}
\end{lemma}

\begin{proof}
Every connected set has a rooted spanning tree.  Explore that tree in depth-
first order, recording for each newly discovered vertex one of at most
\(D\) incident choices and the rooted-tree shape.  The number of rooted tree
shapes with \(m\) vertices is at most \(e^{m-1}\).  Multiplication gives
\eqref{eq:v4v7-animal-count}; overcounting only strengthens the bound.
\end{proof}

The deterministic energy estimate alone does not bound a normalized Gibbs
probability: the partition function and the boundary frustration of a local
component must also be compared.  The required comparison can be isolated
without choosing a global gauge.

\begin{definition}[Uniform local repair ratio]
\label{def:v4v7-repair-ratio}
For a connected plaquette set \(X\), let
\(Z_X^{\rm bad}(\eta)\) be the conditional integral over its internal links,
with exterior links \(\eta\) fixed and with every plaquette in \(X\) bad.
Let \(Z_X^{\rm ref}(\eta)\) be the corresponding integral over a repaired
interior in which those plaquettes are good, with the same exterior data.
A uniform repair inequality is
\begin{equation}
 \frac{Z_X^{\rm bad}(\eta)}{Z_X^{\rm ref}(\eta)}
 \le
 \exp\!\left[-\bigl(c_0g^{-2\alpha}-c_1\bigr)|X|
 +c_2|\partial X|\right]
\label{eq:v4v7-repair-ratio}
\end{equation}
for constants independent of \(X\), shell, volume, and \(\eta\).
\end{definition}

The boundary term is harmless because \(|\partial X|\le C_\partial|X|\).
What matters is that the Wilson energy \(g^{-2\alpha}|X|\) survives
normalization uniformly in the exterior configuration.

\begin{theorem}[Repair ratio implies a defect-polymer KP bound]
\label{thm:v4v7-repair-to-KP}
Assume \eqref{eq:v4v7-repair-ratio} and a bad-component polymer
representation whose component activity is bounded in modulus by the repair
ratio.  Suppose that at most two coupling-normalized background derivatives
multiply that activity by
\begin{equation}
 C_3g^{-M}(1+|X|)^2.
\label{eq:v4v7-derivative-cost}
\end{equation}
Then for every fixed \(\theta>0\) and \(\eta_0>0\), the bad-component
activities satisfy the marked KP estimate
\begin{equation}
 \sup_p\sum_{X\ni p}e^{\theta|X|}
 \max_{0\le m\le2}
 \|\mathcal D_{r,A}^{(m)}w_j(X)\|
 \le\eta_0
\label{eq:v4v7-defect-KP}
\end{equation}
for all sufficiently small \(g\), uniformly in shell and volume.
\end{theorem}

\begin{proof}
Use \(|\partial X|\le C_\partial|X|\) in
\eqref{eq:v4v7-repair-ratio} and define
\[
 \sigma_g=c_0g^{-2\alpha}-c_1-c_2C_\partial.
\]
The repair ratio bounds the unmarked component activity by
\(e^{-\sigma_g|X|}\).  Equations
\eqref{eq:v4v7-derivative-cost} and
\eqref{eq:v4v7-animal-count} therefore give
\begin{align*}
 \sup_p\sum_{X\ni p}e^{\theta|X|}
 \max_{m\le2}\|\mathcal D_{r,A}^{(m)}w_j(X)\|
 &\le C_3g^{-M}
 \sum_{n\ge1}(1+n)^2(eD_4)^{n-1}
 e^{-(\sigma_g-\theta)n}.
\end{align*}
Since \(\sigma_g\to\infty\), the series and its polynomial prefactor tend
to zero faster than any power of \(g\).  It is therefore at most
\(\eta_0\) for sufficiently small \(g\).
\end{proof}

\subsection{The stretched-exponential correction to V6}
\label{sec:v4v7-stretched}

\begin{lemma}[Stretched exponentials absorb all response powers]
\label{lem:v4v7-stretched}
For \(\alpha,\kappa>0\) and every \(M\ge0\),
\begin{equation}
 \lim_{u\downarrow0}u^{-M}e^{-\kappa u^{-\alpha}}=0.
\label{eq:v4v7-stretched-limit}
\end{equation}
\end{lemma}

\begin{proof}
Set \(t=u^{-\alpha}\).  The expression becomes
\(t^{M/\alpha}e^{-\kappa t}\), which tends to zero as
\(t\to\infty\).
\end{proof}

\begin{proposition}[Rooted stretched-exponential replacement for V6]
\label{prop:v4v7-weaken-v6}
Assume Theorem~\ref{thm:v4v7-repair-to-KP}.  Anchor the transverse local
coefficient defining \(\beta^{\rm ex}_{r,j}\) in a fixed finite plaquette
neighborhood \(A_0\).  Then the part of that coefficient containing at least
one bad component satisfies, for sufficiently small \(g\),
\begin{equation}
 |R^{\rm lf}_{r,j}|
 \le C_4g^{-M'}
 \exp\!\left[-\frac{c_0}{2}g^{-2\alpha}\right]
 =C_4u^{-M'/2}
 \exp\!\left[-\frac{c_0}{2}u^{-\alpha}\right],
\label{eq:v4v7-weaker-error}
\end{equation}
uniformly in shell and volume.  This rooted estimate may replace the global
rare-event hypothesis in the application of
Theorem~\ref{thm:v4v6-dominance}.
\end{proposition}

\begin{proof}
Differentiate the connected polymer logarithm at the local background
insertion.  Every surviving cluster containing a bad component is connected
to \(A_0\); components disjoint from the marked cluster cancel against the
logarithmic normalization.  Sum first over the finitely many plaquettes of
\(A_0\), then over the first bad component attached to the marked cluster.
The proof of Theorem~\ref{thm:v4v7-repair-to-KP} bounds a component of size
\(n\) by a polynomial in \(g^{-1}\) and \(n\), times
\[
 \exp[-(c_0g^{-2\alpha}-C_5)n].
\]
Connected-set counting and the convergent good-field cluster attachments
change only \(C_5\) and the polynomial power.  For small \(g\), the sum over
\(n\ge1\) is bounded by its first exponential scale, giving
\eqref{eq:v4v7-weaker-error}.  Lemma~\ref{lem:v4v7-stretched} makes this
smaller than any prescribed positive constant.
\end{proof}

\begin{firewall}[A global bad-event probability is the wrong limit]
\label{fw:v4v7-global-event}
On an increasing lattice, the probability that at least one plaquette is
bad need not be small and normally tends to one.  Thus the global
\(P(\mathsf L_j)\) hypothesis used as a convenient finite-dimensional
criterion in Proposition~\ref{prop:v4v6-large-field} is too strong for a
thermodynamic-limit proof.  Proposition~\ref{prop:v4v7-weaken-v6}
supersedes it for the local beta coefficient: only connected defect clusters
rooted at the response insertion are estimated.
\end{firewall}

For the threshold \eqref{eq:v4v7-threshold}, the Wilson cost in
Proposition~\ref{prop:v4v7-bad-energy} is
\(e^{-c g^{-2\alpha}}=e^{-c u^{-\alpha}}\), exactly of the sufficient rooted
form \eqref{eq:v4v7-weaker-error} once the local repair ratio is available.

\begin{firewall}[Energy is not yet a normalized defect estimate]
\label{fw:v4v7-repair-open}
Proposition~\ref{prop:v4v7-bad-energy} is pointwise.  It does not prove
\eqref{eq:v4v7-repair-ratio}.  In four-dimensional nonabelian lattice gauge
theory, shared links, Bianchi constraints, boundary holonomy, and toron
sectors can change the conditional normalization.  Dividing the bad
Boltzmann weight by an unrelated global lower bound for the partition
function would introduce a volume factor and would not prove the local
polymer estimate.
\end{firewall}

\subsection{V7 result ledger and V8 target}
\label{sec:v4v7-conclusion}

V7 proves that a plaquette-defined sector is the correct gauge-invariant
large-field partition.  Its Wilson energy is
\(c g^{-2\alpha}\) per bad plaquette; spanning-tree gauge has unit Jacobian
and converts the complementary small-curvature event into a canonical local
chart without adding a mass.  Connected-component entropy is only
exponential, while the energy grows as \(g^{-2\alpha}\).  Consequently a
uniform local repair ratio implies the complete derivative-marked defect KP
bound.  The stretched-exponential cost is already sufficient for the V6
one-loop dominance theorem.

The first missing estimate is now \eqref{eq:v4v7-repair-ratio}.  V8 should
attack it on one contractible shell block.  The appropriate upstream route
is an explicit axial-gauge plaquette-coordinate map: compute its exact Haar
Jacobian, separate independent plaquettes from nonabelian Bianchi
constraints, and compare a bad connected component with a repaired
curvature configuration while leaving exterior holonomy fixed.  If a
uniform repair fails, the counterexample should identify whether boundary
frustration or a topological sector must be promoted into the marginal
parent.


\clearpage
\section{V8: free-block repair, boundary frustration, and the variational gap}
\label{sec:v4v8}

V7 isolated a local repair ratio which would turn Wilson plaquette energy
into the derivative-marked defect expansion required by V6.  The ratio was
stated uniformly in arbitrary exterior links.  This note tests that
uniformity.  On a fixed contractible block with free boundary, tree-gauge
Haar coordinates give the desired normalized stretched-exponential bound
directly.  With fixed exterior data, however, the proposed comparison can
fail because the boundary may force curvature and leave the repaired set
empty.  Thus the V7 arbitrary-exterior hypothesis is too strong.

The correct local datum is a variational energy gap relative to the best
boundary-compatible reference configuration, together with an interior
reference ball.  A finite-dimensional Laplace comparison proves that this
gap implies the required repair ratio.  Bad components without such a gap
must be joined to their boundary or ancestor polymer; they cannot be erased
locally.

\subsection{Exact free-block Haar coordinates}
\label{sec:v4v8-free-coordinates}

Let \(Q\) be a fixed finite connected cubical block, with link set \(E_Q\)
and vertex set \(V_Q\).  Choose a root and a spanning tree \(T\).  Put
\begin{equation}
 n_Q=|E_Q|-|V_Q|+1,
 \qquad d_G=\dim SU(N)=N^2-1.
\label{eq:v4v8-coordinate-count}
\end{equation}

\begin{proposition}[Tree-gauge quotient measure]
\label{prop:v4v8-tree-measure}
For every gauge-invariant integrable function \(F\),
\begin{equation}
 \int_{G^{E_Q}}F(U)\,dU
 =\int_{G^{E_Q\setminus T}}F(U^{T}=1,U_{E_Q\setminus T})
 \prod_{e\notin T}dU_e.
\label{eq:v4v8-tree-measure}
\end{equation}
Thus the free-boundary gauge quotient has exactly \(n_Q\) independent Haar
link coordinates and no field-dependent Jacobian.
\end{proposition}

\begin{proof}
Apply Lemma~\ref{lem:v4v7-tree-factorization}.  It factors product Haar
measure into the nonroot vertex-gauge variables and the \(n_Q\) non-tree
links.  Gauge invariance makes the integrand independent of the former.
Their normalized Haar integral is one, leaving
\eqref{eq:v4v8-tree-measure}.
\end{proof}

\begin{lemma}[Haar volume of a small group ball]
\label{lem:v4v8-haar-ball}
There are \(r_1,c_1,c_2>0\) such that for \(0<r<r_1\),
\begin{equation}
 c_1r^{d_G}
 \le \operatorname{Haar}\{U:d_G(U,1)<r\}
 \le c_2r^{d_G}.
\label{eq:v4v8-haar-ball}
\end{equation}
\end{lemma}

\begin{proof}
The exponential map is a diffeomorphism on a neighborhood of zero.  In
these coordinates Haar measure has a smooth strictly positive density whose
value at zero is nonzero.  Compare that density with its minimum and maximum
on a smaller coordinate ball and use Euclidean ball volume.
\end{proof}

Let
\[
 E_Q(U)=\sum_{p\subset Q}\ell(U_p)
\]
be the unscaled Wilson energy.  For a nonempty plaquette set \(X\), define
\begin{align}
 Z^{\rm bad}_{Q,X}(g)
 &=\int 1_{\{d_G(U_p,1)\ge r_g\;\forall p\in X\}}
 e^{-E_Q(U)/g^2}\,dU,
\label{eq:v4v8-free-bad}\\
 Z^{\rm ref}_{Q,X}(g)
 &=\int 1_{\{d_G(U_p,1)< r_g\;\forall p\in X\}}
 e^{-E_Q(U)/g^2}\,dU.
\label{eq:v4v8-free-ref}
\end{align}
The integrals are over the tree-gauge quotient in
\eqref{eq:v4v8-tree-measure}.

\begin{theorem}[Free-block repair bound]
\label{thm:v4v8-free-repair}
For the threshold \(r_g=g^{1-\alpha}\), \(0<\alpha<1\), there are constants
\(C_Q,c_Q,g_Q>0\) such that for every nonempty \(X\) and
\(0<g<g_Q\),
\begin{equation}
 \frac{Z^{\rm bad}_{Q,X}(g)}{Z^{\rm ref}_{Q,X}(g)}
 \le C_Qg^{-d_Gn_Q}
 \exp[-c_Qg^{-2\alpha}|X|].
\label{eq:v4v8-free-repair}
\end{equation}
\end{theorem}

\begin{proof}
Proposition~\ref{prop:v4v7-bad-energy} gives
\[
 Z^{\rm bad}_{Q,X}(g)
 \le\exp[-c_-g^{-2\alpha}|X|]
\]
because normalized Haar volume is at most one.

For the denominator, restrict every one of the \(n_Q\) non-tree links to
the group ball of radius \(c_Q'g\).  Tree links equal one.  Each plaquette is
a product of four links, so bi-invariance and the triangle inequality give
\(d_G(U_p,1)\le4c_Q'g\).  For small \(g\), this is less than
\(r_g=g^{1-\alpha}\), hence the restricted set lies in the reference event.
Lemma~\ref{lem:v4v7-group-comparison} also gives
\(E_Q(U)/g^2\le C_Q'\), since \(Q\) is fixed.  By
Lemma~\ref{lem:v4v8-haar-ball}, the restricted Haar volume is at least
\(c_Q''g^{d_Gn_Q}\).  Therefore
\[
 Z^{\rm ref}_{Q,X}(g)
 \ge c_Q''e^{-C_Q'}g^{d_Gn_Q}.
\]
Divide the two estimates and rename constants.
\end{proof}

The polynomial factor in \eqref{eq:v4v8-free-repair} is harmless by
Lemma~\ref{lem:v4v7-stretched}.  Thus the large-field normalization problem
is solved on a fixed free block.

\subsection{Why arbitrary exterior data invalidate the naive ratio}
\label{sec:v4v8-boundary-counterexample}

\begin{proposition}[Boundary-forced failure of uniform repair]
\label{prop:v4v8-boundary-failure}
There is no repair inequality of the form
\eqref{eq:v4v7-repair-ratio} for all connected \(X\) and all fixed exterior
link configurations if the reference set requires every plaquette of
\(X\) to be good while all of its links may be fixed by that exterior.
\end{proposition}

\begin{proof}
Take \(X\) to consist of one plaquette at the boundary of the integration
region and regard all four of its links as exterior data.  Choose their
ordered product \(H\) with \(d_G(H,1)\ge r_g\).  There are no internal
variables on this cell.  The bad conditional integral is its positive
Boltzmann weight, while the reference condition
\(d_G(H,1)<r_g\) is impossible, so
\(Z_X^{\rm ref}(\eta)=0\).  The ratio is infinite.
\end{proof}

The example is elementary, but its implication is structural.  A local
component cannot be repaired while preserving boundary data which already
carry its required holonomy.  Higher-dimensional Bianchi constraints can
produce the same problem even when some interior links remain.

\begin{firewall}[Correction to the V7 repair hypothesis]
\label{fw:v4v8-v7-correction}
Definition~\ref{def:v4v7-repair-ratio} is a sufficient criterion only on
the class of components whose fixed exterior admits a quantitative repaired
interior.  Its phrase ``arbitrary exterior links'' is false without that
qualification.  Boundary-attached or topologically forced components must
be merged with the exterior ancestry or retained as a separate sector.
Theorem~\ref{thm:v4v7-repair-to-KP} remains valid on any class for which the
qualified ratio is proved.
\end{firewall}

\subsection{The boundary-compatible variational gap}
\label{sec:v4v8-gap}

Let \(x\in G^n\) denote the internal links of a fixed block after all
admissible gauge quotients, let \(\eta\) denote fixed exterior links, and let
\(E_\eta(x)\) include every Wilson plaquette term affected by \(x\).  Let
\(A_{X,\eta}\) be the boundary-compatible bad set and
\(R_{X,\eta}\) a proposed boundary-compatible repaired set.

\begin{definition}[Buffered repair datum]
\label{def:v4v8-buffered-datum}
A buffered repair datum consists of a point
\(x_{*,\eta}\in R_{X,\eta}\), constants \(L,\rho,c>0\), and the properties
\begin{align}
 |E_\eta(x)-E_\eta(y)|
 &\le Ld_{G^n}(x,y),
\label{eq:v4v8-Lipschitz}\\
 B_{G^n}(x_{*,\eta},\rho g^2)
 &\subset R_{X,\eta},
\label{eq:v4v8-reference-ball}\\
 \inf_{x\in A_{X,\eta}}E_\eta(x)-E_\eta(x_{*,\eta})
 &\ge c r_g^2|X|.
\label{eq:v4v8-variational-gap}
\end{align}
The constants must be uniform in the admitted exterior data.
\end{definition}

The ball radius \(g^2\) is chosen only to make its action variation order
one.  Any power giving that property and a polynomial Haar volume would do.

\begin{theorem}[Variational gap implies the conditional repair ratio]
\label{thm:v4v8-gap-to-repair}
Suppose a buffered repair datum exists for a fixed block with \(n\) internal
link coordinates.  Then
\begin{equation}
 \frac{Z_{X}^{\rm bad}(\eta)}{Z_X^{\rm ref}(\eta)}
 \le Cg^{-2d_Gn}
 \exp[-c g^{-2\alpha}|X|]
\label{eq:v4v8-gap-repair}
\end{equation}
for all sufficiently small \(g\), uniformly in its admitted exterior data.
\end{theorem}

\begin{proof}
The numerator is at most
\[
 \exp\!\left[-g^{-2}inf_{A_{X,\eta}}E_\eta\right].
\]
For the denominator, integrate only over the ball in
\eqref{eq:v4v8-reference-ball}.  The Lipschitz bound gives
\[
 E_\eta(x)
 \le E_\eta(x_{*,\eta})+L\rho g^2
\]
there.  A translated version of Lemma~\ref{lem:v4v8-haar-ball} gives ball
volume at least \(C^{-1}g^{2d_Gn}\).  Hence
\[
 Z_X^{\rm ref}(\eta)
 \ge C^{-1}g^{2d_Gn}
 \exp[-g^{-2}E_\eta(x_{*,\eta})-L\rho].
\]
Divide, use \eqref{eq:v4v8-variational-gap}, and insert
\(g^{-2}r_g^2=g^{-2\alpha}\).  Absorb \(e^{L\rho}\) into \(C\).
\end{proof}

\begin{corollary}[Buffered components satisfy the V7 defect bound]
\label{cor:v4v8-buffered-KP}
If every component in a defect class has a uniform buffered repair datum and
the number of internal coordinates is bounded by a constant times
\(|X|\), then its conditional ratios have the form
\eqref{eq:v4v7-repair-ratio} up to powers of \(g^{-1}\) per component.
These powers are absorbed by \(g^{-2\alpha}|X|\) at small coupling, and the
marked KP conclusion \eqref{eq:v4v7-defect-KP} follows.
\end{corollary}

\begin{proof}
Theorem~\ref{thm:v4v8-gap-to-repair} gives a factor
\(g^{-C|X|}=\exp[C|X|\log(g^{-1})]\).  Since
\(g^{-2\alpha}/\log(g^{-1})\to\infty\), half of the variational exponent
absorbs this factor for small \(g\).  The remaining exponent has the V7
repair form.  Apply Theorem~\ref{thm:v4v7-repair-to-KP}.
\end{proof}

\subsection{The corrected ancestry split}
\label{sec:v4v8-ancestry}

The preceding results force a three-way component classification.
\begin{enumerate}
\item A free interior block is paid by Theorem~\ref{thm:v4v8-free-repair}.
\item A boundary-compatible buffered component is paid by
Theorem~\ref{thm:v4v8-gap-to-repair}.
\item A component with no reference ball or no positive variational gap is
not a local fluctuation.  It must be joined to the boundary holonomy,
topological sector, or coarser ancestor which forces it.
\end{enumerate}
This is an ancestry assignment, not a second energy estimate.  Each bad
component enters exactly one of the three classes.

\begin{firewall}[The buffered variational gap is open in four dimensions]
\label{fw:v4v8-gap-open}
No theorem here constructs \(x_{*,\eta}\) or proves
\eqref{eq:v4v8-variational-gap} for every interior four-dimensional
nonabelian defect.  The free-boundary result does not imply it.  The
remaining issue is geometric: preserve the admitted boundary holonomy while
lowering the interior plaquette energy by order \(r_g^2|X|\), with a
uniform reference neighborhood.
\end{firewall}

\subsection{V8 result ledger and V9 target}
\label{sec:v4v8-conclusion}

V8 proves the full normalized stretched-exponential repair estimate on a
fixed free contractible block.  It also gives an explicit boundary
counterexample to the unqualified V7 repair ratio and replaces that ratio
by a precise variational-gap criterion.  The latter is sufficient for the
rooted defect KP expansion, including the polynomial cost of all local Haar
coordinates.  Boundary-forced components are now assigned to their
exterior or topological ancestor instead of being charged a fictitious local
repair.

V9 should attack the buffered gap in the first nontrivial geometry.  It
should use a two-layer contractible cube, fix the outer holonomy, and perform
an explicit radial contraction of the interior plaquette coordinates.  The
calculation must retain the nonabelian Bianchi word.  If the energy drop is
lost to a boundary commutator, that commutator should be isolated as the
coarser curvature coordinate carried by the ancestor polymer.


\clearpage
\section{V9: relative Hodge repair and the small nonabelian boundary chart}
\label{sec:v4v9}

V8 replaced the false arbitrary-exterior repair statement by a variational
gap relative to the best boundary-compatible interior.  This note proves
that gap in two regimes.  For an abelian quadratic curvature energy, relative
Hodge orthogonality gives an exact energy identity.  A bad component is
locally repairable precisely when the boundary-forced minimizing curvature
is small on that component.  For \(SU(N)\), spanning-tree gauge removes the
gauge kernel, the Wilson Hessian is positive on a fixed contractible block,
and finite-dimensional strong convexity persists for sufficiently small
boundary data.  This yields the same gap for bad components which remain in
that chart.

The remaining components are now sharply identified: their minimizing
curvature is already large, their boundary leaves the weak chart, or their
interpolation crosses the cut locus.  Such components must be carried by a
coarser boundary/topological ancestor and are the subject of the V10 audit.

\subsection{The exact abelian relative problem}
\label{sec:v4v9-abelian}

Let \(Q\) be a fixed contractible cubical block.  Give its real cochain
spaces the counting inner product.  Let \(C^1_0(Q)\) be the link cochains
whose boundary values vanish, and let
\[
 d_1:C^1_0(Q)\longrightarrow C^2(Q)
\]
be the relative coboundary.  Fixed exterior data determine an affine link
space \(a_\eta+C^1_0(Q)\).  Define
\begin{equation}
 \mathcal E_\eta(a)=\frac12\|d_1a\|_2^2.
\label{eq:v4v9-abelian-energy}
\end{equation}

\begin{theorem}[Exact relative Hodge energy identity]
\label{thm:v4v9-hodge-identity}
There is a minimizer \(a_{*,\eta}\) of
\eqref{eq:v4v9-abelian-energy}, unique modulo the relative kernel of
\(d_1\).  For every \(v\in C^1_0(Q)\),
\begin{equation}
 \mathcal E_\eta(a_{*,\eta}+v)
 =\mathcal E_\eta(a_{*,\eta})
 +\frac12\|d_1v\|_2^2.
\label{eq:v4v9-hodge-energy}
\end{equation}
Equivalently, the variable exact curvature is orthogonal to the
boundary-forced minimizing curvature.
\end{theorem}

\begin{proof}
The affine space is finite dimensional and the nonnegative quadratic form
attains its infimum after quotienting by \(\ker d_1\).  Its Euler equation is
\begin{equation}
 \langle d_1a_{*,\eta},d_1v\rangle=0
 \qquad(v\in C^1_0(Q)).
\label{eq:v4v9-euler}
\end{equation}
Expand the square
\(\|d_1a_{*,\eta}+d_1v\|_2^2\).  The cross term vanishes by
\eqref{eq:v4v9-euler}, proving \eqref{eq:v4v9-hodge-energy}.  If two
minimizers exist, their difference has zero \(d_1\)-image, which proves the
stated uniqueness.
\end{proof}

For a plaquette set \(X\), call the minimizing curvature buffered when
\begin{equation}
 |(d_1a_{*,\eta})_p|\le\frac{r_g}{2}
 \qquad(p\in X).
\label{eq:v4v9-abelian-buffer}
\end{equation}

\begin{corollary}[Exact abelian buffered gap]
\label{cor:v4v9-abelian-gap}
If \eqref{eq:v4v9-abelian-buffer} holds and
\begin{equation}
 |(d_1a)_p|\ge r_g
 \qquad(p\in X),
\label{eq:v4v9-abelian-bad}
\end{equation}
then
\begin{equation}
 \mathcal E_\eta(a)-\mathcal E_\eta(a_{*,\eta})
 \ge\frac{r_g^2}{8}|X|.
\label{eq:v4v9-abelian-gap}
\end{equation}
\end{corollary}

\begin{proof}
Write \(a=a_{*,\eta}+v\).  The reverse triangle inequality and
\eqref{eq:v4v9-abelian-buffer}--%
\eqref{eq:v4v9-abelian-bad} give
\(|(d_1v)_p|\ge r_g/2\) for every \(p\in X\).  Hence
\(\|d_1v\|_2^2\ge r_g^2|X|/4\).  Insert this in
\eqref{eq:v4v9-hodge-energy}.
\end{proof}

This identity exhibits the exact ancestry split.  The component of curvature
in \(\operatorname{Ran}d_1\) is locally variable and paid by the energy gap.
The orthogonal minimizing component is fixed by the exterior data and must
be charged to the boundary ancestor if it is already large on \(X\).

\subsection{Nonabelian curvature in a tree chart}
\label{sec:v4v9-nonabelian-chart}

Fix tree gauge on the same block.  Let
\(a=(a_e)_{e\notin T}\in\mathfrak{su}(N)^{n_Q}\) be logarithmic coordinates
for the remaining links, and let \(\eta\) be boundary data in a fixed small
chart.  Denote by
\begin{equation}
 \mathcal F_\eta(a)
 =\bigl(\log U_p(a,\eta)\bigr)_{p\subset Q}
\label{eq:v4v9-curvature-map}
\end{equation}
the vector of principal logarithms where defined.

\begin{lemma}[BCH curvature and the retained Bianchi word]
\label{lem:v4v9-BCH}
On a sufficiently small fixed chart,
\begin{equation}
 \mathcal F_\eta(a)
 =d_{1,\eta}a+\mathcal Q_\eta(a,a)+\mathcal R_\eta(a),
\label{eq:v4v9-BCH-curvature}
\end{equation}
where \(d_{1,0}=d_1\), \(\mathcal Q_\eta\) is a bounded bilinear commutator
form, and
\begin{equation}
 \|\mathcal R_\eta(a)\|_2\le C_Q\|a\|_2^3.
\label{eq:v4v9-BCH-remainder}
\end{equation}
The product of oriented plaquette holonomies around every elementary cube is
the exact nonabelian Bianchi word; its logarithm has vanishing linear part
and begins with the commutators contained in \(\mathcal Q_\eta\).
\end{lemma}

\begin{proof}
Each plaquette holonomy is a product of four link exponentials and fixed
boundary factors.  The Baker--Campbell--Hausdorff formula on a common compact
chart gives its linear coboundary, its quadratic commutators, and a uniformly
bounded cubic remainder.  There are finitely many links and plaquettes in
the fixed block, so the component bounds combine into
\eqref{eq:v4v9-BCH-remainder}.  Multiplying the six oriented face holonomies
of a cube with the required parallel transports gives the identity exactly:
each edge path cancels with its inverse.  Applying BCH shows that the linear
coboundary terms cancel and leaves the stated commutator word.  Thus no
Bianchi term has been discarded by linearization.
\end{proof}

Define the boundary-dependent Wilson energy in the tree chart by
\begin{equation}
 E_\eta(a)=\sum_{p\subset Q}\ell(U_p(a,\eta)),
\label{eq:v4v9-chart-energy}
\end{equation}
including all plaquettes affected by the internal coordinates.

\begin{lemma}[Positivity of the gauge-fixed Wilson Hessian]
\label{lem:v4v9-hessian-positive}
At trivial boundary data and \(a=0\), the Hessian of
\eqref{eq:v4v9-chart-energy} is positive definite on the non-tree
coordinates.  Hence there are a chart \(\mathcal U_Q\) and constants
\(m_Q,L_Q>0\) such that, for sufficiently small boundary data,
\begin{align}
 D^2E_\eta(a)&\ge m_Q I
 \qquad(a\in\mathcal U_Q),
\label{eq:v4v9-strong-convexity}\\
 \|\mathcal F_\eta(a)-\mathcal F_\eta(b)\|_2
 &\le L_Q\|a-b\|_2
 \qquad(a,b\in\mathcal U_Q).
\label{eq:v4v9-curvature-Lipschitz}
\end{align}
\end{lemma}

\begin{proof}
Lemma~\ref{lem:v4v7-group-comparison} and
Lemma~\ref{lem:v4v9-BCH} identify the Hessian at zero with a positive
constant times \(d_1^*d_1\).  If a non-tree vector lies in its kernel, its
linear curvature vanishes.  On a contractible block it is a pure gauge; tree
gauge and the fixed root make that vector zero.  Thus the Hessian is positive
definite.  Its least eigenvalue remains positive under a sufficiently small
change of \((a,\eta)\) by continuity.  Smoothness of the finite-dimensional
curvature map on the smaller closed chart gives the Lipschitz constant.
\end{proof}

\begin{theorem}[Small-boundary nonabelian buffered gap]
\label{thm:v4v9-nonabelian-gap}
Let \(a_{*,\eta}\) be an interior minimizer of \(E_\eta\) in the strongly
convex chart \(\mathcal U_Q\).  Suppose
\begin{equation}
 d_G(U_p(a_{*,\eta},\eta),1)
 \le\frac{r_g}{2}
 \qquad(p\in X).
\label{eq:v4v9-minimizer-buffer}
\end{equation}
If \(a\in\mathcal U_Q\) and
\begin{equation}
 d_G(U_p(a,\eta),1)\ge r_g
 \qquad(p\in X),
\label{eq:v4v9-chart-bad}
\end{equation}
then
\begin{equation}
 E_\eta(a)-E_\eta(a_{*,\eta})
 \ge c_Qr_g^2|X|
\label{eq:v4v9-nonabelian-gap}
\end{equation}
for a constant \(c_Q>0\) independent of \(g,X\), and the admitted small
boundary data.
\end{theorem}

\begin{proof}
Strong convexity and stationarity of the interior minimizer give
\begin{equation}
 E_\eta(a)-E_\eta(a_{*,\eta})
 \ge\frac{m_Q}{2}\|a-a_{*,\eta}\|_2^2.
\label{eq:v4v9-convex-gap}
\end{equation}
On the common logarithm chart, group distance and logarithm norm are
uniformly equivalent.  The reverse triangle inequality applied to
\eqref{eq:v4v9-minimizer-buffer}--%
\eqref{eq:v4v9-chart-bad} therefore gives
\begin{equation}
 \|\mathcal F_\eta(a)-
 \mathcal F_\eta(a_{*,\eta})\|_2
 \ge c'_Qr_g|X|^{1/2}.
\label{eq:v4v9-curvature-separation}
\end{equation}
By \eqref{eq:v4v9-curvature-Lipschitz},
\(\|a-a_{*,\eta}\|_2\) is at least the left side of
\eqref{eq:v4v9-curvature-separation} divided by \(L_Q\).  Insert this in
\eqref{eq:v4v9-convex-gap}.
\end{proof}

\begin{corollary}[The V8 buffered repair datum in the weak chart]
\label{cor:v4v9-buffered-datum}
Under the hypotheses of Theorem~\ref{thm:v4v9-nonabelian-gap}, the minimizer
\(a_{*,\eta}\) supplies the point in
Definition~\ref{def:v4v8-buffered-datum}.  For sufficiently small \(g\), a
ball of radius \(\rho g^2\) about it remains in the good reference set.
Consequently Theorem~\ref{thm:v4v8-gap-to-repair} gives the conditional
stretched-exponential repair ratio for this component class.
\end{corollary}

\begin{proof}
The energy is Lipschitz on the closed smaller chart.  Curvature is also
Lipschitz there.  Since
\[
 \frac{g^2}{r_g}=g^{1+\alpha}\longrightarrow0,
\]
one may choose a fixed \(\rho>0\) so that the curvature change on the
\(\rho g^2\)-ball is less than the buffer between
\(r_g/2\) and \(r_g\).  Equation \eqref{eq:v4v9-nonabelian-gap} is the
variational gap \eqref{eq:v4v8-variational-gap}.  Apply
Theorem~\ref{thm:v4v8-gap-to-repair}.
\end{proof}

\subsection{The exact local dichotomy}
\label{sec:v4v9-dichotomy}

\begin{proposition}[Repairable fluctuation or inherited curvature]
\label{prop:v4v9-dichotomy}
For a component \(X\) contained in the weak tree chart, exactly one of the
following disjoint cases holds:
\begin{enumerate}
\item the minimizing curvature satisfies
\(d_G(U_p(a_{*,\eta},\eta),1)\le r_g/2\) for every \(p\in X\), in which
case Theorem~\ref{thm:v4v9-nonabelian-gap} pays the component;
\item at least one \(p\in X\) has minimizing curvature greater than
\(r_g/2\), in which case that plaquette is boundary-forced and is assigned
to the exterior ancestor.
\end{enumerate}
\end{proposition}

\begin{proof}
The two alternatives are complementary and disjoint.  The first has the
proved repair gap.  In the second, the energy-minimizing interior already
contains threshold-scale curvature.  Removing it while fixing the exterior
is not a local fluctuation comparison, so its unique source is the
boundary-compatible minimizing sector.
\end{proof}

\begin{firewall}[The dichotomy is local to the weak chart]
\label{fw:v4v9-chart-scope}
Proposition~\ref{prop:v4v9-dichotomy} does not control boundary data outside
the tree logarithm chart, components containing plaquettes beyond the
injectivity radius, noncontractible toron sectors, or transitions between
topological bundles.  Strong convexity can fail there.  Those sectors have
not yet been bounded by a coarser ancestor measure.
\end{firewall}

\subsection{V9 result ledger and V10 audit target}
\label{sec:v4v9-conclusion}

V9 proves the exact relative Hodge energy identity and buffered repair gap
for the abelian quadratic problem.  It retains the nonabelian Bianchi word
through the BCH curvature map, proves positivity of the tree-gauge Wilson
Hessian on a fixed contractible block, and obtains a uniform strong-convexity
gap for small boundary data.  That gap supplies the V8 reference ball and
therefore the rooted stretched-exponential defect estimate for every weak-
chart component whose minimizing curvature is buffered.

The unresolved components now have a precise ancestry signature: large
minimizing curvature, exit from the logarithm chart, toron holonomy, or
topological bundle change.  V10 is the next compulsory NS/YM audit.  After
that audit it should construct a coarser boundary-curvature variable and
test whether the inherited components telescope once across block scales,
without charging their Wilson energy both locally and to the ancestor.


\clearpage
\section{V10: companion audit and the one-use boundary-energy telescope}
\label{sec:v4v10}

V9 proved local repair when the boundary-compatible minimizer lies in the
weak chart and is buffered on the bad component.  The remaining case is
inherited curvature: enlarging the variable block can release curvature
which a smaller boundary condition forced.  This compulsory audit uses the
current Yang--Mills and Navier--Stokes notes to determine how that release may
be counted.  The answer is an exact nested-minimum telescope.  Its separate
drops are not independent payers; their complete sum is one boundary-
ancestry carrier.

For an abelian quadratic energy, the telescope refines to an orthogonal
relative-Hodge identity.  For a general compact Wilson energy, monotonicity
of nested minima is enough.  A first repairable stopping scale then gives a
single stretched-exponential payment.  Simultaneous components can share
ancestor drops, so a disjoint ancestry forest is still required before this
single-component result can be inserted into the full polymer gas.

\subsection{Compulsory V10 audit}
\label{sec:v4v10-audit}

The exact companion files read were
\begin{enumerate}
\item \texttt{Renewal\_Yang\_Mills\_Mass\_Gap\_Research\_Notes\_V30.tex},
\(579070\) bytes, with SHA--256
\begin{center}\scriptsize\ttfamily
F2874016202D63EEBF987C6E878598A495057EE3938EBABF83A0C82A664F25CB.
\end{center}
\item \texttt{Renewal\_NS\_Stability\_Research\_Notes\_V31.tex},
\(887537\) bytes, with SHA--256
\begin{center}\scriptsize\ttfamily
F1121B62238AD5491BB7CF03635EA9934942EDF573ED8FAAA9EE79E1D38A6696.
\end{center}
\end{enumerate}

Yang--Mills V30 closes its clean discrete-prefix carrier by keeping five
local letters inside one third cumulant and retaining the physical suffix
which supplies the third tree edge.  Its composite symbol is an assembly
device; it does not create that edge.  The SM transfer is exact: the sum of
nested boundary-energy drops may be assembled into one carrier, but the
coarsest minimizing sector must remain present.  Naming an ancestor cannot
manufacture a second curvature payment.

Navier--Stokes V31 still shows that a flat probability mark does not pay its
critical first moment and that a persistent interval may not be charged
again on the recent-entry branch.  Here, a fine bad component may persist
through many block boundaries.  The number of levels is not summable stock.
Only its first repairable stopping scale, or one complete telescoping energy
drop, may be used.

\begin{proposition}[Audit allocation rule]
\label{prop:v4v10-audit-rule}
A boundary-inherited bad component has one ancestry chain.  Along that
chain:
\begin{enumerate}
\item nested minimum drops are assembled before any one is estimated;
\item the component is stopped at its first repairable level;
\item its terminal boundary/topological suffix is retained if no such level
exists; and
\item no persistence level is counted as a new bad-plaquette occurrence.
\end{enumerate}
\end{proposition}

\subsection{Nested Wilson minima}
\label{sec:v4v10-nested-minima}

Fix a finite top block \(Q_L\) and an original link field \(a\).  Let
\[
 Q_0\subset Q_1\subset\cdots\subset Q_L
\]
be its deterministic block-ancestor chain.  Let \(\mathcal A_\ell(a)\) be
the set of fields obtained by varying links in \(Q_\ell\) while retaining
the original exterior links, with gauge-equivalent repetitions quotiented.
Then
\begin{equation}
 \mathcal A_0(a)\subset\mathcal A_1(a)
 \subset\cdots\subset\mathcal A_L(a).
\label{eq:v4v10-nested-sets}
\end{equation}
For the unscaled Wilson energy \(E\) on \(Q_L\), put
\begin{equation}
 m_\ell(a)=\min_{b\in\mathcal A_\ell(a)}E(b),
 \qquad
 \Delta_\ell(a)=m_{\ell-1}(a)-m_\ell(a)
 \quad(1\le\ell\le L).
\label{eq:v4v10-minimum-drops}
\end{equation}

\begin{theorem}[Exact one-use boundary-energy telescope]
\label{thm:v4v10-energy-telescope}
The minima and drops satisfy
\begin{align}
 m_0&\ge m_1\ge\cdots\ge m_L,
\label{eq:v4v10-minimum-monotone}\\
 \Delta_\ell&\ge0,
\label{eq:v4v10-drop-positive}\\
 E(a)-m_k(a)
 &=E(a)-m_0(a)+\sum_{\ell=1}^k\Delta_\ell(a)
 \qquad(0\le k\le L).
\label{eq:v4v10-exact-telescope}
\end{align}
Every energy occurrence on the right belongs to one algebraic difference;
there is no level multiplicity.
\end{theorem}

\begin{proof}
Compactness of the finite link space and continuity give each minimum.
The inclusions \eqref{eq:v4v10-nested-sets} imply
\(m_\ell\le m_{\ell-1}\), proving the first two statements.  Summing the
definitions of \(\Delta_\ell\) gives
\(\sum_{\ell=1}^k\Delta_\ell=m_0-m_k\), which is exactly
\eqref{eq:v4v10-exact-telescope}.
\end{proof}

\begin{firewall}[No dominant-level pigeonhole]
\label{fw:v4v10-no-pigeonhole}
Equation \eqref{eq:v4v10-exact-telescope} does not imply that one
\(\Delta_\ell\) carries a level-independent fraction of the total.  A
pigeonhole choice would lose \(1/L\), reproducing the forbidden scale mark.
The complete sum is the boundary carrier and must be kept assembled, as in
the Yang--Mills V30 cumulant.
\end{firewall}

\subsection{The abelian orthogonal refinement}
\label{sec:v4v10-orthogonal}

For the quadratic relative problem of V9, use a fixed top-block curvature
\(f\).  Let
\begin{equation}
 \mathcal W_\ell
 =d_1C^1_0(Q_\ell)\subset C^2(Q_L)
\label{eq:v4v10-exact-spaces}
\end{equation}
with cochains extended by zero outside \(Q_\ell\), and let \(P_\ell\) be the
orthogonal projection onto \(\mathcal W_\ell\).  The spaces are nested.

\begin{theorem}[Orthogonal ancestry decomposition]
\label{thm:v4v10-orthogonal-ancestry}
For \(h_\ell=(1-P_\ell)f\),
\begin{align}
 m_\ell&=\frac12\|h_\ell\|_2^2,
\label{eq:v4v10-abelian-minimum}\\
 m_{\ell-1}-m_\ell
 &=\frac12\|(P_\ell-P_{\ell-1})f\|_2^2,
\label{eq:v4v10-orthogonal-drop}\\
 \|P_kf\|_2^2
 &=\|P_0f\|_2^2
 +\sum_{\ell=1}^k
 \|(P_\ell-P_{\ell-1})f\|_2^2.
\label{eq:v4v10-orthogonal-sum}
\end{align}
Thus each released exact-curvature layer is orthogonal to every other layer.
\end{theorem}

\begin{proof}
Minimizing \(\|f+w\|_2^2/2\) over \(w\in\mathcal W_\ell\) subtracts
\(P_\ell f\), leaving \(h_\ell\), which proves
\eqref{eq:v4v10-abelian-minimum}.  For nested orthogonal projections,
\(P_\ell-P_{\ell-1}\) is the projection onto
\(\mathcal W_\ell\cap\mathcal W_{\ell-1}^{\perp}\).  Therefore
\[
 P_\ell f=P_{\ell-1}f+(P_\ell-P_{\ell-1})f
\]
is an orthogonal sum.  Pythagoras gives
\eqref{eq:v4v10-orthogonal-drop} and iteration gives
\eqref{eq:v4v10-orthogonal-sum}.
\end{proof}

The nonabelian Wilson telescope in
Theorem~\ref{thm:v4v10-energy-telescope} is not orthogonal, but it retains
the two properties required for ancestry: every drop is nonnegative and the
complete sum has no repetition.

\subsection{First repairable stopping scale}
\label{sec:v4v10-stopping}

Fix a bad component \(X\subset Q_0\).  For each level choose a minimizer
\(a^*_\ell\in\mathcal A_\ell(a)\).  Call level \(\ell\) admissibly buffered
when \(a^*_\ell\) is an interior minimizer, the restrictions of \(a\) and
\(a^*_\ell\) lie in a common weak tree chart, and
\begin{equation}
 d_G(U_p(a^*_\ell),1)\le r_g/2
 \qquad(p\in X).
\label{eq:v4v10-level-buffer}
\end{equation}
Define
\begin{equation}
 \tau(X,a)=\min\{\ell:\ell\text{ is admissibly buffered}\},
\label{eq:v4v10-stopping-time}
\end{equation}
with \(\tau=\infty\) if the set is empty.

\begin{proposition}[Exact stopping partition]
\label{prop:v4v10-stopping-partition}
For every component,
\begin{equation}
 1=\sum_{\ell=0}^L1_{\{\tau=\ell\}}
 +1_{\{\tau=\infty\}}.
\label{eq:v4v10-stopping-partition}
\end{equation}
The events are disjoint, and a component is assigned to one finite repair
carrier or one terminal suffix.
\end{proposition}

\begin{proof}
This is the partition by the first element of a finite ordered set, together
with the empty-set case.  Minimality makes the finite events disjoint.
\end{proof}

\begin{theorem}[One-component multiscale repair]
\label{thm:v4v10-multiscale-repair}
Suppose \(a\) is bad on every plaquette of \(X\) and
\(\tau(X,a)=k<\infty\).  Under the uniform weak-chart constants of V9,
\begin{equation}
 \mathcal C_k(a)
 :=E(a)-m_k(a)
 =E(a)-m_0(a)+\sum_{\ell=1}^k\Delta_\ell(a)
 \ge c_Qr_g^2|X|.
\label{eq:v4v10-complete-carrier}
\end{equation}
Consequently the complete Boltzmann factor associated with this ancestry
chain contains
\begin{equation}
 \exp[-g^{-2}\mathcal C_k(a)]
 \le\exp[-c_Qg^{-2\alpha}|X|].
\label{eq:v4v10-stopped-suppression}
\end{equation}
\end{theorem}

\begin{proof}
At the stopping level, \(a^*_k\) is an interior minimizer in the common weak
chart and is buffered on \(X\).  Theorem~\ref{thm:v4v9-nonabelian-gap}
gives
\(E(a)-E(a^*_k)\ge c_Qr_g^2|X|\).  Since
\(E(a^*_k)=m_k(a)\), the left side is \(\mathcal C_k(a)\).  Its exact
telescope is Theorem~\ref{thm:v4v10-energy-telescope}.  Divide by \(g^2\)
and use \(r_g^2/g^2=g^{-2\alpha}\).
\end{proof}

\begin{firewall}[The carrier is one sum, not \(k+1\) payments]
\label{fw:v4v10-one-carrier}
The factors in \eqref{eq:v4v10-complete-carrier} are an exact decomposition
of \(E(a)-m_k\).  They may be used together in the exponential, but may not
each be attached as an independent plaquette suppression to the same
component.  The terminal weight \(e^{-m_k/g^2}\) also remains in the
conditional reference integral.
\end{firewall}

\subsection{Collision of several ancestry chains}
\label{sec:v4v10-collisions}

For one rooted component, Theorem~\ref{thm:v4v10-multiscale-repair} supplies
the required suppression.  A connected polymer can contain several bad
components.  Their fine plaquette sets are disjoint, but their block
ancestors can merge.  After that merger, the same minimum drop
\(\Delta_\ell\) can occur in more than one formal chain.

\begin{proposition}[Disjoint forest verification criterion]
\label{prop:v4v10-forest-criterion}
Suppose a multi-component configuration has a common nonnegative master
ledger \(\{\delta_\gamma\}_{\gamma\in\mathcal J}\) satisfying
\begin{equation}
 \sum_{\gamma\in\mathcal J}\delta_\gamma
 \le E(a)-m_L(a).
\label{eq:v4v10-master-ledger}
\end{equation}
Suppose each component carrier is the sum of a subset of these entries, a
deterministic ancestry forest assigns every entry to at most one carrier,
and after chains merge their unspent suffixes are replaced by one common
parent suffix.  Then the total assigned carrier is at most
\(E(a)-m_L(a)\), and no Wilson-energy occurrence is duplicated.
\end{proposition}

\begin{proof}
The assignment subsets are disjoint by hypothesis, including after every
merger.  Their total is therefore a sub-sum of the nonnegative master ledger
and is bounded by \eqref{eq:v4v10-master-ledger}.
\end{proof}

\begin{firewall}[The full ancestry forest has not been constructed]
\label{fw:v4v10-forest-open}
Proposition~\ref{prop:v4v10-forest-criterion} is a verification criterion.
It does not construct the common master ledger or a forest compatible
simultaneously with every bad-component geometry, topological sector,
background derivative mark, and polymer incompatibility relation.  Without such a construction, multiplying
the one-component bounds can double count a shared coarse minimum drop.
\end{firewall}

\subsection{V10 result ledger and V11 target}
\label{sec:v4v10-conclusion}

V10 records the required companion audit and proves the exact nested-minimum
energy telescope.  Its abelian specialization is an orthogonal relative-
Hodge decomposition.  The first repairable stopping rule partitions every
single bad component without scale multiplicity, and the complete carrier
at a finite stopping level has the stretched-exponential Wilson cost needed
by V7.  Persistent components retain their boundary/topological suffix.

The remaining combinatorial gate is now precise: several component chains
can merge and share a coarse energy drop.  V11 should construct the
deterministic ancestry forest by ordering block mergers, keeping one common
suffix after each merger, and retaining background-response marks only after
the joint carrier is assembled.  That is the exact analogue of the
Yang--Mills V30 composite-ternary rule and is required before claiming the
full derivative-marked defect KP estimate.


\clearpage
\section{V11: the collar-merger ancestry forest and joint defect payment}
\label{sec:v4v11}

V10 proved a one-component stopping telescope and warned that several
components can share a coarse energy release.  This note constructs a
deterministic collision rule before any repair estimate is multiplied.
Every active packet carries its union of bad plaquettes and a block in a
fixed finite hierarchy.  Nonrepairable packets move to their parent; packets
whose interaction collars meet are merged at their least common ancestor,
including packets which were provisionally repairable.  The process stops
with pairwise disjoint collars.  Repair minimizers on those blocks may then
be applied simultaneously, so their energy drops form a genuine common
master ledger.

This closes the multicomponent normalization for the sector in which every
final packet is repairable.  A packet reaching a root without a weak-chart
buffer remains a terminal boundary/topological sector.  Its weight has not
been bounded here.

\subsection{Blocks, collars, and repairable packets}
\label{sec:v4v11-blocks}

Let \(\mathscr T\) be a finite rooted hierarchy of cubical blocks.  Every
nonroot block \(B\) has a unique parent \(\operatorname{par}B\), and two
blocks have a least common ancestor \(B\vee B'\).  For a block \(B\), let
\begin{equation}
 \operatorname{St}(B)
 =\{p:p\text{ contains a link variable varied inside }B\}
\label{eq:v4v11-star}
\end{equation}
be its plaquette interaction collar.  Thus changing links in \(B\) changes
no Wilson term outside \(\operatorname{St}(B)\).

An active packet is a pair \((B,X)\), where \(X\) is a nonempty union of
original bad plaquette components assigned to \(B\).  Distinct active
packets always carry disjoint sets \(X\).

\begin{definition}[Repairable packet]
\label{def:v4v11-repairable}
A packet \((B,X)\) is repairable when the Wilson minimizer obtained by
varying the links of \(B\) with the current exterior fixed:
\begin{enumerate}
\item is an interior minimizer in the V9 weak tree chart;
\item has plaquette distance at most \(r_g/2\) on every \(p\in X\);
\item has a uniform reference ball of radius \(\rho g^2\); and
\item satisfies the V9 strong-convexity and curvature-Lipschitz constants
uniformly for the packet class.
\end{enumerate}
\end{definition}

By Theorem~\ref{thm:v4v9-nonabelian-gap}, a repairable packet has a local
energy drop at least \(c_Qr_g^2|X|\).

\subsection{The deterministic collar-merger algorithm}
\label{sec:v4v11-algorithm}

Initialize one packet for each fine bad component, using the smallest
hierarchy block containing that component.  Repeatedly perform the first
applicable operation in the following fixed order.
\begin{enumerate}
\item If two active packets \((B,X)\), \((B',X')\) have intersecting
collars, replace them by
\begin{equation}
 (B\vee B',X\sqcup X').
\label{eq:v4v11-merge}
\end{equation}
\item If an active nonroot packet is not repairable, replace
\((B,X)\) by \((\operatorname{par}B,X)\).
\label{item:v4v11-promote}
\item If neither operation applies, stop.  A repairable packet is declared
paid; a nonrepairable root packet is declared terminal.
\end{enumerate}
Ties are resolved by a fixed lexicographic order on hierarchy blocks, so the
algorithm is deterministic.

\begin{theorem}[Termination and exact component partition]
\label{thm:v4v11-termination}
The collar-merger algorithm terminates after finitely many operations.  Its
final packet sets \(X_\nu\) are pairwise disjoint and satisfy
\begin{equation}
 B_{\rm original}=\bigsqcup_\nu X_\nu.
\label{eq:v4v11-bad-partition}
\end{equation}
Their collars \(\operatorname{St}(B_\nu)\) are pairwise disjoint.  Every
final packet is either repairable or a nonrepairable root packet.
\end{theorem}

\begin{proof}
Assign to an active family the ordered pair consisting of its number of
packets and the sum of their distances to a root.  A merger strictly lowers
the first entry.  A promotion leaves the first entry fixed and strictly
lowers the second.  Lexicographic order on nonnegative integer pairs is
well-founded, so the process terminates.

Both operations replace disjoint carried sets by their disjoint union or
leave the carried set unchanged.  Hence the invariant
\eqref{eq:v4v11-bad-partition} holds from initialization to termination.  If
two final collars intersected, the merger operation would still apply.
If a final nonroot packet were nonrepairable, the promotion operation would
still apply.  This proves all claims.
\end{proof}

\begin{firewall}[Provisional repair never survives a later collision]
\label{fw:v4v11-provisional}
A repairable packet is not paid while another active packet may grow into
its collar.  If that occurs, the first operation merges them and discards
both provisional minimizers.  Only the minimizer of the merged packet may
later be used.  Thus no child repair energy is retained beneath a common
parent repair.
\end{firewall}

\subsection{Simultaneous repair on disjoint collars}
\label{sec:v4v11-simultaneous}

Let \(\mathcal P\) be the set of final repairable packets.  For
\(\nu\in\mathcal P\), let \(a^*_{\nu}\) be its boundary-compatible minimizer
and define the local drop
\begin{equation}
 \delta_\nu
 =E_{\operatorname{St}(B_\nu)}(a)
 -E_{\operatorname{St}(B_\nu)}(a^*_{\nu}),
\label{eq:v4v11-local-drop}
\end{equation}
where the local energy contains each plaquette in that collar once.

\begin{theorem}[Disjoint-collar master ledger]
\label{thm:v4v11-master-ledger}
The repair moves for all \(\nu\in\mathcal P\) commute and define a
simultaneously repaired field \(a^{\rm rep}\).  Moreover,
\begin{align}
 E(a)-E(a^{\rm rep})
 &=\sum_{\nu\in\mathcal P}\delta_\nu,
\label{eq:v4v11-energy-additivity}\\
 \delta_\nu&\ge c_Qr_g^2|X_\nu|,
\label{eq:v4v11-packet-gap}\\
 E(a)-E(a^{\rm rep})
 &\ge c_Qr_g^2
 \sum_{\nu\in\mathcal P}|X_\nu|.
\label{eq:v4v11-total-gap}
\end{align}
Thus \(\{\delta_\nu\}\) is a common nonnegative master ledger satisfying
the V10 disjoint forest criterion.
\end{theorem}

\begin{proof}
Pairwise disjoint collars imply that the varied link sets are disjoint and
that no Wilson plaquette term is affected by two repair moves.  Hence the
moves commute.  Every energy term outside their collars is unchanged, and
the changed terms split into the disjoint collar sums.  This proves
\eqref{eq:v4v11-energy-additivity}.  Each final paid packet is repairable, so
Theorem~\ref{thm:v4v9-nonabelian-gap} gives
\eqref{eq:v4v11-packet-gap}.  Sum over \(\nu\) and use additivity.
\end{proof}

\begin{corollary}[Joint reference volume and normalized suppression]
\label{cor:v4v11-joint-ratio}
Assume each repairable block has at most \(C_n|X_\nu|\) internal link
coordinates after gauge quotient.  The product of its V9 reference balls
gives a simultaneous reference set, and the ratio of the paid bad sector to
that reference sector is bounded by
\begin{equation}
 \exp\!\left[-
 \bigl(c_Qg^{-2\alpha}-C\log(g^{-1})\bigr)
 \sum_{\nu\in\mathcal P}|X_\nu|
 \right].
\label{eq:v4v11-joint-ratio}
\end{equation}
For small \(g\), this is at most
\begin{equation}
 \exp\!\left[-\frac{c_Q}{2}g^{-2\alpha}
 \sum_{\nu\in\mathcal P}|X_\nu|\right].
\label{eq:v4v11-joint-stretched}
\end{equation}
\end{corollary}

\begin{proof}
Disjoint collars make the reference balls a Cartesian product.  Haar-ball
volume gives a factor at worst
\(g^{C\sum_\nu|X_\nu|}\), while
Theorem~\ref{thm:v4v11-master-ledger} gives the Wilson exponential in
\eqref{eq:v4v11-total-gap}.  Their ratio is
\eqref{eq:v4v11-joint-ratio}.  Since
\(g^{-2\alpha}/\log(g^{-1})\to\infty\), reduce \(g\) to obtain
\eqref{eq:v4v11-joint-stretched}.
\end{proof}

\subsection{The eventually repairable defect sector}
\label{sec:v4v11-eventual}

Call a bad configuration controlled eventually repairable when the final
forest has no terminal packet and
\begin{equation}
 \sum_{\nu\in\mathcal P}n(B_\nu)
 \le C_n\sum_{\nu\in\mathcal P}|X_\nu|,
\label{eq:v4v11-controlled-volume}
\end{equation}
where \(n(B_\nu)\) is its number of internal link coordinates.  Then the
final sets \(X_\nu\) partition every bad plaquette and their joint reference
volume has only a per-defect polynomial cost.

\begin{theorem}[Multicomponent defect KP bound in the controlled repairable sector]
\label{thm:v4v11-multicomponent-KP}
Suppose the good-field polymer attachments satisfy the V6 bound and at most
two coupling-normalized response derivatives cost a polynomial in
\(g^{-1}\) and total packet size.  On the controlled eventually repairable
sector, the complete bad-component expansion satisfies the rooted marked KP
estimate \eqref{eq:v4v7-defect-KP}, uniformly in shell and volume, for
sufficiently small \(g\).
\end{theorem}

\begin{proof}
The collar-merger output replaces every colliding group of components by one
packet before estimation.  Theorem~\ref{thm:v4v11-master-ledger} and
Corollary~\ref{cor:v4v11-joint-ratio} give a stretched-exponential weight in
the total number of bad plaquettes without duplicated energy.  The forest is a deterministic function of the fine bad set, so promotions
and mergers introduce no choice multiplicity.  Rooted fine defect supports
together with their good-field attachments have the usual exponential
connected-set count.  The controlled-volume condition pays the product
reference ball per bad plaquette, and polynomial derivative marks are
absorbed by \(e^{-c g^{-2\alpha}n/2}\).  The connected-set summation used in
Theorem~\ref{thm:v4v7-repair-to-KP} then converges with arbitrarily small KP
constant for small \(g\).  Good-field attachments are already summable by
the V6 hypothesis.
\end{proof}

\begin{firewall}[Terminal packets are not paid by the repairable sector]
\label{fw:v4v11-terminal}
If a nonrepairable packet reaches a root, none of its bad plaquettes appears
in the exponent \eqref{eq:v4v11-joint-stretched}.  Assigning them the same
Wilson payment would reuse energy already retained in the terminal boundary
weight.  Large boundary holonomy, torons, cut-locus configurations, and
topological bundle changes remain in this terminal sector.  A repairable
packet whose block volume is disproportionate to its bad set is a separate
sparse-promotion sector and is also excluded from
Theorem~\ref{thm:v4v11-multicomponent-KP}.
\end{firewall}

\subsection{V11 result ledger and V12 target}
\label{sec:v4v11-conclusion}

V11 constructs a deterministic ancestry forest for colliding bad
components.  Promotion and least-common-ancestor merger terminate, preserve
an exact partition of the original bad plaquettes, and leave pairwise
disjoint final interaction collars.  Only then are repair minimizers
inserted.  Their moves commute, their Wilson drops add exactly, and their
reference Haar volumes multiply.  This proves the multicomponent derivative-marked defect KP bound on the
controlled eventually repairable sector.

The remaining gauge large-field gates are the terminal forest and sparse-
promotion sectors.
V12 should classify its roots by gauge-invariant boundary holonomy: small
contractible holonomy with a cut-locus artifact, nontrivial toron class, or
quantized topological charge.  The first class should be returned to a
larger logarithm chart; the latter two require explicit finite-dimensional
parent variables rather than another local repair estimate.


\clearpage
\section{V12: flat toron parents and covariant shell spectra}
\label{sec:v4v12}

V11 left terminal packets when no weak-chart buffered minimizer appears or
when promotion volume is disproportionate to the fine defect.  One source
of apparent failure is not a large-field defect at all.  On a periodic
lattice, a flat connection can have nontrivial cycle holonomy and Wilson
action exactly zero.  No plaquette energy can suppress that holonomy, and a
periodic gauge transformation need not move it into the identity link
chart.  This note classifies the flat sector exactly and promotes it into a
finite-dimensional toron parent.

The toron changes the fluctuation spectrum by shifting adjoint momenta.  An
ordinary Fourier shell therefore does not have uniform spectral ancestry
over the parent.  The correct shell is a spectral interval of the covariant
Laplacian conditional on the toron.  Resonance strata, where its centralizer
enlarges, carry additional zero modes and must also remain in the parent.

\subsection{Flat periodic lattice connections}
\label{sec:v4v12-flat-classification}

Let \(\Lambda_M=(\mathbb Z/M\mathbb Z)^d\) and let \(G=SU(N)\).  A link
field is flat when
\begin{equation}
 U_p=1
 \qquad\text{for every plaquette }p.
\label{eq:v4v12-flatness}
\end{equation}
For each coordinate direction, let \(H_\mu\) be the holonomy around the
corresponding fundamental cycle.

\begin{theorem}[Classification of flat torons]
\label{thm:v4v12-flat-classification}
Gauge-equivalence classes of flat periodic \(SU(N)\) link fields are in
bijection with
\begin{equation}
 \mathcal M_{\rm flat}(T^d,SU(N))
 =\{(H_1,\ldots,H_d)\in SU(N)^d:
 [H_\mu,H_\nu]=1\}/SU(N),
\label{eq:v4v12-flat-moduli}
\end{equation}
where the quotient is simultaneous conjugation.  Equivalently,
\begin{equation}
 \mathcal M_{\rm flat}(T^d,SU(N))\cong T^{,d}/W,
\label{eq:v4v12-torus-Weyl}
\end{equation}
with \(T\) a maximal torus and \(W\) the Weyl group acting diagonally.
\end{theorem}

\begin{proof}
Cut the periodic lattice along one coordinate hyperplane in every direction
and choose a spanning tree in the resulting contractible cell.  Tree gauge
makes all tree links one.  Flatness successively forces every remaining
contractible loop holonomy to be one, leaving only one cycle holonomy
\(H_\mu\) per periodic direction.  A plaquette spanning directions
\(\mu,\nu\) across the cuts gives
\(H_\mu H_\nu H_\mu^{-1}H_\nu^{-1}=1\), so the tuple commutes.  The residual
root gauge acts by simultaneous conjugation.

Conversely, a commuting tuple defines a flat connection by taking constant
links whose \(M\)-fold products are the \(H_\mu\); commuting unitary matrices
admit commuting roots in a common maximal torus.  Two such fields are gauge
equivalent exactly when their tuples are simultaneously conjugate.  Finally,
every commuting family of unitary matrices is simultaneously diagonalizable.
Intersection with \(SU(N)\) puts it in a maximal torus, and the remaining
conjugations are the Weyl action.  This proves both descriptions.
\end{proof}

\begin{corollary}[No Wilson payment along the flat moduli]
\label{cor:v4v12-no-Wilson-payment}
For every \(\tau\in\mathcal M_{\rm flat}\),
\begin{equation}
 S_g(U^\tau)=0.
\label{eq:v4v12-zero-action}
\end{equation}
Consequently no estimate of the form
\(S_g(U^\tau)\ge c\,d(\tau,	au_0)^2/g^2\) with \(c>0\) can hold on the
toron moduli.
\end{corollary}

\begin{proof}
Flatness makes every term \(\ell(U_p)\) in
\eqref{eq:v4v7-wilson-action} vanish.  Choosing two distinct toron classes
then contradicts any proposed positive lower bound in their moduli distance.
\end{proof}

\begin{firewall}[A toron is not a large plaquette]
\label{fw:v4v12-toron-not-defect}
A tree-gauge representative of a nontrivial toron can contain links far from
the identity even though every plaquette is exactly good.  Promoting such a
representative through the V11 defect forest and charging its large link
coordinate would violate gauge invariance and Corollary~\ref{cor:v4v12-no-Wilson-payment}.
The toron class is a zero-mode parent coordinate.
\end{firewall}

\subsection{Exact adjoint spectrum in a toron background}
\label{sec:v4v12-spectrum}

Choose commuting logarithms \(\Theta_\mu\in\mathfrak t\) with
\(H_\mu=e^{i\Theta_\mu}\), and use the constant representative
\begin{equation}
 U_\mu=e^{i\Theta_\mu/M}.
\label{eq:v4v12-constant-links}
\end{equation}
On adjoint-valued site fields, define
\begin{align}
 (D_\mu^\tau f)(n)
 &=a^{-1}\bigl(U_\mu f(n+e_\mu)U_\mu^{-1}-f(n)\bigr),
\label{eq:v4v12-covariant-difference}\\
 \Delta_\tau&=\sum_\mu(D_\mu^\tau)^*D_\mu^\tau.
\label{eq:v4v12-covariant-Laplacian}
\end{align}

\begin{theorem}[Toron-shifted lattice spectrum]
\label{thm:v4v12-shifted-spectrum}
Let \(E_\gamma\) be a root vector with
\([\Theta,E_\gamma]=\gamma(\Theta)E_\gamma\).  For
\(k\in\{0,\ldots,M-1\}^d\), the mode
\begin{equation}
 f_{k,\gamma}(n)=e^{2\pi i k\cdot n/M}E_\gamma
\label{eq:v4v12-root-mode}
\end{equation}
has eigenvalue
\begin{equation}
 \lambda_{k,\gamma}(\tau)
 =\frac4{a^2}\sum_{\mu=1}^d
 \sin^2\!\left(
 \frac{2\pi k_\mu+\gamma(\Theta_\mu)}{2M}
 \right).
\label{eq:v4v12-shifted-eigenvalue}
\end{equation}
For Cartan modes, take \(\gamma=0\).  The formula includes every adjoint
eigenmode.
\end{theorem}

\begin{proof}
Conjugation by \(U_\mu\) multiplies \(E_\gamma\) by
\(e^{i\gamma(\Theta_\mu)/M}\), while translation multiplies the Fourier mode
by \(e^{2\pi ik_\mu/M}\).  Thus \(D_\mu^\tau\) has eigenvalue
\[
 a^{-1}\left(
 e^{i(2\pi k_\mu+\gamma(\Theta_\mu))/M}-1
 \right).
\]
Its squared modulus is the corresponding summand in
\eqref{eq:v4v12-shifted-eigenvalue}.  The Cartan and root decomposition of
\(\mathfrak{su}(N)_\mathbb C\) is complete.
\end{proof}

\begin{corollary}[Resonance and centralizer]
\label{cor:v4v12-resonance}
A root zero mode exists exactly when
\begin{equation}
 \gamma(\Theta_\mu)\in2\pi\mathbb Z
 \qquad\text{for every }\mu.
\label{eq:v4v12-root-resonance}
\end{equation}
These roots generate the non-Cartan part of the common centralizer of the
toron tuple.  Away from an \(\varepsilon\)-neighborhood of all resonance
walls, every nonzero root mode has a positive lower spectral bound depending
only on \(a,M,\varepsilon\).
\end{corollary}

\begin{proof}
Every term in \eqref{eq:v4v12-shifted-eigenvalue} is nonnegative.  Their sum
vanishes exactly when each angle is an integer multiple of \(2\pi\), which
is \eqref{eq:v4v12-root-resonance} after choosing the appropriate Fourier
integer.  A root vector commutes with all \(H_\mu\) under the same condition.
On the compact complement of the resonance neighborhoods, the continuous
minimum of the finite set of relevant eigenvalue functions is positive.
\end{proof}

\subsection{Covariant spectral shells}
\label{sec:v4v12-covariant-shells}

Let \(0<K_-<K_+\).  Define the toron-conditional shell projection by
\begin{equation}
 P_{[K_-,K_+]}(\tau)
 =1_{[K_-^2,K_+^2]}(\Delta_\tau).
\label{eq:v4v12-spectral-shell}
\end{equation}

\begin{proposition}[Uniform band gap by construction]
\label{prop:v4v12-band-gap}
On \(\operatorname{Ran}P_{[K_-,K_+]}(\tau)\),
\begin{equation}
 K_-^2\|f\|_2^2
 \le\langle f,\Delta_\tau f\rangle
 \le K_+^2\|f\|_2^2
\label{eq:v4v12-band-gap}
\end{equation}
for every toron \(\tau\).  Hence the shell covariance
\(P_{[K_-,K_+]}(\tau)\Delta_\tau^{-1}\) has norm at most \(K_-^{-2}\),
uniformly over the flat parent.
\end{proposition}

\begin{proof}
This is the spectral theorem applied to the finite self-adjoint matrix
\(\Delta_\tau\).  Every eigenvalue selected by
\eqref{eq:v4v12-spectral-shell} lies in the stated interval, and inversion
on that range gives the covariance bound.
\end{proof}

An ordinary Fourier annulus selects \(k\) independently of \(\tau\).
Equation \eqref{eq:v4v12-shifted-eigenvalue} shows that its covariant
eigenvalue can drift as the toron changes.  The covariant spectral projection
is therefore the correct conditional shell variable.  Its rank can jump at
band crossings; those finite-dimensional crossing modes must be transferred
between adjacent shells by exact spectral projections, not duplicated.

\subsection{Toron parent and centralizer strata}
\label{sec:v4v12-parent}

The resonance walls \eqref{eq:v4v12-root-resonance} form a finite Weyl-
invariant arrangement of subtori in \(T^d\).  Stratify
\(\mathcal M_{\rm flat}\) by the set of roots which vanish on the tuple.
On each stratum the common centralizer type and the number of zero modes are
constant.

\begin{proposition}[Finite centralizer stratification]
\label{prop:v4v12-stratification}
For \(SU(N)\), only finitely many centralizer types occur.  The toron moduli
is a finite disjoint union of smooth orbifold strata on which the zero-mode
rank of \(\Delta_\tau\) is constant.
\end{proposition}

\begin{proof}
The root system of \(SU(N)\) is finite.  A centralizer is determined by the
subset of roots satisfying \eqref{eq:v4v12-root-resonance}; only finitely
many subsets occur.  On the locus where exactly a fixed subset vanishes, the
rank conditions are constant and the quotient by the finite Weyl stabilizer
is an orbifold stratum.  Corollary~\ref{cor:v4v12-resonance} identifies the
zero-mode rank with that centralizer data.
\end{proof}

\begin{definition}[Flat-sector marginal parent]
\label{def:v4v12-flat-parent}
The flat-sector parent consists of:
\begin{enumerate}
\item a toron variable \(\tau\in\mathcal M_{\rm flat}\), stratified by
centralizer type;
\item every covariant zero mode on its stratum; and
\item nonzero fluctuations decomposed by the conditional spectral shells
\eqref{eq:v4v12-spectral-shell}.
\end{enumerate}
\end{definition}

\begin{firewall}[A finite-dimensional parent still needs its measure]
\label{fw:v4v12-parent-measure}
Definition~\ref{def:v4v12-flat-parent} identifies the variables but does not
construct the exact induced toron density after integrating transverse
fluctuations.  Near centralizer jumps, determinant factors can vanish and
the tubular Jacobian changes stratum.  Positivity, normalization, and
reflection positivity of the assembled parent remain to be proved.
\end{firewall}

\subsection{Correction of the terminal forest ledger}
\label{sec:v4v12-terminal-ledger}

\begin{proposition}[Terminal reclassification]
\label{prop:v4v12-terminal-reclassification}
A V11 terminal packet whose plaquettes are flat but whose periodic tree-
gauge links leave the identity chart is removed from the defect ledger and
assigned uniquely to the toron parent of
Definition~\ref{def:v4v12-flat-parent}.  It receives no Wilson-energy
payment.  A packet with nonzero curvature remains in the terminal defect
sector.
\end{proposition}

\begin{proof}
Theorem~\ref{thm:v4v12-flat-classification} assigns every flat packet a
unique simultaneous-conjugacy class \(\tau\).  Corollary~\ref{cor:v4v12-no-Wilson-payment}
forbids a defect payment.  Flat and nonflat plaquette conditions are
disjoint, so the remaining packet stays uniquely in the nonflat terminal
class.
\end{proof}

\subsection{V12 result ledger and V13 target}
\label{sec:v4v12-conclusion}

V12 classifies periodic flat lattice connections as commuting holonomy
tuples modulo conjugacy and proves that their Wilson action is exactly zero.
It computes the complete toron-shifted adjoint Laplacian spectrum, identifies
centralizer resonance strata, and replaces fixed Fourier shells by
toron-conditional covariant spectral shells with an exact band gap.  Flat
terminal packets are thereby removed from the V11 defect forest and promoted
to a finite-dimensional marginal parent without a fictitious energy payer.

The next gate is the induced measure of that parent.  V13 should disintegrate
the compact link Haar measure near each flat stratum, retain its normal
Jacobian and transverse determinant once, and test integrability at
centralizer jumps.  Nonflat terminal packets, sparse promotions, topological
charge sectors, and the hypercharge ultraviolet obstruction remain open.


\clearpage
\section{V13: induced toron measure and soft-mode transfer at centralizer walls}
\label{sec:v4v13}

V12 identified the periodic flat parent and showed that an adjoint shell
spectrum moves with its toron holonomy.  The remaining finite-cutoff question
is measure-theoretic: disintegrate link Haar measure near the flat locus,
retain the gauge-orbit and tubular Jacobians once, and avoid dividing by a
transverse determinant when its eigenvalue vanishes at a centralizer wall.

This note solves that problem on every compact regular stratum and gives an
exact soft-mode transfer valid near a wall.  The resulting finite-cutoff
parent is positive and normalized because it is a pushforward of the compact
Wilson measure.  Uniform continuum estimates, reflection positivity after
the nonlocal toron disintegration, and compatibility with dynamical gravity
remain open.

\subsection{The simultaneous-conjugation Jacobian}
\label{sec:v4v13-orbit-Jacobian}

Represent a flat toron by a commuting tuple
\(H=(H_1,\ldots,H_d)\in T^d\), with
\(H_\mu=e^{i\Theta_\mu}\).  Give \(G^d\) the product bi-invariant metric.
For a positive root \(\gamma\), define
\begin{equation}
 j_\gamma(H)
 =\sum_{\mu=1}^d
 \left|e^{i\gamma(\Theta_\mu)}-1\right|^2.
\label{eq:v4v13-root-Jacobian}
\end{equation}

\begin{proposition}[Orbit-volume factor on the flat locus]
\label{prop:v4v13-orbit-Jacobian}
Up to a constant fixed by Haar normalization, the Riemannian volume density
of the simultaneous-conjugation orbit through a regular commuting tuple is
\begin{equation}
 J_{\rm orb}(H)=\prod_{\gamma>0}j_\gamma(H).
\label{eq:v4v13-orbit-volume}
\end{equation}
It vanishes precisely when the common centralizer is larger than the maximal
torus.
\end{proposition}

\begin{proof}
For an infinitesimal conjugation \(X\in\mathfrak g\), the tangent vector is
\[
 X\longmapsto([X,H_1],\ldots,[X,H_d]).
\]
On the real two-plane associated with a root pair
\(\{\gamma,-\gamma\}\), its squared metric is
\(j_\gamma(H)\|X_\gamma\|^2\).  Distinct root planes are orthogonal and the
Cartan is the stabilizer.  The metric matrix on a root plane is therefore
\(j_\gamma I_2\), whose volume factor is
\(\sqrt{\det(j_\gamma I_2)}=j_\gamma\).  Multiplying over positive roots
gives \eqref{eq:v4v13-orbit-volume}.  A factor vanishes exactly when its
root plane commutes with every \(H_\mu\), which is exactly centralizer
enhancement by Corollary~\ref{cor:v4v12-resonance}.
\end{proof}

\begin{firewall}[The orbit factor is never inverted]
\label{fw:v4v13-no-invert-orbit}
The factor \(J_{\rm orb}\) is part of the pushforward volume density.  At a
centralizer wall it vanishes.  Treating gauge fixing as division by
\(J_{\rm orb}\) while also integrating the enhanced zero modes would count
the same orbit degeneracy twice and create a spurious singularity.
\end{firewall}

\subsection{Tubular disintegration on a regular stratum}
\label{sec:v4v13-tube}

Fix a centralizer stratum \(\Sigma\subset\mathcal M_{\rm flat}\) and a
compact subset \(K\Subset\Sigma\).  After tree gauge and quotient by the
constant stabilizer type, let \(\mathcal C_K\) denote the corresponding
compact piece of the flat configuration locus.  Let \(N\mathcal C_K\) be
its normal bundle inside the finite link-configuration quotient.

\begin{theorem}[Regular tubular Haar disintegration]
\label{thm:v4v13-tubular}
There are \(\varepsilon_K>0\), a tubular map
\begin{equation}
 \Phi_K:\{(\tau,z):\tau\in\mathcal C_K,
 z\in N_\tau\mathcal C_K,\ \|z\|<\varepsilon_K\}
 \longrightarrow\mathcal U_K,
\label{eq:v4v13-tubular-map}
\end{equation}
and a smooth positive density \(J_K(\tau,z)\) such that for every integrable
function supported in \(\mathcal U_K\),
\begin{equation}
 \int F(U)\,dU
 =\int_{\mathcal C_K}\int_{\|z\|<\varepsilon_K}
 F(\Phi_K(\tau,z))J_K(\tau,z)\,dz\,d\tau.
\label{eq:v4v13-tubular-disintegration}
\end{equation}
The density contains \(J_{\rm orb}(\tau)\) exactly once and is bounded above
and below by positive constants on a smaller tube over \(K\).
\end{theorem}

\begin{proof}
On a fixed orbit-type stratum, the gauge quotient is a smooth finite-
dimensional orbifold.  The normal exponential map is a diffeomorphism on a
uniform tube over the compact set \(K\).  The ordinary change-of-variables
formula gives \eqref{eq:v4v13-tubular-disintegration}; its Jacobian is smooth
and positive because the map is a local diffeomorphism.  Separating the
simultaneous-conjugation orbit coordinates gives the metric volume factor
in Proposition~\ref{prop:v4v13-orbit-Jacobian}.  Compactness supplies the
two-sided bounds after shrinking the tube.
\end{proof}

Let \(H_\tau^\perp\) be the Hessian of the unscaled Wilson energy on the
normal fibre.  On a regular stratum all gauge and flat tangent zero modes
have been removed.

\begin{lemma}[Normal positivity away from centralizer jumps]
\label{lem:v4v13-normal-positive}
There is \(c_K>0\) such that
\begin{equation}
 \langle z,H_\tau^\perp z\rangle\ge c_K\|z\|^2
 \qquad(\tau\in K).
\label{eq:v4v13-normal-gap}
\end{equation}
\end{lemma}

\begin{proof}
The Wilson Hessian at a flat connection is a positive constant times the
gauge-covariant complex \(d_\tau^*d_\tau\).  Its kernel consists of gauge
directions and infinitesimal flat deformations.  Both are tangent to the
orbit/toron stratum and have been removed from the normal fibre.  Hence the
normal Hessian is positive definite at every \(\tau\in K\).  Its least
eigenvalue is continuous, and its positive minimum on compact \(K\) is
\(c_K\).
\end{proof}

\subsection{Uniform Morse--Bott Laplace formula}
\label{sec:v4v13-Laplace}

\begin{theorem}[Regular-stratum induced density]
\label{thm:v4v13-Laplace}
Let \(r_K\) be the rank of the normal bundle.  For a smooth test function
\(F\) supported over the smaller tube,
\begin{align}
 &\int_{\mathcal U_K}
 F(U)e^{-E(U)/g^2}\,dU
\notag\\
 &\quad=(2\pi g^2)^{r_K/2}
 \int_{\mathcal C_K}
 F(\tau)J_K(\tau,0)
 \det(H_\tau^\perp)^{-1/2},d\tau
 +O(g^{r_K+1}),
\label{eq:v4v13-Laplace}
\end{align}
where the error is uniform on \(K\).  Thus the leading induced toron density
on a regular stratum is
\begin{equation}
 J_K(\tau,0)\det(H_\tau^\perp)^{-1/2},d\tau.
\label{eq:v4v13-induced-regular}
\end{equation}
\end{theorem}

\begin{proof}
In tubular coordinates, Taylor's theorem and flatness give
\[
 E(\Phi_K(\tau,z))
 =\frac12\langle z,H_\tau^\perp z\rangle+O(\|z\|^3)
\]
uniformly on \(K\).  Put \(z=gy\).  Lemma~\ref{lem:v4v13-normal-positive}
gives a uniform Gaussian majorant after shrinking the tube and splitting off
an exponentially small outer annulus.  The density and test function equal
their values at \(z=0\) plus \(O(g\|y\|)\).  Dominated integration gives the
uniform error, and the fibre Gaussian integral is
\((2\pi)^{r_K/2}\det(H_\tau^\perp)^{-1/2}\).  Restoring the Jacobian
\(g^{r_K}\) proves \eqref{eq:v4v13-Laplace}.
\end{proof}

\begin{firewall}[The regular formula is not uniform at a wall]
\label{fw:v4v13-wall-nonuniform}
As \(K\) approaches a centralizer wall, an eigenvalue of
\(H_\tau^\perp\) can vanish and the rank of the flat tangent space changes.
The determinant in \eqref{eq:v4v13-induced-regular} is then the wrong
coordinate formula.  Extending it through the wall without transferring the
soft eigenmode would create an artificial divergence.
\end{firewall}

\subsection{Exact soft-mode transfer}
\label{sec:v4v13-soft-transfer}

Fix \(\epsilon>0\).  For each toron, let
\begin{align}
 P_{\rm soft}(\tau)&=1_{[0,\epsilon)}(H_\tau),
\label{eq:v4v13-soft-projector}\\
 P_{\rm hard}(\tau)&=1_{[\epsilon,\infty)}(H_\tau).
\label{eq:v4v13-hard-projector}
\end{align}
Gauge and flat tangent directions are included in the soft parent before
this split.  Write a tubular fluctuation as \(z=s+h\) in the corresponding
spectral subspaces.

\begin{theorem}[Hard determinant bound and exact parent transfer]
\label{thm:v4v13-soft-transfer}
On every finite lattice, the tubular contribution has the exact form
\begin{equation}
 \int d\tau\int ds\,
 J_{\rm soft}(\tau,s)
 Z_{\rm hard}(\tau,s;g),
\label{eq:v4v13-exact-parent-integral}
\end{equation}
where
\begin{equation}
 Z_{\rm hard}(\tau,s;g)
 =\int dh\,
 J_{\rm hard}(\tau,s,h)e^{-E(\tau,s,h)/g^2}.
\label{eq:v4v13-hard-integral}
\end{equation}
Every quadratic hard eigenvalue is at least \(\epsilon\), so its Gaussian
determinant factor is bounded by
\begin{equation}
 \det(H_{\rm hard})^{-1/2}
 \le\epsilon^{-r_{\rm hard}/2}.
\label{eq:v4v13-hard-determinant}
\end{equation}
When an eigenvalue crosses \(\epsilon\), the full finite-dimensional mode is
transferred between \(s\) and \(h\); it is never simultaneously integrated
as a parent coordinate and divided out in the hard determinant.
\end{theorem}

\begin{proof}
Insert the two complementary spectral projections into the finite-
dimensional tubular coordinates and apply Fubini's theorem.  Absorb the
Jacobian of this orthogonal change of variables into the two displayed
densities; this gives the exact identity
\eqref{eq:v4v13-exact-parent-integral}.  The spectral theorem gives the lower
bound \(H_{\rm hard}\ge\epsilon I\), and multiplying its eigenvalue bounds
gives \eqref{eq:v4v13-hard-determinant}.  At a crossing, complementary
spectral subspaces still sum to the full fibre.  Assigning the crossing mode
to one side by the half-open interval convention preserves the exact Fubini
integral and prevents duplication.
\end{proof}

The soft rank can jump with \(\tau\), so a single smooth vector bundle need
not cover the wall.  Use the finite centralizer stratification of
Proposition~\ref{prop:v4v12-stratification} and a finite measurable spectral
partition.  Formula \eqref{eq:v4v13-exact-parent-integral} remains an exact
positive disintegration on each piece.

\begin{proposition}[Finite-cutoff parent positivity and normalization]
\label{prop:v4v13-parent-positive}
Let \(\Pi_g\) be the pushforward of the normalized finite-lattice Wilson
measure under the map which records the toron stratum and all transferred
soft coordinates on the chosen tubes, and records one terminal symbol outside
their union.  Then \(\Pi_g\) is a positive probability measure.  Its
conditional hard factor is \eqref{eq:v4v13-hard-integral}, and the orbit,
tubular, and hard normalization factors each occur once.
\end{proposition}

\begin{proof}
A measurable pushforward of a positive probability measure is a positive
probability measure.  Disintegration of the original integral by
Theorem~\ref{thm:v4v13-tubular} and
Theorem~\ref{thm:v4v13-soft-transfer} gives the stated conditional factor.
Both the coordinate changes and Fubini split are identities, so no Jacobian
or determinant can occur a second time.
\end{proof}

\begin{firewall}[Finite positivity is not continuum tightness]
\label{fw:v4v13-no-tightness}
Proposition~\ref{prop:v4v13-parent-positive} does not give a lattice-uniform
bound on the soft rank, the induced density, or its response derivatives.
It also does not show that toron disintegration commutes with reflection or
with the Einstein metric integral.  These are required before the parent can
enter the full V4 projective family.
\end{firewall}

\subsection{V13 result ledger and V14 target}
\label{sec:v4v13-conclusion}

V13 computes the simultaneous-conjugation orbit Jacobian and proves a smooth
positive tubular Haar disintegration on compact regular toron strata.  A
uniform Morse--Bott calculation gives the induced regular-stratum density.
At centralizer walls, the note replaces the singular determinant formula by
an exact spectral transfer: every soft mode becomes a parent coordinate,
while the remaining hard determinant has a fixed lower spectral cutoff.
The finite-cutoff parent is therefore positive and normalized with no double
Jacobian.

V14 should seek the first uniform estimate needed for the continuum: bound
the response derivatives of the hard conditional factor across the finite
stratification and determine whether the number of transferred modes per
physical volume is tight.  Sparse nonflat promotions and genuinely
topological terminal sectors remain outside the toron parent.  Hypercharge,
the coupled Higgs--Yukawa trajectory, and dynamical Einstein fluctuations
also remain open.


\clearpage
\section{V14: toron-uniform soft-mode counting and hard response bounds}
\label{sec:v4v14}

V13 constructed a positive finite-cutoff toron parent by transferring every
soft normal mode before forming the hard determinant.  Two uniform estimates
are needed before this disintegration can enter the marginal shell parent of
V4.  First, the number of transferred modes per physical volume must remain
bounded uniformly in the toron.  Second, first and second derivatives of the
hard conditional logarithm must be controlled without differentiating a
vanishing determinant.

The first estimate follows exactly from the shifted spectrum of V12 and a
uniform lattice-point count.  The second follows from the connected response
identity and Brascamp--Lieb whenever the complete hard potential is uniformly
strongly convex.  A quantitative Hessian criterion reduces that convexity to
domination of all interaction and Jacobian curvatures by the transferred
hard spectral gap.  Proving this domination for the full gauge--ghost--matter
packet remains open.

\subsection{Uniform counting for shifted lattice momenta}
\label{sec:v4v14-counting}

Let \(L=Ma\) be the physical side length of the periodic lattice.  For a
phase vector \(\vartheta\in[-\pi,\pi)^d\), define
\begin{equation}
 \lambda_k(\vartheta)
 =\frac4{a^2}\sum_{\mu=1}^d
 \sin^2\!\left(
 \frac{2\pi k_\mu+\vartheta_\mu}{2M}
 \right),
 \qquad k\in\{0,\ldots,M-1\}^d.
\label{eq:v4v14-shifted-symbol}
\end{equation}

\begin{lemma}[Uniform shifted lattice-point count]
\label{lem:v4v14-lattice-count}
There is \(C_d\) such that, for every shift \(\vartheta\) and
\(0<\epsilon\le a^{-2}\),
\begin{equation}
 \#\{k:\lambda_k(\vartheta)\le\epsilon\}
 \le C_d\bigl(1+L^d\epsilon^{d/2}\bigr).
\label{eq:v4v14-lattice-count}
\end{equation}
\end{lemma}

\begin{proof}
For each coordinate choose the representative
\(x_\mu\in[-\pi M,\pi M]\) of
\(2\pi k_\mu+\vartheta_\mu\) modulo \(2\pi M\).  Since
\(|x_\mu/(2M)|\le\pi/2\),
\(|\sin(x_\mu/(2M))|\ge |x_\mu|/(\pi M)\).  Thus
\(\lambda_k(\vartheta)\le\epsilon\) implies
\begin{equation}
 \sum_\mu x_\mu^2
 \le C M^2a^2\epsilon=C L^2\epsilon.
\label{eq:v4v14-momentum-ball}
\end{equation}
The shifted grid \(2\pi\mathbb Z^d+\vartheta\) has fixed spacing.  The
number of its points in a radius-\(R\) Euclidean ball is bounded uniformly in
the shift by \(C_d(1+R^d)\): cover each grid point by a disjoint cube of
fixed side and compare their union with the ball enlarged by one cube.
Take \(R=C L\sqrt\epsilon\) in
\eqref{eq:v4v14-momentum-ball}.
\end{proof}

\begin{theorem}[Toron-uniform soft-rank density]
\label{thm:v4v14-soft-rank}
Let \(N_{\rm soft}(\tau,\epsilon)\) be the number of adjoint site modes of
\(\Delta_\tau\) with eigenvalue at most \(\epsilon\), counted with
multiplicity.  Then
\begin{align}
 N_{\rm soft}(\tau,\epsilon)
 &\le C_d\dim(G)
 \bigl(1+L^d\epsilon^{d/2}\bigr),
\label{eq:v4v14-soft-rank}\\
 \frac{N_{\rm soft}(\tau,\epsilon)}{L^d}
 &\le C_d\dim(G)
 \bigl(L^{-d}+\epsilon^{d/2}\bigr),
\label{eq:v4v14-soft-density}
\end{align}
uniformly in \(\tau\).  A gauge-fixed adjoint one-form Hessian satisfying a
uniform spectral comparison with at most \(d\) copies of \(\Delta_\tau\)
obeys the same estimate with an additional fixed factor.
\end{theorem}

\begin{proof}
Theorem~\ref{thm:v4v12-shifted-spectrum} gives one shifted scalar lattice
spectrum for each Cartan or root direction.  Apply
Lemma~\ref{lem:v4v14-lattice-count} to every direction and sum over the
\(\dim(G)\) internal modes.  Divide by \(L^d\) for the second formula.  The
one-form statement follows from the assumed spectral comparison and its
fixed component multiplicity.
\end{proof}

In four dimensions the density bound is
\begin{equation}
 L^{-4}N_{\rm soft}(\tau,\epsilon)
 \le C\dim(G)(L^{-4}+\epsilon^2).
\label{eq:v4v14-four-density}
\end{equation}
Thus a fixed physical soft threshold gives a uniformly locally finite parent
field.  The exactly flat toron/centralizer zero modes contribute only a fixed
finite number on each finite-dimensional internal factor.

\begin{firewall}[Bounded density is not finite total rank]
\label{fw:v4v14-density-scope}
At fixed \(\epsilon>0\), the total number of soft modes grows proportionally
to physical volume.  They are the remaining infrared field, not a finite
list of collective torons.  Only the exactly flat holonomy and centralizer
zero modes are finite dimensional.  Claiming a volume-independent total
soft rank would contradict \eqref{eq:v4v14-soft-density}.
\end{firewall}

\subsection{Connected derivatives of the hard conditional factor}
\label{sec:v4v14-hard-response}

Let \(\theta\) denote any finite parent coordinate: a toron coordinate, a
transferred soft amplitude, or a background source.  On a convex hard fibre
\(\mathcal H\), write
\begin{align}
 Z_{\rm h}(\theta)
 &=\int_{\mathcal H}e^{-\mathcal A(\theta,h)}\,dh,
\label{eq:v4v14-hard-Z}\\
 \Gamma_{\rm h}(\theta)&=-\log Z_{\rm h}(\theta),
\label{eq:v4v14-hard-Gamma}\\
 d\mu_\theta(h)&=Z_{\rm h}(\theta)^{-1}
 e^{-\mathcal A(\theta,h)}\,dh.
\label{eq:v4v14-hard-measure}
\end{align}
The real hard potential \(\mathcal A\) includes the Wilson action divided by
\(g^2\), the modulus of the assembled determinant packet, and minus the
logarithm of the tubular Jacobian exactly once.  Any unit-modulus chiral
phase is retained as a separate observable or parent factor; it is not put
inside the positive measure used below.

\begin{theorem}[Hard-factor response theorem]
\label{thm:v4v14-response}
Assume differentiation under the integral is valid and
\begin{equation}
 D_h^2\mathcal A(\theta,h)\ge\kappa_{\rm h}I
\label{eq:v4v14-hard-convexity}
\end{equation}
uniformly on the fibre.  For parent tangent vectors \(v,w\), put
\begin{align}
 X_v(h)&=D_\theta\mathcal A(\theta,h)[v],
\label{eq:v4v14-X}\\
 Y_{v,w}(h)&=D_\theta^2\mathcal A(\theta,h)[v,w].
\label{eq:v4v14-Y}
\end{align}
Then
\begin{align}
 D\Gamma_{\rm h}(\theta)[v]
 &=\mathbb E_{\mu_\theta}X_v,
\label{eq:v4v14-first-hard-response}\\
 D^2\Gamma_{\rm h}(\theta)[v,w]
 &=\mathbb E_{\mu_\theta}Y_{v,w}
 -\langle X_v;X_w\rangle_{\mu_\theta}^{\rm conn}.
\label{eq:v4v14-second-hard-response}
\end{align}
If
\begin{align}
 |X_v|&\le M_0\|v\|,
\label{eq:v4v14-X-bound}\\
 |Y_{v,w}|&\le M_2\|v\|\|w\|,
\label{eq:v4v14-Y-bound}\\
 \|\nabla_hX_v\|&\le M_1\|v\|,
\label{eq:v4v14-gradient-bound}
\end{align}
then
\begin{align}
 |D\Gamma_{\rm h}[v]|
 &\le M_0\|v\|,
\label{eq:v4v14-first-bound}\\
 |D^2\Gamma_{\rm h}[v,w]|
 &\le\left(M_2+\frac{M_1^2}{\kappa_{\rm h}}\right)
 \|v\|\|w\|.
\label{eq:v4v14-second-bound}
\end{align}
\end{theorem}

\begin{proof}
Differentiate the negative logarithm in
\eqref{eq:v4v14-hard-Gamma}; this gives
\eqref{eq:v4v14-first-hard-response}.  Differentiating the normalized
expectation gives the expectation of \(Y_{v,w}\) minus the connected
covariance, proving \eqref{eq:v4v14-second-hard-response}.

Strong convexity gives the Brascamp--Lieb covariance estimate
\begin{equation}
 |\langle f;k\rangle^{\rm conn}|
 \le\frac1{\kappa_{\rm h}}
 \bigl(\mathbb E\|\nabla_hf\|^2\bigr)^{1/2}
 \bigl(\mathbb E\|\nabla_hk\|^2\bigr)^{1/2}.
\label{eq:v4v14-BL}
\end{equation}
Apply it to \(X_v,X_w\), use
\eqref{eq:v4v14-X-bound}--%
\eqref{eq:v4v14-gradient-bound}, and insert the result into the two response
identities.
\end{proof}

\begin{firewall}[The response packet remains assembled]
\label{fw:v4v14-response-assembled}
The expectation and covariance in
\eqref{eq:v4v14-second-hard-response} arise from one derivative of one hard
conditional logarithm.  Bound \eqref{eq:v4v14-second-bound} controls their
complete difference.  It does not provide two independent determinant or
entropy payments.
\end{firewall}

\begin{firewall}[Brascamp--Lieb does not control the phase]
\label{fw:v4v14-phase}
Theorem~\ref{thm:v4v14-response} applies to the positive modulus carrier.
It does not bound derivatives of the chiral determinant phase.  That phase
must use its separately normalized anomaly/holonomy construction and cannot
be absorbed into a fictitious real convex potential.
\end{firewall}

\subsection{A quantitative hard-convexity certificate}
\label{sec:v4v14-convexity-certificate}

In a hard spectral chart, decompose
\begin{equation}
 \mathcal A(\theta,h)
 =\frac1{2g^2}\langle h,H_{\rm h}(\theta)h\rangle
 +\mathcal V(\theta,h),
\label{eq:v4v14-hard-decomposition}
\end{equation}
where every eigenvalue of \(H_{\rm h}\) is at least \(\epsilon\).  The term
\(\mathcal V\) includes nonlinear interactions and the logarithmic Jacobian
with their actual signs.

\begin{proposition}[Spectral domination criterion]
\label{prop:v4v14-spectral-domination}
If
\begin{equation}
 \|D_h^2\mathcal V(\theta,h)_-\|
 \le\frac{\epsilon}{2g^2}
\label{eq:v4v14-negative-Hessian}
\end{equation}
uniformly on the hard fibre, where the subscript denotes the negative part
as a quadratic form, then
\begin{equation}
 D_h^2\mathcal A\ge\frac{\epsilon}{2g^2}I.
\label{eq:v4v14-certified-convexity}
\end{equation}
Consequently Theorem~\ref{thm:v4v14-response} applies with
\(\kappa_{\rm h}=\epsilon/(2g^2)\).
\end{proposition}

\begin{proof}
The quadratic part of \eqref{eq:v4v14-hard-decomposition} has Hessian
\(g^{-2}H_{\rm h}\ge\epsilon g^{-2}I\).  Subtract the largest allowed
negative part in \eqref{eq:v4v14-negative-Hessian}; the remainder is at least
\(\epsilon/(2g^2)I\).
\end{proof}

The criterion is stable under positive Higgs radial curvature but is not
automatic for gauge self-interactions, fermionic logarithms, ghost
cancellations, or the gravitational conformal direction.  Those terms must
first be assembled into their Ward/BRST packet and then tested against
\eqref{eq:v4v14-negative-Hessian}.

\subsection{Stratumwise soft transfer without a singular derivative}
\label{sec:v4v14-strata}

On a compact finite lattice, the toron moduli has the finite centralizer
stratification of V12.  Around each stratum choose a neighborhood and transfer
every mode whose eigenvalue can fall below \(2\epsilon\) in that
neighborhood.  The complementary hard Hessian is then at least
\(2\epsilon\); shrinking the neighborhood leaves a uniform hard lower bound
\(\epsilon\).

\begin{proposition}[Buffered stratum cover]
\label{prop:v4v14-buffered-cover}
There is a finite cover \(\{\mathcal U_i\}_{i=1}^J\) of the finite toron
moduli and, on each member, a fixed-rank soft bundle such that:
\begin{enumerate}
\item every zero or potentially sub-\(\epsilon\) eigenmode belongs to the
soft bundle;
\item the complementary hard Hessian is at least \(\epsilon I\); and
\item a measurable disjoint refinement \(\{\mathcal R_i\}\) with
\(\mathcal R_i\subset\mathcal U_i\) partitions the toron parent, so every
configuration uses exactly one local response formula.
\end{enumerate}
\end{proposition}

\begin{proof}
At a fixed toron, the Hessian is a finite self-adjoint matrix.  Choose a
number \(t_\tau\in[2\epsilon,3\epsilon]\) which is not an eigenvalue, and put
every eigenvector below \(t_\tau\) into the soft space.  The Riesz spectral
projection onto this finite cluster has constant rank on a neighborhood of
\(\tau\).  By continuity, every complementary eigenvalue remains above
\(\epsilon\) after shrinking that neighborhood.  Compactness gives a finite
subcover.
Order it and set
\(\mathcal R_1=\mathcal U_1\),
\(\mathcal R_i=\mathcal U_i\setminus\bigcup_{k<i}\mathcal U_k\).
These measurable sets are disjoint and exhaustive.  On each one use its
containing chart, so no sharp global spectral projector is differentiated.
\end{proof}

\begin{firewall}[Finite cover constants may deteriorate]
\label{fw:v4v14-cover-uniformity}
Proposition~\ref{prop:v4v14-buffered-cover} is uniform on one finite lattice.
It does not show that the number of charts, their spectral buffers, or their
parent-coordinate transition maps remain controlled as volume grows and the
spectrum becomes dense.  The rank density estimate
\eqref{eq:v4v14-soft-density} controls how many modes are transferred, but
not the smoothness of all chart transitions.
\end{firewall}

\subsection{Conditional hard-parent theorem}
\label{sec:v4v14-conditional}

\begin{theorem}[Uniform hard responses under the certified hypotheses]
\label{thm:v4v14-conditional-parent}
Suppose along a cutoff family:
\begin{enumerate}
\item the gauge-fixed one-form Hessian obeys the spectral comparison in
Theorem~\ref{thm:v4v14-soft-rank};
\item the buffered stratum cover has volume-uniform local transition
constants;
\item \eqref{eq:v4v14-negative-Hessian} holds with fixed
\(\epsilon>0\); and
\item for physical volume \(V\), the coupling-normalized parent marks obey
\begin{align*}
 |\mathbb E X_v|&\le M_0V\|v\|,\\
 |\mathbb E Y_{v,w}|&\le M_2V\|v\|\|w\|,\\
 \mathbb E\|\nabla_hX_v\|^2&\le M_1^2V\|v\|^2,
\end{align*}
with constants independent of the cutoff and volume.
\end{enumerate}
Then the transferred soft rank per physical volume and the first two hard
parent responses per physical volume are uniformly bounded.  In particular,
the hard toron factor supplies no new determinant singularity at a
centralizer wall.
\end{theorem}

\begin{proof}
The soft-rank assertion is
Theorem~\ref{thm:v4v14-soft-rank}.  On every stratum chart,
Proposition~\ref{prop:v4v14-spectral-domination} supplies the strong
convexity constant.  Equations \eqref{eq:v4v14-first-hard-response}--%
\eqref{eq:v4v14-BL} and the volume-scaled mark hypotheses give response
bounds proportional to \(V\); division by \(V\) is uniform.  The assumed
chart constants allow the finite disjoint refinement to be summed without a
cutoff-dependent loss.  Soft transfer
removes every sub-\(\epsilon\) eigenvalue before the hard determinant is
formed, so \eqref{eq:v4v13-hard-determinant} remains bounded.
\end{proof}

\begin{firewall}[The conditional hypotheses are not yet proved]
\label{fw:v4v14-conditional-open}
Theorem~\ref{thm:v4v14-conditional-parent} does not establish the negative-
Hessian bound for the interacting Standard Model shell, the uniform
stratum-transition constants, reflection positivity of the resulting parent, or uniform derivatives
of the separately retained chiral phase.  It also keeps the metric fixed.  Therefore it is not a construction
of the SM--Einstein continuum measure.
\end{firewall}

\subsection{V14 result ledger and V15 audit target}
\label{sec:v4v14-conclusion}

V14 proves a toron-uniform lattice Weyl bound: in four dimensions the number
of transferred modes per physical volume is at most
\(C(L^{-4}+\epsilon^2)\).  It proves exact first and second connected
response formulas for the hard factor and bounds them by Brascamp--Lieb under
a complete hard-fibre convexity hypothesis.  A quantitative spectral
domination inequality and a finite buffered stratum cover state precisely
what remains to make those bounds cutoff uniform.

V15 is the next compulsory NS/YM audit.  It should use the newest companion
results to choose between the two open hard-parent gates: proving the
negative-Hessian domination on the assembled gauge--matter packet, or
replacing global convexity by a local cluster/Poincare estimate which remains
valid after V11 defect repair.  Hypercharge, Higgs--Yukawa running,
topological terminal sectors, and dynamical Einstein fluctuations remain
open.


\clearpage
\section{V15: companion audit and the two-mark connected polymer response}
\label{sec:v4v15}

V14 gave a Brascamp--Lieb response bound for a globally convex positive hard
carrier.  The compulsory companion audit shows that this is stronger than
the next local need.  The transverse beta coordinate and toron response
require a genuine second derivative of one connected logarithm.  They do not
require two copies of a first-derivative estimate.  A derivative-marked
polymer expansion supplies exactly that second response and remains valid
for complex activities, so it can retain the separately normalized chiral
phase.

This note proves the two-mark cluster theorem, including the exhaustive
same-polymer/two-polymer marking identity, and combines the good-field and
V11 repaired-defect gases.  It gives a second route to the V14 hard response
bounds which uses local connectedness in place of global convexity.  The
uniform marked activity estimate for the complete SM gauge--matter packet is
still open.

\subsection{Compulsory V15 audit}
\label{sec:v4v15-audit}

The exact companion files read were
\begin{enumerate}
\item \texttt{Renewal\_Yang\_Mills\_Mass\_Gap\_Research\_Notes\_V35.tex},
\(673335\) bytes, with SHA--256
\begin{center}\scriptsize\ttfamily
8B375EBC8A980641DF66EECF2656B33CE5ED45908EBF30FF8D0B9090AD478C88.
\end{center}
\item \texttt{Renewal\_NS\_Stability\_Research\_Notes\_V31.tex},
\(887537\) bytes, with SHA--256
\begin{center}\scriptsize\ttfamily
F1121B62238AD5491BB7CF03635EA9934942EDF573ED8FAAA9EE79E1D38A6696.
\end{center}
\end{enumerate}

Yang--Mills V35 reduces its first new hot obstruction to a star center with
two demand tokens.  One linear endpoint response cannot be squared.  The
valid objects are a proved same-bank second moment or a composite covariance
of the two arms, formed before a positive envelope.  For the present hard
factor, these are exactly the two terms produced when two parent derivatives
hit one polymer activity or two distinct activities in the same connected
cluster.

Navier--Stokes V31 shows that unsigned heat lifting restores the missing
critical scale and that only the complete signed nonlinear correlation could
improve it.  The present transfer is to differentiate the signed/complex
connected cluster first.  Absolute tree bounds are applied only after the
same-polymer and two-polymer terms have been assembled.

\begin{proposition}[Audit response rule]
\label{prop:v4v15-audit-rule}
A second hard-parent response is admissible only if it is represented by:
\begin{enumerate}
\item one actual second derivative of a complete activity; or
\item two actual first derivatives on two vertices of one connected cluster.
\end{enumerate}
Two copies of a bound for \(D\log Z\) do not imply a bound for
\(D^2\log Z\).
\end{proposition}

\subsection{Derivative-marked polymer gas}
\label{sec:v4v15-polymer-gas}

Let \(\mathfrak P_\Lambda\) be a finite polymer set over a finite lattice
\(\Lambda\), with a symmetric incompatibility relation \(X\not\sim Y\).
For a parent parameter \(\theta\) in a Banach space, let
\(w_\theta(X)\in\mathbb C\) be twice Fr\'echet differentiable.  Define
\begin{equation}
 Z_\Lambda(\theta)
 =\sum_{\substack{\mathcal X\subset\mathfrak P_\Lambda\\
 X\sim Y\text{ for distinct }X,Y\in\mathcal X}}
 \prod_{X\in\mathcal X}w_\theta(X).
\label{eq:v4v15-polymer-Z}
\end{equation}
Here \(X\sim Y\) means compatible.

For \(m=0,1,2\), put
\begin{equation}
 q_m(X)
 =\sup_{\|v_1\|,\ldots,\|v_m\|\le1}
 \left|D^mw_\theta(X)[v_1,\ldots,v_m]\right|,
\label{eq:v4v15-mark-size}
\end{equation}
with \(q_0=|w_\theta|\).  For \(a,\rho>0\), define
\begin{equation}
 \|w\|_{a,2;\rho}
 =\sup_{x\in\Lambda}
 \sum_{X\ni x}e^{a|X|}
 \bigl(q_0(X)+\rho q_1(X)+\rho^2q_2(X)\bigr).
\label{eq:v4v15-marked-norm}
\end{equation}

\begin{lemma}[Exact two-mark product identity]
\label{lem:v4v15-two-mark-identity}
For a finite ordered cluster \(\mathbf X=(X_1,\ldots,X_n)\),
\begin{align}
 D^2\prod_{i=1}^nw(X_i)[v_1,v_2]
 &=\sum_{i=1}^n
 D^2w(X_i)[v_1,v_2]
 \prod_{k\ne i}w(X_k)
\notag\\
 &\quad+
 \sum_{\substack{i,j=1\\i\ne j}}^n
 Dw(X_i)[v_1]Dw(X_j)[v_2]
 \prod_{k\ne i,j}w(X_k).
\label{eq:v4v15-two-mark-identity}
\end{align}
The first line is the same-activity second moment; the second line contains
two distinct activity occurrences.  The two classes are disjoint and
exhaustive.
\end{lemma}

\begin{proof}
Apply the product rule once with direction \(v_1\), producing one marked
factor.  Apply it again with direction \(v_2\).  The second derivative either
hits the marked factor again or hits one of the other factors.  These are
respectively the two displayed sums, and every ordered placement occurs
once.
\end{proof}

\begin{firewall}[A repeated mark is not a duplicated activity]
\label{fw:v4v15-repeated-mark}
The first line of \eqref{eq:v4v15-two-mark-identity} differentiates one
activity twice.  It may use a genuine quadratic activity estimate, but not
two independent copies of \(q_1(X)\).  This is the exact counterpart of the
Yang--Mills V35 star-center firewall.
\end{firewall}

\subsection{The connected two-mark theorem}
\label{sec:v4v15-connected-theorem}

Let
\begin{equation}
 \phi^T(X_1,\ldots,X_n)
 =\sum_{G\in\mathcal G_n^{\rm conn}}
 \prod_{\{i,j\}\in E(G)}
 \bigl(-1_{\{X_i\not\sim X_j\}}\bigr)
\label{eq:v4v15-Ursell}
\end{equation}
be the Ursell coefficient.  Whenever it converges absolutely,
\begin{equation}
 \log Z_\Lambda
 =\sum_{n\ge1}\frac1{n!}
 \sum_{X_1,\ldots,X_n}
 \phi^T(X_1,\ldots,X_n)
 \prod_{i=1}^nw(X_i).
\label{eq:v4v15-cluster-log}
\end{equation}

\begin{theorem}[Two-mark connected polymer response]
\label{thm:v4v15-two-mark}
For each finite-range incompatibility geometry and \(a>0\), there is
\(\eta_*(a)>0\) with the following property.  If
\begin{equation}
 \|w\|_{a,2;\rho}\le\eta<\eta_*(a),
\label{eq:v4v15-small-KP}
\end{equation}
then \eqref{eq:v4v15-cluster-log} and its first two parent derivatives
converge absolutely.  If the derivative marks vanish unless the marked
polymer meets a fixed set \(A\subset\Lambda\), then
\begin{align}
 |D\log Z_\Lambda[v]|
 &\le C_1(a,\eta)\rho^{-1}|A|\|v\|,
\label{eq:v4v15-first-cluster-bound}\\
 |D^2\log Z_\Lambda[v_1,v_2]|
 &\le C_2(a,\eta)\rho^{-2}|A|
 \|v_1\|\|v_2\|.
\label{eq:v4v15-second-cluster-bound}
\end{align}
The constants are independent of \(\Lambda\) and remain valid for complex
activities.
\end{theorem}

\begin{proof}
Use the tree-graph bound for \(\phi^T\): the absolute connected-graph sum is
bounded by a sum over rooted spanning trees with a fixed exponential factor
per vertex.  Root a differentiated cluster at the first marked polymer and
then at a site of its intersection with \(A\).  Sum tree leaves inward.
At an unmarked leaf, \eqref{eq:v4v15-marked-norm} pays \(q_0e^{a|X|}\); at
a once-marked leaf it pays \(\rho q_1e^{a|X|}\); at a twice-marked leaf it
pays \(\rho^2q_2e^{a|X|}\).  The residual tree series is geometric for
\(\eta<\eta_*(a)\).

For the first derivative there is one possible mark and one factor
\(\rho^{-1}\).  For the second derivative insert the disjoint identity
\eqref{eq:v4v15-two-mark-identity}.  Its same-activity terms are paid by
\(\rho^2q_2\); its two-activity terms are paid by two factors
\(\rho q_1\) on two actual vertices of the same tree.  In both cases there
is one rooted connected cluster and total factor \(\rho^{-2}|A|\).  Absolute
convergence justifies termwise differentiation.  The proof uses moduli only
after the complex products and Ursell coefficient have been formed, so it
does not require positivity.
\end{proof}

\begin{corollary}[Extensive and local forms]
\label{cor:v4v15-local-extensive}
For a translation-covariant extensive parent source, take
\(A=\Lambda\) in Theorem~\ref{thm:v4v15-two-mark}; the first two derivatives
of \(|\Lambda|^{-1}\log Z_\Lambda\) are uniformly bounded.  For a local
background insertion with fixed support, both derivatives are uniformly
local and independent of volume.
\end{corollary}

\begin{proof}
Divide \eqref{eq:v4v15-first-cluster-bound}--%
\eqref{eq:v4v15-second-cluster-bound} by \(|\Lambda|\) in the extensive
case.  In the local case, \(|A|\) is fixed.
\end{proof}

\subsection{Complex phase and zero-free logarithm}
\label{sec:v4v15-complex}

\begin{proposition}[Cluster-defined phase response]
\label{prop:v4v15-phase}
Under \eqref{eq:v4v15-small-KP}, the right side of
\eqref{eq:v4v15-cluster-log} defines a single analytic branch of
\(\log Z_\Lambda\) connected to the empty-polymer value \(Z=1\).  Its real
and imaginary parts obey the derivative bounds of
Theorem~\ref{thm:v4v15-two-mark}.  Hence a chiral phase represented inside
the complete complex activities has bounded local first and second
responses under the same marked KP hypothesis.
\end{proposition}

\begin{proof}
Absolute convergence of the cluster series defines an analytic logarithm
whose exponential equals the polymer partition function near zero
activities.  Analytic continuation within the connected KP ball fixes the
branch and prevents a zero of \(Z\) there.  The real and imaginary parts are
bounded by the modulus of the complex derivative already estimated in
Theorem~\ref{thm:v4v15-two-mark}.
\end{proof}

\begin{firewall}[Anomaly normalization remains prior data]
\label{fw:v4v15-anomaly}
Proposition~\ref{prop:v4v15-phase} bounds a phase only after the anomaly-
cancelled determinant line and its finite-cutoff branch have been defined.
It does not construct that global line or replace the earlier mapping-torus
holonomy argument.  A local cluster branch cannot erase a global anomaly.
\end{firewall}

\subsection{Composition of good and repaired-defect activities}
\label{sec:v4v15-composition}

Let \(w^{\rm g}\) be the good-field activities from the V6 shell expansion
and \(w^{\rm d}\) the controlled repaired-defect activities of V11.  Regard
them as two species in one incompatibility graph.  Include any irreducible
mixed activity as a third species rather than factoring it artificially.

\begin{lemma}[Marked-norm composition]
\label{lem:v4v15-composition}
For the disjoint union activity family,
\begin{equation}
 \|w^{\rm g}\sqcup w^{\rm d}\sqcup w^{\rm mix}\|_{a,2;\rho}
 \le
 \|w^{\rm g}\|_{a,2;\rho}
 +\|w^{\rm d}\|_{a,2;\rho}
 +\|w^{\rm mix}\|_{a,2;\rho}.
\label{eq:v4v15-composition}
\end{equation}
\end{lemma}

\begin{proof}
For every root site, the defining sum
\eqref{eq:v4v15-marked-norm} over the union of species is the sum of the
three species sums.  Taking the supremum gives the inequality.
\end{proof}

\begin{theorem}[Local hard response without global convexity]
\label{thm:v4v15-hard-response}
Condition on one V13 toron/soft-parent chart.  Suppose the complete hard
factor admits the above polymer representation and
\begin{equation}
 \|w^{\rm g}\|_{a,2;\rho}
 +\|w^{\rm d}\|_{a,2;\rho}
 +\|w^{\rm mix}\|_{a,2;\rho}
 <\eta_*(a)
\label{eq:v4v15-combined-smallness}
\end{equation}
uniformly in shell, volume, and parent coordinate.  Then its local first and
second responses, including the complex phase response, satisfy the uniform
bounds \eqref{eq:v4v15-first-cluster-bound}--%
\eqref{eq:v4v15-second-cluster-bound}.  No global hard-fibre convexity is
required.
\end{theorem}

\begin{proof}
Lemma~\ref{lem:v4v15-composition} verifies the hypothesis of
Theorem~\ref{thm:v4v15-two-mark} for the complete activity family.  Apply
that theorem and Proposition~\ref{prop:v4v15-phase}.
\end{proof}

\begin{corollary}[Conditional completion of the V14 response gate]
\label{cor:v4v15-v14-gate}
If \eqref{eq:v4v15-combined-smallness} and the V14 soft-rank and chart-
transition bounds hold, then the toron-parent hard factor has cutoff-uniform
local first and second responses without assuming
\eqref{eq:v4v14-negative-Hessian}.  If its unmarked and twice-marked
remainder is also small relative to the coefficients
\eqref{eq:v4v6-shell-constant}, the V6 exact one-loop dominance theorem
applies to the nonabelian factors.
\end{corollary}

\begin{proof}
Theorem~\ref{thm:v4v15-hard-response} supplies the response bounds.  The V14
rank and transition estimates allow the conditional charts to be summed.
The final assertion is exactly Theorem~\ref{thm:v4v6-dominance} with the
cluster remainder supplying its uniform small-field error.
\end{proof}

\begin{firewall}[The marked activity hypothesis is open]
\label{fw:v4v15-marked-open}
V6 postulated a derivative-marked good-field polymer estimate; V11 proved
the defect part only on the controlled repairable sector.  No result here
proves \eqref{eq:v4v15-combined-smallness} for the complete chiral
gauge--Higgs--Yukawa shell, its sparse-promotion and terminal sectors, or a
dynamical metric.  The theorem identifies the required second-order object
but does not assume it has already been constructed.
\end{firewall}

\subsection{V15 result ledger and V16 target}
\label{sec:v4v15-conclusion}

V15 records the compulsory companion audit and proves the exact two-mark
product identity.  The connected polymer theorem treats a genuine
same-activity second derivative and two first derivatives on two actual
cluster vertices as disjoint exhaustive cases.  Under one marked KP norm it
gives volume-uniform local and intensive first/second responses for complex
activities.  This provides a local alternative to V14 global convexity and
can retain the anomaly-normalized phase.

V16 should attack the first unproved marked activity estimate.  The smallest
closed subsystem is a fixed toron chart with the gauge shell, ghosts, and one
matter representation assembled into a background-Ward polymer.  The target
is the coupling-normalized \(q_2\) bound for a single complete activity; a
square of its \(q_1\) estimate is explicitly insufficient.  Terminal and
sparse defect sectors, hypercharge UV completion, the coupled Higgs--Yukawa
trajectory, and Einstein fluctuations remain open.


\clearpage
\section{V16: hard trace-log activities and a genuine quadratic response mark}
\label{sec:v4v16}

V15 reduced the hard-parent response problem to a marked activity norm with
one genuine second derivative.  This note proves that estimate for the
Gaussian determinant part of one fixed toron chart.  The proof has three
steps.  An exact resolvent identity gives the first and second derivatives of
one trace logarithm.  A M\"obius inclusion--exclusion extracts its connected
polymer support.  A Combes--Thomas bound for the hard inverse then turns the
connected trace words into exponentially decaying activities.

The construction applies simultaneously to the gauge Hessian determinant,
the ghost determinant, and one hard matter representation, with their actual
packet signs retained until the end.  It does not control the nonlinear
bosonic interaction gas or terminal large-field sectors.

\subsection{One finite-rank trace logarithm}
\label{sec:v4v16-trace-log}

Let \(K(\theta)\) be a twice differentiable finite-rank operator and suppose
\begin{equation}
 \|K(\theta)\|\le\kappa<1.
\label{eq:v4v16-K-small}
\end{equation}
Put
\begin{equation}
 \mathscr L(\theta)=\sigma\operatorname{Tr}\log(1+K(\theta)),
 \qquad Q=(1+K)^{-1},
\label{eq:v4v16-trace-log}
\end{equation}
where \(\sigma\in\mathbb R\) records the determinant multiplicity and sign.

\begin{theorem}[Exact first and second resolvent responses]
\label{thm:v4v16-resolvent-response}
For parent directions \(v,w\),
\begin{align}
 D\mathscr L[v]
 &=\sigma\operatorname{Tr}(QDK[v]),
\label{eq:v4v16-first-resolvent}\\
 D^2\mathscr L[v,w]
 &=\sigma\operatorname{Tr}\!\left(
 QD^2K[v,w]-QDK[w]QDK[v]
 \right).
\label{eq:v4v16-second-resolvent}
\end{align}
If the displayed operators act through a subspace of dimension at most
\(r\), and
\begin{align}
 \|DK[v]\|&\le M_1\|v\|,
\label{eq:v4v16-DK}\\
 \|D^2K[v,w]\|&\le M_2\|v\|\|w\|,
\label{eq:v4v16-D2K}
\end{align}
then
\begin{align}
 |\mathscr L|&\le |\sigma|r[-\log(1-\kappa)],
\label{eq:v4v16-L-bound}\\
 |D\mathscr L[v]|
 &\le\frac{|\sigma|rM_1}{1-\kappa}\|v\|,
\label{eq:v4v16-DL-bound}\\
 |D^2\mathscr L[v,w]|
 &\le |\sigma|r\left(
 \frac{M_2}{1-\kappa}
 +\frac{M_1^2}{(1-\kappa)^2}
 \right)\|v\|\|w\|.
\label{eq:v4v16-D2L-bound}
\end{align}
\end{theorem}

\begin{proof}
The norm-convergent series
\[
 \log(1+K)=\sum_{n\ge1}\frac{(-1)^{n+1}}nK^n
\]
permits termwise differentiation.  Cyclicity of the trace resums the first
derivative to \eqref{eq:v4v16-first-resolvent}.  Since
\(DQ[w]=-QDK[w]Q\), differentiating once more gives
\eqref{eq:v4v16-second-resolvent}.  The Neumann series gives
\(\|Q\|\le(1-\kappa)^{-1}\).  Use
\(|\operatorname{Tr}A|\le r\|A\|\) and
\(\sum_{n\ge1}\kappa^n/n=-\log(1-\kappa)\) to obtain the three bounds.
\end{proof}

For the exponentiated activity
\begin{equation}
 w(\theta)=e^{-\mathscr L(\theta)}-1,
\label{eq:v4v16-exponentiated-activity}
\end{equation}
the genuine quadratic mark is explicit.

\begin{corollary}[Exact quadratic activity mark]
\label{cor:v4v16-q2}
One has
\begin{align}
 Dw[v]&=-e^{-\mathscr L}D\mathscr L[v],
\label{eq:v4v16-Dw}\\
 D^2w[v,w]
 &=e^{-\mathscr L}
 \bigl(D\mathscr L[v]D\mathscr L[w]
 -D^2\mathscr L[v,w]\bigr).
\label{eq:v4v16-D2w}
\end{align}
Consequently the \(q_2\) mark is bounded by the actual two-insertion term in
\eqref{eq:v4v16-second-resolvent}, the actual \(D^2K\) term, and the product
of the two actual first derivatives in \eqref{eq:v4v16-D2w}.
\end{corollary}

\begin{proof}
Differentiate \eqref{eq:v4v16-exponentiated-activity} twice.  The stated
interpretation follows from substituting
\eqref{eq:v4v16-first-resolvent}--%
\eqref{eq:v4v16-second-resolvent}.
\end{proof}

\begin{firewall}[The square in \(D^2w\) is an identity]
\label{fw:v4v16-q2-identity}
The product \(D\mathscr L[v]D\mathscr L[w]\) in
\eqref{eq:v4v16-D2w} is generated by two derivatives of the same exponential
activity.  It is not an inference that a first-response bound may be squared
to replace \(D^2w\).  The separate term \(D^2\mathscr L\) remains present.
\end{firewall}

\subsection{Hard resolvent locality}
\label{sec:v4v16-resolvent-locality}

Let \(A_0\) be a finite-range operator on a block lattice.  Assume its range
is at most \(R\), its row operator norm is at most \(J\), and
\begin{equation}
 \|A_0^{-1}\|\le\delta^{-1}.
\label{eq:v4v16-hard-gap}
\end{equation}

\begin{lemma}[Combes--Thomas hard-inverse bound]
\label{lem:v4v16-Combes-Thomas}
There are \(c,C>0\), depending only on \(R,J/\delta\) and the local fibre
dimension, such that
\begin{equation}
 \|1_xA_0^{-1}1_y\|
 \le C\delta^{-1}e^{-c\delta d(x,y)/(JR)}.
\label{eq:v4v16-Combes-Thomas}
\end{equation}
The constants are independent of volume.
\end{lemma}

\begin{proof}
Let \(\varphi(z)=d(z,x)\) truncated at a large value and conjugate by
\(W_\mu=e^{-\mu\varphi}\).  Finite range gives
\begin{equation}
 \|W_\mu A_0W_\mu^{-1}-A_0\|
 \le CJ(e^{\mu R}-1).
\label{eq:v4v16-conjugation-error}
\end{equation}
Choose \(\mu=c\delta/(JR)\) so that the right side is at most
\(\delta/2\).  The resolvent identity and a Neumann series then give
\(\|(W_\mu A_0W_\mu^{-1})^{-1}\|\le2\delta^{-1}\).  Taking its \((x,y)\) matrix block gives
\(e^{\mu d(x,y)}\|1_xA_0^{-1}1_y\|\le2\delta^{-1}\).
Choose the truncation radius larger than \(d(x,y)\); the constants do not
depend on that radius.
\end{proof}

In a V13 hard toron chart, \(A_0\) is the gauge, ghost, or matter kinetic
operator after every sub-threshold singular mode has been transferred.  The
hard spectral lower bound supplies \eqref{eq:v4v16-hard-gap}.  The toron
appears only in unitary finite-range parallel transports, so \(J,R\) remain
uniform on the chart.

\subsection{Connected inclusion--exclusion}
\label{sec:v4v16-connected}

Let local perturbations be switched by a finite set \(Y\):
\begin{equation}
 K_Y=\sum_{x\in Y}K_x.
\label{eq:v4v16-local-K}
\end{equation}
For a nonempty finite set \(X\), define the connected trace-log coefficient
\begin{equation}
 \mathscr L^c(X)
 =\sum_{Y\subseteq X}(-1)^{|X|-|Y|}
 \sigma\operatorname{Tr}\log(1+K_Y).
\label{eq:v4v16-Mobius}
\end{equation}

\begin{theorem}[Exact connected-word identity]
\label{thm:v4v16-connected-word}
If \(\|K_Y\|<1\) for every \(Y\subseteq X\), then
\begin{equation}
 \mathscr L^c(X)
 =\sigma\sum_{n\ge|X|}\frac{(-1)^{n+1}}n
 \sum_{\substack{x_1,\ldots,x_n\in X\\
 \{x_1,\ldots,x_n\}=X}}
 \operatorname{Tr}(K_{x_1}\cdots K_{x_n}).
\label{eq:v4v16-connected-word}
\end{equation}
The same identity holds after one or two parent derivatives, with the
derivatives placed on actual factors of each trace word according to the
product rule.
\end{theorem}

\begin{proof}
Insert the trace-log series in \eqref{eq:v4v16-Mobius}.  A word whose set of
labels is \(S\subseteq X\) occurs for precisely the subsets \(Y\) containing
\(S\), with total coefficient
\[
 \sum_{S\subseteq Y\subseteq X}(-1)^{|X|-|Y|}
 =(1-1)^{|X\setminus S|}.
\]
This is zero unless \(S=X\), when it is one.  A word using every element of
\(X\) has length at least \(|X|\), proving
\eqref{eq:v4v16-connected-word}.  Absolute convergence permits one or two
derivatives to pass through the sums; the ordinary product rule places them
on the displayed factors.
\end{proof}

\subsection{Weighted trace-word estimate}
\label{sec:v4v16-weighted-words}

Assume there are nonnegative transition majorants
\(k_m(x,y)\), \(m=0,1,2\), such that a cyclic trace word with an
\(m_i\)-fold derivative on its \(i\)-th factor obeys
\begin{equation}
 \left|operatorname{Tr}
 \prod_{i=1}^nD^{m_i}K_{x_i}
ight|
 \le r_0\prod_{i=1}^nk_{m_i}(x_i,x_{i+1}),
 \qquad x_{n+1}=x_1,
\label{eq:v4v16-transition-majorant}
\end{equation}
and, for some \(a,
ho>0\),
\begin{equation}
 \sup_x\sum_y e^{ad(x,y)}
 \bigl(k_0(x,y)+\rho k_1(x,y)+\rho^2k_2(x,y)\bigr)
 \le\zeta<1.
\label{eq:v4v16-walk-smallness}
\end{equation}

\begin{theorem}[Connected trace-log \(q_2\) decay]
\label{thm:v4v16-q2-decay}
Under \eqref{eq:v4v16-transition-majorant}--%
\eqref{eq:v4v16-walk-smallness}, there are constants
\(C_m(a,\zeta)\), independent of volume, such that
\begin{equation}
 \sup_{\|v_i\|\le1}
 |D^m\mathscr L^c(X)[v_1,\ldots,v_m]|
 \le C_mr_0\rho^{-m}
 e^{-a\mathfrak t(X)},
 \qquad m=0,1,2,
\label{eq:v4v16-connected-decay}
\end{equation}
where \(\mathfrak t(X)\) is the length of a shortest lattice tree spanning
\(X\).  In particular, for connected distinct-site polymers,
\(\mathfrak t(X)\ge|X|-1\).
\end{theorem}

\begin{proof}
Use Theorem~\ref{thm:v4v16-connected-word}.  Every surviving cyclic word
visits all sites of \(X\).  The sum of its step lengths is at least the length
of a spanning tree of \(X\); hence the exponential weights in
\eqref{eq:v4v16-walk-smallness} contribute
\(e^{-a\mathfrak t(X)}\).  Root the cycle at one visited site and sum the
remaining transitions successively by the row bound
\eqref{eq:v4v16-walk-smallness}.  The residual length series is bounded by
\(\sum_{n\ge1}\zeta^n/n\).

For one derivative, choose one actual word factor and pay its
\(\rho k_1\) majorant.  For two derivatives, the product rule gives either
one \(\rho^2k_2\) factor or two distinct \(\rho k_1\) factors.  The factors
of \(n\) and \(n(n-1)\) produced by these choices are summable against
\(\zeta^n/n\).  Removing the mark weights gives \(\rho^{-m}\), proving all
three estimates.
\end{proof}

\begin{corollary}[Exponentiated determinant activity]
\label{cor:v4v16-exponentiated-decay}
If \(|\mathscr L^c(X)|\le c_0e^{-a\mathfrak t(X)}\) with \(c_0\) small,
then
\(w^c(X)=e^{-\mathscr L^c(X)}-1\) and its first two derivatives satisfy a
bound of the form \eqref{eq:v4v16-connected-decay}, with changed constants
and any slightly smaller exponential rate.
\end{corollary}

\begin{proof}
Use \eqref{eq:v4v16-Dw}--\eqref{eq:v4v16-D2w} and
\(|e^{-\mathscr L^c}|\le e^{|\mathscr L^c|}\).  Products of two exponential
decays have at least the same decay.  Absorb bounded prefactors by reducing
the rate slightly.
\end{proof}

\subsection{The fixed-chart determinant packet}
\label{sec:v4v16-packet}

For one toron chart, assemble
\begin{equation}
 \mathscr L_{\rm pkt}^c(X)
 =\frac12\mathscr L_{\rm gauge}^c(X)
 -\mathscr L_{\rm ghost}^c(X)
 -\mathscr L_{\rm matter}^c(X),
\label{eq:v4v16-packet}
\end{equation}
with the appropriate real/complex representation multiplicities.  The
minus signs are retained inside the complete packet before taking its
absolute activity norm.

\begin{theorem}[Marked KP bound for the hard Gaussian packet]
\label{thm:v4v16-packet-KP}
Suppose uniformly on a fixed buffered toron chart:
\begin{enumerate}
\item the gauge, ghost, and matter hard kinetic operators are finite range
and have inverse norm at most \(\delta^{-1}\);
\item their local perturbations and first two coupling-normalized parent
derivatives satisfy \eqref{eq:v4v16-transition-majorant}; and
\item their combined weighted walk norm in
\eqref{eq:v4v16-walk-smallness} is \(O(g)\).
\end{enumerate}
Then for sufficiently small \(g\), the connected exponentiated activities of
\eqref{eq:v4v16-packet} satisfy
\begin{equation}
 \|w_{\rm pkt}\|_{a',2;\rho}\le Cg
\label{eq:v4v16-packet-KP}
\end{equation}
for some \(a'>0\), uniformly in volume.  The estimate supplies the genuine
single-activity \(q_2\) input required by V15.
\end{theorem}

\begin{proof}
Lemma~\ref{lem:v4v16-Combes-Thomas} supplies exponential transition
majorants for every hard inverse.  Local vertices and their normalized
derivatives supply the assumed \(O(g)\) row factor, so
\eqref{eq:v4v16-walk-smallness} holds for small \(g\).
Theorem~\ref{thm:v4v16-q2-decay} and
Corollary~\ref{cor:v4v16-exponentiated-decay} give exponential connected
bounds through two marks for each determinant.  Sum the three signed
trace-log packets first and take the marked norm afterwards.  Connected
polymer counting absorbs a small reduction from \(a\) to \(a'\), yielding
\eqref{eq:v4v16-packet-KP}.
\end{proof}

\begin{firewall}[What the determinant theorem does not include]
\label{fw:v4v16-scope}
Theorem~\ref{thm:v4v16-packet-KP} controls the hard Gaussian trace-log
packet on a fixed buffered toron chart.  It does not control nonlinear gauge
vertices left outside the Hessian, Higgs quartic/Yukawa interaction
polymers, chart-transition constants, sparse or terminal defect sectors, or
dynamical metric fluctuations.  It also assumes the globally normalized
chiral determinant branch before local trace-log expansion.
\end{firewall}

\subsection{V16 result ledger and V17 target}
\label{sec:v4v16-conclusion}

V16 proves exact first and second resolvent formulas for a finite-rank trace
logarithm and an exact quadratic derivative of its exponentiated activity.
M\"obius inclusion--exclusion retains precisely the trace words visiting
every polymer site.  A hard-gap Combes--Thomas estimate and a weighted
closed-walk sum give exponential connected decay through two genuine parent
marks.  Under explicit local smallness hypotheses, the assembled gauge,
ghost, and one-representation matter determinant packet therefore satisfies
the V15 marked KP norm on a fixed toron chart.

V17 should add the first nonlinear bosonic activity without losing the
two-mark structure.  The natural subsystem is the canonically normalized
nonabelian cubic and quartic gauge interaction on the V7 good-field chart.
The target is a connected Gaussian cumulant bound with two background marks,
using the hard covariance once and retaining all Ward-related cubic/seagull
terms in one activity.


\clearpage
\section{V17: the complete cubic--seagull packet and its twice-marked remainder}
\label{sec:v4v17}

V16 proves the required quadratic mark for the hard Gaussian determinant
packet.  The first nonlinear bosonic activity contains canonically normalized
Yang--Mills cubic and quartic vertices.  Its one-loop background Hessian is
not the quartic seagull alone and is not two independent cubic responses.  It
is the second derivative of one complete cumulant packet consisting of the
quartic expectation and the two-cubic connected covariance.

This note extracts that packet exactly on one bounded hard block and proves
that the coupling-normalized twice-marked remainder tends to zero on the
mesoscopic chart \(R_g=g^{-\alpha}\) when \(0<\alpha<1/7\).  The cutoff is
kept explicit.  Extending the result to connected many-block activities and
matching the chart cutoff to the invariant V7 plaquette partition remain the
next tasks.

\subsection{A bounded even hard-block law}
\label{sec:v4v17-block-law}

Let \(\nu\) be a centered finite-dimensional hard Gaussian law in one fixed
toron chart.  Choose an even smooth cutoff \(0\le\chi_R\le1\), supported in
\(\{\|\zeta\|\le R\}\), and define
\begin{equation}
 d\nu_R(\zeta)=
 \frac{\chi_R(\zeta)}{\int\chi_R\,d\nu}\,d\nu(\zeta).
\label{eq:v4v17-cutoff-law}
\end{equation}
The law is even and therefore every odd integrable function has zero
expectation.

Let \(P_3(B,\zeta)\) and \(P_4(B,\zeta)\) be the complete local cubic and
quartic gauge vertices in the background split.  Assume they are homogeneous
of degrees three and four in \(B+\zeta\) and, for \(m=0,1,2\),
\begin{align}
 \|D_B^mP_3\|_\infty&\le C_3R^{3-m},
\label{eq:v4v17-P3-bound}\\
 \|D_B^mP_4\|_\infty&\le C_4R^{4-m}
\label{eq:v4v17-P4-bound}
\end{align}
on \(\|B\|\le R/2\) and the cutoff support.  Put
\begin{equation}
 V_g(B,\zeta)=gP_3(B,\zeta)+g^2P_4(B,\zeta)
\label{eq:v4v17-interaction}
\end{equation}
and
\begin{equation}
 W_g(B)=-\log\mathbb E_{\nu_R}e^{-V_g(B,\zeta)}.
\label{eq:v4v17-effective-block}
\end{equation}

\subsection{Bounded cumulants with two background marks}
\label{sec:v4v17-cumulants}

For bounded random variables \(F_1,\ldots,F_n\), let
\(\kappa_n(F_1,\ldots,F_n)\) denote their joint cumulant.

\begin{lemma}[Uniform bounded-cumulant estimate]
\label{lem:v4v17-bounded-cumulant}
There is a numerical \(C>0\) such that
\begin{equation}
 |\kappa_n(F_1,\ldots,F_n)|
 \le n!C^n\prod_{j=1}^n\|F_j\|_\infty.
\label{eq:v4v17-cumulant-bound}
\end{equation}
Moreover, for a twice differentiable bounded family \(F(B)\),
\begin{align}
 D\kappa_n(F^{\otimes n})[v]
 &=n\kappa_n(DF[v],F^{\otimes(n-1)}),
\label{eq:v4v17-cumulant-D}\\
 D^2\kappa_n(F^{\otimes n})[v,w]
 &=n\kappa_n(D^2F[v,w],F^{\otimes(n-1)})
\notag\\
 &\quad+n(n-1)
 \kappa_n(DF[v],DF[w],F^{\otimes(n-2)}).
\label{eq:v4v17-cumulant-D2}
\end{align}
\end{lemma}

\begin{proof}
The moment--cumulant formula is a finite sum over set partitions.  Taking
absolute values bounds each partition term by
\(\prod_j\|F_j\|_\infty\).  The number of partitions, multiplied by the
factorials in the inversion formula, is at most \(n!C^n\) for a numerical
\(C\), proving \eqref{eq:v4v17-cumulant-bound}.  Joint cumulants are
multilinear and symmetric.  Differentiate their \(n\) argument slots once or
twice; choosing the same slot twice gives the first term in
\eqref{eq:v4v17-cumulant-D2}, while choosing two distinct ordered slots gives
the second.
\end{proof}

If
\begin{equation}
 v_g:=C(gR^3+g^2R^4)<1/4,
\label{eq:v4v17-v-small}
\end{equation}
then the cumulant series converges absolutely:
\begin{equation}
 W_g(B)=
 \sum_{n\ge1}\frac{(-1)^{n+1}}{n!}
 \kappa_n(V_g(B)^{\otimes n}).
\label{eq:v4v17-cumulant-series}
\end{equation}
Indeed, Lemma~\ref{lem:v4v17-bounded-cumulant} dominates it by a geometric
series after reducing the numerical constant in
\eqref{eq:v4v17-v-small}.

\subsection{Exact extraction of coupling degree two}
\label{sec:v4v17-one-loop-packet}

Define the total coupling-degree-at-most-two packet by
\begin{equation}
 W_g^{[2]}(B)
 =g\,\mathbb E P_3(B)
 +g^2\mathbb E P_4(B)
 -\frac{g^2}{2}\kappa_2(P_3(B),P_3(B)).
\label{eq:v4v17-degree-two}
\end{equation}

\begin{theorem}[Complete bosonic one-loop Hessian]
\label{thm:v4v17-one-loop-Hessian}
At \(B=0\), evenness of \(\nu_R\) gives
\begin{align}
 g^{-2}D_B^2W_g^{[2]}(0)[v,w]
 &=\mathbb E D_B^2P_4(0)[v,w]
\notag\\
 &\quad-kappa_2(D_B^2P_3(0)[v,w],P_3(0))
\notag\\
 &\quad-kappa_2(D_BP_3(0)[v],D_BP_3(0)[w]).
\label{eq:v4v17-complete-one-loop}
\end{align}
These three terms are one cubic--seagull packet.
\end{theorem}

\begin{proof}
Because \(P_3\) is cubic, \(D_B^2P_3(0)\) is linear in \(\zeta\), so
\(D_B^2\mathbb EP_3(0)=0\) by evenness.  Differentiate the covariance
\(\kappa_2(P_3,P_3)\) twice.  Lemma~\ref{lem:v4v17-bounded-cumulant} gives
\[
 D_B^2\kappa_2(P_3,P_3)[v,w]
 =2\kappa_2(D_B^2P_3[v,w],P_3)
 +2\kappa_2(D_BP_3[v],D_BP_3[w]).
\]
Insert this identity in \eqref{eq:v4v17-degree-two} and divide by \(g^2\).
\end{proof}

\begin{firewall}[Neither cubic term may be detached]
\label{fw:v4v17-one-loop-assembled}
The last two lines of \eqref{eq:v4v17-complete-one-loop} come from the same
two-cubic covariance.  The first is generated when both marks hit one cubic
vertex and the second when they hit the two actual vertices separately.
Dropping either line breaks the exact second derivative and can break the
background Ward identity after the ghost packet is added.
\end{firewall}

\subsection{The twice-marked nonlinear remainder}
\label{sec:v4v17-remainder}

Put
\begin{equation}
 \mathscr R_g(B)=W_g(B)-W_g^{[2]}(B).
\label{eq:v4v17-remainder}
\end{equation}
Every monomial in its cumulant expansion has total coupling degree at least
three.

\begin{theorem}[Mesoscopic twice-marked remainder]
\label{thm:v4v17-remainder}
There are constants \(C,c>0\) such that, when \(v_g<c\),
\begin{equation}
 g^{-2}\|D_B^2\mathscr R_g(B)\|
 \le
 C\frac{gR^7+g^2R^6}{(1-Cv_g)^3}
\label{eq:v4v17-remainder-bound}
\end{equation}
uniformly for \(\|B\|\le R/2\).  With
\begin{equation}
 R=R_g=g^{-\alpha},
 \qquad0<\alpha<\frac17,
\label{eq:v4v17-mesoscopic-radius}
\end{equation}
the right side is
\begin{equation}
 O(g^{1-7\alpha}+g^{2-6\alpha})=o(1).
\label{eq:v4v17-remainder-small}
\end{equation}
\end{theorem}

\begin{proof}
Expand \eqref{eq:v4v17-cumulant-series} by the number \(n_3\) of cubic and
\(n_4\) of quartic entries.  A twice-differentiated term has coupling and
field-size majorant
\begin{equation}
 C^{n_3+n_4}(n_3+n_4)^2
 g^{n_3+2n_4-2}
 R^{3n_3+4n_4-2}.
\label{eq:v4v17-power-majorant}
\end{equation}
This follows from \eqref{eq:v4v17-P3-bound}--%
\eqref{eq:v4v17-P4-bound} and the two placements in
\eqref{eq:v4v17-cumulant-D2}.  The terms of total coupling degree
\(n_3+2n_4\le2\) are exactly those retained in
\eqref{eq:v4v17-degree-two}.  Among the remaining terms, the smallest
coupling degrees are
\begin{align*}
 (n_3,n_4)=(3,0)&:\quad gR^7,\\
 (n_3,n_4)=(1,1)&:\quad gR^5,\\
 (n_3,n_4)=(0,2)&:\quad g^2R^6.
\end{align*}
Since \(R\ge1\), the middle term is bounded by \(gR^7\).  Sum all additional
entries with the geometric parameter
\(Cv_g=C(gR^3+g^2R^4)\); the two mark placements produce at most the third
power of \((1-Cv_g)^{-1}\).  This proves
\eqref{eq:v4v17-remainder-bound}.  Substitution of
\eqref{eq:v4v17-mesoscopic-radius} gives
\eqref{eq:v4v17-remainder-small}; the restriction
\(\alpha<1/7\) also makes \(v_g\to0\).
\end{proof}

\begin{corollary}[Renormalized single-activity \(q_2\)]
\label{cor:v4v17-renormalized-q2}
Let
\begin{equation}
 w_g^{\rm nl}(B)=e^{-\mathscr R_g(B)}-1.
\label{eq:v4v17-nonlinear-activity}
\end{equation}
The analogous zero- and one-mark estimates obtained from the same cumulant
sum, together with Corollary~\ref{cor:v4v16-q2}, imply
\begin{equation}
 g^{-2}\|D_B^2w_g^{\rm nl}\|
 =O(g^{1-7\alpha}+g^{2-6\alpha})
\label{eq:v4v17-activity-q2}
\end{equation}
on the mesoscopic chart, after reducing \(g\).
\end{corollary}

\begin{proof}
The zero- and one-mark versions of
\eqref{eq:v4v17-power-majorant} have fewer removed field powers and no worse
coupling exponent after the degree-two subtraction.  Hence
\(\mathscr R_g\), \(g^{-1}D_B\mathscr R_g\), and
\(g^{-2}D_B^2\mathscr R_g\) tend to zero with a power bounded by the right
side of \eqref{eq:v4v17-activity-q2}.  Apply the exact exponential derivative
identity \eqref{eq:v4v16-D2w}; its product of first derivatives is of higher
smallness order.
\end{proof}

\subsection{Relation to the background Ward packet}
\label{sec:v4v17-Ward}

If the cutoff law, covariance, and cubic/quartic vertices are obtained from
one background-gauge-equivariant finite block integral, then
Theorem~\ref{thm:v4v5-gauge-invariance} applies before the cumulant
expansion.  In the free-covariance convention, all background-dependent
quadratic gauge vertices are retained in \(P_3,P_4\), and
\eqref{eq:v4v17-complete-one-loop} is combined only with disjoint ghost,
matter, and normalization trace logs from V16.  The complete sum is
transverse.  The irrelevant remainder is the complementary part of the same
exact shell logarithm.

\begin{firewall}[Choose one quadratic bookkeeping convention]
\label{fw:v4v17-no-double-quadratic}
If the V16 gauge operator is instead chosen as the full background-dependent
hard Hessian, its trace logarithm already contains the corresponding
quadratic-background seagull and two-vertex terms.  Those terms must then be
removed from \(P_3,P_4\) before applying V17.  The free-covariance and
background-Hessian conventions are equivalent reorganizations of one shell
integral; adding both would double count the gauge one-loop occurrence.
\end{firewall}

\begin{firewall}[A coordinate cutoff must be matched to invariant defects]
\label{fw:v4v17-cutoff-match}
The even cutoff \(\chi_R\) is an abstract local chart cutoff.  It has not
been proved to coincide with a background-gauge-invariant plaquette partition
whose complement is exactly the V11 repaired-defect gas.  Until such a
partition of unity is constructed, the block estimate cannot be inserted
twice, once as a good-field cumulant and again as a defect correction.
\end{firewall}

\begin{firewall}[The full bosonic marked KP norm is not yet proved]
\label{fw:v4v17-KP-open}
Theorem~\ref{thm:v4v17-remainder} is a single bounded-block estimate.  It
does not prove exponential connected decay across many covariance-coupled
blocks, handle chart transitions, or control sparse and terminal defects.
Thus it supplies the local \(q_2\) input but not yet the complete
\eqref{eq:v4v15-combined-smallness} hypothesis.
\end{firewall}

\subsection{V17 result ledger and V18 target}
\label{sec:v4v17-conclusion}

V17 extracts the complete bosonic coupling-degree-two Hessian: the quartic
seagull, the twice-marked cubic occurrence, and the two separately marked
cubic occurrences are retained in one cumulant packet.  For a bounded hard
block, every higher coupling degree is summed absolutely.  On the
mesoscopic chart \(R_g=g^{-\alpha}\), \(\alpha<1/7\), the
coupling-normalized twice-marked remainder is
\(O(g^{1-7\alpha}+g^{2-6\alpha})\).

V18 should turn this single-block estimate into a connected polymer bound.
The first route to test is a finite-range decomposition of the hard toron
covariance: independent separated blocks would reduce connectedness to a
bounded-degree overlap graph.  If exact finite range is unavailable
uniformly in the toron, the upstream replacement is a Gaussian interpolation
tree using the V16 exponential covariance bound.


\clearpage
\section{V18: finite-range Gaussian forests and the bosonic marked KP norm}
\label{sec:v4v18}

V17 supplies a twice-marked nonlinear remainder on one bounded hard block.
This note proves that the estimate propagates to a connected many-block
polymer gas when the hard Gaussian covariance has exact finite range.  The
joint cumulant is written as a Gaussian interpolation forest.  Finite
covariance range bounds both the degree of every admissible tree and the
number of trees, while V17 bounds the finitely many field derivatives at
each vertex.  The result is an exponential connected activity with two
genuine background marks and hence the V15 marked KP norm.

The exact finite-range hypothesis has not been proved uniformly for the
toron-dependent hard covariance.  V16 gives exponential decay instead.  The
next note must decide whether a toron-uniform finite-range decomposition
exists or whether the exponential tail can be inserted as a separate
summable activity.

\subsection{Local Gaussian dependence graph}
\label{sec:v4v18-dependence}

Let \(\zeta=(\zeta_x)_{x\in\Lambda}\) be a centered Gaussian block field
with covariance \(C\).  Assume
\begin{equation}
 C_{xy}=0
 \qquad\text{when }d(x,y)>R_C.
\label{eq:v4v18-finite-range}
\end{equation}
Let \(u_x(B,\zeta)\) depend only on fields in the radius-\(R_u\)
neighborhood of \(x\).  Define the dependence graph \(G_*\) by joining
\(x,y\) when
\begin{equation}
 d(x,y)\le D_0:=2R_u+R_C.
\label{eq:v4v18-dependence-edge}
\end{equation}
Its maximum degree \(D_*\) depends only on \(D_0\) and the lattice
dimension.

\begin{lemma}[Gaussian independence across a disconnected cut]
\label{lem:v4v18-independence}
If finite sets \(X_1,X_2\) have no \(G_*\)-edge between them, then the
families \(\{u_x:x\in X_1\}\) and \(\{u_y:y\in X_2\}\) are independent.
Consequently their joint cumulant vanishes whenever the induced graph
\(G_*|_X\) is disconnected.
\end{lemma}

\begin{proof}
The Gaussian coordinates on which the first family depends have zero
cross-covariance with those supporting the second by
\eqref{eq:v4v18-finite-range}--%
\eqref{eq:v4v18-dependence-edge}.  Joint Gaussianity makes the coordinate
families independent, hence their measurable functions are independent.
A joint cumulant vanishes when its arguments split into two independent
nonempty families, proving the final statement.
\end{proof}

\subsection{Gaussian cumulant forest formula}
\label{sec:v4v18-forest-formula}

For a local function \(f_x\), let \(\nabla_x^{(k)}f_x\) denote its total
\(k\)-th derivative with respect to the Gaussian variables in its support.

\begin{theorem}[Gaussian tree representation of a joint cumulant]
\label{thm:v4v18-Gaussian-tree}
For distinct sites \(X=\{x_1,\ldots,x_n\}\), the joint cumulant has a
representation
\begin{equation}
 \kappa_n(f_{x_1},\ldots,f_{x_n})
 =\sum_{T\in\operatorname{Tree}(G_*|_X)}
 \int_{[0,1]^{n-1}}
 \mathbb E_{C_T(\mathbf s)}
 \mathcal D_T\prod_{i=1}^nf_{x_i}
 \,d\mathbf s,
\label{eq:v4v18-Gaussian-tree}
\end{equation}
where \(C_T(\mathbf s)\) is a positive semidefinite interpolated covariance
with diagonal blocks equal to those of \(C\), and every edge
\(\{i,j\}\in T\) contributes one contraction
\begin{equation}
 \langle\nabla_{x_i},C_{x_ix_j}\nabla_{x_j}\rangle
\label{eq:v4v18-edge-derivative}
\end{equation}
to \(\mathcal D_T\).  A vertex is differentiated exactly
\(\deg_T(i)\) times.
\end{theorem}

\begin{proof}
Introduce independent copies of the Gaussian field for the \(n\) arguments
and interpolate their cross-covariances from zero to the original covariance.
The fundamental theorem of calculus over the partition lattice expands the
connected part as a sum over forests.  Gaussian differentiation with respect
to one cross-covariance gives exactly
\eqref{eq:v4v18-edge-derivative} by integration by parts.  M\"obius
inversion on set partitions cancels every disconnected forest, leaving the
spanning trees and the positive semidefinite forest-interpolated covariance
matrices.  One derivative endpoint lands at each end of every tree edge, so
vertex \(i\) receives \(\deg_T(i)\) derivatives.  By
\eqref{eq:v4v18-finite-range}, an edge outside \(G_*|_X\) has zero
contraction and may be omitted.
\end{proof}

\begin{proposition}[Finite-degree cumulant bound]
\label{prop:v4v18-cumulant-bound}
Suppose \(\|C_{xy}\|\le C_0\) and, for \(0\le k\le D_*\),
\begin{equation}
 \sup_x\|\nabla_x^{(k)}f_x\|_\infty
 \le\varepsilon A^k.
\label{eq:v4v18-field-derivative-bound}
\end{equation}
Then for every nonempty \(X\),
\begin{equation}
 |\kappa(f_x:x\in X)|
 \le\varepsilon^{|X|}
 (2D_*C_0A^2)^{|X|-1},
\label{eq:v4v18-cumulant-decay}
\end{equation}
and the left side is zero unless \(G_*|_X\) is connected.
\end{proposition}

\begin{proof}
The vanishing statement is Lemma~\ref{lem:v4v18-independence}.  In the tree
formula, every edge supplies one covariance factor \(C_0\), and the total
number of field derivatives is \(2(|X|-1)\).  Equation
\eqref{eq:v4v18-field-derivative-bound} therefore bounds each tree term by
\(\varepsilon^{|X|}(C_0A^2)^{|X|-1}\).  A graph of maximum degree \(D_*\)
has at most \((2D_*)^{|X|-1}\) spanning trees: apply the matrix-tree theorem
and bound the nonzero Laplacian eigenvalues by \(2D_*\).  Sum the tree terms.
\end{proof}

\subsection{Two parent marks in the Gaussian forest}
\label{sec:v4v18-marks}

Let \(D_B\) be the background derivative and set the coupling-normalized
mark scale
\begin{equation}
 \rho_g=g^{-1}.
\label{eq:v4v18-mark-scale}
\end{equation}
For \(m=0,1,2\), assume
\begin{equation}
 \sup_x\sup_{0\le k\le D_*}
 \rho_g^mA^{-k}
 \|\nabla_x^{(k)}D_B^mu_x\|_\infty
 \le\varepsilon_g.
\label{eq:v4v18-local-marked-bound}
\end{equation}

\begin{theorem}[Twice-marked Gaussian cumulant]
\label{thm:v4v18-marked-cumulant}
Under \eqref{eq:v4v18-local-marked-bound}, for \(m=0,1,2\),
\begin{equation}
 \|D_B^m\kappa(u_x:x\in X)\|
 \le C_m|X|^m\rho_g^{-m}
 \varepsilon_g^{|X|}
 (2D_*C_0A^2)^{|X|-1}.
\label{eq:v4v18-marked-cumulant}
\end{equation}
For \(m=2\), the derivative placements are exactly:
\begin{enumerate}
\item two marks on one actual vertex \(u_x\); or
\item one mark on each of two distinct actual vertices.
\end{enumerate}
\end{theorem}

\begin{proof}
Background differentiation commutes with the Gaussian expectation because
the reference covariance is fixed in this bookkeeping convention.
Differentiate the product in
\eqref{eq:v4v18-Gaussian-tree}.  There are at most \(|X|\) choices for one
mark and \(|X|^2\) ordered choices for two marks.  Equation
\eqref{eq:v4v18-local-marked-bound} replaces the marked vertex bound by
\(\rho_g^{-1}\varepsilon_gA^k\) or
\(\rho_g^{-2}\varepsilon_gA^k\).  All other field-derivative and tree
factors are unchanged from Proposition~\ref{prop:v4v18-cumulant-bound}.
The two-mark alternatives are the product rule and are disjoint exactly as
in Lemma~\ref{lem:v4v15-two-mark-identity}.
\end{proof}

\begin{firewall}[A covariance derivative is a separate convention]
\label{fw:v4v18-covariance-convention}
Theorem~\ref{thm:v4v18-marked-cumulant} uses a background-independent
Gaussian covariance, matching the free-covariance convention of V17.  If
the covariance depends on the background, its derivative produces actual
resolvent insertions and must be assigned to the V16 trace-log convention.
It may not also be hidden inside \(D_Bu_x\).
\end{firewall}

\subsection{Application to the V17 renormalized local activity}
\label{sec:v4v18-application}

For each elementary interaction block \(x\), retain the original bounded
vertex
\begin{equation}
 V_{g,x}=gP_{3,x}+g^2P_{4,x},
 \qquad
 u_x=e^{-V_{g,x}}-1.
\label{eq:v4v18-local-u}
\end{equation}
For switches \(t=(t_x)\), define the exact connected coefficient
\begin{equation}
 \Omega_g(X)
 =-\left.
 \prod_{x\in X}\frac{\partial}{\partial t_x}
 \log\mathbb E
 \prod_{y\in\Lambda}(1+t_yu_y)
 \right|_{t=0}.
\label{eq:v4v18-Omega}
\end{equation}
Let \(\mathcal T_{\le2}\) retain total powers \(g^p\) with \(p\le2\), and
put
\begin{equation}
 \overline\Omega_g(X)
 =(1-\mathcal T_{\le2})\Omega_g(X).
\label{eq:v4v18-renormalized-Omega}
\end{equation}
The sum of \(\mathcal T_{\le2}\Omega_g(X)\) is the complete marginal packet;
the projection is performed after connected Gaussian averaging.  Since each
site label carries at least one interaction power, only \(|X|\le2\) can
contribute to that packet.

\begin{proposition}[Renormalized connected activity bound]
\label{prop:v4v18-renormalized-connected}
Let \(R_g=g^{-\alpha}\), \(0<\alpha<1/7\), and define
\begin{align}
 v_g&=C(gR_g^3+g^2R_g^4),
\label{eq:v4v18-vg}\\
 \varepsilon_g&=C(gR_g^7+g^2R_g^6).
\label{eq:v4v18-epsilon}
\end{align}
Under the finite-range covariance and fixed-order regulator derivative
bounds of this section, there is \(C_*\) such that, for \(m=0,1,2\),
\begin{equation}
 \rho_g^m\|D_B^m\overline\Omega_g(X)\|
 \le C_*\varepsilon_g
 (C_*v_g)^{|X|-1}
 e^{-a\mathfrak t(X)}.
\label{eq:v4v18-renormalized-bound}
\end{equation}
The coefficient vanishes when \(G_*|_X\) is disconnected.
\end{proposition}

\begin{proof}
Expand every \(u_x\) in its absolutely convergent series in \(V_{g,x}\)
and then use the Gaussian cumulant forest formula.  M\"obius differentiation
in \eqref{eq:v4v18-Omega} forces every label in \(X\) to occur at least
once.  The finite-degree tree estimate of
Proposition~\ref{prop:v4v18-cumulant-bound} supplies one factor \(C_*v_g\)
for each additional label.  For any fixed \(a>0\), every surviving edge
also obeys
\[
 \|C_{xy}\|\le C_0e^{aD_0}e^{-ad(x,y)}
 \qquad(d(x,y)\le D_0).
\]
Multiplying this inequality over a spanning tree and minimizing its total
edge length gives \(e^{-a\mathfrak t(X)}\), with the harmless factor
\(e^{aD_0}\) absorbed into \(C_*\).

After \(\mathcal T_{\le2}\) is removed, every surviving monomial has total
coupling degree at least three.  With two background marks and the factor
\(\rho_g^2=g^{-2}\), the exact power census is the one in
\eqref{eq:v4v17-power-majorant}; its worst terms are
\(gR_g^7\) and \(g^2R_g^6\).  Zero and one marks have no worse normalized
power.  The sums over extra occurrences at already visited labels are
geometric in \(Cv_g\), and the two mark placements are precisely the two
cases of Theorem~\ref{thm:v4v18-marked-cumulant}.  This proves
\eqref{eq:v4v18-renormalized-bound}.  Disconnected vanishing is
Lemma~\ref{lem:v4v18-independence}.
\end{proof}

\begin{theorem}[Finite-range bosonic marked KP estimate]
\label{thm:v4v18-bosonic-KP}
Let \(0<\alpha<1/7\).  Under the exact finite-range covariance hypothesis
\eqref{eq:v4v18-finite-range} and the local regulator bounds above, there are
\(a'>0\) and \(g_0>0\) such that
\begin{equation}
 \|\overline\Omega_g\|_{a',2;\rho_g}
 \le C\varepsilon_g
\label{eq:v4v18-bosonic-KP}
\end{equation}
for \(0<g<g_0\), uniformly in volume.  In particular the nonlinear
bosonic connected remainder supplies its part of
\eqref{eq:v4v15-combined-smallness}.
\end{theorem}

\begin{proof}
Only \(G_*\)-connected sets contribute.  The number of connected
cardinality-\(n\) sets containing a fixed root is at most
\((eD_*)^{n-1}\) by Lemma~\ref{lem:v4v7-connected-count}.  Insert
\eqref{eq:v4v18-renormalized-bound}.  Since both \(v_g\) and
\(\varepsilon_g\) tend to zero, choose \(g\) small and reduce the spatial
rate from \(a\) to \(a'\) so the resulting series is geometric.  Its sum is
at most \(C\varepsilon_g\).
\end{proof}

\begin{corollary}[Determinant plus nonlinear bosonic hard gas]
\label{cor:v4v18-combined-hard}
On a fixed toron chart satisfying both V16 and V18 hypotheses,
\begin{equation}
 \|w_{\rm pkt}^{\rm det}\sqcup\overline\Omega_g\|_{a,2;\rho_g}
 \le Cg+C\varepsilon_g.
\label{eq:v4v18-combined-hard}
\end{equation}
Thus the complete gauge/ghost/matter Gaussian packet and the renormalized
nonlinear gauge remainder obey the V15 two-mark theorem for sufficiently
small \(g\), before Higgs--Yukawa and defect activities are added.
\end{corollary}

\begin{proof}
Apply Theorem~\ref{thm:v4v16-packet-KP},
Theorem~\ref{thm:v4v18-bosonic-KP}, and the marked-norm composition
Lemma~\ref{lem:v4v15-composition}.
\end{proof}

\begin{firewall}[Finite-range toron covariance remains an open input]
\label{fw:v4v18-finite-range-open}
No theorem here constructs an exact finite-range decomposition of the
toron-dependent hard covariance with background-Ward covariance, uniform
range, and uniform derivative constants.  Replacing the V16 exponential
tail by zero without an explicit remainder would change the measure.  Hence
Theorem~\ref{thm:v4v18-bosonic-KP} is not yet available on the full toron
parent.
\end{firewall}

\subsection{V18 result ledger and V19 target}
\label{sec:v4v18-conclusion}

V18 proves the exact Gaussian interpolation-tree representation of a joint
local cumulant.  Under finite covariance range, only bounded-degree spanning
trees survive.  The V17 single-block estimates therefore yield exponential
connected decay through two genuine background marks and the full bosonic
marked KP norm.  Combined with V16, this closes the fixed-chart hard
determinant plus nonlinear gauge packet under the finite-range hypothesis.

V19 should remove that hypothesis.  Starting from the V16
Combes--Thomas bound, truncate the hard covariance at a range \(L_g\), retain
the tail as an explicit interpolation occurrence, and choose \(L_g\) so the
tail marked norm is smaller than the V17 power remainder.  Gauge covariance
and positivity of the covariance split must be checked before using the
truncated field as an independent Gaussian shell.


\clearpage
\section{V19: a positive finite-range Gaussian split and the cutoff interface}
\label{sec:v4v19}

V18 used an exact finite-range covariance as a sufficient input for the
many-block forest estimate.  Entrywise truncation of the exponentially
decaying V16 covariance is not an admissible construction: it need not
preserve positivity, and a difference of two positive matrices need not be
positive.  This note replaces truncation by an exact conditional Gaussian
split of a local precision operator.  The fluctuation covariance is positive,
has strict finite range, and the residual covariance is positive because it
is the covariance of the conditional mean.

Tracing the hypotheses upstream also reveals a cutoff mismatch.  The
normalized V17 law \(d\nu_R\) is not Gaussian, so a Gaussian forest identity
cannot be applied to that conditioned law.  The correct order is to retain
the invariant good-field cutoff as an observable, perform the conditional
Gaussian split with respect to the unconditioned parent, and estimate the
cutoff transition inside the finite-range fluctuation integral.  The exact
split is proved here.  The required tail-sensitive forest norm is isolated
at the end.

\subsection{A typographical correction to V17}
\label{sec:v4v19-erratum}

In \eqref{eq:v4v17-complete-one-loop}, the second and third occurrences
printed as \(-kappa_2\) mean \(-\kappa_2\).  This is a missing TeX escape,
not a change of the formula.  The two terms are the same joint cumulants
defined in Section~\ref{sec:v4v17-cumulants} and used in its proof.

\subsection{Dirichlet interiors and their Schur complement}
\label{sec:v4v19-Schur}

Let \(\Lambda\) be a finite graph and
\(\mathcal H_\Lambda=\bigoplus_{x\in\Lambda}\mathcal H_x\).  Let
\(A=A^*>0\) be a precision operator of range \(r_A\):
\begin{equation}
 A_{xy}=0\qquad\text{if }d(x,y)>r_A.
\label{eq:v4v19-local-precision}
\end{equation}
Choose disjoint interiors \(I_b\subset\Lambda\), \(b\in\mathcal B\), such
that
\begin{equation}
 d(I_b,I_{b'})>r_A\qquad(b\ne b'),
\label{eq:v4v19-separated-interiors}
\end{equation}
and put \(I=\bigsqcup_b I_b\), \(S=\Lambda\setminus I\).  Relative to
\(\mathcal H_I\oplus\mathcal H_S\), write
\begin{equation}
 A=
 \begin{pmatrix}
 A_{II}&A_{IS}\\
 A_{SI}&A_{SS}
 \end{pmatrix},
 \qquad
 \mathcal S_A=A_{SS}-A_{SI}A_{II}^{-1}A_{IS}.
\label{eq:v4v19-Schur-complement}
\end{equation}
Strict positivity of \(A\) implies \(A_{II}>0\) and
\(\mathcal S_A>0\).

Let \(\iota_I:\mathcal H_I\to\mathcal H_\Lambda\) be extension by zero and
define
\begin{align}
 \Gamma&=\iota_I A_{II}^{-1}\iota_I^*,
\label{eq:v4v19-Gamma}\\
 Hq&=
 \binom{-A_{II}^{-1}A_{IS}q}{q},
 \qquad q\in\mathcal H_S.
\label{eq:v4v19-harmonic-extension}
\end{align}

\begin{theorem}[Exact positive finite-range covariance split]
\label{thm:v4v19-positive-split}
With \(C=A^{-1}\),
\begin{equation}
 C=\Gamma+H\mathcal S_A^{-1}H^*.
\label{eq:v4v19-positive-split}
\end{equation}
Both summands are positive semidefinite.  Moreover,
\begin{equation}
 \Gamma_{xy}=0
 \quad\text{unless \(x,y\) lie in the same \(I_b\)}.
\label{eq:v4v19-Gamma-range}
\end{equation}
Thus \(\Gamma\) has range at most
\(R_{\mathcal B}:=\sup_b\operatorname{diam}(I_b)\).

Equivalently, if
\[
 \xi\sim N(0,\Gamma),\qquad
 q\sim N(0,\mathcal S_A^{-1})
\]
are independent and \(\eta=Hq\), then
\begin{equation}
 \zeta=\xi+\eta\sim N(0,C).
\label{eq:v4v19-random-split}
\end{equation}
\end{theorem}

\begin{proof}
Condition \eqref{eq:v4v19-separated-interiors} and locality
\eqref{eq:v4v19-local-precision} make
\(A_{II}=\bigoplus_b A_{I_bI_b}\).  Hence its inverse has the same block
diagonal form, which proves \eqref{eq:v4v19-Gamma-range}.

The block \(LDU\) factorization
\begin{equation}
 A=
 \begin{pmatrix}1&0\\ A_{SI}A_{II}^{-1}&1\end{pmatrix}
 \begin{pmatrix}A_{II}&0\\0&\mathcal S_A\end{pmatrix}
 \begin{pmatrix}1&A_{II}^{-1}A_{IS}\\0&1\end{pmatrix}
\label{eq:v4v19-LDU}
\end{equation}
shows that \(\mathcal S_A>0\).  Inverting
\eqref{eq:v4v19-LDU} and multiplying the blocks gives
\[
 A^{-1}=
 \begin{pmatrix}A_{II}^{-1}&0\\0&0\end{pmatrix}
 +
 \binom{-A_{II}^{-1}A_{IS}}{1}
 \mathcal S_A^{-1}
 \begin{pmatrix}-A_{SI}A_{II}^{-1}&1\end{pmatrix},
\]
which is \eqref{eq:v4v19-positive-split}.  The first summand is
\(\Gamma\); the second is \(H\mathcal S_A^{-1}H^*\).  Their positivity is
now manifest.  Independent centered Gaussian vectors add their covariance
operators, so \eqref{eq:v4v19-random-split} follows.
\end{proof}

\begin{corollary}[Exact determinant and expectation factorization]
\label{cor:v4v19-factorization}
One has
\begin{equation}
 \det A=\det A_{II}\det\mathcal S_A.
\label{eq:v4v19-determinant}
\end{equation}
For every integrable \(F\),
\begin{equation}
 \mathbb E_C F(\zeta)
 =
 \mathbb E_{\mathcal S_A^{-1}}
 \mathbb E_{\Gamma}F(Hq+\xi).
\label{eq:v4v19-expectation-factorization}
\end{equation}
No determinant or covariance remainder is discarded.
\end{corollary}

\begin{proof}
Take determinants in \eqref{eq:v4v19-LDU}; the two triangular factors have
determinant one.  Equation \eqref{eq:v4v19-expectation-factorization} is
Theorem~\ref{thm:v4v19-positive-split} applied to the Gaussian law, first for
bounded continuous \(F\), then for every integrable \(F\) by truncation.
\end{proof}

\begin{firewall}[This is not entrywise covariance truncation]
\label{fw:v4v19-no-truncation}
The finite-range covariance \(\Gamma\) is a Dirichlet conditional covariance
selected by the local precision \(A\).  It is not obtained by setting distant
entries of \(C\) to zero.  The residual
\(H\mathcal S_A^{-1}H^*\) is positive because it is the covariance of the
conditional mean.  These facts are what permit the two summands to be
realized as independent Gaussian fields.
\end{firewall}

\subsection{Conditional application of the V18 forest}
\label{sec:v4v19-conditional-forest}

Let a local factor \(F_x(B,\zeta)\) depend on \(\zeta\) in a fixed
neighborhood of \(x\).  For fixed \(q\), put
\begin{equation}
 u_x(B,q,\xi)=F_x(B,Hq+\xi)-1.
\label{eq:v4v19-conditional-u}
\end{equation}
The \(\xi\)-dependence remains local, and
\eqref{eq:v4v19-Gamma-range} gives it a bounded-degree dependence graph.
Consequently the V18 Gaussian tree formula applies to
\(\mathbb E_\Gamma\prod_x(1+u_x)\) for every fixed \(q\), provided the
required finitely many \(\xi\)-derivatives are integrable uniformly in the
forest interpolation parameters.

\begin{proposition}[Moment form of the finite-range forest bound]
\label{prop:v4v19-shell-before-residual}
Let \(X\) have cardinality \(n\).  Assume the interiors have uniformly
bounded diameter and local cardinality and, for \(m=0,1,2\) and
\(0\le k\le D_*\),
\begin{equation}
 \sup_q\sup_x\sup_{\mathbf s,T}
 \rho_g^m A_0^{-k}
 \left\|
 \nabla_\xi^{(k)}D_B^m u_x(B,q,\xi)
 \right\|_{L^n(\gamma_{\Gamma,T,\mathbf s})}
 \le\epsilon_{g,n}.
\label{eq:v4v19-Ln-input}
\end{equation}
Here \(\gamma_{\Gamma,T,\mathbf s}\) denotes any Gaussian law appearing in
the V18 forest interpolation.  Then the conditional connected coefficient
obeys, uniformly in \(q\),
\begin{equation}
 \rho_g^m
 \left|D_B^m\kappa_\Gamma(u_x:x\in X)\right|
 \le
 C_m n^m\epsilon_{g,n}^{\,n}
 (C_*A_0^2)^{n-1},
 \qquad m\le2,
\label{eq:v4v19-conditional-bound}
\end{equation}
and vanishes unless the finite-range dependence graph on \(X\) is connected.
In particular, the V18 conclusion follows when the left side of
\eqref{eq:v4v19-Ln-input} is bounded uniformly in \(n\), as it is under the
V18 \(L^\infty\) hypothesis.
\end{proposition}

\begin{proof}
Use the V18 tree formula.  A vertex receives
\(k=\deg_T(x)\le D_*\) field derivatives.  Holder's inequality with exponent
\(n\) bounds the expectation of the product of its \(n\) differentiated
vertex factors by \(\epsilon_{g,n}^n A_0^{2(n-1)}\), after the normalized
background marks have been assigned.  Every covariance edge is bounded by
the uniform Dirichlet covariance norm, and the matrix-tree estimate in
Proposition~\ref{prop:v4v18-cumulant-bound} completes the count.  Background
differentiation gives the same one-vertex or two-vertex mark placements as
Theorem~\ref{thm:v4v18-marked-cumulant}.  Finite-range independence gives
the vanishing statement.
\end{proof}

\begin{firewall}[A fixed \(L^p\) tail bound does not close the forest]
\label{fw:v4v19-Lp-caution}
For polymers of unbounded cardinality, Holder requires the \(L^n\) norm in
\eqref{eq:v4v19-Ln-input}; a fixed \(L^{D_*}\) estimate cannot replace it.
Although the original finite-range covariance makes sufficiently separated
local sigma fields independent, a forest-interpolated Gaussian law need not
factor the color classes of the dependence graph.  Thus a coloring argument
does not repair the exponent loss.  V20 must either sum the cutoff
occurrences before Holder, prove a uniform conditional mixing estimate, or
supply an all-order regulator norm.
\end{firewall}

\subsection{Where the invariant cutoff belongs}
\label{sec:v4v19-cutoff-order}

Let \(\chi_{R,x}\) be a smooth local factor constructed from the invariant
V7 plaquette partition after tree gauge and assigned to one owned block.
For the nonlinear vertex \(V_{g,x}\), the good factor is
\begin{equation}
 F_x(B,\zeta)=\chi_{R,x}(B,\zeta)e^{-V_{g,x}(B,\zeta)}.
\label{eq:v4v19-good-factor}
\end{equation}
The exact algebraic decomposition
\begin{equation}
 F_x-1
 =
 \chi_{R,x}(e^{-V_{g,x}}-1)
 -(1-\chi_{R,x})
\label{eq:v4v19-cutoff-decomposition}
\end{equation}
separates the perturbative interaction from the cutoff transition.  On the
support of \(\chi_{R,x}\), the first term is bounded by the V17 mesoscopic
power census.  The second is not small in supremum norm, but it is supported
where a local field leaves the core chart.

\begin{lemma}[Conditional cutoff tail]
\label{lem:v4v19-cutoff-tail}
Suppose the Gaussian fluctuation has uniformly bounded diagonal covariance
and the conditional mean obeys \(\|Hq\|\le R/2\) on the owned block.  Then
\begin{equation}
 \left\|1-\chi_{R,x}(Hq+\xi)\right\|_{L^p(\gamma_\Gamma)}
 \le C_p e^{-cR^2/p}.
\label{eq:v4v19-cutoff-tail}
\end{equation}
The same estimate holds for each fixed derivative order when derivatives of
\(\chi_R\) are supported in the transition annulus and have polynomial
\(R\)-bounds.  With \(R=R_g=g^{-\alpha}\), this is smaller than every power
of \(g\).
\end{lemma}

\begin{proof}
The transition support is contained in
\(\{\|\xi\|\ge cR\}\) when \(\|Hq\|\le R/2\).  A finite-dimensional centered
Gaussian with uniformly bounded covariance has
\(\mathbb P(\|\xi\|\ge cR)\le C e^{-c'R^2}\).  Since the cutoff and each fixed
derivative are bounded by a polynomial in \(R\), Holder's inequality gives
\eqref{eq:v4v19-cutoff-tail}; the polynomial is absorbed after decreasing
\(c\).
\end{proof}

\begin{firewall}[Correction of the V18 application claim]
\label{fw:v4v19-V18-correction}
The V18 forest theorem is correct for a genuine Gaussian reference with
finite-range covariance.  Its stated application to the normalized V17 law
does not follow, because multiplying a Gaussian law by \(\chi_R\) and
renormalizing destroys Gaussianity.  Corollary~\ref{cor:v4v18-combined-hard}
must therefore be read as conditional on an explicit Gaussian-reference
cutoff factorization.  Equations
\eqref{eq:v4v19-expectation-factorization} and
\eqref{eq:v4v19-good-factor} give the correct factorization order, but
\eqref{eq:v4v19-cutoff-tail} still has to be propagated through the full
forest without replacing it by a supremum bound.
\end{firewall}

\begin{firewall}[Coarse means leaving the chart become defects]
\label{fw:v4v19-coarse-defects}
Estimate \eqref{eq:v4v19-cutoff-tail} assumes the harmonic coarse field
\(Hq\) remains in the inner chart.  Configurations violating that condition
cannot be discarded or charged again to the same good factor.  They must be
assigned, through an exact invariant partition of unity, to the repairable,
sparse-promotion, or terminal sectors already separated in V11--V13.
\end{firewall}

\begin{firewall}[Background dependence of the precision]
\label{fw:v4v19-background-precision}
The split above is used in the V17 free-covariance convention, where \(A\)
is fixed under \(D_B\).  If \(A=A(B)\), derivatives of
\(A_{II}^{-1}\), \(H\), and \(\mathcal S_A^{-1}\) are genuine covariance
responses.  They must be retained as the V16 resolvent packet and may not
also be counted as marks on \(u_x\).
\end{firewall}

\subsection{V19 result ledger and compulsory V20 target}
\label{sec:v4v19-conclusion}

V19 proves an exact positivity-preserving replacement for covariance
truncation.  A local precision operator split into separated Dirichlet
interiors has a positive, strict finite-range fluctuation covariance; the
residual is the positive covariance of the harmonic conditional mean.
Determinants and expectations factor exactly.  This supplies the finite-range
Gaussian object required by V18 without altering the measure.

The cutoff audit prevents a false closure.  The V17 normalized cutoff law is
not Gaussian.  The invariant cutoff must remain inside the integrand of the
unconditioned Gaussian factorization, where its transition has a
stretched-Gaussian \(L^p\) cost.  The present \(L^{D_*}\) shortcut is
explicitly rejected for arbitrarily large polymers.  V20 is compulsory:
after reading the newest Yang--Mills and Navier--Stokes notes, it must prove a
valid tail-sensitive finite-range forest estimate, or retreat to a
conditional mixing theorem with hypotheses that can actually be verified.

\clearpage
\section{V20: companion audit and exact cutoff-tail polymerization}
\label{sec:v4v20}

V19 showed that a fixed \(L^p\) estimate cannot simply be multiplied over
polymers of unbounded size.  This compulsory checkpoint reads the newest
companion notes and changes the order of proof.  The Yang--Mills programme
retains actual replacement occurrences before forming a positive square;
the Navier--Stokes programme retains the complete signed transfer before
removing a probability mark.  Here the analogous object is the complete set
of actual cutoff-tail occurrences in one conditional Gaussian shell.

The V19 Dirichlet split makes the fluctuation field a product of independent
block Gaussians.  This note uses that product structure directly.  The
factor expansion reorganizes exactly into connected overlap polymers.
Among \(r\) actual tail occurrences, a greedy independent set has at least
\(r/(D+1)\) members, and their probabilities multiply.  This gives a
stretched-Gaussian payment proportional to the number of tail occurrences,
without a Holder exponent growing with the polymer.

\subsection{Compulsory companion audit}
\label{sec:v4v20-audit}

The exact files inspected were
\begin{enumerate}
\item
\texttt{Renewal\_Yang\_Mills\_Mass\_Gap\_Research\_Notes\_V41.tex},
\(771054\) bytes, SHA--256
\[
 \texttt{B99C89E805430FD4A6C01B7AB5DD2A87A67FF4878C3FB8E3D198585952BAE59D}.
\]
This file is byte-identical to the adjacent V40 file and its final internal
chapter is V40; no distinct V41 mathematical appendix was present at the
time of audit.
\item
\texttt{Renewal\_NS\_Stability\_Research\_Notes\_V31.tex},
\(887537\) bytes, SHA--256
\[
 \texttt{F1121B62238AD5491BB7CF03635EA9934942EDF573ED8FAAA9EE79E1D38A6696}.
\]
\end{enumerate}

The YM lesson is that an exact polarized square can expose two genuine
marks, but the identity \(\|S^*S\|=\|S\|^2\) supplies no small norm by
itself.  A local occurrence also cannot pay mass stored in an unrelated
spectator bank.  The NS lesson is that a flat probability mark is exact,
but unmarking restores the full first-moment weight; heat smoothing alone
does not erase it.  Both forbid the same shortcut in V19: estimating one
cutoff occurrence in a fixed \(L^p\) norm and then reusing that estimate at
every forest vertex would detach the probability payment from its actual
occurrences.

\subsection{Independent base blocks and the overlap graph}
\label{sec:v4v20-overlap}

Let \((\xi_b)_{b\in\mathcal B}\) be independent random variables.  In the
V19 application they are the independent Dirichlet Gaussian interiors.
For each factor label \(x\in\Lambda\), let
\(N(x)\subset\mathcal B\) be the set of base blocks on which \(u_x\)
depends.  Define the overlap graph \(G\) on \(\Lambda\) by
\begin{equation}
 x\sim y
 \quad\Longleftrightarrow\quad
 N(x)\cap N(y)\ne\varnothing.
\label{eq:v4v20-overlap-graph}
\end{equation}
Assume its maximum degree is at most \(D\), uniformly in volume.

For a nonempty \(G\)-connected set \(Y\), define
\begin{equation}
 z(Y)=\mathbb E\prod_{x\in Y}u_x.
\label{eq:v4v20-bare-polymer}
\end{equation}
Two polymers are compatible when they are disjoint and there is no
\(G\)-edge with one endpoint in each of them.

\begin{theorem}[Exact product-block polymerization]
\label{thm:v4v20-exact-polymerization}
For every finite \(\Lambda\),
\begin{equation}
 \mathbb E\prod_{x\in\Lambda}(1+u_x)
 =
 \sum_{\substack{\mathcal Y\ {\rm finite}\\
                  \text{pairwise compatible}}}
 \prod_{Y\in\mathcal Y}z(Y),
\label{eq:v4v20-exact-polymerization}
\end{equation}
where the empty family contributes one.  Thus the conditional shell factor
is exactly a hard-core polymer partition function; no cumulant or covariance
tail has been omitted.
\end{theorem}

\begin{proof}
Expand the product on the left:
\[
 \mathbb E\prod_x(1+u_x)
 =
 \sum_{S\subset\Lambda}\mathbb E\prod_{x\in S}u_x.
\]
Decompose \(S\) into the connected components
\(Y_1,\ldots,Y_k\) of \(G|_S\).  Distinct components depend on disjoint
sets of independent base variables by
\eqref{eq:v4v20-overlap-graph}.  Their products are therefore independent,
and
\[
 \mathbb E\prod_{x\in S}u_x
 =\prod_{j=1}^k
   \mathbb E\prod_{x\in Y_j}u_x
 =\prod_{j=1}^kz(Y_j).
\]
The map from \(S\) to its compatible family of connected components is a
bijection, proving \eqref{eq:v4v20-exact-polymerization}.
\end{proof}

\subsection{A probability payment for every actual tail occurrence}
\label{sec:v4v20-tail-payment}

Write
\begin{equation}
 u_x=p_x+d_x.
\label{eq:v4v20-pd-split}
\end{equation}
The factor \(p_x\) is perturbative.  The factor \(d_x\) is a cutoff
transition supported on an event \(E_x\).  Suppose
\begin{equation}
 |p_x|\le b,\qquad
 |d_x|\le L\,\boldsymbol 1_{E_x},\qquad
 \mathbb P(E_x)\le\tau\le1.
\label{eq:v4v20-local-tail-input}
\end{equation}
If \(x\not\sim y\), the events \(E_x,E_y\) depend on disjoint independent
base blocks.

\begin{lemma}[Independent subset of tail occurrences]
\label{lem:v4v20-independent-tail-set}
Every finite \(R\subset\Lambda\) contains a \(G\)-independent set
\(J\subset R\) satisfying
\begin{equation}
 |J|\ge\frac{|R|}{D+1}.
\label{eq:v4v20-independent-size}
\end{equation}
Consequently
\begin{equation}
 \mathbb P\left(\bigcap_{x\in R}E_x\right)
 \le\tau^{|R|/(D+1)}.
\label{eq:v4v20-joint-tail}
\end{equation}
\end{lemma}

\begin{proof}
Choose a vertex of \(R\), place it in \(J\), and remove it together with
its at most \(D\) neighbors.  Repeat until no vertex remains.  Each chosen
vertex removes at most \(D+1\) vertices, giving
\eqref{eq:v4v20-independent-size}.  The events indexed by \(J\) are
functions of pairwise disjoint families of independent base variables and
are independent.  Since
\(\cap_{x\in R}E_x\subset\cap_{x\in J}E_x\),
\[
 \mathbb P\left(\bigcap_{x\in R}E_x\right)
 \le\prod_{x\in J}\mathbb P(E_x)
 \le\tau^{|J|}
 \le\tau^{|R|/(D+1)}.
\]
\end{proof}

\begin{proposition}[Tail-sensitive bare activity bound]
\label{prop:v4v20-tail-activity}
Under \eqref{eq:v4v20-local-tail-input}, every nonempty
\(G\)-connected \(Y\) obeys
\begin{equation}
 |z(Y)|
 \le\delta^{|Y|},
 \qquad
 \delta:=b+L\tau^{1/(D+1)}.
\label{eq:v4v20-tail-activity}
\end{equation}
\end{proposition}

\begin{proof}
Expand \(\prod_{x\in Y}(p_x+d_x)\) by the subset
\(R\subset Y\) at which \(d_x\) is chosen.  Equations
\eqref{eq:v4v20-local-tail-input} and
\eqref{eq:v4v20-joint-tail} give
\[
 \left|
 \mathbb E
 \prod_{x\in Y\setminus R}p_x
 \prod_{x\in R}d_x
 \right|
 \le
 b^{|Y|-|R|}
 \left(L\tau^{1/(D+1)}\right)^{|R|}.
\]
Sum over \(R\subset Y\) and use the binomial theorem.
\end{proof}

\subsection{Two genuine background marks}
\label{sec:v4v20-marks}

Let \(D_B\) be a background derivative and \(\rho>0\) its normalization.
Assume that for \(m=0,1,2\),
\begin{align}
 \rho^m|D_B^mp_x|&\le b\,,
\label{eq:v4v20-marked-p}\\
 \rho^m|D_B^md_x|&\le
 L\,\boldsymbol 1_{E_x}.
\label{eq:v4v20-marked-d}
\end{align}
The event on the right of \eqref{eq:v4v20-marked-d} is the same actual
cutoff transition for all three derivative orders.

\begin{theorem}[Marked cutoff-tail KP bound]
\label{thm:v4v20-marked-tail-KP}
Under \eqref{eq:v4v20-marked-p}--%
\eqref{eq:v4v20-marked-d}, for \(m=0,1,2\),
\begin{equation}
 \rho^m|D_B^mz(Y)|
 \le C_m|Y|^m\delta^{|Y|},
\label{eq:v4v20-marked-z}
\end{equation}
with \(\delta\) from \eqref{eq:v4v20-tail-activity}.  Hence for every
\(a>0\) there is \(c_a(D)>0\) such that, if
\(\delta<c_a(D)\),
\begin{equation}
 \sup_{x\in\Lambda}
 \sum_{\substack{Y\ni x\\Y\ {\rm connected}}}
 e^{a|Y|}
 \max_{0\le m\le2}\rho^m|D_B^mz(Y)|
 \le C_{a,D}\delta.
\label{eq:v4v20-marked-KP}
\end{equation}
The constants are independent of \(|\Lambda|\).
\end{theorem}

\begin{proof}
Differentiate the product in \eqref{eq:v4v20-bare-polymer}.  One mark has
\(|Y|\) placements.  Two marks either land twice on one actual factor or
once on each of two distinct actual factors, giving at most
\(|Y|^2\) ordered placements.  After multiplying by \(\rho^m\), each
differentiated factor satisfies the same \(b\) or
\(L\boldsymbol 1_{E_x}\) bound as its unmarked version.  The proof of
Proposition~\ref{prop:v4v20-tail-activity} therefore applies to each
placement and proves \eqref{eq:v4v20-marked-z}.

By Lemma~\ref{lem:v4v7-connected-count}, the number of connected
cardinality-\(n\) sets containing a fixed vertex is at most
\((eD)^{n-1}\).  Therefore the left side of
\eqref{eq:v4v20-marked-KP} is bounded by
\[
 C\sum_{n\ge1}n^2(eD)^{n-1}e^{an}\delta^n.
\]
This power series is at most \(C_{a,D}\delta\) when
\(e^{a+1}D\delta\) is sufficiently small.
\end{proof}

\begin{firewall}[The two marks are not a squared tail estimate]
\label{fw:v4v20-no-squared-tail}
When both marks hit distinct cutoff factors, the proof retains the joint
event containing both actual transition occurrences and then selects an
independent subset from the complete transition set.  It never replaces
that joint event by two copies of a one-mark norm.  This is the cutoff
analogue of the actual-occurrence rule in the YM audit.
\end{firewall}

\subsection{Exact surgery of the good-field factor}
\label{sec:v4v20-surgery}

\begin{proposition}[Exact bounded cutoff surgery]
\label{prop:v4v20-cutoff-surgery}
Let \(0\le\chi_{R,x}\le1\) equal one on the inner invariant chart and vanish
outside the outer chart.  Define
\begin{align}
 G_x&=\chi_{R,x}e^{-V_{g,x}},
\label{eq:v4v20-physical-good}\\
 P_x&=e^{-\chi_{R,x}V_{g,x}},
\label{eq:v4v20-bounded-extension}\\
 p_x&=P_x-1,\qquad
 d_x=G_x-P_x.
\label{eq:v4v20-surgery-pd}
\end{align}
Then
\begin{equation}
 G_x=1+p_x+d_x
\label{eq:v4v20-surgery-identity}
\end{equation}
exactly, and \(d_x=0\) wherever \(\chi_{R,x}=1\).  If
\begin{equation}
 |V_{g,x}|\le v_g
 \quad\text{on \(\{\chi_{R,x}>0\}\)},
\label{eq:v4v20-bounded-vertex}
\end{equation}
then, for \(v_g\le1\),
\begin{equation}
 |p_x|\le e v_g,\qquad
 |d_x|\le 2e\,\boldsymbol 1_{\{\chi_{R,x}\ne1\}}.
\label{eq:v4v20-surgery-bounds}
\end{equation}
\end{proposition}

\begin{proof}
Equation \eqref{eq:v4v20-surgery-identity} is the definition.  If
\(\chi=1\), both \(G_x\) and \(P_x\) equal \(e^{-V_{g,x}}\), so
\(d_x=0\).  The elementary inequality
\(|e^{-w}-1|\le e^{|w|}|w|\) gives the first bound in
\eqref{eq:v4v20-surgery-bounds}.  On the transition event, both exponentials
have modulus at most \(e^{v_g}\); outside it their difference is zero.  This
gives the second bound.
\end{proof}

For the conditional V19 Dirichlet shell, assume
\(\|Hq\|\le R/2\) on the owned neighborhood.  Lemma
\ref{lem:v4v19-cutoff-tail} gives
\begin{equation}
 \tau_g
 :=
 \sup_x\mathbb P_\Gamma(\chi_{R,x}(Hq+\xi)\ne1)
 \le Ce^{-cR^2}.
\label{eq:v4v20-tau}
\end{equation}
Equations \eqref{eq:v4v20-surgery-bounds} and
\eqref{eq:v4v20-tail-activity} then give
\begin{equation}
 \delta_g
 \le
 C v_g+C\exp\left[-\frac{cR^2}{D+1}\right].
\label{eq:v4v20-delta}
\end{equation}
For \(R=R_g=g^{-\alpha}\), the cutoff part is smaller than every power of
\(g\).  The factor \(D+1\) is fixed by block geometry and does not grow with
the polymer.

The same conclusion holds through two background marks if derivatives of
\(p_x,d_x\) satisfy \eqref{eq:v4v20-marked-p}--%
\eqref{eq:v4v20-marked-d}.  Polynomial losses from derivatives of the
smooth transition are absorbed by the stretched exponential in
\eqref{eq:v4v20-tau}.

\begin{firewall}[Marginal subtraction must precede the marked norm]
\label{fw:v4v20-marginal-first}
For the raw cubic--quartic factor, the coupling-normalized second derivative
of \(p_x\) need not satisfy \eqref{eq:v4v20-marked-p}: its coupling degrees
one and two are the marginal packet extracted in V17.  Taking the marked
norm before assembling and subtracting that packet would destroy the odd
cubic cancellation and the cubic--seagull Ward combination.  Therefore
Theorem~\ref{thm:v4v20-marked-tail-KP} closes the cutoff transfer only after
the complete degree-at-most-two connected coefficient has been removed.
It does not by itself perform that subtraction.
\end{firewall}

\begin{firewall}[Invariant ownership remains required]
\label{fw:v4v20-invariant-ownership}
The product \(\prod_xG_x\) represents the physical good sector only after
the plaquette transition factors have an exact gauge-invariant ownership
rule and their complement is assigned once to the V11 repairable, sparse,
or terminal sectors.  V20 proves the conditional probability estimate; it
does not manufacture that ownership map or duplicate a defect payment.
\end{firewall}

\subsection{V20 result ledger and V21 target}
\label{sec:v4v20-conclusion}

V20 records the mandatory companion fingerprints and implements their common
lesson.  For a finite-range Dirichlet shell, local factors are functions of
independent base blocks.  Their exact expansion is a hard-core gas of
connected overlap polymers.  Every set of cutoff-tail occurrences contains
a linearly large independent subset, so the tail probability is paid once
for each actual occurrence up to the fixed factor \(D+1\).  This yields a
volume-uniform two-mark KP bound with
\(\delta=b+L\tau^{1/(D+1)}\), and the mesoscopic Gaussian tail is
superpolynomially small.

The remaining obstruction is now algebraic rather than probabilistic.  The
complete coupling-degree-at-most-two cubic--seagull coefficient must be
extracted from the exact connected logarithm of
\eqref{eq:v4v20-exact-polymerization} before the marked norm is taken.
V21 should prove that extraction at the polymer level, show that cutoff
derivatives enter only the stretched-tail species, and then combine the
renormalized perturbative and cutoff-tail gases without duplicating the V16
quadratic packet.

\clearpage
\section{V21: exact marginal extraction before the marked norm}
\label{sec:v4v21}

V20 controls the probability cost of cutoff transitions, but the raw
coupling-normalized second background derivative still contains the
degree-one and degree-two marginal terms.  Those terms cannot be estimated
separately: the quartic seagull and the two-cubic covariance form one
connected Ward packet.  This note performs the required algebra before any
marked norm is taken.

An auxiliary parameter \(s\) records coupling degree while the physical
coupling \(g\), chart radius \(R_g\), and cutoff functions are held fixed.
The tail-free logarithm is differentiated twice at \(s=0\).  The result is
the complete spatial cubic--seagull packet.  The difference between the
physical cutoff logarithm and the tail-free logarithm contains at least one
actual cutoff transition and is controlled by V20.  Finally, replacing the
cutoff by one inside the extracted packet produces only transition-supported
terms, so these too belong to the stretched-tail remainder.

\subsection{Two activity species and an exact logarithmic split}
\label{sec:v4v21-two-species}

For fixed coarse field \(q\), abbreviate the conditional expectation with
respect to the V19 Dirichlet fluctuation by \(\mathbb E_\Gamma\).  Introduce
\begin{equation}
 V_x(s)
 =
 s gP_{3,x}+s^2g^2P_{4,x},
\label{eq:v4v21-scaled-vertex}
\end{equation}
where all functions are evaluated at \((B,Hq+\xi)\).  Let
\begin{align}
 G_x(s)&=\chi_xe^{-V_x(s)},
\label{eq:v4v21-Gs}\\
 P_x(s)&=e^{-\chi_xV_x(s)},
\label{eq:v4v21-Ps}\\
 p_x(s)&=P_x(s)-1,
&
 d_x(s)&=G_x(s)-P_x(s).
\label{eq:v4v21-pd}
\end{align}
The cutoff \(\chi_x=\chi_{R_g,x}\) is independent of the bookkeeping
parameter \(s\), though it may depend on \(B,Hq+\xi\).

Define
\begin{align}
 Z_G(s)&=\mathbb E_\Gamma\prod_xG_x(s),
\label{eq:v4v21-ZG}\\
 Z_p(s)&=\mathbb E_\Gamma\prod_xP_x(s),
\label{eq:v4v21-Zp}\\
 W_G(s)&=-\log Z_G(s),
&
 W_p(s)&=-\log Z_p(s).
\label{eq:v4v21-W}
\end{align}
Whenever the V20 KP bounds hold, both logarithms use the cluster branch
continued from unit activities.

\begin{lemma}[Every logarithmic difference term contains a tail occurrence]
\label{lem:v4v21-tail-log-difference}
In the common KP domain,
\begin{equation}
 W_G(s)=W_p(s)+\Delta_{\rm tail}(s),
\qquad
 \Delta_{\rm tail}(s)
 =
 -\int_0^1
 \frac{d}{dt}
 \log Z(p(s)+t\,d(s))\,dt.
\label{eq:v4v21-log-split}
\end{equation}
Every monomial in the convergent activity expansion of
\(\Delta_{\rm tail}\) contains at least one \(d\)-activity.
\end{lemma}

\begin{proof}
Let
\[
 Z_t=Z(p(s)+t\,d(s))
 =
 \mathbb E_\Gamma
 \prod_x(1+p_x(s)+t\,d_x(s)).
\]
Then \(Z_0=Z_p(s)\) and \(Z_1=Z_G(s)\).  The fundamental theorem of calculus
gives
\[
 -\log Z_1+\log Z_0
 =-\int_0^1\frac{d}{dt}\log Z_t\,dt.
\]
The KP expansion is absolutely and locally uniformly convergent in \(t\),
so it may be differentiated term by term.  The derivative selects one
actual \(d\)-occurrence; all remaining factors are \(p+td\).  Hence every
resulting monomial contains at least one \(d\).
\end{proof}

\begin{lemma}[Marked Lipschitz bound with polynomial mark loss]
\label{lem:v4v21-cluster-Lipschitz}
Let \(z_t=z^0+t h\), \(0\le t\le1\), lie in a common unmarked KP ball
of radius \(\eta<1\).  Write \(\|\cdot\|_{a,0}\) for the corresponding
unmarked activity norm.  Suppose, for some \(M\ge1\) and \(\delta>0\),
\begin{align}
 \sup_{0\le t\le1}\max_{1\le m\le2}
 \|\rho^mD_B^mz_t\|_{a,0}&\le M,
\label{eq:v4v21-path-marks}\\
 \max_{0\le m\le2}
 \|\rho^mD_B^mh\|_{a,0}&\le\delta.
\label{eq:v4v21-difference-marks}
\end{align}
Then, after decomposing the logarithms by their cluster-support union,
\begin{equation}
 \left\|\log Z(z^1)-\log Z(z^0)\right\|_{a',2;\rho}
 \le C_{\eta,a-a'}M^2\delta,
 \qquad0<a'<a.
\label{eq:v4v21-cluster-Lipschitz}
\end{equation}
\end{lemma}

\begin{proof}
Use the absolutely convergent Mayer expansion
\[
 \log Z(z)
 =
 \sum_{n\ge1}\frac1{n!}
 \sum_{Y_1,\ldots,Y_n}
 \varphi^T(Y_1,\ldots,Y_n)
 \prod_{j=1}^nz(Y_j).
\]
Differentiate along \(z_t\).  Exactly one distinguished vertex carries
\(h\).  The tree-graph inequality bounds \(|\varphi^T|\) by a sum over
incompatibility trees.  Root the tree at the distinguished vertex and sum
its branches with the unmarked KP inequality.  Each branch costs at most
\(\eta\), and the loss from \(a\) to \(a'\) sums the support unions.

For one background mark, the product rule puts it either on the distinguished
\(h\)-vertex or on one ordinary \(z_t\)-vertex.  For two marks, they are
both on one actual vertex or on two distinct actual vertices.  Equations
\eqref{eq:v4v21-path-marks}--%
\eqref{eq:v4v21-difference-marks} bound these cases respectively by
\(C\delta\), \(CM\delta\), or \(CM^2\delta\).  The same rooted-tree sums
converge because marks change no incompatibility edge.  Integrate over
\(0\le t\le1\) to prove
\eqref{eq:v4v21-cluster-Lipschitz}.
\end{proof}

Let \(z_p\) be the V20 connected-overlap bare activities formed only from
\(p_x\), and let \(z_G\) be those formed from \(p_x+d_x\).  Expanding their
difference shows that every term of \(z_G-z_p\) has at least one \(d_x\).
The V20 independent-subset proof, with the one- and two-mark product rules,
gives
\begin{equation}
 \max_{0\le m\le2}
 \|\rho_g^mD_B^m(z_G-z_p)\|_{a,0}
 \le L_g\tau_g^{1/(D+1)}
\label{eq:v4v21-tail-activity-difference}
\end{equation}
for a factor \(L_g\) polynomial in \(g^{-1}\) and \(R_g\).  The analogous
marked norms of the interpolation \(z_p+t(z_G-z_p)\) have a polynomial
bound \(M_g\).  Lemma~\ref{lem:v4v21-cluster-Lipschitz}, with \(M_g^2\)
absorbed into \(L_g\), yields
\begin{equation}
 \|\Delta_{\rm tail}\|_{a',2;\rho_g}
 \le
 C L_g\tau_g^{1/(D+1)}.
\label{eq:v4v21-tail-log-bound}
\end{equation}
For \(R_g=g^{-\alpha}\), polynomial \(L_g\), and
\(\tau_g\le Ce^{-cR_g^2}\), this is smaller than every power of \(g\).

\subsection{The exact degree-two coefficient}
\label{sec:v4v21-degree-two}

Put
\begin{equation}
 A_3=\sum_x\chi_xP_{3,x},
 \qquad
 A_4=\sum_x\chi_xP_{4,x}.
\label{eq:v4v21-A34}
\end{equation}
The product in \eqref{eq:v4v21-Zp} is exactly
\begin{equation}
 \prod_xP_x(s)
 =
 \exp\left[-sgA_3-s^2g^2A_4\right].
\label{eq:v4v21-product-exponential}
\end{equation}

\begin{theorem}[Complete conditional cubic--seagull packet]
\label{thm:v4v21-complete-packet}
The coupling-degree-at-most-two Taylor polynomial of \(W_p\) is
\begin{equation}
 W_p^{[2]}(s)
 =
 sg\,\mathbb E_\Gamma A_3
 +
 s^2g^2\,\mathbb E_\Gamma A_4
 -
 \frac{s^2g^2}{2}
 \kappa_{2,\Gamma}(A_3,A_3).
\label{eq:v4v21-complete-packet}
\end{equation}
In particular,
\begin{align}
 g^{-2}D_B^2[s^2]W_p^{[2]}
 &=
 D_B^2\mathbb E_\Gamma A_4
 -\frac12D_B^2
 \kappa_{2,\Gamma}(A_3,A_3)
\notag\\
 &=
 D_B^2\mathbb E_\Gamma A_4
 -
 \kappa_{2,\Gamma}(D_B^2A_3,A_3)
 -
 \kappa_{2,\Gamma}(D_BA_3,D_BA_3),
\label{eq:v4v21-two-mark-packet}
\end{align}
where the final line is evaluated on equal background directions; the
polarized bilinear form has the corresponding two ordered first-derivative
placements.  Thus the seagull, twice-marked cubic occurrence, and two
separately marked cubic occurrences are extracted before a norm is taken.
\end{theorem}

\begin{proof}
Equation \eqref{eq:v4v21-product-exponential} and the bounded chart imply
that \(Z_p(s)\) is analytic near zero.  The cumulant expansion gives
\[
 -\log\mathbb E_\Gamma e^{-sgA_3-s^2g^2A_4}
 =
 sg\,\mathbb E_\Gamma A_3
 +
 s^2g^2\mathbb E_\Gamma A_4
 -
 \frac{s^2g^2}{2}\kappa_{2,\Gamma}(A_3,A_3)
 +O(s^3).
\]
This proves \eqref{eq:v4v21-complete-packet}.  The conditional covariance is
fixed in the free-covariance convention, so background differentiation acts
only on the two argument slots.  Multilinearity and symmetry of
\(\kappa_2\) give
\[
 \frac12D_B^2\kappa_2(A_3,A_3)
 =
 \kappa_2(D_B^2A_3,A_3)
 +\kappa_2(D_BA_3,D_BA_3)
\]
on equal directions, proving \eqref{eq:v4v21-two-mark-packet}.
\end{proof}

\begin{corollary}[Finite-range locality of the marginal packet]
\label{cor:v4v21-packet-locality}
Write \(A_3=\sum_xa_{3,x}\).  If the supports of
\(a_{3,x}\) and \(a_{3,y}\) use disjoint V19 Dirichlet blocks, then
\begin{equation}
 \kappa_{2,\Gamma}(a_{3,x},a_{3,y})=0.
\label{eq:v4v21-covariance-locality}
\end{equation}
Hence \eqref{eq:v4v21-complete-packet} is a sum of uniformly local one-site
and overlapping-pair coefficients.
\end{corollary}

\begin{proof}
The two variables are functions of disjoint families of independent
Dirichlet block fields and are independent.  Their covariance is zero.
The overlap graph has volume-independent degree.
\end{proof}

\subsection{The tail-free irrelevant remainder}
\label{sec:v4v21-irrelevant}

Decompose the convergent Mayer logarithm of \(Z_p(s)\) by its connected
support union:
\begin{equation}
 W_p(s)=\sum_X W_p(X;s).
\label{eq:v4v21-local-log}
\end{equation}
Let \(\mathcal T_{\le2}\) retain powers \(s^0,s^1,s^2\) and define
\begin{equation}
 \overline W_p(X)
 =
 \left.(1-\mathcal T_{\le2})W_p(X;s)\right|_{s=1}.
\label{eq:v4v21-renormalized-local-log}
\end{equation}
Summing \(\mathcal T_{\le2}W_p(X;s)\) over \(X\) gives exactly
\eqref{eq:v4v21-complete-packet}.

\begin{theorem}[Renormalized perturbative marked norm]
\label{thm:v4v21-renormalized-KP}
Assume the V17 local polynomial bounds on the outer chart, the V19
finite-range Dirichlet geometry, and
\[
 R_g=g^{-\alpha},\qquad0<\alpha<\frac17.
\]
Let
\begin{align}
 v_g&=C(gR_g^3+g^2R_g^4),
\label{eq:v4v21-vg}\\
 \epsilon_g&=
 C\frac{gR_g^7+g^2R_g^6}{(1-Cv_g)^3}.
\label{eq:v4v21-epsilon}
\end{align}
For \(g\) sufficiently small, the tail-free coefficients satisfy
\begin{equation}
 \|\overline W_p\|_{a,2;\rho_g}
 \le C\epsilon_g,
 \qquad \rho_g=g^{-1},
\label{eq:v4v21-renormalized-KP}
\end{equation}
uniformly in volume and in coarse fields for which the outer-chart bounds
hold.  Cutoff derivatives supported where \(\chi_x\ne1\) are excluded from
the left side and assigned to the tail correction below.
\end{theorem}

\begin{proof}
Expand each \(p_x(s)\) in powers of the bounded vertex
\(\chi_xV_x(s)\), and use the absolutely convergent Mayer expansion for
\eqref{eq:v4v21-local-log}.  A cluster monomial has a definite number
\(n_3\) of cubic occurrences and \(n_4\) of quartic occurrences.  The
projector in \eqref{eq:v4v21-renormalized-local-log} removes exactly those
with \(n_3+2n_4\le2\).  After two background marks and multiplication by
\(\rho_g^2=g^{-2}\), every surviving monomial obeys the V17 majorant
\[
 C^{n_3+n_4}(n_3+n_4)^2
 g^{n_3+2n_4-2}
 R_g^{3n_3+4n_4-2}.
\]
Its three smallest possibilities are bounded by
\(C(gR_g^7+g^2R_g^6)\); all additional occurrences sum with parameter
\(Cv_g\).

Every bare polymer is connected in the bounded-degree overlap graph, and
the incompatibility graph of every Mayer cluster connects their union.
Root the union at a fixed site and apply the tree-graph inequality, followed
by Lemma~\ref{lem:v4v7-connected-count}.  The factors \(1/n!\) in the
exponential and Mayer series cancel the ordering counts, while a small loss
in the spatial weight sums the rooted trees.  This gives the denominator
\((1-Cv_g)^3\) from the occurrence count and two mark placements, and proves
\eqref{eq:v4v21-renormalized-KP}.  A derivative falling on
\(\chi_x\) is supported on the cutoff transition and is deliberately not
estimated in this tail-free bound.
\end{proof}

\begin{firewall}[Tail-free differentiation is a defined projection]
\label{fw:v4v21-tail-free-derivative}
The last sentence of the proof does not discard cutoff derivatives.  Define
the tail-free background derivative by differentiating \(P_{3,x},P_{4,x}\)
while holding \(\chi_x\) fixed.  The difference from the physical derivative
contains \(D_B\chi_x\) or \(D_B^2\chi_x\), hence is supported where the
cutoff changes.  It is included in
\(\Delta_{\rm tail}^{\rm all}\) below.  This assignment makes the two
derivative classes disjoint and exhaustive.
\end{firewall}

\subsection{Removing the cutoff from the marginal packet}
\label{sec:v4v21-core-packet}

Let \(W_{\rm core}^{[2]}\) denote
\eqref{eq:v4v21-complete-packet} with every \(\chi_x\) replaced by one on
the owned good blocks.  Expanding
\(\chi_x=1-(1-\chi_x)\) shows that
\begin{equation}
 W_p^{[2]}-W_{\rm core}^{[2]}
\label{eq:v4v21-low-tail-difference}
\end{equation}
is a finite sum in which every term contains at least one actual transition
factor \(1-\chi_x\), \(D_B\chi_x\), or \(D_B^2\chi_x\).  The covariance term
contains at most two spatial occurrences, so Lemma
\ref{lem:v4v20-independent-tail-set} and the conditional Gaussian tail give
\begin{equation}
 \|W_p^{[2]}-W_{\rm core}^{[2]}\|_{a,2;\rho_g}
 \le
 C L_g e^{-cR_g^2/(D+1)}
\label{eq:v4v21-low-tail-bound}
\end{equation}
for a polynomial \(L_g\).

Define the complete remainder
\begin{equation}
 \mathscr R_{\rm good}
 =
 \sum_X\overline W_p(X)
 +\Delta_{\rm tail}
 +(W_p^{[2]}-W_{\rm core}^{[2]}),
\label{eq:v4v21-complete-remainder}
\end{equation}
including in \(\Delta_{\rm tail}\) the physical-minus-tail-free derivative
terms from Firewall~\ref{fw:v4v21-tail-free-derivative}.

\begin{corollary}[Conditional good-sector marked bound]
\label{cor:v4v21-good-sector}
Under the hypotheses of Theorems
\ref{thm:v4v20-marked-tail-KP} and
\ref{thm:v4v21-renormalized-KP},
\begin{equation}
 \|\mathscr R_{\rm good}\|_{a',2;\rho_g}
 \le
 C\epsilon_g+
 C L_g e^{-cR_g^2/(D+1)}.
\label{eq:v4v21-good-sector-bound}
\end{equation}
The physical conditional logarithm is exactly
\begin{equation}
 W_G(1)=W_{\rm core}^{[2]}+\mathscr R_{\rm good}.
\label{eq:v4v21-exact-final-split}
\end{equation}
\end{corollary}

\begin{proof}
Combine \eqref{eq:v4v21-log-split},
\eqref{eq:v4v21-renormalized-local-log}, and
\eqref{eq:v4v21-low-tail-difference}; this proves the exact identity.
Apply \eqref{eq:v4v21-renormalized-KP},
\eqref{eq:v4v21-tail-log-bound}, and
\eqref{eq:v4v21-low-tail-bound} and use the triangle inequality only after
the three disjoint classes have been formed.
\end{proof}

\begin{firewall}[Relation to the V16 quadratic convention]
\label{fw:v4v21-V16-convention}
The core packet \(W_{\rm core}^{[2]}\) belongs to the running marginal
parent only in the V17 free-covariance convention.  If the V16 trace-log
already uses the full background-dependent gauge Hessian, the corresponding
quadratic seagull and two-cubic response must be removed from
\(W_{\rm core}^{[2]}\).  Equation
\eqref{eq:v4v21-exact-final-split} does not authorize adding both.
\end{firewall}

\begin{firewall}[What remains conditional]
\label{fw:v4v21-conditional}
The proof assumes a uniformly local invariant ownership of the factors
\(\chi_x\), uniform outer-chart polynomial bounds for every admissible coarse
field, and the V19 fixed-background precision convention.  The V11 sparse
and terminal sectors, chart transitions across toron strata, Higgs--Yukawa
activities, the hypercharge ultraviolet direction, and dynamical metric
fluctuations are not included.
\end{firewall}

\subsection{V21 result ledger and V22 target}
\label{sec:v4v21-conclusion}

V21 proves the algebraic order required by the V20 audit.  The physical
conditional logarithm is split exactly into a tail-free perturbative
logarithm and a logarithmic difference containing at least one actual cutoff
transition.  The complete degree-two coefficient of the first logarithm is
the spatial cubic--seagull packet, including both placements of two genuine
background marks.  Its irrelevant complement satisfies the V17 power
census in a volume-uniform marked polymer norm, while every cutoff-dependent
correction retains a stretched-Gaussian payment.

The next obstruction is the coarse-field condition used in the tail bound:
\(\|Hq\|\le R_g/2\) on every owned good block.  V22 should construct an exact
coarse-good/coarse-defect partition compatible with the V7 and V11
ownership, then determine whether the harmonic extension \(H\) preserves
the plaquette threshold without a scale-dependent loss.  If that pointwise
route fails, the upstream replacement is a local coarse large-deviation
polymer estimate, not a global all-block good event.

\clearpage
\section{V22: multiscale positive Schur shells and first-failure ownership}
\label{sec:v4v22}

V21 controls one finite-range Dirichlet shell when its harmonic coarse mean
lies in the inner chart.  Treating that condition as a global all-block
event would repeat the V7 error.  This note iterates the V19 Schur split.
Every covariance shell is positive and finite range, the remaining
covariance is positive, and all determinant normalizations telescope.
A first-failure identity then assigns each chart exit to one scale only.

The construction is finite-dimensional and exact.  It does not assume that
entrywise covariance truncation is positive.  It also gives a simultaneous
Gaussian excursion estimate for a set of coarse blocks using only an
operator upper bound on the residual covariance; independence of the
residual coordinates is not required.

\subsection{Recursive Schur data}
\label{sec:v4v22-recursion}

Set \(\Lambda_0=\Lambda\), \(\mathcal H_0=\mathcal H_\Lambda\), and
\(A_0=A>0\).  At level \(j\), choose separated interiors
\[
 I_j=\bigsqcup_{b\in\mathcal B_j}I_{j,b}\subset\Lambda_j,
 \qquad
 \Lambda_{j+1}=S_j:=\Lambda_j\setminus I_j,
\]
so distinct interiors are farther apart than the range of \(A_j\).
Relative to \(I_j\oplus S_j\), define
\begin{align}
 A_{j+1}
 &:=
 (A_j)_{S_jS_j}
 -(A_j)_{S_jI_j}(A_j)_{I_jI_j}^{-1}(A_j)_{I_jS_j},
\label{eq:v4v22-next-precision}\\
 \Gamma_j
 &:=
 \iota_{I_j}(A_j)_{I_jI_j}^{-1}\iota_{I_j}^*,
\label{eq:v4v22-shell-covariance}\\
 H_jq
 &:=
 \binom{-(A_j)_{I_jI_j}^{-1}(A_j)_{I_jS_j}q}{q}.
\label{eq:v4v22-level-extension}
\end{align}
Let
\begin{equation}
 E_0=1,\qquad E_j=H_0H_1\cdots H_{j-1}\quad(j\ge1).
\label{eq:v4v22-composed-extension}
\end{equation}

\begin{theorem}[Exact multiscale positive covariance decomposition]
\label{thm:v4v22-multiscale-split}
For every terminal level \(J\),
\begin{equation}
 A_0^{-1}
 =
 \sum_{j=0}^{J-1}E_j\Gamma_jE_j^*
 +
 E_JA_J^{-1}E_J^*.
\label{eq:v4v22-multiscale-covariance}
\end{equation}
Every summand is positive semidefinite.  Consequently there are independent
centered Gaussian fields
\[
 \xi_j\sim N(0,\Gamma_j)\quad(0\le j<J),
 \qquad q_J\sim N(0,A_J^{-1})
\]
such that
\begin{equation}
 \zeta
 =
 \sum_{j=0}^{J-1}E_j\xi_j+E_Jq_J
 \sim N(0,A_0^{-1}).
\label{eq:v4v22-multiscale-field}
\end{equation}
The determinant normalization satisfies
\begin{equation}
 \det A_0
 =
 \left[
 \prod_{j=0}^{J-1}
 \det (A_j)_{I_jI_j}
 \right]\det A_J.
\label{eq:v4v22-telescoping-determinant}
\end{equation}
\end{theorem}

\begin{proof}
The one-step covariance identity
\[
 A_j^{-1}
 =
 \Gamma_j+H_jA_{j+1}^{-1}H_j^*
\]
is Theorem~\ref{thm:v4v19-positive-split} applied at level \(j\).
Substitute it successively into the residual term.  After \(J\) steps this
gives \eqref{eq:v4v22-multiscale-covariance}.  Positivity is preserved by
\(K\mapsto E_jKE_j^*\).  Independent Gaussian summation then proves
\eqref{eq:v4v22-multiscale-field}.

The one-step determinant identity is
\[
 \det A_j
 =
 \det(A_j)_{I_jI_j}\det A_{j+1}.
\]
Multiplication over \(0\le j<J\) cancels every intermediate determinant and
gives \eqref{eq:v4v22-telescoping-determinant}.
\end{proof}

\subsection{Locality under Schur elimination}
\label{sec:v4v22-locality}

Suppose \(A_j\) has range \(r_j\), and every \(I_{j,b}\) has diameter at
most \(\ell_j\).

\begin{proposition}[Deterministic range recursion]
\label{prop:v4v22-range-recursion}
The next precision has range at most
\begin{equation}
 r_{j+1}\le 2r_j+\ell_j.
\label{eq:v4v22-range-recursion}
\end{equation}
The map \(H_j\) has propagation radius at most \(r_j+\ell_j\): the value of
\(H_jq\) in \(I_{j,b}\) depends only on boundary coordinates within distance
\(r_j\) of that same interior.  Hence every shell covariance
\(E_j\Gamma_jE_j^*\) has finite range, bounded by a deterministic expression
in \(r_0,\ldots,r_j\) and \(\ell_0,\ldots,\ell_j\).
\end{proposition}

\begin{proof}
The first term in \eqref{eq:v4v22-next-precision} has range \(r_j\).
A matrix element of the Schur correction between \(x,y\in S_j\) can be
nonzero only if there are \(u,v\) in one common interior \(I_{j,b}\) with
\[
 d(x,u)\le r_j,\qquad d(y,v)\le r_j.
\]
The inverse on \(I_{j,b}\) may be dense inside that block, but
\(d(u,v)\le\ell_j\).  Thus \(d(x,y)\le2r_j+\ell_j\), proving
\eqref{eq:v4v22-range-recursion}.  The same support observation applied to
\eqref{eq:v4v22-level-extension} proves the propagation statement.
Composing finitely many finite-propagation maps with the block-diagonal
finite-range \(\Gamma_j\) preserves finite propagation.
\end{proof}

\begin{firewall}[Finite range grows with scale]
\label{fw:v4v22-growing-range}
Proposition~\ref{prop:v4v22-range-recursion} proves strict finite range at
every fixed scale; it does not give one range uniform in \(j\).  Polymer
adjacency and its degree must therefore be defined in rescaled level-\(j\)
blocks.  Counting all levels in the original fine graph would create a false
scale-dependent degree.
\end{firewall}

\subsection{A correlated simultaneous-excursion estimate}
\label{sec:v4v22-excursion}

The residual field at level \(j\),
\begin{equation}
 \eta_j=E_jq_j,
 \qquad q_j\sim N(0,A_j^{-1}),
\label{eq:v4v22-residual-field}
\end{equation}
need not have finite-range covariance.  Its covariance is the positive
residual
\[
 K_j=E_jA_j^{-1}E_j^*
 \le A_0^{-1}
\]
as quadratic forms, by
\eqref{eq:v4v22-multiscale-covariance}.  Assume the hard gap gives
\begin{equation}
 A_0^{-1}\le\sigma^2\,1
\label{eq:v4v22-hard-upper-covariance}
\end{equation}
uniformly in volume and the buffered toron chart.

\begin{lemma}[Many correlated Gaussian excursions]
\label{lem:v4v22-many-excursions}
Let \(\eta\) be centered Gaussian with covariance \(K\le\sigma^2 1\).
Let \(P_1,\ldots,P_m\) be mutually orthogonal coordinate projections of rank
at most \(r\).  If
\[
 R^2\ge4r\sigma^2\log2,
\]
then
\begin{equation}
 \mathbb P\left(
 \|P_i\eta\|\ge R\text{ for every }i
 \right)
 \le
 \exp\left[-\frac{mR^2}{8\sigma^2}\right].
\label{eq:v4v22-many-excursions}
\end{equation}
\end{lemma}

\begin{proof}
Put \(P=\sum_iP_i\).  On the displayed event,
\(\|P\eta\|^2\ge mR^2\).  For \(t=(4\sigma^2)^{-1}\), the Gaussian
quadratic-moment formula and \(PKP\le\sigma^2P\) give
\[
 \mathbb E e^{t\|P\eta\|^2}
 =
 \det(1-2tPKP)^{-1/2}
 \le
 2^{rm/2}.
\]
Markov's inequality therefore yields
\[
 \mathbb P(\|P\eta\|^2\ge mR^2)
 \le
 \exp\left[
 -\frac{mR^2}{4\sigma^2}+\frac{rm}{2}\log2
 \right].
\]
The assumed lower bound on \(R^2\) makes the second term at most
\(mR^2/(8\sigma^2)\), proving
\eqref{eq:v4v22-many-excursions}.
\end{proof}

Let \(Q_x\) project onto a fixed-radius coordinate neighborhood of a
level-\(j\) block \(x\).  Its rank is uniformly bounded in level coordinates.
Let \(D_j\) bound the overlap degree of these neighborhoods.

\begin{corollary}[Coarse bad-set payment without independence]
\label{cor:v4v22-coarse-bad-payment}
For every finite set \(X\) of level-\(j\) blocks,
\begin{equation}
 \mathbb P\left(
 \|Q_x\eta_j\|\ge R\text{ for every }x\in X
 \right)
 \le
 \exp\left[
 -\frac{R^2}{8\sigma^2(D_j+1)}|X|
 \right]
\label{eq:v4v22-coarse-bad-payment}
\end{equation}
once \(R\) exceeds a level-uniform fixed constant.
\end{corollary}

\begin{proof}
The greedy argument of Lemma
\ref{lem:v4v20-independent-tail-set} selects
\(J\subset X\), \(|J|\ge|X|/(D_j+1)\), whose coordinate neighborhoods are
disjoint.  Their projections are mutually orthogonal.  Apply Lemma
\ref{lem:v4v22-many-excursions} to \(J\).  The event for every member of
\(X\) is contained in the event for every member of \(J\).
\end{proof}

\begin{firewall}[A joint tail is not yet a connected activity]
\label{fw:v4v22-tail-not-connectivity}
Corollary~\ref{cor:v4v22-coarse-bad-payment} pays every actual coarse exit,
but it does not make distant residual events independent.  A connected
coarse polymer estimate must retain either the next finite-range Schur shell
or an exponential residual covariance edge.  Summing
\eqref{eq:v4v22-coarse-bad-payment} over arbitrary unrooted subsets would
introduce a volume divergence.
\end{firewall}

\subsection{Exact ownership by the first failed scale}
\label{sec:v4v22-first-failure}

For a fixed spatial root or owned block, let \(\mathsf G_j\) be the event
that the level-\(j\) residual harmonic field lies in its assigned inner
chart.  No nesting assumption is required.

\begin{lemma}[First-failure partition]
\label{lem:v4v22-first-failure}
For every \(J\ge0\),
\begin{equation}
 1
 =
 \sum_{j=0}^J
 \left(\prod_{k=0}^{j-1}\boldsymbol1_{\mathsf G_k}\right)
 \boldsymbol1_{\mathsf G_j^c}
 +
 \prod_{k=0}^J\boldsymbol1_{\mathsf G_k}.
\label{eq:v4v22-first-failure}
\end{equation}
The empty product at \(j=0\) is one.  Exactly one summand is nonzero.
\end{lemma}

\begin{proof}
If a failure occurs, let \(j\) be its first index.  The \(j\)-th summand is
one and every other summand is zero.  If no failure occurs, only the final
product is one.  This proves both the identity and uniqueness.
\end{proof}

Apply \eqref{eq:v4v22-first-failure} independently to every owned spatial
packet and then group packets by the graph-connected components of their
first-failure labels.  A packet which first fails at level \(j\) is excluded
from every finer good-sector logarithm and is handed once to the level-\(j\)
coarse-defect activity.  Packets in the final all-good term may be integrated
successively using V21 on the independent shells in
\eqref{eq:v4v22-multiscale-field}.

\begin{proposition}[No duplicate cutoff or determinant ancestry]
\label{prop:v4v22-no-duplicate-ancestry}
Under the first-failure partition, every field configuration has:
\begin{enumerate}
\item exactly one Gaussian-shell ancestry from
\eqref{eq:v4v22-multiscale-field};
\item exactly one determinant ancestry from
\eqref{eq:v4v22-telescoping-determinant}; and
\item for each owned packet, either one first-failure label or the terminal
all-good label.
\end{enumerate}
Thus a cutoff transition and a Schur determinant cannot be charged again at
a later scale.
\end{proposition}

\begin{proof}
The Gaussian and determinant statements are identities in
Theorem~\ref{thm:v4v22-multiscale-split}.  The ownership statement is
Lemma~\ref{lem:v4v22-first-failure}.  Taking their product gives a disjoint
partition of every integral term, so no configuration or normalization
factor occurs twice.
\end{proof}

\begin{firewall}[Compatibility with V7 and V11 is still geometric data]
\label{fw:v4v22-V7-compatibility}
The first-failure identity determines scale ownership after local events
\(\mathsf G_j\) have been specified.  It does not prove that their spatial
supports coincide with the invariant plaquette ownership used by V7 and the
repair forest of V11.  That compatibility requires the harmonic extension
to be estimated in plaquette curvature variables, not only in coordinate
norm.
\end{firewall}

\subsection{V22 result ledger and V23 target}
\label{sec:v4v22-conclusion}

V22 proves an exact multiscale decomposition of the hard Gaussian covariance
into independent positive finite-range Schur shells plus a positive terminal
residual.  The determinant factors telescope in the same ancestry.  Schur
elimination preserves locality with an explicit range recursion.  A new
Gaussian quadratic-moment estimate pays a linearly large family of
simultaneous coarse excursions even when the residual coordinates are
correlated.  The first-failure identity assigns every chart exit to one
scale and prevents duplicate cutoff and determinant payments.

V23 should express the harmonic extension in lattice curvature variables.
The target is a deterministic estimate comparing the plaquette curvature of
\(H_jq\) inside a block with boundary/coarse curvature and a collar term.
If a toron or centralizer component lies in the kernel of that estimate, it
must remain in the V12--V14 soft parent rather than be mislabeled as a
large-field defect.

\clearpage
\section{V23: harmonic extension in curvature variables}
\label{sec:v4v23}

V22 assigns the first failure of a coordinate chart to one Schur scale.
The V7 good/defect partition, however, is defined by plaquette holonomy.
A coordinate norm alone cannot bridge the two: an exact gauge direction may
be large with zero curvature, and a periodic flat connection carries toron
data that must remain in the V12--V14 soft parent.

This note proves the required fixed-block bridge.  After tree gauge and
projection of transferred flat modes, finite-dimensional relative Hodge
theory controls a boundary cochain by its boundary curvature and period
obstructions.  Uniform elliptic comparison transfers that control to the
Schur harmonic extension.  A quantitative Baker--Campbell--Hausdorff
estimate then shows that the nonlinear plaquette curvature stays inside a
buffered V7 chart.  The constants are fixed-block constants; uniformity
through all toron strata remains a separate gate.

\subsection{Relative cochains and the hard boundary space}
\label{sec:v4v23-relative}

Let \(Q\) be a fixed finite cubical block.  Write
\[
 C^0(Q)\xrightarrow{d_0}C^1(Q)\xrightarrow{d_1}C^2(Q)
\]
for real Lie-algebra-valued cellular cochains with fixed Euclidean inner
products.  Let
\[
 T:C^1(Q)\longrightarrow C^1(\partial Q)
\]
be the boundary trace.  Tree gauge removes the trace of
\(\operatorname{ran}d_0\).  Let \(\mathcal F_Q=\ker d_1\) and let
\(\mathcal T_Q\subset\mathcal F_Q\) be the finite-dimensional flat sector
which has already been transferred to the toron parent.

Choose a boundary hard space \(\mathcal B_{\rm h}\) complementary to
\[
 T(\operatorname{ran}d_0+\mathcal T_Q)
\]
inside the tree-gauge boundary coordinates.  Let
\begin{equation}
 \mathcal O_\partial:
 \mathcal B_{\rm h}\longrightarrow\mathcal Y_\partial
\label{eq:v4v23-boundary-obstruction-map}
\end{equation}
collect the boundary plaquette coboundary and a basis of relative period
functionals.  Assume it is complete in the precise sense
\begin{equation}
 \ker\mathcal O_\partial
 =
 \mathcal B_{\rm h}\cap T(\mathcal F_Q)
 =
 \{0\}.
\label{eq:v4v23-complete-obstruction}
\end{equation}

\begin{lemma}[Fixed-block relative Hodge estimate]
\label{lem:v4v23-relative-Hodge}
There is a linear extension
\(L:\mathcal B_{\rm h}\to C^1(Q)\), \(TLq=q\), and a constant \(C_Q\) such
that
\begin{equation}
 \|Lq\|+\|d_1Lq\|+\|d_0^*Lq\|
 \le C_Q\|\mathcal O_\partial q\|
\label{eq:v4v23-relative-Hodge}
\end{equation}
for every \(q\in\mathcal B_{\rm h}\).
\end{lemma}

\begin{proof}
Choose any linear right inverse \(L_0\) of the trace on
\(\mathcal B_{\rm h}\), which exists because all spaces are finite
dimensional.  On \(\mathcal B_{\rm h}\), condition
\eqref{eq:v4v23-complete-obstruction} makes
\(q\mapsto\|\mathcal O_\partial q\|\) a norm.  Every two norms on a fixed
finite-dimensional space are equivalent, so
\[
 \|q\|\le C\|\mathcal O_\partial q\|.
\]
The linear maps \(L_0,d_1L_0,d_0^*L_0\) are bounded.  Taking
\(L=L_0\) and combining their operator norms proves
\eqref{eq:v4v23-relative-Hodge}.
\end{proof}

\begin{firewall}[The obstruction map is necessary]
\label{fw:v4v23-obstruction-necessary}
Boundary curvature alone need not be injective on boundary cochains.
Relative periods and transferred flat holonomies can lie in its kernel.
Equation \eqref{eq:v4v23-complete-obstruction} is therefore a theorem
hypothesis, not a consequence of small plaquettes.  Omitting it would
misclassify a toron as a large-field defect.
\end{firewall}

\subsection{The Schur harmonic extension}
\label{sec:v4v23-Schur-harmonic}

Let \(A_Q>0\) be the level precision on the block interior with boundary
trace fixed.  Its Schur harmonic extension \(H_Qq\) is the unique minimizer
of
\begin{equation}
 \frac12\langle a,A_Qa\rangle
 \quad\text{subject to }Ta=q.
\label{eq:v4v23-variational-extension}
\end{equation}
This is the variational form of
\eqref{eq:v4v19-harmonic-extension}.  Assume the hard-sector elliptic
comparison
\begin{equation}
 c_{\rm e}
 \left(\|a\|^2+\|d_1a\|^2+\|d_0^*a\|^2\right)
 \le
 \langle a,A_Qa\rangle
 \le
 C_{\rm e}
 \left(\|a\|^2+\|d_1a\|^2+\|d_0^*a\|^2\right)
\label{eq:v4v23-elliptic-comparison}
\end{equation}
for hard tree-gauge cochains, with constants uniform on the buffered chart.

\begin{theorem}[Curvature control of the harmonic extension]
\label{thm:v4v23-harmonic-curvature}
Under \eqref{eq:v4v23-complete-obstruction} and
\eqref{eq:v4v23-elliptic-comparison},
\begin{equation}
 \|H_Qq\|+\|d_1H_Qq\|+\|d_0^*H_Qq\|
 \le
 C_{\rm H}\|\mathcal O_\partial q\|,
\label{eq:v4v23-harmonic-curvature}
\end{equation}
where
\[
 C_{\rm H}
 =
 C_Q\sqrt{\frac{3C_{\rm e}}{c_{\rm e}}}
\]
is independent of \(q\).
\end{theorem}

\begin{proof}
The competitor \(Lq\) from Lemma
\ref{lem:v4v23-relative-Hodge} has the same trace as \(H_Qq\).  Minimality
and the two sides of \eqref{eq:v4v23-elliptic-comparison} give
\begin{align*}
 c_{\rm e}\left(
 \|H_Qq\|^2+\|d_1H_Qq\|^2+\|d_0^*H_Qq\|^2
 \right)
 &\le
 \langle H_Qq,A_QH_Qq\rangle\\
 &\le
 \langle Lq,A_QLq\rangle\\
 &\le
 C_{\rm e}\left(
 \|Lq\|^2+\|d_1Lq\|^2+\|d_0^*Lq\|^2
 \right)\\
 &\le3C_{\rm e}C_Q^2\|\mathcal O_\partial q\|^2.
\end{align*}
Take square roots.
\end{proof}

\begin{corollary}[Buffered linear chart transfer]
\label{cor:v4v23-linear-buffer}
If
\begin{equation}
 \|\mathcal O_\partial q\|
 \le\frac{R}{4C_{\rm H}},
\label{eq:v4v23-inner-boundary}
\end{equation}
then
\begin{equation}
 \|H_Qq\|\le R/4,\qquad
 \|d_1H_Qq\|\le R/4.
\label{eq:v4v23-inner-extension}
\end{equation}
Thus the V21 hypothesis \(\|Hq\|\le R/2\) holds with a factor-two buffer.
\end{corollary}

\begin{proof}
Insert \eqref{eq:v4v23-inner-boundary} in
\eqref{eq:v4v23-harmonic-curvature}.
\end{proof}

\subsection{From linear curvature to plaquette holonomy}
\label{sec:v4v23-BCH}

Let \(a_e\in\mathfrak g\) be a link cochain on \(Q\) and
\(U_e=\exp(ga_e)\).  For an oriented plaquette \(p\), write
\[
 U_p=U_{e_1}U_{e_2}U_{e_3}^{-1}U_{e_4}^{-1}.
\]

\begin{lemma}[Quantitative plaquette BCH remainder]
\label{lem:v4v23-plaquette-BCH}
There are constants \(r_{\rm BCH},C_{\rm BCH}>0\), depending only on the
compact matrix group and the chosen norm, such that if
\(g\max_{e\subset p}\|a_e\|\le r_{\rm BCH}\), then the principal logarithm
is defined and
\begin{equation}
 \left\|
 \log U_p-g(d_1a)_p
 \right\|
 \le
 C_{\rm BCH}g^2
 \left(\max_{e\subset p}\|a_e\|\right)^2.
\label{eq:v4v23-plaquette-BCH}
\end{equation}
Consequently
\begin{equation}
 d_G(U_p,1)
 \le
 Cg\left(
 \|(d_1a)_p\|
 +g\max_{e\subset p}\|a_e\|^2
 \right).
\label{eq:v4v23-group-curvature}
\end{equation}
\end{lemma}

\begin{proof}
The logarithm of a product of four exponentials is analytic in a fixed
neighborhood of the origin.  Its derivative at the origin is the signed sum
of the four edge variables, namely \((d_1a)_p\).  Taylor's theorem on the
compact ball of radius \(r_{\rm BCH}\) bounds the quadratic remainder by
the right side of \eqref{eq:v4v23-plaquette-BCH}.  Local equivalence of the
group distance and the Lie-algebra norm of the principal logarithm gives
\eqref{eq:v4v23-group-curvature}.
\end{proof}

\begin{theorem}[Nonlinear preservation of the V7 plaquette chart]
\label{thm:v4v23-nonlinear-chart}
Let
\[
 R_g=g^{-\alpha},\qquad
 r_g=g^{1-\alpha},\qquad0<\alpha<1.
\]
Suppose \(q\in\mathcal B_{\rm h}\) obeys
\begin{equation}
 \|\mathcal O_\partial q\|
 \le c_{\rm in}R_g
\label{eq:v4v23-coarse-good}
\end{equation}
with \(c_{\rm in}>0\) small enough in terms of \(C_{\rm H}\).  Then for all
sufficiently small \(g\), every plaquette of the harmonic extension satisfies
\begin{equation}
 d_G\left((\exp(gH_Qq))_p,1\right)
 \le\frac12 r_g.
\label{eq:v4v23-plaquette-buffer}
\end{equation}
The remaining half of the V7 threshold is available for the independent
finite-range fluctuation shell.
\end{theorem}

\begin{proof}
Choose \(c_{\rm in}\) so that
\eqref{eq:v4v23-harmonic-curvature} gives
\[
 \|H_Qq\|+\|d_1H_Qq\|\le c_0R_g
\]
with \(c_0\) as small as required below.  Since
\(gR_g=g^{1-\alpha}\to0\), the BCH hypothesis holds for small \(g\).
Equation \eqref{eq:v4v23-group-curvature} gives
\[
 d_G(U_p,1)
 \le
 Cc_0gR_g+C c_0^2g^2R_g^2
 =
 Cc_0r_g+C c_0^2r_g^2.
\]
First choose \(c_0\) so \(Cc_0\le1/4\), then reduce \(g\) so
\(Cc_0^2r_g\le1/4\).  The result is
\eqref{eq:v4v23-plaquette-buffer}.
\end{proof}

\begin{firewall}[Noncommuting flat coordinates are not free]
\label{fw:v4v23-noncommuting-flat}
The quadratic BCH term shows why an arbitrary large coordinate in
\(\ker d_1\) cannot be ignored: noncommuting link values can create
plaquette curvature at order \(g^2\|a\|^2\).  The transferred toron sector
must consist of genuine flat compact-group holonomies, including their
commutation constraints, as in V12--V13.  The linear projection
\(\mathcal T_Q\) is only the tangent description on a buffered chart.
\end{firewall}

\subsection{Compatibility with first-failure ownership}
\label{sec:v4v23-ownership}

Define the level-\(j\) coarse-good event on an owned block by
\begin{equation}
 \mathsf G_{j,x}
 =
 \left\{
 \|\mathcal O_{\partial,j}q_j\|
 \le c_{\rm in}R_{g_j}
 \right\}.
\label{eq:v4v23-level-good}
\end{equation}
On \(\mathsf G_{j,x}\), Theorem
\ref{thm:v4v23-nonlinear-chart} puts the harmonic extension inside half of
the invariant plaquette threshold.  Its independent Schur fluctuation may
then use the other half in the V20 cutoff surgery.  On
\(\mathsf G_{j,x}^c\), the V22 first-failure identity assigns the block once
to level \(j\).

\begin{proposition}[Gauge directions are excluded from the failure test]
\label{prop:v4v23-gauge-exclusion}
The event \(\mathsf G_{j,x}\) depends only on the hard quotient boundary
class.  Changing \(q_j\) by the trace of an exact cochain or by a transferred
flat toron coordinate does not create a coarse defect label.
\end{proposition}

\begin{proof}
Tree gauge removes the exact trace before
\(\mathcal B_{\rm h}\) is chosen.  The toron trace is projected into
\(\mathcal T_Q\), which is complementary to
\(\mathcal B_{\rm h}\).  The obstruction map acts only on the remaining
hard class, so \eqref{eq:v4v23-level-good} is unchanged.
\end{proof}

\begin{firewall}[Uniform relative Hodge constants remain open]
\label{fw:v4v23-uniformity}
Lemma~\ref{lem:v4v23-relative-Hodge} uses norm equivalence on one fixed
block.  Applying it uniformly along the programme requires:
\begin{enumerate}
\item a rescaled family of block complexes with uniformly bounded
right-inverse constants;
\item a toron-stratum complement whose projection norms stay bounded on the
buffered cover; and
\item background-Ward covariance of the trace, period, and tree-gauge maps.
\end{enumerate}
V23 does not prove these three uniform statements.
\end{firewall}

\subsection{V23 result ledger and V24 target}
\label{sec:v4v23-conclusion}

V23 proves the deterministic curvature bridge on a fixed hard block.
Complete boundary curvature and period data control a relative Hodge
extension after exact gauge and toron directions are removed.  Uniform
elliptic comparison transfers this bound to the Schur harmonic extension.
The nonlinear plaquette BCH remainder is of order
\(g^2R_g^2=r_g^2\), so a sufficiently buffered coarse-good boundary stays
inside half of the V7 invariant plaquette threshold.  First-failure
ownership therefore need not penalize pure-gauge or transferred toron
coordinates.

V24 should attack uniformity of the relative Hodge constants across the
rescaled block family and the finite buffered toron-stratum cover.  The
natural proof is a compactness argument for the normalized obstruction
operator.  Its possible failure mode is an approximate kernel; any such
sequence must be classified as gauge, toron/centralizer soft data, or a
genuine new boundary cohomology obstruction before the programme proceeds.

\clearpage
\section{V24: uniform relative Hodge inversion on buffered toron strata}
\label{sec:v4v24}

V23 reduced harmonic chart transfer to a lower bound for the boundary
curvature--period obstruction operator.  This note proves the needed
uniform compactness theorem and verifies its kernel on ordinary cubical and
periodic blocks.  The only modes invisible to curvature and periods are
exact gauge traces and flat toron traces; both have already been removed
from the hard boundary space.

On a fixed-rank buffered toron chart, the hard projector, obstruction map,
and elliptic precision vary continuously.  A normalized approximate kernel
would converge to a genuine kernel, contradicting the cohomological
classification.  Hence the smallest hard singular value has a positive
uniform lower bound.  The conclusion is uniform over any finite family of
rescaled block types and compact buffered charts.

\subsection{Cohomological completeness of curvature and periods}
\label{sec:v4v24-cohomology}

Let \(Q\) be a finite connected cubical complex and let
\(Z^1(\partial Q)=\ker(d_\partial:C^1(\partial Q)\to C^2(\partial Q))\).
Choose cycles \(\gamma_1,\ldots,\gamma_b\) whose homology classes form a
basis of \(H_1(\partial Q;\mathbb R)\), and define
\begin{equation}
 \operatorname{Per}(q)
 =
 \left(
 \int_{\gamma_1}q,\ldots,\int_{\gamma_b}q
 \right).
\label{eq:v4v24-period-map}
\end{equation}
The coefficients are Lie-algebra valued; the statement applies component by
component.

\begin{lemma}[Boundary de Rham detection]
\label{lem:v4v24-boundary-detection}
If
\begin{equation}
 d_\partial q=0,\qquad \operatorname{Per}(q)=0,
\label{eq:v4v24-zero-curvature-period}
\end{equation}
then \(q=d_\partial\phi\) for a boundary zero-cochain \(\phi\).
\end{lemma}

\begin{proof}
Cellular cohomology identifies
\[
 Z^1(\partial Q)/\operatorname{ran}d_\partial
 \simeq H^1(\partial Q;\mathbb R)\otimes\mathfrak g.
\]
The universal coefficient theorem over \(\mathbb R\) identifies a
cohomology class with its evaluations on a basis of
\(H_1(\partial Q;\mathbb R)\).  The period conditions in
\eqref{eq:v4v24-zero-curvature-period} say that the class of \(q\) evaluates
to zero on that basis, so the class is zero.  Hence \(q\) is exact.
\end{proof}

\begin{corollary}[Ordinary cubical and periodic kernels]
\label{cor:v4v24-block-kernels}
For a four-dimensional cubical ball, \(\partial Q\) has the cellular
cohomology of \(S^3\), so \(H^1(\partial Q)=0\): boundary curvature alone
detects every tree-gauge hard class.  For a periodic four-torus, the four
fundamental periods detect \(H^1(T^4)\); their Lie-algebra values are the
linear toron coordinates.  After exact traces and the transferred toron
space are removed, the V23 obstruction map has zero kernel.
\end{corollary}

\begin{proof}
The first statement follows from \(H^1(S^3;\mathbb R)=0\).  The second
follows from \(H^1(T^4;\mathbb R)\simeq\mathbb R^4\).  Apply Lemma
\ref{lem:v4v24-boundary-detection}; tree gauge removes the exact class and
the V12 soft parent removes the period class.
\end{proof}

\begin{firewall}[Punctured blocks require their actual relative periods]
\label{fw:v4v24-punctured-periods}
If a V11 repair collar or deleted defect packet makes the block boundary
non-spherical, the new first homology generators must be included in
\eqref{eq:v4v24-period-map}.  The cubical-ball corollary cannot be quoted on
a punctured block.  A missing relative period is a genuine boundary
cohomology obstruction, not a toron and not a norm-estimate failure.
\end{firewall}

\subsection{Uniform injectivity on a varying hard bundle}
\label{sec:v4v24-uniform-injectivity}

Let \(K\) be a compact parameter space.  In the application, \(K\) is the
closure of one buffered toron chart together with a compact family of
rescaled local metric coefficients.  Let \(\mathcal B\) and \(\mathcal Y\)
be fixed finite-dimensional Hilbert spaces.  Suppose
\[
 P(\theta)=P(\theta)^*=P(\theta)^2,\qquad \theta\in K,
\]
is a norm-continuous family of fixed-rank projections.  Its range
\(\mathcal B_{\rm h}(\theta)\) is the hard boundary fibre.  Let
\[
 O(\theta):\mathcal B\to\mathcal Y
\]
be norm-continuous and assume
\begin{equation}
 \ker O(\theta)\cap\operatorname{ran}P(\theta)=\{0\}
 \qquad(\theta\in K).
\label{eq:v4v24-pointwise-injectivity}
\end{equation}

\begin{theorem}[Uniform hard obstruction gap]
\label{thm:v4v24-uniform-obstruction}
There is \(\sigma_*>0\) such that
\begin{equation}
 \|O(\theta)q\|
 \ge\sigma_*\|q\|
 \qquad
 (\theta\in K,\ q\in\operatorname{ran}P(\theta)).
\label{eq:v4v24-uniform-obstruction}
\end{equation}
Equivalently, the inverse of \(O(\theta)\) on its hard range has norm at
most \(\sigma_*^{-1}\), uniformly in \(\theta\).
\end{theorem}

\begin{proof}
If no \(\sigma_*\) existed, there would be
\(\theta_n\in K\) and \(q_n\in\operatorname{ran}P(\theta_n)\) with
\[
 \|q_n\|=1,\qquad \|O(\theta_n)q_n\|\longrightarrow0.
\]
Pass to subsequences so \(\theta_n\to\theta_\infty\) and
\(q_n\to q_\infty\); finite dimensionality gives the second convergence.
Norm continuity of \(P\) implies
\[
 P(\theta_\infty)q_\infty
 =
 \lim_nP(\theta_n)q_n
 =
 q_\infty.
\]
Also \(\|q_\infty\|=1\), while continuity of \(O\) gives
\(O(\theta_\infty)q_\infty=0\).  This contradicts
\eqref{eq:v4v24-pointwise-injectivity}.
\end{proof}

\begin{corollary}[Continuous pseudoinverse and covariant extension]
\label{cor:v4v24-pseudoinverse}
Let
\[
 B(\theta)=O(\theta)P(\theta).
\]
On the hard range, the Moore--Penrose inverse
\(B(\theta)^\dagger\) is norm-continuous and
\begin{equation}
 \|B(\theta)^\dagger\|\le\sigma_*^{-1}.
\label{eq:v4v24-pseudoinverse-bound}
\end{equation}
If a compact residual gauge group acts unitarily by
\(U_{\mathcal B}(h)\), \(U_{\mathcal Y}(h)\) and
\begin{align}
 P(h\theta)&=
 U_{\mathcal B}(h)P(\theta)U_{\mathcal B}(h)^{-1},
\label{eq:v4v24-P-equivariant}\\
 O(h\theta)U_{\mathcal B}(h)&=
 U_{\mathcal Y}(h)O(\theta),
\label{eq:v4v24-O-equivariant}
\end{align}
then
\begin{equation}
 B(h\theta)^\dagger
 =
 U_{\mathcal B}(h)B(\theta)^\dagger U_{\mathcal Y}(h)^{-1}.
\label{eq:v4v24-pseudoinverse-equivariant}
\end{equation}
\end{corollary}

\begin{proof}
The nonzero singular values of \(B(\theta)\) are bounded below by
\(\sigma_*\), so its rank is constant and the standard spectral formula for
the pseudoinverse is continuous.  Its norm is the reciprocal of the
smallest nonzero singular value, proving
\eqref{eq:v4v24-pseudoinverse-bound}.  Equations
\eqref{eq:v4v24-P-equivariant}--%
\eqref{eq:v4v24-O-equivariant} conjugate the singular-value decomposition;
uniqueness of the Moore--Penrose inverse gives
\eqref{eq:v4v24-pseudoinverse-equivariant}.
\end{proof}

\subsection{Why the buffered soft transfer is essential}
\label{sec:v4v24-buffer}

Let \(L_\theta\) be the gauge-fixed boundary extension and
\(A_Q(\theta)\) the hard precision.  On one V14 buffered chart, all modes
whose eigenvalue can fall below \(2\epsilon\) are included in the soft
bundle.  The hard complement therefore has a fixed rank and a uniform gap
\(\epsilon\).

\begin{proposition}[Approximate kernels are exactly the transfer test]
\label{prop:v4v24-approximate-kernel}
Suppose \(\theta_n\to\theta_\infty\) and normalized boundary vectors \(q_n\)
satisfy
\[
 \|O(\theta_n)q_n\|\to0.
\]
Then one of the following occurs:
\begin{enumerate}
\item the hard projectors converge and \(q_n\) has a nonzero limit in the
true kernel, which is an exact gauge, toron, or relative cohomology class;
or
\item the projectors do not remain norm-continuous because a spectral mode
crosses the hard threshold.
\end{enumerate}
The V14 buffered transfer excludes the second case on each chart by moving
every crossing mode into the soft bundle.  Corollary
\ref{cor:v4v24-block-kernels} and the complete period list exclude the first
case on the hard quotient of an ordinary block.
\end{proposition}

\begin{proof}
If the projectors converge in norm, repeat the compactness proof of
Theorem~\ref{thm:v4v24-uniform-obstruction}; any convergent subsequence of
\(q_n\) produces a nonzero vector in
\(\ker O(\theta_\infty)\) on the limiting fibre.  Lemma
\ref{lem:v4v24-boundary-detection} classifies that kernel.  If the
projectors do not converge in norm, their fixed-rank spectral subspaces are
not separated from the complementary spectrum; in finite dimension this
requires a spectral crossing at the selected boundary.  The buffered
spectral projector includes the entire crossing cluster on one side of a
fixed gap, making it norm-continuous there.
\end{proof}

\begin{firewall}[A rank jump cannot be repaired by compactness]
\label{fw:v4v24-rank-jump}
Theorem~\ref{thm:v4v24-uniform-obstruction} requires a fixed-rank continuous
hard bundle.  Applying it through a centralizer wall without first
transferring the rank-changing modes would be circular.  The soft transfer
of V13--V14 is logically prior to the compactness argument.
\end{firewall}

\subsection{Uniform harmonic and nonlinear chart constants}
\label{sec:v4v24-uniform-chart}

Assume there are finitely many rescaled block types.  For each type and
buffered chart, suppose:
\begin{enumerate}
\item the cellular trace and coboundary matrices vary in a compact,
shape-regular family;
\item the complete period map of Section
\ref{sec:v4v24-cohomology} is used;
\item the V14 hard projector is norm-continuous and has fixed rank; and
\item the elliptic constants in
\eqref{eq:v4v23-elliptic-comparison} obey
\(0<c_{\rm e}\le C_{\rm e}<\infty\) uniformly.
\end{enumerate}

\begin{theorem}[Uniform buffered harmonic chart transfer]
\label{thm:v4v24-uniform-chart}
Under these hypotheses, there are constants
\(C_{\rm H}^{\rm unif}<\infty\), \(c_{\rm in}>0\), and \(g_0>0\),
independent of block, lattice volume, and toron parameter within the
buffered cover, such that
\begin{equation}
 \|\mathcal O_{\partial,\theta}q\|
 \le c_{\rm in}R_g
 \quad\Longrightarrow\quad
 d_G\left((\exp(gH_{Q,\theta}q))_p,1\right)
 \le\frac12g^{1-\alpha}
\label{eq:v4v24-uniform-chart}
\end{equation}
for every plaquette in the owned block, \(0<g<g_0\), and
\(R_g=g^{-\alpha}\), \(0<\alpha<1\).
The extension may be chosen equivariantly under the residual background
gauge action.
\end{theorem}

\begin{proof}
On each compact chart, Theorem
\ref{thm:v4v24-uniform-obstruction} gives a uniform lower singular value.
The ambient trace right inverses and cellular matrices have uniformly
bounded norms by compactness and shape regularity.  The proof of Lemma
\ref{lem:v4v23-relative-Hodge} therefore gives one uniform \(C_Q\).
Uniform ellipticity then gives a uniform constant in Theorem
\ref{thm:v4v23-harmonic-curvature}.  Take the maximum over the finite
block-type and chart family.

With this constant fixed, the BCH proof of Theorem
\ref{thm:v4v23-nonlinear-chart} chooses one \(c_{\rm in}\) and then one
\(g_0\).  To make the extension equivariant, average an ambient trace right
inverse \(L_0\) over the compact residual gauge group and set
\[
 L_\theta=L_0B(\theta)^\dagger O(\theta)
\]
on the hard fibre.  There
\(B(\theta)^\dagger O(\theta)q=q\), so \(TL_\theta q=q\);
equivariance follows from
\eqref{eq:v4v24-pseudoinverse-equivariant}.
\end{proof}

\begin{corollary}[Closure of the V23 analytic uniformity gate]
\label{cor:v4v24-V23-gate}
Once a volume-uniform V14 buffered cover and a shape-regular finite block
family are supplied, no additional analytic loss occurs in passing from
the boundary curvature--period data to the invariant plaquette cutoff.
The constants in V20--V23 can be chosen uniformly on every certified hard
chart.
\end{corollary}

\begin{proof}
This is Theorem~\ref{thm:v4v24-uniform-chart}, followed by the finite-range
cutoff-tail estimate of V20 and first-failure ownership of V22.
\end{proof}

\begin{firewall}[The cover hypothesis remains unproved along the cutoff]
\label{fw:v4v24-cover-open}
V14 constructs a finite buffered cover on each finite lattice but does not
prove that its number of charts, transition norms, and spectral buffers are
uniform along the continuum cutoff family.  V24 proves that such a uniform
cover is sufficient and that no new fixed-block Hodge degeneration occurs.
It does not prove the volume-uniform cover itself, nor the corresponding
punctured-block period catalogue for all V11 terminal geometries.
\end{firewall}

\subsection{V24 result ledger and compulsory V25 target}
\label{sec:v4v24-conclusion}

V24 proves that curvature and a complete period basis detect every hard
boundary class on ordinary cubical and periodic blocks.  On a fixed-rank
buffered toron chart, compactness turns pointwise injectivity into a uniform
singular-value bound.  Approximate kernels are thereby classified: they are
exact gauge/toron/cohomology modes or threshold-crossing modes, and the
latter must be transferred before compactness is used.  Under the explicit
volume-uniform cover and shape-regularity hypotheses, the relative-Hodge,
Schur harmonic, BCH, and cutoff constants are uniform and background
equivariant.

V25 is the next compulsory companion audit.  It should inspect the newest
Yang--Mills and Navier--Stokes notes, then attack the remaining V14 cover
gate.  The first test is whether the toron-shifted lattice spectrum admits a
single volume-uniform semialgebraic resonance atlas after eigenvalues are
grouped into fixed spectral windows, rather than one chart per individual
mode.

\clearpage
\section{V25: companion audit and a volume-uniform recentered toron atlas}
\label{sec:v4v25}

V24 reduced uniform harmonic control to a volume-uniform buffered toron
decomposition.  A cover by one neighborhood for every crossing Fourier mode
would have cardinality growing with the lattice and would not solve that
gate.  This compulsory audit points to the correct replacement: retain the
complete disjoint partition and prove a pointwise estimate which is
invariant under unitary momentum relabeling.  The number of cells then never
appears as a quantitative factor.

For the explicit V12 toron-shifted covariant Laplacian, each root momentum
can be recentered into one half-open Brillouin cell.  After recentering, all
toron dependence lies in a residual shift of size at most \(\pi/M\).
A single enlarged spectral template contains every mode which can fall
below the soft threshold, while its complement has a uniform hard gap.
The template rank obeys the V14 Weyl bound.  This closes the atlas-count
problem for the flat quadratic Laplacian, though not yet for the complete
interacting Hessian.

\subsection{Compulsory companion audit}
\label{sec:v4v25-audit}

The exact files inspected were
\begin{enumerate}
\item
\texttt{Renewal\_Yang\_Mills\_Mass\_Gap\_Research\_Notes\_V45.tex},
\(835475\) bytes, SHA--256
\[
 \texttt{8FE21478A9D51095B034EEAE6512CC434E707B60A91CA586B15B3570EB9C97C0}.
\]
It was byte-identical to V44 at the audit time, and its last internal
mathematical chapter was V44.
\item
\texttt{Renewal\_NS\_Stability\_Research\_Notes\_V31.tex},
\(887537\) bytes, SHA--256
\[
 \texttt{F1121B62238AD5491BB7CF03635EA9934942EDF573ED8FAAA9EE79E1D38A6696}.
\]
\end{enumerate}

YM V44 proves an exact doubled-prefix square but also a sharp square-root
ceiling: the second copy forced by a positive square does not create a
second response debit.  NS V31 proves that a flat dyadic probability mark
is exact, but removing it restores the full first-moment weight.  The common
lesson for the toron atlas is that multiplicity is neither a new smallness
factor nor a new cost when the cells form one disjoint exact partition.
One must estimate the assembled conditional object uniformly and integrate
over the partition once.

\subsection{Why one global fixed soft subspace is impossible}
\label{sec:v4v25-global-no-go}

On the periodic lattice of side \(M\) and spacing \(a\), the V12 root-sector
eigenvalues are
\begin{equation}
 \lambda_{k,\gamma}(\Theta)
 =
 \frac4{a^2}
 \sum_{\mu=1}^d
 \sin^2\left(
 \frac{\pi k_\mu}{M}+\frac{\gamma(\Theta_\mu)}2
 \right),
 \qquad
 k\in(\mathbb Z/M\mathbb Z)^d.
\label{eq:v4v25-shifted-spectrum}
\end{equation}

\begin{proposition}[No nontrivial global Fourier soft bundle]
\label{prop:v4v25-no-global-bundle}
Fix a root \(\gamma\).  For every momentum \(k\), there is a toron tuple
\(\Theta\) for which
\(\lambda_{k,\gamma}(\Theta)=0\).  Consequently a Fourier subspace,
independent of \(\Theta\), which contains every root mode that can become
soft must contain the entire \(\gamma\)-root Fourier sector.
\end{proposition}

\begin{proof}
Choose each root angle so
\[
 \gamma(\Theta_\mu)
 \equiv-\frac{2\pi k_\mu}{M}\pmod{2\pi}.
\]
Every sine in \eqref{eq:v4v25-shifted-spectrum} then vanishes.  Hence every
Fourier basis vector belongs to the union, over torons, of the zero-mode
spaces.  Any fixed subspace containing that union is the full root sector.
\end{proof}

\begin{firewall}[The global-bundle no-go is not a soft-rank divergence]
\label{fw:v4v25-global-no-go-scope}
At one fixed toron only a Weyl-sized set of modes is soft.  Proposition
\ref{prop:v4v25-no-global-bundle} concerns the union of those sets over the
whole toron moduli.  It does not contradict the V14 soft-rank density bound;
it proves that toron-dependent momentum recentering is necessary.
\end{firewall}

\subsection{Half-open momentum recentering}
\label{sec:v4v25-recentering}

For every root \(\gamma\), direction \(\mu\), and toron tuple, choose the
unique half-open decomposition
\begin{equation}
 \gamma(\Theta_\mu)
 =
 \frac{2\pi}{M}n_{\gamma,\mu}(\Theta)
 +\delta_{\gamma,\mu}(\Theta)
 \pmod{2\pi},
 \qquad
 -\frac{\pi}{M}\le\delta_{\gamma,\mu}<\frac{\pi}{M}.
\label{eq:v4v25-half-open-recenter}
\end{equation}
Here \(n_{\gamma,\mu}\in\mathbb Z/M\mathbb Z\).  Put
\[
 \ell_\mu=k_\mu+n_{\gamma,\mu}(\Theta)
 \pmod M.
\]
Then
\begin{equation}
 \lambda_{k,\gamma}(\Theta)
 =
 \lambda_\ell(\delta_\gamma)
 :=
 \frac4{a^2}
 \sum_{\mu=1}^d
 \sin^2\left(
 \frac{\pi\ell_\mu}{M}+\frac{\delta_{\gamma,\mu}}2
 \right).
\label{eq:v4v25-recentered-spectrum}
\end{equation}
The map \(k\mapsto\ell\) is a unitary permutation of the Fourier basis.
Choose the half-open tie at a cell wall jointly for the pair
\(\{\gamma,-\gamma\}\).  The two permutations then exchange conjugate
Fourier coefficients together, so their direct sum preserves the real
Lie-algebra field inside the complex root decomposition.

Let
\begin{equation}
 \lambda_\ell^0
 =
 \frac4{a^2}\sum_{\mu=1}^d
 \sin^2\left(\frac{\pi\ell_\mu}{M}\right),
 \qquad
 L=aM,
 \qquad
 b_L=\frac{\pi\sqrt d}{L}.
\label{eq:v4v25-central-symbol}
\end{equation}

\begin{lemma}[Uniform square-root displacement]
\label{lem:v4v25-symbol-displacement}
For every allowed residual shift,
\begin{equation}
 \left|
 \sqrt{\lambda_\ell(\delta)}
 -\sqrt{\lambda_\ell^0}
 \right|
 \le b_L.
\label{eq:v4v25-symbol-displacement}
\end{equation}
\end{lemma}

\begin{proof}
Define the vector symbol
\[
 p_\mu(t)=\frac2a\sin(t_\mu/2).
\]
Then \(\sqrt{\lambda_\ell(\delta)}\) is the Euclidean norm of
\(p(2\pi\ell/M+\delta)\), while
\(\sqrt{\lambda_\ell^0}=\|p(2\pi\ell/M)\|\).
Since \(|(2/a)(\sin((t+s)/2)-\sin(t/2))|\le|s|/a\),
\[
 \left\|
 p(2\pi\ell/M+\delta)-p(2\pi\ell/M)
 \right\|
 \le
 \frac{\sqrt d}{a}\max_\mu|\delta_\mu|
 \le\frac{\pi\sqrt d}{aM}=b_L.
\]
The reverse triangle inequality proves
\eqref{eq:v4v25-symbol-displacement}.
\end{proof}

\subsection{One spectral template on every cell}
\label{sec:v4v25-template}

Fix \(\epsilon>0\).  Define the recentered soft label set
\begin{equation}
 \mathcal S_{\epsilon,L}
 =
 \left\{
 \ell:
 \sqrt{\lambda_\ell^0}
 \le\sqrt{2\epsilon}+b_L
 \right\}.
\label{eq:v4v25-soft-template}
\end{equation}
It is independent of \(\gamma\), \(\Theta\), and the recentering cell.

\begin{theorem}[Uniform recentered soft/hard split]
\label{thm:v4v25-recentered-split}
On each half-open recentering cell, transfer the root Fourier modes whose
recentered labels lie in \(\mathcal S_{\epsilon,L}\).  Then:
\begin{enumerate}
\item every mode satisfying
\(\lambda_{k,\gamma}(\Theta)\le2\epsilon\) is transferred;
\item every complementary mode satisfies
\begin{equation}
 \lambda_{k,\gamma}(\Theta)>2\epsilon;
\label{eq:v4v25-hard-gap}
\end{equation}
\item the soft rank is constant on the cell and equal to
\(|\mathcal S_{\epsilon,L}|\) for each root polarization; and
\item changes of recentering cell act by unitary momentum permutations, so
all operator, determinant, and relative-Hodge bounds are unchanged.
\end{enumerate}
\end{theorem}

\begin{proof}
If a recentered label is outside
\(\mathcal S_{\epsilon,L}\), then Lemma
\ref{lem:v4v25-symbol-displacement} gives
\[
 \sqrt{\lambda_\ell(\delta)}
 \ge\sqrt{\lambda_\ell^0}-b_L
 >\sqrt{2\epsilon},
\]
which proves \eqref{eq:v4v25-hard-gap}.  Its contrapositive proves the first
claim.  The template is fixed on each cell, proving constant rank.
Equation \eqref{eq:v4v25-recentered-spectrum} is obtained by a permutation
of Fourier labels; unitary conjugation preserves the stated bounds.
\end{proof}

\begin{corollary}[Uniform template rank density]
\label{cor:v4v25-template-rank}
There is a dimension-dependent constant \(C_d\) such that
\begin{equation}
 |\mathcal S_{\epsilon,L}|
 \le
 C_d\left(1+L^d\epsilon^{d/2}\right).
\label{eq:v4v25-template-rank}
\end{equation}
In four dimensions,
\begin{equation}
 L^{-4}|\mathcal S_{\epsilon,L}|
 \le C\left(L^{-4}+\epsilon^2\right).
\label{eq:v4v25-template-density-four}
\end{equation}
\end{corollary}

\begin{proof}
The elementary inequality
\[
 |\sin(\pi\ell_\mu/M)|
 \ge c\,\frac{|\ell_\mu|_M}{M}
\]
holds for the representative
\(|\ell_\mu|_M\le M/2\).  Thus
\(\ell\in\mathcal S_{\epsilon,L}\) implies
\[
 |\ell|_M
 \le C\left(L\sqrt\epsilon+1\right).
\]
Counting lattice points in this ball gives
\[
 |\mathcal S_{\epsilon,L}|
 \le C_d(1+L\sqrt\epsilon)^d
 \le C_d'(1+L^d\epsilon^{d/2}),
\]
which proves both claims.
\end{proof}

\subsection{The disjoint atlas has no cardinality factor}
\label{sec:v4v25-no-cell-factor}

For all roots and directions together, the nearest-grid choices in
\eqref{eq:v4v25-half-open-recenter} define a finite measurable disjoint
partition
\[
 \mathcal M_{\rm flat}
 =
 \bigsqcup_{\nu\in\mathcal I_M}\mathcal R_{\nu,M}
\]
up to the finite Weyl quotient.  Its cardinality may grow with \(M\).

\begin{lemma}[Pointwise integration over the recentering atlas]
\label{lem:v4v25-pointwise-atlas}
Let \(\mu_M\) be any positive normalized toron-parent measure.  If a
measurable conditional quantity satisfies
\[
 |F_{\nu,M}(\Theta)|\le C
 \qquad
 (\Theta\in\mathcal R_{\nu,M})
\]
with one constant independent of \(\nu,M\), then
\begin{equation}
 \left|
 \sum_{\nu\in\mathcal I_M}
 \int_{\mathcal R_{\nu,M}}
 F_{\nu,M}(\Theta)\,d\mu_M(\Theta)
 \right|
 \le C.
\label{eq:v4v25-pointwise-atlas}
\end{equation}
No factor \(|\mathcal I_M|\) occurs.
\end{lemma}

\begin{proof}
Use the triangle inequality inside the disjoint integral:
\[
 \sum_\nu
 \int_{\mathcal R_{\nu,M}}|F_{\nu,M}|\,d\mu_M
 \le
 C\sum_\nu\mu_M(\mathcal R_{\nu,M})
 =C.
\]
\end{proof}

\begin{firewall}[No artificial gain from small cells]
\label{fw:v4v25-no-small-cell-gain}
The proof uses only normalization and disjointness.  A cell's small toron
measure cannot be reused as a second response or cutoff payment after the
conditional quantity has already been normalized.  This is the precise
atlas analogue of the YM square-root ceiling and the NS unmarking rule.
\end{firewall}

\begin{proposition}[Exact measurable mode ancestry]
\label{prop:v4v25-mode-ancestry}
At every toron, the recentered template and its complement form one
orthogonal decomposition of every root Fourier sector.  The half-open rule
assigns torons on a recentering wall to exactly one cell.  Hence every mode
is integrated either as a transferred soft coordinate or as a hard
coordinate, never both, and Lemma~\ref{lem:v4v25-pointwise-atlas} sums the
conditional estimates without chart multiplicity.
\end{proposition}

\begin{proof}
The map \(k\mapsto\ell\) is bijective, and
\(\mathcal S_{\epsilon,L}\) and its complement partition the recentered
labels.  Half-open intervals make
\eqref{eq:v4v25-half-open-recenter} unique.  Apply the disjoint integral
identity.
\end{proof}

\subsection{Insertion into the V24 uniformity theorem}
\label{sec:v4v25-insertion}

On every recentering cell, the transferred label set is fixed and the hard
gap is \(2\epsilon\).  Thus the hard projector is constant in the recentered
Fourier frame, while the operator coefficients depend continuously on the
residual shifts in the compact cube
\([-\pi/M,\pi/M]^d\).  The V24 compactness argument applies with constants
independent of the cell because all cells use the same template and are
unitarily equivalent.

\begin{theorem}[Quadratic flat-sector atlas gate]
\label{thm:v4v25-flat-atlas-gate}
For the explicit toron-shifted lattice covariant Laplacian of V12, the soft
transfer, hard spectral gap, soft-rank density, and fixed-block relative
Hodge constants can be chosen uniformly in the recentering cell and lattice
volume, subject to the uniform local ellipticity and shape-regularity
hypotheses of V24.  The growing number of resonance cells causes no
multiplicative loss in the normalized toron integral.
\end{theorem}

\begin{proof}
Theorem~\ref{thm:v4v25-recentered-split} gives the fixed template and hard
gap.  Corollary~\ref{cor:v4v25-template-rank} gives the volume-uniform rank
density.  Unitary recentering transfers the V24 obstruction lower bound
between cells without changing it.  Lemma
\ref{lem:v4v25-pointwise-atlas} performs the normalized disjoint integral.
\end{proof}

\begin{firewall}[Toron derivatives across cell walls]
\label{fw:v4v25-toron-derivatives}
The half-open atlas is sufficient for measurable disintegration and for
background marks which hold the toron parent coordinate fixed.  The
recentered projector jumps by a momentum permutation at cell walls.
A classical derivative with respect to the toron coordinate therefore
requires either matching boundary traces between adjacent cells or a weak
derivative formulation in which the wall terms cancel.  V25 does not prove
that matching.
\end{firewall}

\begin{firewall}[The interacting Hessian is not the flat Laplacian]
\label{fw:v4v25-interacting-Hessian}
Theorem~\ref{thm:v4v25-flat-atlas-gate} uses the explicit shifted Fourier
symbol \eqref{eq:v4v25-shifted-spectrum}.  Nonlinear background curvature,
Higgs expectation values, Yukawa mixing, and dynamical metric coefficients
can couple Fourier labels and destroy exact momentum recentering.  Uniform
control of the complete interacting hard Hessian requires a perturbative
spectral comparison or a separate local resolvent argument.
\end{firewall}

\subsection{V25 result ledger and V26 target}
\label{sec:v4v25-conclusion}

V25 records the compulsory companion audit and removes the chart-cardinality
obstruction for the flat quadratic toron sector.  A global
toron-independent soft Fourier space would have to contain every root
momentum.  The correct construction recenters momentum separately on a
half-open disjoint atlas.  One fixed enlarged template then contains every
potentially sub-\(2\epsilon\) mode, its complement has a uniform hard gap,
and its rank density is \(C(L^{-4}+\epsilon^2)\) in four dimensions.
Unitary relabeling makes all conditional estimates cell-independent, and
normalization of the disjoint toron measure prevents a cell-count loss.

V26 should compare the complete interacting hard Hessian with the
recentered flat operator.  The target is an explicit relative form bound
\[
 |\langle u,(K_{\rm int}-K_{\rm flat})u\rangle|
 \le\vartheta\langle u,K_{\rm flat}u\rangle
 +c_\vartheta\|P_{\rm soft}u\|^2
\]
with \(\vartheta<1\), uniformly on the invariant good chart.  Such a bound
would preserve the hard gap without requiring an interacting eigenvector
atlas.

\clearpage
\section{V26: relative form stability of the interacting hard block}
\label{sec:v4v26}

V25 constructs a volume-uniform soft template for the flat toron Laplacian.
The complete interacting Hessian need not share its Fourier eigenvectors.
Building a second eigenvector atlas would duplicate the ancestry and would
be unstable under mode mixing.  The correct question is whether the
interacting quadratic form is a small perturbation on the fixed V25 hard
complement.

This note proves the abstract compression and Schur theorems needed for that
question.  It also proves a local-stencil perturbation bound for the
gauge-fixed Wilson Hessian on the V7 good chart.  If all non-small marginal
quadratic Standard-Model terms are already included in the signed Gaussian
parent, the remaining good-chart perturbation preserves the hard gap and
produces an exact, controlled Schur correction on the soft parent.  Higgs,
Yukawa, hypercharge, and metric residual bounds remain separate inputs.

\subsection{A fixed soft/hard compression}
\label{sec:v4v26-compression}

Let \(\mathcal H=\mathcal H_{\rm s}\oplus\mathcal H_{\rm h}\) be the
orthogonal V25 decomposition, with projections \(P_{\rm s},P_{\rm h}\).
Let \(K_0=K_0^*>0\) be the assembled flat quadratic parent and assume
\begin{equation}
 P_{\rm h}K_0P_{\rm h}\ge2\epsilon P_{\rm h}
\label{eq:v4v26-flat-hard-gap}
\end{equation}
for a fixed \(\epsilon>0\).  Let
\[
 K=K_0+Q,\qquad Q=Q^*,
\]
where \(Q\) is the residual interacting quadratic perturbation.

\begin{lemma}[Hard absorption and soft leakage]
\label{lem:v4v26-form-split}
If \(\|Q\|\le q\), then for
\(u=s+h\), \(s=P_{\rm s}u\), \(h=P_{\rm h}u\), and every
\(0<\vartheta<1\),
\begin{align}
 |\langle u,Qu\rangle|
 &\le
 \left(\frac{q}{2\epsilon}+\vartheta\right)
 \langle h,K_0h\rangle
\notag\\
 &\quad+
 \left(
 q+\frac{q^2}{2\epsilon\vartheta}
 \right)\|s\|^2.
\label{eq:v4v26-form-split}
\end{align}
\end{lemma}

\begin{proof}
The hard--hard term obeys
\[
 |\langle h,Qh\rangle|
 \le q\|h\|^2
 \le\frac{q}{2\epsilon}\langle h,K_0h\rangle.
\]
The soft--soft term is at most \(q\|s\|^2\).  For the two cross terms,
\[
 2|\langle s,Qh\rangle|
 \le2q\|s\|\|h\|
 \le
 \frac{q^2}{2\epsilon\vartheta}\|s\|^2
 +2\epsilon\vartheta\|h\|^2
 \le
 \frac{q^2}{2\epsilon\vartheta}\|s\|^2
 +\vartheta\langle h,K_0h\rangle.
\]
Add the three estimates.
\end{proof}

\begin{theorem}[Stability of the fixed hard gap]
\label{thm:v4v26-hard-gap}
If \(q\le\epsilon\), then
\begin{equation}
 K_{\rm hh}:=P_{\rm h}KP_{\rm h}
 \ge\epsilon P_{\rm h},
 \qquad
 \|K_{\rm hh}^{-1}\|\le\epsilon^{-1}.
\label{eq:v4v26-interacting-hard-gap}
\end{equation}
No interacting spectral projector is required.
\end{theorem}

\begin{proof}
For \(h\in\mathcal H_{\rm h}\),
\[
 \langle h,K_{\rm hh}h\rangle
 \ge
 \langle h,K_0h\rangle-q\|h\|^2
 \ge(2\epsilon-q)\|h\|^2
 \ge\epsilon\|h\|^2.
\]
The inverse bound follows from the spectral theorem.
\end{proof}

\subsection{The soft response is an exact Schur term}
\label{sec:v4v26-Schur}

Write the interacting operator in fixed blocks,
\[
 K=
 \begin{pmatrix}
 K_{\rm ss}&K_{\rm sh}\\
 K_{\rm hs}&K_{\rm hh}
 \end{pmatrix}.
\]
Under Theorem~\ref{thm:v4v26-hard-gap}, define
\begin{equation}
 \mathcal S_K
 =
 K_{\rm ss}-K_{\rm sh}K_{\rm hh}^{-1}K_{\rm hs}.
\label{eq:v4v26-soft-Schur}
\end{equation}

\begin{theorem}[Exact interacting hard integration]
\label{thm:v4v26-exact-hard-integration}
If \(K>0\) and \(q\le\epsilon\), then
\begin{align}
 \det K&=\det K_{\rm hh}\det\mathcal S_K,
\label{eq:v4v26-det-Schur}\\
 K^{-1}
 &=
 \begin{pmatrix}0&0\\0&K_{\rm hh}^{-1}\end{pmatrix}
 +
 \binom{1}{-K_{\rm hh}^{-1}K_{\rm hs}}
 \mathcal S_K^{-1}
 \begin{pmatrix}1&-K_{\rm sh}K_{\rm hh}^{-1}\end{pmatrix}.
\label{eq:v4v26-covariance-Schur}
\end{align}
The response
\begin{equation}
 \mathcal R_{\rm s}
 :=
 K_{\rm sh}K_{\rm hh}^{-1}K_{\rm hs}
\label{eq:v4v26-soft-response}
\end{equation}
is positive semidefinite and is retained once in the soft parent.
\end{theorem}

\begin{proof}
Apply the block \(LDU\) factorization used in
\eqref{eq:v4v19-LDU}, now with the hard block eliminated first.
It gives \eqref{eq:v4v26-det-Schur} and
\eqref{eq:v4v26-covariance-Schur}.  Moreover
\[
 \mathcal R_{\rm s}
 =
 (K_{\rm hh}^{-1/2}K_{\rm hs})^*
 (K_{\rm hh}^{-1/2}K_{\rm hs})\ge0.
\]
\end{proof}

\begin{corollary}[Quadratic size of the off-diagonal response]
\label{cor:v4v26-response-size}
If the flat parent is block diagonal in the V25 split and
\(\|Q\|\le q\le\epsilon\), then
\begin{equation}
 \|K_{\rm sh}\|\le q,
 \qquad
 \|\mathcal R_{\rm s}\|
 \le\frac{q^2}{\epsilon}.
\label{eq:v4v26-response-size}
\end{equation}
\end{corollary}

\begin{proof}
Block diagonality gives
\(K_{\rm sh}=P_{\rm s}QP_{\rm h}\), so its norm is at most \(q\).
Use \eqref{eq:v4v26-interacting-hard-gap} in
\eqref{eq:v4v26-soft-response}.
\end{proof}

\begin{firewall}[Off-diagonal mixing is not discarded]
\label{fw:v4v26-no-drop-mixing}
The fixed V25 subspaces need not reduce \(K\).  The mixed blocks are not
bounded and then omitted: they occur twice in the single positive response
\eqref{eq:v4v26-soft-response}.  Adding that response to the soft action
with the wrong sign, or also retaining it as an independent hard
determinant insertion, would break the exact Schur identity.
\end{firewall}

\subsection{A local-stencil perturbation theorem}
\label{sec:v4v26-local-stencil}

Work on a rescaled lattice of unit spacing and uniformly bounded incidence
degree.  Let \(\tau\) be a flat toron link field and \(U\) a good-chart link
field in tree gauge.  Assume
\begin{equation}
 \max_e d_G(U_e,\tau_e)\le\eta.
\label{eq:v4v26-link-distance}
\end{equation}
Let \(D_U\) be any nearest-neighbor covariant difference operator in a fixed
unitary representation, and let \(D_\tau\) be its flat counterpart.

\begin{lemma}[Covariant difference perturbation]
\label{lem:v4v26-difference-perturbation}
There is a representation- and dimension-dependent constant \(C_D\) such
that
\begin{equation}
 \|D_U-D_\tau\|\le C_D\eta.
\label{eq:v4v26-difference-perturbation}
\end{equation}
Consequently, for every \(0<\vartheta<1\),
\begin{align}
 \left|
 \|D_Uu\|^2-\|D_\tau u\|^2
 \right|
 &\le
 \vartheta\|D_\tau u\|^2
 +
 C_\vartheta\eta^2\|u\|^2.
\label{eq:v4v26-gradient-form-bound}
\end{align}
\end{lemma}

\begin{proof}
Each matrix entry of a unitary representation is Lipschitz on the compact
group.  Thus the coefficient difference on one oriented edge is at most
\(C\eta\).  Every row and column of the lattice difference matrix has a
uniformly bounded number of nonzero entries.  The Schur test proves
\eqref{eq:v4v26-difference-perturbation}.  Put
\(E=D_U-D_\tau\).  Then
\[
 \|D_Uu\|^2-\|D_\tau u\|^2
 =
 2\operatorname{Re}\langle D_\tau u,Eu\rangle+\|Eu\|^2.
\]
Young's inequality and \(\|E\|\le C_D\eta\) prove
\eqref{eq:v4v26-gradient-form-bound}.
\end{proof}

Let \(K_{\rm W}(U)\) be the gauge-fixed Wilson Hessian, including its
background-covariant gauge-fixing block and the matching ghost quadratic
operator, all expressed in one fixed Ward bookkeeping convention.

\begin{proposition}[Good-chart Wilson Hessian perturbation]
\label{prop:v4v26-Wilson-perturbation}
There are \(\eta_0,C_{\rm W}>0\) such that, when
\eqref{eq:v4v26-link-distance} holds with \(\eta\le\eta_0\),
\begin{equation}
 \|K_{\rm W}(U)-K_{\rm W}(\tau)\|
 \le C_{\rm W}\eta.
\label{eq:v4v26-Wilson-operator-bound}
\end{equation}
The same statement holds for a finite direct sum of fixed-representation
covariant scalar and Dirac hopping Hessians.
\end{proposition}

\begin{proof}
On a fixed compact group, the second derivative of the plaquette potential,
parallel transports, gauge-fixing coefficients, and ghost coefficients are
smooth functions of the finitely many links in a plaquette collar.  Taylor's
theorem bounds every local coefficient difference by \(C\eta\).  The
Hessian stencil has uniformly bounded range and a uniformly bounded number
of entries in every row and column.  The Schur test converts the local
coefficient bound into \eqref{eq:v4v26-Wilson-operator-bound}.  Scalar and
Dirac hopping matrices have the same finite-stencil property.
\end{proof}

\begin{corollary}[Invariant good plaquettes give a small quadratic
perturbation]
\label{cor:v4v26-good-plaquette-perturbation}
On a fixed rescaled block, the V7 tree-gauge estimate and the V23 harmonic
buffer give
\begin{equation}
 \eta\le C r_g,\qquad r_g=g^{1-\alpha}.
\label{eq:v4v26-eta-rg}
\end{equation}
Hence the gauge, ghost, and fixed-representation hopping part of the
quadratic perturbation obeys
\begin{equation}
 q_g^{\rm geom}\le Cg^{1-\alpha}=o(1)
\label{eq:v4v26-geometric-q}
\end{equation}
for \(0<\alpha<1\).
\end{corollary}

\begin{proof}
Lemma~\ref{lem:v4v7-small-curvature-chart} bounds the tree-gauge link
distance by a fixed-block multiple of the plaquette threshold.  Theorems
\ref{thm:v4v23-nonlinear-chart} and
\ref{thm:v4v24-uniform-chart} provide the coarse-field buffer.  Apply
Proposition~\ref{prop:v4v26-Wilson-perturbation}.
\end{proof}

\begin{firewall}[Physical scaling is performed before the Schur test]
\label{fw:v4v26-rescaled-lattice}
Equations \eqref{eq:v4v26-difference-perturbation}--%
\eqref{eq:v4v26-geometric-q} are level-rescaled estimates.  In physical
coordinates a covariant difference carries a factor \(a^{-1}\), and both
the flat gap and the perturbation scale accordingly.  Keeping the
\(a^{-1}\) only in the perturbation would create a false ultraviolet
divergence.
\end{firewall}

\subsection{Assembly with the signed Standard-Model parent}
\label{sec:v4v26-SM-assembly}

Let the complete residual quadratic perturbation be
\begin{equation}
 Q_{\rm SM}
 =
 Q_{\rm geom}
 +Q_{\rm Higgs}
 +Q_{\rm Yuk}
 +Q_{\rm hyp}
 +Q_{\rm met}
 +Q_{\rm ct},
\label{eq:v4v26-SM-residual}
\end{equation}
after the non-small running marginal quadratic terms and counterterms have
been inserted in \(K_0\).  Suppose
\begin{equation}
 \|Q_{\rm SM}\|
 \le
 q_g
 :=
 Cg^{1-\alpha}
 +q_{\rm Higgs}
 +q_{\rm Yuk}
 +q_{\rm hyp}
 +q_{\rm met}
 +q_{\rm ct}.
\label{eq:v4v26-total-q}
\end{equation}

\begin{theorem}[Conditional interacting hard-parent stability]
\label{thm:v4v26-SM-hard-stability}
If the recentered flat parent obeys
\eqref{eq:v4v26-flat-hard-gap} and
\begin{equation}
 q_g\le\epsilon,
\label{eq:v4v26-small-total-q}
\end{equation}
then:
\begin{enumerate}
\item the fixed V25 hard block remains positive with inverse norm at most
\(\epsilon^{-1}\);
\item its determinant and covariance factor exactly as in
Theorem~\ref{thm:v4v26-exact-hard-integration};
\item its entire mixing response is the one positive soft Schur term
\(\mathcal R_{\rm s}\), with norm at most \(q_g^2/\epsilon\) when the flat
parent is block diagonal; and
\item the V16 resolvent identities may be applied on the fixed hard
subspace, without differentiating an interacting spectral projector.
\end{enumerate}
\end{theorem}

\begin{proof}
Items one through three are Theorem~\ref{thm:v4v26-hard-gap},
Theorem~\ref{thm:v4v26-exact-hard-integration}, and Corollary
\ref{cor:v4v26-response-size}.  The V16 identities require an invertible
hard operator and differentiate the operator and its inverse on a fixed
space.  Equation \eqref{eq:v4v26-interacting-hard-gap} supplies the required
inverse bound.
\end{proof}

\begin{firewall}[Small residual means after marginal transport]
\label{fw:v4v26-marginal-transport}
Condition \eqref{eq:v4v26-small-total-q} cannot be obtained by declaring a
physical Higgs mass, a Yukawa-induced quadratic term, the hypercharge
polarization, or a metric kinetic response to be small when it is marginal.
Every non-small quadratic term must first be transported into the signed
parent \(K_0\) with its determinant ancestry.  Only its irrelevant or
good-chart variation belongs to \(Q_{\rm SM}\).
\end{firewall}

\begin{firewall}[The open Standard-Model terms]
\label{fw:v4v26-open-SM-terms}
V26 proves \eqref{eq:v4v26-geometric-q} for the finite-stencil geometric
gauge, ghost, and fixed-representation hopping variations.  It does not
prove that
\(q_{\rm Higgs},q_{\rm Yuk},q_{\rm hyp},q_{\rm met},q_{\rm ct}\)
satisfy \eqref{eq:v4v26-small-total-q} along a continuum trajectory.
In particular, the ordinary hypercharge Landau sign remains an ultraviolet
obstruction rather than a perturbative remainder.
\end{firewall}

\subsection{V26 result ledger and V27 target}
\label{sec:v4v26-conclusion}

V26 replaces an interacting eigenvector atlas by a fixed-space relative form
argument.  A residual perturbation smaller than the V25 hard gap preserves
the hard inverse.  Off-diagonal soft--hard mixing is retained exactly as a
single positive Schur response of quadratic size.  On the invariant V7 good
chart, finite-stencil smoothness proves
\(q_g^{\rm geom}=O(g^{1-\alpha})\) for the gauge-fixed Wilson, ghost,
scalar-hopping, and Dirac-hopping geometric variations.

The unresolved inequality is now confined to the non-geometric residuals
in \eqref{eq:v4v26-total-q}.  V27 should begin with the Higgs sector:
separate the running vacuum expectation and mass matrix into \(K_0\), then
prove a marked operator bound for the quartic and gauge--Higgs residual on
the mesoscopic good chart.  Yukawa mixing should be added only after the
bosonic Higgs hard gap and its soft Schur ancestry are fixed.

\clearpage
\section{V27: the marginal Higgs quartic cannot be a small Gaussian residual}
\label{sec:v4v27}

V26 isolates the non-geometric quadratic residuals by an operator norm
\(q_g\).  The first proposed target was the Higgs quartic.  A direct scaling
test shows that this target has the wrong type.  On the mesoscopic chart
\(\|\phi-v\|\le R_g=g^{-\alpha}\), the Hessian variation of a marginal
quartic is of order \(\lambda R_g^2\).  It is not controlled by the gauge
factor \(g^{1-\alpha}\), and for fixed positive \(\lambda\) it grows.

The correct upstream object is the complete marginal Higgs law, exactly as
the V3--V4 entropy diagnosis required.  This note proves the no-go theorem,
constructs the positive finite-cutoff quartic parent in the uniformly convex
hard phase, and derives its exact first and second connected responses and
Brascamp--Lieb bounds.  The broken phase and the coupled gauge--Higgs parent
remain separate gates.

\subsection{Exact quartic Hessian and the mesoscopic no-go}
\label{sec:v4v27-no-go}

Treat the complex Higgs doublet as a vector in \(\mathbb R^4\).  Use the
local potential
\begin{equation}
 V_{\rm H}(\phi)
 =
 \frac{m^2}{2}|\phi|^2+\frac{\lambda}{4}|\phi|^4,
 \qquad\lambda\ge0.
\label{eq:v4v27-Higgs-potential}
\end{equation}
Its gradient and Hessian are
\begin{align}
 \nabla V_{\rm H}(\phi)
 &=
 m^2\phi+\lambda|\phi|^2\phi,
\label{eq:v4v27-Higgs-gradient}\\
 D^2V_{\rm H}(\phi)
 &=
 m^2 1+\lambda\left(|\phi|^2 1+2\phi\otimes\phi\right).
\label{eq:v4v27-Higgs-Hessian}
\end{align}

\begin{lemma}[Hessian variation around a running center]
\label{lem:v4v27-Hessian-variation}
For \(\phi=v+h\),
\begin{equation}
 \left\|
 D^2V_{\rm H}(v+h)-D^2V_{\rm H}(v)
 \right\|
 \le
 C\lambda\left(|v|\,|h|+|h|^2\right).
\label{eq:v4v27-Hessian-variation}
\end{equation}
At \(v=0\) and \(h=Re\) with \(|e|=1\), the norm is exactly
\begin{equation}
 \left\|
 D^2V_{\rm H}(Re)-D^2V_{\rm H}(0)
 \right\|
 =3\lambda R^2.
\label{eq:v4v27-Hessian-lower-test}
\end{equation}
\end{lemma}

\begin{proof}
Use
\[
 \big||v+h|^2-|v|^2\big|
 \le2|v|\,|h|+|h|^2
\]
and
\[
 \|(v+h)\otimes(v+h)-v\otimes v\|
 \le2|v|\,|h|+|h|^2
\]
in \eqref{eq:v4v27-Higgs-Hessian}.  For \(v=0\),
\[
 D^2V_{\rm H}(Re)-D^2V_{\rm H}(0)
 =
 \lambda R^2(1+2e\otimes e),
\]
whose eigenvalue in direction \(e\) is \(3\lambda R^2\) and whose other
eigenvalues are \(\lambda R^2\).
\end{proof}

\begin{theorem}[Marginal-quartic small-perturbation no-go]
\label{thm:v4v27-quartic-no-go}
Let \(R_g=g^{-\alpha}\), \(\alpha>0\).  Any uniform operator bound for the
quartic Hessian variation on \(\{|h|\le R_g\}\) must satisfy
\begin{equation}
 q_{\rm H}(g)\ge3\lambda_gR_g^2
 =3\lambda_g g^{-2\alpha}.
\label{eq:v4v27-quartic-no-go}
\end{equation}
Therefore a vanishing perturbation
\(q_{\rm H}(g)=o(1)\) requires
\begin{equation}
 \lambda_g g^{-2\alpha}\longrightarrow0.
\label{eq:v4v27-extra-quartic-scaling}
\end{equation}
Even the weaker V26 hard-gap condition \(q_{\rm H}(g)\le\epsilon\) requires
\(\lambda_gg^{-2\alpha}\le\epsilon/3\) on this test.  Neither restriction
follows from small gauge coupling for a fixed positive or merely
logarithmically decreasing marginal quartic.
\end{theorem}

\begin{proof}
The test field \(h=R_ge\) lies in the chart.  Equation
\eqref{eq:v4v27-Hessian-lower-test} forces
\eqref{eq:v4v27-quartic-no-go}.  Requiring that lower bound to
vanish gives \eqref{eq:v4v27-extra-quartic-scaling}; requiring it merely to
be at most \(\epsilon\) gives
\(\lambda_gg^{-2\alpha}\le\epsilon/3\).  Any inverse power of
\(\log g^{-1}\) decays more slowly than the positive power \(g^{2\alpha}\),
so logarithmic decrease is insufficient for either bound at small \(g\).
\end{proof}

\begin{firewall}[Shrinking the chart would remove the gauge estimate]
\label{fw:v4v27-no-chart-cheat}
The radius \(R_g=g^{-\alpha}\) is what makes the complement of the
small-field chart stretched-exponentially rare in V7 and V20.  Replacing it
by a fixed or vanishing radius solely to make
\(\lambda R_g^2\) small would forfeit that cutoff-tail estimate.  The
marginal quartic must be reorganized rather than hidden by a different
chart.
\end{firewall}

\subsection{Promotion to a quartic marginal parent}
\label{sec:v4v27-quartic-parent}

On a finite lattice let
\begin{equation}
 S_{\rm H}(\phi;B)
 =
 \frac12\langle\phi,K_{\rm H}(B)\phi\rangle
 +
 \frac{\lambda}{4}\sum_x|\phi_x|^4
 -
 \langle J(B),\phi\rangle,
\label{eq:v4v27-Higgs-action}
\end{equation}
where \(\phi_x\in\mathbb R^4\), \(\lambda\ge0\), and
\(K_{\rm H}(B)=K_{\rm H}(B)^*\).  The linear term permits a running center
or boundary source.  Assume
\begin{equation}
 K_{\rm H}(B)\ge\kappa\,1,
 \qquad \kappa>0,
\label{eq:v4v27-Higgs-gap}
\end{equation}
uniformly in the hard background chart.  Define
\begin{equation}
 d\mu_{\rm H,B}(\phi)
 =
 Z_{\rm H}(B)^{-1}e^{-S_{\rm H}(\phi;B)}\,d\phi.
\label{eq:v4v27-Higgs-parent}
\end{equation}

\begin{theorem}[Positive strongly convex Higgs parent]
\label{thm:v4v27-convex-parent}
Under \eqref{eq:v4v27-Higgs-gap}, the finite-volume normalization
\(Z_{\rm H}(B)\) is finite and positive, and
\begin{equation}
 D_\phi^2S_{\rm H}(\phi;B)
 =
 K_{\rm H}(B)
 +
 \lambda\,\operatorname{diag}_x
 \left(|\phi_x|^2 1+2\phi_x\otimes\phi_x\right)
 \ge\kappa\,1.
\label{eq:v4v27-action-Hessian}
\end{equation}
Every conditional block law obtained by fixing exterior Higgs variables has
the same lower Hessian bound \(\kappa\).
\end{theorem}

\begin{proof}
The quartic term is nonnegative and grows as \(\|\phi\|^4\), while the
quadratic part has lower bound
\(\kappa\|\phi\|^2/2-\|J\|\|\phi\|\).  Thus the density is integrable and
strictly positive.  Equation \eqref{eq:v4v27-action-Hessian} follows from
\eqref{eq:v4v27-Higgs-Hessian}; the quartic Hessian is positive
semidefinite.  Fixing exterior variables changes the interior action by
linear terms and by the principal interior compression of
\(K_{\rm H}\).  A principal compression of an operator bounded below by
\(\kappa\) has the same lower bound.
\end{proof}

\begin{corollary}[Brascamp--Lieb and Poincare bounds]
\label{cor:v4v27-BL}
For every smooth \(F\),
\begin{equation}
 \operatorname{Var}_{\mu_{\rm H,B}}(F)
 \le
 \mathbb E_{\mu_{\rm H,B}}
 \left\langle
 \nabla F,
 (D_\phi^2S_{\rm H})^{-1}\nabla F
 \right\rangle
 \le
 \kappa^{-1}
 \mathbb E_{\mu_{\rm H,B}}\|\nabla F\|^2.
\label{eq:v4v27-BL}
\end{equation}
The same inequalities hold for every conditional block law, uniformly in
its exterior field.
\end{corollary}

\begin{proof}
The first inequality is the finite-dimensional Brascamp--Lieb variance
inequality for a \(C^2\) strictly convex density.  The second follows from
\eqref{eq:v4v27-action-Hessian}.  The conditional assertion follows from
the last statement of Theorem~\ref{thm:v4v27-convex-parent}.
\end{proof}

\subsection{Exact responses of the non-Gaussian parent}
\label{sec:v4v27-responses}

For a background direction \(v\), put
\[
 X_v=D_BS_{\rm H}[v],
 \qquad
 Y_{v,w}=D_B^2S_{\rm H}[v,w].
\]

\begin{theorem}[Connected Higgs-parent response identities]
\label{thm:v4v27-response-identities}
Whenever differentiation under the finite-dimensional integral is valid,
\begin{align}
 D_B[-\log Z_{\rm H}][v]
 &=
 \mathbb E_{\mu_{\rm H,B}}X_v,
\label{eq:v4v27-first-response}\\
 D_B^2[-\log Z_{\rm H}][v,w]
 &=
 \mathbb E_{\mu_{\rm H,B}}Y_{v,w}
 -
 \operatorname{Cov}_{\mu_{\rm H,B}}(X_v,X_w).
\label{eq:v4v27-second-response}
\end{align}
If
\begin{equation}
 \|\nabla_\phi X_v\|\le M_1\|v\|,
 \qquad
 |\mathbb E Y_{v,w}|\le M_2\|v\|\|w\|
\label{eq:v4v27-response-input}
\end{equation}
uniformly, then
\begin{equation}
 \left|
 D_B^2[-\log Z_{\rm H}][v,w]
 \right|
 \le
 \left(M_2+\kappa^{-1}M_1^2\right)
 \|v\|\|w\|.
\label{eq:v4v27-response-bound}
\end{equation}
\end{theorem}

\begin{proof}
Differentiate \(Z_{\rm H}=\int e^{-S_{\rm H}}d\phi\) once and twice and
divide by \(Z_{\rm H}\).  The disconnected product generated by the second
derivative of the logarithm subtracts the product of first moments, giving
the covariance in \eqref{eq:v4v27-second-response}.  The polarized
Brascamp--Lieb inequality and Cauchy--Schwarz give
\[
 |\operatorname{Cov}(X_v,X_w)|
 \le
 \kappa^{-1}
 \left(\mathbb E\|\nabla X_v\|^2\right)^{1/2}
 \left(\mathbb E\|\nabla X_w\|^2\right)^{1/2}.
\]
Insert \eqref{eq:v4v27-response-input}.
\end{proof}

\begin{firewall}[The covariance is one genuine quadratic response]
\label{fw:v4v27-one-covariance}
The last term in \eqref{eq:v4v27-second-response} is the covariance of two
actual first-derivative occurrences inside one second derivative of the
logarithm.  It may not be replaced by the square of an unrelated one-mark
bound or counted once in the Higgs parent and again in a gauge--Higgs
seagull packet.
\end{firewall}

\subsection{A conditional influence estimate}
\label{sec:v4v27-influence}

Partition the Higgs variables into an interior block \(I\) and exterior
variables \(\eta\).  Suppose the quadratic hopping part is
\[
 \frac12\langle\phi_I,K_{II}\phi_I\rangle
 +\langle\phi_I,K_{I\partial}\eta\rangle
 +S_\partial(\eta).
\]
For an interior observable \(F\), let
\(\mathbb E_I^\eta F\) denote conditional expectation under the quartic
parent.

\begin{proposition}[Boundary influence under strong convexity]
\label{prop:v4v27-boundary-influence}
For an exterior direction \(h\),
\begin{equation}
 \left|
 D_\eta\mathbb E_I^\eta F[h]
 \right|
 \le
 \kappa^{-1}
 \left(
 \mathbb E_I^\eta\|\nabla_{\phi_I}F\|^2
 \right)^{1/2}
 \|K_{I\partial}h\|.
\label{eq:v4v27-boundary-influence}
\end{equation}
\end{proposition}

\begin{proof}
Differentiating the normalized conditional expectation gives
\[
 D_\eta\mathbb E_I^\eta F[h]
 =
 -\operatorname{Cov}_I^\eta
 \left(
 F,\langle\phi_I,K_{I\partial}h\rangle
 \right).
\]
Apply the polarized form of
\eqref{eq:v4v27-BL}.  The gradient of the second argument is the constant
vector \(K_{I\partial}h\), which gives
\eqref{eq:v4v27-boundary-influence}.
\end{proof}

\begin{corollary}[A concrete Higgs mixing criterion]
\label{cor:v4v27-Dobrushin-criterion}
If the block decomposition satisfies
\begin{equation}
 \sup_I
 \sum_{J\ne I}
 \kappa^{-1}\|K_{IJ}\|
 <1,
\label{eq:v4v27-Dobrushin-criterion}
\end{equation}
then the matrix of single-block Lipschitz influences is a strict
contraction.  Establishing the corresponding connected polymer decay still
requires iterating this contraction with the actual local observables.
\end{corollary}

\begin{proof}
Proposition~\ref{prop:v4v27-boundary-influence} bounds the influence from
block \(J\) to block \(I\) by
\(\kappa^{-1}\|K_{IJ}\|\).  The displayed row sum is therefore below one.
\end{proof}

\begin{firewall}[A row-sum criterion is not yet the full cluster theorem]
\label{fw:v4v27-Dobrushin-open}
Corollary~\ref{cor:v4v27-Dobrushin-criterion} proves contraction of one
conditional response.  It does not by itself establish the all-orders
connected-cumulant estimate, the two-background-mark KP norm, or uniformity
when the hopping is marginal rather than diagonally dominant.  Those steps
must retain the complete quartic conditional law.
\end{firewall}

\subsection{Gauge--Higgs scaling after quartic promotion}
\label{sec:v4v27-gauge-Higgs}

Once the entire quartic law is in the parent, the remaining local
gauge--Higgs variations have different scaling.  Expanding
\(\|D_U\phi\|^2\) around a running center \(v\) produces a gauge-boson
quadratic term of size \(g^2|v|^2\), which belongs in \(K_0\), a mixed
gauge--Higgs term of size \(g|h|\), and a residual gauge quadratic term of
size \(g^2|h|^2\).  On \(|h|\le R_g\),
\begin{equation}
 gR_g=g^{1-\alpha}\to0,
 \qquad
 g^2R_g^2=g^{2-2\alpha}\to0
\label{eq:v4v27-gauge-Higgs-small}
\end{equation}
for \(0<\alpha<1\).  These geometric residuals are compatible with V26.
The pure quartic Hessian is absent from \(Q_{\rm SM}\) because it is now
part of the non-Gaussian marginal parent.

\begin{firewall}[Yukawa is not a positive quartic perturbation]
\label{fw:v4v27-Yukawa-separate}
A Yukawa factor is Grassmann-valued before fermion integration and generally
complex after taking a chiral determinant.  Positivity and
Brascamp--Lieb for \(\mu_{\rm H,B}\) do not control its phase.  It must be
inserted as the anomaly-normalized determinant/interaction packet of
V15--V16 and tested against the quartic parent without changing the latter
into a signed probability measure.
\end{firewall}

\begin{firewall}[Broken symmetry is not covered by global convexity]
\label{fw:v4v27-broken-phase}
Condition \eqref{eq:v4v27-Higgs-gap} is natural for a uniformly convex hard
unbroken sector.  A negative running mass can create a Mexican-hat
potential, Goldstone directions, and several gauge-related local charts.
V27 does not prove a global convex parent in that regime.  The radial hard
mode, gauge orbit, and soft vacuum manifold must first be disintegrated
without double counting, analogously to the toron construction.
\end{firewall}

\subsection{V27 result ledger and V28 target}
\label{sec:v4v27-conclusion}

V27 proves that a marginal Higgs quartic cannot be placed in a vanishing V26
Gaussian residual on the mesoscopic chart unless the extra, unjustified
scaling \(\lambda_g=o(g^{2\alpha})\) holds; preserving only a fixed hard gap
still requires \(\lambda_g=O(g^{2\alpha})\) with a sufficiently small
coefficient.  It promotes the complete
positive quartic law into the finite-cutoff marginal parent.  In the
uniformly convex hard phase, that parent is strongly log-concave, all
conditional block laws retain the same Hessian gap, and exact first and
second background responses obey Brascamp--Lieb bounds.  Gauge--Higgs
geometric residuals remain small after the non-small mass and quartic pieces
are transported.

V28 should turn the one-block influence estimate into a connected
two-marked polymer theorem for the quartic parent under
\eqref{eq:v4v27-Dobrushin-criterion}.  The proof must use one coupled
conditional-resampling chain so the two marks remain actual occurrences.
If diagonal dominance fails at a marginal hopping strength, the programme
must instead block to a scale where the hard conditional gap dominates the
boundary influence.

\clearpage
\section{V28: a two-mark hopping expansion over the quartic Higgs parent}
\label{sec:v4v28}

V27 proves one-block influence under strong convexity, but a Dobrushin row
sum alone was not promoted to an all-orders cumulant theorem.  This note
takes a more explicit route in the diagonally dominant regime.  The exact
onsite quartic laws form a product reference.  Every off-diagonal Higgs
hopping is retained as one bond activity, and selected bonds decompose
uniquely into connected components.

The fields are unbounded, so a supremum estimate is unavailable.  Strong
convexity gives uniform local subgaussian moments.  Since the lattice degree
is bounded, Finner's product-measure inequality uses only one fixed moment
order, independent of polymer size.  Background differentiation of the
normalized onsite laws creates an actual centered score at one site or an
actual second-density score at one site; two first scores at distinct sites
are retained separately.  These placements yield the V15 marked KP norm.

\subsection{Onsite product parent and hopping bonds}
\label{sec:v4v28-product-parent}

Let \(G=(\Lambda,E)\) be a finite graph of maximum degree \(D\).  At each
site \(x\), let \(\phi_x\in\mathbb R^r\), with \(r=4\) for the Standard
Model Higgs doublet.  Let
\begin{equation}
 d\mu_{x,B}(\phi_x)
 =
 z_x(B)^{-1}e^{-V_x(B,\phi_x)}\,d\phi_x,
\label{eq:v4v28-onsite-law}
\end{equation}
where
\begin{equation}
 D_{\phi_x}^2V_x(B,\phi_x)\ge\kappa\,1,
 \qquad\kappa>0.
\label{eq:v4v28-onsite-convexity}
\end{equation}
Put
\[
 \mu_{0,B}=\bigotimes_{x\in\Lambda}\mu_{x,B}.
\]
For an unoriented edge \(e=\{x,y\}\), let
\begin{equation}
 H_e(B,\phi_x,\phi_y)
 =
 \langle\phi_x,J_e(B)\phi_y\rangle,
 \qquad
 b_e=e^{-H_e}-1.
\label{eq:v4v28-bond-activity}
\end{equation}
The hopping partition function relative to the onsite parent is
\begin{equation}
 \mathcal Z_{\rm hop}(B)
 =
 \mathbb E_{0,B}
 \prod_{e\in E}(1+b_e)
 =
 \mathbb E_{0,B}
 e^{-\sum_{e\in E}H_e}.
\label{eq:v4v28-hopping-partition}
\end{equation}

\subsection{Uniform local moments}
\label{sec:v4v28-moments}

Assume the onsite minimizers \(m_x(B)\) obey
\begin{equation}
 \sup_{x,B}|m_x(B)|\le M_0.
\label{eq:v4v28-minimizer-bound}
\end{equation}

\begin{lemma}[Subgaussian onsite moment]
\label{lem:v4v28-subgaussian}
There are \(c_r,C_r>0\), depending only on \(r\), such that
\begin{equation}
 \sup_{x,B}
 \mathbb E_{\mu_{x,B}}
 \exp\left(c_r\kappa|\phi_x-m_x(B)|^2\right)
 \le C_r.
\label{eq:v4v28-subgaussian}
\end{equation}
Consequently, for every fixed \(p<\infty\),
\begin{equation}
 \|\phi_x\|_{L^p(\mu_{x,B})}
 \le C_{r,p}\left(M_0+\kappa^{-1/2}\right).
\label{eq:v4v28-polynomial-moment}
\end{equation}
\end{lemma}

\begin{proof}
The Bakry--Emery criterion applied to
\eqref{eq:v4v28-onsite-convexity} gives the logarithmic Sobolev inequality
with constant \(\kappa^{-1}\), hence Gaussian concentration for every
one-Lipschitz function:
\[
 \mathbb P(F-\mathbb EF\ge t)
 \le e^{-\kappa t^2/2}.
\]
The integration-by-parts identity
\[
 r
 =
 \mathbb E\langle\phi_x-m_x,\nabla V_x(\phi_x)\rangle
\]
and strong monotonicity of \(\nabla V_x\) give
\[
 \mathbb E|\phi_x-m_x|^2\le r/\kappa.
\]
Apply concentration to
\(F(\phi)=|\phi-m_x|\), whose mean is at most
\(\sqrt{r/\kappa}\).  Integrating the resulting Gaussian tail proves
\eqref{eq:v4v28-subgaussian} for a sufficiently small \(c_r>0\).
Polynomial moments follow from the exponential moment and
\(|\phi|\le|m_x|+|\phi-m_x|\).
\end{proof}

\begin{firewall}[Scope of the onsite moment estimate]
\label{fw:v4v28-onsite-moment-scope}
Lemma~\ref{lem:v4v28-subgaussian} uses the full uniform strong-convexity
bound \eqref{eq:v4v28-onsite-convexity} and the minimizer bound
\eqref{eq:v4v28-minimizer-bound}.  It does not apply across an untransferred
Goldstone or negative-mass direction.  No upper Hessian bound or small
quartic coupling is required in the convex hard sector.
\end{firewall}

Fix
\begin{equation}
 p_*=4(D+2).
\label{eq:v4v28-fixed-moment-order}
\end{equation}

\begin{lemma}[Small bond norm]
\label{lem:v4v28-small-bond}
Suppose
\begin{equation}
 p_*\|J_e\|
 \left(M_0^2+\kappa^{-1}\right)
 \le c_*
\label{eq:v4v28-small-hopping}
\end{equation}
for a sufficiently small \(c_*=c_*(r)\).  Then
\begin{equation}
 \|b_e\|_{L^{p_*}(\mu_{x,B}\otimes\mu_{y,B})}
 \le
 C_*
 \|J_e\|
 \left(M_0^2+\kappa^{-1}\right)
 =:\jmath_e.
\label{eq:v4v28-bond-Lp}
\end{equation}
\end{lemma}

\begin{proof}
Use
\[
 |e^{-t}-1|\le |t|e^{|t|}
\]
and
\[
 |H_e|
 \le
 \frac{\|J_e\|}{2}
 \left(|\phi_x|^2+|\phi_y|^2\right).
\]
Holder's inequality, Lemma~\ref{lem:v4v28-subgaussian}, and
\eqref{eq:v4v28-small-hopping} bound the required exponential quadratic
moments.  The remaining factor
\(|\phi_x||\phi_y|\) is bounded by
\eqref{eq:v4v28-polynomial-moment}.  This gives
\eqref{eq:v4v28-bond-Lp}.
\end{proof}

\subsection{Exact connected-bond polymerization}
\label{sec:v4v28-polymerization}

For a nonempty edge set \(A\subset E\), let \(V(A)\) be its incident
vertices.  Call \(A\) connected when the graph \((V(A),A)\) is connected.
For \(|X|\ge2\), define
\begin{equation}
 w_{\rm H}(X)
 =
 \sum_{\substack{A\subset E:\ A\ {\rm connected}\\V(A)=X}}
 \mathbb E_{0,B}\prod_{e\in A}b_e.
\label{eq:v4v28-Higgs-polymer}
\end{equation}
Two such polymers are compatible when their vertex sets are disjoint.

\begin{theorem}[Exact hopping polymer gas]
\label{thm:v4v28-exact-hopping-gas}
One has
\begin{equation}
 \mathcal Z_{\rm hop}(B)
 =
 \sum_{\substack{\mathcal X\ {\rm finite}\\
                  \text{pairwise vertex-disjoint}}}
 \prod_{X\in\mathcal X}w_{\rm H}(X),
\label{eq:v4v28-exact-hopping-gas}
\end{equation}
where the empty family contributes one.
\end{theorem}

\begin{proof}
Expand \(\prod_{e\in E}(1+b_e)\) by the selected edge set
\(A\subset E\).  Decompose \(A\) uniquely into its graph-connected
components \(A_1,\ldots,A_k\).  Their incident vertex sets are disjoint.
Because \(\mu_{0,B}\) is a site-product measure, the corresponding products
of bond functions are independent, and
\[
 \mathbb E_{0,B}\prod_{e\in A}b_e
 =
 \prod_{j=1}^k
 \mathbb E_{0,B}\prod_{e\in A_j}b_e.
\]
Grouping connected components with the same incident set gives
\eqref{eq:v4v28-exact-hopping-gas}.
\end{proof}

\begin{lemma}[Finner bound for one connected edge set]
\label{lem:v4v28-Finner}
If \(\jmath=\sup_e\jmath_e\le1\), then
\begin{equation}
 \left|
 \mathbb E_{0,B}\prod_{e\in A}b_e
 \right|
 \le
 \prod_{e\in A}\jmath_e
 \le
 \jmath^{|A|}.
\label{eq:v4v28-Finner}
\end{equation}
\end{lemma}

\begin{proof}
Finner's product-measure inequality states that
\[
 \mathbb E\prod_{e\in A}|f_e(\phi_x,\phi_y)|
 \le
 \prod_{e\in A}\|f_e\|_{L^{p_e}}
\]
whenever, at every site \(x\),
\(\sum_{e\ni x}p_e^{-1}\le1\).
Take every \(p_e=p_*\).  Since at most \(D\) selected edges meet a site and
\(D/p_*<1\), the condition holds.  Use
\eqref{eq:v4v28-bond-Lp} and monotonicity of probability-space \(L^p\)
norms.
\end{proof}

\begin{proposition}[Exponential unmarked activity]
\label{prop:v4v28-unmarked-activity}
There is \(C_D\) such that
\begin{equation}
 |w_{\rm H}(X)|
 \le
 (C_D\jmath)^{|X|-1}
\label{eq:v4v28-unmarked-activity}
\end{equation}
for every connected \(X\) with \(|X|\ge2\), provided
\(C_D\jmath<1\).
\end{proposition}

\begin{proof}
A connected edge set on \(X\) contains at least \(|X|-1\) edges.  The
induced graph has at most \(D|X|/2\) edges, so the number of its edge subsets
is at most \(2^{D|X|/2}\).  Sum
\eqref{eq:v4v28-Finner} over connected edge sets and absorb this exponential
count into \(C_D^{|X|-1}\).
\end{proof}

\subsection{Density scores and actual background marks}
\label{sec:v4v28-scores}

For a unit background direction \(v\), define the centered onsite score
\begin{equation}
 \sigma_x[v]
 =
 D_BV_x[v]
 -
 \mathbb E_{\mu_{x,B}}D_BV_x[v].
\label{eq:v4v28-first-score}
\end{equation}
For two directions \(v_1,v_2\), define the genuine second density score
\begin{equation}
 r_x[v_1,v_2]
 =
 \sigma_x[v_1]\sigma_x[v_2]
 -
 D_B\sigma_x[v_1][v_2].
\label{eq:v4v28-second-score}
\end{equation}

\begin{lemma}[Exact first and second density derivatives]
\label{lem:v4v28-density-derivatives}
The normalized onsite law satisfies
\begin{align}
 D_Bd\mu_{x,B}[v]
 &=
 -\sigma_x[v]\,d\mu_{x,B},
\label{eq:v4v28-density-D1}\\
 D_B^2d\mu_{x,B}[v_1,v_2]
 &=
 r_x[v_1,v_2]\,d\mu_{x,B}.
\label{eq:v4v28-density-D2}
\end{align}
In particular,
\(\mathbb E\sigma_x[v]=0\) and
\(\mathbb E r_x[v_1,v_2]=0\).
\end{lemma}

\begin{proof}
Differentiate
\(z_x^{-1}e^{-V_x}\).  The derivative of \(\log z_x\) is
\(-\mathbb E D_BV_x\), giving
\eqref{eq:v4v28-density-D1}.  Differentiate that identity again:
\[
 D_{v_2}(-\sigma_{v_1}d\mu)
 =
 \left(
 \sigma_{v_1}\sigma_{v_2}-D_{v_2}\sigma_{v_1}
 \right)d\mu.
\]
This is \eqref{eq:v4v28-density-D2}.  Integrating either derivative of the
identity \(\int d\mu_{x,B}=1\) gives the zero means.
\end{proof}

Let \(D_B\) also act on the bond matrices \(J_e(B)\).  Fix a mark scale
\(\rho>0\).  Assume that for \(m=0,1,2\),
\begin{equation}
 \rho^m
 \|D_B^mb_e\|_{L^{p_*}}
 \le\jmath
\label{eq:v4v28-marked-bond}
\end{equation}
uniformly in edges and unit directions, and that
\begin{equation}
 \rho\|\sigma_x\|_{L^{p_*}}
 +
 \rho^2\|r_x\|_{L^{p_*}}
 \le M_*
\label{eq:v4v28-score-bound}
\end{equation}
uniformly in sites and unit directions.

\begin{theorem}[Two-mark Higgs-polymer estimate]
\label{thm:v4v28-two-mark-polymer}
Under \eqref{eq:v4v28-marked-bond}--%
\eqref{eq:v4v28-score-bound}, there is
\(C=C(D,M_*)\) such that, for \(m=0,1,2\),
\begin{equation}
 \rho^m q_m^{\rm H}(X)
 \le
 C|X|^m(C\jmath)^{|X|-1},
\label{eq:v4v28-two-mark-activity}
\end{equation}
where \(q_m^{\rm H}(X)\) is the V15 derivative size of
\(w_{\rm H}(X)\).  The second derivative consists exactly of:
\begin{enumerate}
\item two derivatives on one actual bond;
\item one derivative on each of two actual bonds;
\item one genuine second density score \(r_x\);
\item two first density scores \(\sigma_x,\sigma_y\) on distinct sites; or
\item one bond derivative and one density score.
\end{enumerate}
\end{theorem}

\begin{proof}
Differentiate each term in \eqref{eq:v4v28-Higgs-polymer}.  Derivatives of
the integrand land on actual bond factors.  Derivatives of the product
reference law use Lemma~\ref{lem:v4v28-density-derivatives}; the product rule
gives one \(r_x\) or two scores at distinct sites.  A bond mark and a density
mark give the mixed case.  These five cases are disjoint and exhaustive.

A connected edge set with vertex set \(X\) has at most \(D|X|/2\) bond
slots and \(|X|\) density slots.  Thus there are at most \(C|X|\) one-mark
placements and \(C|X|^2\) two-mark placements.  At any site, at most \(D\)
bond factors and two score factors occur.  The choice
\(p_*=4(D+2)\) leaves enough Holder budget for all of them.  Finner's
inequality, \eqref{eq:v4v28-marked-bond}, and
\eqref{eq:v4v28-score-bound} bound each normalized marked term by
\(C(M_*)\jmath^{|A|}\).  Sum the connected edge sets exactly as in
Proposition~\ref{prop:v4v28-unmarked-activity}.
\end{proof}

\begin{firewall}[The second density score is not \(\sigma_x^2\)]
\label{fw:v4v28-second-score}
When both marks hit one normalized onsite law, the exact factor is
\(r_x=\sigma_x\sigma_x-D_B\sigma_x\), with polarized directions as in
\eqref{eq:v4v28-second-score}.  Dropping the derivative of the centered
score would violate normalization and turn a genuine second derivative into
the square of a first-moment estimate.
\end{firewall}

\subsection{The marked KP norm}
\label{sec:v4v28-KP}

\begin{theorem}[Quartic-parent hopping KP theorem]
\label{thm:v4v28-Higgs-KP}
For every \(a>0\), there is
\(\jmath_*(a,D,M_*)>0\) such that
\(\jmath<\jmath_*\) implies
\begin{equation}
 \|w_{\rm H}\|_{a,2;\rho}
 \le C_{a,D,M_*}\jmath.
\label{eq:v4v28-Higgs-KP}
\end{equation}
The hopping logarithm and its first two background derivatives therefore
converge absolutely and have the local and intensive response bounds of
Theorem~\ref{thm:v4v15-two-mark}.
\end{theorem}

\begin{proof}
A polymer vertex set is connected in \(G\).  Lemma
\ref{lem:v4v7-connected-count} bounds the number of cardinality-\(n\)
sets containing a fixed root by \((eD)^{n-1}\).  Insert
\eqref{eq:v4v28-two-mark-activity}:
\[
 \|w_{\rm H}\|_{a,2;\rho}
 \le
 C\sum_{n\ge2}
 n^2(eD)^{n-1}e^{an}(C\jmath)^{n-1}.
\]
This is \(O(\jmath)\) when \(\jmath\) is small enough.  Apply
Theorem~\ref{thm:v4v15-two-mark}.
\end{proof}

\begin{corollary}[Relation to the V27 row-sum criterion]
\label{cor:v4v28-Dobrushin-relation}
If the graph degree is bounded and
\[
 \sup_x\sum_{e\ni x}
 \|J_e\|\left(M_0^2+\kappa^{-1}\right)
\]
is sufficiently small in terms of
\(a,D,M_*\), then
\eqref{eq:v4v28-Higgs-KP} holds.  This is an explicit strengthened form of
the V27 diagonal-dominance condition, with the moment constants needed for
the unbounded quartic field displayed.
\end{corollary}

\begin{proof}
Lemma~\ref{lem:v4v28-small-bond} bounds \(\jmath_e\) by the corresponding
dimensionless hopping strength.  Bounded degree converts the row sum to the
uniform \(\jmath\) bound required in
Theorem~\ref{thm:v4v28-Higgs-KP}.
\end{proof}

\begin{firewall}[What the product expansion does not cover]
\label{fw:v4v28-scope}
The theorem requires a uniformly convex onsite hard law, bounded local
score moments, and sufficiently small hopping relative to that law.  It
does not cover a broken Higgs phase with Goldstone directions, a marginal
hopping outside the convergence ball, the chiral Yukawa determinant,
defect-sector hopping across a V11 repair collar, or metric-dependent
unbounded block degree.
\end{firewall}

\subsection{V28 result ledger and V29 target}
\label{sec:v4v28-conclusion}

V28 constructs an exact connected hopping polymer gas over the non-Gaussian
quartic Higgs parent.  Strong convexity supplies the fixed local moments
needed for unbounded fields, and bounded lattice degree makes the Finner
exponent independent of polymer size.  Exact derivatives of each normalized
onsite density produce one centered score, one genuine second-density score,
or two scores on distinct actual sites.  Together with actual bond
derivatives, these give a volume-uniform V15 two-mark KP norm at sufficiently
small dimensionless hopping.

The next gate is diagonal dominance.  V29 should test whether one rescaled
Schur blocking step improves the ratio of boundary hopping to the hard
conditional Higgs gap.  If it does, iterate until
\(\jmath<\jmath_*\); if the ratio is scale invariant, the complete
nearest-neighbor kinetic term must be promoted into a correlated quartic
parent and a different mixing theorem is required.

\clearpage
\section{V29: blocking amplifies contraction but does not create it}
\label{sec:v4v29}

V28 gives a complete two-mark hopping expansion when the dimensionless
bond strength is inside a strict convergence ball.  Its proposed fallback
was to block until boundary hopping became small relative to the conditional
hard gap.  This note tests that fallback exactly.

A massless nearest-neighbor Gaussian chain is already decisive.  On a block
of length \(\ell\), the Dirichlet gap is of order \(\ell^{-2}\), while
harmonic boundary influence remains order one and has an
\(\ell^{1/2}\) operator norm into the full block.  Naive blocking therefore
does not create diagonal dominance.  A valid positive result remains:
when a one-step Dobrushin matrix already has row sum \(\vartheta<1\), a
buffer of width \(L\) suppresses influence on the core by
\(\vartheta^L/(1-\vartheta)\).  The collar must remain an actual parent
variable; discarding it would merely hide the original boundary coupling.

\subsection{The exact massless chain test}
\label{sec:v4v29-chain}

Consider the real Gaussian action on sites
\(\{0,1,\ldots,\ell+1\}\),
\begin{equation}
 S(\phi)
 =
 \frac12\sum_{i=0}^{\ell}
 (\phi_{i+1}-\phi_i)^2,
\label{eq:v4v29-chain-action}
\end{equation}
with boundary values
\(\phi_0=a\), \(\phi_{\ell+1}=b\).  The interior precision is the
Dirichlet tridiagonal matrix
\begin{equation}
 A_\ell=
 \begin{pmatrix}
 2&-1&&0\\
 -1&2&\ddots&\\
 &\ddots&\ddots&-1\\
 0&&-1&2
 \end{pmatrix}.
\label{eq:v4v29-Dirichlet-matrix}
\end{equation}

\begin{lemma}[Exact conditional mean and gap]
\label{lem:v4v29-chain-mean-gap}
The conditional mean is
\begin{equation}
 m_i(a,b)
 =
 \frac{\ell+1-i}{\ell+1}a
 +
 \frac{i}{\ell+1}b,
 \qquad1\le i\le\ell,
\label{eq:v4v29-chain-mean}
\end{equation}
and
\begin{equation}
 \lambda_{\min}(A_\ell)
 =
 4\sin^2\left(\frac{\pi}{2(\ell+1)}\right)
 \sim\frac{\pi^2}{(\ell+1)^2}.
\label{eq:v4v29-chain-gap}
\end{equation}
\end{lemma}

\begin{proof}
The Euler--Lagrange equation is
\(2m_i-m_{i-1}-m_{i+1}=0\) with the prescribed endpoint values.  Its unique
solution is the affine interpolation
\eqref{eq:v4v29-chain-mean}.  The normalized sine vectors
\[
 v_i^{(k)}
 =
 \sin\left(\frac{\pi ki}{\ell+1}\right),
 \qquad1\le k\le\ell,
\]
diagonalize \(A_\ell\) with eigenvalues
\(4\sin^2(\pi k/(2(\ell+1)))\).  The smallest is \(k=1\), proving
\eqref{eq:v4v29-chain-gap}.
\end{proof}

\begin{proposition}[No contraction from naive super-blocking]
\label{prop:v4v29-blocking-no-go}
For the left boundary response,
\begin{align}
 \max_{1\le i\le\ell}
 \left|\frac{\partial m_i}{\partial a}\right|
 &=
 \frac{\ell}{\ell+1}\longrightarrow1,
\label{eq:v4v29-sup-influence}\\
 \left\|
 \frac{\partial m}{\partial a}
 \right\|_{\ell^2(\{1,\ldots,\ell\})}^2
 &=
 \frac{\ell(2\ell+1)}{6(\ell+1)}
 \sim\frac{\ell}{3}.
\label{eq:v4v29-L2-influence}
\end{align}
At every interior site,
\begin{equation}
 \frac{\partial m_i}{\partial a}
 +
 \frac{\partial m_i}{\partial b}
 =1.
\label{eq:v4v29-total-influence}
\end{equation}
Thus increasing \(\ell\) neither produces a strict sup-norm row contraction
nor makes the boundary-to-full-block operator small.  Simultaneously, the
conditional gap decreases as \(\ell^{-2}\).
\end{proposition}

\begin{proof}
Differentiate \eqref{eq:v4v29-chain-mean}.  The largest left coefficient is
at \(i=1\).  Also
\[
 \sum_{i=1}^{\ell}
 \left(\frac{\ell+1-i}{\ell+1}\right)^2
 =
 \frac1{(\ell+1)^2}\sum_{j=1}^{\ell}j^2
 =
 \frac{\ell(2\ell+1)}{6(\ell+1)}.
\]
The two affine coefficients sum to one, proving
\eqref{eq:v4v29-total-influence}.  Combine with
\eqref{eq:v4v29-chain-gap}.
\end{proof}

\begin{firewall}[A shrinking Dirichlet gap is not a generated infrared mass]
\label{fw:v4v29-no-generated-mass}
The block boundary condition removes the constant mode on one finite
interior and creates the gap
\eqref{eq:v4v29-chain-gap}.  That gap tends to zero with block size and
cannot be used as a uniform Higgs or gauge mass.  Calling it a mass would
repeat the V7 shell-gap error.
\end{firewall}

\subsection{The massive comparison}
\label{sec:v4v29-massive}

Add \(m^2\sum_{i=1}^{\ell}\phi_i^2/2\) to
\eqref{eq:v4v29-chain-action}.  Let \(\rho>0\) solve
\begin{equation}
 2+m^2=2\cosh\rho.
\label{eq:v4v29-rho}
\end{equation}

\begin{lemma}[Massive boundary decay]
\label{lem:v4v29-massive-decay}
With right boundary zero, the conditional mean is
\begin{equation}
 m_i(a,0)
 =
 a\,
 \frac{\sinh(\rho(\ell+1-i))}
 {\sinh(\rho(\ell+1))}.
\label{eq:v4v29-massive-mean}
\end{equation}
For sites at distance at least \(L\) from the left boundary,
\begin{equation}
 \left|\frac{\partial m_i}{\partial a}\right|
 \le C_\rho e^{-\rho L}.
\label{eq:v4v29-massive-buffer}
\end{equation}
However the first interior layer has response tending to
\(e^{-\rho}\), not to zero, as \(\ell\to\infty\).
\end{lemma}

\begin{proof}
The recurrence
\((2+m^2)m_i-m_{i-1}-m_{i+1}=0\) has solutions
\(e^{\pm\rho i}\).  Enforcing the two boundary values gives
\eqref{eq:v4v29-massive-mean}.  The ratio of hyperbolic sines is bounded by
a constant times \(e^{-\rho i}\), proving
\eqref{eq:v4v29-massive-buffer}.  At \(i=1\), the ratio tends to
\(e^{-\rho}\).
\end{proof}

The massive test identifies the correct geometry: a protected core gains
decay with its distance from the boundary, while the collar does not.  A
super-site containing both core and collar cannot claim the core decay as a
bound for the whole site.

\subsection{Abstract path-sum amplification}
\label{sec:v4v29-path-sum}

Let \(C=(c_{xy})\) be a nonnegative influence matrix on a graph, with
\begin{equation}
 \sup_x\sum_yc_{xy}\le\vartheta<1.
\label{eq:v4v29-strict-row-sum}
\end{equation}
Assume \(c_{xy}=0\) unless \(d(x,y)\le R_0\).  Define its resolvent
\begin{equation}
 D_C=(1-C)^{-1}=\sum_{n\ge0}C^n.
\label{eq:v4v29-influence-resolvent}
\end{equation}

\begin{theorem}[Buffered influence amplification]
\label{thm:v4v29-buffer-amplification}
If \(L\ge0\) is an integer and sets \(A,B\) satisfy
\(d(A,B)>LR_0\), then
\begin{equation}
 \sup_{x\in A}\sum_{y\in B}(D_C)_{xy}
 \le
 \frac{\vartheta^{L+1}}{1-\vartheta}.
\label{eq:v4v29-buffer-amplification}
\end{equation}
\end{theorem}

\begin{proof}
A nonzero term \((C^n)_{xy}\) is a sum over paths of \(n\) influence edges,
each of range at most \(R_0\).  If \(x\in A\), \(y\in B\), then
\(n\ge L+1\).  The row-sum bound gives
\[
 \sup_x\sum_y(C^n)_{xy}\le\vartheta^n
\]
by induction.  Sum the geometric tail from \(n=L+1\).
\end{proof}

\begin{corollary}[Core response under a Dobrushin coupling]
\label{cor:v4v29-core-response}
Suppose a coupled conditional-resampling construction has one-step
Lipschitz matrix \(C\).  An exterior perturbation supported beyond a collar
of width \(LR_0\) changes a unit-Lipschitz core observable by at most
\begin{equation}
 C\frac{\vartheta^{L+1}}{1-\vartheta}
\label{eq:v4v29-core-response}
\end{equation}
times the size of the exterior perturbation.
\end{corollary}

\begin{proof}
The standard coupling telescope propagates a discrepancy along successive
single-site resamplings.  A path of \(n\) updates contributes the matrix
element \(C^n\).  Summing all paths from the exterior support to the core
gives \(D_C\); apply
Theorem~\ref{thm:v4v29-buffer-amplification}.
\end{proof}

\begin{firewall}[The theorem assumes the strict inequality it amplifies]
\label{fw:v4v29-no-create-contraction}
Equation \eqref{eq:v4v29-buffer-amplification} is useful only after
\(\vartheta<1\) has been proved.  At the critical value
\(\vartheta=1\), its geometric series diverges and the massless chain
saturates the failure.  Blocking cannot be cited as a proof of the strict
one-step inequality.
\end{firewall}

\subsection{Core, collar, and one-use ancestry}
\label{sec:v4v29-core-collar}

Let a block \(Q\) be decomposed into a core \(Q^\circ\) and collar
\(\partial_LQ\) of width \(LR_0\).  Keep the collar variables in the coarser
parent and integrate only the core conditional law.

\begin{proposition}[Exact core/collar factorization]
\label{prop:v4v29-core-collar}
For every finite-volume Higgs measure,
\begin{equation}
 \int F(\phi)\,d\mu(\phi)
 =
 \int
 \left[
 \int F(\phi_{Q^\circ},\eta_{\partial_LQ},\phi_{\rm ext})
 \,d\mu_{Q^\circ}^{\eta,\phi_{\rm ext}}(\phi_{Q^\circ})
 \right]
 d\mu_{\partial_LQ,\rm ext}(\eta,\phi_{\rm ext}).
\label{eq:v4v29-core-collar-factorization}
\end{equation}
If the one-step influence matrix satisfies
\eqref{eq:v4v29-strict-row-sum}, dependence of a core observable on variables
beyond the collar obeys
\eqref{eq:v4v29-core-response}.  The collar itself receives no such factor.
\end{proposition}

\begin{proof}
Equation \eqref{eq:v4v29-core-collar-factorization} is disintegration of a
finite-dimensional probability measure.  Apply Corollary
\ref{cor:v4v29-core-response} to the inner conditional expectation.
A collar observable has distance zero from the retained boundary, so the
path-length hypothesis gives no decay for it.
\end{proof}

\begin{firewall}[The collar is not a second defect or hopping payment]
\label{fw:v4v29-collar-one-use}
The collar variables are retained once in the next parent measure.  They may
not also be charged as a rare large-field component or as a second copy of
the same hopping bond.  Their role is to create geometric separation between
the integrated core and the next exterior response.
\end{firewall}

\subsection{Consequence for the V28 Higgs expansion}
\label{sec:v4v29-consequence}

\begin{theorem}[Exact scope of blocking for the quartic parent]
\label{thm:v4v29-blocking-scope}
For the V28 quartic Higgs parent:
\begin{enumerate}
\item if the microscopic dimensionless influence matrix has row sum
\(\vartheta<1\), buffered core integration improves exterior response by
\(\vartheta^{L+1}/(1-\vartheta)\) and can make derived long-range polymer
attachments arbitrarily small;
\item if the microscopic criterion is at the critical boundary or fails,
naive finite blocking does not prove a strict contraction; and
\item in the massless marginal kinetic model, the exact chain
\eqref{eq:v4v29-chain-mean} rules out any general theorem claiming that
block size alone creates the missing inequality.
\end{enumerate}
\end{theorem}

\begin{proof}
The first item is Theorem~\ref{thm:v4v29-buffer-amplification} and
Proposition~\ref{prop:v4v29-core-collar}.  The second follows from
Firewall~\ref{fw:v4v29-no-create-contraction}.  The third is Proposition
\ref{prop:v4v29-blocking-no-go}.
\end{proof}

\begin{firewall}[The Standard-Model Higgs trajectory is not classified]
\label{fw:v4v29-Higgs-trajectory}
V29 determines what blocking can prove once the microscopic influence ratio
is known.  It does not establish whether the running Standard-Model Higgs
mass, quartic, and hopping parameters stay in the uniformly convex
\(\vartheta<1\) region at all ultraviolet scales.  The broken phase and
critical surface can reach \(\vartheta=1\), where the present expansion does
not close.
\end{firewall}

\subsection{V29 result ledger and compulsory V30 target}
\label{sec:v4v29-conclusion}

V29 disproves the proposed claim that blocking by itself creates diagonal
dominance.  The massless chain has a Dirichlet gap of order
\(\ell^{-2}\), unit total boundary influence, and a boundary-to-block
\(\ell^2\) norm growing as \(\ell^{1/2}\).  With a genuine positive mass,
only a protected core gains exponential boundary decay.  Abstractly, a
strict one-step influence matrix yields the exact buffered factor
\(\vartheta^{L+1}/(1-\vartheta)\); the collar remains a one-use parent
variable.

V30 is the next compulsory companion audit.  It should then decide between
the two remaining Higgs routes: prove that the ultraviolet running
trajectory lies uniformly inside the convex Dobrushin region, or construct
a correlated quartic kinetic parent on the critical surface.  The second
route requires a mixing estimate for the full gradient-plus-quartic measure,
not a product hopping expansion.

\clearpage
\section{V30: companion audit and exponential response of the correlated
quartic parent}
\label{sec:v4v30}

V29 proves that blocking cannot manufacture diagonal dominance at a
massless critical point.  This compulsory checkpoint selects the upstream
alternative: retain the complete gradient-plus-quartic Higgs measure as one
positive parent before estimating its response.  In the uniformly convex
hard phase, small hopping is unnecessary for two-point locality.

The proof uses the finite-dimensional Helffer--Sjostrand identity.  The
covariance is the inverse of the Witten one-form operator
\(\mathcal L_1=\mathcal L\otimes1+D^2S\) between two actual gradients.
A spatial exponential weight commutes with the scalar configuration-space
generator.  Only the finite-range off-diagonal Hessian is changed by
conjugation, so a Combes--Thomas argument gives an exponentially decaying
one-form resolvent.  This yields a volume-uniform local second background
response for the full correlated positive Higgs parent.

\subsection{Compulsory companion audit}
\label{sec:v4v30-audit}

The exact files inspected were
\begin{enumerate}
\item
\texttt{Renewal\_Yang\_Mills\_Mass\_Gap\_Research\_Notes\_V51.tex},
\(938029\) bytes, SHA--256
\[
 \texttt{85998AD9C785A8DA59A10AD87CBD7EC14FA5B2D472620F753FC6BA086DFEA59D}.
\]
Its newest result computes finite-heat \(SU(2)\) character moments and
exhibits a genuine four-active-row composite cumulant
\(\kappa_3(x_j^2,x_j,x_j)\ge (2d_j^3)^{-1}\) for large \(j\).  Thus an
isolated two-inverse-Casimir estimate fails; any remaining suppression must
come from the complete prefix, payer/suffix carrier, coordinate
cancellation, or a different legal tree.
\item
\texttt{Renewal\_NS\_Stability\_Research\_Notes\_V31.tex},
\(887537\) bytes, SHA--256
\[
 \texttt{F1121B62238AD5491BB7CF03635EA9934942EDF573ED8FAAA9EE79E1D38A6696}.
\]
Its newest result retains the stopping-time-correlated formation occupation
and proves that heat smoothing or a polarization hinge alone cannot supply
the missing total-sector first moment.
\end{enumerate}

The common lesson is structural.  YM V51 shows quantitatively that the
isolated four-row joining letter cannot pay the desired second debit; its
prefix and carrier must remain attached.  NS V31 leaves the complete signed
triad correlation intact because separate positive estimates return the
known critical stock.  For the Higgs sector, splitting the kinetic term into
small product-reference bonds is valid inside the V28 convergence ball, but
it is not the only parent.  Outside that ball one must retain the correlated
convex measure and estimate its complete response operator.

\subsection{The correlated convex Higgs measure}
\label{sec:v4v30-measure}

Let \(\phi=(\phi_x)_{x\in\Lambda}\in(\mathbb R^r)^\Lambda\), and let
\begin{equation}
 d\mu_B(\phi)
 =
 Z(B)^{-1}e^{-S_B(\phi)}\,d\phi.
\label{eq:v4v30-measure}
\end{equation}
Assume:
\begin{enumerate}
\item \(S_B\in C^3\) and
\begin{equation}
 D_\phi^2S_B(\phi)\ge\kappa\,1
\label{eq:v4v30-uniform-convexity}
\end{equation}
for a volume-independent \(\kappa>0\);
\item the off-diagonal blocks of \(D^2S_B\) have range \(R_0\) and
\begin{equation}
 \sup_{x,\phi,B}
 \sum_{y\ne x}
 \left\|
 (D_\phi^2S_B(\phi))_{xy}
 \right\|
 \le J_0;
\label{eq:v4v30-offdiagonal-row}
\end{equation}
and
\item all local observables below and their first field gradients are
square-integrable uniformly in the stated background chart.
\end{enumerate}
For the Higgs action, the quartic Hessian is onsite and positive; the
off-diagonal blocks in
\eqref{eq:v4v30-offdiagonal-row} come from the finite-range covariant
kinetic stencil.

Define the nonnegative scalar generator
\begin{equation}
 \mathcal L
 =
 -\Delta_\phi+\nabla_\phi S_B\cdot\nabla_\phi
\label{eq:v4v30-generator}
\end{equation}
on \(L^2(\mu_B)\), and the Witten one-form operator
\begin{equation}
 \mathcal L_1
 =
 \mathcal L\otimes1+D_\phi^2S_B
\label{eq:v4v30-Witten}
\end{equation}
on
\(L^2(\mu_B;\bigoplus_x\mathbb R^r)\).

\begin{lemma}[One-form gap]
\label{lem:v4v30-one-form-gap}
Under \eqref{eq:v4v30-uniform-convexity},
\begin{equation}
 \mathcal L_1\ge\kappa\,1,
 \qquad
 \|\mathcal L_1^{-1}\|\le\kappa^{-1}.
\label{eq:v4v30-one-form-gap}
\end{equation}
\end{lemma}

\begin{proof}
For a smooth vector field \(u\), integration by parts gives
\[
 \langle u,(\mathcal L\otimes1)u\rangle_{L^2(\mu_B)}
 =
 \sum_i\|\partial_{\phi_i}u\|_{L^2(\mu_B)}^2\ge0.
\]
The Hessian term is at least
\(\kappa\|u\|^2\) by
\eqref{eq:v4v30-uniform-convexity}.  Add the two forms and apply the spectral
theorem.
\end{proof}

\subsection{The exact Helffer--Sjostrand covariance}
\label{sec:v4v30-HS}

\begin{theorem}[Helffer--Sjostrand identity]
\label{thm:v4v30-HS}
For smooth square-integrable \(F,G\),
\begin{equation}
 \operatorname{Cov}_{\mu_B}(F,G)
 =
 \left\langle
 \nabla_\phi F,
 \mathcal L_1^{-1}\nabla_\phi G
 \right\rangle_{L^2(\mu_B)}.
\label{eq:v4v30-HS}
\end{equation}
\end{theorem}

\begin{proof}
Replace \(G\) by \(G-\mathbb EG\).  Uniform convexity gives a Poincare
inequality, so \(\mathcal L\) is invertible on centered functions.  Let
\(u=\mathcal L^{-1}G\).  Integration by parts gives
\[
 \operatorname{Cov}(F,G)
 =
 \langle F,\mathcal Lu\rangle
 =
 \langle\nabla F,\nabla u\rangle.
\]
Differentiating \(\mathcal Lu=G\) yields the exact commutation identity
\[
 \nabla(\mathcal Lu)
 =
 (\mathcal L\otimes1+D^2S_B)\nabla u
 =
 \mathcal L_1\nabla u
 =
 \nabla G.
\]
Lemma~\ref{lem:v4v30-one-form-gap} gives
\(\nabla u=\mathcal L_1^{-1}\nabla G\), proving
\eqref{eq:v4v30-HS}.  Approximation extends the identity to the stated
square-integrable class.
\end{proof}

\begin{firewall}[The covariance is already the assembled carrier]
\label{fw:v4v30-HS-carrier}
Equation \eqref{eq:v4v30-HS} contains the complete correlated Hessian inside
one inverse.  Replacing it by a product-reference bond series and then also
retaining \(\mathcal L_1^{-1}\) would count the same hopping response twice.
V28 and V30 are alternative representations in their common domain.
\end{firewall}

\subsection{A spatial Combes--Thomas estimate}
\label{sec:v4v30-CT}

For \(A\subset\Lambda\), put
\[
 \psi_A(x)=d(x,A).
\]
Let \(W_\eta\) multiply the \(x\)-component of a one-form by
\(e^{-\eta\psi_A(x)}\).  It acts only on the spatial component index and
therefore commutes with \(\mathcal L\otimes1\).

\begin{lemma}[Weighted Hessian commutator]
\label{lem:v4v30-weighted-Hessian}
There is \(C=C(R_0)\) such that, for \(0<\eta\le1\),
\begin{equation}
 \left\|
 W_\eta D^2S_B W_\eta^{-1}-D^2S_B
 \right\|
 \le
 C J_0\eta e^{\eta R_0}.
\label{eq:v4v30-weighted-Hessian}
\end{equation}
\end{lemma}

\begin{proof}
The \(xy\) block is multiplied by
\(e^{-\eta(\psi_A(x)-\psi_A(y))}-1\).  It vanishes for \(x=y\); for a
nonzero off-diagonal block, \(d(x,y)\le R_0\), so
\[
 \left|
 e^{\eta(\psi_A(x)-\psi_A(y))}-1
 \right|
 \le
 \eta R_0e^{\eta R_0}.
\]
Use \eqref{eq:v4v30-offdiagonal-row} and the symmetric column bound in the
Schur test.
\end{proof}

Choose \(\eta_*>0\) so
\begin{equation}
 C J_0\eta_*e^{\eta_*R_0}\le\kappa/2.
\label{eq:v4v30-eta-choice}
\end{equation}

\begin{theorem}[Exponential one-form resolvent]
\label{thm:v4v30-Witten-decay}
For subsets \(A,C\subset\Lambda\),
\begin{equation}
 \left\|
 P_A\mathcal L_1^{-1}P_C
 \right\|
 \le
 \frac2\kappa
 e^{-\eta_*d(A,C)}.
\label{eq:v4v30-Witten-decay}
\end{equation}
The projections act on spatial one-form components; the norm is on
\(L^2(\mu_B;\bigoplus_x\mathbb R^r)\).
\end{theorem}

\begin{proof}
By Lemma~\ref{lem:v4v30-weighted-Hessian},
\[
 W_{\eta_*}\mathcal L_1W_{\eta_*}^{-1}
 =
 \mathcal L_1+E_{\eta_*},
 \qquad
 \|E_{\eta_*}\|\le\kappa/2.
\]
Lemma~\ref{lem:v4v30-one-form-gap} and a Neumann series give
\[
 \left\|
 (W_{\eta_*}\mathcal L_1W_{\eta_*}^{-1})^{-1}
 \right\|
 \le2/\kappa.
\]
The weight is one on \(A\) and at most
\(e^{-\eta_*d(A,C)}\) on \(C\).  Conjugating the inverse back and inserting
the two spatial projections proves
\eqref{eq:v4v30-Witten-decay}.
\end{proof}

\begin{corollary}[Exponential covariance of local observables]
\label{cor:v4v30-covariance-decay}
If \(\nabla F\) is supported in \(A\) and \(\nabla G\) in \(C\), then
\begin{equation}
 |\operatorname{Cov}_{\mu_B}(F,G)|
 \le
 \frac2\kappa e^{-\eta_*d(A,C)}
 \|\nabla F\|_{L^2(\mu_B)}
 \|\nabla G\|_{L^2(\mu_B)}.
\label{eq:v4v30-covariance-decay}
\end{equation}
\end{corollary}

\begin{proof}
Insert the support projections in
\eqref{eq:v4v30-HS}, apply
\eqref{eq:v4v30-Witten-decay}, and use Cauchy--Schwarz.
\end{proof}

\subsection{A genuine local second background response}
\label{sec:v4v30-response}

Let \(b=(b_x)_{x\in\Lambda}\) be local background coordinates.  Put
\begin{equation}
 X_x=D_{b_x}S_B,
 \qquad
 Y_{xy}=D_{b_x}D_{b_y}S_B.
\label{eq:v4v30-local-scores}
\end{equation}
Assume \(\nabla_\phi X_x\) is supported in a fixed-radius neighborhood
\(A_x\), and
\begin{equation}
 |\mathbb E_{\mu_B}X_x|\le M_0,
 \qquad
 \|\nabla_\phi X_x\|_{L^2(\mu_B)}\le M_1,
 \qquad
 |\mathbb E_{\mu_B}Y_{xy}|
 \le
 M_2\boldsymbol1_{\{d(x,y)\le R_1\}}.
\label{eq:v4v30-score-locality}
\end{equation}

\begin{theorem}[Exponential two-mark response kernel]
\label{thm:v4v30-two-mark-response}
The exact first derivative and Hessian of the correlated Higgs free energy
are
\begin{equation}
 D_{b_x}[-\log Z(B)]=\mathbb E_{\mu_B}X_x,
\label{eq:v4v30-exact-first-response}
\end{equation}
and
\begin{equation}
 D_{b_x}D_{b_y}[-\log Z(B)]
 =
 \mathbb E_{\mu_B}Y_{xy}
 -
 \operatorname{Cov}_{\mu_B}(X_x,X_y).
\label{eq:v4v30-exact-response}
\end{equation}
There are \(C,\eta'>0\), independent of volume, such that
\begin{equation}
 \left|
 D_{b_x}D_{b_y}[-\log Z(B)]
 \right|
 \le
 C(M_2+\kappa^{-1}M_1^2)
 e^{-\eta'd(x,y)}.
\label{eq:v4v30-response-decay}
\end{equation}
\end{theorem}

\begin{proof}
Differentiate the normalized finite-dimensional integral once and twice, as in
Theorem~\ref{thm:v4v27-response-identities}; this proves
\eqref{eq:v4v30-exact-response}.  The contact term has range \(R_1\).
For the covariance, apply Corollary
\ref{cor:v4v30-covariance-decay} to \(A_x,A_y\).  Their fixed radii change
the distance by a constant, which is absorbed into \(C\).  Reduce
\(\eta_*\) if necessary to dominate the finite-range contact term.
\end{proof}

\begin{firewall}[The two marks are actual source occurrences]
\label{fw:v4v30-actual-two-marks}
The covariance in \eqref{eq:v4v30-exact-response} contains the two actual
score occurrences \(X_x,X_y\).  The contact term is the genuine same-action
second derivative \(Y_{xy}\).  These are the two exhaustive placements in
one differentiated logarithm; no square of a first-response upper bound is
used.
\end{firewall}

\begin{corollary}[Local and intensive correlated-parent response]
\label{cor:v4v30-local-intensive}
For a background perturbation supported in a fixed set, the first response
is volume-uniform and the second-response kernel is localized
exponentially.  For a
translation-covariant extensive perturbation, the second response divided
by \(|\Lambda|\) is uniformly bounded.
\end{corollary}

\begin{proof}
Use \eqref{eq:v4v30-exact-first-response} and the \(M_0\) bound, then sum
\eqref{eq:v4v30-response-decay}.  A finite-dimensional lattice has
uniformly summable exponential balls.  A local source has finitely many
roots; an extensive source has \(|\Lambda|\) roots, and division by volume
removes that factor.
\end{proof}

\subsection{Comparison with the V28--V29 route}
\label{sec:v4v30-comparison}

\begin{theorem}[Two valid Higgs hard-phase routes]
\label{thm:v4v30-two-routes}
In the uniformly convex hard phase:
\begin{enumerate}
\item if the dimensionless hopping satisfies the V28 small-bond criterion,
the product-reference expansion gives a full two-mark KP gas and may be
combined with complex activity species through V15;
\item without small hopping, assumptions
\eqref{eq:v4v30-uniform-convexity}--%
\eqref{eq:v4v30-offdiagonal-row} give uniform local first responses and an
exponential second-response kernel for the complete positive correlated
Higgs parent.
\end{enumerate}
The second route proves the response needed for the positive Higgs modulus
but does not by itself produce an all-orders polymer expansion for attached
complex Yukawa activities.
\end{theorem}

\begin{proof}
The first item is Theorem~\ref{thm:v4v28-Higgs-KP}.  The second is
Theorem~\ref{thm:v4v30-two-mark-response} and Corollary
\ref{cor:v4v30-local-intensive}.  Their scopes differ exactly as stated.
\end{proof}

\begin{firewall}[Uniform convexity remains the phase boundary]
\label{fw:v4v30-phase-boundary}
The Witten gap is inherited from
\eqref{eq:v4v30-uniform-convexity}.  At a broken or critical Higgs point
where the hard Hessian develops an untransferred zero mode, neither
\eqref{eq:v4v30-one-form-gap} nor the exponential decay theorem follows.
The vacuum manifold and Goldstone/gauge directions must be transferred
before this route can resume.
\end{firewall}

\begin{firewall}[Yukawa and chiral phase remain outside the positive parent]
\label{fw:v4v30-Yukawa-open}
The Helffer--Sjostrand identity applies to the positive Higgs probability
measure.  A chiral Yukawa determinant can be complex and can change the
measure's sign or phase.  It must remain an attached anomaly-normalized
activity.  V30 does not prove its marked norm over the correlated quartic
parent.
\end{firewall}

\subsection{V30 result ledger and V31 target}
\label{sec:v4v30-conclusion}

V30 records the compulsory companion audit and implements its shared
assembly lesson.  The complete correlated gradient-plus-quartic Higgs law
is retained before estimation.  Uniform convexity gives a gap for its Witten
one-form operator.  A spatial Combes--Thomas conjugation proves exponential
decay of the exact Helffer--Sjostrand covariance, and hence of the genuine
two-mark background-response kernel, without a small-hopping assumption.

The next target is the Yukawa attachment.  V31 should condition the chiral
fermion determinant on the correlated Higgs parent, use the V16 resolvent
word to localize one and two actual Higgs/Yukawa marks, and determine whether
their exponential moments are controlled by the quartic tail.  The
anomaly-normalized phase must remain separate from the positive carrier, and
no determinant modulus may be inserted twice.

\clearpage
\section{V31: gap-free local Yukawa attachment before fermion integration}
\label{sec:v4v31}

V30 controls the exact two-point response of the full correlated positive
Higgs parent, but it deliberately leaves the complex Yukawa determinant
outside that probability measure.  Applying the V16 trace-log formula
directly to a Higgs-dependent fermion matrix would require a uniform inverse.
That requirement fails at a Yukawa zero and is stronger than local
integrability.

This note goes upstream of the determinant.  At finite cutoff the Yukawa
vertex is an even element of a finite Grassmann algebra, so its exponential
is a finite polynomial.  We prove a volume-uniform twice-marked bound for
that polynomial in \(L^2\) of the correlated quartic Higgs parent.  No
fermion inverse and no choice of determinant phase is used.  The remaining
step is a connected fermionic contraction theorem that combines these local
coefficients with the hard fermion covariance exactly once.

\subsection{Local exponential moments of the correlated convex parent}
\label{sec:v4v31-Higgs-moments}

Retain the measure \(\mu_B\) and the uniform convexity constant \(\kappa\)
from V30.  For a finite set \(X\subset\Lambda\), write
\[
 \phi_X=(\phi_x)_{x\in X}\in\mathbb R^{r|X|}.
\]
Assume throughout this note that
\begin{equation}
 \sup_{B,x}\left|\mathbb E_{\mu_B}\phi_x\right|\le m_*.
\label{eq:v4v31-mean-bound}
\end{equation}
In the centered unbroken Higgs chart the action is even in \(\phi\), and
then \(m_*=0\).  The displayed form also permits a bounded recentered chart.

\begin{lemma}[Correlated-parent subgaussian estimate]
\label{lem:v4v31-subgaussian}
If \(D_\phi^2S_B\ge\kappa 1\), then for every finite \(X\), every
\(a\in\mathbb R^{r|X|}\), and every \(t\in\mathbb R\),
\begin{equation}
 \mathbb E_{\mu_B}
 \exp\!\left(
 t\,a\mathbin{\cdot}
 [\phi_X-\mathbb E_{\mu_B}\phi_X]
 \right)
 \le
 \exp\!\left(\frac{t^2|a|^2}{2\kappa}\right).
\label{eq:v4v31-linear-subgaussian}
\end{equation}
Consequently, if \(n_X=r|X|\) and \(0<\tau<\kappa/4\), then
\begin{equation}
 \mathbb E_{\mu_B}e^{\tau|\phi_X|^2}
 \le
 e^{2\tau |X|m_*^2}
 \left(1-\frac{4\tau}{\kappa}\right)^{-n_X/2}.
\label{eq:v4v31-quadratic-exponential}
\end{equation}
Both estimates are independent of the total volume.
\end{lemma}

\begin{proof}
The Bakry--Emery criterion for the generator
\(\mathcal L\) in \eqref{eq:v4v30-generator} gives the logarithmic Sobolev
inequality
\begin{equation}
 \operatorname{Ent}_{\mu_B}(f^2)
 \le\frac2\kappa
 \mathbb E_{\mu_B}|\nabla_\phi f|^2.
\label{eq:v4v31-LSI}
\end{equation}
For completeness, apply \eqref{eq:v4v31-LSI} to
\(f=e^{tF/2}\), where
\(F=a\mathbin{\cdot}(\phi_X-\mathbb E\phi_X)\).
If \(H(t)=\log\mathbb E e^{tF}\), division by
\(\mathbb E e^{tF}\) gives
\[
 tH'(t)-H(t)\le\frac{t^2|a|^2}{2\kappa}.
\]
Since \(H(0)=H'(0)=0\), integration of
\((H(t)/t)'\le |a|^2/(2\kappa)\), with the same argument for
negative \(t\), proves \eqref{eq:v4v31-linear-subgaussian}.

Put \(Z=\phi_X-\mathbb E\phi_X\), and let
\(G\) be a standard Gaussian vector in \(\mathbb R^{n_X}\), independent of
\(Z\).  The Gaussian identity and Tonelli's theorem give
\[
 \mathbb E_Ze^{2\tau|Z|^2}
 =
 \mathbb E_G\mathbb E_Z
 e^{2\sqrt{\tau}\,G\cdot Z}
 \le
 \mathbb E_Ge^{2\tau|G|^2/\kappa}
 =
 \left(1-\frac{4\tau}{\kappa}\right)^{-n_X/2}.
\]
Finally,
\[
 |\phi_X|^2\le2|Z|^2+
 2|\mathbb E\phi_X|^2
 \le2|Z|^2+2|X|m_*^2.
\]
Multiplication by the deterministic exponential proves
\eqref{eq:v4v31-quadratic-exponential}.
\end{proof}

\begin{corollary}[Uniform fixed-cell polynomial moments]
\label{cor:v4v31-polynomial-moments}
For every fixed cell-size bound \(b\) and every \(p<\infty\),
\begin{equation}
 \sup_{\substack{B,\Lambda\\X\subset\Lambda,\ |X|\le b}}
 \left\|1+|\phi_X|\right\|_{L^p(\mu_B)}
 <\infty.
\label{eq:v4v31-polynomial-moments}
\end{equation}
Moreover, for every \(c<\infty\) there is
\(\epsilon_0=\epsilon_0(c,b,r,\kappa)>0\) such that
\[
 \sup_{\substack{0\le\epsilon\le\epsilon_0\\
 B,\Lambda,\ |X|\le b}}
 \mathbb E_{\mu_B}
 \exp\!\left(4c\epsilon[1+|\phi_X|]\right)
 <\infty.
\]
\end{corollary}

\begin{proof}
The first assertion follows by integrating the tail implied by
\eqref{eq:v4v31-quadratic-exponential}.  For the second, choose
\(0<\tau<\kappa/4\) and use
\[
 4c\epsilon|\phi_X|
 \le\tau|\phi_X|^2+\frac{4c^2\epsilon^2}{\tau}.
\]
Equation \eqref{eq:v4v31-quadratic-exponential} gives the uniform bound.
\end{proof}

\subsection{Finite Grassmann coefficient norm}
\label{sec:v4v31-Grassmann}

Let \(\mathfrak G_q\) be the complex Grassmann algebra on a fixed ordered
list \(\xi_1,\ldots,\xi_q\).  Every \(A\in\mathfrak G_q\) has a unique
expansion
\[
 A=\sum_{I\subset\{1,\ldots,q\}}a_I\xi_I,
 \qquad
 \|A\|_{\mathfrak G}:=\sum_I|a_I|,
\]
where the monomial \(\xi_I\) is written in increasing order.

\begin{lemma}[Submultiplicativity and finite exponential]
\label{lem:v4v31-Grassmann-norm}
For \(A,C\in\mathfrak G_q\),
\begin{equation}
 \|AC\|_{\mathfrak G}
 \le\|A\|_{\mathfrak G}\|C\|_{\mathfrak G}.
\label{eq:v4v31-Grassmann-submultiplicative}
\end{equation}
If \(A\) has vanishing scalar part, then \(A^{q+1}=0\), and
\begin{align}
 \|e^A\|_{\mathfrak G}&\le e^{\|A\|_{\mathfrak G}},
\label{eq:v4v31-Grassmann-exponential}\\
 \|e^A-1\|_{\mathfrak G}
 &\le\|A\|_{\mathfrak G}e^{\|A\|_{\mathfrak G}}.
\label{eq:v4v31-Grassmann-activity}
\end{align}
\end{lemma}

\begin{proof}
On multiplying two ordered monomials, an intersecting pair gives zero and a
disjoint pair gives, up to sign, their union.  The triangle inequality after
collecting equal unions proves
\eqref{eq:v4v31-Grassmann-submultiplicative}.  Every nonzero product of more
than \(q\) positive-degree monomials repeats a generator, so the exponential
terminates.  Apply submultiplicativity term by term.  For the final bound use
\(\sum_{n\ge1}u^n/n!\le ue^u\).
\end{proof}

\subsection{The local Yukawa packet and its two actual marks}
\label{sec:v4v31-Yukawa-packet}

Let \(X\) be contained in a cell of at most \(b\) sites.  Write the complete
local Yukawa vertex, including its Hermitian-conjugate term when present, as
\begin{equation}
 \mathcal V_X(B,\phi,\xi)
 =
 \sum_{\alpha=1}^{N_Y}
 y_\alpha c_{\alpha,X}(B,\phi_X)\Xi_\alpha.
\label{eq:v4v31-Yukawa-vertex}
\end{equation}
Here each \(\Xi_\alpha\) is an even positive-degree Grassmann monomial and
the complex coefficients \(y_\alpha\) include the finite internal
multiplicities.  Put
\[
 \epsilon_Y:=\sum_{\alpha=1}^{N_Y}|y_\alpha|.
\]
For \(j=0,1,2\), assume that the background derivatives act only on the
coefficient in \eqref{eq:v4v31-Yukawa-vertex} and obey
\begin{equation}
 \sup_{\|v_i\|\le1}
 \left|
 D_B^j c_{\alpha,X}[v_1,\ldots,v_j]
 \right|
 \le L_j(1+|\phi_X|).
\label{eq:v4v31-coefficient-bound}
\end{equation}
This is the natural finite-cell bound for a Yukawa coefficient affine in
the Higgs field and smooth in a bounded gauge/metric chart.

Define the unintegrated local activity
\begin{equation}
 w_X=e^{-\mathcal V_X}-1\in\mathfrak G_q.
\label{eq:v4v31-unintegrated-activity}
\end{equation}
Since \(\mathcal V_X\) is even, it commutes with its background derivatives
in the supercommutative algebra.  Therefore
\begin{align}
 D_Bw_X[v]
 &=-e^{-\mathcal V_X}D_B\mathcal V_X[v],
\label{eq:v4v31-first-mark}\\
 D_B^2w_X[v_1,v_2]
 &=e^{-\mathcal V_X}
 \left(
 D_B\mathcal V_X[v_1]D_B\mathcal V_X[v_2]
 -D_B^2\mathcal V_X[v_1,v_2]
 \right).
\label{eq:v4v31-second-mark}
\end{align}

\begin{theorem}[Gap-free twice-marked local Yukawa bound]
\label{thm:v4v31-Yukawa-L2}
Under \eqref{eq:v4v31-mean-bound} and
\eqref{eq:v4v31-coefficient-bound}, there are
\(\epsilon_0>0\) and constants \(C_0,C_1,C_2<\infty\), independent of
volume and of the cell location, such that, for
\(0\le\epsilon_Y\le\epsilon_0\),
\begin{equation}
 \sup_{\|v_i\|\le1}
 \left\|
 D_B^j w_X[v_1,\ldots,v_j]
 \right\|_{L^2(\mu_B;\mathfrak G_q)}
 \le C_j\epsilon_Y,
 \qquad j=0,1,2.
\label{eq:v4v31-Yukawa-L2}
\end{equation}
No lower bound on a Higgs-dependent fermion determinant or singular value is
required.
\end{theorem}

\begin{proof}
By \eqref{eq:v4v31-coefficient-bound} and the definition of
\(\epsilon_Y\),
\begin{equation}
 \sup_{\|v_i\|\le1}
 \|D_B^j\mathcal V_X[v_1,\ldots,v_j]\|_{\mathfrak G}
 \le\epsilon_YL_j(1+|\phi_X|),
 \qquad j=0,1,2.
\label{eq:v4v31-vertex-mark-bound}
\end{equation}
Lemma~\ref{lem:v4v31-Grassmann-norm},
\eqref{eq:v4v31-first-mark}, and
\eqref{eq:v4v31-second-mark} give the pointwise estimates
\begin{align*}
 \|w_X\|_{\mathfrak G}
 &\le
 \epsilon_YL_0(1+|\phi_X|)
 e^{\epsilon_YL_0(1+|\phi_X|)},\\
 \|D_Bw_X[v]\|_{\mathfrak G}
 &\le
 \epsilon_YL_1(1+|\phi_X|)
 e^{\epsilon_YL_0(1+|\phi_X|)},\\
 \|D_B^2w_X[v_1,v_2]\|_{\mathfrak G}
 &\le
 \left[
 \epsilon_Y^2L_1^2(1+|\phi_X|)^2+
 \epsilon_YL_2(1+|\phi_X|)
 \right]
 e^{\epsilon_YL_0(1+|\phi_X|)}.
\end{align*}
Square these inequalities.  Corollary
\ref{cor:v4v31-polynomial-moments}, followed by Cauchy--Schwarz, bounds every
resulting polynomial times exponential uniformly for
\(\epsilon_Y\le\epsilon_0\).  The first two lines are
\(O(\epsilon_Y)\) in \(L^2\).  The last is
\(O(\epsilon_Y+\epsilon_Y^2)\), which is
\(O(\epsilon_Y)\) after reducing \(\epsilon_0\) to at most one.  This proves
\eqref{eq:v4v31-Yukawa-L2}.
\end{proof}

\begin{firewall}[The quadratic mark is not a squared estimate]
\label{fw:v4v31-two-marks}
The first term in \eqref{eq:v4v31-second-mark} contains two actual
first-derivative occurrences generated by differentiating the same finite
Grassmann exponential.  The second term contains the actual second
derivative of the Yukawa coefficient.  Both are retained in
\eqref{eq:v4v31-Yukawa-L2}; no first-mark upper bound replaces the second
derivative.
\end{firewall}

\begin{firewall}[No inverse determinant and no phase claim]
\label{fw:v4v31-no-inverse}
Theorem~\ref{thm:v4v31-Yukawa-L2} is proved before Berezin integration.
It remains valid when a later Higgs-dependent fermion determinant vanishes,
because it never divides by that determinant.  It neither chooses nor
controls the anomaly-normalized chiral phase.
\end{firewall}

\subsection{Relation to the existing determinant ancestry}
\label{sec:v4v31-ancestry}

\begin{theorem}[Exact placement of the new packet]
\label{thm:v4v31-placement}
The V31 activity is disjoint from the V16 hard trace-log packet provided the
fermion quadratic covariance is placed in the Gaussian parent and
\(\mathcal V_X\) contains only Yukawa vertices left outside that quadratic
form.  Under this convention, subsequent Berezin expectation contracts the
V31 Grassmann generators with the V16 matter covariance once, while the
V16 determinant normalization remains in its existing Schur ancestry.
\end{theorem}

\begin{proof}
A finite-dimensional fermion action splits algebraically as
\[
 \bar\psi A_0\psi+\sum_X\mathcal V_X.
\]
The normalized Gaussian Berezin identity is
\[
 \frac{\int e^{-\bar\psi A_0\psi}F(\psi,\bar\psi)
       \,d\bar\psi\,d\psi}
      {\int e^{-\bar\psi A_0\psi}\,d\bar\psi\,d\psi}
 =
 \mathbb E_{C_0}^{\rm F}F,
 \qquad C_0=A_0^{-1}.
\]
The denominator is exactly the determinant of the quadratic parent; the
numerator contains each factor \(e^{-\mathcal V_X}=1+w_X\) once.  Thus
contracting the \(w_X\)'s cannot create a second quadratic determinant.
Conversely, moving any part of \(\mathcal V_X\) into \(A_0\) changes both
objects and requires redoing the split rather than multiplying the old
packets.  This proves the stated ancestry.
\end{proof}

\begin{firewall}[Local \(L^2\) control is not a connected KP theorem]
\label{fw:v4v31-not-KP}
Equation \eqref{eq:v4v31-Yukawa-L2} controls each fixed-cell coefficient
through two marks.  It does not yet sum connected fermionic contractions
over arbitrarily many cells.  That step needs the antisymmetric Wick
determinants, the decay of \(C_0\), and the correlated Higgs coefficients in
one common connected expansion.
\end{firewall}

\subsection{V31 result ledger and V32 target}
\label{sec:v4v31-conclusion}

V31 proves a correlated-parent subgaussian estimate and volume-uniform
fixed-cell exponential moments.  It introduces the finite Grassmann
coefficient norm, retains the unintegrated Yukawa exponential as an exact
finite polynomial, and proves its local \(L^2\) bounds through two genuine
background marks.  The construction is gap-free at Yukawa determinant
zeros and fits the existing quadratic determinant ancestry without
duplication.

V32 should perform the fermionic contraction.  The first target is a Gram
bound for the antisymmetric Wick minors with two marked Yukawa cells,
followed by an exact connected-support extraction.  The proof must identify
whether the resulting connected sum closes under a small
coupling-normalized walk norm.  If it does not, the complete marginal
Higgs--Yukawa shell law, rather than isolated Yukawa cells, must be promoted
to the parent.


\clearpage
\section{V32: connected fermionic contraction of the local Yukawa packet}
\label{sec:v4v32}

V31 gives volume-uniform Higgs \(L^2\) bounds for one unintegrated Yukawa
cell, but a continuum construction needs a sum over arbitrarily many cells.
This note performs the first missing global step.  A fermionic forest
interpolation extracts only contractions whose cell graph is connected.
Gram's inequality controls every unextracted Wick minor without a
factorial, while decay of the hard fermion covariance pays a spanning tree.
The joint exponential moment from V31 controls all correlated Higgs
coefficients at once.

The result is a marked, absolutely summable expansion of the conditional
fermionic effective action for sufficiently small normalized Yukawa
coupling.  The background directions treated here act on the Yukawa
coefficients.  Marks of the hard covariance and Higgs connections between
distinct fermion-connected components remain for the next step.

\subsection{Fermionic Gaussian expectation and Gram data}
\label{sec:v4v32-Gram-data}

Let cells be indexed by a block lattice \(\mathbb L\).  Cell \(x\) contains
at most \(q\) fermion generators.  Let
\(\mathbb E_C^{\rm F}\) be the normalized fermionic Gaussian expectation
with covariance \(C\).  Thus, up to the sign fixed by the chosen ordering,
\begin{equation}
 \mathbb E_C^{\rm F}
 \left[
 \bar\psi_{i_1}\psi_{j_1}\cdots
 \bar\psi_{i_m}\psi_{j_m}
 \right]
 =
 \det\!\left(C_{j_a i_b}\right)_{a,b=1}^m.
\label{eq:v4v32-Wick-determinant}
\end{equation}

Assume that \(C\) has a Gram representation
\begin{equation}
 C_{ji}=\langle f_j,g_i\rangle_{\mathcal H},
 \qquad
 \|f_j\|,\|g_i\|\le\gamma,
\label{eq:v4v32-Gram-representation}
\end{equation}
and a cell majorant
\begin{equation}
 |C_{ji}|
 \le\gamma^2 c(x(j),x(i)),
 \qquad
 \chi_a:=
 \sup_x\sum_y e^{a d(x,y)}c(x,y)<\infty
\label{eq:v4v32-covariance-majorant}
\end{equation}
for some \(a>0\).  We take \(c(x,x)\ge1\).  The V16 hard-gap
Combes--Thomas estimate supplies such an exponential majorant for the hard
matter covariance.

\begin{lemma}[Gram bound for every Wick minor]
\label{lem:v4v32-Gram}
Under \eqref{eq:v4v32-Gram-representation},
\begin{equation}
 \left|
 \det(C_{j_a i_b})_{a,b=1}^m
 \right|
 \le\gamma^{2m}.
\label{eq:v4v32-Gram-bound}
\end{equation}
The same estimate holds after replacing \(C\) by
\(C^s_{ji}=s_{x(j)x(i)}C_{ji}\), whenever
\(s_{xy}=\langle u_x,u_y\rangle\) is positive semidefinite and
\(\|u_x\|\le1\).
\end{lemma}

\begin{proof}
The first assertion is Gram's determinant inequality applied to the two
families \(f_{j_a}\) and \(g_{i_b}\).  For the second, use the Gram
representation
\[
 C^s_{ji}
 =
 \langle f_j\otimes u_{x(j)},g_i\otimes u_{x(i)}\rangle_{
 \mathcal H\otimes\mathcal U}.
\]
The tensor-product vectors have norms at most \(\gamma\), so the same
inequality applies.
\end{proof}

\subsection{Exact connected support}
\label{sec:v4v32-connected-support}

For even Grassmann polynomials \(P_1,\ldots,P_n\), supported in distinct
cells \(x_1,\ldots,x_n\), define their truncated fermionic expectation by
\begin{equation}
 \left\langle P_1;\ldots;P_n\right\rangle_C^{\rm T}
 =
 \sum_{\pi\in\Pi_n}
 (|\pi|-1)!(-1)^{|\pi|-1}
 \prod_{B\in\pi}
 \mathbb E_C^{\rm F}\prod_{i\in B}P_i,
\label{eq:v4v32-truncated}
\end{equation}
where \(\Pi_n\) is the set of partitions of
\(\{1,\ldots,n\}\).  This is exactly the coefficient of
\(t_1\cdots t_n\) in
\[
 \log\mathbb E_C^{\rm F}
 \prod_{i=1}^n(1+t_iP_i).
\]

\begin{lemma}[Fermionic cell-tree bound]
\label{lem:v4v32-cell-tree}
There is a constant \(K_{\rm F}=K_{\rm F}(q,\gamma)\) such that
\begin{equation}
 \left|
 \left\langle P_1;\ldots;P_n\right\rangle_C^{\rm T}
 \right|
 \le
 K_{\rm F}^{\,n}
 \prod_{i=1}^n\|P_i\|_{\mathfrak G}
 \sum_{T\in\mathcal T_n}
 \prod_{\{i,j\}\in E(T)}c(x_i,x_j).
\label{eq:v4v32-cell-tree}
\end{equation}
For \(n=1\), the tree sum is interpreted as one.
\end{lemma}

\begin{proof}
Introduce independent interpolation parameters on the off-diagonal cell
blocks of \(C\).  Apply the Brydges--Kennedy forest formula to the
multilinear coefficient defining
\eqref{eq:v4v32-truncated}.  Its connected part contains exactly the trees
on \(\{1,\ldots,n\}\).  For a fixed tree \(T\), differentiation of one
covariance block for each edge \(\{i,j\}\) uses fermionic Gaussian
integration by parts.  It removes one generator at each endpoint and
produces one cross-cell covariance entry.  There are at most \(q^2\)
choices for that entry, and
\eqref{eq:v4v32-covariance-majorant} bounds it by
\(\gamma^2c(x_i,x_j)\).

The remaining Gaussian expectation is a Wick determinant for the
interpolated covariance.  The forest interpolation matrix is a correlation
matrix: it is positive semidefinite, has diagonal one, and admits vectors
\(u_x\) of norm at most one.  Lemma~\ref{lem:v4v32-Gram} therefore bounds
every remaining minor by a power of \(\gamma\).  At most \(qn\) generators
were present initially, so all powers of \(\gamma\), all endpoint choices,
and the coefficient-norm expansion of the \(P_i\)'s are bounded by
\(K_{\rm F}^n\).  Multiplication over the tree edges and summation over
\(T\) proves \eqref{eq:v4v32-cell-tree}.

The same forest formula also shows the cancellation directly: a Wick term
whose cross-cell graph is disconnected is independent of at least one
interpolation edge needed to join its components, and hence contributes
zero to the connected coefficient.
\end{proof}

\begin{firewall}[Why Gram's inequality is needed]
\label{fw:v4v32-no-Wick-factorial}
Expanding the determinant in \eqref{eq:v4v32-Wick-determinant} into
permutations and taking absolute values would introduce a factorial at every
order.  Lemma~\ref{lem:v4v32-Gram} retains the antisymmetric determinant
until after estimation.  Only the tree endpoint choices remain, and those
are exponential in the number of cells.
\end{firewall}

\subsection{A joint bound for correlated Higgs coefficients}
\label{sec:v4v32-joint-Higgs}

Let \(X_i\) be the Higgs support of the Yukawa cell at \(x_i\).  Assume
\begin{equation}
 |X_i|\le b,
 \qquad
 \sup_{z\in\Lambda}
 \#\{i:z\in X_i\}\le\nu.
\label{eq:v4v32-bounded-overlap}
\end{equation}
For \(m=0,1,2\), set
\[
 A_{m,i}(\phi)
 =
 \sup_{\|v_k\|\le1}
 \|D_B^m w_{X_i}[v_1,\ldots,v_m]\|_{\mathfrak G}.
\]

\begin{lemma}[All-cell coefficient product]
\label{lem:v4v32-all-cell-product}
Under the V31 hypotheses, there are
\(\epsilon_0>0\) and \(K_{\rm H}<\infty\), independent of \(n\), the
volume, and the cell locations, such that whenever
\(\epsilon_Y\le\epsilon_0\) and
\(m_i\in\{0,1,2\}\),
\begin{equation}
 \mathbb E_{\mu_B}
 \prod_{i=1}^n A_{m_i,i}
 \le
 (K_{\rm H}\epsilon_Y)^n.
\label{eq:v4v32-all-cell-product}
\end{equation}
\end{lemma}

\begin{proof}
The pointwise estimates in the proof of
Theorem~\ref{thm:v4v31-Yukawa-L2} imply, after increasing a constant
\(K_0\),
\[
 A_{m,i}
 \le
 K_0\epsilon_Y
 (1+|\phi_{X_i}|)^2
 e^{K_0\epsilon_Y(1+|\phi_{X_i}|)}
\]
for all \(m=0,1,2\) and \(\epsilon_Y\le1\).  Fix
\(0<\delta<\kappa/(4\nu)\).  The elementary supremum
\[
 K_\delta:=
 \sup_{u\ge0}
 (1+u)^2e^{K_0\epsilon_0(1+u)-\delta u^2}
 <\infty
\]
is finite.  Hence
\[
 \prod_{i=1}^n A_{m_i,i}
 \le
 (K_0K_\delta\epsilon_Y)^n
 \exp\!\left(\delta\sum_{i=1}^n|\phi_{X_i}|^2\right).
\]
Let \(U=\bigcup_iX_i\).  Bounded overlap gives
\[
 \sum_i|\phi_{X_i}|^2\le\nu|\phi_U|^2,
 \qquad |U|\le bn.
\]
Apply \eqref{eq:v4v31-quadratic-exponential} with
\(\tau=\delta\nu<\kappa/4\).  Its right side is a fixed constant to the
power \(|U|\), hence at most \(K_1^n\).  Absorbing \(K_1\) proves
\eqref{eq:v4v32-all-cell-product}.
\end{proof}

\subsection{Marked fermion-connected Yukawa coefficients}
\label{sec:v4v32-marked-connected}

For a nonempty finite set \(S=\{x_1,\ldots,x_n\}\) of cells, define the
conditional fermion-connected coefficient
\begin{equation}
 \mathcal A_{\rm F}(S;\phi,B)
 =
 \left\langle
 w_{X_1};\ldots;w_{X_n}
 \right\rangle_C^{\rm T}.
\label{eq:v4v32-connected-coefficient}
\end{equation}
In this subsection, \(D_B\) acts on the Yukawa coefficients in the
\(w_{X_i}\)'s while the hard covariance \(C\) is held fixed.

\begin{theorem}[Twice-marked fermion-connected tree estimate]
\label{thm:v4v32-marked-tree}
There is \(K<\infty\) such that, for \(j=0,1,2\),
\begin{equation}
 \mathbb E_{\mu_B}
 \sup_{\|v_k\|\le1}
 \left|
 D_B^j\mathcal A_{\rm F}(S)[v_1,\ldots,v_j]
 \right|
 \le
 n^j(K\epsilon_Y)^n
 \sum_{T\in\mathcal T_n}
 \prod_{\{i,l\}\in E(T)}c(x_i,x_l).
\label{eq:v4v32-marked-tree}
\end{equation}
Every derivative on the left is placed on an actual Yukawa cell according
to the product rule.
\end{theorem}

\begin{proof}
Differentiate the multilinear expression
\eqref{eq:v4v32-truncated}.  For one derivative there are \(n\) placements.
For two derivatives there are \(n\) same-cell second derivatives and at
most \(n(n-1)\) ordered two-cell placements.  Thus there are at most \(n^j\)
terms for order \(j\).

Apply Lemma~\ref{lem:v4v32-cell-tree} to each term.  Its product of
Grassmann norms is one of the products covered by
Lemma~\ref{lem:v4v32-all-cell-product}.  Taking the Higgs expectation and
combining \(K_{\rm F}\) with \(K_{\rm H}\) yields
\eqref{eq:v4v32-marked-tree}.
\end{proof}

\begin{corollary}[Absolute summability of the conditional effective action]
\label{cor:v4v32-summability}
Let \(\mathfrak t(S)\) be the length of a shortest lattice tree spanning
\(S\).  For every \(0<a'<a\) and \(h>0\), there is
\(\epsilon_*(a',h,q,\gamma,\chi_a,b,\nu,\kappa)>0\) such that, for
\(\epsilon_Y<\epsilon_*\) and \(j=0,1,2\),
\begin{equation}
 \sup_{x\in\mathbb L}
 \sum_{\substack{S\Subset\mathbb L\\x\in S}}
 e^{h|S|+a'\mathfrak t(S)}
 \mathbb E_{\mu_B}
 \sup_{\|v_k\|\le1}
 |D_B^j\mathcal A_{\rm F}(S)[v_1,\ldots,v_j]|
 <\infty.
\label{eq:v4v32-summability}
\end{equation}
The left side is \(O(\epsilon_Y)\) as
\(\epsilon_Y\downarrow0\).
\end{corollary}

\begin{proof}
For \(|S|=n\), replace the set sum by an ordered sum over
\(x_2,\ldots,x_n\), divided by \((n-1)!\).  For every tree \(T\) spanning
\(S\),
\[
 \mathfrak t(S)\le
 \sum_{\{i,l\}\in E(T)}d(x_i,x_l).
\]
Thus the factor \(e^{a'\mathfrak t(S)}\) can be paid edge by edge.
Root every labeled tree at vertex one and sum its other cell positions
successively from the leaves.  Equation
\eqref{eq:v4v32-covariance-majorant}, with \(a'<a\), bounds each weighted
edge sum by \(\chi_a\).  Cayley's formula gives \(n^{n-2}\) labeled trees.
Therefore the order-\(n\) contribution is bounded by
\[
 e^{hn}n^j(K\epsilon_Y)^n
 \frac{n^{n-2}}{(n-1)!}\chi_a^{n-1}.
\]
Stirling's lower bound implies
\(n^{n-2}/(n-1)!\le K_2^n\) for a universal \(K_2\).
The resulting series, including the polynomial factor \(n^j\), is
geometric after reducing \(\epsilon_Y\).  Its first term is
\(O(\epsilon_Y)\).
\end{proof}

\begin{firewall}[This is the conditional fermionic logarithm]
\label{fw:v4v32-conditional-log}
The coefficients \(\mathcal A_{\rm F}(S;\phi,B)\) form the logarithm after
fermionic Gaussian expectation at fixed \(\phi\).  Averaging these
coefficients absolutely over \(\mu_B\) proves
\eqref{eq:v4v32-summability}, but it is not the identity
\[
 \log\mathbb E_{\mu_B}e^{\sum_S\mathcal A_{\rm F}(S;\phi,B)}
 =
 \sum_S\mathbb E_{\mu_B}\mathcal A_{\rm F}(S;\phi,B).
\]
Higgs correlations between several fermion-connected components remain in
the logarithm on the left.
\end{firewall}

\begin{firewall}[Covariance marks remain in their resolvent ancestry]
\label{fw:v4v32-covariance-marks}
Theorem~\ref{thm:v4v32-marked-tree} holds \(C\) fixed.  A gauge or metric
background derivative also differentiates \(C=A_0^{-1}\), producing the
actual V16 resolvent insertion
\(-C(D_BA_0)C\), and at second order its two-insertion and contact terms.
Those marks must be interleaved with the Yukawa-cell marks before the final
marked norm is claimed.
\end{firewall}

\subsection{V32 result ledger and V33 target}
\label{sec:v4v32-conclusion}

V32 retains each antisymmetric Wick determinant, proves a Gram bound for all
forest-interpolated minors, and extracts the exact connected cell support.
The correlated all-cell Higgs moment costs only a constant per cell because
V31 applies to the union of supports and the overlap is bounded.  These
facts give an absolutely summable conditional fermionic effective action
through two actual Yukawa-coefficient marks for small normalized coupling.

V33 should assemble the two omitted connection mechanisms.  First it should
differentiate the forest formula when the hard covariance depends on the
background and assign every resolvent insertion to V16.  Second it should
take the logarithm with respect to the correlated Higgs parent, using the
V30 Witten resolvent for links between distinct fermion-connected
components.  A valid theorem must use one mixed spanning tree whose edges
are typed as fermion covariance or Higgs response; separate exponentiation
of the two connected series would miss their common connectivity.


\clearpage
\section{V33: covariance marks in the fermion-connected Yukawa action}
\label{sec:v4v33}

V32 held the hard fermion covariance fixed while differentiating the local
Yukawa coefficients.  Gauge and metric backgrounds also change that
covariance.  This note closes that conditional omission.  Weighted kernel
norms turn the exact inverse derivatives into exponentially decaying marked
lines.  Differentiating a Wick determinant selects actual covariance
entries and leaves complementary minors, so Gram's inequality continues to
remove the factorial.  The result is a twice-marked conditional fermionic
effective action including both coefficient and covariance marks.

The Higgs probability measure is still evaluated at the chosen background;
its own derivative is not taken in this note.  That derivative produces the
V30 Higgs score and is the remaining mixed-connection problem.

\subsection{Weighted resolvent derivatives}
\label{sec:v4v33-resolvent}

For a block kernel \(K\), define
\begin{equation}
 \|K\|_a
 =
 \sup_i\sum_j
 e^{a d(x(i),x(j))}|K_{ij}|.
\label{eq:v4v33-weighted-kernel}
\end{equation}
Finite internal fibres are included in the indices \(i,j\).

\begin{lemma}[Weighted kernel algebra]
\label{lem:v4v33-kernel-algebra}
If \(K,L\) have finite weighted norms, then
\begin{equation}
 \|KL\|_a\le\|K\|_a\|L\|_a.
\label{eq:v4v33-kernel-algebra}
\end{equation}
In particular,
\[
 |K_{ij}|\le\|K\|_a e^{-a d(x(i),x(j))}.
\]
\end{lemma}

\begin{proof}
The triangle inequality for the block metric gives
\[
 e^{a d(x(i),x(k))}
 \le
 e^{a d(x(i),x(j))}
 e^{a d(x(j),x(k))}.
\]
Insert this into the row sum for \(KL\), sum first over \(k\), then over
\(j\), and take the supremum in \(i\).  The pointwise estimate is one term
of the defining row sum.
\end{proof}

Let \(A_0(B)\) be the V16 hard fermion operator and
\(C(B)=A_0(B)^{-1}\).  For two background directions \(v_1,v_2\), direct
differentiation of \(A_0C=1\) gives
\begin{align}
 D_BC[v_1]
 &=-C\,D_BA_0[v_1]\,C,
\label{eq:v4v33-first-resolvent}\\
 D_B^2C[v_1,v_2]
 &=
 C\,D_BA_0[v_2]\,C\,D_BA_0[v_1]\,C
 +C\,D_BA_0[v_1]\,C\,D_BA_0[v_2]\,C
 \notag\\
 &\quad
 -C\,D_B^2A_0[v_1,v_2]\,C.
\label{eq:v4v33-second-resolvent}
\end{align}

\begin{theorem}[Marked hard-covariance decay]
\label{thm:v4v33-marked-covariance}
Assume that for some \(a_0>0\),
\begin{equation}
 \|C\|_{a_0}\le M_C
\label{eq:v4v33-C-bound}
\end{equation}
and, for \(m=1,2\),
\begin{equation}
 \sup_{\|v_k\|\le1}
 \|D_B^mA_0[v_1,\ldots,v_m]\|_{a_0}
 \le u_m\rho^{-m}.
\label{eq:v4v33-A-marks}
\end{equation}
Then
\begin{align}
 \sup_{\|v\|\le1}\|D_BC[v]\|_{a_0}
 &\le M_C^2u_1\rho^{-1},
\label{eq:v4v33-DC-bound}\\
 \sup_{\|v_k\|\le1}\|D_B^2C[v_1,v_2]\|_{a_0}
 &\le
 \left(2M_C^3u_1^2+M_C^2u_2\right)\rho^{-2}.
\label{eq:v4v33-D2C-bound}
\end{align}
The bounds are uniform in volume.
\end{theorem}

\begin{proof}
Apply Lemma~\ref{lem:v4v33-kernel-algebra} to every product in
\eqref{eq:v4v33-first-resolvent} and
\eqref{eq:v4v33-second-resolvent}.  The two ordered first-derivative terms
give the factor two in \eqref{eq:v4v33-D2C-bound}.
\end{proof}

\begin{remark}[Input from the V16 hard gap]
\label{rem:v4v33-V16-input}
The V16 Combes--Thomas estimate gives
\eqref{eq:v4v33-C-bound}, with any sufficiently small exponential rate,
from the hard inverse bound and finite range.  Local stencil derivatives of
\(A_0\) give \eqref{eq:v4v33-A-marks}.  Thus
Theorem~\ref{thm:v4v33-marked-covariance} is an explicit consequence of the
existing fixed-chart hard hypotheses, rather than a new fermion-gap
assumption.
\end{remark}

\subsection{Marked minors}
\label{sec:v4v33-minors}

For row and column lists \(J=(j_1,\ldots,j_m)\) and
\(I=(i_1,\ldots,i_m)\), put
\(\Delta_{J,I}(C)=\det(C_{j_ai_b})\).

\begin{lemma}[Exact first and second determinant marks]
\label{lem:v4v33-determinant-marks}
For any differentiable matrix path \(C(B)\),
\begin{align}
 D_B\Delta_{J,I}[v]
 &=
 \sum_{a,b}
 \operatorname{cof}_{ab}(C_{J,I})
 (D_BC[v])_{j_ai_b},
\label{eq:v4v33-first-minor}\\
 D_B^2\Delta_{J,I}[v_1,v_2]
 &=
 \sum_{a,b}
 \operatorname{cof}_{ab}(C_{J,I})
 (D_B^2C[v_1,v_2])_{j_ai_b}
 \notag\\
 &\quad+
 \sum_{\substack{a\ne c\\b\ne d}}
 \varepsilon_{ac;bd}
 \det C_{J\setminus\{j_a,j_c\},
         I\setminus\{i_b,i_d\}}
 (D_BC[v_1])_{j_ai_b}
 (D_BC[v_2])_{j_ci_d}.
\label{eq:v4v33-second-minor}
\end{align}
Here \(\varepsilon_{ac;bd}\in\{-1,1\}\) is the cofactor sign.  These
identities are valid even when \(C_{J,I}\) is singular.
\end{lemma}

\begin{proof}
Use the permutation formula for the determinant.  One derivative selects
one actual matrix entry in each permutation and leaves its complementary
cofactor.  Two derivatives either both hit the same selected entry, giving
the first line of \eqref{eq:v4v33-second-minor}, or hit two entries with
different rows and columns, giving the second line.  This derivation uses no
inverse matrix, so singular minors cause no exception.
\end{proof}

\begin{corollary}[Gram control after two covariance marks]
\label{cor:v4v33-marked-Gram}
Under the V32 Gram hypothesis and
Theorem~\ref{thm:v4v33-marked-covariance}, a first marked \(m\)-pair Wick
determinant is bounded by \(m^2\) choices of one
\(D_BC\) entry times a Gram bound for an
\((m-1)\)-pair minor.  A second marked determinant is bounded by
\(m^2\) choices of one \(D_B^2C\) entry and at most \(m^4\) choices of two
\(D_BC\) entries, times the corresponding Gram minors.
\end{corollary}

\begin{proof}
Apply Lemma~\ref{lem:v4v33-determinant-marks}.  Every complementary
determinant is a minor of the unmarked covariance and is bounded by
Lemma~\ref{lem:v4v32-Gram}.  Bound the marked entries with
\eqref{eq:v4v33-DC-bound}--\eqref{eq:v4v33-D2C-bound}.
\end{proof}

\begin{firewall}[A covariance mark is an actual inverse derivative]
\label{fw:v4v33-actual-resolvent}
Equations \eqref{eq:v4v33-first-resolvent} and
\eqref{eq:v4v33-second-resolvent} retain the ordered two-insertion terms and
the genuine \(D_B^2A_0\) contact term.  A bound for
\(D_BC\) is not squared in place of \(D_B^2C\).
\end{firewall}

\subsection{The fully marked conditional fermion tree}
\label{sec:v4v33-full-tree}

Define nonnegative cell kernels
\begin{align}
 c_0(x,y)
 &:=
 \gamma^{-2}
 \max_{\substack{i:x(i)=x\\j:x(j)=y}}|C_{ij}|,
\label{eq:v4v33-c0}\\
 c_m(x,y)
 &:=
 \rho^m
 \sup_{\|v_k\|\le1}
 \max_{\substack{i:x(i)=x\\j:x(j)=y}}
 |D_B^mC[v_1,\ldots,v_m]_{ij}|,
 \quad m=1,2,
\label{eq:v4v33-cm}\\
 \widehat c(x,y)&:=c_0(x,y)+c_1(x,y)+c_2(x,y).
\label{eq:v4v33-chat}
\end{align}
After multiplying by a fixed fibre constant, the weighted row sum
\begin{equation}
 \widehat\chi_{a_0}
 :=
 \sup_x\sum_y e^{a_0d(x,y)}\widehat c(x,y)
 <\infty
\label{eq:v4v33-chat-row}
\end{equation}
follows from Theorem~\ref{thm:v4v33-marked-covariance}.

For the local Yukawa activity, set
\[
 \widehat A_i
 =
 \sum_{m=0}^2\rho^m
 \sup_{\|v_k\|\le1}
 \|D_B^mw_{X_i}[v_1,\ldots,v_m]\|_{\mathfrak G}.
\]

\begin{lemma}[Forest bound with coefficient and covariance marks]
\label{lem:v4v33-fully-marked-forest}
There is \(K_0=K_0(q,\gamma,M_C,u_1,u_2)\) such that, for
\(j=0,1,2\),
\begin{equation}
 \rho^j
 \sup_{\|v_k\|\le1}
 \left|
 D_B^j
 \left\langle
 w_{X_1};\ldots;w_{X_n}
 \right\rangle_{C(B)}^{\rm T}
 [v_1,\ldots,v_j]
 \right|
 \le
 n^{2j}K_0^n
 \prod_{i=1}^n\widehat A_i
 \sum_{T\in\mathcal T_n}
 \prod_{\{i,l\}\in E(T)}
 \widehat c(x_i,x_l).
\label{eq:v4v33-fully-marked-forest}
\end{equation}
\end{lemma}

\begin{proof}
Differentiate the same positive-semidefinite forest interpolation used in
Lemma~\ref{lem:v4v32-cell-tree}, holding its interpolation parameters fixed.
A derivative has two exhaustive placements: on a local polynomial or on a
covariance entry.  Polynomial placements are bounded by
\(\widehat A_i\).  Covariance placements are described exactly by
Lemma~\ref{lem:v4v33-determinant-marks}, both for extracted tree lines and
for entries inside the remaining Wick minor.

An extracted marked line is bounded by \(c_1\) or \(c_2\); an unmarked line
is bounded by \(c_0\).  If a mark occurs inside the remaining determinant,
Corollary~\ref{cor:v4v33-marked-Gram} leaves one or two marked entries and
an unmarked Gram minor.  There are at most \((qn)^2\) choices per
derivative.  For \(j\le2\), enlarge \(K_0\) and bound all choice factors by
\(n^{2j}K_0^n\).  Multiplication by \(\rho^j\) distributes the mark weights
among the selected factors.  Replacing every edge kernel by
\(\widehat c\) and every local marked norm by \(\widehat A_i\) then proves
\eqref{eq:v4v33-fully-marked-forest}.
\end{proof}

\begin{theorem}[Full twice-marked conditional Yukawa summability]
\label{thm:v4v33-full-summability}
Assume the V31 coefficient hypotheses, the bounded-overlap condition
\eqref{eq:v4v32-bounded-overlap}, and the hard-covariance hypotheses
\eqref{eq:v4v33-C-bound}--\eqref{eq:v4v33-A-marks}.  For every
\(0<a'<a_0\) and \(h>0\), there is
\(\epsilon_*>0\) such that, when
\(\epsilon_Y<\epsilon_*\),
\begin{equation}
 \sup_{x\in\mathbb L}
 \sum_{\substack{S\Subset\mathbb L\\x\in S}}
 e^{h|S|+a'\mathfrak t(S)}
 \sum_{j=0}^2\rho^j
 \mathbb E_{\mu_B}
 \sup_{\|v_k\|\le1}
 \left|
 D_B^j\mathcal A_{\rm F}(S;\phi,B)
 [v_1,\ldots,v_j]
 \right|
 <\infty.
\label{eq:v4v33-full-summability}
\end{equation}
The entire left side is \(O(\epsilon_Y)\).
\end{theorem}

\begin{proof}
The pointwise V31 estimates and the proof of
Lemma~\ref{lem:v4v32-all-cell-product} apply to
\(\widehat A_i\), giving
\[
 \mathbb E_{\mu_B}\prod_{i=1}^n\widehat A_i
 \le(K_{\rm H}\epsilon_Y)^n.
\]
Apply Lemma~\ref{lem:v4v33-fully-marked-forest}.  Root each labeled tree,
pay \(e^{a'\mathfrak t(S)}\) along its edges, and sum cell locations using
\eqref{eq:v4v33-chat-row}.  The remaining series is bounded by
\[
 \sum_{n\ge1}
 e^{hn}n^4(K_1\epsilon_Y)^n
 \frac{n^{n-2}}{(n-1)!}
 \widehat\chi_{a_0}^{\,n-1}.
\]
As in V32, the tree/factorial ratio is at most exponential.  The displayed
series converges and is \(O(\epsilon_Y)\) for sufficiently small
\(\epsilon_Y\).
\end{proof}

\begin{corollary}[No missing conditional background placement]
\label{cor:v4v33-no-missing-placement}
Before differentiating the Higgs expectation, the order-zero, first, and
second background responses of the conditional fermionic logarithm have
all their placements represented in
\eqref{eq:v4v33-full-summability}: coefficient marks occur in V31 local
polynomials, and covariance marks occur in the exact resolvent words
\eqref{eq:v4v33-first-resolvent}--%
\eqref{eq:v4v33-second-resolvent}.
\end{corollary}

\begin{proof}
The background dependence of
\(\mathcal A_{\rm F}(S;\phi,B)\) enters only through its local polynomials
and the Gaussian covariance.  The ordinary product rule through order two
has precisely the placements listed in the statement, and
Lemma~\ref{lem:v4v33-fully-marked-forest} includes each one.
\end{proof}

\begin{firewall}[The quadratic determinant remains single]
\label{fw:v4v33-determinant-once}
Differentiating the normalized covariance \(C=A_0^{-1}\) does not create a
new factor \(\det A_0\).  That determinant is already the V16/V22 Schur
normalization.  V33 differentiates only the normalized Gaussian
expectation of the V31 Yukawa polynomial.
\end{firewall}

\begin{firewall}[The outer Higgs logarithm remains open]
\label{fw:v4v33-outer-Higgs-open}
Equation \eqref{eq:v4v33-full-summability} controls the conditional
fermionic logarithm and all its intrinsic marks.  Differentiating
\[
 \log\mathbb E_{\mu_B}
 \exp\!\left(\sum_S\mathcal A_{\rm F}(S;\phi,B)\right)
\]
also differentiates \(\mu_B\), producing Higgs scores and connected
correlations between several \(\mathcal A_{\rm F}\)'s.  V30 supplies their
two-point carrier, but an all-orders mixed tree has not yet been proved.
\end{firewall}

\subsection{V33 result ledger and V34 target}
\label{sec:v4v33-conclusion}

V33 proves a weighted kernel algebra, exact first and second hard-covariance
resolvent words, and their volume-uniform exponential decay.  Exact
differentiation of Wick determinants leaves Gram-controlled complementary
minors.  Inserting those marked lines into the V32 forest closes the
conditional Yukawa effective action through two background marks, including
both coefficient and covariance placements.

The remaining outer-Higgs logarithm should not be attacked by asserting
that two-point decay implies every cumulant bound.  V34 should instead
construct a good-field/tail split adapted to the hard fermion gap.  On a
cell satisfying \(\epsilon_Y|\phi_X|\ll\delta_{\rm F}\), the Yukawa
perturbation enters the convergent V16 trace-word calculus.  Simultaneous
large Higgs cells should be paid directly by
\eqref{eq:v4v31-quadratic-exponential}, with one actual tail occurrence per
failed cell.  If this split closes, it avoids an unproved higher-cumulant
theorem for the correlated Higgs law.


\clearpage
\section{V34: smooth good-field Yukawa promotion and rare-tail ownership}
\label{sec:v4v34}

V33 closes the conditional fermionic logarithm, but integrating its
exponential against the correlated Higgs law would require more than the
two-point estimate of V30.  This note goes upstream.  We split the local
Yukawa multiplication operator with a smooth Higgs-amplitude cutoff.  The
good part is uniformly small relative to the hard fermion covariance, so
its determinant is zero-free.  Its modulus can be absorbed into the
positive Higgs parent.  A direct Hessian estimate shows that this promoted
parent remains uniformly convex.  The complementary part owns an actual
large-Higgs occurrence, and the V31 exponential moment pays every
simultaneous tail cell.

This produces an exact modulus/phase/tail ancestry.  It does not infer an
all-orders Higgs cumulant theorem from covariance decay.

\subsection{A background-independent smooth split}
\label{sec:v4v34-split}

Fix \(0<\alpha<1\) and, for \(0<\epsilon_Y<1\), put
\begin{equation}
 R_Y=\epsilon_Y^{-\alpha}.
\label{eq:v4v34-threshold}
\end{equation}
Let \(\chi\in C^\infty([0,\infty);[0,1])\) satisfy
\[
 \chi(t)=1\quad(t\le1),
 \qquad
 \chi(t)=0\quad(t\ge2).
\]
For a Yukawa cell \(X_x\), define the smooth radius
\[
 r_x(\phi)=\left(1+|\phi_{X_x}|^2\right)^{1/2},
 \qquad
 g_x(\phi)=\chi(r_x/R_Y),
 \qquad
 t_x=1-g_x.
\]
The cutoff depends on the integration variable, but not on the external
background coordinate \(B\).

Let \(V_x(B,\phi)\) be the finite-fibre fermion multiplication operator
corresponding to the local Yukawa bilinear.  Assume the physical affine
bounds
\begin{align}
 \|V_x\|&\le L\epsilon_Y(1+r_x),
\label{eq:v4v34-V0}\\
 \|D_\phi V_x[h]\|&\le L\epsilon_Y|h_{X_x}|,
\label{eq:v4v34-V1}\\
 D_\phi^2V_x&=0
\label{eq:v4v34-V2}
\end{align}
and the bounded-support/overlap hypotheses of V32.  Set
\begin{equation}
 V_{\rm G}=\sum_xg_xV_x,
 \qquad
 V_{\rm T}=\sum_xt_xV_x.
\label{eq:v4v34-GT}
\end{equation}
Thus \(V_{\rm G}+V_{\rm T}=\sum_xV_x\) exactly.

\begin{lemma}[Good and tail derivative bounds]
\label{lem:v4v34-GT-derivatives}
There is \(C_\chi<\infty\), independent of volume and
\(\epsilon_Y\), such that
\begin{align}
 \|g_xV_x\|
 &\le C_\chi L\epsilon_Y^{1-\alpha},
\label{eq:v4v34-good-size}\\
 \|D_\phi(g_xV_x)[h]\|
 &\le C_\chi L\epsilon_Y|h_{X_x}|,
\label{eq:v4v34-good-first}\\
 \|D_\phi^2(g_xV_x)[h,k]\|
 &\le
 C_\chi L\epsilon_Y^{1+\alpha}
 |h_{X_x}||k_{X_x}|.
\label{eq:v4v34-good-second}
\end{align}
Moreover, \(t_xV_x\) vanishes unless \(r_x\ge R_Y\).
\end{lemma}

\begin{proof}
On the support of \(g_x\), \(r_x\le2R_Y\), which gives
\eqref{eq:v4v34-good-size}.  The radius satisfies
\[
 \|D_\phi r_x\|\le1,
 \qquad
 \|D_\phi^2r_x\|\le r_x^{-1}\le1.
\]
Hence
\[
 \|D_\phi g_x\|\le C_\chi R_Y^{-1},
 \qquad
 \|D_\phi^2g_x\|\le C_\chi R_Y^{-2}.
\]
Apply the product rule, use \eqref{eq:v4v34-V0}--%
\eqref{eq:v4v34-V2}, and use \(r_x\le2R_Y\) wherever a derivative of
\(g_x\) is nonzero.  The first derivative is
\(O(\epsilon_Y)\).  The second is
\(O(\epsilon_Y/R_Y)=O(\epsilon_Y^{1+\alpha})\).
Finally, \(t_x=0\) for \(r_x\le R_Y\).
\end{proof}

\subsection{Zero-free good determinant}
\label{sec:v4v34-good-determinant}

Let \(A_0\) and \(C=A_0^{-1}\) be the hard fermion operator and covariance.
Use the weighted norm \eqref{eq:v4v33-weighted-kernel} at an exponent
\(a_0>0\), and assume \(\|C\|_{a_0}\le M_C\).  Since the cells have bounded
overlap and \(V_{\rm G}\) is a local multiplication operator,
Lemma~\ref{lem:v4v34-GT-derivatives} gives
\begin{equation}
 \|CV_{\rm G}\|_{a_0}
 \le C_0M_CL\epsilon_Y^{1-\alpha}
 =:\zeta_Y.
\label{eq:v4v34-good-walk}
\end{equation}

\begin{theorem}[Exact good/tail determinant factorization]
\label{thm:v4v34-factorization}
For sufficiently small \(\epsilon_Y\), \(\zeta_Y<1/2\),
\[
 Q_{\rm G}:=(1+CV_{\rm G})^{-1}
\]
exists in the weighted kernel algebra and
\begin{equation}
 C_{\rm G}:=(A_0+V_{\rm G})^{-1}=Q_{\rm G}C,
 \qquad
 \|C_{\rm G}\|_{a_0}\le2M_C.
\label{eq:v4v34-CG}
\end{equation}
At every finite cutoff,
\begin{equation}
 \frac{\det(A_0+V_{\rm G}+V_{\rm T})}{\det A_0}
 =
 \underbrace{\det(1+CV_{\rm G})}_{\mathcal D_{\rm G}}
 \underbrace{\det(1+C_{\rm G}V_{\rm T})}_{\mathcal D_{\rm T}}.
\label{eq:v4v34-determinant-factorization}
\end{equation}
The good determinant \(\mathcal D_{\rm G}\) never vanishes.
\end{theorem}

\begin{proof}
The Neumann series converges in the Banach algebra of
Lemma~\ref{lem:v4v33-kernel-algebra}; it gives
\(\|Q_{\rm G}\|_{a_0}\le(1-\zeta_Y)^{-1}\le2\).
Since
\[
 A_0+V_{\rm G}=A_0(1+CV_{\rm G}),
\]
inversion gives \eqref{eq:v4v34-CG}.  Also
\[
 A_0+V_{\rm G}+V_{\rm T}
 =
 (A_0+V_{\rm G})(1+C_{\rm G}V_{\rm T}).
\]
Taking finite-dimensional determinants proves
\eqref{eq:v4v34-determinant-factorization}.  Invertibility of
\(1+CV_{\rm G}\) proves the last assertion.
\end{proof}

\begin{firewall}[The split changes no determinant]
\label{fw:v4v34-exact-split}
Equation \eqref{eq:v4v34-determinant-factorization} is an algebraic
factorization of the one original fermion determinant.  The first factor
contains the good multiplication operator and the second contains its
exact complement with the updated covariance \(C_{\rm G}\).  Neither
factor may be multiplied by a second copy of the V16 quadratic
normalization.
\end{firewall}

\subsection{Convexity after promoting the good modulus}
\label{sec:v4v34-convexity}

Because \(\mathcal D_{\rm G}\ne0\), define on the fixed anomaly-normalized
branch
\begin{equation}
 \ell_{\rm G}(\phi,B)
 =
 \operatorname{Tr}\log(1+CV_{\rm G}),
 \qquad
 \log|\mathcal D_{\rm G}|=\operatorname{Re}\ell_{\rm G}.
\label{eq:v4v34-good-log}
\end{equation}
For a field direction \(h\), the exact derivatives are
\begin{align}
 D_\phi\ell_{\rm G}[h]
 &=
 \operatorname{Tr}\!\left(
 Q_{\rm G}C\,D_\phi V_{\rm G}[h]\right),
\label{eq:v4v34-good-first-response}\\
 D_\phi^2\ell_{\rm G}[h,h]
 &=
 \operatorname{Tr}\!\left(
 Q_{\rm G}C\,D_\phi^2V_{\rm G}[h,h]\right)
 \notag\\
 &\quad-
 \operatorname{Tr}\!\left(
 Q_{\rm G}C\,D_\phi V_{\rm G}[h]\,
 Q_{\rm G}C\,D_\phi V_{\rm G}[h]\right).
\label{eq:v4v34-good-second-response}
\end{align}

\begin{lemma}[Extensive and spatially weighted Hessian bound]
\label{lem:v4v34-log-Hessian}
There is a volume-independent \(C_1<\infty\) such that
\begin{equation}
 \left|
 D_\phi^2\log|\mathcal D_{\rm G}|[h,h]
 \right|
 \le
 C_1\left(
 \epsilon_Y^{1+\alpha}+\epsilon_Y^2
 \right)\|h\|_{\ell^2(\Lambda)}^2.
\label{eq:v4v34-log-Hessian}
\end{equation}
Moreover, for some \(0<a_1<a_0\), the field-block Hessian
\(H_{\rm G}=D_\phi^2\log|\mathcal D_{\rm G}|\) satisfies
\begin{equation}
 \sup_x\sum_y e^{a_1d(x,y)}
 \|(H_{\rm G})_{xy}\|
 \le
 C_1\left(
 \epsilon_Y^{1+\alpha}+\epsilon_Y^2
 \right).
\label{eq:v4v34-weighted-log-Hessian}
\end{equation}
\end{lemma}

\begin{proof}
Put \(R_{\rm G}=Q_{\rm G}C\).  Its weighted row norm and, by applying the
same estimate to the adjoint, its weighted column norm are uniformly
bounded.  The first trace in
\eqref{eq:v4v34-good-second-response} contains a local diagonal insertion
bounded by
\eqref{eq:v4v34-good-second}.  Summing its diagonal kernel and using bounded
overlap gives
\[
 \left|
 \operatorname{Tr}(R_{\rm G}D_\phi^2V_{\rm G}[h,h])
 \right|
 \le C\epsilon_Y^{1+\alpha}\sum_x|h_x|^2.
\]
For the second trace, the Hilbert--Schmidt inequality and the weighted
Schur bounds give
\[
 \left|
 \operatorname{Tr}
 (R_{\rm G}D_\phi V_{\rm G}[h])^2
 \right|
 \le
 \|R_{\rm G}D_\phi V_{\rm G}[h]\|_{\rm HS}^2
 \le C\epsilon_Y^2\sum_x|h_x|^2.
\]
At the kernel level, the first term joins only field variables in one
bounded cell.  The second joins cells \(x,y\) through a product of two
kernels of \(R_{\rm G}\), and hence is bounded by
\(C\epsilon_Y^2e^{-2a_1d(x,y)}\) after choosing \(a_1<a_0\).
Summing the weighted rows proves
\eqref{eq:v4v34-weighted-log-Hessian}.  Take real parts and combine the
estimates.
\end{proof}

Define the promoted positive measure
\begin{equation}
 d\nu_{\rm G,B}(\phi)
 =
 Z_{\rm G}(B)^{-1}
 |\mathcal D_{\rm G}(\phi,B)|\,d\mu_B(\phi).
\label{eq:v4v34-promoted-measure}
\end{equation}

\begin{theorem}[The promoted good-modulus parent remains convex]
\label{thm:v4v34-promoted-convexity}
For sufficiently small \(\epsilon_Y\), the action of
\(\nu_{\rm G,B}\) satisfies
\begin{equation}
 D_\phi^2
 \left(S_B-\log|\mathcal D_{\rm G}|\right)
 \ge\frac\kappa2\,1.
\label{eq:v4v34-promoted-convexity}
\end{equation}
Consequently the V30 Witten-resolvent covariance theorem, with the
exponentially decaying Hessian extension described in the proof below, and
the V31 subgaussian argument apply to \(\nu_{\rm G,B}\), with \(\kappa\)
replaced by \(\kappa/2\), provided its recentered mean remains uniformly
bounded.
\end{theorem}

\begin{proof}
By V30, \(D_\phi^2S_B\ge\kappa1\).
Lemma~\ref{lem:v4v34-log-Hessian} bounds the operator norm of the subtracted
Hessian by
\(C_1(\epsilon_Y^{1+\alpha}+\epsilon_Y^2)\).
Choose \(\epsilon_Y\) so this is at most \(\kappa/2\).

It remains to check the locality hypothesis in the Witten argument.
The original off-diagonal Hessian has finite range by V30, while
\eqref{eq:v4v34-weighted-log-Hessian} gives an exponentially summable
additional kernel.  If \(H\) obeys
\(\sup_x\sum_y e^{a_1d(x,y)}\|H_{xy}\|\le J_{a_1}\), then for
\(0<\eta<a_1\) the V30 weight satisfies
\[
 \|W_\eta H W_\eta^{-1}-H\|
 \le
 C\eta(a_1-\eta)^{-1}J_{a_1}.
\]
Indeed, its \(xy\) block gains
\(e^{-\eta(\psi(x)-\psi(y))}-1\); use
\[
 |e^{-\eta(\psi(x)-\psi(y))}-1|
 \le\eta d(x,y)e^{\eta d(x,y)}
\]
and the Schur test.  The right side tends to zero with \(\eta\).
Together with the one-form gap \(\kappa/2\), the same Neumann argument as
V30 yields exponential resolvent and covariance decay.

The Bakry--Emery argument also depends only on the lower Hessian bound.
The mean condition is stated separately because strong convexity controls
centered fluctuations, while the chosen field chart must control its
center.
\end{proof}

\begin{firewall}[The phase is not absorbed into the positive parent]
\label{fw:v4v34-phase-separate}
Only \(|\mathcal D_{\rm G}|\) enters
\eqref{eq:v4v34-promoted-measure}.  The unit-modulus factor
\[
 \exp(i\,\operatorname{Im}\ell_{\rm G})
 =
 \mathcal D_{\rm G}/|\mathcal D_{\rm G}|
\]
remains on the globally anomaly-normalized determinant branch.  Convexity
of the modulus does not prove phase convergence or reflection positivity.
\end{firewall}

\subsection{One exponential payer for every tail cell}
\label{sec:v4v34-tail}

The next lemma is stated for either \(\mu_B\), with convexity
\(\kappa\), or the promoted law \(\nu_{\rm G,B}\), with convexity
\(\kappa/2\).  Denote the applicable convexity constant by
\(\kappa_{\rm p}\), and assume
\[
 \sup_x|\mathbb E\phi_x|\le m_{\rm p}.
\]

\begin{lemma}[Simultaneous tail payment]
\label{lem:v4v34-simultaneous-tail}
Let \(x_1,\ldots,x_n\) be distinct cells satisfying the overlap bound
\eqref{eq:v4v32-bounded-overlap}.  There are \(c_{\rm T},K_{\rm T}>0\),
independent of \(n\), volume, and cell positions, such that
\begin{equation}
 \mathbb P\!\left(
 r_{x_i}\ge R_Y\ \hbox{for every }i
 \right)
 \le
 \left[
 K_{\rm T}e^{-c_{\rm T}R_Y^2}
 \right]^n.
\label{eq:v4v34-simultaneous-tail}
\end{equation}
More generally, for fixed \(p,c<\infty\),
\begin{equation}
 \mathbb E
 \prod_{i=1}^n
 \left[
 \boldsymbol1_{\{r_{x_i}\ge R_Y\}}
 (1+r_{x_i})^p e^{c\epsilon_Y(1+r_{x_i})}
 \right]
 \le
 \left[
 K_{p,c}e^{-c_{\rm T}R_Y^2}
 \right]^n.
\label{eq:v4v34-marked-tail}
\end{equation}
Thus the tail activity per actual failed cell is
\(O(e^{-c_{\rm T}\epsilon_Y^{-2\alpha}})\).
\end{lemma}

\begin{proof}
Choose \(\delta>0\) so small that
\(3\delta\nu<\kappa_{\rm p}/4\).  For \(u\ge R_Y\), polynomial growth and
the linear exponential can be absorbed as
\[
 (1+u)^pe^{c\epsilon_Y(1+u)}
 \le K_{p,c,\delta}e^{\delta u^2}.
\]
Also
\[
 \boldsymbol1_{\{u\ge R_Y\}}
 \le e^{-\delta R_Y^2}e^{\delta u^2}.
\]
Multiplying these inequalities gives at most
\(K e^{-\delta R_Y^2}e^{2\delta u^2}\) per cell.
Bounded overlap implies
\[
 \sum_i r_{x_i}^2
 \le n+\nu|\phi_U|^2,
 \qquad U=\bigcup_iX_{x_i},
 \qquad |U|\le bn.
\]
Apply the quadratic exponential estimate
\eqref{eq:v4v31-quadratic-exponential}, with
\(\kappa_{\rm p}\) in place of \(\kappa\), to the exponent
\(2\delta\nu|\phi_U|^2\).  Our choice of \(\delta\) leaves room below its
threshold.  The result is a fixed constant to the power \(|U|\), hence to
the power \(n\), times \(e^{-\delta nR_Y^2}\).  Reduce the exponent and
absorb the fixed constants to obtain
\eqref{eq:v4v34-marked-tail}.  The probability bound is the case
\(p=c=0\).
\end{proof}

For any finite cell set \(S\), the pointwise identity
\begin{equation}
 1
 =
 \sum_{T\subseteq S}
 \prod_{x\in T}\boldsymbol1_{\{r_x\ge R_Y\}}
 \prod_{x\in S\setminus T}\boldsymbol1_{\{r_x<R_Y\}}
\label{eq:v4v34-pattern-partition}
\end{equation}
assigns each actual tail cell to the set \(T\).  Smooth transition cells
are treated as tail cells by using the support condition \(r_x\ge R_Y\).

\begin{theorem}[Tail determinant has a super-small local parameter]
\label{thm:v4v34-tail-small}
Represent \(\mathcal D_{\rm T}\) in
\eqref{eq:v4v34-determinant-factorization} as the normalized Berezin
expectation with covariance \(C_{\rm G}\), leaving every \(t_xV_x\) as a
local Grassmann polynomial.  Under the V31--V33 finite-fibre hypotheses,
every connected coefficient containing \(n_{\rm T}\ge1\) distinct tail
cells has, through two background marks, an additional factor
\begin{equation}
 \left[
 K e^{-c_{\rm T}\epsilon_Y^{-2\alpha}}
 \right]^{n_{\rm T}}.
\label{eq:v4v34-tail-parameter}
\end{equation}
For small \(\epsilon_Y\), these coefficients satisfy the V33 marked tree
sum with an arbitrarily small tail activity parameter.
\end{theorem}

\begin{proof}
The covariance \(C_{\rm G}\) has the uniform weighted bound
\eqref{eq:v4v34-CG}; its first two background derivatives follow by the
same inverse identities as V33 because
\(\|CV_{\rm G}\|_{a_0}<1/2\).
Each tail Grassmann coefficient and its first two marks is bounded by a
fixed polynomial times
\(e^{c\epsilon_Y(1+r_x)}\), exactly as in V31, and vanishes unless
\(r_x\ge R_Y\).  Apply
\eqref{eq:v4v34-marked-tail} jointly to the distinct tail cells.  The
fermionic connected-support extraction and Gram estimate of V32--V33 then
supply the same tree product, with
\eqref{eq:v4v34-tail-parameter} multiplying the ordinary exponential
constant per tail cell.  Since the displayed parameter tends to zero
faster than any power of \(\epsilon_Y\), the marked tree sum converges after
reducing \(\epsilon_Y\).
\end{proof}

\begin{firewall}[No global good-field event]
\label{fw:v4v34-no-global-event}
The proof never conditions on the event that every cell in the volume is
good; that event would have a volume-dependent probability.  The exact
operator split \eqref{eq:v4v34-GT} is defined for every field.  Each
connected coefficient pays only the tail cells that actually occur in its
support.
\end{firewall}

\subsection{V34 result ledger and compulsory V35 target}
\label{sec:v4v34-conclusion}

V34 constructs a smooth, background-independent good/tail decomposition of
the Yukawa multiplication operator.  The good walk norm is
\(O(\epsilon_Y^{1-\alpha})\), giving an exact zero-free determinant and a
uniformly local updated covariance.  The Hessian of the good log modulus is
\(O(\epsilon_Y^{1+\alpha}+\epsilon_Y^2)\), so its promotion preserves
strong convexity.  Every complementary tail cell receives the simultaneous
payment \(e^{-c\epsilon_Y^{-2\alpha}}\), and the one original determinant
factors without duplicate normalization.

V35 is a compulsory companion audit.  It should then address the two
remaining pieces of this Yukawa subsystem: a quantitative bound on the
recentered mean of the promoted positive parent, and a connected,
anomaly-normalized estimate for the good determinant phase.  The tail
factor is already assigned; it must not be reintroduced when the phase is
treated.


\clearpage
\section{V35: companion audit, recentering, and the extensive phase test}
\label{sec:v4v35}

V34 promotes the modulus of the smooth good-field determinant into a
positive uniformly convex Higgs parent and leaves its phase separate.  This
compulsory checkpoint first records the active companion files.  It then
closes the mean bound left in V34, constructs the connected local expansion
of the good phase, and tests the correct infinite-volume normalization on
an exactly soluble model.

The test forces a change of target.  An extensive phase can have an
expectation tending to zero exponentially in volume while every normalized
local ratio has a stable limit.  The continuum programme must therefore
control a finite-volume zero-free phase pressure and local ratios after its
vacuum contribution is subtracted.  Requiring the unrenormalized phase
expectation itself to have a nonzero infinite-volume limit would reject even
the elementary Gaussian source.

\subsection{Compulsory companion audit}
\label{sec:v4v35-audit}

The exact files inspected at this checkpoint were
\begin{enumerate}
\item
\texttt{Renewal\_Yang\_Mills\_Mass\_Gap\_Research\_Notes\_V55.tex},
\(1001521\) bytes, SHA--256
\[
 \texttt{0DA16A5EA97FA36DE18FDECE43B459BA8E5B7DBCC22AE4219A29B811E9D8A647}.
\]
Its newest theorem combines the V54 factorization
\(p(xy+xz+yz)\) with a complete-word prefix response linear in the actual
prefix stock.  Squaring the actual operator word closes the pure
prefix-depletion face.  The three quotient-cut faces and their
cut-specific analytic response estimates remain open.
\item
\texttt{Renewal\_NS\_Stability\_Research\_Notes\_V31.tex},
\(887537\) bytes, SHA--256
\[
 \texttt{F1121B62238AD5491BB7CF03635EA9934942EDF573ED8FAAA9EE79E1D38A6696}.
\]
Its newest result proves that flat dyadic marks and heat lifting recover the
critical first moment only at the critical source scale.  The remaining
route must retain the correlation among actual recent-entry times, the
complete signed triad work, and the active high-parent tail; the
polarization hinge has an exact balanced-sector nullspace.
\end{enumerate}

The usable lesson is about proof order, not transfer of estimates.  YM V55
closes one classified face only after retaining and squaring its complete
operator word; its other cut faces still require their own analytic
responses.  NS V31 shows that separate positive bounds erase the
correlation that must supply a gain.  Accordingly, the SM determinant is
kept in the exact modulus/phase/tail ancestry of V34.  Convexity estimates
apply only to the modulus, while the complete signed phase is retained
until its normalized local ratios are formed.

\begin{firewall}[Companion type boundary]
\label{fw:v4v35-companion-types}
No Wilson tree currency or Navier--Stokes heat estimate is used as a
Standard-Model inequality.  The audit changes only the admissible proof
order: exact channel decomposition precedes the channel-specific analytic
estimate, and signed correlations are not replaced by their absolute
envelopes.
\end{firewall}

\subsection{Uniform recentering of the promoted parent}
\label{sec:v4v35-recentering}

Write
\[
 u_{\rm G}(\phi,B)=\log|\mathcal D_{\rm G}(\phi,B)|
\]
and interpolate
\begin{equation}
 d\nu_{t,B}
 =
 Z_t(B)^{-1}e^{t u_{\rm G}}\,d\mu_B,
 \qquad 0\le t\le1.
\label{eq:v4v35-interpolation}
\end{equation}
Thus \(\nu_{0,B}=\mu_B\) and
\(\nu_{1,B}=\nu_{\rm G,B}\).

\begin{lemma}[Local gradient of the good log modulus]
\label{lem:v4v35-local-gradient}
Under the V34 hypotheses, there is a volume-independent \(C_u\) such that
\begin{equation}
 \sup_{\phi,B,y}
 \left|\nabla_{\phi_y}u_{\rm G}(\phi,B)\right|
 \le C_u\epsilon_Y.
\label{eq:v4v35-local-gradient}
\end{equation}
Every interpolating action
\(S_B-tu_{\rm G}\) has Hessian at least \(\kappa/2\) and an exponentially
summable off-diagonal Hessian, uniformly in \(t\).
\end{lemma}

\begin{proof}
Equation \eqref{eq:v4v34-good-first-response} gives
\[
 D_{\phi_y}u_{\rm G}
 =
 \operatorname{Re}\operatorname{Tr}
 \left(Q_{\rm G}C\,D_{\phi_y}V_{\rm G}\right).
\]
The insertion \(D_{\phi_y}V_{\rm G}\) is supported in the bounded number of
cells containing \(y\) and has norm \(O(\epsilon_Y)\) by
\eqref{eq:v4v34-good-first}.  The weighted row and column norms of
\(Q_{\rm G}C\) are uniform by
\eqref{eq:v4v34-CG}; its relevant diagonal trace is therefore
\(O(\epsilon_Y)\), proving
\eqref{eq:v4v35-local-gradient}.

The Hessian assertion is the proof of
Theorem~\ref{thm:v4v34-promoted-convexity}, with the perturbing Hessian
multiplied by \(t\in[0,1]\).  The weighted estimate
\eqref{eq:v4v34-weighted-log-Hessian} is uniform in \(t\).
\end{proof}

\begin{lemma}[Uniform local Witten kernel along the interpolation]
\label{lem:v4v35-interpolated-Witten}
There are \(C_{\rm W},\eta_{\rm W}>0\), independent of volume and
\(t\in[0,1]\), such that the Witten one-form operator
\(\mathcal L_{1,t}\) of \(\nu_{t,B}\) satisfies
\begin{equation}
 \|P_x\mathcal L_{1,t}^{-1}P_y\|
 \le C_{\rm W}e^{-\eta_{\rm W}d(x,y)}.
\label{eq:v4v35-interpolated-Witten}
\end{equation}
\end{lemma}

\begin{proof}
Lemma~\ref{lem:v4v35-local-gradient} gives the uniform one-form gap
\(\mathcal L_{1,t}\ge\kappa/2\).  Its off-diagonal Hessian is the sum of the
V30 finite-range kernel and \(t\) times the exponentially summable V34
kernel.  The weighted-commutator estimate in the proof of
Theorem~\ref{thm:v4v34-promoted-convexity} tends to zero with the
conjugation exponent, uniformly in \(t\).  Choose that exponent so the
commutator norm is at most \(\kappa/4\), and apply the V30 Neumann argument.
\end{proof}

\begin{theorem}[Quantitative recentering]
\label{thm:v4v35-recentering}
There is a volume-independent \(C_m\) such that
\begin{equation}
 \sup_{B,x}
 \left|
 \mathbb E_{\nu_{\rm G,B}}\phi_x
 -
 \mathbb E_{\mu_B}\phi_x
 \right|
 \le C_m\epsilon_Y.
\label{eq:v4v35-recentering}
\end{equation}
In particular, the V34 promoted parent satisfies
\[
 \sup_{B,x}
 |\mathbb E_{\nu_{\rm G,B}}\phi_x|
 \le m_*+C_m\epsilon_Y,
\]
so the mean hypothesis in
Theorem~\ref{thm:v4v34-promoted-convexity} is closed.
\end{theorem}

\begin{proof}
Differentiate the normalized interpolation:
\[
 \frac{d}{dt}\mathbb E_{\nu_{t,B}}\phi_x
 =
 \operatorname{Cov}_{\nu_{t,B}}(\phi_x,u_{\rm G}).
\]
Apply the Helffer--Sjostrand identity for \(\nu_{t,B}\), insert spatial
projections, and use
\eqref{eq:v4v35-interpolated-Witten}:
\begin{align*}
 \left|
 \operatorname{Cov}_{\nu_{t,B}}(\phi_x,u_{\rm G})
 \right|
 &\le
 \sum_y
 \|P_x\mathcal L_{1,t}^{-1}P_y\|
 \|\nabla_{\phi_x}\phi_x\|_{L^2(\nu_{t,B})}
 \|\nabla_{\phi_y}u_{\rm G}\|_{L^2(\nu_{t,B})}\\
 &\le
 C_{\rm W}C_u\epsilon_Y
 \sum_y e^{-\eta_{\rm W}d(x,y)}
 \le C_m\epsilon_Y.
\end{align*}
Exponential balls are uniformly summable.  Integrate from \(t=0\) to one.
\end{proof}

\subsection{Connected local expansion of the good phase}
\label{sec:v4v35-phase-expansion}

Put
\[
 K_{\rm G}=CV_{\rm G},
 \qquad
 \vartheta_{\rm G}
 =
 \operatorname{Im}\operatorname{Tr}\log(1+K_{\rm G}).
\]
Since \(\|K_{\rm G}\|_{a_0}\le\zeta_Y<1/2\),
\begin{equation}
 \vartheta_{\rm G}
 =
 \sum_{n\ge1}\frac{(-1)^{n+1}}n
 \operatorname{Im}\operatorname{Tr}K_{\rm G}^n
\label{eq:v4v35-phase-series}
\end{equation}
converges absolutely.  Let \(K_{\rm G}=\sum_xK_x\), and define
\(\vartheta_{\rm G}^c(S)\) by the V16 Mobius extraction: only cyclic words
whose set of cell labels is exactly \(S\) remain.

\begin{theorem}[Marked connected phase norm]
\label{thm:v4v35-phase-norm}
For every \(0<a'<a_0\) and \(h>0\), there are
\(\epsilon_*>0\), \(\rho>0\), and \(C<\infty\) such that
\begin{equation}
 \sup_x
 \sum_{\substack{S\Subset\mathbb L\\x\in S}}
 e^{h|S|+a'\mathfrak t(S)}
 \sum_{j=0}^2\rho^j
 \sup_{\|v_k\|\le1}
 \left|
 D_B^j\vartheta_{\rm G}^c(S)
 [v_1,\ldots,v_j]
 \right|
 \le
 C\epsilon_Y^{1-\alpha}
\label{eq:v4v35-phase-norm}
\end{equation}
for \(\epsilon_Y<\epsilon_*\).  Pointwise in \(\phi\),
\begin{equation}
 \frac{\mathcal D_{\rm G}}{|\mathcal D_{\rm G}|}
 =
 \exp\!\left(
 i\sum_{\varnothing\ne S\Subset\mathbb L}
 \vartheta_{\rm G}^c(S)
 \right).
\label{eq:v4v35-phase-factorization}
\end{equation}
\end{theorem}

\begin{proof}
The good cutoff gives the pointwise local size
\(\|K_x\|=O(\epsilon_Y^{1-\alpha})\).
The hard covariance and its first two background derivatives have the V33
weighted exponential bounds.  Background derivatives of \(g_x\) vanish,
and those of \(V_x\) satisfy the V31 coefficient bounds on its capped
support.  Therefore the transition majorants in V16 have weighted walk norm
\(O(\epsilon_Y^{1-\alpha})<1\).

Apply Theorem~\ref{thm:v4v16-connected-word} to the imaginary part of
\eqref{eq:v4v35-phase-series}, retaining all one- and two-mark placements.
The V16 weighted trace-word proof, strengthened by the marked-minor
bookkeeping of V33, gives tree decay.  Rooted-tree summation as in V32
proves \eqref{eq:v4v35-phase-norm}.  Summing the exact connected Mobius
coefficients reconstructs \eqref{eq:v4v35-phase-series}; exponentiation and
\(\mathcal D_{\rm G}=e^{\operatorname{Tr}\log(1+K_{\rm G})}\) prove
\eqref{eq:v4v35-phase-factorization}.
\end{proof}

\begin{corollary}[Small local phase activities]
\label{cor:v4v35-phase-activities}
Define
\[
 z_{\rm ph}(S)=e^{i\vartheta_{\rm G}^c(S)}-1.
\]
After reducing \(\epsilon_*\), the left side of
\eqref{eq:v4v35-phase-norm} remains bounded by
\(C'\epsilon_Y^{1-\alpha}\) with
\(\vartheta_{\rm G}^c\) replaced by \(z_{\rm ph}\) and with its exact first
two derivatives.
\end{corollary}

\begin{proof}
Use
\[
 |e^{iu}-1|\le |u|e^{|\operatorname{Im}u|}
\]
and differentiate the exponential exactly through order two.  The phase
coefficients are real, so the exponential has modulus one.  Products of two
marked coefficients are controlled by the convolution algebra of the
rooted tree norm after reducing its small norm.
\end{proof}

\begin{firewall}[A small local phase norm is not a lower bound on its mean]
\label{fw:v4v35-phase-mean-open}
Theorem~\ref{thm:v4v35-phase-norm} is a deterministic connected expansion
at fixed Higgs field.  It does not imply a volume-uniform positive lower
bound for
\(\left|\mathbb E_{\nu_{\rm G,B}}e^{i\vartheta_{\rm G}}\right|\).
The phase has an extensive vacuum contribution, and its expectation may
decay with volume even when every local coefficient is small.
\end{firewall}

\subsection{The exact extensive Gaussian phase test}
\label{sec:v4v35-Gaussian-test}

\begin{proposition}[Vanishing phase mean with stable local ratios]
\label{prop:v4v35-Gaussian-phase}
Let \(\mu_N\) be the product standard Gaussian law on
\(\mathbb R^N\), and let
\[
 \vartheta_N(\phi)=\epsilon\sum_{x=1}^N\phi_x
\]
with fixed nonzero \(\epsilon\).  Then
\begin{equation}
 \mathbb E_{\mu_N}e^{i\vartheta_N}
 =
 e^{-\epsilon^2N/2}
 \longrightarrow0.
\label{eq:v4v35-Gaussian-denominator}
\end{equation}
Nevertheless, for every bounded observable \(F\) depending only on a fixed
set \(A\subset\{1,\ldots,N\}\),
\begin{equation}
 \frac{
 \mathbb E_{\mu_N}[F(\phi_A)e^{i\vartheta_N}]
 }{
 \mathbb E_{\mu_N}e^{i\vartheta_N}
 }
 =
 e^{\epsilon^2|A|/2}
 \mathbb E_{\mu_A}
 \left[
 F(\phi_A)e^{i\epsilon\sum_{x\in A}\phi_x}
 \right],
\label{eq:v4v35-Gaussian-local-ratio}
\end{equation}
which is independent of \(N\) once \(A\subset\{1,\ldots,N\}\).
\end{proposition}

\begin{proof}
For one standard Gaussian \(Z\),
\(\mathbb E e^{i\epsilon Z}=e^{-\epsilon^2/2}\).
Independence gives
\eqref{eq:v4v35-Gaussian-denominator}.  In the numerator, factor the
coordinates outside \(A\); they contribute
\(e^{-\epsilon^2(N-|A|)/2}\).  Division by the denominator leaves
\eqref{eq:v4v35-Gaussian-local-ratio}.
\end{proof}

\begin{corollary}[Corrected phase-convergence target]
\label{cor:v4v35-corrected-phase-target}
A nonzero infinite-volume limit of the unrenormalized phase expectation is
not a necessary condition for convergence of normalized local Schwinger
ratios.  The appropriate target is:
\begin{enumerate}
\item finite-volume zero-freeness on the anomaly-normalized branch;
\item existence of the phase pressure
\[
 p_{\rm ph}
 =
 \lim_{\Lambda\uparrow\mathbb L}
 |\Lambda|^{-1}
 \log\mathbb E_{\nu_{\rm G,B}}e^{i\vartheta_{\rm G}};
\]
and
\item convergence of local numerator/denominator ratios after the common
vacuum phase pressure cancels.
\end{enumerate}
\end{corollary}

\begin{proof}
Proposition~\ref{prop:v4v35-Gaussian-phase} violates the nonzero-limit
condition while satisfying exact volume-independent local ratios.  Its
logarithmic denominator per site is the finite pressure
\(-\epsilon^2/2\), and the exterior vacuum factor cancels from every local
ratio.  Hence the three stated properties capture the mechanism that the
stronger condition misses.
\end{proof}

\begin{firewall}[Correction of the V2 sufficient criterion]
\label{fw:v4v35-V2-correction}
Proposition~\ref{prop:v4v2-modulus-phase} remains a valid sufficient
criterion when its nonzero limiting expectation hypothesis holds.  V35
shows that the hypothesis is not necessary for an extensive phase and must
not be used as the required continuum gate.  No conclusion about the
physical CKM or gravitational phase follows from the Gaussian example.
\end{firewall}

\subsection{V35 result ledger and V36 target}
\label{sec:v4v35-conclusion}

V35 records the compulsory YM/NS audit and applies its proof-order lesson.
A localized Helffer--Sjostrand interpolation proves that promotion of the
good determinant modulus shifts every Higgs one-point mean by only
\(O(\epsilon_Y)\), closing the recentering hypothesis of V34.  The complete
good determinant phase has a connected, spatially decaying expansion
through two background marks of size
\(O(\epsilon_Y^{1-\alpha})\).  An exact Gaussian test then proves that its
unrenormalized expectation need not stay nonzero in the thermodynamic
limit; the correct objects are the phase pressure and normalized local
ratios.

V36 should attack finite-volume zero-freeness and the phase pressure for a
strongly convex, exponentially mixing parent carrying the small connected
phase potentials of
Theorem~\ref{thm:v4v35-phase-norm}.  The theorem must retain the full signed
phase.  If a general zero-free statement fails, the next step is to build
the smallest strongly convex counterexample and identify the additional
charge-conjugation, sector-pairing, or determinant-line hypothesis that
excludes it.


\clearpage
\section{V36: zero-free phase pressure in the product-reference regime}
\label{sec:v4v36}

V35 identifies the correct phase target but leaves it open over the general
correlated parent.  There is, however, one regime in which the missing
all-orders connection is already available: the V28 small-hopping
expansion over the onsite quartic product measure.  This note combines that
expansion with the V34 good/tail determinant split and the V35 connected
phase coefficients in one polymer gas.  The V15 complex cluster theorem
then gives finite-volume zero-freeness, a thermodynamic residual phase
pressure, and convergent normalized local ratios.

This is a closure theorem for the uniformly convex, small-hopping,
small-Yukawa, hard-fermion-gap chart.  It does not extend the result to the
arbitrary-hopping correlated-parent route of V30.

\subsection{Connectedization over a product measure}
\label{sec:v4v36-connectedization}

Let
\[
 d\mu_0(\phi)=\prod_{x\in\Lambda}d\mu_x(\phi_x)
\]
be the onsite quartic product parent used in V28.  An elementary factor
\(f_\gamma(\phi)\) has a finite support
\(\operatorname{supp}\gamma\subset\Lambda\), and write
\(f_\gamma=1+u_\gamma\).  Two labels are incompatible when their supports
intersect or are adjacent in a fixed finite-range dependency graph.

Expanding a finite family gives
\begin{equation}
 \mathbb E_{\mu_0}\prod_\gamma(1+u_\gamma)
 =
 \sum_{\mathcal F}
 \mathbb E_{\mu_0}\prod_{\gamma\in\mathcal F}u_\gamma.
\label{eq:v4v36-elementary-expansion}
\end{equation}
For a finite connected label family \(\mathcal C\), define
\begin{equation}
 W(\mathcal C)
 =
 \mathbb E_{\mu_0}
 \prod_{\gamma\in\mathcal C}u_\gamma,
 \qquad
 X(\mathcal C)=
 \bigcup_{\gamma\in\mathcal C}\operatorname{supp}\gamma.
\label{eq:v4v36-connected-weight}
\end{equation}

\begin{lemma}[Exact product-parent polymerization]
\label{lem:v4v36-product-polymerization}
Grouping every \(\mathcal F\) in
\eqref{eq:v4v36-elementary-expansion} into the connected components of its
label-overlap graph gives exactly
\begin{equation}
 \mathbb E_{\mu_0}\prod_\gamma(1+u_\gamma)
 =
 \sum_{\substack{\{\mathcal C_1,\ldots,\mathcal C_k\}\\
 X(\mathcal C_i)\cap X(\mathcal C_j)=\varnothing}}
 \prod_{i=1}^k W(\mathcal C_i).
\label{eq:v4v36-polymer-gas}
\end{equation}
If the elementary marked rooted norm is at most \(\eta\), then there are
\(C,a'>0\) such that the induced polymer activities satisfy
\begin{equation}
 \|W\|_{a',2;\rho}\le C\eta
\label{eq:v4v36-induced-norm}
\end{equation}
whenever \(\eta\) is sufficiently small.
\end{lemma}

\begin{proof}
Distinct connected components have disjoint dependency supports.  The
product measure therefore factors their expectations, proving
\eqref{eq:v4v36-polymer-gas} term by term.

For the norm bound, root a connected label family at a label meeting a
fixed site.  Every connected graph contains a spanning tree.  The standard
tree-graph bound replaces the sum of connected graphs by a fixed
exponential factor times a sum of rooted trees.  Sum leaves inward using
the elementary rooted norm.  The resulting series is geometric for small
\(\eta\), and its first term is \(O(\eta)\).  One background derivative
marks one actual elementary factor.  Two derivatives have the exhaustive
placements of Lemma~\ref{lem:v4v15-two-mark-identity}: one genuine second
mark or two first marks.  The same rooted-tree sum, with the factors
\(\rho,\rho^2\), proves \eqref{eq:v4v36-induced-norm}.
\end{proof}

\begin{firewall}[Where product structure is used]
\label{fw:v4v36-product-use}
The exact factorization in
\eqref{eq:v4v36-polymer-gas} uses independence of disjoint onsite
coordinates.  Exponential two-point covariance decay of a correlated
measure does not imply this identity.  Thus
Lemma~\ref{lem:v4v36-product-polymerization} cannot be inserted into the
V30 route without an additional all-orders connected theorem.
\end{firewall}

\subsection{The complete elementary species}
\label{sec:v4v36-species}

Use the exact V34 determinant factorization
\[
 \frac{\det(A_0+V_{\rm G}+V_{\rm T})}{\det A_0}
 =
 \mathcal D_{\rm G}\mathcal D_{\rm T}.
\]
Resolve its factors into the following elementary species over \(\mu_0\):
\begin{enumerate}
\item the V28 hopping-bond factors, with marked norm
      \(\eta_{\rm hop}\);
\item the connected real coefficients of
      \(\log|\mathcal D_{\rm G}|\), exponentiated as local modulus
      activities;
\item the phase activities
      \(z_{\rm ph}(S)=e^{i\vartheta_{\rm G}^c(S)}-1\) from V35;
\item the V34 tail Grassmann/determinant coefficients, expanded with
      \(C_{\rm G}\); and
\item mixed good-tail words generated by the Neumann series for
      \(C_{\rm G}=(1+CV_{\rm G})^{-1}C\).
\end{enumerate}
The fifth species is not optional: it records the dependence of every tail
contraction on the good operator.

\begin{lemma}[Complete marked elementary norm]
\label{lem:v4v36-complete-norm}
Under the V28, V31--V35 hypotheses, the five species above have a common
finite-range incompatibility graph and satisfy
\begin{equation}
 \eta_{\rm elem}
 \le
 C\left[
 \eta_{\rm hop}
 +\epsilon_Y^{1-\alpha}
 +e^{-c\epsilon_Y^{-2\alpha}}
 \right]
\label{eq:v4v36-complete-norm}
\end{equation}
through two background marks.  Every factor in the original hopping and
Yukawa integrand occurs exactly once.
\end{lemma}

\begin{proof}
V28 gives the first term.  The real and imaginary parts of the same
good-field trace-log words have identical absolute transition majorants;
Theorem~\ref{thm:v4v35-phase-norm} therefore bounds both the modulus and
phase species by \(C\epsilon_Y^{1-\alpha}\).  V34 gives the
super-small tail parameter.

For the mixed species, expand
\[
 C_{\rm G}
 =
 \sum_{n\ge0}(-CV_{\rm G})^nC.
\]
The weighted norm of \(CV_{\rm G}\) is
\(\zeta_Y=O(\epsilon_Y^{1-\alpha})<1/2\).
A mixed word containing a tail occurrence therefore has its tail parameter
from V34 and a summable geometric string of good factors.  Differentiating
through order two places marks on actual covariance or vertex factors as in
V33; the marked geometric series remains bounded after reducing
\(\epsilon_Y\).  This proves
\eqref{eq:v4v36-complete-norm}.

The hopping bonds come only from the Higgs kinetic action.  The good and
tail operators partition each Yukawa multiplication operator, and
Theorem~\ref{thm:v4v34-factorization} partitions the one determinant.
The Neumann words only expand the updated covariance inside its tail
factor.  Hence no original factor or determinant normalization is
duplicated.
\end{proof}

\subsection{Two zero-free partition functions}
\label{sec:v4v36-two-partitions}

Let \(Z_\Lambda^{+}\) be the product-reference polymer representation
containing the hopping and good-modulus species.  It represents the
positive good-modulus partition function.  Let \(Z_\Lambda^{\rm c}\) be the
same gas with the anomaly-normalized phase, tail, and mixed species also
inserted.  It represents the complete complex Yukawa partition function in
the fixed hard chart.  Both retain the same V16/V22 quadratic determinant
normalization outside their normalized polymer factors.

\begin{theorem}[Finite-volume zero-freeness and marked logarithms]
\label{thm:v4v36-zero-free}
There is \(\eta_*>0\) such that if the right side of
\eqref{eq:v4v36-complete-norm} is below \(\eta_*\), then
\[
 Z_\Lambda^{+}\ne0,
 \qquad
 Z_\Lambda^{\rm c}\ne0
\]
for every finite volume.  The V15 cluster branches
\(\log Z_\Lambda^{+}\) and
\(\log Z_\Lambda^{\rm c}\) are analytic continuations from zero activities
and have volume-uniform intensive first and second background responses.
Thus
\begin{equation}
 \mathfrak p_{\Lambda}^{\rm res}
 =
 |\Lambda|^{-1}
 \left(
 \log Z_\Lambda^{\rm c}
 -
 \log Z_\Lambda^{+}
 \right)
\label{eq:v4v36-residual-pressure}
\end{equation}
is defined on one zero-free branch.
\end{theorem}

\begin{proof}
Lemma~\ref{lem:v4v36-product-polymerization} converts the complete
elementary family into polymer activities with marked norm
\(C\eta_{\rm elem}\).  Reduce the couplings so that both the positive
subfamily and the complete family lie inside the V15 KP ball.  Apply
Theorem~\ref{thm:v4v15-two-mark} and
Proposition~\ref{prop:v4v15-phase} to each gas.  Their cluster logarithms
prevent zeros on the branches connected to the empty gas, and
Corollary~\ref{cor:v4v15-local-extensive} supplies the intensive response
bounds.  Subtraction gives \eqref{eq:v4v36-residual-pressure}.
\end{proof}

\subsection{Thermodynamic pressure and local ratios}
\label{sec:v4v36-thermodynamic}

Assume translation-covariant bulk activities and take a van Hove sequence
\(\Lambda_n\uparrow\mathbb L\).

\begin{theorem}[Residual phase pressure and local-ratio limit]
\label{thm:v4v36-pressure-ratios}
Under the hypotheses of
Theorem~\ref{thm:v4v36-zero-free}, the limit
\begin{equation}
 \mathfrak p^{\rm res}
 =
 \lim_{n\to\infty}\mathfrak p_{\Lambda_n}^{\rm res}
\label{eq:v4v36-pressure-limit}
\end{equation}
exists and is independent of the van Hove sequence.  If a bounded local
observable \(F\) admits the same finite marked polymer insertion, then
\begin{equation}
 \frac{
 \mathbb E_{\mu_0}
 \left[F\,(\hbox{complete normalized integrand})\right]
 }{
 Z_{\Lambda_n}^{\rm c}
 }
\label{eq:v4v36-local-ratio}
\end{equation}
converges as soon as \(\Lambda_n\) contains the fixed support of \(F\).
The exterior vacuum clusters cancel between numerator and denominator.
The same conclusions hold through two local background derivatives.
\end{theorem}

\begin{proof}
The absolutely convergent cluster expansion writes
\(\log Z_\Lambda\) as a sum of connected clusters.  Assign each
translation class to one distinguished root site, for example the
lexicographically first site of its support.  The rooted absolute sum is
finite by the marked KP bound.  Clusters whose translate lies wholly in
\(\Lambda_n\) contribute their translation density; clusters meeting the
boundary have total contribution bounded by the rooted norm times the
boundary-to-volume ratio.  The latter tends to zero along a van Hove
sequence.  Dominated convergence proves the pressure limit for each gas and
hence \eqref{eq:v4v36-pressure-limit}.

For the local ratio, use the cluster expansion with one distinguished
\(F\)-root.  Every cluster disconnected from the support of \(F\) occurs
identically in numerator and denominator and cancels.  The remaining rooted
clusters are absolutely summable uniformly in volume and stabilize by
dominated convergence.  The two-mark version follows from the exhaustive
marked placements in V15.
\end{proof}

\begin{corollary}[Closure of the V35 phase target in the KP regime]
\label{cor:v4v36-phase-closure}
In the small-hopping product-reference regime, the three corrected V35
phase requirements are satisfied: the finite-volume normalized complex
partition is zero-free, the residual phase/tail pressure exists, and local
normalized ratios converge after the common vacuum pressure cancels.
\end{corollary}

\begin{proof}
These are Theorems~\ref{thm:v4v36-zero-free} and
\ref{thm:v4v36-pressure-ratios}.  The anomaly-normalized determinant branch
is prior data by the V15 firewall and is retained by construction.
\end{proof}

\begin{firewall}[No double use of the correlated parent]
\label{fw:v4v36-alternative-route}
V36 uses the V28 onsite product parent and its hopping activities.  It uses
the algebraic determinant split and local estimates of V31--V35, but it
does not also integrate with the V34 promoted correlated measure.
Conversely, the arbitrary-hopping V30/V34 route cannot cite the exact
product factorization \eqref{eq:v4v36-polymer-gas}.  The two constructions
are alternative routes on their common domain.
\end{firewall}

\subsection{V36 result ledger and V37 target}
\label{sec:v4v36-conclusion}

V36 proves exact connectedization of local factors over the onsite quartic
product parent and bounds the complete five-species hopping/Yukawa marked
norm.  The V15 complex cluster theorem then supplies two finite-volume
zero-free logarithms, a thermodynamic residual phase pressure, and
convergent normalized local ratios through two marks.  This closes the V35
phase target inside the V28 small-hopping convergence ball.

The arbitrary-hopping strongly convex route remains open because disjoint
supports do not factor under its correlated parent.  V37 should test the
smallest proposed replacement: a block decoupling interpolation for the
promoted measure whose derivative is the exact cross-block covariance.
The V30 Witten kernel supplies one decaying bridge.  The proof must
determine whether repeated interpolation derivatives yield a tree with one
bridge per block connection without uncontrolled derivatives of the
Witten resolvent.  If they do not, the route should stay restricted to the
V36 KP domain rather than asserting complete analyticity.


\clearpage
\section{V37: the block-decoupling hierarchy stops beyond two Witten bridges}
\label{sec:v4v37}

V36 closes the phase-pressure problem when the Higgs kinetic term is a
small activity over a product reference.  The proposed extension was to
decouple blocks inside the fully correlated convex parent and use the V30
Witten kernel for every forest edge.  This note tests that proposal
directly.

The first two interpolation derivatives work: they are an expectation and
an exact covariance of local boundary interactions, and the latter decays
exponentially.  The third derivative is a genuine third cumulant.
Differentiating the Helffer--Sjostrand formula exposes both a covariance of
a resolvent density and an insertion
\((\nabla U)\cdot\nabla+D^2U\) inside the Witten inverse.  Neither is
controlled by the two-point operator norm alone.  A uniformly convex
Gaussian example shows factorial higher interpolation cumulants despite a
uniform Witten gap.  Thus the proposed one-bridge iteration is not a proof
of complete analyticity.

\subsection{A convex block interpolation}
\label{sec:v4v37-interpolation}

Let \(\mathcal B\) be a finite block partition.  Suppose the promoted
positive action has an exact decomposition
\begin{equation}
 S_{\boldsymbol s}(\phi)
 =
 \sum_{B\in\mathcal B}S_B^{\rm int}(\phi_B)
 +
 \sum_{e\in\mathcal E}s_eU_e(\phi),
 \qquad 0\le s_e\le1,
\label{eq:v4v37-action}
\end{equation}
where \(U_e\) is supported in a fixed-width neighborhood of the two blocks
incident to \(e\).  Assume, uniformly in
\(\boldsymbol s\),
\begin{equation}
 D_\phi^2S_{\boldsymbol s}\ge\kappa_{\rm d}1,
\label{eq:v4v37-uniform-convexity}
\end{equation}
and that its off-diagonal Hessian has the finite-range or exponentially
summable form admitted by V34.  For the covariant Higgs kinetic energy,
each edge term is a positive square, so multiplying it by
\(s_e\in[0,1]\) preserves its positivity.  Any good-modulus contribution
must be interpolated in a decomposition for which
\eqref{eq:v4v37-uniform-convexity} is verified; it is not automatic from
convexity at the endpoint.

Write
\[
 Z(\boldsymbol s)=\int e^{-S_{\boldsymbol s}(\phi)}\,d\phi,
 \qquad
 \mathbb E_{\boldsymbol s}F
 =
 Z(\boldsymbol s)^{-1}\int Fe^{-S_{\boldsymbol s}}\,d\phi.
\]

\begin{lemma}[Exact interpolation derivative]
\label{lem:v4v37-first-derivative}
For every differentiable observable \(F_{\boldsymbol s}\),
\begin{equation}
 \partial_{s_e}\mathbb E_{\boldsymbol s}F_{\boldsymbol s}
 =
 \mathbb E_{\boldsymbol s}\partial_{s_e}F_{\boldsymbol s}
 -
 \operatorname{Cov}_{\boldsymbol s}(F_{\boldsymbol s},U_e).
\label{eq:v4v37-expectation-derivative}
\end{equation}
In particular,
\begin{equation}
 \partial_{s_e}\log Z=-\mathbb E_{\boldsymbol s}U_e.
\label{eq:v4v37-log-first}
\end{equation}
\end{lemma}

\begin{proof}
Differentiate the numerator and denominator of the normalized finite
dimensional integral.  The derivative of \(e^{-S_{\boldsymbol s}}\) is
\(-U_ee^{-S_{\boldsymbol s}}\).  The quotient rule subtracts
\(\mathbb EF\,\mathbb EU_e\), giving the covariance.
\end{proof}

\begin{theorem}[Cumulant hierarchy of the decoupling parameters]
\label{thm:v4v37-cumulant-hierarchy}
For pairwise not necessarily distinct edges
\(e_1,\ldots,e_n\),
\begin{equation}
 \partial_{s_{e_1}}\cdots\partial_{s_{e_n}}\log Z
 =
 (-1)^n
 \kappa_{\boldsymbol s}
 (U_{e_1},\ldots,U_{e_n}),
 \qquad n\ge1.
\label{eq:v4v37-cumulant-hierarchy}
\end{equation}
In particular,
\begin{align}
 \partial_{s_e}\partial_{s_f}\log Z
 &=
 \operatorname{Cov}_{\boldsymbol s}(U_e,U_f),
\label{eq:v4v37-log-second}\\
 \partial_{s_e}\partial_{s_f}\partial_{s_g}\log Z
 &=
 -\kappa_{3,\boldsymbol s}(U_e,U_f,U_g).
\label{eq:v4v37-log-third}
\end{align}
\end{theorem}

\begin{proof}
Introduce sources \(t_e\) and use
\[
 \log Z(\boldsymbol s+\boldsymbol t)-\log Z(\boldsymbol s)
 =
 \log\mathbb E_{\boldsymbol s}
 \exp\!\left(-\sum_et_eU_e\right).
\]
By the defining power series of joint cumulants, the coefficient of
\(t_{e_1}\cdots t_{e_n}\) is
\((-1)^n\kappa(U_{e_1},\ldots,U_{e_n})\), with repeated labels handled by
ordinary differentiation.  This proves
\eqref{eq:v4v37-cumulant-hierarchy}; the two displayed cases follow.
\end{proof}

\subsection{What the V30 kernel actually proves}
\label{sec:v4v37-two-bridge}

Let \(A_e\) denote the field-variable support of
\(\nabla_\phi U_e\).  Assume
\begin{equation}
 \sup_{\boldsymbol s,e}
 \|\nabla_\phi U_e\|_{L^2(\mu_{\boldsymbol s})}
 \le M_U.
\label{eq:v4v37-edge-gradient}
\end{equation}

\begin{theorem}[Exponential two-edge interpolation response]
\label{thm:v4v37-two-edge}
There are \(C,\eta>0\), independent of volume and the block locations,
such that
\begin{equation}
 \left|
 \partial_{s_e}\partial_{s_f}\log Z
 \right|
 \le
 C M_U^2e^{-\eta d(A_e,A_f)}.
\label{eq:v4v37-two-edge}
\end{equation}
More generally, if \(F\) is local with gradient support \(A_F\), then
\begin{equation}
 \left|
 \partial_{s_e}\mathbb E_{\boldsymbol s}F
 \right|
 \le
 C
 e^{-\eta d(A_F,A_e)}
 \|\nabla_\phi F\|_{L^2(\mu_{\boldsymbol s})}M_U
\label{eq:v4v37-local-first}
\end{equation}
when \(F\) has no explicit \(s_e\)-dependence.
\end{theorem}

\begin{proof}
Uniform convexity and the uniform off-diagonal Hessian hypothesis give the
V30/V34 one-form resolvent estimate, uniformly in
\(\boldsymbol s\).  Apply the Helffer--Sjostrand covariance identity to
\((U_e,U_f)\) and insert the support projections
\(P_{A_e},P_{A_f}\).  Equation
\eqref{eq:v4v37-edge-gradient} gives
\eqref{eq:v4v37-two-edge}.  Apply the same argument to \((F,U_e)\) and use
Lemma~\ref{lem:v4v37-first-derivative} for
\eqref{eq:v4v37-local-first}.
\end{proof}

\begin{firewall}[Two-edge decay is not a block forest]
\label{fw:v4v37-two-not-all}
Theorem~\ref{thm:v4v37-two-edge} supplies one decaying carrier between two
actual boundary interactions.  A connected expansion on \(n\) blocks
requires a joint \(n\)-edge cumulant or an exact forest representation that
bounds it.  Reusing
\eqref{eq:v4v37-two-edge} independently on the edges of a proposed tree
would square or multiply a two-point estimate without deriving the
corresponding higher derivative.
\end{firewall}

\subsection{Differentiating the Witten carrier}
\label{sec:v4v37-Witten-derivative}

Let
\[
 \mathcal L_{\boldsymbol s}
 =
 -\Delta_\phi+\nabla S_{\boldsymbol s}\cdot\nabla_\phi,
 \qquad
 \mathcal L_{1,\boldsymbol s}
 =
 \mathcal L_{\boldsymbol s}\otimes1+D_\phi^2S_{\boldsymbol s}.
\]
Direct differentiation gives
\begin{equation}
 \partial_{s_e}\mathcal L_{1,\boldsymbol s}
 =
 [\nabla_\phi U_e\cdot\nabla_\phi]\otimes1
 +D_\phi^2U_e
 =:\mathcal R_e.
\label{eq:v4v37-Witten-insertion}
\end{equation}

\begin{theorem}[Exact derivative of a Witten covariance]
\label{thm:v4v37-Witten-derivative}
Let \(F,G\) have no explicit \(s_e\)-dependence and put
\[
 v_{\boldsymbol s}
 =
 \mathcal L_{1,\boldsymbol s}^{-1}\nabla_\phi G,
 \qquad
 h_{\boldsymbol s}(\phi)
 =
 \nabla_\phi F\cdot v_{\boldsymbol s}.
\]
Then
\begin{align}
 \partial_{s_e}
 \operatorname{Cov}_{\boldsymbol s}(F,G)
 &=
 -\operatorname{Cov}_{\boldsymbol s}(h_{\boldsymbol s},U_e)
 \notag\\
 &\quad-
 \left\langle
 \nabla_\phi F,
 \mathcal L_{1,\boldsymbol s}^{-1}
 \mathcal R_e
 \mathcal L_{1,\boldsymbol s}^{-1}
 \nabla_\phi G
 \right\rangle_{L^2(\mu_{\boldsymbol s})}.
\label{eq:v4v37-Witten-derivative}
\end{align}
The sum equals
\(-\kappa_{3,\boldsymbol s}(F,G,U_e)\).
\end{theorem}

\begin{proof}
Start from
\[
 \operatorname{Cov}_{\boldsymbol s}(F,G)
 =
 \int
 \nabla F\cdot
 \mathcal L_{1,\boldsymbol s}^{-1}\nabla G
 \,d\mu_{\boldsymbol s}.
\]
Differentiation of the normalized measure gives the first covariance in
\eqref{eq:v4v37-Witten-derivative}.  The inverse identity
\[
 \partial_{s_e}\mathcal L_{1,\boldsymbol s}^{-1}
 =
 -\mathcal L_{1,\boldsymbol s}^{-1}
 \mathcal R_e
 \mathcal L_{1,\boldsymbol s}^{-1}
\]
gives the second line.  Lemma~\ref{lem:v4v37-first-derivative}, applied
directly to \(\operatorname{Cov}(F,G)\), identifies their sum with the
third cumulant.
\end{proof}

\begin{assessment}[The new analytic input required at order three]
\label{ass:v4v37-order-three-input}
The operator \(\mathcal R_e\) is not merely the bounded field Hessian
\(D^2U_e\).  It contains the first-order differential operator
\((\nabla U_e)\cdot\nabla\), whose coefficient is unbounded for a quadratic
kinetic edge.  The V30 estimate controls
\(\mathcal L_1^{-1}\) on one-forms in \(L^2\), but supplies neither:
\begin{enumerate}
\item a spatial Sobolev bound for
      \(\mathcal L_1^{-1}\mathcal R_e\mathcal L_1^{-1}\); nor
\item a covariance estimate for the resolvent density
      \(h_{\boldsymbol s}\), whose field derivative itself differentiates
      the inverse.
\end{enumerate}
Thus \eqref{eq:v4v37-Witten-derivative} is an exact acquisition map for the
missing theorem, not a bound derived from V30.
\end{assessment}

\subsection{A sharp uniformly convex Gaussian test}
\label{sec:v4v37-Gaussian-test}

\begin{proposition}[Factorial interpolation cumulants with a fixed gap]
\label{prop:v4v37-Gaussian-factorial}
Let
\[
 d\mu_0(x,y)=(2\pi)^{-1}e^{-(x^2+y^2)/2}\,dx\,dy,
 \qquad
 U(x,y)=\frac12(x-y)^2,
\]
and
\[
 Z(s)=\mathbb E_{\mu_0}e^{-sU},
 \qquad s\ge0.
\]
Then
\begin{align}
 Z(s)&=(1+2s)^{-1/2},
\label{eq:v4v37-Gaussian-Z}\\
 \frac{d^n}{ds^n}\log Z(s)\bigg|_{s=0}
 &=
 (-1)^n2^{n-1}(n-1)!,
\label{eq:v4v37-Gaussian-derivatives}\\
 \kappa_n(U,\ldots,U)
 &=
 2^{n-1}(n-1)!.
\label{eq:v4v37-Gaussian-cumulants}
\end{align}
Meanwhile the interpolated action
\[
 \frac12(x^2+y^2)+\frac s2(x-y)^2
\]
has Hessian at least \(1\) for every \(s\ge0\).
\end{proposition}

\begin{proof}
The normalized difference
\(z=(x-y)/\sqrt2\) is a standard Gaussian and \(U=z^2\).  Hence
\[
 Z(s)=\mathbb E e^{-sz^2}=(1+2s)^{-1/2}.
\]
Differentiate
\[
 \log Z(s)=-\frac12\log(1+2s)
\]
at zero to obtain
\eqref{eq:v4v37-Gaussian-derivatives}.  Compare with
Theorem~\ref{thm:v4v37-cumulant-hierarchy} to get
\eqref{eq:v4v37-Gaussian-cumulants}.  The Hessian has eigenvalue one in the
constant direction and \(1+2s\) in the difference direction.
\end{proof}

\begin{corollary}[A Witten gap does not supply an exponential cumulant
hierarchy]
\label{cor:v4v37-no-gap-hierarchy}
Uniform convexity and the resulting two-point Witten bound do not by
themselves imply a bound
\[
 |\kappa_n(U,\ldots,U)|\le C^n
\]
with \(C\) independent of \(n\).  They also do not justify a block cluster
expansion at arbitrary edge strength by repeated use of a two-point
estimate.
\end{corollary}

\begin{proof}
Equation \eqref{eq:v4v37-Gaussian-cumulants} grows as
\(2^{n-1}(n-1)!\), while the Hessian gap is uniformly one.
\end{proof}

\begin{remark}[What the Gaussian test does not exclude]
\label{rem:v4v37-Gaussian-scope}
Cluster formulae contain factorial denominators, so factorial cumulants can
still be summable when another small parameter or exact forest
cancellation is present.  The proposition does not disprove complete
analyticity of every uniformly convex model.  It disproves only the
proposed inference from the V30 two-point bound to an all-orders
exponential cumulant estimate.
\end{remark}

\subsection{Decision and V38 target}
\label{sec:v4v37-conclusion}

V37 proves the exact block-interpolation cumulant hierarchy and exponential
locality of its first nontrivial two-edge response.  It then differentiates
the Witten carrier exactly and identifies the two new order-three objects.
The Gaussian test shows that a fixed convexity gap permits factorial higher
edge cumulants, so V30 cannot be iterated as one independent bridge per
forest edge.

The arbitrary-hopping phase route therefore remains open.  V38 should
replace global edge interpolation by buffered core/collar elimination.
Conditioned on a collar, separated cores factor exactly.  V30 gives
exponentially small response from a collar through a buffer into a core,
while V29 supplies one-use ancestry for the collar variables.  The target
is a single elimination step with an exact boundary kernel and a contraction
factor \(e^{-\eta L}\).  It must keep the first collar layer in the retained
skeleton, since V29 proves that blocking cannot make that layer weak.


\clearpage
\section{V38: exact buffered core elimination for the finite-range Higgs carrier}
\label{sec:v4v38}

V37 shows that global edge interpolation does not turn the V30 two-point
kernel into an all-orders forest.  The upstream replacement is to use
conditional independence where it is exact.  This note performs one
buffered elimination step for the finite-range positive Higgs carrier.
Separated interiors factor after the retained skeleton is fixed.  Uniform
convexity survives conditioning and marginalization, while the Witten
kernel makes the response of a deep-core observable to skeleton data
exponentially small in the collar width.

The result concerns the positive Higgs action before the nonlocal fermion
determinant is integrated.  Fermion lines joining different cores remain a
separate, exponentially decaying activity.

\subsection{Core, collar, and skeleton geometry}
\label{sec:v4v38-geometry}

Assume the positive Higgs action is a sum of terms of interaction range at
most \(R_0\):
\begin{equation}
 S_B(\phi)=\sum_{Y:\operatorname{diam}Y\le R_0}S_{B,Y}(\phi_Y),
 \qquad
 D_\phi^2S_B\ge\kappa1.
\label{eq:v4v38-finite-range-action}
\end{equation}
Choose disjoint block interiors \(I_1,\ldots,I_N\) whose mutual distance is
greater than \(R_0\), and put
\[
 \Sigma=\Lambda\setminus\bigcup_{j=1}^NI_j.
\]
The set \(\Sigma\) is the retained skeleton.  For a width \(L>R_0\), define
the deep core
\[
 I_j^\circ(L)=\{x\in I_j:d(x,\Sigma)\ge L\}
\]
and the interior collar \(K_j(L)=I_j\setminus I_j^\circ(L)\).
For fixed skeleton field \(\eta=\phi_\Sigma\), let
\(\mu_I^\eta\) be the conditional Higgs law on \(I=\bigcup_jI_j\).

\begin{theorem}[Exact conditional product of cores]
\label{thm:v4v38-conditional-product}
For every \(\eta\),
\begin{equation}
 \mu_I^\eta=\bigotimes_{j=1}^N\mu_{I_j}^\eta.
\label{eq:v4v38-conditional-product}
\end{equation}
Each conditional action has Hessian at least \(\kappa1\) on its interior
variables.
\end{theorem}

\begin{proof}
No interaction set \(Y\) in \eqref{eq:v4v38-finite-range-action} meets two
distinct interiors.  After \(\eta\) is fixed, the action is a sum of one
function of \(\phi_{I_j}\) for each \(j\), plus a constant.  The normalized
conditional density factors.  Restricting the full Hessian inequality to
directions supported in \(I_j\) gives the convexity bound.
\end{proof}

\begin{firewall}[Why the promoted determinant is not inside this factor]
\label{fw:v4v38-determinant-outside}
The V34 good determinant has an exponentially local but nonzero kernel at
arbitrary separation.  If its modulus were inserted into
\eqref{eq:v4v38-finite-range-action}, the exact product
\eqref{eq:v4v38-conditional-product} would fail.  V38 applies the split to
the finite-range Higgs carrier and leaves all fermion determinant
connections in their V32--V36 activity ancestry.
\end{firewall}

\subsection{Conditional Witten decay}
\label{sec:v4v38-conditional-Witten}

Let \(\mathcal L_{1,I_j}^{\eta}\) be the Witten one-form operator of
\(\mu_{I_j}^{\eta}\), with Dirichlet restriction to interior variables.

\begin{lemma}[Uniform conditional one-form kernel]
\label{lem:v4v38-conditional-Witten}
There are \(C,\gamma>0\), independent of \(j,\eta,L\), and total volume,
such that
\begin{equation}
 \|P_A(\mathcal L_{1,I_j}^{\eta})^{-1}P_C\|
 \le C e^{-\gamma d(A,C)}
\label{eq:v4v38-conditional-Witten}
\end{equation}
for \(A,C\subset I_j\).
\end{lemma}

\begin{proof}
The conditional Hessian has the same lower bound \(\kappa\), interaction
range \(R_0\), and off-diagonal row bound as the full action.  The spatial
weight in V30 acts only on the retained interior one-form indices.
Consequently the V30 commutator estimate and Neumann series apply with the
same constants, uniformly in the boundary value \(\eta\).
\end{proof}

\subsection{Exponential boundary-to-core response}
\label{sec:v4v38-boundary-response}

Let \(F\) depend only on fields in \(A\subset I_j^\circ(L)\).  A skeleton
direction \(h\) changes the conditional action through
\[
 G_h(\phi_{I_j})=D_\eta S_{B,I_j}^\eta[h].
\]
Finite range implies that \(\nabla_{\phi_{I_j}}G_h\) is supported in
\[
 \partial_{R_0}I_j=\{x\in I_j:d(x,\Sigma)\le R_0\}.
\]
Assume the local stencil bound
\begin{equation}
 \|\nabla_{\phi_{I_j}}G_h\|_{L^2(\mu_{I_j}^\eta)}
 \le J_\partial\|h\|_{\ell^2(\Sigma)}.
\label{eq:v4v38-boundary-stencil}
\end{equation}

\begin{theorem}[Buffered conditional influence]
\label{thm:v4v38-buffered-influence}
If \(F\) has no explicit dependence on \(\eta\), then
\begin{equation}
 |D_\eta\mathbb E_{\mu_{I_j}^\eta}F[h]|
 \le CJ_\partial e^{-\gamma(L-R_0)}
 \|\nabla_{\phi_{I_j}}F\|_{L^2(\mu_{I_j}^\eta)}
 \|h\|_{\ell^2(\Sigma)}.
\label{eq:v4v38-buffered-influence}
\end{equation}
\end{theorem}

\begin{proof}
Differentiating the normalized conditional integral gives
\[
 D_\eta\mathbb E^\eta F[h]
 =-\operatorname{Cov}_{\mu_{I_j}^\eta}(F,G_h).
\]
Use the conditional Helffer--Sjostrand identity and insert the gradient
support projections \(P_A\) and \(P_{\partial_{R_0}I_j}\).
Their distance is at least \(L-R_0\).  Apply
Lemma~\ref{lem:v4v38-conditional-Witten},
\eqref{eq:v4v38-boundary-stencil}, and Cauchy--Schwarz.
\end{proof}

\begin{corollary}[A strict core response for a wide buffer]
\label{cor:v4v38-strict-core}
For every \(0<\theta<1\), a width
\[
 L\ge R_0+\gamma^{-1}\log(CJ_\partial/\theta)
\]
makes the normalized boundary-to-deep-core response at most \(\theta\).
The width is independent of the number of blocks and total volume.
\end{corollary}

\begin{proof}
Insert the stated value of \(L\) in
\eqref{eq:v4v38-buffered-influence}.
\end{proof}

\subsection{The exact skeleton marginal}
\label{sec:v4v38-skeleton}

Define
\begin{equation}
 e^{-S_{\rm eff}(\eta)}
 =\int e^{-S_B(\phi_I,\eta)}\,d\phi_I.
\label{eq:v4v38-effective-action}
\end{equation}
By Theorem~\ref{thm:v4v38-conditional-product},
\[
 S_{\rm eff}(\eta)=S_\Sigma(\eta)-\sum_{j=1}^N\log Z_{I_j}(\eta),
\]
with one conditional normalization per eliminated core.

\begin{theorem}[Strong convexity survives core elimination]
\label{thm:v4v38-marginal-convexity}
The skeleton marginal satisfies
\begin{equation}
 D_\eta^2S_{\rm eff}\ge\kappa1.
\label{eq:v4v38-marginal-convexity}
\end{equation}
Its exact Hessian is
\begin{equation}
 D_\eta^2S_{\rm eff}
 =
 \mathbb E^\eta D_\eta^2S_B
 -
 \operatorname{Cov}_{\mu_I^\eta}
 (D_\eta S_B,D_\eta S_B),
\label{eq:v4v38-effective-Hessian}
\end{equation}
with polarization understood.
\end{theorem}

\begin{proof}
Write
\[
 S_B(\phi_I,\eta)=
 \frac\kappa2(|\phi_I|^2+|\eta|^2)+V(\phi_I,\eta),
\]
where \(V\) is convex.  Then
\[
 e^{-S_{\rm eff}(\eta)}
 =
 e^{-\kappa|\eta|^2/2}
 \int e^{-\kappa|\phi_I|^2/2-V(\phi_I,\eta)}\,d\phi_I.
\]
The integrand after the first factor is jointly log-concave.
Prekopa's theorem says its \(\phi_I\)-marginal is log-concave.  Hence
\(S_{\rm eff}(\eta)-\kappa|\eta|^2/2\) is convex, proving
\eqref{eq:v4v38-marginal-convexity}.  Differentiating
\eqref{eq:v4v38-effective-action} twice gives
\eqref{eq:v4v38-effective-Hessian}.
\end{proof}

\begin{firewall}[The covariance subtraction is already used]
\label{fw:v4v38-effective-Hessian-once}
The second term in \eqref{eq:v4v38-effective-Hessian} is the one exact
response generated by eliminating the core.  It may not be added again as
a separate boundary polymer or as a second Schur correction.
\end{firewall}

\subsection{One-use ancestry and the collar obstruction}
\label{sec:v4v38-ancestry}

\begin{theorem}[One-use buffered elimination]
\label{thm:v4v38-one-use}
After one elimination step:
\begin{enumerate}
\item every deep-core factor is integrated in exactly one conditional law
      \(\mu_{I_j}^\eta\);
\item every response from that factor to the skeleton crosses its retained
      collar once and receives the factor in
      \eqref{eq:v4v38-buffered-influence};
\item the conditional normalizations occur once in \(S_{\rm eff}\); and
\item the first collar layer remains an active skeleton variable.
\end{enumerate}
No block-cardinality factor appears in a response rooted at one fixed deep
core.
\end{theorem}

\begin{proof}
Items one and three are the exact disintegration
\eqref{eq:v4v38-conditional-product}--%
\eqref{eq:v4v38-effective-action}.  Item two is
Theorem~\ref{thm:v4v38-buffered-influence}; a rooted response belongs to one
conditional factor, so no sum over unrelated cores is present.  The action
couples the boundary layer directly to \(\eta\), and no positive separation
exists there.  It must therefore remain in the skeleton.
\end{proof}

\begin{firewall}[Blocking does not weaken the first collar layer]
\label{fw:v4v38-first-collar}
Corollary~\ref{cor:v4v38-strict-core} applies only after a positive buffer.
For an observable in \(\partial_{R_0}I_j\), the exponential factor is
order one.  This is the exact V29 boundary obstruction in the present
correlated-parent setting; the strict core response cannot be quoted as a
strict response of the entire block.
\end{firewall}

\subsection{V38 result ledger and V39 target}
\label{sec:v4v38-conclusion}

V38 proves exact conditional factorization of separated interiors for the
finite-range positive Higgs carrier, uniform conditional Witten decay, and
an exponentially small skeleton response for observables behind a buffer.
It also proves that strong convexity survives exact core integration and
records one-use ancestry for the conditional normalization and response.
The retained first collar layer remains strongly coupled.

V39 should attach the local Yukawa Grassmann factors before fermion
integration.  Factors wholly inside distinct deep cores inherit the exact
Higgs conditional product, while hard fermion covariance lines between
cores are exponentially small by V16.  The target is a typed connected tree
whose Higgs edges are exact conditional-core links and whose fermion edges
are Gram-controlled covariance links.  Skeleton and collar factors must
remain in a separate retained species.  The next compulsory companion audit
is V40.


\clearpage
\section{V39: deep-core Yukawa supervertices and a typed conditional tree}
\label{sec:v4v39}

V38 gives exact conditional independence of separated Higgs interiors but
leaves the fermions outside.  This note attaches the V31 local Yukawa
Grassmann factors in the part of each interior protected by a buffer.
Higgs integration produces one finite Grassmann supervertex per core.
Different cores are then connected only by hard fermion covariance lines
at fixed skeleton.  The V32 Gram forest controls those lines without a
factorial.  Differentiation with respect to the retained skeleton creates a
different edge: an exact conditional Higgs covariance, paid by the V38
buffer factor.

The resulting tree has two edge types with separate ancestry.  It closes
the deep-core conditional gas and its first skeleton response.  Collar and
skeleton Yukawa factors, and the second skeleton response, remain open.

\subsection{Finite central supports}
\label{sec:v4v39-supports}

Use the V38 interiors \(I_j\) and skeleton \(\Sigma\).  Choose central
Yukawa cell sets \(\mathcal D_j\) such that
\begin{equation}
 \bigcup_{x\in\mathcal D_j}X_x
 \subset I_j^\circ(L),
 \qquad
 |\mathcal D_j|\le M_L<\infty.
\label{eq:v4v39-central-cells}
\end{equation}
The number \(M_L\) depends on the chosen fixed block geometry but not on the
total volume.  Assume that distinct central supports satisfy
\begin{equation}
 d\!\left(
 \bigcup_{x\in\mathcal D_j}X_x,
 \bigcup_{y\in\mathcal D_k}X_y
 \right)
 \ge 2L+c_{\rm b}d_{\rm b}(j,k),
 \qquad j\ne k,
\label{eq:v4v39-core-separation}
\end{equation}
for a block metric \(d_{\rm b}\) and \(c_{\rm b}>0\).

Let \(w_x=e^{-\mathcal V_x}-1\) be the unintegrated Yukawa activity of V31.
At fixed skeleton field \(\eta\), define
\begin{align}
 P_j(\eta)
 &:=
 \mathbb E_{\mu_{I_j}^{\eta}}
 \prod_{x\in\mathcal D_j}(1+w_x),
\label{eq:v4v39-core-factor}\\
 U_j(\eta)&:=P_j(\eta)-1.
\label{eq:v4v39-core-activity}
\end{align}
The expectation acts only on Higgs coefficients; the Grassmann generators
remain unintegrated.  Hence \(U_j\) has zero scalar part and belongs to a
Grassmann algebra with at most \(q_L=qM_L\) generators.

\begin{theorem}[Exact conditional Higgs factorization of core vertices]
\label{thm:v4v39-core-factorization}
For every finite core set \(J\) and fixed \(\eta\),
\begin{equation}
 \mathbb E_{\mu_I^\eta}
 \prod_{j\in J}
 \prod_{x\in\mathcal D_j}(1+w_x)
 =
 \prod_{j\in J}P_j(\eta).
\label{eq:v4v39-core-factorization}
\end{equation}
Every local Yukawa factor on the left occurs in exactly one \(P_j\).
\end{theorem}

\begin{proof}
Theorem~\ref{thm:v4v38-conditional-product} gives a product of the
conditional Higgs measures on the distinct interiors.  The factors indexed
by \(\mathcal D_j\) depend only on \(\phi_{I_j}\) and the fixed
\(\eta\).  Fubini's theorem gives
\eqref{eq:v4v39-core-factorization}.  Disjointness of the interiors assigns
each central cell to one factor.
\end{proof}

\subsection{Marked supervertex norm}
\label{sec:v4v39-supervertex-norm}

Background derivatives in this subsection hold \(\eta\) fixed and act on
the local Yukawa coefficients.  For \(m=0,1,2\), put
\[
 A_{m,j}(\eta)
 =
 \sup_{\|v_k\|\le1}
 \|D_B^mU_j(\eta)[v_1,\ldots,v_m]\|_{\mathfrak G}.
\]

\begin{lemma}[One core costs one Yukawa parameter]
\label{lem:v4v39-one-core}
There are \(\epsilon_0>0\) and
\(K_L<\infty\), independent of volume and the core label, such that
\begin{equation}
 \mathbb E_{\mu_\Sigma}
 \prod_{j\in J}
 \left(
 \sum_{m=0}^2\rho^mA_{m,j}(\eta)
 \right)
 \le
 (K_L\epsilon_Y)^{|J|}
\label{eq:v4v39-core-product-bound}
\end{equation}
for distinct cores \(J\) and \(\epsilon_Y\le\epsilon_0\).
Here \(\mu_\Sigma\) is the exact V38 skeleton marginal.
\end{lemma}

\begin{proof}
Expand \(U_j\) from \eqref{eq:v4v39-core-factor}:
\[
 U_j
 =
 \sum_{\varnothing\ne E\subseteq\mathcal D_j}
 \mathbb E_{\mu_{I_j}^{\eta}}\prod_{x\in E}w_x.
\]
After one or two background derivatives, the product rule places marks on
actual local factors, with the same-activity and two-activity alternatives
of V15.  For a fixed choice of nonempty \(E_j\) in every core, positivity of
the coefficient norm, conditional Jensen, and the exact conditional
product give
\begin{align*}
 &\mathbb E_{\mu_\Sigma}
 \prod_{j\in J}
 \left\|
 \mathbb E_{\mu_{I_j}^{\eta}}
 [\hbox{marked product on }E_j]
 \right\|_{\mathfrak G}\\
 &\qquad\le
 \mathbb E_{\mu_B}
 \prod_{j\in J}\prod_{x\in E_j}
 [\hbox{corresponding local marked norms}].
\end{align*}
The all-cell exponential-moment argument of
Lemma~\ref{lem:v4v32-all-cell-product} bounds the right side by one constant
times \(\epsilon_Y\) for every selected local activity.  Since every
\(E_j\) is nonempty, at least one factor \(\epsilon_Y\) is gained per core.
Summing the at most \(2^{M_L}-1\) nonempty subsets and the polynomially many
mark placements gives \eqref{eq:v4v39-core-product-bound}, after all
\(L\)-dependent finite constants are absorbed into \(K_L\).
\end{proof}

\begin{firewall}[Finite block size precedes coupling smallness]
\label{fw:v4v39-order-limits}
The constant \(K_L\) may grow with the fixed core size.  V39 first chooses a
finite buffer width and block geometry, then requires
\(K_L\epsilon_Y\) to be small.  It does not claim a bound uniform as
\(L\to\infty\) at fixed Yukawa coupling.
\end{firewall}

\subsection{Fermion edges between core supervertices}
\label{sec:v4v39-fermion-edges}

Let \(C\) be the V16 hard fermion covariance.  Define
\begin{equation}
 c_L(j,k)
 =
 \gamma^{-2}
 \max_{\substack{x\in\mathcal D_j,\ y\in\mathcal D_k\\
                 a\in X_x,\ b\in X_y}}
 |C_{ab}|,
 \qquad j\ne k.
\label{eq:v4v39-block-majorant}
\end{equation}

\begin{lemma}[Buffered fermion row sum]
\label{lem:v4v39-fermion-row}
For every \(0<a'<a_0c_{\rm b}\), there is a finite \(C_L\) such that
\begin{equation}
 \sup_j\sum_{k\ne j}
 e^{a'd_{\rm b}(j,k)}c_L(j,k)
 \le C_Le^{-2a_0L}.
\label{eq:v4v39-fermion-row}
\end{equation}
\end{lemma}

\begin{proof}
The V16 Combes--Thomas bound gives
\[
 |C_{ab}|\le C\gamma^2e^{-a_0d(a,b)}.
\]
Use \eqref{eq:v4v39-core-separation}.  The number of blocks at block
distance \(n\) grows only polynomially, while
\(e^{-(a_0c_{\rm b}-a')n}\) is summable.  The finite internal cell choices
are absorbed into \(C_L\), leaving \(e^{-2a_0L}\).
\end{proof}

At fixed \(\eta\), define the fermion-truncated core coefficient
\begin{equation}
 \mathcal K_{\rm core}(J;\eta)
 =
 \left\langle U_j(\eta):j\in J\right\rangle_C^{\rm T}.
\label{eq:v4v39-core-cumulant}
\end{equation}

\begin{theorem}[Deep-core fermion tree]
\label{thm:v4v39-fermion-tree}
For \(m=0,1,2\), with background derivatives acting on the Yukawa
coefficients and the hard covariance as in V33,
\begin{align}
 &\mathbb E_{\mu_\Sigma}
 \sup_{\|v_k\|\le1}
 |D_B^m\mathcal K_{\rm core}(J;\eta)
   [v_1,\ldots,v_m]|
\notag\\
 &\quad\le
 |J|^{2m}(K'_L\epsilon_Y)^{|J|}
 \sum_{T\in\mathcal T_J}
 \prod_{\{j,k\}\in E(T)}\widehat c_L(j,k),
\label{eq:v4v39-fermion-tree}
\end{align}
where \(\widehat c_L\) also includes the first two marked covariance
majorants of V33.  For sufficiently small \(\epsilon_Y\), the rooted
core-polymer sum has exponential block-tree decay and is
\(O(\epsilon_Y)\), uniformly in the number of cores.
\end{theorem}

\begin{proof}
Apply the marked fermion forest
Lemma~\ref{lem:v4v33-fully-marked-forest} with one supervertex \(U_j\) per
core and local generator bound \(q_L\).  Average its product of marked
supervertex norms with
Lemma~\ref{lem:v4v39-one-core}.  The marked resolvent identities of V33
give the same exponentially decaying block majorant as
Lemma~\ref{lem:v4v39-fermion-row}, with changed constants.  Rooted-tree
summation is now identical to V32 on the block lattice.  The factors
\(|J|^{2m}\) are harmless in the resulting geometric series when
\(K'_L\epsilon_Y\) is sufficiently small.
\end{proof}

\subsection{The distinct Higgs-response edge}
\label{sec:v4v39-Higgs-edge}

Let \(h\) be a skeleton direction.  Because the local Yukawa factors have
no explicit dependence on \(\eta\),
\begin{equation}
 D_\eta U_j[h]
 =
 -\operatorname{Cov}_{\mu_{I_j}^{\eta}}
 \left(
 \prod_{x\in\mathcal D_j}(1+w_x),
 G_{j,h}
 \right),
\label{eq:v4v39-skeleton-derivative}
\end{equation}
where \(G_{j,h}=D_\eta S_{B,I_j}^{\eta}[h]\).

\begin{lemma}[Buffered Higgs mark on a core supervertex]
\label{lem:v4v39-Higgs-mark}
Let
\[
 \widehat A_k=\sum_{m=0}^2\rho^mA_{m,k}.
\]
There is \(K_L^{\rm H}<\infty\) such that, for every finite core set
\(J\) containing \(j\),
\begin{align}
 &\mathbb E_{\mu_\Sigma}
 \left[
 \|D_\eta U_j[h]\|_{\mathfrak G}
 \prod_{k\in J\setminus\{j\}}\widehat A_k
 \right]
\notag\\
 &\qquad\le
 K_L^{\rm H}\epsilon_Y
 e^{-\gamma(L-R_0)}
 \|h\|_{\ell^2(\Sigma)}
 (K_L\epsilon_Y)^{|J|-1}.
\label{eq:v4v39-Higgs-mark}
\end{align}
The covariance in \eqref{eq:v4v39-skeleton-derivative} is one actual
conditional Higgs connection.
\end{lemma}

\begin{proof}
Expand the nonconstant part of the product in
\eqref{eq:v4v39-skeleton-derivative}.  Its field gradient is supported in
the union of the central Yukawa cells, at distance at least \(L-R_0\) from
the boundary stencil.  Apply
Theorem~\ref{thm:v4v38-buffered-influence} coefficientwise in the finite
Grassmann algebra.  The gradient of a product places one field derivative
on an actual Yukawa factor.

Multiply by the displayed factors from the other cores before integrating
over \(\eta\).  Conditional Jensen and the exact product
\eqref{eq:v4v39-core-factorization} convert the resulting conditional
product into a moment under the original full Higgs law.  Cauchy--Schwarz
for the one differentiated core and the all-cell exponential-moment
argument of Lemma~\ref{lem:v4v32-all-cell-product} then give one
\(\epsilon_Y\) for every nonempty core, with a finite constant depending on
\(M_L\).  The selected core also carries the V38 spatial factor
\(e^{-\gamma(L-R_0)}\).  Summing its finitely many local subsets and mark
placements proves \eqref{eq:v4v39-Higgs-mark}.
\end{proof}

\begin{theorem}[Typed conditional tree with one skeleton mark]
\label{thm:v4v39-typed-tree}
Differentiate a fermion-connected coefficient
\(\mathcal K_{\rm core}(J;\eta)\) once in a skeleton direction.  Every term
has:
\begin{enumerate}
\item a fermion spanning tree on \(J\), with each edge paid by
      \(\widehat c_L\); and
\item one Higgs-response edge from an actual core vertex to the skeleton,
      paid by \(K_L^{\rm H}\epsilon_Ye^{-\gamma(L-R_0)}\).
\end{enumerate}
After averaging over the skeleton, the rooted sum of these typed trees is
finite for small \(K'_L\epsilon_Y\), and its skeleton-response norm is
\begin{equation}
 O\!\left(
 \epsilon_Ye^{-\gamma(L-R_0)}
 \right).
\label{eq:v4v39-typed-tree-smallness}
\end{equation}
\end{theorem}

\begin{proof}
The fermion cumulant is multilinear in the supervertices.  Its skeleton
derivative therefore selects one actual \(U_j\) and replaces it by
\(D_\eta U_j[h]\); the covariance \(C\) is independent of the Higgs
skeleton variable.  Apply the V33 fermion forest to the remaining
multilinear expression.  Lemma~\ref{lem:v4v39-Higgs-mark} pays the selected
Higgs edge, while Lemma~\ref{lem:v4v39-one-core} pays every other
supervertex.  Lemma~\ref{lem:v4v39-fermion-row} and rooted-tree summation
finish the estimate.  The two edge types arise from different exact
derivatives and are never interchanged.
\end{proof}

\begin{firewall}[The typed tree has one Higgs edge, not an all-orders
skeleton forest]
\label{fw:v4v39-one-Higgs-edge}
Theorem~\ref{thm:v4v39-typed-tree} proves the first response of a
fermion-connected deep-core coefficient to the retained skeleton.  A second
skeleton derivative differentiates a conditional covariance and introduces
the V37 third-cumulant/Witten-insertion hierarchy.  V39 does not obtain that
second response by squaring
\eqref{eq:v4v39-Higgs-mark}.
\end{firewall}

\begin{firewall}[Collar factors remain retained]
\label{fw:v4v39-collar-retained}
Yukawa cells outside \(\mathcal D_j\), including those in the first collar
layer, are not included in \(U_j\).  Their response to the skeleton is
order one by V29/V38.  They must be retained as skeleton/collar activities
and cannot be assigned the deep-core factor
\(e^{-\gamma(L-R_0)}\).
\end{firewall}

\subsection{V39 result ledger and compulsory V40 target}
\label{sec:v4v39-conclusion}

V39 integrates all central Yukawa factors into one exact finite Grassmann
supervertex per Higgs core.  At least one local Yukawa parameter is gained
per nonempty core.  Gram-controlled hard fermion lines give a connected
core tree with an additional separation factor, and one skeleton
derivative adds a distinct conditional-Higgs edge carrying the buffer
decay.  Determinant and conditional-normalization ancestries remain
single.

V40 is a compulsory companion audit.  It should then decide how to treat
the retained collar species.  The first candidate is a colored elimination:
integrate one family of mutually separated collars while keeping the other
colors fixed, then rotate colors.  The audit must determine whether this
creates a genuine second skeleton response or merely cycles the V37
third-cumulant obstruction.  No deep-core estimate may be charged to a
first-layer collar factor.


\clearpage
\section{V40: companion audit and the exact angle missing from colored elimination}
\label{sec:v4v40}

V39 closes the deep-core conditional Yukawa tree but leaves the first-layer
collar in the skeleton.  A proposed next step is to color the collar blocks,
integrate one separated color at a time, and rotate colors.  This compulsory
audit shows why the norm of each conditional expectation is not enough.
The relevant quantity is the angle between successive conditional
subspaces.  A two-block Gaussian calculation gives that angle exactly,
proves a strict contraction under a relative Schur margin, and shows that a
fixed lower Hessian bound alone does not make the contraction uniform.

\subsection{Compulsory companion audit}
\label{sec:v4v40-audit}

The exact active companion snapshots inspected were
\begin{enumerate}
\item
\texttt{Renewal\_Yang\_Mills\_Mass\_Gap\_Research\_Notes\_V58.tex},
\(1051969\) bytes, SHA--256
\[
 \texttt{EF9BFBC3142C38B4A0EA64F8274D3C32A6960FCB791EEBD5359F027B94B5B536}.
\]
Its newest theorem identifies the cut-projective operator norm that
preserves product-state capacity through a separator.  It proves a sharp
counterexample: an ordinary norm-one unitary can amplify product capacity
by a representation-dimension factor.  The physical separator still needs
a projective bound or an exact commutator-square estimate.
\item
\texttt{Renewal\_NS\_Stability\_Research\_Notes\_V31.tex},
\(887537\) bytes, SHA--256
\[
 \texttt{F1121B62238AD5491BB7CF03635EA9934942EDF573ED8FAAA9EE79E1D38A6696}.
\]
Its terminal result still shows that heat lifting and a polarization hinge
do not create the missing critical first moment.  Any gain must retain the
actual stopping-time-correlated signed nonlinear source.
\end{enumerate}

The transfer is structural.  YM V58 distinguishes ordinary contraction
from contraction in the norm compatible with the desired tensor cut.  The
colored Higgs step similarly needs a strict angle between conditional
subspaces; the fact that each conditional expectation has operator norm
one is insufficient.  NS V31 again forbids manufacturing the missing gain
after the correlated object has been split into positive pieces.

\begin{firewall}[Companion type boundary]
\label{fw:v4v40-companion-types}
No Wilson capacity, projective tensor norm, Navier--Stokes heat estimate, or
helical identity is used as a Higgs inequality.  The audit selects the
quantity that must be computed before colored conditional maps are
composed.
\end{firewall}

\subsection{Conditional expectations as projections}
\label{sec:v4v40-projections}

Let \(\mu\) be a positive probability measure on variables
\((x,y)\).  On \(L^2(\mu)\), define
\[
 P_xF=\mathbb E[F\mid x],
 \qquad
 P_yF=\mathbb E[F\mid y],
 \qquad
 P_0F=\mathbb EF.
\]

\begin{lemma}[Projection facts]
\label{lem:v4v40-projections}
The maps \(P_x,P_y,P_0\) are orthogonal projections,
\begin{equation}
 \|P_x\|=\|P_y\|=1,
 \qquad
 P_xP_0=P_yP_0=P_0.
\label{eq:v4v40-projection-norms}
\end{equation}
The round-trip operator on centered \(y\)-observables is
\begin{equation}
 Q_y=(P_yP_xP_y)|_{L^2_0(\sigma(y))}.
\label{eq:v4v40-round-trip}
\end{equation}
A colored two-step elimination is a strict \(L^2\) contraction on that
sector exactly when
\begin{equation}
 q:=\|Q_y\|<1.
\label{eq:v4v40-angle}
\end{equation}
\end{lemma}

\begin{proof}
Conditional expectation is the orthogonal projection onto the closed
subspace of functions measurable with respect to the conditioning sigma
field.  Constants belong to both subspaces, proving
\eqref{eq:v4v40-projection-norms}.  Starting with a \(y\)-measurable
centered function, one color update applies \(P_x\), and the return update
applies \(P_y\).  Inserting the identity \(P_yF=F\) gives
\eqref{eq:v4v40-round-trip}.  Its norm is the definition of strict
round-trip contraction.
\end{proof}

\begin{firewall}[Norm-one color maps do not provide the sweep gap]
\label{fw:v4v40-norm-one}
Submultiplicativity gives only
\(\|P_yP_xP_y\|\le1\).  A bound \(q<1\) requires the relative position of
the two conditional subspaces after constants are removed.  It cannot be
obtained by multiplying the two norm-one bounds.
\end{firewall}

\subsection{The exact Gaussian angle}
\label{sec:v4v40-Gaussian-angle}

Let
\[
 H=
 \begin{pmatrix}A&B\\B^*&D\end{pmatrix}>0
\]
be a finite-dimensional real Gaussian precision matrix, and let
\(\mu_H\) be the centered Gaussian law proportional to
\(\exp[-\frac12\langle(x,y),H(x,y)\rangle]\).
Put
\begin{equation}
 R=A^{-1/2}BD^{-1/2},
 \qquad
 \rho=\|R\|<1.
\label{eq:v4v40-canonical-correlation}
\end{equation}

\begin{theorem}[Gaussian conditional angle]
\label{thm:v4v40-Gaussian-angle}
For the Gaussian law,
\begin{equation}
 \|P_xP_y-P_0\|=\rho,
 \qquad
 \|Q_y\|=\rho^2.
\label{eq:v4v40-Gaussian-angle}
\end{equation}
On linear observables the round trip is represented, after whitening, by
\(R^*R\).
\end{theorem}

\begin{proof}
Completing the square gives
\[
 \mathbb E[x\mid y]=-A^{-1}By,
 \qquad
 \mathbb E[y\mid x]=-D^{-1}B^*x.
\]
After whitening the conditional linear maps are \(-R\) and \(-R^*\);
their round trip is \(R^*R\).

For the full \(L^2\) statement, use the orthogonal Hermite-chaos
decomposition of a Gaussian space.  Conditional expectation maps the
\(n\)-th homogeneous chaos by the \(n\)-fold symmetric tensor power of the
linear canonical-correlation map.  Its norm is \(\rho^n\).  Constants form
the zeroth chaos, and the largest norm on the positive chaoses is attained
at \(n=1\), namely \(\rho\).  The round trip has norm \(\rho^2\).
Positivity of the Schur complements gives \(\rho<1\).
\end{proof}

\begin{theorem}[A relative Schur margin gives a quantitative sweep]
\label{thm:v4v40-Schur-margin}
If
\begin{equation}
 H\ge\kappa1,
 \qquad
 A\le K1,
\label{eq:v4v40-two-sided-curvature}
\end{equation}
then
\begin{equation}
 \rho^2\le1-\frac{\kappa}{K},
 \qquad
 q\le1-\frac{\kappa}{K}<1.
\label{eq:v4v40-Schur-contraction}
\end{equation}
The same conclusion follows with \(D\le K1\).
\end{theorem}

\begin{proof}
For any \(u\), test the lower Hessian bound on
\[
 (u,-D^{-1}B^*u).
\]
This gives
\[
 \langle u,(A-BD^{-1}B^*)u\rangle
 \ge\kappa|u|^2.
\]
Conjugating by \(A^{-1/2}\) yields
\[
 RR^*
 =
 A^{-1/2}BD^{-1}B^*A^{-1/2}
 \le
 1-\kappa A^{-1}
 \le
 \left(1-\frac{\kappa}{K}\right)1.
\]
Take norms and use
Theorem~\ref{thm:v4v40-Gaussian-angle}.  The argument with \(A,D\)
interchanged is identical.
\end{proof}

\subsection{Why lower convexity alone is insufficient}
\label{sec:v4v40-sharpness}

\begin{proposition}[Fixed gap with an arbitrarily slow color sweep]
\label{prop:v4v40-slow-sweep}
For \(M>0\), let
\begin{equation}
 H_M=
 \begin{pmatrix}M+1&-M\\-M&M+1\end{pmatrix}.
\label{eq:v4v40-slow-precision}
\end{equation}
Then \(H_M\ge1\), but
\begin{equation}
 \rho_M=\frac{M}{M+1},
 \qquad
 q_M=\left(\frac{M}{M+1}\right)^2
 \longrightarrow1.
\label{eq:v4v40-slow-angle}
\end{equation}
Thus no \(q<1\) depending only on the lower convexity constant can control
colored elimination.
\end{proposition}

\begin{proof}
The eigenvalues of \(H_M\) are \(1\) and \(2M+1\), so its lower gap is one.
Here \(A=D=M+1\) and \(B=-M\), making
\(\rho=|B|/\sqrt{AD}=M/(M+1)\).  Apply
Theorem~\ref{thm:v4v40-Gaussian-angle}.
\end{proof}

\begin{firewall}[This does not contradict V30 spatial decay]
\label{fw:v4v40-no-contradiction}
The V30 decay exponent depends on the off-diagonal row strength relative to
the convexity gap.  In \eqref{eq:v4v40-slow-precision}, that strength grows
like \(M\), so the available decay rate degenerates as
\(q_M\to1\).  The example rules out dependence on \(\kappa\) alone, not a
bound using the complete curvature data.
\end{firewall}

\subsection{Colored collar decision}
\label{sec:v4v40-decision}

Let \(P_1,\ldots,P_r\) be the conditional projections for a finite coloring
of collar blocks, and let \(P_0\) project onto the variables retained after
a full sweep.  Define the centered sweep
\begin{equation}
 \mathcal Q=(P_r\cdots P_1-P_0)|_{L^2_0}.
\label{eq:v4v40-colored-sweep}
\end{equation}

\begin{theorem}[Sufficient colored-sweep interface]
\label{thm:v4v40-sweep-interface}
If
\begin{equation}
 \|\mathcal Q\|\le q<1
\label{eq:v4v40-sweep-gap}
\end{equation}
uniformly in volume, boundary data, and the good-field chart, then:
\begin{enumerate}
\item the response propagated through repeated color sweeps is bounded by
      the convergent resolvent
      \((1-\mathcal Q)^{-1}=\sum_{n\ge0}\mathcal Q^n\);
\item every propagated source is used once per sweep step; and
\item the second skeleton response can be organized by the exact
      same-block second derivative and two distinct first-derivative
      placements before applying the geometric bound.
\end{enumerate}
The response constant is at most \((1-q)^{-1}\) times the one-sweep marked
constant.
\end{theorem}

\begin{proof}
The norm hypothesis makes the Neumann series converge on the centered
subspace.  Expanding it records the exact number and order of sweep
applications.  Differentiate the finite product
\(P_r\cdots P_1\) twice before summing the series.  The ordinary product rule
gives one twice-differentiated conditional kernel or two actual
once-differentiated kernels, exactly as in V15.  Taking norms only after
these placements are assembled gives the stated factor.
\end{proof}

\begin{assessment}[What remains to verify for the nonlinear collar]
\label{ass:v4v40-nonlinear-gap}
The V38 lower Hessian bound proves that every color conditional is
well-defined and each \(P_c\) is an orthogonal projection.  It does not
prove \eqref{eq:v4v40-sweep-gap}.  A sufficient acquisition is a uniform
relative curvature bound on the good-field chart, analogous to
\eqref{eq:v4v40-two-sided-curvature}, in a conditional norm compatible with
the collar/skeleton cut.  The quartic Hessian is unbounded globally, so this
must be combined with a smooth good-curvature/tail partition.  Every failed
curvature cell must receive its own V31/V34 exponential payer.
\end{assessment}

\begin{firewall}[The sweep interface is conditional]
\label{fw:v4v40-sweep-open}
Theorem~\ref{thm:v4v40-sweep-interface} is a sufficient interface, not a
proof that the nonlinear Higgs collar satisfies it.  The Gaussian
calculation identifies the missing angle exactly; substituting the
individual norms \(\|P_c\|=1\) would repeat the error excluded by
\ref{fw:v4v40-norm-one}.
\end{firewall}

\subsection{V40 result ledger and V41 target}
\label{sec:v4v40-conclusion}

V40 records the compulsory YM/NS audit, identifies colored conditional
expectations as orthogonal projections, and proves that their relative
angle is the required sweep parameter.  For a Gaussian two-block action the
round-trip constant is the square of the canonical correlation.  A
two-sided relative curvature bound yields a strict quantitative
contraction, while a fixed lower Hessian bound alone allows the contraction
to approach one.  The conditional sweep theorem then shows exactly how a
proved angle would organize the second response.

V41 should construct a smooth good-curvature/tail partition for the
nonlinear Higgs collar.  On the good chart, bound both the diagonal block
curvature and the relative cross block Hessian to obtain a uniform
conditional angle.  On the complement, use the simultaneous V34 tail
estimate with one payer per failed collar cell.  The first task is to prove
that the cutoff derivatives introduced by this partition do not erase the
sweep gap or duplicate the V34 Yukawa tail.


\clearpage
\section{V41: a nonlinear two-color sweep gap from conditional Poincare}
\label{sec:v4v41}

V40 proposed a good-curvature cutoff because the quartic Higgs Hessian is
unbounded above.  That proposal is stronger than necessary.  The angle
between two conditional subspaces can be bounded by differentiating a
conditional expectation and applying Poincare twice: once in the
conditioned color and once in its marginal.  The bound uses the lower
convexity and the mixed color Hessian, but no upper bound on the onsite
Hessian.  Since the quartic contribution is onsite, it does not enter the
mixed Hessian at all.

This yields a uniform strict even/odd sweep for the positive bipartite Higgs
carrier, including arbitrary hopping strength as long as the lower
convexity and mixed-stencil norm stay fixed.  It does not absorb the
fermion determinant or prove a complex all-orders phase theorem.

\subsection{A nonlinear maximal-correlation theorem}
\label{sec:v4v41-maximal-correlation}

Let
\[
 d\mu(x,y)=Z^{-1}e^{-S(x,y)}\,dx\,dy
\]
on finite-dimensional Euclidean spaces \(X\oplus Y\).  Assume
\begin{equation}
 D^2S\ge\kappa1,
 \qquad
 \sup_{x,y}\|D_xD_yS(x,y)\|\le J.
\label{eq:v4v41-hypotheses}
\end{equation}
Let \(\mu_X\) be the \(x\)-marginal and \(\mu^x_Y\) the conditional
\(y\)-law.  For a centered \(g\in L^2(\mu_Y)\), define
\[
 Tg(x)=\mathbb E[g(Y)\mid X=x].
\]

\begin{theorem}[Nonlinear conditional-angle bound]
\label{thm:v4v41-angle}
Under \eqref{eq:v4v41-hypotheses},
\begin{equation}
 \|T\|_{L^2_0(\mu_Y)\to L^2_0(\mu_X)}^2
 \le
 \frac{J^2}{\kappa^2+J^2}
 =:q_{\kappa,J}<1.
\label{eq:v4v41-angle}
\end{equation}
Equivalently, with the conditional projections of V40,
\begin{equation}
 \|(P_yP_xP_y)|_{L^2_0(\sigma(y))}\|
 \le q_{\kappa,J}.
\label{eq:v4v41-round-trip}
\end{equation}
The constants are independent of the dimensions of \(X,Y\).
\end{theorem}

\begin{proof}
It is enough first to take smooth \(g\).  Put
\[
 h(x)=Tg(x)=\int g(y)\,d\mu_Y^x(y).
\]
Differentiating the normalized conditional integral in a direction
\(u\in X\) gives
\begin{equation}
 u\cdot\nabla_xh(x)
 =
 -\operatorname{Cov}_{\mu_Y^x}
 \left(g,u\cdot\nabla_xS\right).
\label{eq:v4v41-conditional-derivative}
\end{equation}
The conditional \(y\)-action has Hessian at least \(\kappa1\).  Its
Poincare inequality and Cauchy--Schwarz imply
\begin{align*}
 |u\cdot\nabla_xh(x)|^2
 &\le
 \operatorname{Var}_{\mu_Y^x}(g)
 \operatorname{Var}_{\mu_Y^x}
 (u\cdot\nabla_xS)\\
 &\le
 \frac1\kappa
 \operatorname{Var}_{\mu_Y^x}(g)
 \mathbb E_{\mu_Y^x}
 \left|
 \nabla_y(u\cdot\nabla_xS)
 \right|^2\\
 &\le
 \frac{J^2}{\kappa}|u|^2
 \operatorname{Var}_{\mu_Y^x}(g).
\end{align*}
Taking the supremum over unit \(u\) yields
\begin{equation}
 |\nabla_xh(x)|^2
 \le
 \frac{J^2}{\kappa}
 \operatorname{Var}_{\mu_Y^x}(g).
\label{eq:v4v41-gradient-bound}
\end{equation}

By the strong-convexity marginal theorem proved in V38, the \(x\)-marginal
has Hessian at least \(\kappa1\).  Apply its Poincare inequality and
\eqref{eq:v4v41-gradient-bound}:
\begin{equation}
 \operatorname{Var}_{\mu_X}(h)
 \le
 \frac{J^2}{\kappa^2}
 \mathbb E_{\mu_X}
 \operatorname{Var}_{\mu_Y^x}(g).
\label{eq:v4v41-pre-angle}
\end{equation}
The conditional variance identity is
\[
 \operatorname{Var}_{\mu_Y}(g)
 =
 \mathbb E_{\mu_X}\operatorname{Var}_{\mu_Y^x}(g)
 +
 \operatorname{Var}_{\mu_X}(h).
\]
Write \(r=J^2/\kappa^2\).  Substitution into
\eqref{eq:v4v41-pre-angle} gives
\[
 (1+r)\operatorname{Var}_{\mu_X}(h)
 \le r\operatorname{Var}_{\mu_Y}(g),
\]
which is \eqref{eq:v4v41-angle}.  Smooth functions are dense in the
relevant \(L^2\) spaces.  Finally, the round trip on \(y\)-observables is
\(T^*T\), so its norm is \(\|T\|^2\).
\end{proof}

\begin{remark}[Relation to the exact Gaussian answer]
\label{rem:v4v41-Gaussian-comparison}
For the Gaussian model of V40, the exact round-trip norm is the squared
canonical correlation.  Theorem~\ref{thm:v4v41-angle} gives a possibly
weaker but dimension-free upper bound using only \(\kappa\) and \(J\).
For
\[
 H_M=\begin{pmatrix}M+1&-M\\-M&M+1\end{pmatrix},
\]
it gives \(q\le M^2/(1+M^2)\), which tends to one consistently with the
exact value \([M/(M+1)]^2\).
\end{remark}

\begin{assessment}[Nonduplication relative to V27 and V40]
\label{ass:v4v41-nonduplication}
V27's Dobrushin criterion sums single-block Lipschitz influences and
requires a row sum below one.  V40 computes the exact Gaussian conditional
angle and obtains a sufficient bound from two-sided curvature.  V41 proves
a different nonlinear \(L^2\) statement: the conditional variance identity
converts a finite mixed-Hessian bound into a strict maximal-correlation
constant even when the onsite upper Hessian is unbounded.
\end{assessment}

\subsection{Uniformity under exterior conditioning}
\label{sec:v4v41-conditioning}

Let \(z\) be an exterior skeleton variable and consider a conditional action
\(S_z(x,y)\).  Suppose the full positive carrier has
\(D^2S\ge\kappa1\), and the \(x\)-\(y\) mixed Hessian has norm at most \(J\).

\begin{corollary}[Boundary-uniform angle]
\label{cor:v4v41-boundary-uniform}
For every fixed \(z\), the conditional two-color round trip satisfies
\begin{equation}
 \|Q_y^z\|\le q_{\kappa,J}<1
\label{eq:v4v41-boundary-angle}
\end{equation}
with the same constant, independently of the boundary value, block size,
and total volume.
\end{corollary}

\begin{proof}
Fixing \(z\) restricts the Hessian to the principal \(x\oplus y\) block, so
the lower bound remains \(\kappa\).  It changes the conditional action by
terms that do not enlarge the prescribed \(x\)-\(y\) mixed Hessian norm.
Apply Theorem~\ref{thm:v4v41-angle}.
\end{proof}

\subsection{Application to the quartic Higgs carrier}
\label{sec:v4v41-Higgs}

Assume the physical finite-range Higgs stencil is bipartite after the V38
block choice.  Write its two retained colors as \(E,O\).  The Hessian
\eqref{eq:v4v27-action-Hessian} has the block form
\[
 D_\phi^2S_{\rm H}
 =
 \begin{pmatrix}
 A_E(\phi_E)&K_{EO}(B)\\
 K_{OE}(B)&A_O(\phi_O)
 \end{pmatrix}.
\]
The quartic Hessian contributes only to \(A_E,A_O\).  Put
\begin{equation}
 J_{\rm H}:=
 \sup_B\|K_{EO}(B)\|<\infty.
\label{eq:v4v41-Higgs-mixed}
\end{equation}

\begin{theorem}[Strict nonlinear Higgs color sweep]
\label{thm:v4v41-Higgs-sweep}
Under the V27 lower bound \(D^2S_{\rm H}\ge\kappa1\) and
\eqref{eq:v4v41-Higgs-mixed}, the even/odd conditional sweep, with arbitrary
exterior skeleton fixed, satisfies
\begin{equation}
 \boxed{
 \|Q_{\rm H}\|
 \le
 \frac{J_{\rm H}^2}{\kappa^2+J_{\rm H}^2}
 <1.}
\label{eq:v4v41-Higgs-sweep}
\end{equation}
No upper bound on \(|\phi|\), no good-curvature cutoff, and no small-hopping
hypothesis is required.
\end{theorem}

\begin{proof}
The quartic potential and all onsite mass/source terms have zero mixed
\(E\)-\(O\) Hessian.  The only mixed block is the bounded covariant kinetic
stencil \(K_{EO}\).  Apply
Corollary~\ref{cor:v4v41-boundary-uniform} with \(J=J_{\rm H}\).
\end{proof}

\begin{firewall}[Bipartite means the complete positive stencil]
\label{fw:v4v41-bipartite}
Theorem~\ref{thm:v4v41-Higgs-sweep} applies when every interaction in the
positive Higgs carrier lies either onsite or between the two colors.  An
improved action with same-color or longer-range terms requires a finite
multicolor extension; those terms may not be dropped to force a
checkerboard form.
\end{firewall}

\begin{firewall}[The determinant remains outside]
\label{fw:v4v41-determinant-outside}
The theorem concerns the positive finite-range Higgs carrier.  The good
fermion determinant modulus of V34 is nonlocal after fermion integration.
It is not inserted into \(S_{\rm H}\) here.  The V39 Grassmann
supervertices and fermion edges remain external activities.
\end{firewall}

\subsection{Second response of the positive collar sweep}
\label{sec:v4v41-second-response}

Let
\[
 \mathcal Q_{\rm H}=P_OP_E-P_0
\]
on the centered two-color sector, with the external skeleton fixed.
Theorem~\ref{thm:v4v41-Higgs-sweep} gives
\begin{equation}
 \|(1-\mathcal Q_{\rm H})^{-1}\|
 \le
 \frac1{1-q_{\kappa,J_{\rm H}}}
 =
 1+\frac{J_{\rm H}^2}{\kappa^2}.
\label{eq:v4v41-sweep-resolvent}
\end{equation}

\begin{corollary}[Conditional closure of the twice-marked collar response]
\label{cor:v4v41-collar-response}
Suppose the one-sweep conditional kernels and their first two background
derivatives have the local \(L^2\) score bounds established in
V30--V33.  Then the repeated even/odd positive-carrier response has a
volume-uniform second derivative.  Its exact numerator consists of:
\begin{enumerate}
\item one actual second derivative of one color-conditional kernel; or
\item two actual first derivatives at their ordered sweep positions.
\end{enumerate}
The geometric propagation costs at most the factor in
\eqref{eq:v4v41-sweep-resolvent}.
\end{corollary}

\begin{proof}
Theorem~\ref{thm:v4v41-Higgs-sweep} verifies the strict-gap hypothesis of
Theorem~\ref{thm:v4v40-sweep-interface}.  That theorem differentiates one
finite sweep before summing the Neumann series and retains the two
exhaustive mark placements.  Insert
\eqref{eq:v4v41-sweep-resolvent}.
\end{proof}

\begin{firewall}[No good-curvature tail is created]
\label{fw:v4v41-no-curvature-tail}
V41 withdraws the proposed good-curvature/tail split of V40 for the
positive quartic carrier.  Introducing it would create cutoff derivatives
and a second exceptional-field ancestry without necessity.  The V34 tail
remains solely the large-field complement of the Yukawa good-determinant
split.
\end{firewall}

\subsection{V41 result ledger and V42 target}
\label{sec:v4v41-conclusion}

V41 proves a dimension-free nonlinear conditional-angle theorem from two
Poincare inequalities and the conditional variance identity.  Applied to
the bipartite quartic Higgs action, it gives a strict even/odd sweep with
constant \(J_{\rm H}^2/(\kappa^2+J_{\rm H}^2)\), despite the unbounded
onsite quartic Hessian and without small hopping.  Together with the V40
sweep interface, this closes the positive collar's conditional second
response under the existing local marked-score bounds.

V42 should compose this sweep with the V39 typed core tree.  The required
object is the conditional expectation channel acting on a finite
Grassmann coefficient space.  The scalar \(L^2\) angle controls each
coefficient, but the full Grassmann norm is a projective sum over
coefficients.  Following the V40 audit lesson, V42 must prove that the
finite core fibre costs only a fixed \(q_L\)-dependent constant and does
not grow with the number of cores or sweep steps.


\clearpage
\section{V42: finite Grassmann fibres do not degrade the color-sweep gap}
\label{sec:v4v42}

V41 proves a strict scalar \(L^2\) angle for the two-color positive Higgs
carrier.  V39 core activities take values in a finite Grassmann algebra.
Estimating every coefficient separately in the \(\ell^1\) algebra norm
would appear to introduce a fibre factor at each sweep.  This note proves
that such repeated loss is artificial.  In the Hilbert coefficient norm,
conditional expectation is the tensor product of the scalar channel with
the identity, so its centered contraction is unchanged.  Conversion to the
submultiplicative Grassmann norm is needed only at actual interaction
vertices.

\subsection{Two coefficient norms}
\label{sec:v4v42-coefficient-norms}

Let \(\mathfrak G_q\) have the ordered monomial basis
\(\{\xi_I:I\subseteq\{1,\ldots,q\}\}\), of dimension
\begin{equation}
 N_q=2^q.
\label{eq:v4v42-fibre-dimension}
\end{equation}
For \(A=\sum_Ia_I\xi_I\), define
\[
 \|A\|_1=\sum_I|a_I|,
 \qquad
 \|A\|_2=\left(\sum_I|a_I|^2\right)^{1/2}.
\]

\begin{lemma}[Finite-fibre norm interface]
\label{lem:v4v42-norm-interface}
For all \(A,B\in\mathfrak G_q\),
\begin{align}
 \|A\|_2&\le\|A\|_1\le\sqrt{N_q}\|A\|_2,
\label{eq:v4v42-norm-equivalence}\\
 \|AB\|_2&\le\sqrt{N_q}\|A\|_2\|B\|_2.
\label{eq:v4v42-Hilbert-product}
\end{align}
Extraction of any fixed Grassmann coefficient has operator norm one in
\(\|\cdot\|_2\).
\end{lemma}

\begin{proof}
The first line is the ordinary comparison of \(\ell^2\) and \(\ell^1\) on
\(N_q\) coefficients.  By V31 submultiplicativity,
\[
 \|AB\|_2\le\|AB\|_1
 \le\|A\|_1\|B\|_1
 \le N_q\|A\|_2\|B\|_2.
\]
This gives a valid constant \(N_q\).  A sharper use of convolution on the
finite group of subsets, followed by Young's inequality, gives
\[
 \|AB\|_2\le\|A\|_1\|B\|_2
 \le\sqrt{N_q}\|A\|_2\|B\|_2,
\]
where the Grassmann signs and zero products can only reduce the absolute
convolution.  A coefficient extraction is an inner product with one
orthonormal basis vector.
\end{proof}

\subsection{Vector-valued conditional expectation}
\label{sec:v4v42-vector-channel}

For a probability law \(\mu\), identify
\[
 L^2(\mu;\mathfrak G_q,\|\cdot\|_2)
 =
 L^2(\mu)\widehat\otimes\mathbb C^{N_q}.
\]
If \(T:L^2(\mu_Y)\to L^2(\mu_X)\) is the V41 conditional channel, define
\(T_q=T\otimes1_{\mathbb C^{N_q}}\).

\begin{theorem}[The fibre preserves the maximal-correlation constant]
\label{thm:v4v42-fibre-angle}
On the coefficientwise centered subspace,
\begin{equation}
 \|T_q\|=\|T\|,
 \qquad
 \|T_q^*T_q\|=\|T\|^2
 \le q_{\kappa,J}.
\label{eq:v4v42-fibre-angle}
\end{equation}
Consequently
\begin{equation}
 \|(1-T_q^*T_q)^{-1}\|
 \le\frac1{1-q_{\kappa,J}}
 =
 1+\frac{J^2}{\kappa^2},
\label{eq:v4v42-fibre-resolvent}
\end{equation}
with no dependence on \(q\).
\end{theorem}

\begin{proof}
Write \(F=\sum_If_I\xi_I\).  Conditional expectation acts coefficientwise:
\[
 T_qF=\sum_I(Tf_I)\xi_I.
\]
Therefore
\[
 \|T_qF\|_{L^2(\mu_X;\ell^2)}^2
 =
 \sum_I\|Tf_I\|_{L^2(\mu_X)}^2
 \le
 \|T\|^2\sum_I\|f_I\|_{L^2(\mu_Y)}^2.
\]
Taking a vector with one nonzero coefficient gives the reverse norm
inequality.  The adjoint and round-trip statements follow from Hilbert
tensor products.  Apply Theorem~\ref{thm:v4v41-angle} and sum the Neumann
series for \eqref{eq:v4v42-fibre-resolvent}.
\end{proof}

\begin{firewall}[Why the \(\ell^1\) estimate must not be iterated]
\label{fw:v4v42-no-fibre-per-sweep}
Applying
\(\|A\|_1\le\sqrt{N_q}\|A\|_2\) before every color update would produce
\((\sqrt{N_q})^n\) after \(n\) sweeps and would destroy the Neumann series.
Theorem~\ref{thm:v4v42-fibre-angle} propagates in the Hilbert coefficient
space for the entire sweep history.  Norm conversion occurs only when an
actual Grassmann product or final coefficient extraction is performed.
\end{firewall}

\subsection{Insertion of a finite core supervertex}
\label{sec:v4v42-core-insertion}

For a V39 core, \(q=q_L\) is fixed once the buffer geometry is fixed.  Let
\(\widehat U_j=U_j-\mathbb E_{\mu_\Sigma}U_j\), with centering performed
coefficientwise.

\begin{lemma}[One fibre constant per actual core vertex]
\label{lem:v4v42-one-fibre-cost}
Consider a multilinear expression containing \(n\) distinct core
supervertices and any number of color-sweep propagators between them.
When it is estimated by:
\begin{enumerate}
\item the Hilbert coefficient norm along every sweep segment;
\item inequality~\eqref{eq:v4v42-Hilbert-product} at every actual Grassmann
      multiplication; and
\item coefficient extraction only at the end,
\end{enumerate}
its total fibre factor is at most
\begin{equation}
 (C_{q_L})^n,
\label{eq:v4v42-one-fibre-cost}
\end{equation}
independently of the number of sweeps and total number of cores.
\end{lemma}

\begin{proof}
A sweep segment has norm bounded by
\eqref{eq:v4v42-fibre-resolvent} and carries no fibre factor.  A tree with
\(n\) vertices has at most \(n-1\) binary multiplications.  Each costs at
most \(\sqrt{N_{q_L}}\) by
\eqref{eq:v4v42-Hilbert-product}.  The last coefficient extraction costs
one.  Absorb
\((\sqrt{N_{q_L}})^{n-1}\) into \(C_{q_L}^n\).
\end{proof}

The reference in item two above means
Lemma~\ref{lem:v4v42-norm-interface}, specifically
\eqref{eq:v4v42-Hilbert-product}.

\begin{theorem}[Core-tree stability through the nonlinear sweep]
\label{thm:v4v42-core-sweep}
Assume the V39 core-tree hypotheses and the bipartite V41 Higgs carrier.
Then every fermion-connected core coefficient can be propagated through the
complete even/odd collar sweep with:
\begin{enumerate}
\item the same fermion spanning-tree factors as
      \eqref{eq:v4v39-fermion-tree};
\item the scalar sweep factor
      \(1+J_{\rm H}^2/\kappa^2\);
\item one finite factor \(C_{q_L}\) per actual core supervertex; and
\item no factor depending on the number of sweep iterations.
\end{enumerate}
For sufficiently small
\begin{equation}
 C_{q_L}K'_L\epsilon_Y,
\label{eq:v4v42-effective-smallness}
\end{equation}
the rooted typed-tree sum remains convergent.
\end{theorem}

\begin{proof}
Center every coefficient function before it enters a repeated sweep.  The
constant part is assigned to the effective skeleton activity once.  Apply
Theorem~\ref{thm:v4v42-fibre-angle} to the centered part along the entire
sweep history, and use
\eqref{eq:v4v42-fibre-resolvent}.  The fermion forest and Higgs expectation
are still those of V39; only their coefficient carrier norm has changed.
Lemma~\ref{lem:v4v42-one-fibre-cost} supplies the vertexwise finite factor.
It changes \(K'_L\) to \(C_{q_L}K'_L\), so the same rooted-tree proof
converges under \eqref{eq:v4v42-effective-smallness}.
\end{proof}

\subsection{Marked sweep placements}
\label{sec:v4v42-marked}

Let a background derivative act on a local Yukawa coefficient, the hard
fermion covariance, or one color-conditional Higgs kernel.  Through second
order, differentiate the finite core tree and one finite color sweep before
expanding the sweep resolvent.

\begin{theorem}[No fibre amplification of the quadratic mark]
\label{thm:v4v42-two-marks}
Under the marked hypotheses of V33, V39, and
Corollary~\ref{cor:v4v41-collar-response}, the propagated core tree has a
volume-uniform second background response.  Its exhaustive placements are:
\begin{enumerate}
\item one actual second derivative on a Yukawa supervertex, fermion
      covariance line, or conditional Higgs kernel; and
\item two actual first derivatives on two ordered occurrences among those
      same objects.
\end{enumerate}
The fibre constants remain vertexwise as in
\eqref{eq:v4v42-one-fibre-cost}; no additional \(N_{q_L}\) is generated by
the mark positions or sweep length.
\end{theorem}

\begin{proof}
The ordinary product rule gives the two stated classes before any norm is
taken.  V33 controls covariance-line marks, V39 controls core-vertex and
one typed Higgs-edge marks, and V41 controls their propagation through the
positive sweep.  Every conditional map continues to act diagonally on the
fixed monomial basis after differentiation.  Apply the Hilbert tensor norm
through sweep segments and charge
Lemma~\ref{lem:v4v42-norm-interface} only at actual multiplications.
The polynomial number of mark placements is absorbed by the convergent
tree series, exactly as in V33.
\end{proof}

\begin{firewall}[Finite fibre is not uniform block size]
\label{fw:v4v42-fixed-fibre}
The constant \(C_{q_L}\) grows with the number of Grassmann generators in a
core.  V42 proves independence from total volume and sweep count after a
fixed block geometry is chosen.  It does not permit \(L\to\infty\) at fixed
\(\epsilon_Y\) without rechecking
\eqref{eq:v4v42-effective-smallness}.
\end{firewall}

\begin{firewall}[The complex phase remains outside the positive sweep]
\label{fw:v4v42-phase-outside}
The sweep gap is a theorem for the positive Higgs conditional law.  Complex
phase and tail coefficients may be carried as finite Grassmann-valued
activities, but they are not inserted into the measure defining
\(P_E,P_O\).  V42 therefore does not prove zero-freeness of the
arbitrary-hopping complex partition.
\end{firewall}

\subsection{V42 result ledger and V43 target}
\label{sec:v4v42-conclusion}

V42 proves that the finite Grassmann coefficient space preserves the exact
scalar maximal-correlation constant in its Hilbert norm.  The
submultiplicative coefficient norm costs one fixed factor per actual core
vertex, not per color sweep.  Consequently the V39 fermion/Higgs typed tree
and its genuine quadratic marks remain convergent through the V41 nonlinear
collar sweep for sufficiently small Yukawa coupling at fixed block geometry.

V43 should analyze the constant coefficient assigned to the effective
skeleton after each core sweep.  Its logarithm contains the vacuum and
local phase-pressure terms.  The target is an exact telescoping identity
showing that constants from successive block eliminations occur once, while
local insertions cancel all exterior constants in normalized ratios.  This
is the correlated-parent analogue of the V35 Gaussian cancellation and the
V22 determinant ancestry.


\clearpage
\section{V43: the exact normalization cocycle under successive eliminations}
\label{sec:v4v43}

V42 propagates centered finite-Grassmann coefficients through the nonlinear
Higgs sweep.  The complementary constant coefficient must be handled
without confusing three different objects: a field-independent scalar, a
Grassmann coefficient independent of the Higgs field, and a function of the
retained skeleton.  Only the first is a vacuum normalization.  This note
keeps the full elimination maps until a scalar normalization functional is
actually applied.

The resulting normalization cocycle telescopes exactly.  Every vacuum
constant occurs once, and the same product cancels identically from a
normalized local numerator.  A thermodynamic phase-pressure sum is obtained
only when the scalar cocycle has a proved local block decomposition; that
locality remains a separate analytic input in the arbitrary-hopping route.

\subsection{Normalized elimination maps}
\label{sec:v4v43-maps}

For levels \(0\le k\le J\), let \(\mathcal A_k\) be the finite-dimensional
coefficient space of functions of the level-\(k\) retained variables,
including any finite Grassmann coefficients.  Let
\begin{equation}
 \mathcal T_k:\mathcal A_k\longrightarrow\mathcal A_{k+1}
\label{eq:v4v43-elimination-map}
\end{equation}
be the exact linear map that integrates the variables eliminated at level
\(k\).  It may include the V38 conditional Higgs expectation, the V41 color
sweep, and normalized fermion contraction, but not a second copy of the
V22 quadratic determinant.  Assume
\[
 \mathcal T_k1=1.
\]
Let \(\omega_k:\mathcal A_k\to\mathbb C\) be a scalar normalization
functional with \(\omega_k(1)=1\).

Starting with \(F_0\in\mathcal A_0\), define recursively
\begin{align}
 \widetilde F_{k+1}&=\mathcal T_kF_k,
\label{eq:v4v43-raw-step}\\
 c_k&=\omega_{k+1}(\widetilde F_{k+1}),
\label{eq:v4v43-step-constant}\\
 F_{k+1}&=c_k^{-1}\widetilde F_{k+1},
\qquad\omega_{k+1}(F_{k+1})=1,
\label{eq:v4v43-normalized-step}
\end{align}
whenever \(c_k\ne0\).

\begin{theorem}[Exact normalization cocycle]
\label{thm:v4v43-cocycle}
For every \(1\le n\le J\),
\begin{equation}
 \mathcal T_{n-1}\cdots\mathcal T_0F_0
 =
 \left(\prod_{k=0}^{n-1}c_k\right)F_n.
\label{eq:v4v43-cocycle}
\end{equation}
Consequently the logarithm on any branch fixed continuously from the
zero-interaction point is
\begin{equation}
 \log C_n=\sum_{k=0}^{n-1}\log c_k,
 \qquad
 C_n:=\prod_{k=0}^{n-1}c_k.
\label{eq:v4v43-log-cocycle}
\end{equation}
No intermediate normalization appears more than once.
\end{theorem}

\begin{proof}
For \(n=1\), equations
\eqref{eq:v4v43-raw-step}--\eqref{eq:v4v43-normalized-step} give
\(\mathcal T_0F_0=c_0F_1\).  If the identity holds at \(n\), linearity of
\(\mathcal T_n\) gives
\[
 \mathcal T_n\cdots\mathcal T_0F_0
 =
 \left(\prod_{k<n}c_k\right)\mathcal T_nF_n
 =
 \left(\prod_{k\le n}c_k\right)F_{n+1}.
\]
Induction proves \eqref{eq:v4v43-cocycle}.  Products of nonzero scalars have
the additive logarithm on the branch obtained by continuous continuation
from \(c_k=1\).
\end{proof}

\begin{firewall}[A skeleton function is not a vacuum constant]
\label{fw:v4v43-skeleton-not-constant}
The conditional factor \(P_j(\eta)\) of V39 is a Grassmann polynomial whose
coefficients depend on the retained skeleton.  It remains inside
\(\widetilde F_{k+1}\) until \(\omega_{k+1}\) is applied.  Pulling
\(P_j(\eta)\) out of the skeleton integral and recording it as \(c_k\)
would be a false factorization.
\end{firewall}

\subsection{Exact cancellation in local ratios}
\label{sec:v4v43-local-ratios}

Let \(O\) be a local observable and carry its numerator through the same
unnormalized maps:
\begin{equation}
 N_n(O)=
 \mathcal T_{n-1}\cdots\mathcal T_0(OF_0).
\label{eq:v4v43-local-numerator}
\end{equation}
Normalize it using the denominator cocycle,
\[
 \widehat N_n(O)=C_n^{-1}N_n(O).
\]

\begin{theorem}[Vacuum constants cancel from every normalized numerator]
\label{thm:v4v43-local-cancellation}
At every finite level,
\begin{equation}
 \frac{
 \omega_n(N_n(O))
 }{
 \omega_n(\mathcal T_{n-1}\cdots\mathcal T_0F_0)
 }
 =
 \omega_n(\widehat N_n(O)).
\label{eq:v4v43-local-cancellation}
\end{equation}
The full product \(C_n\), including its modulus and phase, cancels exactly.
The numerator must use the constants defined by the denominator recursion;
recomputing them in the presence of \(O\) would change the ratio.
\end{theorem}

\begin{proof}
By Theorem~\ref{thm:v4v43-cocycle} and
\(\omega_n(F_n)=1\), the denominator is \(C_n\).  The numerator is
\(C_n\widehat N_n(O)\) by definition.  Apply the linear functional
\(\omega_n\) and divide by the nonzero scalar \(C_n\).
\end{proof}

\begin{corollary}[Extensive decay of a phase mean is harmless to the
algebraic ratio]
\label{cor:v4v43-extensive-phase}
The identity \eqref{eq:v4v43-local-cancellation} remains valid when
\(|C_n|\) decays exponentially with volume, as in the V35 Gaussian phase
test.  What is required at finite volume is \(C_n\ne0\), not a
volume-uniform positive lower bound.
\end{corollary}

\begin{proof}
No estimate of \(|C_n|^{-1}\) is used after the common factor is cancelled
algebraically.  The finite-volume division is legitimate whenever
\(C_n\ne0\).
\end{proof}

\subsection{Exact first and second marks of the cocycle}
\label{sec:v4v43-marks}

Suppose the maps, functionals, and constants depend twice differentiably on
a background parameter \(B\), and remain nonzero on the chosen branch.

\begin{theorem}[Marked logarithmic normalization]
\label{thm:v4v43-marked-cocycle}
For background directions \(v,w\),
\begin{align}
 D_B\log C_n[v]
 &=
 \sum_{k<n}\frac{D_Bc_k[v]}{c_k},
\label{eq:v4v43-first-log-mark}\\
 D_B^2\log C_n[v,w]
 &=
 \sum_{k<n}
 \left[
 \frac{D_B^2c_k[v,w]}{c_k}
 -
 \frac{D_Bc_k[v]D_Bc_k[w]}{c_k^2}
 \right].
\label{eq:v4v43-second-log-mark}
\end{align}
The two terms at level \(k\) are the genuine second mark of that
normalization and the two actual first marks generated by differentiating
its logarithm.
\end{theorem}

\begin{proof}
Differentiate the finite sum
\eqref{eq:v4v43-log-cocycle} once and twice.  No cross-level term appears
in the second derivative of the logarithm because the product has already
been converted to the exact sum of level logarithms.
\end{proof}

\begin{firewall}[No square of a first-mark estimate]
\label{fw:v4v43-two-marks}
The second term in \eqref{eq:v4v43-second-log-mark} is an identity involving
two derivatives of the same \(c_k\).  It does not permit
\(D_B^2c_k\) to be replaced by a square of a bound for \(D_Bc_k\).
Both terms must be controlled by the V42 marked elimination.
\end{firewall}

\subsection{Block-local scalar constants and pressure}
\label{sec:v4v43-pressure}

The following theorem states the extra locality needed for a phase pressure.
Assume that at level \(k\),
\begin{equation}
 c_k=\prod_{b\in\mathcal B_k}c_{k,b},
 \qquad
 |c_{k,b}-1|\le\delta_k<\frac12,
\label{eq:v4v43-local-constants}
\end{equation}
where each \(c_{k,b}\) is a scalar attached to one exact eliminated block,
not a skeleton-dependent function.  Let the block volumes satisfy
\(v_k\ge v_0b^{dk}\), \(b>1\), and
\[
 |\mathcal B_k|\le C|\Lambda|/v_k.
\]

\begin{theorem}[Sufficient local pressure ledger]
\label{thm:v4v43-pressure}
If
\begin{equation}
 \sum_{k\ge0}v_k^{-1}\delta_k<\infty,
\label{eq:v4v43-pressure-sum}
\end{equation}
then the scalar normalization satisfies the uniform bound
\begin{equation}
 \frac1{|\Lambda|}|\log C_J|
 \le
 2C\sum_{k<J}v_k^{-1}\delta_k.
\label{eq:v4v43-pressure-bound}
\end{equation}
For translation-covariant bulk blocks along a van Hove sequence, the limit
of \(|\Lambda|^{-1}\log C_J\) exists after \(J\to\infty\), up to the
separately controlled terminal remainder.  Its imaginary part is the
vacuum phase-pressure contribution of the cocycle.
\end{theorem}

\begin{proof}
The disk \(|z-1|<1/2\) is zero-free and carries the logarithm continued from
one.  There
\[
 |\log z|\le2|z-1|.
\]
Use \eqref{eq:v4v43-local-constants} and sum over blocks:
\[
 |\log C_J|
 \le
 \sum_{k<J}\sum_{b\in\mathcal B_k}|\log c_{k,b}|
 \le
 2C|\Lambda|\sum_{k<J}v_k^{-1}\delta_k.
\]
This proves \eqref{eq:v4v43-pressure-bound}.  Under translation covariance,
root each bulk block at one lattice site.  Boundary blocks have
vanishing density along a van Hove sequence, and the absolutely summable
majorant \eqref{eq:v4v43-pressure-sum} permits dominated convergence over
levels.
\end{proof}

\begin{remark}[Geometric block density pays bounded level constants]
\label{rem:v4v43-geometric-density}
If \(\delta_k\le\delta<1/2\), then
\(\sum_kv_k^{-1}\delta_k<\infty\) solely from geometric growth of \(v_k\).
The local constant need not decay with scale for its vacuum pressure to be
finite.  This statement concerns the pressure density; convergence of a
local insertion's residual ancestry is a separate estimate.
\end{remark}

\begin{firewall}[The local scalar factorization is not yet proved]
\label{fw:v4v43-locality-open}
The V38 conditional product gives block factors at fixed skeleton, while
V41--V42 control their centered propagation.  They do not yet prove the
scalar product \eqref{eq:v4v43-local-constants} after the retained skeleton
is integrated.  Thus Theorem~\ref{thm:v4v43-pressure} is a sufficient
pressure ledger, not a claim that the arbitrary-hopping phase pressure is
closed.
\end{firewall}

\subsection{Relation to determinant ancestry}
\label{sec:v4v43-ancestry}

\begin{theorem}[Distinct normalization ancestries]
\label{thm:v4v43-distinct-ancestry}
The scalar cocycle \(C_n\) records interaction and phase constants produced
by normalized Higgs/Yukawa elimination.  The V22 product
\[
 \prod_k\det(A_k)_{I_kI_k}
\]
records the quadratic Gaussian Schur normalization.  With the normalized
maps \(\mathcal T_k\) defined after the Gaussian parent is divided out, the
two products are disjoint and each occurs once.
\end{theorem}

\begin{proof}
The V22 determinants arise when the quadratic covariance is factorized,
before normalized interaction expectations are taken.  By definition,
\(\mathcal T_k1=1\); hence it contains no remaining Gaussian normalization.
The constants \(c_k\) are obtained by applying \(\omega_{k+1}\) to the
interaction factor \(\mathcal T_kF_k\).  These operations act on disjoint
factors of the exact normalized integral.
\end{proof}

\subsection{V43 result ledger and V44 target}
\label{sec:v4v43-conclusion}

V43 constructs the exact scalar normalization cocycle for arbitrary
successive elimination maps, proves its one-use telescope, and proves exact
cancellation of its modulus and phase from every normalized local ratio.
It also records the genuine first and second logarithmic marks and separates
the interaction cocycle from the V22 Gaussian determinant ancestry.  Under
an explicit block-local scalar factorization, geometric block density gives
a convergent vacuum pressure ledger.

V44 should prove quasi-locality of the scalar normalization functional
generated by one V41 sweep.  The target is a comparison of two boundary
conditions that agree near an eliminated block: their local scalar
constants should differ by the V38 exponential buffer factor followed by
the V41 sweep resolvent.  Such a comparison would turn the conditional
block factors into local scalar activities and verify the missing input
\eqref{eq:v4v43-local-constants} without assuming skeleton independence.


\clearpage
\section{V44: buffered quasi-locality of one scalar normalization}
\label{sec:v4v44}

V43 isolates the scalar normalization cocycle but leaves open the locality
needed for its pressure.  There are two different questions.  First, does a
normalization produced in one eliminated core depend only weakly on remote
skeleton data?  Second, does the skeleton expectation of the product of all
such normalizations factor into scalar block constants?  The first statement
is a boundary-response theorem.  The second is a connected-correlation
problem and does not follow from the first.

This note proves the first statement in a form compatible with the V38
conditional Witten kernel and the V41 nonlinear color sweep.  On a zero-free
local chart, the logarithm of one conditional normalization has an exact
shell telescoping expansion whose increments decay exponentially.  The only
cost of resolving repeated even/odd sweeps is
\((1-q_{\rm H})^{-1}=1+J_{\rm H}^2/\kappa^2\).  The final firewall identifies
the remaining multi-block cumulant problem precisely.

\subsection{One-core scalar and its zero-free chart}
\label{sec:v4v44-one-core}

Use the V38 geometry.  Fix one interior (I), a deep set
(A\subset I^\circ(L)), and the conditional positive Higgs law
(\mu_I^\eta), where (\eta) is the retained skeleton field.  Let
(U(\phi_I)) be a complex (C^1) scalar insertion supported in (A), and
put
\begin{equation}
 Z_U(\eta):=\mathbb E_{\mu_I^\eta}(1+U).
\label{eq:v4v44-local-normalization}
\end{equation}
This is a scalar only after the conditional expectation has been taken.  It
is still a function of (\eta), so it is not yet one of the constants
(c_{k,b}) in \eqref{eq:v4v43-local-constants}.

\begin{lemma}[A uniform local zero-free disk]
\label{lem:v4v44-zero-free-disk}
If
\begin{equation}
 \|U\|_{L^\infty(\mu_I^\eta)}\le\epsilon<1
\label{eq:v4v44-small-insertion}
\end{equation}
uniformly on a connected boundary chart (\mathfrak C), then
\begin{equation}
 |Z_U(\eta)-1|\le\epsilon,
 \qquad
 |Z_U(\eta)|\ge\rho:=1-\epsilon>0
\label{eq:v4v44-zero-free-bound}
\end{equation}
on (\mathfrak C).  Consequently (\log Z_U) has the unique branch obtained
by continuation from (U=0).
\end{lemma}

\begin{proof}
Normalized expectation is a contraction on (L^\infty), so

\[
 |Z_U(\eta)-1|
 =|\mathbb E_{\mu_I^\eta}U|
 \le\epsilon.
\]

The disk (|z-1|<1) omits zero and is simply connected.  The lower bound and
the logarithmic branch follow.
\end{proof}

\begin{firewall}[Coefficient smallness must precede the logarithm]
\label{fw:v4v44-zero-free-input}
The finite Grassmann norm of V42 controls a scalar projection only when the
chosen scalar functional is contractive in that norm.  Lemma
\ref{lem:v4v44-zero-free-disk} therefore applies after such a projection has
been specified and \eqref{eq:v4v44-small-insertion} has been verified.  A
small Grassmann polynomial cannot be called a nonzero complex number before
this step.
\end{firewall}

\subsection{Direct buffered logarithmic response}
\label{sec:v4v44-direct-response}

For a skeleton direction (h), retain the V38 score

\[
 G_h=D_\eta S_{B,I}^\eta[h].
\]

Assume the boundary stencil estimate
\eqref{eq:v4v38-boundary-stencil}.  The insertion (U) has no explicit
dependence on (\eta).

\begin{theorem}[Buffered logarithmic boundary derivative]
\label{thm:v4v44-log-derivative}
On the chart of Lemma~\ref{lem:v4v44-zero-free-disk},
\begin{equation}
 \left|D_\eta\log Z_U(\eta)[h]\right|
 \le
 \frac{CJ_\partial}{\rho}
 e^{-\gamma(L-R_0)}
 \|\nabla_{\phi_I}U\|_{L^2(\mu_I^\eta)}
 \|h\|_{\ell^2(\Sigma)}.
\label{eq:v4v44-log-derivative}
\end{equation}
The constants (C,\gamma) are the uniform constants of V38 and are
independent of the total volume and number of other cores.
\end{theorem}

\begin{proof}
Differentiating the normalized conditional integral gives

\[
 D_\eta Z_U(\eta)[h]
 =-\operatorname{Cov}_{\mu_I^\eta}(U,G_h).
\]

The constant part of (1+U) has zero covariance.  Apply Theorem
\ref{thm:v4v38-buffered-influence} to the real and imaginary parts of (U),
or equivalently use the complexification of the same Hilbert-space bound.
This gives

\[
 |D_\eta Z_U(\eta)[h]|
 \le CJ_\partial e^{-\gamma(L-R_0)}
 \|\nabla_{\phi_I}U\|_{L^2(\mu_I^\eta)}
 \|h\|_{\ell^2(\Sigma)}.
\]

Since (D\log Z_U=(DZ_U)/Z_U), divide by the lower bound
(|Z_U|\ge\rho) in \eqref{eq:v4v44-zero-free-bound}.
\end{proof}

\begin{corollary}[Finite comparison along a boundary segment]
\label{cor:v4v44-finite-comparison}
Let (\eta_t=\eta_0+t(\eta_1-\eta_0)), (0\le t\le1), remain in
(\mathfrak C).  If

\[
 M_U:=\sup_{0\le t\le1}
 \|\nabla_{\phi_I}U\|_{L^2(\mu_I^{\eta_t})}<\infty,
\]

then, on the continued logarithmic branch,
\begin{equation}
 |\log Z_U(\eta_1)-\log Z_U(\eta_0)|
 \le
 \frac{CJ_\partial M_U}{\rho}
 e^{-\gamma(L-R_0)}
 \|\eta_1-\eta_0\|_{\ell^2(\Sigma)}.
\label{eq:v4v44-finite-comparison}
\end{equation}
\end{corollary}

\begin{proof}
Integrate \eqref{eq:v4v44-log-derivative} along the segment (\eta_t).
The branch is continuous because the entire segment is zero-free.
\end{proof}

\begin{remark}[Why a boundary norm is unavoidable]
\label{rem:v4v44-boundary-norm}
The Higgs skeleton is unbounded.  Agreement of two boundary conditions near
(I) cannot by itself give a uniform deterministic comparison independent
of the size of their remote values.  Equation
\eqref{eq:v4v44-finite-comparison} records the necessary weighted boundary
size.  Uniform moment estimates will replace this size by an integrable
majorant below.
\end{remark}

\subsection{The exact sweep-resolvent cost}
\label{sec:v4v44-sweep}

The preceding derivative is the seed response across one buffer.  When the
retained collar is represented by repeated even/odd conditional updates, the
seed can recur only through the centered V41 sweep.  The following abstract
statement separates the exact response identity from its estimate.

Let (\mathcal H_\eta) be the centered coefficient Hilbert space used in
V42, and let (\mathcal Q_\eta:\mathcal H_\eta\to\mathcal H_\eta) be one
complete positive-carrier color round trip.  Assume
\begin{equation}
 \|\mathcal Q_\eta\|\le q_{\rm H}
 :=\frac{J_{\rm H}^2}{\kappa^2+J_{\rm H}^2}<1.
\label{eq:v4v44-sweep-contraction}
\end{equation}
For a boundary direction (h), let (r_{\eta,h}\in\mathcal H_\eta) be the
one-use marked seed and let (\ell_\eta\in\mathcal H_\eta^*) be the scalar
readout.  Suppose the differentiated finite sweep, followed by its convergent
limit, gives the exact identity
\begin{equation}
 D_\eta Z_U(\eta)[h]
 =\ell_\eta(1-\mathcal Q_\eta)^{-1}r_{\eta,h}.
\label{eq:v4v44-response-identity}
\end{equation}

\begin{theorem}[A buffered seed survives all color sweeps]
\label{thm:v4v44-sweep-resolved}
If
\begin{equation}
 \|\ell_\eta\|\le L_U,
 \qquad
 \|r_{\eta,h}\|\le
 B_Ue^{-\gamma(L-R_0)}\|h\|_{\ell^2(\Sigma)},
\label{eq:v4v44-seed-bound}
\end{equation}
then
\begin{equation}
 |D_\eta\log Z_U(\eta)[h]|
 \le
 \frac{L_UB_U}{\rho}
 \left(1+\frac{J_{\rm H}^2}{\kappa^2}\right)
 e^{-\gamma(L-R_0)}
 \|h\|_{\ell^2(\Sigma)}.
\label{eq:v4v44-sweep-log-bound}
\end{equation}
No coefficient-fibre dimension and no number of sweep iterations occurs in
this estimate.
\end{theorem}

\begin{proof}
The Neumann series converges in operator norm by
\eqref{eq:v4v44-sweep-contraction}, and

\[
 \|(1-\mathcal Q_\eta)^{-1}\|
 \le\sum_{m\ge0}q_{\rm H}^m
 =\frac1{1-q_{\rm H}}
 =1+\frac{J_{\rm H}^2}{\kappa^2}.
\]

Insert this in \eqref{eq:v4v44-response-identity}, use
\eqref{eq:v4v44-seed-bound}, and divide by
(|Z_U(\eta)|\ge\rho).  V42 shows that the conditional maps act diagonally
on the finite Grassmann coefficient fibre, so
\eqref{eq:v4v44-sweep-contraction} holds there with the same scalar norm.
\end{proof}

\begin{assessment}[Acquisition map for the sweep identity]
\label{ass:v4v44-sweep-acquisition}
Equation \eqref{eq:v4v44-response-identity} is obtained by differentiating
the finite sequence of conditional color kernels before summing the sweep
history.  Each differentiated kernel supplies one actual score insertion;
the undifferentiated intervals between its creation and the scalar readout
are powers of (\mathcal Q_\eta).  Corollary
\ref{cor:v4v41-collar-response} proves the same placement rule through two
marks.  Thus \eqref{eq:v4v44-response-identity} is the one-mark specialization
of the established interface, provided the local score is the V38 buffered
seed.  It is not an additional independence assumption.
\end{assessment}

\begin{firewall}[The buffer factor is charged once]
\label{fw:v4v44-buffer-once}
The exponential in \eqref{eq:v4v44-seed-bound} is paid when the marked source
first crosses from the remote boundary into the deep core.  Subsequent color
sweeps are bounded by powers of (q_{\rm H}).  Attaching a second V38 factor
to every power of (\mathcal Q_\eta) would duplicate the same geometric
separation.
\end{firewall}

\subsection{Exact shell telescoping of the local logarithm}
\label{sec:v4v44-shells}

Fix a reference skeleton field (\bar\eta).  Let (B_r(I)) be the skeleton
sites at graph distance at most (r) from (I), and define the truncated
boundary
\begin{equation}
 \eta^{[r]}_x=
 \begin{cases}
  \eta_x,&x\in B_r(I),\\
  \bar\eta_x,&x\notin B_r(I).
 \end{cases}
\label{eq:v4v44-truncated-boundary}
\end{equation}
For a finite skeleton choose (R_\Lambda) with
(B_{R_\Lambda}(I)=\Sigma).  Suppose every segment between
(\eta^{[r-1]}) and (\eta^{[r]}) stays in the zero-free chart.  Set
\begin{align}
 a_{I,0}(\eta)&:=\log Z_U(\eta^{[0]}),
\label{eq:v4v44-shell-zero}\\
 a_{I,r}(\eta)&:=
 \log Z_U(\eta^{[r]})-
 \log Z_U(\eta^{[r-1]}),\qquad r\ge1.
\label{eq:v4v44-shell-increment}
\end{align}

\begin{theorem}[Quasi-local logarithmic activity]
\label{thm:v4v44-quasilocal-activity}
Assume the sweep-resolved hypotheses of Theorem
\ref{thm:v4v44-sweep-resolved} uniformly for shell perturbations.  Then
\begin{equation}
 \log Z_U(\eta)
 =\sum_{r=0}^{R_\Lambda}a_{I,r}(\eta)
\label{eq:v4v44-shell-telescope}
\end{equation}
exactly, and for (r>R_0),
\begin{equation}
 |a_{I,r}(\eta)|
 \le
 K_U e^{-\gamma(r-R_0)}
 \|\eta-\bar\eta\|_{\ell^2(S_r(I))},
\label{eq:v4v44-shell-decay}
\end{equation}
where (S_r(I)=B_r(I)\setminus B_{r-1}(I)) and
\begin{equation}
 K_U:=\frac{L_UB_U}{\rho}
 \left(1+\frac{J_{\rm H}^2}{\kappa^2}\right).
\label{eq:v4v44-KU}
\end{equation}
Consequently
\begin{equation}
 Z_U(\eta)
 =\exp(a_{I,0}(\eta))
 \prod_{r=1}^{R_\Lambda}\exp(a_{I,r}(\eta))
\label{eq:v4v44-multiplicative-shells}
\end{equation}
is an exact multiplicative decomposition into increasingly supported
quasi-local scalar activities.
\end{theorem}

\begin{proof}
Equations \eqref{eq:v4v44-shell-zero}--%
\eqref{eq:v4v44-shell-increment} telescope, and
(\eta^{[R_\Lambda]}=\eta), proving
\eqref{eq:v4v44-shell-telescope}.  The two boundaries in the (r)-th
increment differ only on (S_r(I)).  Interpolate between them.  Their marked
seed must cross distance at least (r-R_0) before reaching the supported
insertion.  Apply \eqref{eq:v4v44-sweep-log-bound} along that segment with
(h=(\eta-\bar\eta)1_{S_r(I)}) and integrate.  This proves
\eqref{eq:v4v44-shell-decay}.  Exponentiating the exact logarithmic telescope
on its continued branch gives \eqref{eq:v4v44-multiplicative-shells}.
\end{proof}

\begin{corollary}[Absolute shell summability in the positive skeleton law]
\label{cor:v4v44-shell-summability}
Let (\nu_\Sigma) be any skeleton law for which, uniformly in volume,
\begin{equation}
 \mathbb E_{\nu_\Sigma}
 \|\eta-\bar\eta\|_{\ell^2(S_r(I))}
 \le M(1+r)^p
\label{eq:v4v44-shell-moment}
\end{equation}
for fixed (M,p<\infty).  Then
\begin{equation}
 \sum_{r>R_0}
 \mathbb E_{\nu_\Sigma}|a_{I,r}|
 \le
 K_UM\sum_{r>R_0}e^{-\gamma(r-R_0)}(1+r)^p
 <\infty
\label{eq:v4v44-absolute-shell-sum}
\end{equation}
uniformly in the total volume.
\end{corollary}

\begin{proof}
Take expectations in \eqref{eq:v4v44-shell-decay}, use
\eqref{eq:v4v44-shell-moment}, and sum the exponentially weighted
polynomial.  For a finite-dimensional lattice field with a uniform second
moment per site, \eqref{eq:v4v44-shell-moment} follows from Cauchy--Schwarz
and polynomial shell growth, with (p=(d-1)/2) enlarged if necessary.
\end{proof}

\begin{remark}[Thermodynamic stabilization of one block]
\label{rem:v4v44-one-block-limit}
The summable majorant in \eqref{eq:v4v44-absolute-shell-sum} implies that the
logarithm attached to one fixed block stabilizes in (L^1) as the remote
boundary recedes, provided the finite-volume skeleton laws are locally
consistent.  This is the locality property required of an individual block
activity.  It is stronger than a pointwise statement on bounded fields and
is compatible with the unbounded positive Higgs carrier.
\end{remark}

\subsection{What quasi-locality does and does not factor}
\label{sec:v4v44-factorization}

For several eliminated cores, exact V38 conditioning gives skeleton
functions (Z_{U_b}(\eta)).  Their contribution after the skeleton is
integrated has the form
\begin{equation}
 \mathcal Z_{\mathcal B}
 =\mathbb E_{\nu_\Sigma}
 \prod_{b\in\mathcal B}Z_{U_b}(\eta).
\label{eq:v4v44-correlated-product}
\end{equation}
Theorem~\ref{thm:v4v44-quasilocal-activity} gives a quasi-local shell
decomposition for each factor in the integrand.

\begin{proposition}[The first connected obstruction is an ordinary
covariance]
\label{prop:v4v44-two-block-obstruction}
For two cores (b_1,b_2),
\begin{equation}
 \mathcal Z_{\{b_1,b_2\}}
 =
 \mathbb E Z_{U_{b_1}}\,
 \mathbb E Z_{U_{b_2}}
 +
 \operatorname{Cov}_{\nu_\Sigma}
 (Z_{U_{b_1}},Z_{U_{b_2}}).
\label{eq:v4v44-two-block-covariance}
\end{equation}
Therefore individual quasi-locality implies the scalar product form in
\eqref{eq:v4v43-local-constants} only if the connected skeleton term is also
controlled and assigned its own activity ancestry.
\end{proposition}

\begin{proof}
Equation \eqref{eq:v4v44-two-block-covariance} is the definition of
covariance applied to \eqref{eq:v4v44-correlated-product}.  The second term
is generally nonzero for a correlated skeleton law, even when the two
functions are local or exponentially quasi-local.
\end{proof}

\begin{firewall}[Quasi-local factors are not independent scalar factors]
\label{fw:v4v44-no-false-factorization}
The product in \eqref{eq:v4v44-correlated-product} occurs inside one
correlated expectation.  Replacing it by
(\prod_b\mathbb E Z_{U_b}) discards the covariance in
\eqref{eq:v4v44-two-block-covariance} and all higher connected cumulants.
Thus V44 does not claim that the hypothesis
\eqref{eq:v4v43-local-constants} has been proved.  It proves the quasi-local
input from which a connected normalization expansion must start.
\end{firewall}

\subsection{V44 result ledger and V45 target}
\label{sec:v4v44-conclusion}

V44 proves a zero-free buffered derivative for one conditional scalar
normalization, composes it with the exact V41 sweep resolvent, and constructs
an exact shell telescope for its logarithm.  Under the uniform positive-law
moment bound, the shell activities are absolutely summable in (L^1),
uniformly in volume.  This supplies a rigorous quasi-locality theorem for
each conditional block normalization without calling a skeleton function a
vacuum constant.

The multi-block scalar normalization remains correlated.  V45 must first
perform the scheduled YM/NS audit, then attack the connected term in
\eqref{eq:v4v44-correlated-product}.  The immediate target is a two-block
bound obtained by inserting the V44 shell expansions into the skeleton
Witten covariance.  Only after the resulting decay and ancestry are fixed
should one decide whether a tree-cumulant criterion closes all block
orders, or whether the programme must move further upstream to a positive
skeleton polymerization.


\clearpage
\section{V45: companion audit and the two-block connected normalization}
\label{sec:v4v45}

V44 proves that one conditional block normalization is quasi-local in the
retained skeleton.  That statement controls the truncation error of one
factor.  It does not control the connected expectation of several factors.
This note performs the scheduled companion audit and then closes the first
nontrivial connected case.  Two block factors are truncated to disjoint
neighbourhoods, the skeleton Witten kernel controls the covariance of the
truncations, and the V44 tails restore the original factors.  The resulting
pair activity decays exponentially in the block separation.

The proof deliberately stops at order two.  Pair covariance decay does not
imply an all-order cumulant tree estimate.  The all-order normalization
polymer is the next gate.

\subsection{Scheduled read-only companion audit}
\label{sec:v4v45-audit}

The audit was performed on the newest active companion notes present in the
shared \texttt{latex\_notes} folder at the start of V45.  Their exact
fingerprints were
\begin{align*}
 &\texttt{Renewal\_Yang\_Mills\_Mass\_Gap\_Research\_Notes\_V61.tex},\\
 &\qquad 1\,098\,122\ \text{bytes},\qquad
 \texttt{SHA--256 }\
 \texttt{53ECBD3348DF95FDE8CA89B7E80581132821842FB82DDCB4FFA1082161F16E66},\\
 &\texttt{Renewal\_NS\_Stability\_Research\_Notes\_V31.tex},\\
 &\qquad 887\,537\ \text{bytes},\qquad
 \texttt{SHA--256 }\
 \texttt{F1121B62238AD5491BB7CF03635EA9934942EDF573ED8FAAA9EE79E1D38A6696}.
\end{align*}

The Yang--Mills note proves that its coefficient-one moment/covariance
restriction preserves genuine response factors but creates no missing
partner-ownership fraction.  Its spectator-extension test rules out
relabeling an ordinary moment as the required normalized response carrier.
The Navier--Stokes note proves that flat probability marks remove stopping
family multiplicity, while removal of the mark costs the exact dyadic first
moment.  A uniform bound on each normalized mark does not pay that moment.

\begin{assessment}[Type-correct transfer from the companion audit]
\label{ass:v4v45-audit-transfer}
Only the proof order transfers to the present programme.
\begin{enumerate}
\item The V44 shell restriction has coefficient one.  It preserves an
      existing skeleton response but cannot create a connected covariance
      factor or an inverse block-density weight.
\item A bound uniform in the shell or block label does not pay the sum over
      those labels.  The exact exponential shell weight and the geometric
      block density must remain visible until summation.
\item Any missing connected smallness must come from an actual skeleton
      covariance or higher cumulant.  It may not be attached after taking
      norms merely because such a factor would close the ledger.
\end{enumerate}
No Yang--Mills character estimate and no Navier--Stokes heat estimate is
imported into the Standard Model carrier.
\end{assessment}

\subsection{Skeleton mixing interface}
\label{sec:v4v45-skeleton-interface}

Let (\nu_\Sigma) be the positive skeleton law after the chosen V38 core
elimination and before complex block insertions are averaged.  Its Witten
one-form operator is denoted by (\mathcal L_{1,\Sigma}).  Assume the
uniform spatial kernel estimate
\begin{equation}
 \|P_X\mathcal L_{1,\Sigma}^{-1}P_Y\|
 \le C_\Sigma e^{-m_\Sigma d(X,Y)}
\label{eq:v4v45-skeleton-Witten}
\end{equation}
for skeleton sets (X,Y).  This is the skeleton version of the V30/V38
finite-range uniformly convex Witten estimate.  Strong convexity survives
the V38 marginalization by Theorem
\ref{thm:v4v38-marginal-convexity}; the finite-range row and commutator
bounds needed for \eqref{eq:v4v45-skeleton-Witten} are retained here as an
explicit hypothesis because an arbitrary multilevel effective Hessian need
not remain finite range.

For complex (C^1) functions (F,G), the complexified
Helffer--Sjostrand identity is
\begin{equation}
 \operatorname{Cov}_{\nu_\Sigma}(F,G)
 =\mathbb E_{\nu_\Sigma}
 \left\langle\nabla F,
 \mathcal L_{1,\Sigma}^{-1}\nabla G\right\rangle,
\label{eq:v4v45-HS-covariance}
\end{equation}
with bilinear covariance and with real and imaginary parts treated
separately when taking absolute values.

\begin{lemma}[Separated local skeleton functions]
\label{lem:v4v45-separated-covariance}
If (F) and (G) depend only on skeleton coordinates in (X) and (Y),
respectively, then
\begin{equation}
 |\operatorname{Cov}_{\nu_\Sigma}(F,G)|
 \le C_\Sigma e^{-m_\Sigma d(X,Y)}
 \|\nabla F\|_{L^2(\nu_\Sigma)}
 \|\nabla G\|_{L^2(\nu_\Sigma)}.
\label{eq:v4v45-separated-covariance}
\end{equation}
\end{lemma}

\begin{proof}
Insert (P_X) and (P_Y) on the two gradients in
\eqref{eq:v4v45-HS-covariance}, apply
\eqref{eq:v4v45-skeleton-Witten}, and then use Cauchy--Schwarz in the field
measure.  Complexification changes only a fixed harmless factor, absorbed
in (C_\Sigma).
\end{proof}

\begin{firewall}[The effective-range hypothesis is not automatic]
\label{fw:v4v45-effective-range}
The exact Hessian in \eqref{eq:v4v38-effective-Hessian} contains a covariance
subtraction and can acquire long tails after repeated eliminations.  Strong
convexity alone supplies a spectral gap but not the spatial kernel
\eqref{eq:v4v45-skeleton-Witten}.  At a single finite-range skeleton stage,
the V30 commutator proof applies.  At later stages, exponential Hessian tails
must be propagated or the argument must be restarted from the preceding
finite-range parent.
\end{firewall}

\subsection{Local truncations of a block normalization}
\label{sec:v4v45-truncations}

For a block (b), write
\begin{equation}
 F_b(\eta):=Z_{U_b}(\eta)
\label{eq:v4v45-block-factor}
\end{equation}
for the V44 scalar normalization, and use the truncated boundary
\eqref{eq:v4v44-truncated-boundary} to define
\begin{equation}
 F_b^{[R]}(\eta):=Z_{U_b}(\eta^{[R]}).
\label{eq:v4v45-truncated-factor}
\end{equation}
The reference field outside (B_R(b)) is fixed.  Hence
(F_b^{[R]}) depends only on (B_R(b)).

The V44 first-response estimate and its shell moment argument give the
following quantitative interface.  Assume constants (M_0,A_0,G_0,\alpha>0),
uniform in (b,R), and total volume, such that
\begin{align}
 \|F_b\|_{L^\infty(\nu_\Sigma)}
 +\|F_b^{[R]}\|_{L^\infty(\nu_\Sigma)}
 &\le M_0,
\label{eq:v4v45-uniform-factor}\\
 \|F_b-F_b^{[R]}\|_{L^2(\nu_\Sigma)}
 &\le A_0e^{-\alpha R},
\label{eq:v4v45-L2-tail}\\
 \|\nabla F_b^{[R]}\|_{L^2(\nu_\Sigma)}
 &\le G_0.
\label{eq:v4v45-gradient-ledger}
\end{align}

\begin{lemma}[Acquisition of the truncation interfaces]
\label{lem:v4v45-truncation-acquisition}
Suppose the V44 sweep-resolved directional derivative bound holds shell by
shell, and suppose the positive skeleton law has the uniform second shell
moment
\begin{equation}
 \left\|
  \|\eta-\bar\eta\|_{\ell^2(S_r(b))}
 \right\|_{L^2(\nu_\Sigma)}
 \le M_2(1+r)^p.
\label{eq:v4v45-second-shell-moment}
\end{equation}
Then \eqref{eq:v4v45-L2-tail} holds for every
(0<\alpha<\gamma), with a uniform (A_0).  If the near-block directional
derivatives are uniformly square integrable, then
\eqref{eq:v4v45-gradient-ledger} also holds.  Under
\eqref{eq:v4v44-small-insertion}, one may take (M_0=1+\epsilon).
\end{lemma}

\begin{proof}
Telescope (F_b-F_b^{[R]}) over shells (r>R) by changing one shell at a
time.  The unlogged form of the V44 response estimate gives

\[
 \|F_b^{[r]}-F_b^{[r-1]}\|_{L^2}
 \le C_0e^{-\gamma(r-R_0)}M_2(1+r)^p.
\]

Minkowski's inequality and the exponential tail show that the sum over
(r>R) is at most (A_0e^{-\alpha R}) for any
(\alpha<\gamma).  For the gradient, split a skeleton direction into
orthogonal shells.  The squared operator norm of the derivative is bounded
by the finite near-block contribution plus

\[
 C_0^2\sum_{r>R_0}e^{-2\gamma(r-R_0)},
\]

uniformly in (R).  Integrating this pointwise or conditional bound under
(\nu_\Sigma) gives \eqref{eq:v4v45-gradient-ledger}.  Finally,
(|F_b|=|\mathbb E(1+U_b)|\le1+\epsilon), and the same estimate applies to
the truncated boundary.
\end{proof}

\begin{firewall}[The (L^2) tail is a new, explicit moment requirement]
\label{fw:v4v45-L2-not-L1}
Corollary~\ref{cor:v4v44-shell-summability} proves an (L^1) shell sum from
first moments.  A covariance approximation needs
\eqref{eq:v4v45-L2-tail}.  V45 therefore assumes the second shell moment
\eqref{eq:v4v45-second-shell-moment} and derives the stronger tail; it does
not silently use \(L^1\) convergence as \(L^2\) convergence.
\end{firewall}

\subsection{Exponential decay of the connected pair}
\label{sec:v4v45-pair}

Let (D=d(b,c)) denote the distance between the reference supports of two
blocks.  Enlarge the definition of (D) by a fixed block-diameter constant
so that
\begin{equation}
 d(B_R(b),B_R(c))\ge(D-2R)_+.
\label{eq:v4v45-support-separation}
\end{equation}

\begin{theorem}[Two-block connected normalization]
\label{thm:v4v45-two-block}
Under \eqref{eq:v4v45-skeleton-Witten} and
\eqref{eq:v4v45-uniform-factor}--%
\eqref{eq:v4v45-gradient-ledger}, there are constants
(C_2,m_2>0), independent of the block pair and total volume, such that
\begin{equation}
 \boxed{
 |\operatorname{Cov}_{\nu_\Sigma}(F_b,F_c)|
 \le C_2e^{-m_2D}.}
\label{eq:v4v45-two-block}
\end{equation}
One may take any
\begin{equation}
 m_2<\frac13\min\{\alpha,m_\Sigma\}.
\label{eq:v4v45-pair-rate}
\end{equation}
\end{theorem}

\begin{proof}
For (D\ge6), choose (R=\lfloor D/3\rfloor) and abbreviate
(F_R=F_b^{[R]}), (G_R=F_c^{[R]}).  Add and subtract the local
truncations:
\begin{align*}
 \operatorname{Cov}(F_b,F_c)
 &=\operatorname{Cov}(F_R,G_R)\\
 &\quad+\operatorname{Cov}(F_b-F_R,F_c)
 +\operatorname{Cov}(F_R,F_c-G_R).
\end{align*}
For complex random variables, Cauchy--Schwarz gives

\[
 |\operatorname{Cov}(X,Y)|
 \le\|X-\mathbb EX\|_2\|Y-\mathbb EY\|_2
 \le\|X\|_2\|Y\|_2.
\]

Hence the last two terms have total modulus at most
(2M_0A_0e^{-\alpha R}).  The local term is controlled by Lemma
\ref{lem:v4v45-separated-covariance},
\eqref{eq:v4v45-support-separation}, and
\eqref{eq:v4v45-gradient-ledger}:

\[
 |\operatorname{Cov}(F_R,G_R)|
 \le C_\Sigma G_0^2e^{-m_\Sigma(D-2R)}.
\]

Since (R\le D/3), both exponents are at least a fixed multiple of (D/3),
up to an absorbed rounding constant.  This proves
\eqref{eq:v4v45-two-block} at every rate in
\eqref{eq:v4v45-pair-rate}.  Increase (C_2) to cover the finitely many
separations (D<6), using
(|\operatorname{Cov}(F_b,F_c)|\le M_0^2).
\end{proof}

\begin{corollary}[A zero-free far-pair logarithmic activity]
\label{cor:v4v45-pair-log}
Assume also (|F_b-1|\le\epsilon<1), so that

\[
 m_b:=\mathbb E_{\nu_\Sigma}F_b,
 \qquad |m_b|\ge\rho:=1-\epsilon.
\]

If (D) is large enough that (C_2e^{-m_2D}\le\rho^2/2), define
\begin{equation}
 \Psi_{bc}:=
 \log\mathbb E(F_bF_c)-\log m_b-\log m_c
\label{eq:v4v45-pair-activity}
\end{equation}
on the branch continued from zero insertion.  Then
\begin{equation}
 |\Psi_{bc}|
 \le\frac{2C_2}{\rho^2}e^{-m_2D}.
\label{eq:v4v45-pair-activity-bound}
\end{equation}
\end{corollary}

\begin{proof}
By \eqref{eq:v4v44-two-block-covariance},

\[
 \mathbb E(F_bF_c)
 =m_bm_c\left(1+w_{bc}\right),
 \qquad
 w_{bc}:=\frac{\operatorname{Cov}(F_b,F_c)}{m_bm_c}.
\]

The hypotheses give (|w_{bc}|\le1/2), so the continued logarithm satisfies
(|\log(1+w_{bc})|\le2|w_{bc}|).  Apply
Theorem~\ref{thm:v4v45-two-block}.
\end{proof}

\begin{remark}[Pair summation pays the exact geometric density]
\label{rem:v4v45-pair-sum}
On a lattice of polynomial volume growth, the number of blocks (c) at
distance in ([n,n+1)) from a fixed (b) is at most (C(1+n)^{d-1}).
Therefore \eqref{eq:v4v45-pair-activity-bound} gives

\[
 \sup_b\sum_{c:\,d(b,c)\ge D_0}|\Psi_{bc}|
 \le\frac{2C_2C}{\rho^2}
 \sum_{n\ge D_0}(1+n)^{d-1}e^{-m_2n}<\infty.
\]

This is the pair analogue of the NS audit lesson: the exponential activity
is retained until the exact shell multiplicity is summed.
\end{remark}

\subsection{Why the pair theorem does not close the pressure}
\label{sec:v4v45-higher}

For a finite block set (X), define the connected skeleton cumulant
\begin{equation}
 \kappa_\Sigma(F_b:b\in X)
 :=
 \left.
 \prod_{b\in X}\frac{\partial}{\partial t_b}
 \log\mathbb E_{\nu_\Sigma}
 \exp\left(\sum_{b\in X}t_bF_b\right)
 \right|_{t=0}.
\label{eq:v4v45-connected-cumulant}
\end{equation}
For (|X|=2), this is the covariance controlled in
\eqref{eq:v4v45-two-block}.  A pressure expansion for
\eqref{eq:v4v44-correlated-product} requires a uniform summable tree bound
for all (|X|).

\begin{firewall}[Two-point Witten decay is not an all-order tree]
\label{fw:v4v45-pair-not-tree}
The Helffer--Sjostrand identity
\eqref{eq:v4v45-HS-covariance} has two gradient slots.  Iterating it
differentiate the Witten carrier and creates the higher insertions identified
in V37.  Theorem~\ref{thm:v4v45-two-block} therefore proves the pair activity
only.  It does not imply

\[
 |\kappa_\Sigma(F_b:b\in X)|
 \le C^{|X|}\sum_{T\text{ on }X}
 \prod_{\{b,c\}\in T}e^{-m_2d(b,c)}
\]

for (|X|\ge3).  Attaching the proved pair factors to an arbitrary spanning
tree would repeat the precise inference rejected in V37 and by the YM V61
audit.
\end{firewall}

\begin{assessment}[Exact location of the remaining scalar gate]
\label{ass:v4v45-scalar-gate}
The V43 pressure factorization may now be replaced through order two by
actual connected objects: one-block means (m_b) and pair activities
(\Psi_{bc}).  Both have volume-uniform local ledgers, and the pair ledger is
summable over relative positions.  The first unproved scalar object is the
three-block cumulant of the quasi-local normalizations.  Its proof must
either retain an actual three-slot response generated by the positive
skeleton law or use a polymer representation whose connectedness is exact
before estimation.
\end{assessment}

\subsection{V45 result ledger and V46 target}
\label{sec:v4v45-conclusion}

V45 records the exact active YM V61 and NS V31 fingerprints and imports only
their proof-order discipline.  It upgrades V44's shell control to an
(L^2) local truncation under a stated second shell moment, proves a local
skeleton covariance bound from the Witten kernel, and combines the two to
obtain exponential decay of the connected two-block normalization.  The
corresponding far-pair logarithmic activity is zero-free and summable over
relative block positions.

V46 should attack the three-block connected normalization before proposing
an all-order compiler.  The first calculation is the exact derivative of
(\operatorname{Cov}_{\nu_\Sigma}(F_b,F_c)) under a third local insertion,
with the derivative of the Witten inverse retained.  If that calculation
recreates the uncontrolled V37 carrier, the direction must move upstream to
an exact positive-skeleton polymerization or a finite-dependence
decomposition; it must not infer a tree bound from
\eqref{eq:v4v45-two-block} alone.


\clearpage
\section{V46: the exact three-block derivative exposes the Witten carrier}
\label{sec:v4v46}

V45 reduces the first open scalar object to the three-block connected
cumulant.  This note performs the exact calculation before applying any
absolute value.  Tilting the positive skeleton law by the third insertion
and differentiating the Helffer--Sjostrand covariance produces two terms.
One is a two-resolvent insertion with the geometry of a three-vertex path.
The other is the covariance of the Witten energy density with the third
insertion.  Its gradient differentiates the Witten inverse with respect to
the skeleton field and recreates the V37 higher carrier.

Thus the two-block theorem cannot simply be iterated.  The calculation also
identifies a sufficient carrier norm which would close order three.  Since
that norm is not supplied by lower convexity and two-point decay alone, the
next version will move upstream to the exact conditional-cumulant algebra of
the bipartite positive carrier.

\subsection{Third cumulant as a tilted covariance derivative}
\label{sec:v4v46-tilt}

Let

\[
 d\nu=Z^{-1}e^{-S(\eta)}\,d\eta
\]

be a finite-dimensional positive skeleton law, and let \(F,G,H\) be smooth
complex functions with sufficient integrability for the differentiations
below.  Define the \(H\)-tilted law
\begin{equation}
 d\nu_t=
 \frac{e^{tH}}{\mathbb E_\nu e^{tH}}\,d\nu,
 \qquad S_t=S-tH.
\label{eq:v4v46-tilted-law}
\end{equation}
For real \(t\) and real \(H\) this is a positive law near zero.  The complex
formula follows by polarization or by treating \(t\) as a formal analytic
parameter wherever the moment generating function is nonzero.

The third joint cumulant is
\begin{align}
 \kappa_3(F,G,H)
 &:={\mathbb E}(FGH)
 -{\mathbb E}(FG){\mathbb E}H
 -{\mathbb E}(FH){\mathbb E}G
 -{\mathbb E}(GH){\mathbb E}F
 +2{\mathbb E}F{\mathbb E}G{\mathbb E}H.
\label{eq:v4v46-third-cumulant}
\end{align}

\begin{lemma}[Tilt derivative identity]
\label{lem:v4v46-tilt-derivative}
One has
\begin{equation}
 \kappa_3(F,G,H)
 =\left.
 \frac{d}{dt}\operatorname{Cov}_{\nu_t}(F,G)
 \right|_{t=0}.
\label{eq:v4v46-tilt-derivative}
\end{equation}
\end{lemma}

\begin{proof}
For every integrable \(A\), normalized differentiation gives

\[
 \left.\frac d{dt}\mathbb E_{\nu_t}A\right|_{t=0}
 =\operatorname{Cov}_\nu(A,H).
\]

Apply this to
\(\mathbb E_t(FG)-\mathbb E_tF\,\mathbb E_tG\) and expand.  The five terms
are exactly \eqref{eq:v4v46-third-cumulant}.
\end{proof}

\subsection{Differentiating the one-form inverse}
\label{sec:v4v46-inverse}

Use the one-form convention
\begin{equation}
 \mathcal L_{1,t}
 =-\Delta+\nabla S_t\mathbin\cdot\nabla+D^2S_t,
 \qquad R_t:=\mathcal L_{1,t}^{-1}.
\label{eq:v4v46-Witten-operator}
\end{equation}
Define the actual tilt insertion on one-forms
\begin{equation}
 \mathcal B_H
 :=\nabla H\mathbin\cdot\nabla+D^2H.
\label{eq:v4v46-tilt-insertion}
\end{equation}
Then
\begin{equation}
 \left.\frac d{dt}\mathcal L_{1,t}\right|_{t=0}
 =-\mathcal B_H,
 \qquad
 \left.\frac d{dt}R_t\right|_{t=0}
 =R\mathcal B_HR,
 \quad R:=R_0.
\label{eq:v4v46-inverse-derivative}
\end{equation}

\begin{theorem}[Exact three-cumulant Witten decomposition]
\label{thm:v4v46-three-cumulant}
Let
\begin{equation}
 K_{F,G}(\eta):=
 \langle\nabla F(\eta),R\nabla G(\eta)\rangle.
\label{eq:v4v46-energy-density}
\end{equation}
Then
\begin{equation}
 \boxed{
 \kappa_3(F,G,H)
 =
 \operatorname{Cov}_\nu(K_{F,G},H)
 +\mathbb E_\nu
 \langle\nabla F,R\mathcal B_HR\nabla G\rangle.}
\label{eq:v4v46-three-cumulant-decomposition}
\end{equation}
Both terms are actual derivatives of the same covariance; neither may be
discarded or counted twice.
\end{theorem}

\begin{proof}
For each \(t\), the Helffer--Sjostrand identity gives

\[
 \operatorname{Cov}_{\nu_t}(F,G)
 =\mathbb E_{\nu_t}
 \langle\nabla F,R_t\nabla G\rangle.
\]

Differentiate at zero.  Differentiating the normalized expectation produces
\(\operatorname{Cov}_\nu(K_{F,G},H)\).  Differentiating \(R_t\) produces the
second term by \eqref{eq:v4v46-inverse-derivative}.  Invoke
Lemma~\ref{lem:v4v46-tilt-derivative}.
\end{proof}

\begin{firewall}[The insertion term is not the complete third cumulant]
\label{fw:v4v46-insertion-not-all}
The second term in
\eqref{eq:v4v46-three-cumulant-decomposition} visibly contains two Witten
bridges and is tempting to retain as the whole answer.  The first term is the
change of the normalized carrier measure and is generally nonzero.  Omitting
it would differentiate an unnormalized integral while calling the result a
cumulant.
\end{firewall}

\subsection{The term with manifest three-vertex geometry}
\label{sec:v4v46-manifest-tree}

Let \(X_F,X_G,X_H\) be the coordinate supports of the three functions.  The
operator \(\mathcal B_H\) is supported on \(X_H\).  To state the exact norm
needed for its insertion, let \(\mathcal D_H\) be a one-form space supported
on \(X_H\) and assume
\begin{equation}
 |\mathbb E_\nu\langle u,\mathcal B_Hv\rangle|
 \le B_H
 \|u\|_{L^2(\nu;\mathcal D_H)}
 \|v\|_{L^2(\nu;\mathcal D_H)}
\label{eq:v4v46-insertion-form-bound}
\end{equation}
for the resolvent outputs used below.  This is a form bound on the actual
first-order-plus-Hessian insertion, not merely a bound on \(H\).

\begin{proposition}[The explicit insertion carries a path tree]
\label{prop:v4v46-path-term}
Assume the projected Witten estimate
\eqref{eq:v4v45-skeleton-Witten} and
\eqref{eq:v4v46-insertion-form-bound}.  Then
\begin{align}
 &\left|
 \mathbb E_\nu
 \langle\nabla F,R\mathcal B_HR\nabla G\rangle
 \right|\\
 &\qquad\le
 C_\Sigma^2B_H
 e^{-m_\Sigma[d(X_F,X_H)+d(X_H,X_G)]}
 \|\nabla F\|_2\|\nabla G\|_2.
\label{eq:v4v46-path-term}
\end{align}
\end{proposition}

\begin{proof}
Because \(\mathcal B_H\) is supported in \(X_H\), insert \(P_{X_H}\) on
both sides of it.  Apply \eqref{eq:v4v45-skeleton-Witten} to
\(P_{X_H}R P_{X_F}\) and \(P_{X_H}R P_{X_G}\), then apply the form bound
\eqref{eq:v4v46-insertion-form-bound}.  The product of the two exponential
bridges gives \eqref{eq:v4v46-path-term}.
\end{proof}

\begin{remark}[All three labelled paths occur by symmetry]
\label{rem:v4v46-three-paths}
The cumulant is symmetric in \(F,G,H\), although
\eqref{eq:v4v46-three-cumulant-decomposition} singles out the tilt variable.
Repeating the identity with each variable in the middle exposes the other
two labelled path trees.  This symmetry does not estimate the remaining
energy-density covariance; it only gives three equivalent decompositions of
the same scalar.
\end{remark}

\subsection{Field differentiation of the Witten energy density}
\label{sec:v4v46-field-carrier}

To apply \eqref{eq:v4v45-HS-covariance} to
\(\operatorname{Cov}(K_{F,G},H)\), one needs \(\nabla K_{F,G}\).  Let \(v\)
be a constant skeleton tangent direction and write \(\partial_v\) for the
corresponding directional derivative on functions.  Direct differentiation
of the coefficients of the untilted one-form operator gives the commutator
\begin{equation}
 [\partial_v,\mathcal L_1]
 =\mathcal A_v,
 \qquad
 \mathcal A_v
 :=(D^2S\,v)\mathbin\cdot\nabla+D^3S[v].
\label{eq:v4v46-field-insertion}
\end{equation}
Since \(\mathcal L_1R=1\), commuting \(\partial_v\) through this identity
gives
\begin{equation}
 [\partial_v,R]=-R\mathcal A_vR.
\label{eq:v4v46-field-inverse}
\end{equation}

\begin{proposition}[Exact gradient of the Witten energy density]
\label{prop:v4v46-energy-gradient}
For every direction \(v\),
\begin{align}
 D_\eta K_{F,G}[v]
 &=\langle D^2F[v],R\nabla G\rangle
 +\langle\nabla F,RD^2G[v]\rangle\\
 &\quad-
 \langle\nabla F,R\mathcal A_vR\nabla G\rangle.
\label{eq:v4v46-energy-gradient}
\end{align}
The final summand contains the full drift insertion
\((D^2S\,v)\cdot\nabla\) and the third derivative \(D^3S[v]\).
\end{proposition}

\begin{proof}
Apply the product rule to
\eqref{eq:v4v46-energy-density}.  The two endpoint derivatives give the
first line.  To differentiate \(R\nabla G\), commute \(\partial_v\) through
\(R\) with \eqref{eq:v4v46-field-inverse}; this gives the second endpoint
term and the carrier term.
\end{proof}

For the quartic positive Higgs carrier, \(D^3S[v]\) grows linearly in the
field and \(D^2S\,v\) contains the unbounded onsite Hessian.  Strong
convexity controls \(R\) in \(L^2\), but it does not by itself give the
localized operator-form bounds for \(\mathcal A_v\) needed after the two
resolvents in \eqref{eq:v4v46-energy-gradient}.
\begin{firewall}[The V37 carrier reappears before order-three estimation]
\label{fw:v4v46-V37-carrier}
Applying a two-point covariance estimate to
\(\operatorname{Cov}(K_{F,G},H)\) requires the three terms in
\eqref{eq:v4v46-energy-gradient}.  Bounding only the endpoint Hessians
omits the derivative of the inverse.  Replacing that derivative by a second
copy of the pair bound omits \(\mathcal A_v\).  This is exactly the
higher-carrier obstruction isolated in V37, now derived for the V45
three-block normalization.
\end{firewall}

\subsection{A sufficient order-three carrier interface}
\label{sec:v4v46-sufficient}

The open estimate can be stated without ambiguity.  Let
\(\mathfrak t(F,G,H)\) be the sum of the three labelled path weights
\begin{align}
 \mathfrak t(F,G,H)
 :=&\ e^{-m[d(X_F,X_G)+d(X_G,X_H)]}\\
 &+e^{-m[d(X_G,X_H)+d(X_H,X_F)]}\\
 &+e^{-m[d(X_H,X_F)+d(X_F,X_G)]}.
\label{eq:v4v46-three-tree-weight}
\end{align}

\begin{theorem}[Sufficient three-block carrier bound]
\label{thm:v4v46-sufficient-carrier}
Suppose the insertion term satisfies
\eqref{eq:v4v46-path-term}, and suppose the actual energy-density covariance
obeys
\begin{equation}
 |\operatorname{Cov}_\nu(K_{F,G},H)|
 \le C_3A_FA_GA_H\,\mathfrak t(F,G,H),
\label{eq:v4v46-open-carrier-bound}
\end{equation}
where the \(A\)'s are volume-uniform local derivative norms.  Then
\begin{equation}
 |\kappa_3(F,G,H)|
 \le C'_3A_FA_GA_H\,\mathfrak t(F,G,H).
\label{eq:v4v46-three-cumulant-tree}
\end{equation}
For V44 block factors, local truncation at a fixed fraction of the minimal
spanning-tree length adds only an exponentially smaller error, as in the
proof of Theorem~\ref{thm:v4v45-two-block}.
\end{theorem}

\begin{proof}
Use the exact decomposition
\eqref{eq:v4v46-three-cumulant-decomposition}.  Bound its first term by
\eqref{eq:v4v46-open-carrier-bound}.  Proposition
\ref{prop:v4v46-path-term} bounds the second term by one of the three
summands in \eqref{eq:v4v46-three-tree-weight}.  Enlarge the constant.  For
quasi-local factors, add and subtract local truncations.  Every error
contains at least one V45 \(L^2\) truncation tail; generalized
Hölder with the bounded remaining factors pays that tail.  Choosing the
truncation radius proportional to the minimal tree length preserves an
exponential rate.
\end{proof}

\begin{assessment}[The sufficient interface has not been acquired]
\label{ass:v4v46-interface-open}
Equation \eqref{eq:v4v46-open-carrier-bound} is the exact missing estimate,
not a result of V46.  Neither V38 conditional Witten decay, V41 maximal
correlation, nor V45 pair covariance controls the field insertion
\(\mathcal A_v\) in \eqref{eq:v4v46-energy-gradient}.  Theorem
\ref{thm:v4v46-sufficient-carrier} records what would suffice and prevents a
future proof from hiding the missing derivative inside a generic tree
constant.
\end{assessment}

\subsection{V46 result ledger and upstream V47 target}
\label{sec:v4v46-conclusion}

V46 proves the exact tilted-Witten formula
\eqref{eq:v4v46-three-cumulant-decomposition}.  Its explicit inverse
derivative has two Witten bridges and a local middle insertion, hence a
genuine path-tree estimate under a stated form bound.  The normalized-measure
derivative produces a second term whose gradient contains the complete
field insertion \(\mathcal A_v\).  This identifies, rather than assumes
away, the first unproved three-block carrier.

The direct differentiation route has therefore returned to the V37
obstruction.  V47 should go upstream to the bipartite conditional law before
the Witten inverse is differentiated.  The target is the exact law of total
cumulance under one color conditioning, including every set partition and
every conditional connected block.  Conditional independence may then
annihilate whole terms algebraically.  Only the surviving conditional
cumulants should be estimated; no pair bound should be promoted to an
all-order tree without this incidence calculation.


\clearpage
\section{V47: conditional cumulance and a three-block bound without differentiating the resolvent}
\label{sec:v4v47}

V46 shows that differentiating the Witten representation of a covariance
splits the third cumulant into a manifest two-bridge term and a difficult
field derivative of the Witten carrier.  That split is exact, but it is not
the only exact organization of the same cumulant.  At order three, one may
first group two local factors into their complete product and then use a
single covariance across a geometric cut.  This keeps the cancellation
between the two V46 terms intact and never differentiates the resolvent.

This note develops that upstream route.  It first records the law of total
cumulance under one color conditioning and identifies the terms annihilated
by conditional independence.  It then proves the product-regrouping identity
and applies the V45 skeleton Witten estimate to local truncations of three
V44 block factors.  The resulting connected three-block coefficient decays
exponentially in the diameter of the block triple and is summable over the
two free block positions.

\subsection{The law of total cumulance}
\label{sec:v4v47-total-cumulance}

Let \(\mathscr G\) be a sub-\(\sigma\)-field.  For a nonempty finite set
\(B\) of indices, write
\[
 \kappa(X_i:i\in B\mid\mathscr G)
\]
for the joint cumulant computed in the conditional law given
\(\mathscr G\).  It is a \(\mathscr G\)-measurable random variable.  If
\(\pi\) is a partition of an index set \(I\), its blocks are denoted by
\(B\in\pi\).

\begin{theorem}[Law of total cumulance]
\label{thm:v4v47-total-cumulance}
For random variables with a joint moment generating function in a
neighbourhood of the origin,
\begin{equation}
 \kappa(X_i:i\in I)
 =
 \sum_{\pi\in\Pi(I)}
 \kappa\left(
   \kappa(X_i:i\in B\mid\mathscr G):B\in\pi
 \right).
\label{eq:v4v47-total-cumulance}
\end{equation}
In the outer cumulant, the arguments are the conditional cumulants indexed
by the blocks of \(\pi\).  A one-argument outer cumulant is expectation.
\end{theorem}

\begin{proof}
Introduce formal variables \(t_i\) and the conditional cumulant generating
function
\[
 K_{\mathscr G}(t)
 :=
 \log\mathbb E\left[
  \exp\left(\sum_{i\in I}t_iX_i\right)
  \middle|\mathscr G
 \right].
\]
The unconditional cumulant generating function is
\[
 K(t)=\log\mathbb E\exp(K_{\mathscr G}(t)).
\]
Differentiate once in every \(t_i\) and evaluate at zero.  The multivariate
chain rule partitions the derivatives among the copies of
\(K_{\mathscr G}\).  A block \(B\) placed on one copy gives
\(\kappa(X_i:i\in B\mid\mathscr G)\), while differentiation of the outer
logarithm gives the joint cumulant of those block variables.  Summing over
all set partitions gives \eqref{eq:v4v47-total-cumulance}.  Analyticity
justifies the calculation; truncation gives the usual finite-moment
extension whenever all displayed moments exist.
\end{proof}

For three variables, Theorem~\ref{thm:v4v47-total-cumulance} reads
\begin{align}
 \kappa_3(X,Y,Z)
 &=\mathbb E\,\kappa_3(X,Y,Z\mid\mathscr G)
 +\kappa_3(\mathbb E[X\mid\mathscr G],
           \mathbb E[Y\mid\mathscr G],
           \mathbb E[Z\mid\mathscr G])\nonumber\\
 &\quad
 +\operatorname{Cov}\bigl(
    \mathbb E[X\mid\mathscr G],
    \operatorname{Cov}(Y,Z\mid\mathscr G)\bigr)\nonumber\\
 &\quad
 +\operatorname{Cov}\bigl(
    \mathbb E[Y\mid\mathscr G],
    \operatorname{Cov}(X,Z\mid\mathscr G)\bigr)\nonumber\\
 &\quad
 +\operatorname{Cov}\bigl(
    \mathbb E[Z\mid\mathscr G],
    \operatorname{Cov}(X,Y\mid\mathscr G)\bigr).
\label{eq:v4v47-total-cumulance-three}
\end{align}

\begin{corollary}[Conditional independence relocates the complete cumulant]
\label{cor:v4v47-conditional-relocation}
If \(X_1,\ldots,X_n\) belong to mutually independent components of the
conditional law given \(\mathscr G\), then
\begin{equation}
 \kappa(X_1,\ldots,X_n)
 =
 \kappa\bigl(
  \mathbb E[X_1\mid\mathscr G],\ldots,
  \mathbb E[X_n\mid\mathscr G]
 \bigr).
\label{eq:v4v47-conditional-relocation}
\end{equation}
\end{corollary}

\begin{proof}
Every conditional cumulant involving variables from two distinct
conditional-independent components vanishes.  In
\eqref{eq:v4v47-total-cumulance}, every partition except the singleton
partition contains such a block.  The singleton conditional cumulants are
the conditional means, giving \eqref{eq:v4v47-conditional-relocation}.
\end{proof}

\begin{firewall}[Conditioning moves order; it does not lower order]
\label{fw:v4v47-order-not-lowered}
Equation \eqref{eq:v4v47-conditional-relocation} replaces each variable by a
conditional mean but leaves the cumulant order equal to \(n\).  Conditional
independence therefore removes false internal connections; it does not turn
an unproved \(n\)-point bound into a product of pair bounds.
\end{firewall}

\subsection{What the V41 sweep gives before supports meet}
\label{sec:v4v47-sweep-cumulant}

For the bipartite positive Higgs carrier, one color is conditionally a
product when the other color and the exterior skeleton are fixed.  Suppose
three coefficient functions lie in distinct conditional components.  After
conditional expectation, their supports expand to neighbouring sites of the
opposite color.  Repeat only while the three expanded supports remain in
distinct components at the next conditioning.

\begin{proposition}[Pre-collision cumulant contraction]
\label{prop:v4v47-precollision}
Suppose \(m\) successive color conditionings satisfy the distinct-component
hypothesis of Corollary~\ref{cor:v4v47-conditional-relocation}.  Suppose each
centered conditional channel has \(L^2\) norm at most
\(\sqrt{q_{\rm H}}\), with \(q_{\rm H}\) from
\eqref{eq:v4v44-sweep-contraction}.  Then
\begin{align}
 |\kappa_3(F_1,F_2,F_3)|
 &\le
 2q_{\rm H}^{m}
 \min_{\{i,j,k\}=\{1,2,3\}}
 \|F_i-\mathbb EF_i\|_2
 \|F_j-\mathbb EF_j\|_2
 \|F_k\|_\infty .
\label{eq:v4v47-precollision}
\end{align}
The same estimate holds coefficientwise in the V42 finite Grassmann Hilbert
fibre, with no fibre factor per conditioning.
\end{proposition}

\begin{proof}
Apply Corollary~\ref{cor:v4v47-conditional-relocation} at each step.  This is
an exact equality until the conditional components meet.  Centering commutes
with normalized conditional expectation.  After \(m\) channels, the
\(L^2\) norm of each centered image is at most
\(q_{\rm H}^{m/2}\) times its initial centered norm, while conditional
expectation is an \(L^\infty\) contraction.  Since
\[
 \kappa_3(A,B,C)
 =\mathbb E[(A-\mathbb EA)(B-\mathbb EB)(C-\mathbb EC)],
\]
Cauchy--Schwarz on two factors and
\(\|C-\mathbb EC\|_\infty\le2\|C\|_\infty\) prove
\eqref{eq:v4v47-precollision}.  V42 identifies the coefficient Hilbert
channel with the scalar channel tensored with the identity.
\end{proof}

\begin{assessment}[Role of the pre-collision estimate]
\label{ass:v4v47-precollision-role}
Proposition~\ref{prop:v4v47-precollision} gives genuine exponential
suppression when geometry permits a number of conditional steps
proportional to separation.  Once two support light cones meet, conditional
independence no longer gives \eqref{eq:v4v47-conditional-relocation}.  The
estimate must stop there; continuing the factor \(q_{\rm H}^m\) beyond
collision would erase an actual connected interaction.
\end{assessment}

\subsection{An exact product regrouping at order three}
\label{sec:v4v47-regrouping}

\begin{lemma}[Third cumulant as covariances across one cut]
\label{lem:v4v47-covariance-regrouping}
For integrable \(F,G,H\),
\begin{equation}
 \boxed{
 \kappa_3(F,G,H)
 =
 \operatorname{Cov}(FG,H)
 -\mathbb EF\,\operatorname{Cov}(G,H)
 -\mathbb EG\,\operatorname{Cov}(F,H).}
\label{eq:v4v47-covariance-regrouping}
\end{equation}
\end{lemma}

\begin{proof}
Expand the right side:
\[
 \mathbb E(FGH)-\mathbb E(FG)\mathbb EH
 -\mathbb EF\,\mathbb E(GH)
 -\mathbb EG\,\mathbb E(FH)
 +2\mathbb EF\,\mathbb EG\,\mathbb EH.
\]
This is \eqref{eq:v4v46-third-cumulant}.
\end{proof}

\begin{proposition}[A local three-point cut estimate]
\label{prop:v4v47-local-cut}
Let \(F,G,H\) be supported in skeleton sets \(X_F,X_G,X_H\), respectively,
and suppose
\[
 d(X_H,X_F\cup X_G)\ge\delta.
\]
If \(\|F\|_\infty,\|G\|_\infty\le M\) and
\(\|\nabla F\|_2,\|\nabla G\|_2,\|\nabla H\|_2\le A\), then under
\eqref{eq:v4v45-skeleton-Witten},
\begin{equation}
 |\kappa_3(F,G,H)|
 \le4C_\Sigma MA^2e^{-m_\Sigma\delta}.
\label{eq:v4v47-local-cut}
\end{equation}
\end{proposition}

\begin{proof}
The grouped product is supported in \(X_F\cup X_G\), and
\[
 \|\nabla(FG)\|_2
 \le\|G\|_\infty\|\nabla F\|_2
   +\|F\|_\infty\|\nabla G\|_2
 \le2MA.
\]
Apply Lemma~\ref{lem:v4v45-separated-covariance} to
\(\operatorname{Cov}(FG,H)\), then to the other two covariances in
\eqref{eq:v4v47-covariance-regrouping}.  Since
\(|\mathbb EF|,|\mathbb EG|\le M\) and both pair separations are at least
\(\delta\), the three contributions total at most
\[
 2C_\Sigma MA^2e^{-m_\Sigma\delta}
 +C_\Sigma MA^2e^{-m_\Sigma\delta}
 +C_\Sigma MA^2e^{-m_\Sigma\delta}.
\]
\end{proof}

\begin{remark}[Where the V46 carrier cancellation went]
\label{rem:v4v47-carrier-cancellation}
Lemma~\ref{lem:v4v47-covariance-regrouping} estimates the complete grouped
observable \(FG\) before resolving its covariance.  If that covariance were
instead differentiated as in V46, its measure derivative and inverse
derivative would reappear separately.  Their difficult pieces cancel or
combine inside the complete product identity.  V47 preserves this
cancellation rather than estimating the two V46 summands independently.
\end{remark}

\subsection{Three quasi-local block normalizations}
\label{sec:v4v47-quasilocal-three}

Let \(b_1,b_2,b_3\) have reference centers \(x_1,x_2,x_3\) in the skeleton
metric.  The untruncated factors \(F_i:=F_{b_i}\) satisfy the V45 interfaces
\eqref{eq:v4v45-uniform-factor}--\eqref{eq:v4v45-gradient-ledger}.  Put
\begin{equation}
 D_3:=\max_{i,j}d(x_i,x_j).
\label{eq:v4v47-triple-diameter}
\end{equation}
Assume every unexpanded reference block lies within radius \(r_*\) of its
center, with \(r_*\) fixed by the block geometry.

\begin{lemma}[One member of a triple is isolated at half the diameter]
\label{lem:v4v47-isolated-center}
For three points in a metric space, there is an index \(i\) such that
\begin{equation}
 d(x_i,\{x_j,x_k\})\ge\frac{D_3}{2},
 \qquad \{i,j,k\}=\{1,2,3\}.
\label{eq:v4v47-isolated-center}
\end{equation}
\end{lemma}

\begin{proof}
Choose a diameter pair, say \(d(x_1,x_3)=D_3\).  The triangle inequality
gives
\[
 d(x_1,x_2)+d(x_2,x_3)\ge D_3.
\]
If the first summand is at least \(D_3/2\), then \(x_1\) is at distance at
least \(D_3/2\) from both \(x_2\) and \(x_3\).  Otherwise the second summand
is at least \(D_3/2\), and \(x_3\) has the stated property.
\end{proof}

\begin{theorem}[Exponential three-block connected normalization]
\label{thm:v4v47-three-block}
Under the V45 skeleton mixing and truncation interfaces, there are
constants \(C_3,m_3>0\), independent of the block triple and total volume,
such that
\begin{equation}
 \boxed{
 |\kappa_3(F_{b_1},F_{b_2},F_{b_3})|
 \le C_3e^{-m_3D_3}.}
\label{eq:v4v47-three-block}
\end{equation}
One may take any
\begin{equation}
 m_3<\frac18\min\{\alpha,m_\Sigma\}.
\label{eq:v4v47-three-rate}
\end{equation}
\end{theorem}

\begin{proof}
For large \(D_3\), choose \(R=\lfloor D_3/8\rfloor\) and replace every
\(F_i\) by \(F_i^{[R]}\).  By trilinearity of the cumulant, the difference
is a sum of three terms, each with one factor
\(F_i-F_i^{[R]}\) and two bounded factors.  For a variable \(X\) in the
first slot and bounded \(Y,Z\),
\[
 |\kappa_3(X,Y,Z)|
 \le
 \|X-\mathbb EX\|_2
 \|Y-\mathbb EY\|_2
 \|Z-\mathbb EZ\|_\infty
 \le2\|X\|_2\|Y\|_\infty\|Z\|_\infty.
\]
Equations \eqref{eq:v4v45-uniform-factor} and
\eqref{eq:v4v45-L2-tail} therefore give
\begin{equation}
 \left|
  \kappa_3(F_1,F_2,F_3)
  -\kappa_3(F_1^{[R]},F_2^{[R]},F_3^{[R]})
 \right|
 \le6M_0^2A_0e^{-\alpha R}.
\label{eq:v4v47-three-tail}
\end{equation}

Choose the isolated center supplied by
Lemma~\ref{lem:v4v47-isolated-center}.  Its truncated support is separated
from the union of the other two truncated supports by at least
\[
 \delta_R\ge\frac{D_3}{2}-2(R+r_*).
\]
For \(D_3\) larger than a fixed multiple of \(r_*\), this is at least
\(D_3/8\).  Apply Proposition~\ref{prop:v4v47-local-cut} with
\(M=M_0\) and \(A=G_0\):
\[
 |\kappa_3(F_1^{[R]},F_2^{[R]},F_3^{[R]})|
 \le4C_\Sigma M_0G_0^2e^{-m_\Sigma D_3/8}.
\]
Combine this with \eqref{eq:v4v47-three-tail}.  Rounding changes only the
constant.  Enlarge \(C_3\) to cover bounded \(D_3\), using the centered
product formula and \(\|F_i\|_\infty\le M_0\).
\end{proof}

\begin{corollary}[Absolute summability of the normalized triple coefficient]
\label{cor:v4v47-three-sum}
Assume the one-block means obey \(|m_b|\ge\rho>0\), and define
\begin{equation}
 w_{bcd}:=
 \frac{\kappa_3(F_b,F_c,F_d)}{m_bm_cm_d}.
\label{eq:v4v47-normalized-triple}
\end{equation}
On a block lattice of polynomial volume growth,
\begin{equation}
 \sup_b\sum_{c,d}|w_{bcd}|<\infty.
\label{eq:v4v47-three-sum}
\end{equation}
\end{corollary}

\begin{proof}
Theorem~\ref{thm:v4v47-three-block} gives
\[
 |w_{bcd}|\le\rho^{-3}C_3e^{-m_3D_3}.
\]
For fixed \(b\), the number of ordered pairs \((c,d)\) whose triple diameter
lies in \([n,n+1)\) is at most \(C(1+n)^{2d-1}\).  Sum the exponential
majorant.
\end{proof}

\begin{firewall}[Triple summability is not an all-order pressure theorem]
\label{fw:v4v47-three-not-all}
Equation \eqref{eq:v4v47-three-sum} closes the connected coefficient at
order three.  It supplies no bound uniform in the cumulant order and no
control of the partition factorials in
\eqref{eq:v4v47-total-cumulance}.  A pressure logarithm requires those
combinatorial weights to be summed together with block density.  V47 does
not extrapolate \(C_3\) to \(C^n\) or attach pair decay along an arbitrary
tree.
\end{firewall}

\subsection{V47 result ledger and V48 target}
\label{sec:v4v47-conclusion}

V47 proves the exact law of total cumulance, its conditional-independence
relocation rule, and the finite pre-collision contraction supplied by the
V41 color channel.  At order three it avoids the V46 differentiated carrier
by regrouping the complete product on one side of a geometric cut.  The
V45 Witten covariance then proves exponential diameter decay for the
three-block connected normalization.  V44 truncation tails restore the
quasi-local factors, and polynomial block growth makes the normalized triple
coefficient absolutely summable.

V48 should derive the general cut identity for an \(n\)-point cumulant in
terms of covariances of complete products and lower cumulants, with every
partition coefficient explicit.  The decisive question is whether recursive
choice of a separated cut yields a tree-graph majorant with constants small
enough for a Kotecky--Preiss sum after the local Yukawa factor is inserted.
If the partition factorial defeats that ledger, the route must move further
upstream to a direct polymer generating function rather than estimate
cumulants one order at a time.


\clearpage
\section{V48: an exact cumulant cut identity and its factorial ledger}
\label{sec:v4v48}

V47 proves the three-block connected bound by grouping two factors on one
side of a geometric cut and applying one actual covariance across that cut.
The same cancellation has an exact all-order form.  Compare the original
moment functional with the product functional that keeps every correlation
inside the two sides of a chosen bipartition but removes all correlations
across it.  The cumulant of the product functional vanishes.  Telescoping the
difference of the two moment--cumulant formulae exposes one genuine
cross-cut covariance in every summand.

This note proves that identity and tracks all partition coefficients.  For
bounded small local insertions, the result gives factorially controlled
connected decay across any prescribed cut.  The \(1/n!\) in the logarithmic
generating function pays the factorial.  A separate geometric issue remains:
choosing only one cut gives decay at the longest available separation, not a
product along every edge of a spanning tree.  At fixed order this is
summable, but the decay rate deteriorates with the order.  V48 records this
boundary rather than calling the one-cut estimate a Kotecky--Preiss theorem.

\subsection{Moment notation on a bipartition}
\label{sec:v4v48-notation}

Let \(I\) be a finite set with \(|I|=n\ge2\), and choose a nontrivial
bipartition
\[
 I=A\mathbin{\dot\cup}B.
\]
For \(C\subseteq I\), put
\[
 X_C:=\prod_{i\in C}X_i,\qquad X_\varnothing:=1,
 \qquad m_C:=\mathbb E X_C.
\]
Define the cut-product moment
\begin{equation}
 m_C^\otimes:=
 \mathbb E X_{C\cap A}\,
 \mathbb E X_{C\cap B}.
\label{eq:v4v48-product-moment}
\end{equation}
For a partition \(\pi\in\Pi(I)\), let
\[
 \pi_\times:=
 \{C\in\pi:C\cap A\ne\varnothing,\ C\cap B\ne\varnothing\}
\]
be its crossing blocks.  Fix once and for all an order
\(\pi_\times=(C_1,\ldots,C_r)\).  The final identity is independent of this
auxiliary order.

The moment--cumulant formula is
\begin{equation}
 \kappa(X_i:i\in I)
 =
 \sum_{\pi\in\Pi(I)}
 \mu(\pi)\prod_{C\in\pi}m_C,
 \qquad
 \mu(\pi)=(-1)^{|\pi|-1}(|\pi|-1)!.
\label{eq:v4v48-moment-cumulant}
\end{equation}

\subsection{The exact cut telescope}
\label{sec:v4v48-cut-telescope}

\begin{theorem}[Exact one-covariance cut identity]
\label{thm:v4v48-cut-identity}
For every nontrivial cut \(I=A\dot\cup B\),
\begin{align}
 \kappa(X_i:i\in I)
 &=
 \sum_{\pi\in\Pi(I)}\mu(\pi)
 \prod_{C\in\pi\setminus\pi_\times}m_C
 \sum_{j=1}^{r}
 \operatorname{Cov}\left(
   X_{C_j\cap A},X_{C_j\cap B}
 \right)\nonumber\\
 &\qquad\qquad\qquad\qquad\times
 \prod_{\ell<j}m_{C_\ell}
 \prod_{\ell>j}m_{C_\ell}^\otimes .
\label{eq:v4v48-cut-identity}
\end{align}
The inner sum is zero when \(r=0\).  Every nonzero summand contains exactly
one displayed covariance crossing the chosen cut.
\end{theorem}

\begin{proof}
Let \(\mathbb E^\otimes\) be the product of the joint marginal law of
\((X_i)_{i\in A}\) and the joint marginal law of
\((X_i)_{i\in B}\).  Its moments are exactly
\(\mathbb E^\otimes X_C=m_C^\otimes\).  Since the two nonempty families are
independent under this law,
\[
 \kappa^\otimes(X_i:i\in I)=0.
\]
Subtract its moment--cumulant formula from
\eqref{eq:v4v48-moment-cumulant}:
\begin{equation}
 \kappa(X_i:i\in I)
 =
 \sum_{\pi\in\Pi(I)}\mu(\pi)
 \left[
  \prod_{C\in\pi}m_C-\prod_{C\in\pi}m_C^\otimes
 \right].
\label{eq:v4v48-cut-difference}
\end{equation}
For a block lying wholly on one side, \(m_C=m_C^\otimes\).  Factor those
common moments.  For the ordered crossing blocks, use the elementary
product telescope
\[
 \prod_{j=1}^{r}a_j-\prod_{j=1}^{r}b_j
 =
 \sum_{j=1}^{r}(a_j-b_j)
 \prod_{\ell<j}a_\ell\prod_{\ell>j}b_\ell
\]
with \(a_j=m_{C_j}\), \(b_j=m_{C_j}^\otimes\).  Finally,
\[
 m_{C_j}-m_{C_j}^\otimes
 =
 \operatorname{Cov}\left(
 X_{C_j\cap A},X_{C_j\cap B}
 \right).
\]
Substitution in \eqref{eq:v4v48-cut-difference} proves
\eqref{eq:v4v48-cut-identity}.
\end{proof}

\begin{remark}[The V47 identity is the \(n=3\) compression]
\label{rem:v4v48-V47}
At \(n=3\), expand \eqref{eq:v4v48-cut-identity} for the cut
\(\{F,G\}\dot\cup\{H\}\) and recombine its moment terms.  The result is
exactly \eqref{eq:v4v47-covariance-regrouping}.  Thus V48 adds no fictitious
response: it exposes the same complete-product covariance at every order.
\end{remark}

\begin{firewall}[The crossing factor is a covariance, not a desired edge]
\label{fw:v4v48-actual-covariance}
The small factor in \eqref{eq:v4v48-cut-identity} is the covariance of the
two complete products \(X_{C_j\cap A}\) and \(X_{C_j\cap B}\).  It must be
estimated with their actual supports and derivatives.  It may not be
replaced by a selected pair covariance between convenient representatives
of the two sides.
\end{firewall}

\subsection{Explicit partition growth}
\label{sec:v4v48-partitions}

Let \(S(n,k)\) denote a Stirling number of the second kind and set
\begin{equation}
 a_n:=\sum_{k=1}^{n}k!S(n,k).
\label{eq:v4v48-ordered-Bell}
\end{equation}

\begin{lemma}[A usable ordered-partition bound]
\label{lem:v4v48-ordered-partition}
With \(r_0=\log(3/2)\),
\begin{equation}
 a_n\le2n!r_0^{-n}.
\label{eq:v4v48-ordered-partition-bound}
\end{equation}
\end{lemma}

\begin{proof}
The exponential generating function is
\[
 \sum_{n\ge0}a_n\frac{z^n}{n!}
 =
 \sum_{k\ge0}(e^z-1)^k
 =
 \frac1{2-e^z}.
\]
On \(|z|=r_0\),
\[
 |e^z-1|\le e^{|z|}-1=\frac12,
 \qquad
 |2-e^z|=|1-(e^z-1)|\ge\frac12.
\]
Cauchy's coefficient estimate gives
\(a_n/n!\le2r_0^{-n}\).
\end{proof}

\subsection{A local all-order cut estimate}
\label{sec:v4v48-local-bound}

Suppose \(X_i\) is supported in a skeleton set \(S_i\), and assume
\begin{equation}
 \|X_i\|_\infty\le\epsilon,
 \qquad
 \|\nabla X_i\|_{L^2(\nu_\Sigma)}\le g\epsilon
\label{eq:v4v48-small-local-factor}
\end{equation}
for fixed \(\epsilon,g>0\).  For a cut \(A\dot\cup B\), put
\begin{equation}
 \delta(A,B):=
 d\left(\bigcup_{i\in A}S_i,\bigcup_{j\in B}S_j\right).
\label{eq:v4v48-cut-distance}
\end{equation}

\begin{theorem}[Factorial cut clustering from one actual covariance]
\label{thm:v4v48-factorial-cut}
Under the skeleton Witten estimate
\eqref{eq:v4v45-skeleton-Witten},
\begin{equation}
 \boxed{
 |\kappa(X_i:i\in I)|
 \le
 \frac{C_\Sigma g^2}{2}\,
 n^2a_n\,\epsilon^n
 e^{-m_\Sigma\delta(A,B)}.}
\label{eq:v4v48-factorial-cut}
\end{equation}
Consequently
\begin{equation}
 \frac1{n!}|\kappa(X_i:i\in I)|
 \le
 C_\Sigma g^2n^2
 \left(\frac{\epsilon}{r_0}\right)^n
 e^{-m_\Sigma\delta(A,B)}.
\label{eq:v4v48-factorial-cancelled}
\end{equation}
\end{theorem}

\begin{proof}
For \(D\subseteq I\), the product rule and
\eqref{eq:v4v48-small-local-factor} give
\begin{equation}
 \|\nabla X_D\|_2
 \le |D|g\epsilon^{|D|}.
\label{eq:v4v48-product-gradient}
\end{equation}
Indeed each differentiated factor contributes at most
\(g\epsilon\), and the remaining \(|D|-1\) factors contribute
\(\epsilon^{|D|-1}\).

Fix a partition \(\pi\) and one crossing block \(C\).  Lemma
\ref{lem:v4v45-separated-covariance}, applied to the two complete products,
gives
\begin{align}
 \left|
 \operatorname{Cov}(X_{C\cap A},X_{C\cap B})
 \right|
 &\le
 C_\Sigma e^{-m_\Sigma\delta(A,B)}
 |C\cap A|\,|C\cap B|\,g^2\epsilon^{|C|}.
\label{eq:v4v48-cross-product-covariance}
\end{align}
Every other moment in the same telescoping summand is bounded by the
corresponding power of \(\epsilon\).  Thus its total power is
\(\epsilon^n\).

The crossing blocks are disjoint, so
\[
 \sum_{C\in\pi_\times}|C\cap A|\,|C\cap B|
 \le
 \left(\sum_C|C\cap A|\right)
 \left(\sum_C|C\cap B|\right)
 \le |A||B|\le\frac{n^2}{4}.
\]
For partitions with \(k\) blocks, the absolute Möbius coefficient is
\((k-1)!\le k!\).  Summing
\eqref{eq:v4v48-cut-identity} in absolute value therefore gives
\[
 |\kappa(X_i:i\in I)|
 \le
 \frac{C_\Sigma g^2n^2}{4}\epsilon^n
 e^{-m_\Sigma\delta(A,B)}
 \sum_{k=1}^n(k-1)!S(n,k).
\]
Bound the final sum by \(a_n\).  This is slightly stronger than
\eqref{eq:v4v48-factorial-cut}; the displayed constant leaves room for the
fixed real/imaginary polarization factor in
\eqref{eq:v4v45-HS-covariance}.  Insert
\eqref{eq:v4v48-ordered-partition-bound} and divide by \(n!\) to obtain
\eqref{eq:v4v48-factorial-cancelled}.
\end{proof}

\begin{corollary}[Application to local normalization perturbations]
\label{cor:v4v48-local-normalizations}
Let \(F_b=1+Y_b\) be strictly local scalar block normalizations.  If the
\(Y_b\) satisfy \eqref{eq:v4v48-small-local-factor}, then for every
\(n\ge2\),
\begin{equation}
 \kappa(F_{b_1},\ldots,F_{b_n})
 =
 \kappa(Y_{b_1},\ldots,Y_{b_n}),
\label{eq:v4v48-constant-removal}
\end{equation}
and the latter obeys
\eqref{eq:v4v48-factorial-cut} for every cut of the block labels.
\end{corollary}

\begin{proof}
A joint cumulant of order at least two vanishes when one argument is
constant.  Multilinearity proves \eqref{eq:v4v48-constant-removal}; then
apply Theorem~\ref{thm:v4v48-factorial-cut}.
\end{proof}

\subsection{Fixed-order geometry and the all-order loss}
\label{sec:v4v48-geometry}

Let \(x_1,\ldots,x_n\) be block centers and let \(T_{\min}\) be a
minimum-length spanning tree in the complete metric graph on these centers.
Write \(L_{\min}\) for its total length and \(\ell_{\max}\) for its longest
edge.

\begin{lemma}[A longest-tree-edge cut]
\label{lem:v4v48-tree-cut}
Deleting a longest edge of \(T_{\min}\) gives a cut \(A\dot\cup B\) such that
\begin{equation}
 d(\{x_i:i\in A\},\{x_j:j\in B\})
 \ge\ell_{\max}
 \ge\frac{L_{\min}}{n-1}.
\label{eq:v4v48-tree-cut}
\end{equation}
\end{lemma}

\begin{proof}
If a cross-cut pair had distance smaller than the deleted edge, inserting
that pair would reconnect the two components and produce a shorter spanning
tree, contradicting minimality.  The second inequality follows because
\(T_{\min}\) has \(n-1\) edges.
\end{proof}

For blocks of fixed radius \(r_*\), the support distance of this cut is at
least \(\ell_{\max}-2r_*\).  Thus
Theorem~\ref{thm:v4v48-factorial-cut} gives, after changing its constant,
\begin{equation}
 |\kappa(X_1,\ldots,X_n)|
 \le C_ng^2\epsilon^n
 \exp\left(-\frac{m_\Sigma}{n-1}L_{\min}\right)
\label{eq:v4v48-fixed-order-tree-length}
\end{equation}
for each fixed \(n\).  Quasi-local V44 factors obey the same conclusion at
fixed \(n\): truncate every factor at a fixed fraction of
\(\ell_{\max}\), use the V45 \(L^2\) tail for the replacement error, and
apply the local estimate to the truncations.

\begin{corollary}[Fixed-order thermodynamic summability]
\label{cor:v4v48-fixed-order-sum}
On a block lattice of polynomial volume growth, for every fixed \(n\),
\begin{equation}
 \sup_{b_1}
 \sum_{b_2,\ldots,b_n}
 |\kappa(Y_{b_1},\ldots,Y_{b_n})|
 <\infty.
\label{eq:v4v48-fixed-order-sum}
\end{equation}
\end{corollary}

\begin{proof}
The minimal tree length controls the maximal distance from \(b_1\).  At
fixed \(n\), the number of \((n-1)\)-tuples with maximal radius in
\([R,R+1)\) grows polynomially in \(R\).  Equation
\eqref{eq:v4v48-fixed-order-tree-length} has a positive exponential rate
for that fixed \(n\), so the sum converges.
\end{proof}

\begin{firewall}[The one-cut rate deteriorates with cumulant order]
\label{fw:v4v48-rate-loss}
The exponent in \eqref{eq:v4v48-fixed-order-tree-length} is
\(m_\Sigma/(n-1)\).  Summing positions before fixing \(n\) then produces
spatial constants growing roughly like a power of \(n\) for each of the
\(n-1\) free vertices.  The factorial cancellation in
\eqref{eq:v4v48-factorial-cancelled} does not automatically pay that
geometric growth in four dimensions.  Hence V48 proves fixed-order
summability and explicit analytic smallness, but not an all-order
Kotecky--Preiss sum.
\end{firewall}

\begin{assessment}[What a uniform tree compiler must improve]
\label{ass:v4v48-tree-compiler}
A sufficient next bound would retain one exponential factor for every edge
of a spanning tree,
\[
 |\kappa(Y_{b_1},\ldots,Y_{b_n})|
 \le n!C^n\epsilon^n
 \sum_{T\text{ on }[n]}
 \prod_{\{i,j\}\in T}e^{-m d(b_i,b_j)}.
\]
After rooting a tree, each free block position would then be summed against
one fixed exponential kernel, uniformly in \(n\).  V48 supplies one actual
cross-cut covariance and the exact partition coefficients.  The missing
step is a recursion which exposes the remaining \(n-2\) actual connections
without differentiating the Witten inverse or duplicating a covariance.
\end{assessment}

\subsection{V48 result ledger and V49 target}
\label{sec:v4v48-conclusion}

V48 proves the exact all-order cut identity
\eqref{eq:v4v48-cut-identity}.  Each term contains one covariance of
complete products across the chosen cut.  The skeleton Witten estimate then
gives factorially controlled cut clustering for small bounded local
normalization perturbations, with the complete ordered-partition constant
made explicit.  The logarithmic \(1/n!\) pays that factorial, and every
fixed cumulant order is thermodynamically summable, including after V44--V45
quasi-local truncation.

V49 should seek an exact recursive form of the cut telescope which assigns
successive actual cross-cut covariances to distinct edges of a labelled
tree.  The recursion must preserve the complete product on each side until
its covariance is used.  If such a recursion necessarily reuses a Witten
carrier, the next upstream alternative is a direct forest interpolation of
the positive skeleton generating function with the V48 product functional
as its decoupled endpoint.


\clearpage
\section{V49: conditional collision partitions and the first exact tree recursion}
\label{sec:v4v49}

V48 exposes one actual covariance across any chosen cut, but estimating the
remaining moment factors independently loses the other edges of a connected
tree.  The V47 law of total cumulance contains a different recursion.  Under
one color conditioning, supports lying in distinct conditional components
have no joint conditional cumulant.  When components first meet, the
conditional cumulant created inside the merged component becomes one new
argument of an outer cumulant.  Its order is smaller by the number of
internal connections already acquired.

This note makes that collision algebra exact.  It proves the general
component-restricted partition formula, computes every fourth-order merger
case, and identifies the analytic carrier norm needed at a binary collision.
The result is an exact ancestry recursion: a binary conditional covariance
is created once and reduces the outer order by one.  A full spatial tree
bound still requires a uniform estimate for each such collision carrier
through later color sweeps.

\subsection{Conditional component partitions}
\label{sec:v4v49-components}

Let \(I\) be a finite label set and let
\[
 \mathcal C=\{C^{(1)},\ldots,C^{(s)}\}
\]
be a partition of \(I\).  Assume that, conditionally on a
\(\sigma\)-field \(\mathscr G\), the random families
\[
 (X_i)_{i\in C^{(1)}},\ldots,(X_i)_{i\in C^{(s)}}
\]
are mutually independent.  Write \(\pi\preceq\mathcal C\) when every block
of the partition \(\pi\in\Pi(I)\) is contained in one component of
\(\mathcal C\).

\begin{theorem}[Conditional collision-partition identity]
\label{thm:v4v49-collision-partition}
Under the preceding conditional independence,
\begin{equation}
 \boxed{
 \kappa(X_i:i\in I)
 =
 \sum_{\substack{\pi\in\Pi(I)\\\pi\preceq\mathcal C}}
 \kappa\left(
   \kappa(X_i:i\in B\mid\mathscr G):B\in\pi
 \right).}
\label{eq:v4v49-collision-partition}
\end{equation}
Thus only conditional cumulants contained in actual conditional components
survive.
\end{theorem}

\begin{proof}
Start from the law of total cumulance
\eqref{eq:v4v47-total-cumulance}.  If a block \(B\) meets two components of
\(\mathcal C\), its arguments split into two nonempty conditionally
independent families, so
\[
 \kappa(X_i:i\in B\mid\mathscr G)=0.
\]
Every partition not refining \(\mathcal C\) therefore contributes zero.
The remaining partitions give \eqref{eq:v4v49-collision-partition}.
\end{proof}

\begin{corollary}[No collision means pure projection]
\label{cor:v4v49-no-collision}
If every component of \(\mathcal C\) is a singleton, then
\begin{equation}
 \kappa(X_i:i\in I)
 =
 \kappa(M_i:i\in I),
 \qquad
 M_i:=\mathbb E[X_i\mid\mathscr G].
\label{eq:v4v49-no-collision}
\end{equation}
\end{corollary}

\begin{proof}
Only the singleton partition refines a singleton component partition.
\end{proof}

\begin{firewall}[A merger is determined by conditional incidence]
\label{fw:v4v49-incidence}
A block enters the sum in \eqref{eq:v4v49-collision-partition} only if all
of its labels lie in one actual conditional component.  Spatial closeness
after an estimate, overlap of coarse envelopes, or a convenient spanning
tree does not create such a block.  The component partition must be fixed
before the conditional cumulants are expanded.
\end{firewall}

\subsection{Complete fourth-order merger formulae}
\label{sec:v4v49-fourth}

For four variables put
\[
 M_i:=\mathbb E[X_i\mid\mathscr G],\qquad
 C_{ij}:=\operatorname{Cov}(X_i,X_j\mid\mathscr G).
\]

\begin{proposition}[One binary collision at order four]
\label{prop:v4v49-one-pair}
If the conditional component partition is
\(\{\{1,2\},\{3\},\{4\}\}\), then
\begin{equation}
 \boxed{
 \kappa_4(X_1,X_2,X_3,X_4)
 =
 \kappa_3(C_{12},M_3,M_4)
 +
 \kappa_4(M_1,M_2,M_3,M_4).}
\label{eq:v4v49-one-pair}
\end{equation}
\end{proposition}

\begin{proof}
Exactly two partitions of \(\{1,2,3,4\}\) refine the stated component
partition while having no zero conditional block:
\[
 \{\{1,2\},\{3\},\{4\}\},
 \qquad
 \{\{1\},\{2\},\{3\},\{4\}\}.
\]
Their outer cumulants are the two terms in
\eqref{eq:v4v49-one-pair}.
\end{proof}

\begin{proposition}[Two simultaneous binary collisions]
\label{prop:v4v49-two-pairs}
If the conditional component partition is
\(\{\{1,2\},\{3,4\}\}\), then
\begin{align}
 \kappa_4(X_1,X_2,X_3,X_4)
 &=
 \operatorname{Cov}(C_{12},C_{34})\nonumber\\
 &\quad+\kappa_3(C_{12},M_3,M_4)
 +\kappa_3(M_1,M_2,C_{34})\nonumber\\
 &\quad+\kappa_4(M_1,M_2,M_3,M_4).
\label{eq:v4v49-two-pairs}
\end{align}
\end{proposition}

\begin{proof}
The refining partitions are
\[
 \{12,34\},\qquad
 \{12,3,4\},\qquad
 \{1,2,34\},\qquad
 \{1,2,3,4\},
\]
where commas separate blocks.  Insert their conditional cumulants in
\eqref{eq:v4v49-collision-partition}.
\end{proof}

\begin{proposition}[A three-label component and one spectator]
\label{prop:v4v49-triple-component}
If the conditional component partition is
\(\{\{1,2,3\},\{4\}\}\), write
\[
 C_{123}:=\kappa_3(X_1,X_2,X_3\mid\mathscr G).
\]
Then
\begin{align}
 \kappa_4(X_1,X_2,X_3,X_4)
 &=
 \operatorname{Cov}(C_{123},M_4)\nonumber\\
 &\quad+
 \kappa_3(C_{12},M_3,M_4)
 +\kappa_3(C_{13},M_2,M_4)
 +\kappa_3(C_{23},M_1,M_4)\nonumber\\
 &\quad+
 \kappa_4(M_1,M_2,M_3,M_4).
\label{eq:v4v49-triple-component}
\end{align}
\end{proposition}

\begin{proof}
The partitions refining the component split are obtained by choosing an
arbitrary partition of \(\{1,2,3\}\) and adjoining the singleton
\(\{4\}\).  The one-block, three pair-plus-singleton, and all-singleton
partitions inside the triple give exactly the displayed terms.
\end{proof}

\begin{assessment}[A binary collision acquires one connection]
\label{ass:v4v49-binary-acquisition}
In \eqref{eq:v4v49-one-pair}, the carrier \(C_{12}\) is an actual
conditional covariance created inside the first merged component.  It
replaces two arguments by one and therefore lowers the outer order from four
to three.  If that outer third cumulant later acquires two further binary
collision carriers, the resulting ancestry has three distinct connections,
the edge count of a tree on four labels.  The second term in
\eqref{eq:v4v49-one-pair} is the no-collision branch and must continue
through the next conditional projection.
\end{assessment}

\subsection{The exact recursive compiler}
\label{sec:v4v49-recursion}

Consider alternating conditional stages \(r=0,1,\ldots\).  At stage \(r\),
let \(\mathscr G_r\) be the retained color and let
\(\mathcal C_r(\pi)\) be the actual conditional component partition of the
current carrier arguments indexed by the blocks of a partition \(\pi\) of
the original labels.  For a current block \(B\), define
\[
 K_{r,B}:=
 \kappa(X_i^{(r)}:i\in B\mid\mathscr G_r),
\]
where singleton \(K_{r,\{i\}}\) is the conditional mean.

\begin{theorem}[Finite collision-history expansion]
\label{thm:v4v49-collision-history}
After any finite number \(R\) of conditional stages, the original cumulant
is an exact finite sum over collision histories.  At each stage a history:
\begin{enumerate}
\item chooses a partition refining the actual conditional component
      partition;
\item replaces every chosen block by its complete conditional cumulant; and
\item takes the outer cumulant of the resulting block carriers.
\end{enumerate}
Every history term is indexed by a nested sequence of partitions
\[
 \pi_0\preceq\pi_1\preceq\cdots\preceq\pi_R
\]
of the original labels.  Whenever a block of size \(s\ge2\) is created, the
number of current arguments drops by \(s-1\).  A history ending in one
argument has total order drop
\begin{equation}
 \sum_{\text{created blocks }B}(|B|-1)=n-1.
\label{eq:v4v49-order-drop}
\end{equation}
\end{theorem}

\begin{proof}
Apply Theorem~\ref{thm:v4v49-collision-partition} at the first stage.  Each
outer cumulant is again a cumulant of \(\mathscr G_0\)-measurable carrier
variables.  Apply the same theorem to their actual component partition at
the next color conditioning.  Iteration for \(R\) stages gives a finite sum
because every stage has finitely many set partitions.

Each conditional block is a union of blocks already present, so the
partitions coarsen and form a nested sequence.  Replacing \(s\) arguments by
one lowers their number by \(s-1\).  If a history starts with \(n\) singleton
arguments and ends with one, telescoping the argument count gives
\eqref{eq:v4v49-order-drop}.
\end{proof}

\begin{corollary}[Binary histories have the exact tree edge count]
\label{cor:v4v49-binary-tree-count}
If every nonsingleton block created in a terminal history has size two, then
the history contains exactly \(n-1\) binary conditional covariance
carriers.  Joining the two child clusters at each creation gives a labelled
tree on the original \(n\) labels.
\end{corollary}

\begin{proof}
Every creation contributes one to
\eqref{eq:v4v49-order-drop}, so there are \(n-1\) creations.  Starting with
singleton vertices, joining two previously disjoint clusters never creates
a cycle.  A cycle-free graph on \(n\) vertices with \(n-1\) edges is a tree.
\end{proof}

\begin{firewall}[A multi-label carrier is a hyperedge, not several free
pair edges]
\label{fw:v4v49-hyperedge}
If a block of size \(s>2\) is created in one conditional component, its
carrier is the complete conditional \(s\)-cumulant.  Its order drop is
\(s-1\), but this does not license replacing it by \(s-1\) pair
covariances.  Such a replacement requires its own proved internal tree
bound.  The collision-history expansion preserves the hyperedge until that
bound is acquired.
\end{firewall}

\subsection{The analytic carrier interface}
\label{sec:v4v49-carrier-interface}

The algebra becomes a spatial estimate only after every created carrier is
bounded in the norm used by the next conditional channel.  For a binary
collision between current carriers \(U,V\), define
\[
 \mathcal C_r(U,V):=
 \operatorname{Cov}(U,V\mid\mathscr G_r).
\]

\begin{definition}[Binary conditional collision carrier bound]
\label{def:v4v49-BCCC}
A stage has \(BCCC(\Gamma,w)\) if every actual binary collision satisfies
\begin{align}
 \|\mathcal C_r(U,V)\|_{2}
 &\le
 \Gamma\,w(\operatorname{supp}U,\operatorname{supp}V)
 \|U\|_{\mathfrak h}\|V\|_{\mathfrak h},
\label{eq:v4v49-BCCC-L2}\\
 \|\mathcal C_r(U,V)\|_{\infty}
 &\le
 \Gamma\,w(\operatorname{supp}U,\operatorname{supp}V)
 \|U\|_{\mathfrak h,\infty}\|V\|_{\mathfrak h,\infty},
\label{eq:v4v49-BCCC-Linfty}
\end{align}
where \(w(X,Y)\) is a summable spatial kernel and the carrier norms are
stable under the subsequent normalized conditional projections.
\end{definition}

\begin{theorem}[Sufficient binary collision-tree compiler]
\label{thm:v4v49-binary-compiler}
Assume:
\begin{enumerate}
\item every terminal collision history is binary;
\item every stage satisfies \(BCCC(\Gamma,w)\);
\item for disjoint label clusters \(A,B\), the collision kernel satisfies
\begin{equation}
 w\left(\bigcup_{i\in A}S_i,\bigcup_{j\in B}S_j\right)
 \le\sum_{i\in A}\sum_{j\in B}w(S_i,S_j);
\label{eq:v4v49-kernel-subadditivity}
\end{equation}
\item singleton projections obey the V41 centered contraction and the V42
      coefficient-fibre stability; and
\item the sum of absolute partition coefficients generated at each merger
      is bounded by \(K^n\).
\end{enumerate}
Then
\begin{equation}
 |\kappa(X_1,\ldots,X_n)|
 \le
 K^n\Gamma^{n-1}
 \left(\prod_{i=1}^n\|X_i\|_{\mathfrak h,\infty}\right)
 \sum_{T\text{ on }[n]}
 \prod_{\{i,j\}\in T}w(S_i,S_j),
\label{eq:v4v49-binary-compiler}
\end{equation}
with additional powers of \(q_{\rm H}\) for every pre-collision sweep.
\end{theorem}

\begin{proof}
Expand by Theorem~\ref{thm:v4v49-collision-history}.  By
Corollary~\ref{cor:v4v49-binary-tree-count}, each terminal history has
\(n-1\) distinct binary creations and determines a labelled tree.  Apply
\eqref{eq:v4v49-BCCC-L2} or
\eqref{eq:v4v49-BCCC-Linfty} once at each creation and stability at every
later projection.  At a merger of two clusters, expand its kernel by
\eqref{eq:v4v49-kernel-subadditivity} and select one original-label edge
crossing that merger.  The selected edges connect previously disjoint
clusters, so the \(n-1\) selections form a labelled tree.  This produces one
factor \(\Gamma w(S_i,S_j)\) for each tree edge.  No carrier is used twice.
Sum histories producing the same labelled tree and use the coefficient
hypothesis.  Before a collision, apply
Proposition~\ref{prop:v4v47-precollision} to retain the corresponding
centered sweep powers.
\end{proof}

\begin{assessment}[The compiler is sufficient but not yet physical]
\label{ass:v4v49-BCCC-open}
Theorem~\ref{thm:v4v49-binary-compiler} proves that the collision recursion
would supply the uniform tree geometry missing in V48.  Its hypotheses are
not yet fully acquired.  In particular, physical support collisions can
join more than two clusters at the same stage, and V41 gives an \(L^2\)
contraction for conditional means rather than the carrier bounds
\eqref{eq:v4v49-BCCC-L2}--\eqref{eq:v4v49-BCCC-Linfty}.  V49 therefore
establishes the exact algebra and the required analytic interface, not the
all-order Standard Model pressure.
\end{assessment}

\subsection{V49 result ledger and V50 audit target}
\label{sec:v4v49-conclusion}

V49 proves that the law of total cumulance restricts exactly to partitions
inside actual conditional components.  It computes all fourth-order merger
types and proves the finite collision-history expansion.  A binary collision
creates one actual conditional covariance, lowers the outer order by one,
and terminal binary histories have exactly the edge count and acyclic
ancestry of a labelled tree.  This is the exact recursive structure absent
from the one-cut V48 estimate.

The analytic binary-carrier norm and the treatment of simultaneous
multi-label collisions remain open.  V50 must begin with the scheduled
read-only YM/NS audit.  It should then test \(BCCC(\Gamma,w)\) at the first
physical even/odd collision using the exact conditional product law.  If an
\(L^\infty\) carrier bound fails for the unbounded Higgs field, the route
should retain an \(L^2\)-\(L^p\) norm ladder and determine its exact cost
before attempting the all-order history sum.


\clearpage
\section{V50: companion audit and the first physical collision carrier}
\label{sec:v4v50}

V49 derives an exact collision-history recursion and isolates the binary
conditional collision carrier bound needed to turn that recursion into a
uniform tree estimate.  This scheduled audit version first records the
current companion files.  It then tests the carrier bound at one physical
even/odd collision.  The value of the conditional covariance has the
required uniform exponential estimate.  Its derivative with respect to the
retained color, however, is a complete conditional third cumulant.  Hence
the value bound supplies one edge but does not by itself put the new carrier
in the derivative norm required for the next edge.

\subsection{Scheduled read-only companion audit}
\label{sec:v4v50-audit}

The newest active companion files present when the V50 direction was fixed
were
\begin{align*}
 &\texttt{Renewal\_Yang\_Mills\_Mass\_Gap\_Research\_Notes\_V66.tex},\\
 &\qquad 1\,171\,436\ \text{bytes},\qquad
 \texttt{SHA--256 }\\
 &\qquad
 \texttt{6C32338BEB3C96FAF36C03BC01C17E9C3DD81B0E02CB976FDF76D529AF97B9D3},
 \\
 &\texttt{Renewal\_NS\_Stability\_Research\_Notes\_V31.tex},\\
 &\qquad 887\,537\ \text{bytes},\qquad
 \texttt{SHA--256 }\\
 &\qquad
 \texttt{F1121B62238AD5491BB7CF03635EA9934942EDF573ED8FAAA9EE79E1D38A6696}.
\end{align*}

The Yang--Mills note has advanced from the V61 file audited in V45.  V66
proves source occurrence and central commutation for its surviving bank
chains before applying their estimates.  For the remaining lateral branch,
its V67 programme requires one exact conditional square or covariance
identity containing both relevant occurrences and explicitly forbids
multiplying their separate capacities.  The Navier--Stokes file is unchanged
from V45.  Its terminal V31 ledger continues to distinguish a normalized
probability mark from the exact first moment required to remove that mark.

\begin{assessment}[Type-correct V50 transfer]
\label{ass:v4v50-audit-transfer}
Only proof order transfers.
\begin{enumerate}
\item A collision edge is paid only by the actual conditional cumulant
      created in \eqref{eq:v4v49-collision-partition}.  Separate decay
      capacities of its child clusters may not be multiplied unless they
      occur together in that carrier.
\item A value estimate on one normalized carrier does not pay the derivative
      norm required when that carrier participates in a later collision.
\item The exponential spatial kernel and the exact number of later carrier
      uses must remain visible until the collision history is summed.
\end{enumerate}
No Yang--Mills row estimate and no Navier--Stokes heat estimate is imported
into the Higgs law.
\end{assessment}

\subsection{One conditional covariance inside an actual component}
\label{sec:v4v50-first-carrier}

Fix one color and condition on the opposite color and the exterior skeleton,
collectively denoted by \(\eta\).  Let \(\mu_D^\eta\) be the conditional
positive Higgs law in one actual conditional component \(D\).  Assume its
one-form inverse obeys, uniformly in \(\eta\),
\begin{equation}
 \|P_X(\mathcal L_{1,D}^{\eta})^{-1}P_Y\|
 \le C_D e^{-m_Dd(X,Y)}.
\label{eq:v4v50-conditional-Witten}
\end{equation}
This is the conditional V38 Witten estimate on the chosen component.

For two coefficient functions \(U,V\) supported in \(X,Y\subseteq D\), set
\begin{equation}
 \mathcal C^\eta(U,V)
 :=
 \operatorname{Cov}_{\mu_D^\eta}(U,V).
\label{eq:v4v50-conditional-carrier}
\end{equation}

\begin{theorem}[The first collision has a uniform value bound]
\label{thm:v4v50-first-BCCC}
Suppose
\begin{equation}
 \sup_\eta\|\nabla U\|_{L^2(\mu_D^\eta)}\le A_U,
 \qquad
 \sup_\eta\|\nabla V\|_{L^2(\mu_D^\eta)}\le A_V.
\label{eq:v4v50-uniform-gradients}
\end{equation}
Then
\begin{equation}
 \boxed{
 \|\mathcal C^\eta(U,V)\|_{L^\infty(d\eta)}
 \le
 C_De^{-m_Dd(X,Y)}A_UA_V.}
\label{eq:v4v50-first-BCCC}
\end{equation}
If \(U,V\) belong to distinct conditional components, their carrier is
identically zero.
\end{theorem}

\begin{proof}
For fixed \(\eta\), apply the conditional Helffer--Sjostrand identity:
\[
 \mathcal C^\eta(U,V)
 =
 \mathbb E_{\mu_D^\eta}
 \left\langle
  \nabla U,(\mathcal L_{1,D}^{\eta})^{-1}\nabla V
 \right\rangle.
\]
Insert the support projections, use
\eqref{eq:v4v50-conditional-Witten}, and apply Cauchy--Schwarz in
\(\mu_D^\eta\).  The constants and gradient bounds are uniform in the
conditioned field, so taking the supremum proves
\eqref{eq:v4v50-first-BCCC}.  Distinct components are independent in the
conditional product law, so their covariance vanishes exactly.
\end{proof}

\begin{corollary}[One acquired collision edge]
\label{cor:v4v50-one-edge}
For the value component of the V49 carrier norm, Theorem
\ref{thm:v4v50-first-BCCC} verifies
\[
 \Gamma=C_DA_UA_V,\qquad
 w(X,Y)=e^{-m_Dd(X,Y)}
\]
at the first binary collision.  The same bound holds in every finite
Grassmann coefficient component and in its Hilbert direct sum, as in V42.
\end{corollary}

\begin{firewall}[Uniform conditional gradients are an actual hypothesis]
\label{fw:v4v50-gradient-hypothesis}
The lower Hessian bound controls conditional variances, but it does not make
\(\sup_\eta\mathbb E_{\mu_D^\eta}|\nabla U|^2\) finite for an arbitrary
polynomial insertion under unbounded boundary fields.  Theorem
\ref{thm:v4v50-first-BCCC} applies directly to the bounded-derivative
good-field insertion.  Polynomial Yukawa coefficients require their V34
good/tail split or a separate uniform moment estimate; neither is hidden in
\eqref{eq:v4v50-uniform-gradients}.
\end{firewall}

\subsection{The derivative needed at the next collision}
\label{sec:v4v50-carrier-derivative}

Let a retained-color direction \(h\) change the conditional action through
the score
\[
 G_h:=D_\eta S_D^\eta[h].
\]
Allow \(U,V\) also to have explicit \(\eta\)-dependence.  Write
\(D_\eta U[h]\) and \(D_\eta V[h]\) for those explicit derivatives.

\begin{theorem}[Exact derivative of a conditional covariance carrier]
\label{thm:v4v50-carrier-derivative}
Whenever differentiation under the conditional integral is valid,
\begin{align}
 D_\eta\mathcal C^\eta(U,V)[h]
 &=
 \operatorname{Cov}_{\mu_D^\eta}(D_\eta U[h],V)
 +\operatorname{Cov}_{\mu_D^\eta}(U,D_\eta V[h])\nonumber\\
 &\quad-
 \kappa_{3,\mu_D^\eta}(U,V,G_h).
\label{eq:v4v50-carrier-derivative}
\end{align}
In particular, if \(U,V\) have no explicit retained-color dependence,
\begin{equation}
 D_\eta\mathcal C^\eta(U,V)[h]
 =-\kappa_{3,\mu_D^\eta}(U,V,G_h).
\label{eq:v4v50-pure-carrier-derivative}
\end{equation}
\end{theorem}

\begin{proof}
Normalized differentiation gives, for any conditional observable \(A\),
\[
 D_\eta\mathbb E^\eta A[h]
 =
 \mathbb E^\eta D_\eta A[h]
 -\operatorname{Cov}_{\mu_D^\eta}(A,G_h).
\]
Apply this identity to
\(\mathbb E^\eta(UV)-\mathbb E^\eta U\,\mathbb E^\eta V\) and collect the
two explicit derivatives.  The remaining score terms are
\[
 -\left[
 \operatorname{Cov}(UV,G_h)
 -\mathbb EU\,\operatorname{Cov}(V,G_h)
 -\mathbb EV\,\operatorname{Cov}(U,G_h)
 \right].
\]
The bracket is the third cumulant by
\eqref{eq:v4v47-covariance-regrouping}.
\end{proof}

\begin{corollary}[The value bound does not close the recursive carrier norm]
\label{cor:v4v50-value-not-derivative}
The estimate \eqref{eq:v4v50-first-BCCC} controls
\(\mathcal C^\eta(U,V)\) as a retained-color function.  To use that function
as an input of a later Witten covariance, one needs its retained-color
gradient.  Equation \eqref{eq:v4v50-pure-carrier-derivative} shows that this
gradient is a conditional third cumulant and is not bounded by
\eqref{eq:v4v50-first-BCCC} alone.
\end{corollary}

\begin{proof}
The next Helffer--Sjostrand estimate contains the \(L^2\) norm of the
gradient of each input carrier.  Substitute
\eqref{eq:v4v50-carrier-derivative}; the third-cumulant term remains even
when the two explicit derivatives vanish.
\end{proof}

\begin{firewall}[No multiplication of two child capacities]
\label{fw:v4v50-no-capacity-product}
The right side of \eqref{eq:v4v50-pure-carrier-derivative} is one complete
conditional cumulant with the actual score \(G_h\).  Replacing it by the
product of the value bound \eqref{eq:v4v50-first-BCCC} and a separate
boundary-response bound would assert the joint occurrence which must still
be proved.  This is the precise SM analogue of the residual YM V66
instruction to derive one covariance containing both banks.
\end{firewall}

\subsection{Why an \(L^2\)-\(L^p\) ladder is not yet enough}
\label{sec:v4v50-norm-ladder}

Hölder gives, for centered conditional variables,
\[
 |\kappa_3(U,V,G_h)|
 \le
 \|U-\mathbb E^\eta U\|_{p}
 \|V-\mathbb E^\eta V\|_{q}
 \|G_h-\mathbb E^\eta G_h\|_{r},
 \qquad
 \frac1p+\frac1q+\frac1r=1.
\]
Strong log-concavity and polynomial local scores can provide finite moments
at every fixed exponent on a good chart.  This yields a finite derivative
norm at any fixed collision depth.  It does not yet give:
\begin{enumerate}
\item an exponential spatial factor for each new collision;
\item constants with controlled growth as the collision order increases;
      or
\item stability of the resulting carrier under the next derivative.
\end{enumerate}

\begin{assessment}[The first-edge result and the remaining recursive gate]
\label{ass:v4v50-first-edge-status}
V50 acquires the value part of \(BCCC\) for one physical binary collision.
The derivative part returns exactly the next conditional cumulant, with no
missing term.  A closed all-order norm therefore needs either a uniform
conditional strong-mixing estimate for bounded carriers which does not ask
for their gradients, or a generating-function estimate that controls all
score derivatives simultaneously.  A fixed-\(p\) Hölder estimate cannot be
iterated without tracking its order growth.
\end{assessment}

\subsection{V50 result ledger and V51 target}
\label{sec:v4v50-conclusion}

V50 records the exact active YM V66 and unchanged NS V31 fingerprints and
imports only their proof-order lessons.  It proves a uniform exponentially
decaying \(L^\infty\) value bound for the first actual conditional covariance
carrier, including exact vanishing between distinct conditional components.
It then proves the full retained-color derivative identity
\eqref{eq:v4v50-carrier-derivative}.  That identity shows why the value
bound alone does not satisfy the stable carrier norm assumed by the V49
tree compiler.

V51 should move upstream from individual carrier derivatives to a
conditional exponential generating function.  The target is a uniform
analytic bound on
\[
 \log\mathbb E_{\mu_D^\eta}
 \exp\left(\sum_j t_jU_j\right)
\]
in a polydisc whose radius is independent of the collision depth, together
with a boundary derivative bound before coefficient extraction.  Cauchy
estimates would then control every conditional cumulant and its next score
derivative in one norm.  If the radius shrinks with the number of supports,
the all-order tree route remains open and must return to a direct positive
polymer expansion.


\clearpage
\section{V51: a component-local analytic cumulant generator}
\label{sec:v4v51}

V50 shows that differentiating one conditional covariance carrier creates
the next complete conditional cumulant.  Estimating carriers one derivative
at a time would therefore reproduce the same problem at every order.  The
upstream object is the logarithm of the conditional exponential generating
function.  One analytic estimate on that logarithm controls all conditional
cumulants, and one boundary derivative estimate controls the next score
mark at every order.

The important locality is not a polydisc with a budget summed over the whole
lattice.  Under the exact conditional product law, the logarithm is a sum
over conditional components.  Its zero-free domain is therefore governed
by the largest insertion load in a single component.  The analytic radius
is independent of the total number of components and total volume.

\subsection{Conditional product and component loads}
\label{sec:v4v51-product}

Fix one color and the exterior skeleton.  Denote the retained data by
\(\eta\), and suppose the variables integrated at this stage split into
conditional components \(D_a\), \(a\in\mathcal A\), with
\begin{equation}
 \mu^\eta=\bigotimes_{a\in\mathcal A}\mu_a^\eta.
\label{eq:v4v51-conditional-product}
\end{equation}
Let \(U_1,\ldots,U_N\) be bounded complex scalar insertions.  Each \(U_j\)
depends on exactly one component, denoted \(a(j)\).  Put
\begin{align}
 u_j&:=\sup_\eta\|U_j\|_{L^\infty(\mu_{a(j)}^\eta)},\label{eq:v4v51-uj}\\
 W_a(t)&:=\sum_{j:a(j)=a}t_jU_j,\label{eq:v4v51-Wa}\\
 A_a(t)&:=\sum_{j:a(j)=a}|t_j|u_j.\label{eq:v4v51-load}
\end{align}
For a fixed number
\begin{equation}
 0<a_0<\log2,
\label{eq:v4v51-a0}
\end{equation}
define the component-load domain
\begin{equation}
 \mathfrak D_{a_0}:=
 \left\{t\in\mathbb C^N:
 \sup_{a\in\mathcal A}A_a(t)\le a_0\right\}.
\label{eq:v4v51-domain}
\end{equation}

\begin{theorem}[Exact component factorization of the generator]
\label{thm:v4v51-generator-factorization}
Define
\begin{align}
 Z(t,\eta)&:=
 \mathbb E_{\mu^\eta}
 \exp\left(\sum_{j=1}^Nt_jU_j\right),\label{eq:v4v51-Z}\\
 Z_a(t,\eta)&:=
 \mathbb E_{\mu_a^\eta}e^{W_a(t)}.\label{eq:v4v51-Za}
\end{align}
Then
\begin{equation}
 Z(t,\eta)=\prod_{a\in\mathcal A}Z_a(t,\eta).
\label{eq:v4v51-Z-factorization}
\end{equation}
On \(\mathfrak D_{a_0}\), every factor is zero-free and obeys
\begin{align}
 |Z_a-1|&\le e^{a_0}-1<1,\label{eq:v4v51-zero-disk}\\
 |Z_a|&\ge\delta_0:=2-e^{a_0}>0.\label{eq:v4v51-delta}
\end{align}
Consequently the branch continued from \(t=0\) satisfies
\begin{equation}
 K(t,\eta):=\log Z(t,\eta)
 =\sum_{a\in\mathcal A}K_a(t,\eta),
 \qquad K_a:=\log Z_a.
\label{eq:v4v51-log-factorization}
\end{equation}
\end{theorem}

\begin{proof}
Equation \eqref{eq:v4v51-conditional-product} and the fact that each
insertion belongs to one component give
\[
 e^{\sum_jt_jU_j}=\prod_ae^{W_a}
\]
and hence \eqref{eq:v4v51-Z-factorization}.  On
\(\mathfrak D_{a_0}\), \(\|W_a\|_\infty\le a_0\).  Therefore
\[
 |Z_a-1|
 \le\mathbb E|e^{W_a}-1|
 \le e^{\|W_a\|_\infty}-1
 \le e^{a_0}-1.
\]
The disk in \eqref{eq:v4v51-zero-disk} omits zero, and the reverse triangle
inequality gives \eqref{eq:v4v51-delta}.  Each factor has the logarithmic
branch continued from one.  Taking the logarithm of the finite product on
that branch proves \eqref{eq:v4v51-log-factorization}.
\end{proof}

\begin{corollary}[Mixed components have zero conditional cumulant]
\label{cor:v4v51-mixed-zero}
If a label set \(J\) meets two distinct conditional components, then
\begin{equation}
 \left.
 \prod_{j\in J}\partial_{t_j}K(t,\eta)
 \right|_{t=0}
 =0.
\label{eq:v4v51-mixed-zero}
\end{equation}
\end{corollary}

\begin{proof}
Every summand \(K_a\) in \eqref{eq:v4v51-log-factorization} depends only on
parameters whose insertions belong to \(D_a\).  A mixed derivative meeting
two components annihilates every summand.
\end{proof}

\begin{firewall}[The relevant load is local, not global]
\label{fw:v4v51-local-load}
The zero-free estimate does not require
\(\sum_{j=1}^N|t_j|u_j\le a_0\).  It requires the supremum of the component
loads in \eqref{eq:v4v51-domain}.  Replacing that supremum by the global sum
would force the analytic radius to shrink with the number of remote,
conditionally independent insertions and would discard the exact product
structure.
\end{firewall}

\subsection{All conditional cumulants from one Cauchy estimate}
\label{sec:v4v51-cauchy}

Put
\begin{equation}
 L_0:=-\log(2-e^{a_0}).
\label{eq:v4v51-L0}
\end{equation}
Since \(e^{a_0}-1<1\), the power series for
\(\log(1+(Z_a-1))\) and \eqref{eq:v4v51-zero-disk} give
\begin{equation}
 |K_a(t,\eta)|\le L_0
\label{eq:v4v51-log-bound}
\end{equation}
throughout \(\mathfrak D_{a_0}\), uniformly in \(a,\eta\), and volume.

\begin{theorem}[Uniform all-order conditional cumulants]
\label{thm:v4v51-all-cumulants}
Let \(J=\{j_1,\ldots,j_n\}\) consist of labels assigned to one conditional
component.  Their conditional joint cumulant satisfies
\begin{align}
 \left|
 \kappa_{\mu_a^\eta}(U_{j_1},\ldots,U_{j_n})
 \right|
 &\le
 L_0\left(\frac{n}{a_0}\right)^n
 \prod_{\ell=1}^nu_{j_\ell}\label{eq:v4v51-cumulant-nn}\\
 &\le
 L_0n!\left(\frac{e}{a_0}\right)^n
 \prod_{\ell=1}^nu_{j_\ell}.
\label{eq:v4v51-cumulant-factorial}
\end{align}
The constants do not depend on the total number \(N\) of insertions, the
number of conditional components, or the total volume.
\end{theorem}

\begin{proof}
If some \(u_{j_\ell}=0\), then that insertion vanishes almost surely and both bounds are immediate.  Otherwise, for every \(\ell\), choose the Cauchy radius
\[
 r_{j_\ell}:=\frac{a_0}{nu_{j_\ell}}.
\]
On the resulting distinguished boundary,
\(\sum_{\ell=1}^nr_{j_\ell}u_{j_\ell}=a_0\), so the polydisc lies in
\(\mathfrak D_{a_0}\).  Apply the multivariate Cauchy estimate to \(K_a\)
and use \eqref{eq:v4v51-log-bound}.  This proves
\eqref{eq:v4v51-cumulant-nn}.  The elementary lower bound
\(n!\ge(n/e)^n\) gives \eqref{eq:v4v51-cumulant-factorial}.
\end{proof}

\begin{remark}[Repeated labels]
\label{rem:v4v51-repeated-labels}
The same proof applies to a multi-index with repeated differentiation in
some \(t_j\).  Replace \(n\) by the total derivative order and repeat the
corresponding radius in the Cauchy product.  Thus the estimate controls
self-collisions as well as cumulants of distinct insertions.
\end{remark}

\subsection{One boundary derivative before coefficient extraction}
\label{sec:v4v51-boundary-derivative}

Let the conditional component action be \(S_a^\eta\).  A retained-data
direction \(h\) produces the score
\begin{equation}
 G_{a,h}:=D_\eta S_a^\eta[h].
\label{eq:v4v51-score}
\end{equation}
Assume
\begin{equation}
 \sup_\eta
 \|G_{a,h}-\mathbb E_{\mu_a^\eta}G_{a,h}\|_{L^2(\mu_a^\eta)}
 \le\sigma_a\|h\|
\label{eq:v4v51-score-L2}
\end{equation}
and, for explicit retained-data dependence of the insertions,
\begin{equation}
 \sup_\eta\|D_\eta U_j[h]\|_\infty
 \le d_0u_j\|h\|.
\label{eq:v4v51-explicit-derivative}
\end{equation}

\begin{theorem}[Uniform analytic boundary derivative]
\label{thm:v4v51-analytic-derivative}
On \(\mathfrak D_{a_0}\),
\begin{align}
 |D_\eta K_a(t,\eta)[h]|
 &\le
 M_{1,a}\|h\|,\label{eq:v4v51-analytic-derivative}\\
 M_{1,a}
 &:=
 \frac{
  e^{a_0}d_0a_0+(e^{a_0}-1)\sigma_a
 }{2-e^{a_0}}.
\label{eq:v4v51-M1}
\end{align}
Consequently, for labels \(j_1,\ldots,j_n\) in the component,
\begin{equation}
 \left|
 D_\eta
 \kappa_{\mu_a^\eta}(U_{j_1},\ldots,U_{j_n})[h]
 \right|
 \le
 M_{1,a}n!\left(\frac e{a_0}\right)^n
 \left(\prod_{\ell=1}^nu_{j_\ell}\right)\|h\|.
\label{eq:v4v51-marked-cumulant}
\end{equation}
\end{theorem}

\begin{proof}
Differentiate the normalized conditional expectation before taking the
logarithm.  With \(W=W_a(t)\), \(G=G_{a,h}\), and
\(\widetilde G=G-\mathbb E^\eta G\), one obtains the exact identity
\begin{equation}
 D_\eta K_a[h]
 =
 \frac{\mathbb E^\eta(e^WD_\eta W[h])}
      {\mathbb E^\eta e^W}
 -
 \frac{\mathbb E^\eta(e^W\widetilde G)}
      {\mathbb E^\eta e^W}.
\label{eq:v4v51-log-derivative-identity}
\end{equation}
The explicit derivative assumption and the component load give
\[
 \|D_\eta W[h]\|_\infty
 \le d_0A_a(t)\|h\|
 \le d_0a_0\|h\|.
\]
Use \(|e^W|\le e^{a_0}\) and the denominator lower bound
\eqref{eq:v4v51-delta} on the first fraction.  Since
\(\mathbb E^\eta\widetilde G=0\), the numerator of the second fraction is
\(\mathbb E^\eta[(e^W-1)\widetilde G]\).  Cauchy--Schwarz and
\eqref{eq:v4v51-score-L2} give
\[
 \left|
 \mathbb E^\eta[(e^W-1)\widetilde G]
 \right|
 \le(e^{a_0}-1)\sigma_a\|h\|.
\]
This proves \eqref{eq:v4v51-analytic-derivative}.  The function
\(D_\eta K_a[h]\) is analytic in the same polydisc.  Apply the Cauchy radii
used in Theorem~\ref{thm:v4v51-all-cumulants} and the bound
\eqref{eq:v4v51-analytic-derivative}; then use
\(n^n\le e^nn!\).
\end{proof}

\begin{corollary}[All score-marked conditional cumulants]
\label{cor:v4v51-score-marked}
If the \(U_j\) have no explicit \(\eta\)-dependence, then
\begin{equation}
 D_\eta
 \kappa_{\mu_a^\eta}(U_{j_1},\ldots,U_{j_n})[h]
 =
 -\kappa_{\mu_a^\eta}
 (U_{j_1},\ldots,U_{j_n},G_{a,h}),
\label{eq:v4v51-score-cumulant-identity}
\end{equation}
and the right side is bounded by
\eqref{eq:v4v51-marked-cumulant}.
\end{corollary}

\begin{proof}
Differentiate the conditional moment--cumulant formula, or take the
coefficient of \(t_{j_1}\cdots t_{j_n}\) in
\eqref{eq:v4v51-log-derivative-identity}.  With no explicit derivatives,
only the centered score term remains.  Its sign is negative because the
conditional density is \(e^{-S_a^\eta}\).
\end{proof}

\begin{firewall}[The scalar body is fixed before taking the logarithm]
\label{fw:v4v51-scalar-body}
Theorem~\ref{thm:v4v51-generator-factorization} concerns the scalar
normalization after the V43 functional has been applied.  A finite
Grassmann-valued generator can be treated by expanding around a zero-free
scalar body, but its nilpotent coefficients do not themselves define a
complex zero-free disk.  The scalar branch and the coefficient norm must be
fixed in that order.
\end{firewall}

\subsection{What the analytic estimate acquires for collision histories}
\label{sec:v4v51-collision-use}

\begin{theorem}[Component-local analytic carrier stability]
\label{thm:v4v51-carrier-stability}
At one conditional stage satisfying
\eqref{eq:v4v51-conditional-product}, all collision carriers created inside
one component obey the factorial value bound
\eqref{eq:v4v51-cumulant-factorial}.  Their next retained-data derivative,
including the complete score carrier required by
\eqref{eq:v4v50-carrier-derivative}, obeys
\eqref{eq:v4v51-marked-cumulant}.  Carriers mixing distinct components
vanish exactly.
\end{theorem}

\begin{proof}
Conditional collision carriers are precisely coefficients of \(K_a\), by
the definition of joint cumulants.  Apply
Theorem~\ref{thm:v4v51-all-cumulants} to their values and
Theorem~\ref{thm:v4v51-analytic-derivative} to their retained-data
derivatives.  Corollary~\ref{cor:v4v51-mixed-zero} handles distinct
components.
\end{proof}

\begin{assessment}[The analytic radius is uniform but the spatial tree is
not yet assembled]
\label{ass:v4v51-status}
V51 closes the order-growth problem for all bounded insertions assigned to
one exact conditional-product stage.  The constants grow as \(n!C^n\), not
with the total number of remote components.  What remains is to combine
these local collision constants with the pre-collision spatial factors of
V47 and the collision-history algebra of V49.  Simultaneous multi-label
collisions must be charged by their complete conditional cumulant once; they
cannot be replaced by free pair edges.
\end{assessment}

\subsection{V51 result ledger and V52 target}
\label{sec:v4v51-conclusion}

V51 proves exact factorization and zero-freeness of the conditional
cumulant generator on a component-load domain independent of volume.  One
multivariate Cauchy estimate gives every component cumulant with
\(n!C^n\) growth.  Differentiating the logarithm before coefficient
extraction gives a uniform analytic bound for every next score mark, thereby
closing the derivative defect found in V50 for bounded good-field
insertions.

V52 should combine the V51 local analytic carrier with the V49
collision-history expansion.  The required ledger must assign V47
pre-collision sweep decay to the distances between child clusters, charge
one V51 cumulant at each actual merger, and sum binary and simultaneous
merger histories without exceeding a factorial-times-exponential constant.
The V34 large-field tail remains a separate ancestry and must not be folded
into the bounded component generator.


\clearpage
\section{V52: merger-history summation and the multiway spatial obstruction}
\label{sec:v4v52}

V51 controls every conditional merger cumulant and its next score derivative
by \(k!C^k\), uniformly in volume and in the number of remote conditional
components.  Two questions remain before the V49 collision histories become
a pressure expansion.  First, do the nested set partitions introduce growth
worse than factorial times exponential?  Second, does every merger acquire
enough distinct spatial decay factors to form a tree?

This note answers the first question affirmatively by an analytic generating
function for merger hierarchies.  It also proves a sufficient weighted
merger theorem: local spanning-tree weights at each merger flatten to one
global spanning tree.  The second question fails if one uses only the V41
\(L^2\) contraction at a simultaneous \(k\)-way collision.  A centered rare
event has arbitrarily small \(L^2\) norm but all of its fixed-order cumulants
are of the same small order, rather than the \((k-1)\)-st power required by
a tree.  The missing input is therefore a multiway spatial carrier estimate,
not merger-history combinatorics.

\subsection{Abstract merger hierarchies}
\label{sec:v4v52-hierarchies}

A merger hierarchy on a finite label set \(I\) is a rooted nonplane tree
whose leaves are the elements of \(I\) and whose internal vertices have
arity at least two.  An internal vertex represents one conditional
cumulant of its child carriers.  Let \(A_k\ge0\) be a majorant for a local
\(k\)-child merger, including its analytic coefficient but not its child
weights.

Let \(H_n\) be the total weight of all merger hierarchies on \(n\) labelled
leaves, with weight \(\prod_vA_{\deg(v)}\).  Introduce the exponential
generating function
\begin{equation}
 H(z):=\sum_{n\ge1}H_n\frac{z^n}{n!}.
\label{eq:v4v52-H}
\end{equation}

\begin{lemma}[Exact hierarchy generating equation]
\label{lem:v4v52-hierarchy-equation}
As a formal power series,
\begin{equation}
 H(z)=z+\sum_{k\ge2}\frac{A_k}{k!}H(z)^k.
\label{eq:v4v52-hierarchy-equation}
\end{equation}
\end{lemma}

\begin{proof}
A hierarchy is either one labelled leaf, contributing \(z\), or an
unordered set of \(k\ge2\) child hierarchies attached to a root of weight
\(A_k\).  The labelled product rule gives \(H^k/k!\) for the unordered
children.  Summing the possible arities proves
\eqref{eq:v4v52-hierarchy-equation}.
\end{proof}

\begin{theorem}[Factorial-exponential hierarchy count]
\label{thm:v4v52-hierarchy-bound}
Suppose
\begin{equation}
 A_k\le A_*k!C_*^k,\qquad k\ge2.
\label{eq:v4v52-local-merger-growth}
\end{equation}
Then there is a finite constant \(C_H=C_H(A_*,C_*)\) such that
\begin{equation}
 H_n\le n!C_H^n,\qquad n\ge1.
\label{eq:v4v52-hierarchy-bound}
\end{equation}
\end{theorem}

\begin{proof}
The nonlinear part of \eqref{eq:v4v52-hierarchy-equation} is coefficientwise
majorized by
\begin{equation}
 \Phi(h):=
 A_*\sum_{k\ge2}(C_*h)^k
 =
 A_*\frac{(C_*h)^2}{1-C_*h}.
\label{eq:v4v52-Phi}
\end{equation}
The implicit equation
\[
 \widehat H=z+\Phi(\widehat H)
\]
has a unique analytic solution near \(z=0\): at the origin the derivative
of \(\widehat H-z-\Phi(\widehat H)\) with respect to \(\widehat H\) is one.
Choose a positive circle \(|z|=r_H\) contained in this analytic
neighbourhood and let \(M_H=\sup_{|z|=r_H}|\widehat H(z)|\).
Coefficientwise induction in the two formal equations gives
\([z^n]H\le[z^n]\widehat H\).  Cauchy's estimate therefore yields
\[
 \frac{H_n}{n!}\le M_Hr_H^{-n}.
\]
Enlarge \(C_H\) to dominate both \(r_H^{-1}\) and the finite factor \(M_H\).
\end{proof}

\begin{corollary}[V51 pays the abstract history combinatorics]
\label{cor:v4v52-V51-history}
The V51 bound
\[
 A_k=L_0k!\left(\frac e{a_0}\right)^k
\]
satisfies \eqref{eq:v4v52-local-merger-growth}.  Hence nested simultaneous
mergers and their set-partition histories have total coefficient growth at
most \(n!C_H^n\), before spatial weights are inserted.
\end{corollary}

\begin{firewall}[A factorial at each vertex does not multiply into
\((n!)^{n}\)]
\label{fw:v4v52-local-factorials}
The factors \(k!\) in the V51 local cumulants are paired with the
\(1/k!\) for the unordered child set in
\eqref{eq:v4v52-hierarchy-equation}.  Bounding every internal vertex by the
global order \(n!\) before summing histories would discard this cancellation
and create a false superfactorial obstruction.
\end{firewall}

\subsection{When local merger trees flatten to a global tree}
\label{sec:v4v52-flatten}

For disjoint child label sets \(I_1,\ldots,I_k\), let
\(\mathscr T(I_1,\ldots,I_k)\) be the set of trees on the \(k\) child
clusters.  An edge between \(I_r\) and \(I_s\) is assigned the kernel
\begin{equation}
 W(I_r,I_s):=
 \sum_{i\in I_r}\sum_{j\in I_s}w(i,j).
\label{eq:v4v52-cluster-kernel}
\end{equation}
This is the cluster subadditivity required in
\eqref{eq:v4v49-kernel-subadditivity}.

\begin{definition}[Spatially paid local merger]
\label{def:v4v52-paid-merger}
A \(k\)-child merger is spatially paid if its norm is bounded by
\begin{equation}
 A_*k!C_*^k
 \sum_{\tau\in\mathscr T(I_1,\ldots,I_k)}
 \prod_{\{r,s\}\in\tau}W(I_r,I_s)
 \prod_{r=1}^k\mathcal N(I_r),
\label{eq:v4v52-paid-merger}
\end{equation}
where \(\mathcal N(I_r)\) is the inherited norm of the child carrier.
\end{definition}

\begin{theorem}[Flattening of paid merger histories]
\label{thm:v4v52-flattening}
Suppose every internal vertex of a collision history satisfies
\eqref{eq:v4v52-paid-merger}, with the scalar merger weights obeying
\eqref{eq:v4v52-local-merger-growth}.  Then
\begin{equation}
 |\kappa(X_1,\ldots,X_n)|
 \le
 n!C_{\rm hist}^n
 \left(\prod_{i=1}^n\mathcal N(\{i\})\right)
 \sum_{T\text{ on }[n]}
 \prod_{\{i,j\}\in T}w(i,j).
\label{eq:v4v52-global-tree}
\end{equation}
\end{theorem}

\begin{proof}
At an internal vertex choose one child-cluster tree \(\tau\) from
\eqref{eq:v4v52-paid-merger}.  Expand each cluster edge with
\eqref{eq:v4v52-cluster-kernel} and choose one original-label edge between
the two child clusters.  By induction, each child hierarchy already carries
a tree on its own leaves.  The union of the \(k\) child trees and the
\(k-1\) selected edges of \(\tau\) is connected and has
\[
 \sum_{r=1}^k(|I_r|-1)+(k-1)
 =\left|\bigcup_rI_r\right|-1
\]
edges.  It is therefore a tree on the leaves below that vertex.

At the root this construction gives a global labelled tree.  Now fix both a
merger hierarchy and a global labelled tree \(T\).  If they are compatible,
the edges of \(T\) between the child clusters at an internal vertex determine
the quotient tree at that vertex uniquely; the original-label endpoint of
every quotient edge is also fixed by \(T\).  Thus a fixed hierarchy
contributes at most once to a fixed \(T\).  The total coefficient collected
at \(T\) is consequently bounded by the scalar weighted hierarchy sum in
Theorem~\ref{thm:v4v52-hierarchy-bound}, namely
\(n!C_{\rm hist}^n\).  Summing the resulting global trees gives
\eqref{eq:v4v52-global-tree}.
\end{proof}

\begin{remark}[Why no local Cayley factor enters the scalar hierarchy]
\label{rem:v4v52-local-tree-count}
The sum over the \(k^{k-2}\) trees at a \(k\)-child merger remains an edge
sum in \eqref{eq:v4v52-paid-merger}.  It is not inserted into the scalar
coefficient \(A_k\).  After flattening, those choices are collected in the
single explicit sum over global labelled trees in
\eqref{eq:v4v52-global-tree}.  Charging a Cayley factor both locally and in
the final tree sum would count the same edge choice twice and could create a
false zero-radius majorant.
\end{remark}

\subsection{What \(L^2\) pre-collision contraction can pay}
\label{sec:v4v52-L2}

For a binary collision, the two centered child carriers can both be placed
in \(L^2\).  After \(m\) pre-collision conditional channels, V41 gives one
factor \(q_{\rm H}^{m/2}\) from each child and hence
\(q_{\rm H}^{m}\).  Finite propagation speed converts this into one
exponential edge weight.  This is exactly the binary case acquired in
V47--V50.

For \(k>2\), an \(L^2\)-\(L^\infty\) estimate can place only two children in
\(L^2\).  Interpolating every child into \(L^k\) does not improve the total
power.  Indeed, if \(P^m\) is an \(L^\infty\) contraction and contracts
centered \(L^2\) norm by \(q_{\rm H}^{m/2}\), then
\begin{equation}
 \|P^mf\|_k
 \le
 q_{\rm H}^{m/k}
 \|f\|_2^{2/k}\|f\|_\infty^{1-2/k}.
\label{eq:v4v52-interpolation}
\end{equation}
Multiplying \(k\) such estimates yields only \(q_{\rm H}^{m}\), not
\(q_{\rm H}^{m(k-1)}\).

\begin{proposition}[Small \(L^2\) norm does not imply a multiway tree power]
\label{prop:v4v52-rare-event}
Let \(E\) be an event of probability \(p\in(0,1/2)\) and put
\[
 Y_p:=\mathbf1_E-p.
\]
Then \(\|Y_p\|_\infty\le1\), \(\|Y_p\|_2^2=p(1-p)\), and for every fixed
\(k\ge2\),
\begin{equation}
 \kappa_k(Y_p,\ldots,Y_p)=p+O_k(p^2)
 \qquad(p\downarrow0).
\label{eq:v4v52-Bernoulli-cumulant}
\end{equation}
Consequently there is no constant \(C_k\), independent of \(p\), such that
\begin{equation}
 |\kappa_k(Y_p,\ldots,Y_p)|
 \le C_k\|Y_p\|_2^{\,2(k-1)}
\label{eq:v4v52-false-tree-power}
\end{equation}
for \(k\ge3\).
\end{proposition}

\begin{proof}
The centered and uncentered Bernoulli variables have the same cumulants of
order at least two.  Their cumulant generating function is
\[
 \log\mathbb E e^{t\mathbf1_E}
 =\log\bigl(1+p(e^t-1)\bigr)
 =p(e^t-1)+O(p^2)
\]
uniformly for \(t\) in a fixed neighbourhood of zero.  The \(k\)-th
derivative at zero is \(p+O_k(p^2)\), proving
\eqref{eq:v4v52-Bernoulli-cumulant}.  The proposed right side in
\eqref{eq:v4v52-false-tree-power} is of order \(p^{k-1}\), which is
smaller than order \(p\) for \(k\ge3\).
\end{proof}

\begin{firewall}[The rare-event test is a norm no-go, not a Higgs
counterexample]
\label{fw:v4v52-rare-scope}
Proposition~\ref{prop:v4v52-rare-event} does not assert that the conditional
Higgs carrier contains a Bernoulli spike.  It proves that boundedness plus
\(L^2\) contraction alone cannot establish the \((k-1)\)-edge power required
at a simultaneous merger.  A physical proof must use additional
log-concave regularity, resolve the merger into actual binary incidences, or
derive the complete multiway spatial carrier directly.
\end{firewall}

\subsection{V52 result ledger and V53 target}
\label{sec:v4v52-conclusion}

V52 proves that V51's \(k!C^k\) local cumulant bounds make the complete
merger-hierarchy sum factorial-exponential rather than superfactorial.  It
also proves that locally paid merger trees flatten without cycles or
duplicate edges to one global labelled tree.  Thus history combinatorics is
not the current bottleneck.

The missing hypothesis is the spatial payment
\eqref{eq:v4v52-paid-merger} for a simultaneous multiway collision.  V41
\(L^2\) contraction pays the binary case but cannot, as a matter of norm
logic, manufacture \(k-1\) edge factors at order \(k\).

V53 should return to the exact finite-range bipartite incidence at the first
multiway collision.  It should determine whether the common conditional
component admits a canonical sequence of binary conditional covariances
before the component cumulant is estimated.  If the exact law of total
cumulance produces an irreducible \(k\)-carrier, the alternative target is
a log-Sobolev or exponential-Orlicz strong data-processing estimate that
contracts each child load separately.  The V52 rare-event test must be used
to reject any argument based only on \(L^2\) smallness.


\clearpage
\section{V53: log-Sobolev analyticity pays every child of a multiway collision}
\label{sec:v4v53}

V52 proves that boundedness and \(L^2\) contraction alone cannot supply the
\(k-1\) spatial factors required at a simultaneous \(k\)-way collision.
The positive good-field carrier has more structure.  Its strong convexity
gives a log-Sobolev inequality, and the V34 smooth cutoff gives globally
bounded derivatives for the good Yukawa insertion.  A centered Lipschitz
function under a log-Sobolev law has sub-Gaussian concentration.  Therefore
its complex exponential expectation remains close to one on a domain
controlled by the sum of Lipschitz loads.

This note proves the corresponding multivariate zero-free theorem.  Cauchy
extraction gives a joint cumulant bound proportional to the product of all
individual Lipschitz constants.  At a physical multiway collision, those
constants contain the separate pre-collision response distances of the
children.  Dropping the smallest one and using the triangle inequality
turns the remaining factors into a child spanning tree.  This acquires the
spatial value estimate missing in V52 for the smooth good-field branch.

\subsection{A concentration input from strong convexity}
\label{sec:v4v53-concentration}

Let \(\nu\) be a probability law on a finite-dimensional Euclidean field
space satisfying the log-Sobolev inequality
\begin{equation}
 \operatorname{Ent}_\nu(f^2)
 \le\frac{2}{\kappa_*}\mathbb E_\nu|\nabla f|^2.
\label{eq:v4v53-LSI}
\end{equation}
Every law with \(D^2S\ge\kappa_*1\) satisfies
\eqref{eq:v4v53-LSI} by the Bakry--Emery criterion.  The V38 conditional and
marginal positive Higgs laws have this lower Hessian bound; after V34
promotion of the good determinant modulus one uses the reduced constant
\(\kappa_*=\kappa/2\).

\begin{lemma}[Real and complex Lipschitz concentration]
\label{lem:v4v53-complex-concentration}
Let \(F\) be complex valued, \(\mathbb E_\nu F=0\), and
\begin{equation}
 \|\nabla\operatorname{Re}F\|_\infty+
 \|\nabla\operatorname{Im}F\|_\infty\le L.
\label{eq:v4v53-complex-Lipschitz}
\end{equation}
There are constants \(c_*,C_*>0\), depending only on \(\kappa_*\), such that
\begin{align}
 \nu(|F|\ge r)&\le4e^{-c_*r^2/L^2},\label{eq:v4v53-tail}\\
 \mathbb E_\nu\bigl(|F|e^{|F|}\bigr)
 &\le C_*L
\label{eq:v4v53-exponential-first-moment}
\end{align}
whenever \(0<L\le1\).
\end{lemma}

\begin{proof}
The Herbst argument applied to \eqref{eq:v4v53-LSI} gives, for every
centered real \(L\)-Lipschitz function \(G\),
\[
 \nu(|G|\ge r)\le2e^{-\kappa_*r^2/(2L^2)}.
\]
Apply this to the real and imaginary parts.  Since
\(|F|\ge r\) implies that one part has modulus at least \(r/\sqrt2\), a
union bound gives \eqref{eq:v4v53-tail} with
\(c_*=\kappa_*/4\).

For the second estimate, use tail integration with
\(\psi(r)=re^r\):
\[
 \mathbb E\psi(|F|)
 =\int_0^\infty\psi'(r)\nu(|F|\ge r)\,dr.
\]
Insert \eqref{eq:v4v53-tail}, put \(r=Ls\), and use \(L\le1\).
The integral of
\((1+s)e^s e^{-c_*s^2}\) is finite and depends only on
\(\kappa_*\), proving \eqref{eq:v4v53-exponential-first-moment}.
\end{proof}

\subsection{A Lipschitz-load zero-free generator}
\label{sec:v4v53-generator}

Let \(F_1,\ldots,F_k\) be complex functions with
\(\mathbb E_\nu F_i=0\).  Put
\begin{equation}
 L_i:=
 \|\nabla\operatorname{Re}F_i\|_\infty+
 \|\nabla\operatorname{Im}F_i\|_\infty.
\label{eq:v4v53-Li}
\end{equation}
For \(t\in\mathbb C^k\), define
\[
 W(t):=\sum_{i=1}^kt_iF_i,
 \qquad
 \Lambda(t):=\sum_{i=1}^k|t_i|L_i.
\]

\begin{theorem}[Zero-free centered generator from Lipschitz load]
\label{thm:v4v53-zero-free-generator}
There is \(s_0=s_0(\kappa_*)>0\) such that, whenever
\(\Lambda(t)\le s_0\),
\begin{align}
 \left|\mathbb E_\nu e^{W(t)}-1\right|&\le\frac12,
\label{eq:v4v53-generator-disk}\\
 \left|\log\mathbb E_\nu e^{W(t)}\right|&\le\log2
\label{eq:v4v53-generator-log}
\end{align}
on the logarithmic branch continued from \(t=0\).  The domain depends on
the sum of individual Lipschitz loads and not on the number of variables.
\end{theorem}

\begin{proof}
The function \(W(t)\) is centered.  For complex \(t_i\), the sum of its
real and imaginary Lipschitz constants is at most \(2\Lambda(t)\).  Choose
\(s_0\) so that \(2s_0\le1\) and \(2C_*s_0\le1/2\), with
\(C_*\) from
Lemma~\ref{lem:v4v53-complex-concentration}.  Since
\[
 |e^z-1|\le |z|e^{|z|},
\]
\eqref{eq:v4v53-exponential-first-moment} gives
\[
 \left|\mathbb E e^{W(t)}-1\right|
 \le\mathbb E|W(t)|e^{|W(t)|}
 \le2C_*\Lambda(t)\le\frac12.
\]
This proves \eqref{eq:v4v53-generator-disk}.  The disk
\(|z-1|\le1/2\) is zero-free, and the logarithmic power series gives
\[
 |\log z|
 \le\sum_{n\ge1}\frac{|z-1|^n}{n}
 \le\log2,
\]
which proves \eqref{eq:v4v53-generator-log}.
\end{proof}

\begin{firewall}[Centering is essential for the Lipschitz domain]
\label{fw:v4v53-centering}
The Lipschitz seminorm does not control an additive constant.  Theorem
\ref{thm:v4v53-zero-free-generator} applies to
\(F_i-\mathbb EF_i\).  Joint cumulants of order at least two are invariant
under this centering, but the uncentered exponential generator must first
have its scalar means factored out.
\end{firewall}

\subsection{All-order product of individual response scales}
\label{sec:v4v53-cumulants}

\begin{theorem}[Lipschitz product bound for a joint cumulant]
\label{thm:v4v53-Lipschitz-cumulant}
For \(k\ge2\),
\begin{align}
 |\kappa_\nu(F_1,\ldots,F_k)|
 &\le
 (\log2)\left(\frac{k}{s_0}\right)^k
 \prod_{i=1}^kL_i\label{eq:v4v53-Lipschitz-kk}\\
 &\le
 (\log2)k!\left(\frac e{s_0}\right)^k
 \prod_{i=1}^kL_i.
\label{eq:v4v53-Lipschitz-factorial}
\end{align}
If one \(L_i\) vanishes, the cumulant vanishes.
\end{theorem}

\begin{proof}
When every \(L_i>0\), use the Cauchy radii
\[
 r_i=\frac{s_0}{kL_i}.
\]
Their distinguished boundary satisfies
\(\sum_i r_iL_i=s_0\).  Apply the multivariate Cauchy estimate to the
analytic logarithm in
Theorem~\ref{thm:v4v53-zero-free-generator}.  This gives
\eqref{eq:v4v53-Lipschitz-kk}.  Use \(k^k\le e^kk!\) for
\eqref{eq:v4v53-Lipschitz-factorial}.  If \(L_i=0\), then \(F_i\) is
constant on the connected field space and its centered version is zero.
\end{proof}

\begin{assessment}[Why the V52 rare event is excluded]
\label{ass:v4v53-rare-event-excluded}
The centered Bernoulli spike in
Proposition~\ref{prop:v4v52-rare-event} has small \(L^2\) norm but no
correspondingly small Lipschitz seminorm on a strongly log-concave
Euclidean field space.  V53 uses the latter seminorm and the log-Sobolev
concentration it generates.  Thus
\eqref{eq:v4v53-Lipschitz-factorial} adds a genuine hypothesis absent from
the no-go example.
\end{assessment}

\subsection{Application at a physical multiway collision}
\label{sec:v4v53-collision}

Let child carriers \(F_i\) approach one actual conditional component
centered at a lattice site or fixed-size block \(x\).  Let \(S_i\) be the
reference support of child \(i\).  Assume the buffered response and the
finite-speed pre-collision sweep give
\begin{equation}
 L_i\le\lambda_i e^{-m_{\rm c}d(S_i,x)}
\label{eq:v4v53-child-response}
\end{equation}
with volume-uniform \(\lambda_i\).  For the smooth V34 good insertion,
\eqref{eq:v4v34-good-first} supplies the required uniform local derivative;
the V38 conditional Witten estimate and V41 pre-collision propagation
supply the spatial factor.  Equation
\eqref{eq:v4v53-child-response} is retained as the explicit acquisition
interface for carriers created at later stages.

\begin{theorem}[A multiway collision carries a child spanning tree]
\label{thm:v4v53-multiway-tree}
Under \eqref{eq:v4v53-child-response}, the collision cumulant satisfies
\begin{equation}
 |\kappa_\nu(F_1,\ldots,F_k)|
 \le
 (\log2)k!\left(\frac e{s_0}\right)^k
 \left(\prod_{i=1}^k\lambda_i\right)
 \sum_{r=1}^k
 \prod_{i\ne r}
 e^{-\frac12m_{\rm c}d(S_i,S_r)}.
\label{eq:v4v53-multiway-tree}
\end{equation}
Each product on the right is the star tree rooted at child \(r\).
\end{theorem}

\begin{proof}
Choose \(r\) for which \(d(S_r,x)\) is minimal.  From
\eqref{eq:v4v53-Lipschitz-factorial} and
\eqref{eq:v4v53-child-response},
\[
 |\kappa_\nu(F_1,\ldots,F_k)|
 \le
 (\log2)k!\left(\frac e{s_0}\right)^k
 \left(\prod_i\lambda_i\right)
 \prod_i e^{-m_{\rm c}d(S_i,x)}.
\]
Drop the factor indexed by \(r\), which is at most one.  For \(i\ne r\),
the triangle inequality and the choice of \(r\) give
\[
 d(S_i,S_r)
 \le d(S_i,x)+d(S_r,x)
 \le2d(S_i,x),
\]
with fixed block radii absorbed in the constant.  Hence
\[
 e^{-m_{\rm c}d(S_i,x)}
 \le e^{-\frac12m_{\rm c}d(S_i,S_r)}.
\]
This proves the term rooted at the minimizing child.  Enlarging the right
side to the sum over all possible roots gives
\eqref{eq:v4v53-multiway-tree}.
\end{proof}

\begin{corollary}[The V52 spatial merger interface on the good branch]
\label{cor:v4v53-paid-merger}
For bounded smooth good-field child carriers satisfying
\eqref{eq:v4v53-child-response}, the local collision obeys
\eqref{eq:v4v52-paid-merger} with
\[
 w(S_i,S_j)=e^{-\frac12m_{\rm c}d(S_i,S_j)}
\]
and with the star subfamily of child trees.  Therefore the V52 flattening
theorem applies to their value norms.
\end{corollary}

\begin{proof}
The star trees in \eqref{eq:v4v53-multiway-tree} belong to
\(\mathscr T(I_1,\ldots,I_k)\).  Their coefficient has the required
\(k!C^k\) form.  Insert the inherited child amplitudes into the
\(\lambda_i\).
\end{proof}

\begin{firewall}[The V34 tail remains a separate species]
\label{fw:v4v53-tail-separate}
The uniform derivative input uses the smooth good operator
\(g_xV_x\) and \eqref{eq:v4v34-good-first}.  The tail operator
\(t_xV_x\) is unbounded but already receives one independent factor
\[
 e^{-c_{\rm T}\epsilon_Y^{-2\alpha}}
\]
per actual tail cell from
Theorem~\ref{thm:v4v34-tail-small}.  V53 neither inserts the tail into the
Lipschitz generator nor charges its rare-cell payment again.
\end{firewall}

\begin{firewall}[Later-stage spatial derivatives remain an interface]
\label{fw:v4v53-later-gradient}
V51 controls the total retained-data derivative of every local collision
carrier with factorial-exponential order growth.  To apply V53 recursively,
that derivative must also satisfy the spatially resolved child estimate
\eqref{eq:v4v53-child-response}.  The first good-field generation has this
estimate from V34/V38.  V53 does not yet prove its preservation after an
arbitrary merger history.
\end{firewall}

\subsection{V53 result ledger and V54 target}
\label{sec:v4v53-conclusion}

V53 proves a zero-free centered exponential generator on a domain controlled
only by total Lipschitz load.  Multivariate Cauchy extraction gives every
joint cumulant with \(k!C^k\) growth times the product of the individual
Lipschitz response scales.  At a physical multiway collision, those scales
produce a star tree between the child clusters.  This verifies the V52
spatial merger interface for the value norm of the first smooth good-field
collision and bypasses the \(L^2\) rare-event obstruction.

V54 should differentiate the Lipschitz-load generator with respect to the
next retained color before coefficient extraction.  The target is a
spatially resolved score estimate showing that the derivative of a merger
carrier retains one response scale for every child and one additional scale
for the new boundary mark.  Such a theorem would propagate
\eqref{eq:v4v53-child-response} through arbitrary collision histories and
close the analytic tree compiler for the bounded good-field sector.


\clearpage
\section{V54: the marked Lipschitz generator propagates the collision tree}
\label{sec:v4v54}

V53 proves that a multiway collision of smooth good-field carriers has a
value bound containing one Lipschitz response scale from every child.  To
reuse the resulting carrier at the next color step, its derivative with
respect to the newly retained field must have the same spatial structure.
V50 identifies that derivative as the complete conditional cumulant with
one score insertion.

The score should therefore be inserted into the V53 centered exponential
generator before coefficient extraction.  It becomes one additional
Lipschitz variable.  The resulting Cauchy estimate contains the product of
all child response scales and the mixed-derivative scale of the score.
When these scales are resolved spatially around the collision component,
they form a star tree on the child clusters together with the new mark.
This closes the first-derivative recursion for the bounded good-field
collision hierarchy.

\subsection{The spatial score seminorm}
\label{sec:v4v54-score}

Let \(\mu_D^\eta\) be one conditional positive component law and let a
retained-field direction \(h_Y\) be supported in a set \(Y\).  Its centered
score is
\begin{equation}
 \widetilde G_{Y,h}
 :=
 D_\eta S_D^\eta[h_Y]
 -\mathbb E_{\mu_D^\eta}D_\eta S_D^\eta[h_Y].
\label{eq:v4v54-centered-score}
\end{equation}
Centering does not change its gradient in the integrated field.  Define
\begin{equation}
 B_D(Y):=
 \sup_{\eta,\ \|h_Y\|\le1}
 \left(
  \|\nabla\operatorname{Re}\widetilde G_{Y,h}\|_\infty+
  \|\nabla\operatorname{Im}\widetilde G_{Y,h}\|_\infty
 \right).
\label{eq:v4v54-score-seminorm}
\end{equation}

For the finite-range positive Higgs action, \(B_D(Y)=0\) unless \(Y\) meets
the fixed interaction neighbourhood of \(D\).  After promotion of the V34
good determinant modulus, the mixed Hessian has the exponentially weighted
row bound \eqref{eq:v4v34-weighted-log-Hessian}.  Thus the combined
good-carrier score has the interface
\begin{equation}
 B_D(Y)
 \le J_*e^{-m_*d(D,Y)}
\label{eq:v4v54-score-decay}
\end{equation}
after reducing \(m_*>0\) below the finite-range/weighted-Hessian rate.

\begin{lemma}[Acquisition of the score kernel]
\label{lem:v4v54-score-kernel}
At the first V34/V38 good-field conditional stage,
\eqref{eq:v4v54-score-decay} holds with constants independent of volume and
the conditioned field.
\end{lemma}

\begin{proof}
For the original Higgs action,
\(\nabla_\phi D_\eta S_D^\eta[h_Y]\) is the mixed Hessian applied to
\(h_Y\).  Finite interaction range gives exact vanishing outside a fixed
neighbourhood and a uniform local stencil norm inside it.  The promoted good
modulus changes this mixed Hessian by a kernel whose exponentially weighted
row and column sums are bounded by
\eqref{eq:v4v34-weighted-log-Hessian}.  Add the two kernels and decrease the
exponential rate.  The V34 bounds are pointwise in the field and uniform in
volume.
\end{proof}

\begin{firewall}[The score mean has no Lipschitz cost]
\label{fw:v4v54-score-mean}
The conditional mean subtracted in
\eqref{eq:v4v54-centered-score} depends on the retained data but is constant
in the integrated variable.  It therefore has zero integrated-field
gradient and contributes nothing to
\eqref{eq:v4v54-score-seminorm}.  It must still be retained when
differentiating the normalized conditional law, as in
\eqref{eq:v4v51-log-derivative-identity}.
\end{firewall}

\subsection{The score as one generator variable}
\label{sec:v4v54-extended-generator}

Let \(F_1,\ldots,F_k\) be centered complex child carriers in one conditional
component, with Lipschitz seminorms \(L_i\) from
\eqref{eq:v4v53-Li}.  For a fixed direction \(h_Y\), introduce an additional
complex parameter \(s\) and define
\begin{equation}
 \mathcal K(t,s)
 :=
 \log\mathbb E_{\mu_D^\eta}
 \exp\left(
  \sum_{i=1}^kt_iF_i+s\widetilde G_{Y,h}
 \right).
\label{eq:v4v54-extended-generator}
\end{equation}

\begin{theorem}[All-order score-marked Lipschitz bound]
\label{thm:v4v54-score-marked}
For every \(k\ge1\),
\begin{align}
 &\left|
 \kappa_{\mu_D^\eta}
 (F_1,\ldots,F_k,\widetilde G_{Y,h})
 \right|\nonumber\\
 &\qquad\le
 (\log2)(k+1)!
 \left(\frac e{s_0}\right)^{k+1}
 B_D(Y)\|h_Y\|
 \prod_{i=1}^kL_i.
\label{eq:v4v54-score-marked}
\end{align}
The constants are uniform in volume, the conditioned field, and the number
of remote conditional components.
\end{theorem}

\begin{proof}
The generator \eqref{eq:v4v54-extended-generator} is the V53 generator for
\(k+1\) centered variables.  Its total Lipschitz load is
\[
 \sum_i|t_i|L_i+|s|B_D(Y)\|h_Y\|.
\]
Apply Theorem~\ref{thm:v4v53-Lipschitz-cumulant} at order \(k+1\).
\end{proof}

\begin{corollary}[Spatial derivative of a merger carrier]
\label{cor:v4v54-merger-derivative}
If the children have no explicit \(\eta\)-dependence at the current stage,
then
\begin{align}
 &\left|
 D_\eta\kappa_{\mu_D^\eta}(F_1,\ldots,F_k)[h_Y]
 \right|\nonumber\\
 &\qquad\le
 (\log2)(k+1)!
 \left(\frac e{s_0}\right)^{k+1}
 B_D(Y)\|h_Y\|
 \prod_{i=1}^kL_i.
\label{eq:v4v54-merger-derivative}
\end{align}
\end{corollary}

\begin{proof}
The normalized derivative identity is
\[
 D_\eta\kappa_{\mu_D^\eta}(F_1,\ldots,F_k)[h_Y]
 =
 -\kappa_{\mu_D^\eta}
 (F_1,\ldots,F_k,\widetilde G_{Y,h}),
\]
the all-order version of
\eqref{eq:v4v51-score-cumulant-identity}.  Apply
Theorem~\ref{thm:v4v54-score-marked}.
\end{proof}

\subsection{The children and mark form one tree}
\label{sec:v4v54-marked-tree}

Let \(x\) denote the collision component.  Suppose the child response scales
and score scale obey
\begin{align}
 L_i&\le\lambda_i e^{-m_{\rm c}d(S_i,x)},\label{eq:v4v54-child-scales}\\
 B_D(Y)&\le J_*e^{-m_{\rm c}d(Y,x)},\label{eq:v4v54-mark-scale}
\end{align}
after taking a common rate \(m_{\rm c}>0\).

\begin{theorem}[Marked multiway collision tree]
\label{thm:v4v54-marked-tree}
Under \eqref{eq:v4v54-child-scales}--\eqref{eq:v4v54-mark-scale},
\begin{align}
 &\left|
 D_\eta\kappa_{\mu_D^\eta}(F_1,\ldots,F_k)[h_Y]
 \right|\nonumber\\
 &\quad\le
 C_*^{k+1}(k+1)!
 J_*\left(\prod_{i=1}^k\lambda_i\right)\|h_Y\|
 \sum_{r\in\{Y,1,\ldots,k\}}
 \prod_{\substack{v\in\{Y,1,\ldots,k\}\\v\ne r}}
 e^{-\frac12m_{\rm c}d(S_v,S_r)},
\label{eq:v4v54-marked-tree}
\end{align}
where \(S_Y:=Y\), and \(C_*\) depends only on the V53 analytic radius.
Every product is a star tree on the \(k\) children and the mark.
\end{theorem}

\begin{proof}
Insert \eqref{eq:v4v54-child-scales} and
\eqref{eq:v4v54-mark-scale} into
\eqref{eq:v4v54-merger-derivative}.  This gives the product of
\(k+1\) exponentials measuring the distances of the nodes
\(Y,S_1,\ldots,S_k\) from \(x\).  Choose a node \(r\) of minimal distance
from \(x\), drop its exponential factor, and use
\[
 d(S_v,S_r)\le d(S_v,x)+d(S_r,x)\le2d(S_v,x)
\]
for every \(v\ne r\).  The remaining \(k\) factors give the star rooted at
\(r\).  Enlarge to the sum over all roots and absorb
\((\log2)(e/s_0)^{k+1}\) into \(C_*^{k+1}\).
\end{proof}

\begin{corollary}[Recursive spatial Lipschitz stability]
\label{cor:v4v54-recursive-stability}
For the first smooth good-field merger, the retained-field derivative of
the new carrier is bounded by a factorial-exponential coefficient times a
tree containing:
\begin{enumerate}
\item every tree already carried by the child histories;
\item \(k-1\) edges joining the \(k\) child clusters at their collision; and
\item one edge joining the new retained-field mark.
\end{enumerate}
No edge or score is used twice.
\end{corollary}

\begin{proof}
The V53 value theorem supplies a child collision tree.  The stronger marked
estimate \eqref{eq:v4v54-marked-tree} supplies a tree on the children and
mark directly; retain this latter tree rather than multiplying the two
estimates.  Flatten it with the inherited child trees by
Theorem~\ref{thm:v4v52-flattening}.  The edge count is
\[
 \sum_i(|I_i|-1)+k
 =\left|\bigcup_iI_i\right|,
\]
which is the tree edge count on the original leaves plus one marked vertex.
\end{proof}

\begin{firewall}[Use the marked tree instead of value tree times response]
\label{fw:v4v54-no-double-tree}
Corollary~\ref{cor:v4v54-recursive-stability} uses the complete marked
cumulant once.  Multiplying the V53 value tree by an independent score
response would retain \(2k-1\) newly created edges where the marked tree has
only \(k\), and would double count the child collision.  The marked
generator replaces that product.
\end{firewall}

\subsection{Explicit child derivatives}
\label{sec:v4v54-explicit}

If a child \(F_i(\eta)\) has explicit retained-field dependence, normalized
differentiation gives
\begin{align}
 D_\eta\kappa(F_1,\ldots,F_k)[h_Y]
 &=
 \sum_{i=1}^k
 \kappa(F_1,\ldots,D_\eta F_i[h_Y],\ldots,F_k)\nonumber\\
 &\quad-
 \kappa(F_1,\ldots,F_k,\widetilde G_{Y,h}).
\label{eq:v4v54-full-derivative}
\end{align}

\begin{proposition}[One-use ancestry of explicit marks]
\label{prop:v4v54-explicit-marks}
Suppose each explicit derivative \(D_\eta F_i[h_Y]\) already has a marked
tree bound obtained at its creation stage.  Then every term in
\eqref{eq:v4v54-full-derivative} has exactly one mark ancestry:
the first line inherits the mark from one child, while the second line
creates the current score mark through
\eqref{eq:v4v54-marked-tree}.
\end{proposition}

\begin{proof}
Multilinearity places an explicit derivative in exactly one child slot.
Apply the V53 value generator to the remaining undifferentiated children and
flatten with the marked child tree.  For the last term, use
Theorem~\ref{thm:v4v54-marked-tree}.  These alternatives are the exhaustive
product-rule placements in \eqref{eq:v4v54-full-derivative}.
\end{proof}

\begin{firewall}[The large-field tail and complex phase remain separate]
\label{fw:v4v54-sector-boundary}
The log-Sobolev law is the positive good carrier.  V34 tail cells retain
their super-small local parameter and existing marked fermion tree.  The
unit-modulus determinant phase remains in the V35--V36 phase ancestry.
Neither is inserted into the positive conditional score in V54.
\end{firewall}

\subsection{V54 result ledger and V55 audit target}
\label{sec:v4v54-conclusion}

V54 inserts the centered retained-field score into the V53 Lipschitz-load
generator before coefficient extraction.  It proves an all-order marked
cumulant bound containing every child response scale and the score's mixed
Hessian scale.  Spatial resolution produces a tree on all child clusters
plus the new mark.  Together with exact product-rule placement of explicit
child derivatives, this propagates the first-derivative tree norm through
arbitrary good-field merger histories without multiplying separate value
and response capacities.

V55 must begin with the scheduled YM/NS audit.  It should then derive the
twice-marked generator identity needed for the continuum Hessian and Ward
interfaces.  The two genuine alternatives are one second score/child mark
or two ordered first marks.  Their spatial ledger must contain a tree on the
children and both marks, with no square substituted for an actual second
derivative.


\clearpage
\section{V55: companion audit and the exact twice-marked collision identity}
\label{sec:v4v55}

V54 propagates one retained-field mark through the good-field collision
generator.  A continuum Hessian or twice-marked Ward interface requires the
second derivative.  There are two genuinely different sources: one second
derivative of the conditional action or insertion, and two first
derivatives at their actual slots.  This note derives the complete identity
at the cumulant level before estimating either source.

The two-first-score channel is controlled by the V53 Lipschitz generator.
The one-second-score channel requires the integrated-field Lipschitz
seminorm of \(D_\eta^2S\), hence a third field derivative of the conditional
action.  V34 proves a weighted Hessian bound for the promoted good modulus,
but it does not state this third-derivative bound.  V55 therefore isolates
that acquisition rather than replacing it by a square of the first-score
estimate.

\subsection{Scheduled read-only companion audit}
\label{sec:v4v55-audit}

The newest active companion files at the completed audit were
\begin{align*}
 &\texttt{Renewal\_Yang\_Mills\_Mass\_Gap\_Research\_Notes\_V71.tex},\\
 &\qquad 1\,250\,532\ \text{bytes},\qquad
 \texttt{SHA--256 }\\
 &\qquad
 \texttt{8F57CA5FEF72B9E470965CDC1E600BED15017AE51490D8678E02C20989F13CB7},
 \\
 &\texttt{Renewal\_NS\_Stability\_Research\_Notes\_V31.tex},\\
 &\qquad 887\,537\ \text{bytes},\qquad
 \texttt{SHA--256 }\\
 &\qquad
 \texttt{F1121B62238AD5491BB7CF03635EA9934942EDF573ED8FAAA9EE79E1D38A6696}.
\end{align*}

The Yang--Mills note has advanced from V66 to V71.  Its V70 audit and V71
ledger preserve exact local incidence: two bank labels selecting the same
protected output still give one idempotent projector, and distinct
coordinate labels do not imply orthogonality or a second response.  Its next
off-diagonal calculation keeps both complete joining words until their
actual cross term is computed.  The Navier--Stokes note is unchanged; its
one-use persistence and exact unmarking rules remain the relevant proof
discipline.

\begin{assessment}[Type-correct V55 transfer]
\label{ass:v4v55-audit-transfer}
The present two marks are counted by differentiating one exact conditional
cumulant.  If both derivatives hit the same score or insertion, there is one
second-mark occurrence.  If they hit distinct first-order positions, there
are two actual first-mark occurrences.  Labels alone do not create a second
factor, and the first-score bound may not be squared to stand in for the
second score.  No numerical YM or NS estimate is imported.
\end{assessment}

\subsection{A general derivative rule for conditional cumulants}
\label{sec:v4v55-derivative-rule}

Let \(\mu^\eta\) have conditional density proportional to
\(e^{-S(\phi,\eta)}\).  For a retained direction \(v\), set
\begin{equation}
 G_v:=D_\eta S[v],
 \qquad
 \widetilde G_v:=G_v-\mathbb E^\eta G_v.
\label{eq:v4v55-first-score}
\end{equation}

\begin{lemma}[One derivative of a conditional cumulant]
\label{lem:v4v55-one-derivative}
For \(\eta\)-dependent conditional observables \(X_1,\ldots,X_m\),
\begin{align}
 D_\eta\kappa^\eta(X_1,\ldots,X_m)[v]
 &=
 \sum_{j=1}^m
 \kappa^\eta(X_1,\ldots,D_\eta X_j[v],\ldots,X_m)\nonumber\\
 &\quad-
 \kappa^\eta(X_1,\ldots,X_m,G_v).
\label{eq:v4v55-one-derivative}
\end{align}
The final argument may equivalently be \(\widetilde G_v\).
\end{lemma}

\begin{proof}
Differentiate the conditional cumulant generating function
\[
 \log\mathbb E^\eta
 \exp\left(\sum_jt_jX_j(\eta)\right)
\]
before extracting the coefficient of \(t_1\cdots t_m\).  Differentiating
the insertions gives the first line.  Differentiating the normalized density
gives minus the joint cumulant with \(G_v\).  Since a cumulant of order at
least two is unchanged when one argument is shifted by a constant,
\(G_v\) may be centered.
\end{proof}

\subsection{The exact two-mark formula}
\label{sec:v4v55-two-mark}

Let \(U_1,\ldots,U_n\) be conditional insertions.  Write
\[
 U_{i;v}:=D_\eta U_i[v],
 \qquad
 U_{i;vw}:=D_\eta^2U_i[v,w],
 \qquad
 H_{vw}:=D_\eta^2S[v,w].
\]

\begin{theorem}[Complete twice-marked conditional cumulant]
\label{thm:v4v55-two-mark}
For retained directions \(v,w\),
\begin{align}
 D_\eta^2\kappa^\eta(U_1,\ldots,U_n)[v,w]
 &=
 \sum_i
 \kappa^\eta(U_1,\ldots,U_{i;vw},\ldots,U_n)\nonumber\\
 &\quad+
 \sum_{i\ne j}
 \kappa^\eta(U_1,\ldots,U_{i;v},\ldots,U_{j;w},\ldots,U_n)\nonumber\\
 &\quad-
 \sum_i
 \kappa^\eta(U_1,\ldots,U_{i;v},\ldots,U_n,G_w)\nonumber\\
 &\quad-
 \sum_j
 \kappa^\eta(U_1,\ldots,U_{j;w},\ldots,U_n,G_v)\nonumber\\
 &\quad-
 \kappa^\eta(U_1,\ldots,U_n,H_{vw})\nonumber\\
 &\quad+
 \kappa^\eta(U_1,\ldots,U_n,G_v,G_w).
\label{eq:v4v55-two-mark}
\end{align}
If the \(U_i\) have no explicit retained-field dependence, this reduces to
\begin{equation}
 \boxed{
 D_\eta^2\kappa^\eta(U_1,\ldots,U_n)[v,w]
 =
 -\kappa^\eta(U_1,\ldots,U_n,H_{vw})
 +\kappa^\eta(U_1,\ldots,U_n,G_v,G_w).}
\label{eq:v4v55-pure-two-mark}
\end{equation}
\end{theorem}

\begin{proof}
Apply Lemma~\ref{lem:v4v55-one-derivative} first in direction \(v\), then
differentiate every resulting cumulant in direction \(w\).  When \(w\)
hits \(U_{i;v}\), it gives \(U_{i;vw}\); when it hits another insertion, it
gives the ordered pair \(U_{i;v},U_{j;w}\).  Differentiation of the
conditional law appends \(G_w\) with a minus sign.  Finally,
\[
 D_\eta G_v[w]=H_{vw}.
\]
Differentiating the first score term therefore gives the last three lines,
with the double-score sign positive.  This proves
\eqref{eq:v4v55-two-mark}; setting all explicit insertion derivatives to
zero proves \eqref{eq:v4v55-pure-two-mark}.
\end{proof}

\begin{firewall}[A second score is not two first scores]
\label{fw:v4v55-second-not-square}
The last two terms of \eqref{eq:v4v55-pure-two-mark} are distinct
coefficients of the twice-differentiated normalized density.  The contact
term contains the actual Hessian \(H_{vw}\).  The separated term contains
the two actual scores \(G_v,G_w\).  Neither term can be omitted, and no
bound on \(G_v\) and \(G_w\) proves a bound on \(H_{vw}\).
\end{firewall}

\subsection{Spatial seminorms for the two mark species}
\label{sec:v4v55-two-seminorms}

For directions supported in \(Y,Z\), retain the first-score seminorms
\(B_D^{(1)}(Y)\), \(B_D^{(1)}(Z)\) from
\eqref{eq:v4v54-score-seminorm}.  Define the centered second score
\[
 \widetilde H_{Y,Z;v,w}
 :=
 D_\eta^2S[v_Y,w_Z]
 -\mathbb E^\eta D_\eta^2S[v_Y,w_Z]
\]
and its integrated-field Lipschitz seminorm
\begin{equation}
 B_D^{(2)}(Y,Z):=
 \sup_{\eta,\ \|v_Y\|,\|w_Z\|\le1}
 \left(
  \|\nabla\operatorname{Re}\widetilde H_{Y,Z;v,w}\|_\infty+
  \|\nabla\operatorname{Im}\widetilde H_{Y,Z;v,w}\|_\infty
 \right).
\label{eq:v4v55-second-score-seminorm}
\end{equation}

\begin{theorem}[Sufficient twice-marked Lipschitz tree interface]
\label{thm:v4v55-sufficient-two-tree}
Suppose child carriers obey \eqref{eq:v4v54-child-scales}, first scores obey
\eqref{eq:v4v54-score-decay}, and the second score has a junction
decomposition
\begin{equation}
 B_D^{(2)}(Y,Z)
 \le
 J_2\sum_x
 e^{-m_2[d(D,x)+d(x,Y)+d(x,Z)]}.
\label{eq:v4v55-second-score-tree}
\end{equation}
Then both terms in \eqref{eq:v4v55-pure-two-mark} have
factorial-exponential bounds times a tree on all children and the two marked
supports.  The two-first-score term is obtained by applying
Theorem~\ref{thm:v4v53-Lipschitz-cumulant} to
\[
 U_1,\ldots,U_n,\widetilde G_v,\widetilde G_w,
\]
and the second-score term by applying it to
\[
 U_1,\ldots,U_n,\widetilde H_{vw}.
\]
\end{theorem}

\begin{proof}
For two first scores, Theorem
\ref{thm:v4v53-Lipschitz-cumulant} gives the product of all child Lipschitz
scales and both first-score scales.  Repeat the nearest-junction star
argument of Theorem~\ref{thm:v4v54-marked-tree} on the children and the two
marks.

For the second score, insert
\eqref{eq:v4v55-second-score-tree} into the V53 product bound.  For each
junction \(x\), the two factors joining \(Y\) and \(Z\) to \(x\), together
with the child-to-\(D\) factors, form a connected graph.  Delete redundant
edges to obtain a spanning tree and decrease the exponential rate if
necessary.  The sum over \(x\) is finite by exponential convolution.
\end{proof}

\begin{assessment}[The first-score channel is acquired; the second-score
seminorm is open]
\label{ass:v4v55-second-score-open}
The two-first-score term is controlled by V53--V54 and
\eqref{eq:v4v54-score-decay}.  For the original finite-range quadratic
cross-color Higgs stencil, \(D_\eta^2S[v,w]\) is independent of the
integrated field in the mixed directions, so its centered second score
vanishes.  The promoted V34 good modulus contributes a nontrivial second
score.  Bounding \eqref{eq:v4v55-second-score-seminorm} differentiates its
field Hessian once more.  V34 proves the weighted Hessian estimate
\eqref{eq:v4v34-weighted-log-Hessian}, but not the required weighted third
derivative.  Thus \eqref{eq:v4v55-second-score-tree} remains an explicit
good-modulus acquisition.
\end{assessment}

\begin{firewall}[The V34 Hessian bound cannot be differentiated for free]
\label{fw:v4v55-no-free-third}
A uniform bound on \(D_\phi^2\log|\mathcal D_{\rm G}|\) does not imply a
bound on \(D_\phi^3\log|\mathcal D_{\rm G}|\).  The latter contains one
genuine third derivative of the smooth local cutoff insertion, mixed
products of its first and second derivatives, and three ordered first
derivatives separated by good resolvents.  These words must be written and
estimated before \eqref{eq:v4v55-second-score-tree} is claimed.
\end{firewall}

\subsection{V55 result ledger and V56 target}
\label{sec:v4v55-conclusion}

V55 records the exact active YM V71 and unchanged NS V31 fingerprints and
imports their one-occurrence proof order.  It proves the general derivative
rule for conditional cumulants and the complete twice-marked identity
\eqref{eq:v4v55-two-mark}.  The two-first-score channel has the V54 marked
tree bound.  The one-second-score channel is reduced to the precise
weighted third-derivative estimate
\eqref{eq:v4v55-second-score-tree}.

V56 should return to the V34 good determinant
\(\ell_{\rm G}=\operatorname{Tr}\log(1+CV_{\rm G})\), differentiate it three
times without symmetrically duplicating ordered words, and prove an
exponentially weighted trilinear kernel bound.  The local derivative sizes
\[
 D V_{\rm G}=O(\epsilon_Y),\qquad
 D^2V_{\rm G}=O(\epsilon_Y^{1+\alpha})
\]
are already established.  The missing local input is the third derivative
of the smooth cutoff insertion; after it is acquired, the resolvent words
must be assigned their actual three mark positions.


\clearpage
\section{V56: the weighted third derivative of the promoted good determinant}
\label{sec:v4v56}

V55 reduces the one-second-score branch to a weighted third field derivative
of the promoted V34 good determinant modulus.  This note acquires that
derivative.  The smooth cutoff contributes one additional inverse power of
the good-field radius at each derivative.  The trace logarithm then has six
exact word types: one third-derivative contact, three mixed first/second
words, and two ordered three-first-derivative words.

Keeping these words before taking absolute values gives an exponentially
weighted trilinear kernel.  Its local sizes are respectively
\[
 \epsilon_Y^{1+2\alpha},\qquad
 \epsilon_Y^{2+\alpha},\qquad
 \epsilon_Y^3.
\]
This proves the V55 second-score junction interface and closes the
twice-marked collision tree for the smooth good-field sector.

\subsection{Third derivative of the smooth local insertion}
\label{sec:v4v56-local-third}

Recall \(R_Y=\epsilon_Y^{-\alpha}\),
\(g_x=\chi(r_x/R_Y)\), and the affine Yukawa insertion \(V_x\) from V34.
Write
\[
 V_{{\rm G},x}:=g_xV_x,
 \qquad
 V_{\rm G}=\sum_xV_{{\rm G},x}.
\]

\begin{lemma}[Third derivative of one good cell]
\label{lem:v4v56-good-third}
There is \(C_{\chi,3}<\infty\), independent of volume, \(x\), and
\(\epsilon_Y\), such that
\begin{equation}
 \|D_\phi^3V_{{\rm G},x}[h,k,l]\|
 \le
 C_{\chi,3}L\epsilon_Y^{1+2\alpha}
 |h_{X_x}||k_{X_x}||l_{X_x}|.
\label{eq:v4v56-good-third}
\end{equation}
The derivative is supported where \(R_Y\le r_x\le2R_Y\).
\end{lemma}

\begin{proof}
Every derivative of \(\chi(r_x/R_Y)\) is supported in the transition
annulus.  There \(r_x\asymp R_Y\), and direct differentiation of
\(r_x=(1+|\phi_{X_x}|^2)^{1/2}\) gives
\[
 \|Dr_x\|\le1,\qquad
 \|D^2r_x\|\le C R_Y^{-1},\qquad
 \|D^3r_x\|\le C R_Y^{-2}.
\]
The chain rule therefore gives
\[
 \|D^2g_x\|\le C_\chi R_Y^{-2},
 \qquad
 \|D^3g_x\|\le C_\chi R_Y^{-3}.
\]
Since \(V_x\) is affine in \(\phi\), \(D^2V_x=D^3V_x=0\), and
\begin{align*}
 D^3(g_xV_x)[h,k,l]
 &=
 D^3g_x[h,k,l]V_x\\
 &\quad+
 D^2g_x[h,k]DV_x[l]
 +D^2g_x[h,l]DV_x[k]
 +D^2g_x[k,l]DV_x[h].
\end{align*}
On the transition annulus,
\(\|V_x\|\le CL\epsilon_YR_Y\) and
\(\|DV_x[q]\|\le L\epsilon_Y|q_{X_x}|\).
Every displayed term is therefore at most
\[
 CL\epsilon_YR_Y^{-2}
 =CL\epsilon_Y^{1+2\alpha}
\]
times the three local direction norms.
\end{proof}

\begin{assessment}[The local derivative hierarchy]
\label{ass:v4v56-local-hierarchy}
Together with Lemma~\ref{lem:v4v34-GT-derivatives}, V56 gives
\begin{align}
 DV_{\rm G}&=O(\epsilon_Y),\label{eq:v4v56-D1}\\
 D^2V_{\rm G}&=O(\epsilon_Y^{1+\alpha}),\label{eq:v4v56-D2}\\
 D^3V_{\rm G}&=O(\epsilon_Y^{1+2\alpha})\label{eq:v4v56-D3}
\end{align}
as local multilinear insertions with bounded overlap.  These are derivative
sizes, not three independent activity factors.
\end{assessment}

\subsection{The exact third trace-log word}
\label{sec:v4v56-trace-word}

Put
\[
 Q_{\rm G}:=(1+CV_{\rm G})^{-1},
 \qquad
 R_{\rm G}:=Q_{\rm G}C,
 \qquad
 \ell_{\rm G}:=\operatorname{Tr}\log(1+CV_{\rm G}).
\]
Field differentiation gives
\begin{equation}
 D_\phi R_{\rm G}[h]
 =-R_{\rm G}D_\phi V_{\rm G}[h]R_{\rm G}.
\label{eq:v4v56-R-derivative}
\end{equation}

\begin{theorem}[Exact third derivative of the good trace logarithm]
\label{thm:v4v56-exact-third}
For field directions \(h,k,l\),
\begin{align}
 D_\phi^3\ell_{\rm G}[h,k,l]
 &=
 \operatorname{Tr}(R_{\rm G}V_{hkl})\nonumber\\
 &\quad-
 \operatorname{Tr}(R_{\rm G}V_hR_{\rm G}V_{kl})
 -\operatorname{Tr}(R_{\rm G}V_kR_{\rm G}V_{hl})
 -\operatorname{Tr}(R_{\rm G}V_lR_{\rm G}V_{hk})\nonumber\\
 &\quad+
 \operatorname{Tr}(R_{\rm G}V_hR_{\rm G}V_kR_{\rm G}V_l)
 +\operatorname{Tr}(R_{\rm G}V_hR_{\rm G}V_lR_{\rm G}V_k),
\label{eq:v4v56-exact-third}
\end{align}
where \(V_h=D_\phi V_{\rm G}[h]\),
\(V_{hk}=D_\phi^2V_{\rm G}[h,k]\), and
\(V_{hkl}=D_\phi^3V_{\rm G}[h,k,l]\).
\end{theorem}

\begin{proof}
Differentiate the exact V34 second derivative
\[
 D^2\ell_{\rm G}[h,k]
 =
 \operatorname{Tr}(R_{\rm G}V_{hk})
 -\operatorname{Tr}(R_{\rm G}V_kR_{\rm G}V_h)
\]
in direction \(l\).  Use \eqref{eq:v4v56-R-derivative} at each of the two
resolvent positions and the product rule at each insertion.  Cyclicity of
the finite trace rewrites the result in the displayed symmetric form.  The
two final words are the two orders of \(k,l\) between the fixed first
insertion \(h\); they must both be retained.
\end{proof}

\begin{firewall}[Six word types, with no symmetric overcount]
\label{fw:v4v56-six-words}
The three mixed terms correspond to the three choices of the singleton
first derivative.  The last line contains the two cyclically inequivalent
orders of three first derivatives.  Writing all six permutations in the
trace would count each cyclic word three times.  Replacing the last line by
the square of a second-derivative estimate would lose the ordered
three-mark ancestry.
\end{firewall}

\subsection{Weighted trilinear kernel}
\label{sec:v4v56-weighted-kernel}

Let \(T_{xyz}\) be the block kernel of the trilinear form
\(D_\phi^3\ell_{\rm G}\), so that
\[
 D_\phi^3\ell_{\rm G}[h,k,l]
 =
 \sum_{x,y,z}T_{xyz}[h_x,k_y,l_z].
\]
Use the tree weight
\begin{equation}
 \mathfrak d(x;y,z):=
 \min\{
  d(x,y)+d(x,z),\
  d(y,x)+d(y,z),\
  d(z,x)+d(z,y)
 \}.
\label{eq:v4v56-tree-distance}
\end{equation}

\begin{theorem}[Exponentially weighted third-derivative bound]
\label{thm:v4v56-weighted-third}
There are \(a_2>0\) and \(C_3<\infty\), independent of volume, such that
\begin{align}
 \sup_x\sum_{y,z}
 e^{a_2\mathfrak d(x;y,z)}
 \|T_{xyz}\|
 &\le
 C_3\left(
  \epsilon_Y^{1+2\alpha}
  +\epsilon_Y^{2+\alpha}
  +\epsilon_Y^3
 \right).
\label{eq:v4v56-weighted-third}
\end{align}
The same estimate holds for
\(D_\phi^3\log|\mathcal D_{\rm G}|=\operatorname{Re}D_\phi^3\ell_{\rm G}\).
\end{theorem}

\begin{proof}
The good resolvent \(R_{\rm G}\) has the uniform exponentially weighted row
and column bounds established in V34.  The contact word in
\eqref{eq:v4v56-exact-third} contains one local insertion
\(V_{hkl}\).  Lemma~\ref{lem:v4v56-good-third} and bounded cell overlap give
the first power on the right of
\eqref{eq:v4v56-weighted-third}.

Each mixed word contains one \(V_h=O(\epsilon_Y)\), one
\(V_{kl}=O(\epsilon_Y^{1+\alpha})\), and two good-resolvent kernels.  Weighted
Schur convolution joins their two local cells and gives
\(O(\epsilon_Y^{2+\alpha})\).  Each final word contains three
\(O(\epsilon_Y)\) local insertions and three resolvent kernels.  Choose two
of the three resolvent paths forming a spanning tree of their cells and
sum the remaining kernel with its unweighted Schur norm.  This gives the
tree distance \eqref{eq:v4v56-tree-distance} and
\(O(\epsilon_Y^3)\).

Choose \(a_2\) below the V34 resolvent decay rate.  The weighted convolution
algebra and bounded overlap make all row sums uniform in volume.  Taking
real parts cannot increase the absolute kernel norm.
\end{proof}

\begin{corollary}[Acquisition of the V55 second-score junction]
\label{cor:v4v56-second-score}
For the promoted good-modulus contribution to the conditional action,
\eqref{eq:v4v55-second-score-tree} holds with
\begin{equation}
 J_2
 \le
 C_3\left(
  \epsilon_Y^{1+2\alpha}
  +\epsilon_Y^{2+\alpha}
  +\epsilon_Y^3
 \right)
\label{eq:v4v56-J2}
\end{equation}
and an exponential rate below \(a_2\).  Hence both terms of the pure
twice-marked identity \eqref{eq:v4v55-pure-two-mark} carry a tree on the
children and the two marks.
\end{corollary}

\begin{proof}
The second score is the restriction of
\(D_\phi^2[-\log|\mathcal D_{\rm G}|]\) to the two retained directions.
Its integrated-field gradient is the corresponding restriction of the
third derivative bounded in
\eqref{eq:v4v56-weighted-third}.  Resolve its three block indices into the
conditional component and the two marked supports.  The tree-distance
weight gives exactly the junction convolution in
\eqref{eq:v4v55-second-score-tree}.  Invoke
Theorem~\ref{thm:v4v55-sufficient-two-tree}.
\end{proof}

\begin{firewall}[The phase and tail are outside the third derivative]
\label{fw:v4v56-sector-boundary}
V56 differentiates only the promoted positive modulus of
\(\mathcal D_{\rm G}\).  The unit-modulus phase is not part of the
log-Sobolev action.  Tail cells remain in the V34 tail determinant with
their separate super-small parameter.  No derivative word in
\eqref{eq:v4v56-exact-third} is charged to either species.
\end{firewall}

\subsection{V56 result ledger and V57 target}
\label{sec:v4v56-conclusion}

V56 proves the local third derivative
\eqref{eq:v4v56-good-third}, the exact six-type trace-log identity
\eqref{eq:v4v56-exact-third}, and the volume-uniform weighted trilinear
kernel bound \eqref{eq:v4v56-weighted-third}.  This acquires the one-second-
score junction left open in V55.  Together with the two-first-score channel,
the bounded positive good-field collision hierarchy is now stable through
two retained-field marks.

V57 should insert the resulting zero-, one-, and two-mark collision trees
into the V43 normalization cocycle and prove a Kotecky--Preiss bound for the
complete good-field scalar pressure species.  The proof must keep the V39
fermion tree, V53 Higgs collision tree, V34 tail species, and residual phase
species as distinct typed edges until their separate row sums are combined.


\clearpage
\section{V57: a rooted star-merger sum closes the positive good-field pressure species}
\label{sec:v4v57}

V53--V56 give exact collision histories whose local \(k\)-child merger is
paid by a star tree and has coefficient \(k!C^k\), through two retained-field
marks.  A coarse replacement by \(n!C^n\) times the sum of all Cayley trees
would be too large for a polymer sum.  The correct summation retains the
hierarchy.  At each merger the children are an unordered set, producing
\(1/k!\); this cancels the local \(k!\).  The \(k\) choices of star center
leave only a polynomial branching factor.

This note proves the rooted majorant and applies it to the scalar
normalization of the bounded positive good-field species.  Its pressure,
local ratios, and first two marked pressure derivatives have
volume-uniform absolutely convergent expansions.  The V34 tail and the
residual complex phase remain separate species and are not claimed closed by
this positive argument.

\subsection{Abstract rooted star-merger species}
\label{sec:v4v57-abstract}

Let \(\mathbb B\) be the block lattice and let
\(w(b,c)=w(c,b)\ge0\) satisfy
\begin{equation}
 W:=\sup_{b\in\mathbb B}\sum_{c\ne b}w(b,c)<\infty.
\label{eq:v4v57-row-sum}
\end{equation}
A leaf at block \(b\) has norm at most \(\varepsilon\).  An internal merger
with an unordered set of \(k\ge2\) child species has norm at most
\begin{equation}
 A_*k!C_*^k
 \sum_{r=1}^k
 \prod_{i\ne r}w(b_i,b_r)
 \prod_{i=1}^kN_i,
\label{eq:v4v57-local-star}
\end{equation}
where \(b_i\) is the anchor of child \(i\) and \(N_i\) is its inherited
rooted norm.

Let \(\mathcal R_b\) be the sum of absolute weights of all rooted merger
histories whose distinguished leaf is at \(b\), including a factor
\(e^{a|X|}\) for their total leaf support.  Put
\[
 R:=\sup_b\mathcal R_b.
\]

\begin{theorem}[Rooted star-merger majorant]
\label{thm:v4v57-rooted-majorant}
The rooted sum is majorized by the minimal nonnegative solution of
\begin{equation}
 R
 \le
 e^a\varepsilon+
 A_*\sum_{k\ge2}
 kC_*^kW^{k-1}R^k.
\label{eq:v4v57-rooted-recursion}
\end{equation}
There is \(\varepsilon_*=\varepsilon_*(a,A_*,C_*,W)>0\) such that, for
\(\varepsilon<\varepsilon_*\), this solution is finite and
\begin{equation}
 R\le2e^a\varepsilon.
\label{eq:v4v57-rooted-small}
\end{equation}
\end{theorem}

\begin{proof}
At a merger with \(k\) children, labelled-set expansion of the unordered
child family contributes \(1/k!\).  This cancels the \(k!\) in
\eqref{eq:v4v57-local-star}.  Choose the star center in \(k\) ways.  Hold
its anchor fixed and sum every other child anchor with
\eqref{eq:v4v57-row-sum}; this gives \(W^{k-1}\).  Each child history is
bounded by \(R\).  The support weight is already multiplicative across the child norms; only
the one-leaf term changes from \(\varepsilon\) to \(e^a\varepsilon\).  No
second factor \(e^{ak}\) is charged at a merger.  Summing \(k\) proves
\eqref{eq:v4v57-rooted-recursion}.

The nonlinear series is
\begin{equation}
 \sum_{k\ge2}kx^k
 =\frac{x^2(2-x)}{(1-x)^2},
 \qquad |x|<1,
\label{eq:v4v57-branch-series}
\end{equation}
with \(x=C_*WR\).  It is analytic and quadratic at the origin.  Choose
\(r>0\) so that \(C_*Wr<1/2\) and the nonlinear part of the right side of
\eqref{eq:v4v57-rooted-recursion} is at most \(r/2\) for \(R\le r\).
Then choose \(e^a\varepsilon\le r/2\).  Monotone iteration from zero stays
in \([0,r]\) and converges to a finite fixed point.  Reducing
\(\varepsilon_*\) once more makes the nonlinear part at
\(R=2e^a\varepsilon\) at most \(e^a\varepsilon\), proving
\eqref{eq:v4v57-rooted-small}.
\end{proof}

\begin{firewall}[Do not flatten before summing the branching factorial]
\label{fw:v4v57-no-coarse-Cayley}
The proof uses the local \(k!\), the unordered-child \(1/k!\), and the
\(k\) star roots at the same merger.  Replacing the hierarchy first by
\[
 n!C^n\sum_{T\text{ on }[n]}\prod_{\{i,j\}\in T}w(i,j)
\]
would retain both the global factorial and the Cayley count and would not
imply \eqref{eq:v4v57-rooted-recursion}.  The hierarchical star sum is the
stronger information and must be retained until the rooted block positions
are summed.
\end{firewall}

\subsection{Polynomial mark factors preserve the radius}
\label{sec:v4v57-marks}

The V54 one-mark merger has a \((k+1)!\) coefficient and the V55--V56
two-mark mergers have \((k+2)!\) coefficients.  After division by the
unordered-child \(k!\), these add respectively
\((k+1)\) and \((k+1)(k+2)\).

\begin{lemma}[One and two marks have the same convergence radius]
\label{lem:v4v57-marked-series}
For \(j=0,1,2\), the branching series
\begin{equation}
 \sum_{k\ge2}
 k(1+k)^j x^k
\label{eq:v4v57-marked-series}
\end{equation}
is analytic for \(|x|<1\).  Consequently the rooted sums with at most two
marks converge for sufficiently small \(\varepsilon\), with the same
strict smallness condition \(C_*WR<1\).
\end{lemma}

\begin{proof}
Apply \((x\partial_x)^j\) to the rational series
\eqref{eq:v4v57-branch-series}.  This changes the pole order at \(x=1\) but
not the radius of convergence.  The finite choices of whether a mark is a
child mark, first score, second score, or ordered pair of first scores add
only another fixed polynomial in \(k\), already covered by increasing the
power in \eqref{eq:v4v57-marked-series}.
\end{proof}

\begin{theorem}[Rooted good-field histories through two marks]
\label{thm:v4v57-rooted-marked}
Assume the V53 value, V54 one-mark, and V55--V56 two-mark local star bounds.
For sufficiently small leaf activity \(\varepsilon\), the unmarked, once
marked, and twice marked rooted history sums are finite uniformly in volume.
Their bounds are respectively
\begin{equation}
 O(\varepsilon),\qquad O(\varepsilon),\qquad O(\varepsilon),
\label{eq:v4v57-three-rooted-bounds}
\end{equation}
with mark-direction norms factored out.
\end{theorem}

\begin{proof}
Use Theorem~\ref{thm:v4v57-rooted-majorant} for the unmarked subtrees.
At the unique topmost marked merger, use the appropriate marked series from
Lemma~\ref{lem:v4v57-marked-series}; below it, every branch is bounded by
the unmarked fixed point \(R=O(\varepsilon)\).  If two marks lie in distinct
child branches, their ordered placement is one of the explicit terms in
\eqref{eq:v4v55-two-mark}; if they meet at one merger, use either the V56
second-score word or the two-first-score word.  Each history has exactly one
of these ancestries.  The analytic marked series and
\eqref{eq:v4v57-rooted-small} give
\eqref{eq:v4v57-three-rooted-bounds}.
\end{proof}

\subsection{Good-field scalar normalization}
\label{sec:v4v57-pressure}

Let \(Y_b\) be the scalar good-field block perturbation after the V43
normalization functional is applied and its scalar mean is separated.  For
finite \(\Lambda\subset\mathbb B\), introduce independent bookkeeping
variables \(z_b\):
\begin{equation}
 \mathcal Z_{\Lambda}^{\rm G}(z)
 :=
 \mathbb E_{\nu_{\rm G}}
 \prod_{b\in\Lambda}(1+z_bY_b).
\label{eq:v4v57-good-Z}
\end{equation}
The coefficient of \(\prod_{b\in X}z_b\) in
\(\log\mathcal Z_{\Lambda}^{\rm G}(z)\) is the joint cumulant
\begin{equation}
 \zeta_{\rm G}(X)
 :=
 \kappa_{\nu_{\rm G}}(Y_b:b\in X).
\label{eq:v4v57-good-polymer}
\end{equation}
The V49 collision-history identity decomposes this coefficient exactly into
the merger histories summed above.

\begin{theorem}[Positive good-field scalar pressure and two marks]
\label{thm:v4v57-good-pressure}
Assume:
\begin{enumerate}
\item the scalar leaf norm satisfies
      \(\|Y_b\|_{\rm leaf}\le\varepsilon_{\rm G}\), with
      \(\varepsilon_{\rm G}\le C\epsilon_Y^{1-\alpha}\);
\item the collision kernels have the uniform row sum
      \eqref{eq:v4v57-row-sum}; and
\item the V53--V56 local star bounds hold at every conditional stage.
\end{enumerate}
For sufficiently small \(\epsilon_Y\),
\begin{equation}
 \sup_b
 \sum_{\substack{X\ni b\\X\ {\rm finite}}}
 e^{a|X|}|\zeta_{\rm G}(X)|
 \le C_a\varepsilon_{\rm G}.
\label{eq:v4v57-KP}
\end{equation}
Consequently:
\begin{enumerate}
\item \(\mathcal Z_\Lambda^{\rm G}(1)\ne0\) on the branch continued from
      zero coupling;
\item \(|\Lambda|^{-1}\log\mathcal Z_\Lambda^{\rm G}(1)\) has a
      translation-covariant van Hove limit;
\item normalized local good-field ratios have thermodynamic limits; and
\item the first two background derivatives of the good-field pressure and
      local ratios are obtained by termwise differentiation and are
      volume-uniform.
\end{enumerate}
\end{theorem}

\begin{proof}
The exact collision-history decomposition and
Theorem~\ref{thm:v4v57-rooted-marked} give
\eqref{eq:v4v57-KP}; choose \(\epsilon_Y\) so that the leaf activity lies
below the rooted threshold.  Absolute convergence of the rooted polymer
sum gives a zero-free logarithmic branch by continuation from
\(\epsilon_Y=0\).  Translation covariance and the exponential support
weight imply that boundary polymers have vanishing density along a van Hove
sequence, while dominated convergence gives the pressure limit.

For a local insertion, root every connected coefficient at its support.
The same bound makes the exterior tail summable and yields a local limit.
Lemma~\ref{lem:v4v57-marked-series} and
Theorem~\ref{thm:v4v57-rooted-marked} justify the first two termwise
derivatives.  The V43 normalization cocycle assigns every scalar constant
once, and Theorem~\ref{thm:v4v43-local-cancellation} cancels those constants
from normalized local ratios.
\end{proof}

\begin{assessment}[Scope of the acquired pressure]
\label{ass:v4v57-pressure-scope}
Theorem~\ref{thm:v4v57-good-pressure} closes the scalar pressure species
generated by the bounded positive V34 good-field carrier under the stated
all-stage collision interface.  The first physical stage and its first two
marks are proved in V53--V56.  Preservation at every later rescaled stage
requires the same uniform block geometry and row sums; this is an explicit
renormalization-scale hypothesis, not a theorem of V57.
\end{assessment}

\begin{firewall}[The complete complex normalization is not yet closed]
\label{fw:v4v57-complex-open}
The positive good-field pressure is only one typed species.  The V34 tail
has its already proved super-small local parameter and fermion tree, but its
joint composition with every good-field collision scale has not yet been
inserted into \eqref{eq:v4v57-rooted-recursion}.  The residual determinant
phase is complex and lies outside the positive log-Sobolev law.  Therefore
V57 does not claim the full V43 complex cocycle pressure or the full
Standard Model continuum.
\end{firewall}

\subsection{V57 result ledger and V58 target}
\label{sec:v4v57-conclusion}

V57 proves that the actual hierarchical star structure cancels local merger
factorials before block positions are summed.  The resulting rooted
branching series is analytic, and polynomial factors from one or two marks
do not reduce its radius.  Under the explicit all-scale good-field collision
interface, this yields a Kotecky--Preiss bound, thermodynamic scalar pressure,
local ratios, and two marked derivatives for the bounded positive
good-field species.

V58 should compose the V34 rare-tail species with the good-field rooted
recursion.  Every mixed history must choose its first tail cell, pay the
factor \(e^{-c_{\rm T}\epsilon_Y^{-2\alpha}}\) there once, and then use the
ordinary typed fermion/Higgs row sums.  Only after that mixed positive
modulus is closed should the programme return to the residual complex phase.


\clearpage
\section{V58: tail-owned rooted composition and a quantitative rare-sector correction}
\label{sec:v4v58}

V57 summed the hierarchical collision histories of the bounded good-field
species.  V34 had already supplied a super-small activity for every
connected hard-scale coefficient containing an actual Yukawa tail cell,
and V36 had absorbed the good Neumann strings inside the updated covariance
\(C_{\rm G}\).  What was missing was a proof that inserting such tail
coefficients into the V57 hierarchy does not spend their rare-event factor
again at each later collision.

This note supplies the combinatorial composition theorem.  A formal
variable counts primitive tail components.  Differentiating the rooted
fixed point marks one such component; the derivative is controlled by the
susceptibility of the same subcritical branching map.  Every mixed history
is thereby assigned at least one tail owner and is \(O(\varepsilon_{\rm
T})\), while all other good or tail descendants are paid by the ordinary
rooted sum.  The argument also covers the first two background marks.

The application is conditional only at one clearly identified point:
V34 proves the primitive hard-scale tail norm, whereas the local star
estimate of V53--V56 was proved for bounded good-field carriers.  Its typed
extension to an unbounded tail carrier at every conditional Higgs collision
is stated below as the remaining analytic interface.  V58 proves that no
additional global combinatorial obstruction occurs after that interface is
available.

\subsection{Primitive good and tail families}
\label{sec:v4v58-primitives}

A primitive activity \(q_\sigma(X;b)\), of type
\(\sigma\in\{{\rm G},{\rm T}\}\), has finite support \(X\) and a
distinguished incidence block \(b\in X\).  Incidences are counted rather
than choosing a single arbitrary anchor; this directly controls local
ratios.  Let \(\|\cdot\|_{(j)}\), \(0\le j\le2\), denote the value norm and
the first two background-marked norms, with the direction norms divided
out.  Assume
\begin{align}
 \gamma_j
 &:=
 \sup_b\sum_{X\ni b}e^{a|X|}
 \|q_{\rm G}(X;b)\|_{(j)}
 \le C_j\varepsilon_{\rm G},
\label{eq:v4v58-good-primitive}\\
 \tau_j
 &:=
 \sup_b\sum_{X\ni b}e^{a|X|}
 \|q_{\rm T}(X;b)\|_{(j)}
 \le C_j\varepsilon_{\rm T}.
\label{eq:v4v58-tail-primitive}
\end{align}
For the Yukawa split, a tail primitive includes its complete
\(C_{\rm G}\)-Neumann string.  Thus a good vertex occurring inside
\(C_{\rm G}\) is not also introduced as a separate good primitive.

Suppose a conditional collision of an unordered set of \(k\ge2\) child
histories satisfies the typed version of
\eqref{eq:v4v57-local-star}:
\begin{equation}
 A_*k!C_*^k
 \sum_{r=1}^k
 \prod_{i\ne r}w_{\sigma_i\sigma_r}(b_i,b_r)
 \prod_{i=1}^kN_i,
\label{eq:v4v58-typed-star}
\end{equation}
where
\begin{equation}
 W_*:=
 \max_{\sigma,\sigma'}
 \sup_b\sum_c w_{\sigma\sigma'}(b,c)<\infty.
\label{eq:v4v58-typed-row}
\end{equation}
Through two marks, allow the polynomial branching factors proved in
V54--V56.  In particular, the \(j\)-marked local majorant is bounded by a
fixed constant times \(k!(1+k)^{j+1}C_*^k\) times the same star product.

\subsection{The tail-counting fixed point}
\label{sec:v4v58-fixed-point}

Set
\begin{equation}
 \Phi(r)
 :=
 A_*\sum_{k\ge2}kC_*^kW_*^{k-1}r^k.
\label{eq:v4v58-Phi}
\end{equation}
This is analytic for \(C_*W_*r<1\) and is quadratic at zero.  Introduce
\(s\in[0,1]\), multiplying every primitive tail activity by \(s\), and let
\(R(s)\) be the unmarked rooted absolute history sum.

\begin{theorem}[Tail-owned rooted composition]
\label{thm:v4v58-tail-owned}
Choose \(r_0>0\) so that \(C_*W_*r_0<1\) and
\begin{equation}
 q_*:=\sup_{0\le r\le r_0}\Phi'(r)<1.
\label{eq:v4v58-contraction}
\end{equation}
If
\begin{equation}
 \gamma_0+\tau_0+\Phi(r_0)\le r_0,
\label{eq:v4v58-ball}
\end{equation}
then, uniformly in volume and \(s\in[0,1]\),
\begin{equation}
 R(s)\le\gamma_0+s\tau_0+\Phi(R(s))\le r_0.
\label{eq:v4v58-Rs}
\end{equation}
Moreover, the sum \(M_{\rm T}\) of all rooted histories containing at
least one primitive tail component satisfies
\begin{equation}
 M_{\rm T}
 \le
 \frac{\tau_0}{1-q_*}.
\label{eq:v4v58-tail-susceptibility}
\end{equation}
Thus the rare factor is paid once, independently of the number of later
good-field collisions.
\end{theorem}

\begin{proof}
At every unordered merger, \(1/k!\) from the child set cancels the
\(k!\) in \eqref{eq:v4v58-typed-star}.  There are \(k\) star centers, and
summing the remaining anchors gives \(W_*^{k-1}\).  Since the exponential
support weight is already present in every primitive and inherited child
norm, it creates no additional merger factor.  This proves
\eqref{eq:v4v58-Rs}.  Monotone iteration from zero stays in
\([0,r_0]\) by \eqref{eq:v4v58-ball}.

For a finite-depth history sum, differentiate with respect to \(s\).
The derivative chooses one actual primitive tail component.  If
\(U(s)\) is the resulting one-tail-marked rooted sum, the product rule at
the top merger and the preceding star summation give
\begin{equation}
 U(s)\le\tau_0+\Phi'(R(s))U(s).
\label{eq:v4v58-U}
\end{equation}
Hence \(U(s)\le\tau_0/(1-q_*)\), uniformly in depth.  Every history with
\(n_{\rm T}\ge1\) tail primitives occurs in \(U(1)\) exactly
\(n_{\rm T}\) times and therefore at least once.  Its unmarked absolute
weight is consequently bounded by \(U(1)\).  Passing monotonically to
unbounded depth proves \eqref{eq:v4v58-tail-susceptibility}.

Equivalently, order primitive incidences deterministically and call the
first tail incidence the owner.  The derivative proof merely overcounts
the possible owner by \(n_{\rm T}\); it never asks a single tail incidence
to supply more than its one primitive factor.
\end{proof}

\begin{corollary}[Lipschitz dependence on the tail species]
\label{cor:v4v58-tail-Lipschitz}
Let \(R_{\rm G}\) be the rooted sum with \(\tau_0=0\) and let
\(R_{\rm G+T}\) be the sum with both types.  Under
\eqref{eq:v4v58-contraction}--\eqref{eq:v4v58-ball},
\begin{equation}
 0\le R_{\rm G+T}-R_{\rm G}
 \le\frac{\tau_0}{1-q_*}.
\label{eq:v4v58-rooted-difference}
\end{equation}
\end{corollary}

\begin{proof}
Classify histories according to whether they contain a tail primitive.
The first class is precisely the good history sum; the second is bounded
by Theorem~\ref{thm:v4v58-tail-owned}.
\end{proof}

\subsection{Two background marks}
\label{sec:v4v58-two-marks}

\begin{lemma}[Tail ownership survives two marks]
\label{lem:v4v58-marked-tail-owned}
Assume \eqref{eq:v4v58-good-primitive}--%
\eqref{eq:v4v58-tail-primitive} through \(j=2\) and the marked polynomial
version of \eqref{eq:v4v58-typed-star}.  After reducing
\(\gamma_0+\tau_0\), the once- and twice-marked rooted sums are finite.
The sub-sums containing a tail primitive obey
\begin{equation}
 M_{\rm T}^{(j)}
 \le C_j'\varepsilon_{\rm T},
 \qquad j=0,1,2,
\label{eq:v4v58-marked-tail-bound}
\end{equation}
with constants independent of volume.
\end{lemma}

\begin{proof}
For a fixed collision history, background differentiation has the exact
placements already classified in V54--V56: one first mark; one second
mark; or two first marks on distinct or coincident ancestors.  After the
unordered-child factorial is cancelled, these placements multiply the
\(k\)-th branching coefficient by a fixed polynomial in \(k\).  Therefore
all required series are obtained from \(\Phi\) by finitely many
applications of \(r\partial_r\), and have the same radius
\(C_*W_*r<1\).

Now add the tail-counting variable \(s\).  The value, one-mark, and
two-mark majorants form a triangular system: the value equation involves
only the value sum; the one-mark equation is linear in the one-mark sum;
the two-mark equation is linear in the two-mark sum with a source
quadratic in already bounded one-mark sums.  Differentiate this system in
\(s\).  In each equation the coefficient of its highest marked unknown is
\(\Phi'(R)\), plus a quantity tending to zero with \(R\) from the marked
polynomial series.  Reduce the activity so that this total coefficient is
less than one.  The inhomogeneous terms contain either
\(\tau_j=O(\varepsilon_{\rm T})\) or a lower marked
tail-owned sum.  Induction on \(j=0,1,2\), beginning with
\eqref{eq:v4v58-U}, proves
\eqref{eq:v4v58-marked-tail-bound}.
\end{proof}

\begin{firewall}[A tail owner is an actual primitive component]
\label{fw:v4v58-owner}
The factor \(\varepsilon_{\rm T}\) belongs to a V34 connected coefficient
containing at least one actual cell on which \(t_xV_x\ne0\).  A collision
vertex, a covariance edge, or a good \(C_{\rm G}\)-Neumann factor is not a
second tail event.  Conversely, a primitive containing several actual tail
cells retains all the per-cell factors proved in
\eqref{eq:v4v34-tail-parameter}; the present theorem uses only one of them
to control the mixed hierarchy and does not erase the others.
\end{firewall}

\subsection{Application to the Yukawa normalization}
\label{sec:v4v58-application}

Theorem~\ref{thm:v4v34-tail-small} and the V36 Neumann-word estimate give,
through two background marks,
\begin{equation}
 \varepsilon_{\rm T}
 \le
 K\exp\!\left(-c_{\rm T}\epsilon_Y^{-2\alpha}\right)
\label{eq:v4v58-Yukawa-tail}
\end{equation}
for the complete primitive tail family.  The V57 good family has
\(\varepsilon_{\rm G}=O(\epsilon_Y^{1-\alpha})\).  Combining these facts
with the typed collision interface gives the following consequence.

\begin{theorem}[Rare-tail correction to the good-field branch]
\label{thm:v4v58-tail-pressure}
Assume the all-stage good collision hypotheses of
Theorem~\ref{thm:v4v57-good-pressure} and the typed extension
\eqref{eq:v4v58-typed-star} for collisions having at least one V34 tail
descendant.  Then the branch obtained by inserting the exact tail
determinant into the good-modulus parent has an absolutely convergent
connected logarithm for sufficiently small \(\epsilon_Y\).  If
\(p_{\rm G}\) and \(p_{\rm G+T}\) denote its finite-volume intensive
logarithms, then through two background marks
\begin{equation}
 \left\|
 D_B^j(p_{\rm G+T}-p_{\rm G})
 \right\|
 \le
 C_j
 \exp\!\left(-c_{\rm T}\epsilon_Y^{-2\alpha}\right),
 \qquad j=0,1,2.
\label{eq:v4v58-pressure-difference}
\end{equation}
The same estimate holds for normalized local ratios with an admissible
local insertion, uniformly before the thermodynamic limit.  Hence their
van Hove limits exist whenever the translation-covariant typed interface
is uniform.
\end{theorem}

\begin{proof}
Use V34--V36 to take the complete hard-scale tail coefficients, including
their \(C_{\rm G}\)-Neumann words, as the primitive tail family.  Use the
V57 good coefficients as the good family.  For small \(\epsilon_Y\), their
sum satisfies \eqref{eq:v4v58-ball}.  Lemma
\ref{lem:v4v58-marked-tail-owned} bounds every connected coefficient in
the logarithmic difference by the right side of
\eqref{eq:v4v58-Yukawa-tail}.  Summing its root incidence gives the
extensive bound \eqref{eq:v4v58-pressure-difference}.  Rooting at a local
insertion gives the local-ratio bound.  The exponential support weight
makes boundary-rooted coefficients have vanishing density along a van
Hove sequence.
\end{proof}

\begin{assessment}[What V58 closes and what it leaves]
\label{ass:v4v58-scope}
V58 closes the entire abstract composition after a primitive tail activity
has entered a subcritical typed collision hierarchy.  It proves that later
branching does not destroy the super-small V34 factor.  At the physical
level, \eqref{eq:v4v58-pressure-difference} is conditional on the typed
tail collision estimate.  V34 proves rare hard-scale coefficients and
V53--V56 prove good-carrier collisions, but neither theorem by itself
proves \eqref{eq:v4v58-typed-star} for an unbounded tail descendant under
successive conditional Higgs laws.
\end{assessment}

\begin{firewall}[The good determinant phase remains separate]
\label{fw:v4v58-phase-open}
The tail determinant may be complex; Theorem
\ref{thm:v4v58-tail-pressure} controls it as a small connected perturbation
of the good-modulus branch.  The unit-modulus factor
\(\mathcal D_{\rm G}/|\mathcal D_{\rm G}|\) from
\eqref{fw:v4v34-phase-separate} has not been inserted into the correlated
V57 collision hierarchy.  No claim about reflection positivity or the full
SM--Einstein continuum follows from this tail composition.
\end{firewall}

\subsection{V58 result ledger and V59 target}
\label{sec:v4v58-conclusion}

V58 proves a two-species rooted fixed-point theorem.  The mixed part is the
tail derivative of the subcritical rooted map and is bounded by the
primitive tail norm divided by \(1-\Phi'(R)\).  Polynomial branching
factors preserve the conclusion through two background marks.  Applied to
V34, it shows that the full later collision hierarchy changes the
good-modulus pressure and local ratios by only
\(O(e^{-c_{\rm T}\epsilon_Y^{-2\alpha}})\), provided the typed tail
collision interface holds.

V59 should go upstream to that interface.  It should prove a one-tail
conditional covariance estimate directly from strong convexity and the
simultaneous tail exponential moment, using a smooth tail carrier and
paying the first failed cell before applying the Witten resolvent.  The
bound must be stable under one and two background marks and must not replace
an unbounded tail observable by a false uniform Lipschitz norm.

\clearpage
\section{V59: the first rare-tail collision and the failure of a conditional supremum}
\label{sec:v4v59}

V58 reduces the mixed good/tail hierarchy to a typed collision estimate.
The most tempting estimate would condition on one color, take a supremum
over the retained field, and reuse the V34 tail probability.  That route is
false: an unbounded conditioned boundary can shift the conditional mean
past the fixed Yukawa cutoff.  This note first records the obstruction, then
proves the correct first-collision estimate before taking such a supremum.

The proof uses three existing ingredients without changing their ancestry.
The V34 smooth cutoff ensures that a tail carrier and its field derivative
own an actual event \(r_x\ge R_Y\).  The V31 quadratic exponential moment
pays that event in weighted \(L^2\).  The V30 Witten inverse then supplies
the spatial edge.  Product regrouping and the V48 cut identity extend this
one rare edge through the two genuine background marks.  The result closes
the first full-law typed collision, but it does not yet transfer tail
ownership through an arbitrary conditional boundary.

\subsection{Why uniform conditional rarity is false}
\label{sec:v4v59-conditional-obstruction}

\begin{proposition}[Strong convexity does not give a uniform conditional
tail]
\label{prop:v4v59-no-conditional-sup}
For every \(\delta>0\), the probability measure on \(\mathbb R^2\) with
action
\begin{equation}
 S(x,y)
 =
 \frac12(x-y)^2+\frac\delta2(x^2+y^2)
\label{eq:v4v59-Gaussian-action}
\end{equation}
is uniformly strongly log-concave, but for every \(R<\infty\),
\begin{equation}
 \sup_{y\in\mathbb R}
 \mathbb P(|x|\ge R\mid y)=1.
\label{eq:v4v59-conditional-tail-fails}
\end{equation}
\end{proposition}

\begin{proof}
The Hessian of \eqref{eq:v4v59-Gaussian-action} is
\[
 \begin{pmatrix}1+\delta&-1\\-1&1+\delta\end{pmatrix},
\]
whose eigenvalues are \(\delta\) and \(2+\delta\).  Thus the convexity
constant is \(\delta\).  Conditional on \(y\),
\begin{equation}
 x\mid y
 \sim
 N\!\left(\frac{y}{1+\delta},\frac1{1+\delta}\right).
\label{eq:v4v59-conditional-Gaussian}
\end{equation}
Letting \(|y|\to\infty\) makes the conditional probability of
\(|x|\ge R\) tend to one, which proves
\eqref{eq:v4v59-conditional-tail-fails}.
\end{proof}

\begin{firewall}[No supremum over an unowned boundary field]
\label{fw:v4v59-no-tail-sup}
The unconditional V34 estimate cannot be inserted into
\(\sup_\eta\|\nabla T\|_{L^2(\mu_D^\eta)}\).  If the boundary field
\(\eta\) is large, it must itself be retained as a tail owner, or the
conditional estimate must remain averaged over its skeleton law.  A proof
which drops that ownership before taking the supremum loses the rare
factor exactly as in Proposition~\ref{prop:v4v59-no-conditional-sup}.
\end{firewall}

\subsection{A rare weighted gradient}
\label{sec:v4v59-rare-gradient}

Work under the promoted positive law \(\nu_{\rm G,B}\) of
\eqref{eq:v4v34-promoted-measure}, with the uniform convexity and bounded
mean hypotheses of Theorem~\ref{thm:v4v34-promoted-convexity}.  A
tail-owned local carrier \(T_x\) is one whose value, first two background
derivatives, and one field derivative obey, for \(0\le j\le2\),
\begin{align}
 |D_B^jT_x|
 &\le
 K_j\boldsymbol1_{\{r_x\ge R_Y\}}
 (1+|\phi_{A_x}|)^{p_j}
 e^{c_j\epsilon_Y(1+|\phi_{A_x}|)},
\label{eq:v4v59-tail-envelope}\\
 |\nabla_\phi D_B^jT_x|
 &\le
 K_j\boldsymbol1_{\{r_x\ge R_Y\}}
 (1+|\phi_{A_x}|)^{p_j}
 e^{c_j\epsilon_Y(1+|\phi_{A_x}|)}.
\label{eq:v4v59-tail-gradient-envelope}
\end{align}
Here \(A_x\) has uniformly bounded size and contains the failed Yukawa cell.
The derivative of the smooth cutoff is allowed: its support
\(R_Y\le r_x\le2R_Y\) is contained in the displayed tail event.

\begin{lemma}[The tail event pays the \(L^p\) gradient norm]
\label{lem:v4v59-tail-gradient}
For every fixed \(q<\infty\), there are \(c_q,C_{j,q}>0\) such that
\begin{equation}
 \|D_B^jT_x\|_{L^q(\nu_{\rm G,B})}
 +
 \|\nabla_\phi D_B^jT_x\|_{L^q(\nu_{\rm G,B})}
 \le
 C_{j,q}e^{-c_qR_Y^2},
 \qquad 0\le j\le2.
\label{eq:v4v59-tail-gradient}
\end{equation}
The constants are uniform in volume and \(x\).
\end{lemma}

\begin{proof}
Raise \eqref{eq:v4v59-tail-envelope} and
\eqref{eq:v4v59-tail-gradient-envelope} to the \(q\)-th power.  Their
right sides are a fixed polynomial and a fixed linear exponential on an
actual tail event.  Apply
\eqref{eq:v4v34-marked-tail} with \(n=1\), enlarged \(p\), and enlarged
\(c\).  Taking the \(q\)-th root only changes the positive exponential
constant.  Since \(R_Y=\epsilon_Y^{-\alpha}\), it remains
\(e^{-c_q\epsilon_Y^{-2\alpha}}\).
\end{proof}

The hard-scale Grassmann coefficients of
Theorem~\ref{thm:v4v34-tail-small} satisfy
\eqref{eq:v4v59-tail-envelope}--%
\eqref{eq:v4v59-tail-gradient-envelope}.  Indeed, differentiating a finite
Grassmann coefficient or a convergent \(C_{\rm G}\)-Neumann word adds a
fixed polynomial, while differentiating \(t_x\) remains in the transition
annulus.  Bounded fibre dimension and the V33 weighted kernel algebra keep
the constants volume independent.

\begin{lemma}[A complete product retains its tail payment]
\label{lem:v4v59-tail-product}
Let \(F_1,\ldots,F_m\), \(m\le3\), be local good carriers or actual first
or second score insertions.  Suppose their values and field gradients have
uniform fixed \(L^p\) bounds for every fixed \(p<\infty\).  Then
\begin{equation}
 \left\|
 D_B^jT_x\prod_{\ell=1}^mF_\ell
 \right\|_{L^1(\nu_{\rm G,B})}
 +
 \left\|
 \nabla_\phi\!\left(
 D_B^jT_x\prod_{\ell=1}^mF_\ell
 \right)
 \right\|_{L^2(\nu_{\rm G,B})}
 \le
 C_j e^{-cR_Y^2},
 \quad 0\le j\le2.
\label{eq:v4v59-tail-product}
\end{equation}
\end{lemma}

\begin{proof}
Holder's inequality, Lemma~\ref{lem:v4v59-tail-gradient}, and the
ordinary moment bounds first prove the \(L^1\) term.  Apply the product
rule to the gradient term.  In every summand either the tail derivative
\(\nabla D_B^jT_x\) or the tail value \(D_B^jT_x\) remains.  Use Holder's
inequality with one fixed exponent for each of the at most four factors,
Lemma~\ref{lem:v4v59-tail-gradient} for the tail factor, and the assumed
moment bounds for all other factors.  Renaming the reduced positive
exponent gives \eqref{eq:v4v59-tail-product}.
\end{proof}

\subsection{The first typed collision}
\label{sec:v4v59-first-collision}

\begin{theorem}[One tail--good Witten edge]
\label{thm:v4v59-tail-Witten-edge}
Let \(T_x\) satisfy
\eqref{eq:v4v59-tail-envelope}--%
\eqref{eq:v4v59-tail-gradient-envelope}, and let \(F_y\) be a local carrier
with \(\|\nabla F_y\|_{L^2}\le L_F\).  If their field-gradient supports are
contained in uniformly bounded neighborhoods of \(x,y\), then
\begin{equation}
 \left|
 \operatorname{Cov}_{\nu_{\rm G,B}}(T_x,F_y)
 \right|
 \le
 C L_F
 e^{-cR_Y^2}e^{-m d(x,y)}.
\label{eq:v4v59-tail-Witten-edge}
\end{equation}
\end{theorem}

\begin{proof}
The promoted law has the exponentially decaying one-form inverse proved in
Theorem~\ref{thm:v4v34-promoted-convexity}.  Apply the
Helffer--Sjostrand identity \eqref{eq:v4v30-HS}, insert the two
field-gradient support projections, and use the corresponding version of
\eqref{eq:v4v30-Witten-decay}.  Cauchy--Schwarz gives the product of the
two \(L^2\) gradient norms.  Lemma~\ref{lem:v4v59-tail-gradient} pays the
tail norm.  The bounded support radii only change the constant and the
decay exponent.
\end{proof}

This is stronger than multiplying an ordinary covariance estimate by an
unrelated tail probability: the rare event is inside the actual gradient
on one side of the Witten carrier.

\subsection{Fixed marked order by a geometric cut}
\label{sec:v4v59-marked-cut}

Let \(H_1=T_x,H_2,\ldots,H_n\), \(2\le n\le4\), where the other \(H_i\)
are local good carriers or score insertions satisfying the hypotheses of
Lemma~\ref{lem:v4v59-tail-product}.  Write \(b_i\) for their support
centers and
\[
 {\rm diam}(b_1,\ldots,b_n)
 =
 \max_{i,j}d(b_i,b_j).
\]

\begin{theorem}[One-tail cumulants through order four]
\label{thm:v4v59-tail-cumulants}
There are \(c,m,C_n>0\), uniform in volume, such that
\begin{equation}
 \left|
 \kappa_{\nu_{\rm G,B}}(H_1,\ldots,H_n)
 \right|
 \le
 C_n e^{-cR_Y^2}
 \exp\!\left(
 -\frac{m}{n-1}{\rm diam}(b_1,\ldots,b_n)
 \right),
 \qquad 2\le n\le4.
\label{eq:v4v59-tail-cumulants}
\end{equation}
\end{theorem}

\begin{proof}
Take a minimum spanning tree on the \(n\) support centers and remove a
longest edge.  Its length is at least
\({\rm diam}(b_1,\ldots,b_n)/(n-1)\).  The fundamental cut has separation
equal to that edge length; otherwise a shorter crossing edge would replace
it in the spanning tree.

Apply the exact V48 cut identity
\eqref{eq:v4v48-cut-identity} across this bipartition.  Every nonzero term contains one displayed covariance of complete products,
with a finite ordered-partition coefficient depending only on \(n\).  If
that covariance contains \(T_x\), Lemma~\ref{lem:v4v59-tail-product} makes
the gradient on its tail side \(O(e^{-cR_Y^2})\), while the other gradient
is uniformly bounded.  If \(T_x\) instead lies in one of the moment factors
outside the displayed covariance, the first term on the left of
\eqref{eq:v4v59-tail-product} pays that moment; the displayed covariance
still supplies the cut decay with two ordinary gradient bounds.  Thus every
term has both the one tail payment and the actual Witten edge.  Uniformly
bounded support radii are absorbed into the constants.  There are only
finitely many ordered partitions for \(n\le4\), so their absolute sum is
absorbed into \(C_n\).
\end{proof}

\begin{theorem}[The first tail collision through two genuine background
marks]
\label{thm:v4v59-two-mark-tail-collision}
Let
\[
 \mathcal C_B(T_x,F_y)
 =
 \operatorname{Cov}_{\nu_{\rm G,B}}(T_x,F_y).
\]
Assume the explicit derivatives of \(T_x,F_y\) through order two and the
first and second density scores \(G_v,H_{vw}\) satisfy the local moment and
gradient hypotheses above.  Then, after factoring out the background
direction norms,
\begin{equation}
 \left|D_B^j\mathcal C_B(T_x,F_y)\right|
 \le
 C_j e^{-cR_Y^2}e^{-m_jd(x,y)},
 \qquad j=0,1,2.
\label{eq:v4v59-two-mark-tail-collision}
\end{equation}
\end{theorem}

\begin{proof}
The value case is Theorem~\ref{thm:v4v59-tail-Witten-edge}.  Normalized
differentiation gives exactly
\begin{equation}
 D_v\kappa(T,F)
 =
 \kappa(T_v,F)+\kappa(T,F_v)-\kappa(T,F,G_v).
\label{eq:v4v59-first-mark-identity}
\end{equation}
A second differentiation gives
\begin{align}
 D_wD_v\kappa(T,F)
 &=
 \kappa(T_{vw},F)+\kappa(T_v,F_w)
 -\kappa(T_v,F,G_w)\nonumber\\
 &\quad+
 \kappa(T_w,F_v)+\kappa(T,F_{vw})
 -\kappa(T,F_v,G_w)\nonumber\\
 &\quad-
 \kappa(T_w,F,G_v)-\kappa(T,F_w,G_v)
 -\kappa(T,F,H_{vw})\nonumber\\
 &\quad+
 \kappa(T,F,G_v,G_w).
\label{eq:v4v59-second-mark-identity}
\end{align}
These are the V50 and V55 score identities with all explicit derivative
placements retained.  Every term contains \(T\), \(T_v\), \(T_w\), or
\(T_{vw}\), hence one actual tail owner.  Apply
Theorem~\ref{thm:v4v59-tail-cumulants} at orders two, three, and four.
For a fixed number of local supports, exponential diameter decay implies
exponential decay in \(d(x,y)\), after reducing the exponent and summing
the finitely many possible local score positions.
\end{proof}

\begin{assessment}[Acquired interface and remaining transfer problem]
\label{ass:v4v59-interface}
Theorem~\ref{thm:v4v59-two-mark-tail-collision} acquires the physical
first tail--good collision, with the same super-small form required in V58
and with both actual background marks.  It is an averaged statement under
the complete promoted law.  It does not imply the conditional
\(L^\infty(d\eta)\) carrier required by the present V49--V57 color
recursion, because Proposition~\ref{prop:v4v59-no-conditional-sup} rules
out that implication for an absolute Higgs cutoff.
\end{assessment}

\begin{firewall}[Fixed marked order is not an all-order typed star theorem]
\label{fw:v4v59-fixed-order}
The V48 cut identity is used here only at cumulant orders \(2,3,4\), which
are exactly the orders generated by differentiating one covariance twice.
Its constants depend on that fixed order.  This does not prove
\eqref{eq:v4v58-typed-star} for arbitrary merger valence or at every later
conditional stage.
\end{firewall}

\subsection{V59 result ledger and compulsory V60 target}
\label{sec:v4v59-conclusion}

V59 proves that the smooth V34 tail cutoff pays the value and field
gradient of a primitive tail coefficient in every fixed \(L^p\) norm.  The
actual Witten carrier then yields a spatial tail--good covariance with the
super-small factor inside the covariance.  A minimum-spanning-tree cut and
the exact V48 identity propagate this factor through the third and fourth
cumulants required by one and two background marks.

V59 also proves that uniform conditional rarity is false even for a
two-variable strongly convex Gaussian.  V60 must therefore begin with the
scheduled YM/NS audit and then replace the impossible conditional supremum
by boundary-tail transfer: a large conditional mean must be charged to an
actual retained-field tail cell, while the complementary bounded-boundary
sector may use the conditional Witten estimate.

\clearpage
\section{V60: companion audit and boundary-tail transfer before the conditional supremum}
\label{sec:v4v60}

V59 proves the first tail collision under the full promoted Higgs law and
shows that a fixed absolute Yukawa cutoff is not rare uniformly over an
unbounded conditioned boundary.  The correct repair is to split the
boundary before taking a conditional supremum.  On the bounded-boundary
piece, exponential response of the conditional mean restores the V34 tail
payment.  On its complement, a unique retained boundary cell owns a
super-small guard factor and remains attached to the same Witten
covariance.  This is a one-use transfer of rarity, not a second determinant
tail.

This note performs the scheduled companion audit, constructs an exact
smooth first-owner partition, and proves the resulting one-step conditional
tail edge through two background marks and one retained-field derivative.
The carrier norm is averaged over the retained skeleton on the owner branch;
Proposition~\ref{prop:v4v59-no-conditional-sup} shows why that change is
necessary.

\subsection{Compulsory companion audit}
\label{sec:v4v60-audit}

The exact files inspected were:
\begin{enumerate}
\item
\texttt{Renewal\_Yang\_Mills\_Mass\_Gap\_Research\_Notes\_V75.tex},
\(1\,312\,134\) bytes, SHA--256
\[
 \texttt{E741DC0EBED6B41047ED6873111E7EEB80781F104296B1E055EB5301C4A1DF8B}.
\]
Its newest result closes every nondepleted \(2+1+1\) quotient geometry and
reduces the residual cell to simultaneous prefix and quotient depletion.
It proves that separate estimates \(O(p)\) and \(O(q)\) give only
\(O(\min\{p,q\})\), not the required product \(O(pq)\).  The missing card
must retain the prefix covariance and unsplit joining cumulant in one word.
\item
\texttt{Renewal\_NS\_Stability\_Research\_Notes\_V31.tex},
\(887\,537\) bytes, SHA--256
\[
 \texttt{F1121B62238AD5491BB7CF03635EA9934942EDF573ED8FAAA9EE79E1D38A6696}.
\]
Its current endpoint proves that flat dyadic stopping marks are exact, but
removing a mark costs the missing first-moment weight.  Heat lifting pays
that weight only in the critical source norm, while weaker norms return the
known continuation stock.  The remaining opportunity is the undivided
correlation among the actual recent-entry event, signed triad work, and its
high parent.
\end{enumerate}

The transferable lesson is proof order.  For the present SM tail problem,
one may not multiply an unconditional tail estimate by a conditional
Witten estimate after taking separate suprema; V59 gives an explicit
counterexample.  The rare boundary guard and the conditional covariance
must remain in one carrier, exactly as the YM prefix and quotient factors
must remain in one word.  The NS audit also requires a unique owner before
a family of guard marks is summed.  The construction below therefore uses
an exact first-owner partition with thresholds increasing in the distance
from the tail cell.

\subsection{Response of a conditional mean}
\label{sec:v4v60-conditional-mean}

Split the promoted positive Higgs field into an integrated color
\(\phi_I\) and a retained skeleton \(\eta\).  Write
\(\nu_I^\eta\) for the conditional law.  Assume its Hessian has the
uniform lower bound \(\kappa_{\rm p}>0\), its one-form inverse satisfies
\begin{equation}
 \|P_A(\mathcal L_{1,I}^{\eta})^{-1}P_C\|
 \le C_{\rm W}e^{-m_{\rm W}d(A,C)},
\label{eq:v4v60-conditional-Witten}
\end{equation}
and changing one retained cell \(y\) creates a score
\(K_y^\eta=D_{\eta_y}S_I^\eta\) with
\begin{equation}
 \sup_y\sum_{u\in I}e^{m_\partial d(u,y)}
 \|P_u\nabla_{\phi_I}K_y^\eta\|_{L^2(\nu_I^\eta)}
 \le J_{\partial}
\label{eq:v4v60-boundary-score}
\end{equation}
for some \(m_\partial>0\).  These are the conditional V38 bounds for the
finite-range Higgs stencil, with the exponentially summable V34 modulus
Hessian extension where it is present.

For an integrated Yukawa cell \(x\), put
\[
 m_x(\eta)=\mathbb E_{\nu_I^\eta}\phi_{A_x}.
\]

\begin{lemma}[Exponential boundary response of the conditional mean]
\label{lem:v4v60-mean-response}
There are \(C_m,m_m>0\), independent of volume and \(\eta\), such that
\begin{equation}
 |m_x(\eta)-m_x(0)|
 \le
 C_m\sum_{y\in\Sigma}
 e^{-m_md(x,y)}|\eta_y|.
\label{eq:v4v60-mean-response}
\end{equation}
\end{lemma}

\begin{proof}
Along the line \(t\eta\), normalized differentiation gives
\[
 D_{\eta_y}m_x(t\eta)
 =
 -\operatorname{Cov}_{\nu_I^{t\eta}}
 \left(\phi_{A_x},K_y^{t\eta}\right).
\]
The field gradient of \(\phi_{A_x}\) is a fixed local projection.  Apply
the conditional Helffer--Sjostrand identity,
\eqref{eq:v4v60-conditional-Witten}, and
\eqref{eq:v4v60-boundary-score}, then convolve the two exponential kernels.
After reducing the smaller exponent, this gives
\[
 |D_{\eta_y}m_x(t\eta)|
 \le C_me^{-m_md(x,y)}
\]
uniformly in \(t,\eta\).  Integrate from \(t=0\) to \(t=1\) and sum the
retained coordinates.
\end{proof}

Assume also the centered chart has
\begin{equation}
 \sup_x|m_x(0)|\le M_0.
\label{eq:v4v60-zero-boundary-mean}
\end{equation}
Since \(R_Y\to\infty\), this fixed number is at most \(R_Y/8\) for small
\(\epsilon_Y\).

\subsection{A smooth first-owner boundary partition}
\label{sec:v4v60-owner-partition}

Fix \(0<\beta<m_m\), and let
\[
 \mathcal S_\beta
 =
 \sup_x\sum_{y\in\Sigma}
 e^{-(m_m-\beta)d(x,y)}<\infty.
\]
Choose \(\theta>0\) so
\begin{equation}
 2C_m\theta\mathcal S_\beta\le\frac18.
\label{eq:v4v60-theta}
\end{equation}
For the tail cell \(x\), define distance-dependent guard thresholds
\begin{equation}
 L_{xy}
 =
 \theta R_Ye^{\beta d(x,y)}.
\label{eq:v4v60-guard-threshold}
\end{equation}
Let the smooth radius of the retained cell be
\(r_y^\Sigma=(1+|\eta_y|^2)^{1/2}\), and use the V34 cutoff \(\chi\) to set
\[
 g_{xy}(\eta)=\chi(r_y^\Sigma/L_{xy}).
\]
Order the finite retained cells deterministically and define
\begin{align}
 G_x(\eta)
 &:=
 \prod_{y\in\Sigma}g_{xy}(\eta),
\label{eq:v4v60-good-guard}\\
 B_{xy}(\eta)
 &:=
 (1-g_{xy}(\eta))
 \prod_{z<y}g_{xz}(\eta).
\label{eq:v4v60-first-owner}
\end{align}

\begin{lemma}[Exact owner partition and its support]
\label{lem:v4v60-owner-partition}
Pointwise in every finite volume,
\begin{equation}
 1=G_x+\sum_{y\in\Sigma}B_{xy}.
\label{eq:v4v60-owner-identity}
\end{equation}
On the support of \(G_x\),
\begin{equation}
 |m_x(\eta)|\le\frac12R_Y
\label{eq:v4v60-good-mean}
\end{equation}
for sufficiently small \(\epsilon_Y\).  On the support of \(B_{xy}\),
\begin{equation}
 r_y^\Sigma\ge L_{xy}
 =\theta R_Ye^{\beta d(x,y)}.
\label{eq:v4v60-owner-support}
\end{equation}
\end{lemma}

\begin{proof}
The finite product telescope
\[
 1-\prod_yg_{xy}
 =
 \sum_y(1-g_{xy})\prod_{z<y}g_{xz}
\]
proves \eqref{eq:v4v60-owner-identity}.  If \(G_x\ne0\), every
\(r_y^\Sigma<2L_{xy}\).  Lemma~\ref{lem:v4v60-mean-response},
\eqref{eq:v4v60-theta}, and
\eqref{eq:v4v60-zero-boundary-mean} give
\[
 |m_x(\eta)|
 \le
 \frac18R_Y+
 2C_m\theta R_Y
 \sum_y e^{-(m_m-\beta)d(x,y)}
 \le\frac38R_Y,
\]
which implies the stated weaker bound.  If \(B_{xy}\ne0\), then
\(1-g_{xy}\ne0\), and the defining property of \(\chi\) gives
\(r_y^\Sigma\ge L_{xy}\).
\end{proof}

\begin{lemma}[The owner weights are super-summable]
\label{lem:v4v60-owner-small}
Let \(A_y(\eta)\) be any fixed-degree cylinder polynomial times
\(e^{c\epsilon_Y(1+r_y^\Sigma)}\), or the same kind of envelope with
exponentially summable boundary coefficients, including one field derivative of
\(B_{xy}\).  For every fixed \(p<\infty\), there are \(c_p,C_p>0\) such
that
\begin{equation}
 \|B_{xy}A_y\|_{L^p(\nu_\Sigma)}
 +
 \|\nabla_{\eta}B_{xy}\,A_y\|_{L^p(\nu_\Sigma)}
 \le
 C_p
 \exp\!\left[
 -c_pR_Y^2e^{2\beta d(x,y)}
 \right].
\label{eq:v4v60-owner-small}
\end{equation}
Consequently
\begin{equation}
 \sup_x\sum_y
 \exp\!\left[
 -c_pR_Y^2e^{2\beta d(x,y)}
 \right]
 \le
 C_p'e^{-c_p'R_Y^2}.
\label{eq:v4v60-owner-row}
\end{equation}
\end{lemma}

\begin{proof}
The support statement \eqref{eq:v4v60-owner-support} supplies an actual
retained-field event at threshold \(L_{xy}\).  Apply the correlated
subgaussian estimate \eqref{eq:v4v31-quadratic-exponential}, absorbing the
fixed polynomial and linear exponential as in
Lemma~\ref{lem:v4v34-simultaneous-tail}.  A derivative of \(g_{xy}\)
adds \(L_{xy}^{-1}\) and remains in the annulus
\(L_{xy}\le r_y^\Sigma\le2L_{xy}\); derivatives of earlier good factors
are treated by assigning their own first differentiated index, which only
improves the corresponding threshold.  Holder's inequality handles the
finite local polynomial factors.

The number of cells at distance \(n\) grows at most polynomially, whereas
\(\exp[-cR_Y^2e^{2\beta n}]\) decays super-exponentially in \(n\).
Summing shells proves \eqref{eq:v4v60-owner-row}.
\end{proof}

\subsection{Conditional rarity on the guarded sector}
\label{sec:v4v60-guarded-rarity}

\begin{lemma}[Uniform conditional tail payment after guarding]
\label{lem:v4v60-guarded-tail}
Let \(T_x\) be a V59 tail-owned carrier supported in \(A_x\).  On the
support of \(G_x\), for \(0\le j\le2\),
\begin{equation}
 \|D_B^jT_x\|_{L^p(\nu_I^\eta)}
 +
 \|\nabla_{\phi_I}D_B^jT_x\|_{L^p(\nu_I^\eta)}
 \le
 C_{j,p}e^{-c_pR_Y^2},
\label{eq:v4v60-guarded-tail}
\end{equation}
uniformly in the retained field.
\end{lemma}

\begin{proof}
The conditional action remains \(\kappa_{\rm p}\)-strongly convex.
Therefore its local field is subgaussian around the conditional mean, by
the same Bakry--Emery argument as V31.  Lemma
\ref{lem:v4v60-owner-partition} places that mean within \(R_Y/2\) of the
origin on every cell of the fixed support after reducing constants for the
bounded cell radius.  The event \(r_x\ge R_Y\) then forces a centered
fluctuation of size \(cR_Y\).  Raise the V59 envelope to the \(p\)-th
power, absorb its polynomial and linear exponential into a small quadratic
exponential, and integrate the conditional Gaussian tail.  This proves
\eqref{eq:v4v60-guarded-tail}.
\end{proof}

\subsection{The transferred Witten carrier}
\label{sec:v4v60-transferred-carrier}

Let \(F_z\) be an ordinary local good carrier in the integrated color and
put
\[
 \mathcal C^\eta(T_x,F_z)
 =
 \operatorname{Cov}_{\nu_I^\eta}(T_x,F_z).
\]
Insert \eqref{eq:v4v60-owner-identity} without taking an absolute value:
\begin{equation}
 \mathcal C^\eta(T_x,F_z)
 =
 G_x\mathcal C^\eta(T_x,F_z)
 +
 \sum_yB_{xy}\mathcal C^\eta(T_x,F_z).
\label{eq:v4v60-carrier-split}
\end{equation}

\begin{theorem}[One-step boundary-tail transfer]
\label{thm:v4v60-boundary-transfer}
Assume the local coefficient and score envelopes of V34, V50, and V59.
There are \(c,m,C>0\) such that the guarded part obeys
\begin{equation}
 \sup_\eta
 \left|
 G_x\mathcal C^\eta(T_x,F_z)
 \right|
 \le
 Ce^{-cR_Y^2}e^{-md(x,z)},
\label{eq:v4v60-guarded-carrier}
\end{equation}
while each owner part obeys, in every fixed skeleton \(L^p\) norm,
\begin{equation}
 \left\|
 B_{xy}\mathcal C^\eta(T_x,F_z)
 \right\|_{L^p(\nu_\Sigma)}
 \le
 C_p
 e^{-md(x,z)}
 \exp\!\left[
 -c_pR_Y^2e^{2\beta d(x,y)}
 \right].
\label{eq:v4v60-owner-carrier}
\end{equation}
The same estimates hold for the exact placements generated by at most two
background marks and by one retained-field derivative.  In particular,
after summing the owner,
\begin{equation}
 \sup_x\sum_y
 \left\|
 B_{xy}\mathcal C^\eta(T_x,F_z)
 \right\|_{L^p(\nu_\Sigma)}
 \le
 C_p'e^{-c_p'R_Y^2}e^{-md(x,z)}.
\label{eq:v4v60-summed-owner-carrier}
\end{equation}
\end{theorem}

\begin{proof}
For \eqref{eq:v4v60-guarded-carrier}, apply the conditional
Helffer--Sjostrand identity and
\eqref{eq:v4v60-conditional-Witten}.  Lemma
\ref{lem:v4v60-guarded-tail} pays the tail gradient uniformly; the good
gradient is uniformly bounded.

On an owner sector, apply the same Witten identity before integrating the
skeleton.  Its spatial inverse still supplies \(e^{-md(x,z)}\).
Conditional moments of the two gradients grow at most by a fixed
polynomial and linear exponential in the exponentially weighted nearby
boundary field.  Keep \(B_{xy}\) attached, use Holder's inequality under
the skeleton law, and apply Lemma~\ref{lem:v4v60-owner-small}.  This proves
\eqref{eq:v4v60-owner-carrier}; its row sum is
\eqref{eq:v4v60-owner-row}.

A background derivative produces the exact V50--V55 alternatives: an
explicit derivative, a first score, a second score, or two first scores.
A retained-field derivative either hits a guard or differentiates the
conditional covariance.  In the latter case the normalized derivative is
the complete third cumulant of
\eqref{eq:v4v50-carrier-derivative}, plus explicit derivatives.  A second
background mark raises the fixed order at most to four, as displayed in
\eqref{eq:v4v59-second-mark-identity}.  On the guarded sector every
complete product containing \(T_x\) has the conditional rare-gradient
bound.  On an owner sector the factor \(B_{xy}\), including any derivative
of it, remains attached and supplies
\eqref{eq:v4v60-owner-small}.  The V48 cut identity provides the actual
spatial covariance in every fixed-order term.  Finite mark placements
change only the constants.
\end{proof}

\begin{corollary}[A two-edge typed tail packet]
\label{cor:v4v60-two-edge-packet}
Every term in \eqref{eq:v4v60-carrier-split} has one of two exact
ancestries:
\begin{enumerate}
\item the physical tail cell \(x\) pays
      \(e^{-cR_Y^2}\), and the Witten inverse joins \(x\) to \(z\); or
\item the first retained owner \(y\) pays
      \(e^{-cR_Y^2e^{2\beta d(x,y)}}\), and the same unsplit conditional
      covariance joins \(x\) to \(z\).
\end{enumerate}
Both kernels have finite rooted row sums.  No product is inferred from two
separately compressed estimates.
\end{corollary}

\begin{proof}
This is precisely
\eqref{eq:v4v60-guarded-carrier}--%
\eqref{eq:v4v60-summed-owner-carrier}.  The owner kernel is stronger than
an exponential edge and hence is row summable.
\end{proof}

\begin{firewall}[The guard is transferred ownership, not a duplicate tail]
\label{fw:v4v60-no-double-tail}
The factors \(B_{xy}\) form a partition of the complement of \(G_x\).
Exactly one first retained cell owns each configuration in that complement.
If \(y\) also supports a physical \(t_yV_y\) insertion, both ancestries
must be merged and its field event paid once.  The guard partition is
introduced only inside the estimate of the existing \(x\)-tail carrier; it
does not add a second determinant vertex or a second V34 tail activity.
\end{firewall}

\begin{assessment}[What changes in the collision compiler]
\label{ass:v4v60-compiler}
The impossible norm
\(\sup_\eta\|\mathcal C^\eta\|\) is replaced by a direct sum of a guarded
\(L^\infty\) branch and first-owner skeleton \(L^p\) branches.  V60 proves
that this mixed norm is stable for one physical conditional collision and
its two marks.  The V58 abstract sum was stated for a scalar typed norm.
To close all collision histories, the V49--V58 compiler must be lifted to
this mixed guarded/owner norm and must show that later owner collisions
preserve the unique first-owner partition.  That all-depth lifting is not
proved here.
\end{assessment}

\subsection{V60 result ledger and V61 target}
\label{sec:v4v60-conclusion}

V60 records the exact YM V75 and NS V31 fingerprints and imports their
common proof-order lesson: the two required gains must remain in one
carrier, and a family of marks needs a unique owner before summation.  It
proves exponential response of the conditional Higgs mean, constructs an
exact smooth first-owner boundary partition with distance-growing
thresholds, and shows that its owner weights are super-summable.  The
guarded conditional sector recovers the physical V34 tail payment; the
complement transfers that payment to one retained cell without detaching
the Witten edge.  This gives the one-step typed tail packet through two
background marks and one retained derivative.

V61 should lift the rooted star-merger recursion from a scalar norm to the
guarded/owner direct-sum norm.  Its central algebraic requirement is
idempotent ownership: multiplying two first-owner partitions at a later
collision must choose one least common owner and must not charge both
partitions as independent rare events.  If that algebra does not have a
submultiplicative row norm, the programme must replace color conditioning
by a full-law forest interpolation, where the V59 rare gradient remains
inside every interpolated Witten edge.

\clearpage
\section{V61: an inserted log-Sobolev generator closes the one-tail
multiway sector}
\label{sec:v4v61}

V60 transfers a tail payment across one conditional boundary, but its
fixed-order proof does not yet cover a collision with arbitrarily many good
children.  The V53 generator does cover arbitrary valence, although it
requires every child to have a uniform Lipschitz seminorm.  A tail carrier
does not satisfy that hypothesis.

The correct modification is to introduce the tail carrier only as one
linear insertion in the good-field exponential generator.  The denominator
is the zero-free V53 generator of the Lipschitz children.  Cauchy--Schwarz
bounds the inserted numerator by the rare \(L^2\) norm of the tail times a
uniform exponential moment of the good load.  Multivariate Cauchy
extraction then gives a factorial cumulant bound containing the tail factor
once and every good-child Lipschitz scale separately.

This closes arbitrary collision valence in the sector with one tail
descendant, including the V60 guarded/owner alternatives and two background
marks.  Histories containing several distinct tail descendants require an
additional consolidation theorem and are left explicit.

\subsection{A generator with one \(L^2\) insertion}
\label{sec:v4v61-inserted-generator}

Let \(\nu\) satisfy the log-Sobolev inequality
\eqref{eq:v4v53-LSI}.  Let \(F_1,\ldots,F_m\) be centered complex
observables with Lipschitz seminorms \(L_i\) as in
\eqref{eq:v4v53-Li}.  Put
\[
 W(t)=\sum_{i=1}^m t_iF_i,
 \qquad
 M(t)=\mathbb E_\nu e^{W(t)},
 \qquad
 \Lambda(t)=\sum_i|t_i|L_i.
\]
For \(T\in L^2(\nu)\), define
\begin{equation}
 H_T(t)
 =
 \frac{\mathbb E_\nu[T e^{W(t)}]}{M(t)}.
\label{eq:v4v61-inserted-generator}
\end{equation}

\begin{lemma}[Uniform exponential square moment]
\label{lem:v4v61-square-moment}
There are \(s_1,C_1>0\), depending only on the log-Sobolev constant,
such that
\begin{equation}
 \Lambda(t)\le s_1
 \quad\Longrightarrow\quad
 |M(t)|\ge\frac12,
 \qquad
 \mathbb E_\nu|e^{W(t)}|^2\le C_1.
\label{eq:v4v61-square-moment}
\end{equation}
\end{lemma}

\begin{proof}
Take \(s_1\le s_0/4\), where \(s_0\) is the V53 zero-free radius.
Theorem~\ref{thm:v4v53-zero-free-generator} gives
\(|M(t)-1|\le1/2\).  Since every \(F_i\) is centered,
\(\operatorname{Re}W(t)\) is centered, and its real Lipschitz seminorm is
at most \(2\Lambda(t)\).  The Herbst exponential-moment bound following
from \eqref{eq:v4v53-LSI} gives
\[
 \mathbb E_\nu e^{2\operatorname{Re}W(t)}
 \le
 \exp(C\Lambda(t)^2).
\]
This is uniformly bounded on \(\Lambda(t)\le s_1\), and it equals
\(\mathbb E|e^W|^2\).
\end{proof}

\begin{theorem}[One-insertion factorial cumulant bound]
\label{thm:v4v61-inserted-cumulant}
For \(m\ge1\),
\begin{equation}
 \boxed{
 |\kappa_\nu(T,F_1,\ldots,F_m)|
 \le
 C_2m!\left(\frac{e}{s_1}\right)^m
 \|T\|_{L^2(\nu)}
 \prod_{i=1}^mL_i.}
\label{eq:v4v61-inserted-cumulant}
\end{equation}
The constant \(C_2\) depends only on the log-Sobolev constant.
\end{theorem}

\begin{proof}
By Lemma~\ref{lem:v4v61-square-moment} and Cauchy--Schwarz,
\begin{equation}
 |H_T(t)|
 \le
 2C_1^{1/2}\|T\|_2
 \qquad(\Lambda(t)\le s_1).
\label{eq:v4v61-H-bound}
\end{equation}
The numerator and denominator are analytic on this load domain, so \(H_T\)
is analytic there.  The quotient rule, or equivalently differentiating
once at \(s=0\) in
\[
 \log\mathbb E_\nu
 \exp\!\left(sT+\sum_it_iF_i\right),
\]
as a formal derivative in \(s\), gives
\begin{equation}
 \partial_{t_1}\cdots\partial_{t_m}H_T(0)
 =
 \kappa_\nu(T,F_1,\ldots,F_m).
\label{eq:v4v61-H-cumulant}
\end{equation}
No exponential moment of \(T\) is required for this one derivative.

When every \(L_i>0\), choose Cauchy radii
\[
 r_i=\frac{s_1}{mL_i}.
\]
Their distinguished polytorus lies in \(\Lambda(t)\le s_1\).
Equations \eqref{eq:v4v61-H-bound}--%
\eqref{eq:v4v61-H-cumulant} give
\[
 |\kappa_\nu(T,F_1,\ldots,F_m)|
 \le
 2C_1^{1/2}\|T\|_2
 \left(\frac m{s_1}\right)^m
 \prod_iL_i.
\]
Use \(m^m\le e^mm!\).  If some \(L_i=0\), its centered observable vanishes
on the connected field space and so does the cumulant.
\end{proof}

\begin{remark}[The tail needs no exponential generating variable]
\label{rem:v4v61-no-tail-mgf}
A V34 tail coefficient may grow like a polynomial times
\(e^{c\epsilon_Yr}\); an exponential of that coefficient need not be
integrable.  The proof never forms \(e^{sT}\) for nonzero \(s\).
Equation \eqref{eq:v4v61-H-cumulant} uses only the derivative at \(s=0\),
which is justified by \(T\in L^2\) and the square exponential moment of the
good load.
\end{remark}

\subsection{A star at a one-tail collision}
\label{sec:v4v61-one-tail-star}

Let one tail child \(T_0\) and \(k-1\) good children
\(F_1,\ldots,F_{k-1}\) approach a conditional component centered at \(x\).
Let \(S_i\) be their reference supports.  Assume the pre-collision channels
give
\begin{align}
 \|T_0\|_{L^2}
 &\le
 \tau\lambda_0e^{-m_cd(S_0,x)},
\label{eq:v4v61-tail-response}\\
 L_i
 &\le
 \lambda_i e^{-m_cd(S_i,x)},
 \qquad 1\le i\le k-1.
\label{eq:v4v61-good-responses}
\end{align}
In the guarded V60 branch,
\(\tau=Ce^{-cR_Y^2}\).  In an owner branch it is the corresponding
owner weight before its skeleton norm is taken.

\begin{theorem}[One-tail multiway star]
\label{thm:v4v61-one-tail-star}
Under \eqref{eq:v4v61-tail-response}--%
\eqref{eq:v4v61-good-responses},
\begin{align}
 &|\kappa(T_0,F_1,\ldots,F_{k-1})|\nonumber\\
 &\quad\le
 k!C_*^k\tau
 \left(\prod_{i=0}^{k-1}\lambda_i\right)
 \sum_{r=0}^{k-1}
 \prod_{\substack{0\le i\le k-1\\i\ne r}}
 e^{-\frac12m_cd(S_i,S_r)}.
\label{eq:v4v61-one-tail-star}
\end{align}
Thus the rare factor and all \(k-1\) spatial edges occur in one complete
collision estimate.
\end{theorem}

\begin{proof}
Apply Theorem~\ref{thm:v4v61-inserted-cumulant} with \(m=k-1\), then insert
\eqref{eq:v4v61-tail-response}--%
\eqref{eq:v4v61-good-responses}.  This produces \(\tau\) times the product
of all \(k\) pre-collision exponential factors.  Choose \(r\) for which
\(d(S_r,x)\) is minimal and drop its factor, which is at most one.  For
\(i\ne r\),
\[
 d(S_i,S_r)
 \le d(S_i,x)+d(S_r,x)
 \le2d(S_i,x).
\]
The remaining product therefore dominates the \(r\)-rooted star on the
right of \eqref{eq:v4v61-one-tail-star}.  Enlarge a fixed exponential
constant to replace \((k-1)!\) by \(k!\), and then sum over all possible
roots.
\end{proof}

\begin{firewall}[The tail \(L^2\) response is an actual child scale]
\label{fw:v4v61-tail-response}
The factor \(e^{-m_cd(S_0,x)}\) in
\eqref{eq:v4v61-tail-response} must come from the actual pre-collision
propagation of the tail child, as in the V59--V60 Witten carrier.  It may
not be inserted merely because the tail event is rare.  Rarity supplies
\(\tau\); propagation supplies the child response distance.
\end{firewall}

\subsection{The guarded/owner norm at arbitrary valence}
\label{sec:v4v61-owner-lift}

For a one-tail carrier based at \(x\), define
\begin{equation}
 \|T\|_{\mathfrak O_x}
 :=
 \|G_xT\|_{L^\infty(d\eta;L^2(d\nu_I^\eta))}
 +
 \sum_y
 \|B_{xy}T\|_{L^p(d\nu_\Sigma;L^2(d\nu_I^\eta))}.
\label{eq:v4v61-owner-norm}
\end{equation}
The owner sum may include the stronger distance weight in
\eqref{eq:v4v60-owner-carrier}.  V60 proves
\begin{equation}
 \|T\|_{\mathfrak O_x}
 \le C e^{-cR_Y^2}
\label{eq:v4v61-owner-tail-size}
\end{equation}
through the relevant fixed marks.

\begin{theorem}[One-tail owner module is stable at a multiway merger]
\label{thm:v4v61-owner-module}
Suppose one child of a \(k\)-way collision lies in
\(\mathfrak O_x\), the other \(k-1\) children satisfy the V53 Lipschitz
response bounds, and the conditional log-Sobolev and response constants are
uniform.  Then the collision carrier lies in the corresponding output
owner space and satisfies
\begin{equation}
 \|\mathcal M_k(T_0,F_1,\ldots,F_{k-1})\|_{\mathfrak O}
 \le
 k!C_*^k
 \|T_0\|_{\mathfrak O_x}
 \left(\prod_{i=1}^{k-1}\lambda_i\right)
 \sum_r\prod_{i\ne r}w(S_i,S_r).
\label{eq:v4v61-owner-module}
\end{equation}
The same assertion holds through two background marks.
\end{theorem}

\begin{proof}
On the guarded component of
\eqref{eq:v4v61-owner-norm}, use
Lemma~\ref{lem:v4v60-guarded-tail} and
Theorem~\ref{thm:v4v61-one-tail-star} pointwise in \(\eta\).  On the
\(y\)-owner component, apply the inserted-generator estimate pointwise with
the conditional tail norm \(A_y(\eta)\).  Its right side is linear in
\(A_y\).  Keep \(B_{xy}\) attached, take the skeleton \(L^p\) norm, and use
\eqref{eq:v4v60-owner-small}.  Summing \(y\) uses
\eqref{eq:v4v60-owner-row}.  The same owner label remains attached to the
output; no new partition is introduced at the merger.

A background derivative has the exhaustive V54--V56 placements.  If it
hits the tail child, use its marked \(\mathfrak O_x\) norm.  If it creates
a first or second score, include that score among the Lipschitz loads in
Theorem~\ref{thm:v4v61-inserted-cumulant}.  Two first scores are two actual
loads.  Their polynomial choice factors in \(k\) change the coefficient but
not the convergence radius.
\end{proof}

\subsection{All-depth summation with one primitive tail}
\label{sec:v4v61-one-tail-histories}

Let \(R\) be the unmarked good-field rooted sum of V57.  A merger with one
distinguished tail-descended child has one choice of that child in addition
to the \(k\) star-center choices.  Define
\begin{equation}
 \Psi(R)
 :=
 A_{\mathfrak O}
 \sum_{k\ge2}
 k^2C_{\mathfrak O}^kW^{k-1}R^{k-1}.
\label{eq:v4v61-Psi}
\end{equation}
It is analytic for \(C_{\mathfrak O}WR<1\) and tends to zero with \(R\).

\begin{theorem}[Rooted histories with exactly one primitive tail]
\label{thm:v4v61-one-tail-histories}
If \(R\) is sufficiently small that \(\Psi(R)<1\), the rooted sum
\(M_1\) of collision histories containing exactly one primitive tail cell
obeys
\begin{equation}
 M_1
 \le
 \frac{C e^{-cR_Y^2}}{1-\Psi(R)}.
\label{eq:v4v61-one-tail-histories}
\end{equation}
The once- and twice-marked sums obey the same estimate with changed
constants.  The bound is uniform in volume.
\end{theorem}

\begin{proof}
At the first merger above the unique tail primitive, exactly one child
contains its ancestry.  Apply Theorem~\ref{thm:v4v61-owner-module} to that
child and the V57 norm to all other children.  The unordered-child \(1/k!\)
cancels the local \(k!\); \(k\) choices of star center and at most \(k\)
choices of the tail child give the series
\eqref{eq:v4v61-Psi}.  Therefore
\[
 M_1
 \le
 Ce^{-cR_Y^2}+\Psi(R)M_1.
\]
This proves \eqref{eq:v4v61-one-tail-histories}.  Marking differentiates
the same analytic branching series a fixed number of times, as in
Lemma~\ref{lem:v4v57-marked-series}; it does not change its radius.
\end{proof}

\begin{corollary}[The linear tail coefficient of the pressure]
\label{cor:v4v61-linear-tail-pressure}
Introduce a scalar \(s\) multiplying every primitive tail activity and let
\(p(s)\) be the connected pressure branch whenever its coefficients are
defined.  Under the hypotheses above,
\begin{equation}
 \left|D_B^j p'(0)\right|
 \le
 C_j e^{-cR_Y^2},
 \qquad j=0,1,2.
\label{eq:v4v61-linear-tail-pressure}
\end{equation}
The same bound holds for the first tail coefficient of a normalized local
ratio.
\end{corollary}

\begin{proof}
Differentiation at \(s=0\) selects exactly one primitive tail cell.
Root its collision history and apply
Theorem~\ref{thm:v4v61-one-tail-histories}.  Translation covariance gives
the intensive coefficient; rooting at the local insertion gives the local
ratio.
\end{proof}

\begin{assessment}[The unresolved multi-tail sector]
\label{ass:v4v61-multi-tail-open}
V61 closes every valence and every later merger for histories containing
one primitive tail cell.  It does not yet prove the full difference
\(p(1)-p(0)\).  A merger may contain several tail descendants.  If they own
distinct failed cells, the V34 simultaneous-tail theorem should supply a
product of activities.  If two descendants refer to the same failed cell,
the event is idempotent and can be paid only once.  A valid compiler must
first consolidate equal physical cell labels and then apply the
simultaneous-tail estimate to the distinct owner set.
\end{assessment}

\begin{firewall}[One inserted generator is not a many-tail moment
generator]
\label{fw:v4v61-one-insertion}
Theorem~\ref{thm:v4v61-inserted-cumulant} differentiates once in the
unbounded tail direction.  It does not assert analyticity of
\(\mathbb E e^{\sum_qs_qT_q}\), which may fail for V34 tail coefficients.
Several tail insertions must be handled by the exact finite-cell pattern
partition and simultaneous tail moments, not by inventing a tail
exponential generator.
\end{firewall}

\subsection{V61 result ledger and V62 target}
\label{sec:v4v61-conclusion}

V61 proves an all-order cumulant estimate with one rare \(L^2\) insertion
and arbitrarily many centered Lipschitz loads.  The rare norm appears once,
each good load contributes its own response scale, and the physical
collision therefore carries a complete star.  The V60 guarded/owner
direct-sum norm is stable under this merger.  Consequently every rooted
history with exactly one primitive tail cell is summable through two
background marks, and the coefficient \(p'(0)\) is
\(O(e^{-cR_Y^2})\).

V62 should construct the distinct-tail-set algebra.  It must consolidate
all occurrences of the same physical failed cell before applying
\eqref{eq:v4v34-simultaneous-tail}, prove a projective norm proportional to
\(\tau^{|S|}\) for a distinct owner set \(S\), and show that star mergers
act submultiplicatively on disjoint sets while overlap is resolved by one
local consolidation coefficient.  Only then can the linear result be
summed to the full tail correction asserted conditionally in V58.

\clearpage
\section{V62: distinct tail-cell sets and a tilted multi-tail cumulant bound}
\label{sec:v4v62}

V61 controls an arbitrary-valence collision with one tail descendant by
using the tail only as a linear insertion in the V53 good-field generator.
For several tail descendants, one must first decide whether repeated
occurrences refer to distinct failed cells.  Paying the same event twice
would be false.

This note resolves that algebra before estimating the cumulant.  The local
Grassmann exponential is finite, so all occurrences at one physical cell
are consolidated into one elementary cell packet.  The V49 collision
history partitions this elementary label set; different children at every
merger therefore carry disjoint sets of tail-cell labels.  The V34
simultaneous-tail theorem then gives one activity for every element of the
union.

For the local cumulant estimate, tilt only by the bounded-Lipschitz good
children.  Under that zero-free tilted law, express the cumulant of the
tail variables by the ordinary finite moment--cumulant formula.  Joint
\(L^2\) tail products replace a nonexistent tail moment-generating
function.  A final Cauchy extraction in the good variables yields a
factorial-exponential coefficient with every distinct tail activity and
every good Lipschitz scale.

\subsection{One elementary packet per physical tail cell}
\label{sec:v4v62-local-consolidation}

At a Yukawa cell \(x\), let
\[
 \mathcal V_x^{\rm T}
 =
 \overline\psi\,(t_xV_x)\psi .
\]
Because the local Grassmann fibre has fixed finite dimension \(d_{\rm F}\),
define the complete nonconstant local packet
\begin{equation}
 \Theta_x
 :=
 \exp(-\mathcal V_x^{\rm T})-1
 =
 \sum_{j=1}^{d_{\rm F}}
 \frac{(-\mathcal V_x^{\rm T})^j}{j!}.
\label{eq:v4v62-local-tail-packet}
\end{equation}
All repeated powers at \(x\) are internal terms of \(\Theta_x\), rather
than distinct tail labels.

\begin{lemma}[Exact cell consolidation]
\label{lem:v4v62-cell-consolidation}
At every finite cutoff,
\begin{equation}
 \prod_x\exp(-\mathcal V_x^{\rm T})
 =
 \prod_x(1+\Theta_x).
\label{eq:v4v62-cell-product}
\end{equation}
After expansion, a selected family is indexed by a set of distinct
physical cells.  Background differentiation through order two changes the
coefficient of \(\Theta_x\) but never creates a second label at \(x\).
\end{lemma}

\begin{proof}
Equation \eqref{eq:v4v62-local-tail-packet} is the exact finite Taylor
series in the local Grassmann algebra, so
\(1+\Theta_x=\exp(-\mathcal V_x^{\rm T})\).  Multiplication over cells proves
\eqref{eq:v4v62-cell-product}.  The product rule places a first or second
mark on an actual local coefficient in the same finite polynomial.  It
does not alter the cell index of that factor.
\end{proof}

Let \(S(h)\) be the set of selected tail-cell labels below a collision
history \(h\).

\begin{lemma}[Child tail sets are disjoint]
\label{lem:v4v62-disjoint-children}
At every internal vertex with child histories \(h_1,\ldots,h_k\),
\begin{equation}
 S(h_i)\cap S(h_j)=\varnothing\quad(i\ne j),
 \qquad
 S(h)=\mathbin{\dot\bigcup}_{i=1}^kS(h_i).
\label{eq:v4v62-disjoint-union}
\end{equation}
\end{lemma}

\begin{proof}
The V49 collision-history expansion is obtained by successive partitions
of the original elementary label set.  Sibling child blocks in a partition
are disjoint.  Intersecting each child block with the tail-cell subset
preserves disjointness, and their union is the parent's tail subset.
Lemma~\ref{lem:v4v62-cell-consolidation} ensures that the original label
set contains \(x\) at most once.
\end{proof}

\begin{firewall}[Consolidate before assigning a tail activity]
\label{fw:v4v62-consolidate-first}
The equality
\(\boldsymbol1_{\{r_x\ge R_Y\}}^j
=\boldsymbol1_{\{r_x\ge R_Y\}}\)
does not provide \(j\) rare factors.  All same-cell Grassmann powers must
first be combined in \(\Theta_x\); only then may
\eqref{eq:v4v34-simultaneous-tail} be applied to the resulting distinct
cell set.
\end{firewall}

\subsection{Joint product norms from simultaneous tail payment}
\label{sec:v4v62-product-norms}

Let \(T_1,\ldots,T_\ell\) be tail carriers whose physical owner sets
\(S_1,\ldots,S_\ell\) are pairwise disjoint.  Write
\[
 n_B=\left|\mathbin{\dot\bigcup}_{q\in B}S_q\right|
 \qquad(B\subseteq\{1,\ldots,\ell\}).
\]

\begin{lemma}[Distinct-tail product bound]
\label{lem:v4v62-tail-product-bound}
Under the V31--V34 finite-fibre and bounded-overlap hypotheses, there are
\(c_{\rm T},K_{\rm T}>0\) such that, for every nonempty \(B\),
\begin{equation}
 \left\|\prod_{q\in B}T_q\right\|_{L^2(\nu_{\rm G,B})}
 \le
 K_{\rm T}^{n_B}
 e^{-c_{\rm T}R_Y^2n_B}
 \prod_{q\in B}a_q.
\label{eq:v4v62-tail-product-bound}
\end{equation}
Here \(a_q\) contains the ordinary fermion-tree and pre-collision response
norm of carrier \(q\), with no tail probability.  One or two actual
background marks multiply the right side by at most a polynomial in
\(n_B\), which is absorbed by replacing \(K_{\rm T}\) by a larger constant.
\end{lemma}

\begin{proof}
Expand each carrier in its finite Grassmann coefficient norm and keep the
local packets \(\Theta_x\) unintegrated until the cell set is fixed.  The
square of a coefficient has a fixed polynomial and linear-exponential
envelope at every selected cell and vanishes unless
\(r_x\ge R_Y\).  Apply
\eqref{eq:v4v34-marked-tail} simultaneously to the distinct union of
\(n_B\) cells, with doubled polynomial and exponential constants, and then
take the square root.  Reduce \(c_{\rm T}\) after the square root.

The V33 Gram/tree estimate and all already acquired response kernels are
the factors \(a_q\).  A mark has one actual local placement; two marks have
one second placement or two first placements.  There are at most
\(C(1+n_B)^2\) choices.  Since
\((1+n)^2\le C_0^n\) for \(n\ge1\), this is absorbed into
\(K_{\rm T}^{n_B}\).
\end{proof}

For later use, write
\begin{equation}
 \tau_q
 :=
 K_{\rm T}^{|S_q|}
 e^{-c_{\rm T}R_Y^2|S_q|}a_q.
\label{eq:v4v62-tauq}
\end{equation}
Then \eqref{eq:v4v62-tail-product-bound} reads
\[
 \left\|\prod_{q\in B}T_q\right\|_2
 \le\prod_{q\in B}\tau_q.
\]

\subsection{Tail cumulants under a good-field tilt}
\label{sec:v4v62-tilted-tail}

Let \(F_1,\ldots,F_m\) be centered good observables with Lipschitz loads
\(L_i\).  On the V61 polydisc \(\Lambda(t)\le s_1\), define the normalized
complex tilt
\begin{equation}
 \mathbb E_tA
 :=
 \frac{\mathbb E_\nu[Ae^{W(t)}]}{M(t)}.
\label{eq:v4v62-tilted-expectation}
\end{equation}
It is a normalized linear functional; positivity is not needed below.
Define its tail cumulant by the finite moment--cumulant formula
\begin{equation}
 K_{\rm T}(t)
 :=
 \sum_{\pi\in\Pi([\ell])}
 (-1)^{|\pi|-1}(|\pi|-1)!
 \prod_{B\in\pi}
 \mathbb E_t\prod_{q\in B}T_q.
\label{eq:v4v62-tilted-tail-cumulant}
\end{equation}

\begin{lemma}[Uniform tilted tail cumulant]
\label{lem:v4v62-tilted-tail-cumulant}
If all nonempty tail products satisfy
\[
 \left\|\prod_{q\in B}T_q\right\|_2
 \le\prod_{q\in B}\tau_q,
\]
then, uniformly for \(\Lambda(t)\le s_1\),
\begin{equation}
 |K_{\rm T}(t)|
 \le
 \ell!C_{\rm T}^{\ell}
 \prod_{q=1}^{\ell}\tau_q.
\label{eq:v4v62-tilted-tail-bound}
\end{equation}
\end{lemma}

\begin{proof}
Lemma~\ref{lem:v4v61-square-moment}, Cauchy--Schwarz, and
\(|M(t)|\ge1/2\) give, for each nonempty \(B\),
\[
 \left|
 \mathbb E_t\prod_{q\in B}T_q
 \right|
 \le
 2C_1^{1/2}\prod_{q\in B}\tau_q.
\]
In a partition product, every tail index occurs exactly once, so the product
of the \(\tau_q\)'s is independent of the partition.  Taking absolute
values in \eqref{eq:v4v62-tilted-tail-cumulant} leaves
\[
 \sum_{j=1}^{\ell}(j-1)!S(\ell,j)C_0^j.
\]
With \(C_0\ge1\), this is at most
\(C_0^\ell a_\ell\), where \(a_\ell\) is the ordered Bell number of
\eqref{eq:v4v48-ordered-Bell}.  Apply
\eqref{eq:v4v48-ordered-partition-bound} and absorb the fixed constants
into \(C_{\rm T}^{\ell}\).
\end{proof}

\begin{theorem}[Several tail insertions and arbitrary good valence]
\label{thm:v4v62-multi-tail-cumulant}
Let \(k=\ell+m\), with \(\ell\ge1\) and \(m\ge0\).  Under the hypotheses
above,
\begin{equation}
 \boxed{
 |\kappa_\nu(T_1,\ldots,T_\ell,F_1,\ldots,F_m)|
 \le
 k!C_*^k
 \left(\prod_{q=1}^{\ell}\tau_q\right)
 \left(\prod_{i=1}^{m}L_i\right).}
\label{eq:v4v62-multi-tail-cumulant}
\end{equation}
\end{theorem}

\begin{proof}
Normalized differentiation of a cumulant under an exponential tilt inserts
the corresponding centered load.  Iterating this identity gives
\begin{equation}
 \partial_{t_1}\cdots\partial_{t_m}K_{\rm T}(0)
 =
 \kappa_\nu(T_1,\ldots,T_\ell,F_1,\ldots,F_m).
\label{eq:v4v62-tilt-derivative}
\end{equation}
This can also be checked directly by differentiating the finite formula
\eqref{eq:v4v62-tilted-tail-cumulant}; no exponential of a tail variable is
formed.

For \(m\ge1\), take the Cauchy radii
\(r_i=s_1/(mL_i)\) and apply
Lemma~\ref{lem:v4v62-tilted-tail-cumulant}.  This gives
\[
 \ell!C_{\rm T}^{\ell}
 \left(\frac m{s_1}\right)^m
 \left(\prod_q\tau_q\right)\left(\prod_iL_i\right).
\]
Use \(m^m\le e^mm!\) and
\(\ell!m!\le(\ell+m)!=k!\).  The case \(m=0\) is
Lemma~\ref{lem:v4v62-tilted-tail-cumulant} at \(t=0\).  Enlarging one
constant proves \eqref{eq:v4v62-multi-tail-cumulant}.
\end{proof}

\begin{firewall}[The tilt contains only good Lipschitz loads]
\label{fw:v4v62-good-tilt-only}
The complex tilt \eqref{eq:v4v62-tilted-expectation} is zero-free because
\(W(t)\) contains only the V53 good loads.  Tail variables enter the finite
moment--cumulant polynomial
\eqref{eq:v4v62-tilted-tail-cumulant}.  None is exponentiated, and no tail
moment-generating function is assumed.
\end{firewall}

\subsection{The distinct-set projective ledger}
\label{sec:v4v62-projective-ledger}

For a coefficient family \(Q(S)\) indexed by finite distinct tail-cell
sets, define the weighted set norm
\begin{equation}
 \|Q\|_{\mathfrak T(a,\tau)}
 :=
 \sup_b
 \sum_{\substack{S\ni b\\S\ {\rm finite}}}
 e^{a|S|}\tau^{-|S|}\|Q(S;b)\|.
\label{eq:v4v62-set-norm}
\end{equation}
Products of sibling histories use the disjoint union from
\eqref{eq:v4v62-disjoint-union}.

\begin{lemma}[Disjoint union is projective]
\label{lem:v4v62-projective-union}
If \(S_1,\ldots,S_k\) are the tail sets of sibling histories, then
\begin{equation}
 e^{a|\dot\cup_iS_i|}\tau^{-|\dot\cup_iS_i|}
 =
 \prod_i e^{a|S_i|}\tau^{-|S_i|}.
\label{eq:v4v62-projective-union}
\end{equation}
Consequently the distinct-set weight adds no branching coefficient beyond
the ordinary unordered-child and star factors of V57.
\end{lemma}

\begin{proof}
Disjointness gives
\(|\dot\cup_iS_i|=\sum_i|S_i|\).  Exponentiation proves
\eqref{eq:v4v62-projective-union}.  The sum over set labels is the ordinary
labelled-set product already accompanied by \(1/k!\); no additional choice
of an owner is made.
\end{proof}

Theorem~\ref{thm:v4v62-multi-tail-cumulant} and
Lemma~\ref{lem:v4v62-projective-union} remove the local factorial and
same-cell idempotence obstructions.  A spatial condition is still required.

\begin{theorem}[Sufficient joint-response interface]
\label{thm:v4v62-joint-response-interface}
Suppose that, at a \(k\)-child collision centered at \(x\), every nonempty
product of tail-child loads satisfies the strengthened form
\begin{equation}
 \left\|\prod_{q\in B}T_q\right\|_2
 \le
 \prod_{q\in B}
 \left[
 \tau_q\lambda_qe^{-m_cd(S_q,x)}
 \right],
\label{eq:v4v62-joint-response}
\end{equation}
while every good child satisfies
\(L_i\le\lambda_i e^{-m_cd(S_i,x)}\).  Then the complete mixed cumulant has
the typed star bound
\begin{equation}
 |\kappa(T_1,\ldots,T_\ell,F_1,\ldots,F_m)|
 \le
 k!C_*^k
 \left(\prod_q\tau_q\right)
 \left(\prod_{i=1}^k\lambda_i\right)
 \sum_r\prod_{i\ne r}
 e^{-\frac12m_cd(S_i,S_r)}.
\label{eq:v4v62-mixed-star}
\end{equation}
\end{theorem}

\begin{proof}
Insert \eqref{eq:v4v62-joint-response} and the good response bounds into
Theorem~\ref{thm:v4v62-multi-tail-cumulant}.  The result contains the
product of all \(k\) pre-collision exponentials.  Drop the factor of the
child closest to \(x\) and use the same triangle inequality as in
Theorem~\ref{thm:v4v61-one-tail-star}.
\end{proof}

\begin{assessment}[The remaining analytic gate is joint spatial response]
\label{ass:v4v62-joint-response-open}
V34 proves the tail product factor in
\eqref{eq:v4v62-joint-response} at the hard scale.  V60 proves the required
spatial response and owner transfer for one tail child.  Neither result
currently proves the product of the separate pre-collision spatial factors
for an arbitrary family of tail descendants.  V52 warns that individual
\(L^2\) contractions cannot simply be multiplied to obtain this power.
Thus \eqref{eq:v4v62-joint-response} is a sufficient interface, not an
acquired theorem.
\end{assessment}

\subsection{V62 result ledger and V63 target}
\label{sec:v4v62-conclusion}

V62 consolidates every physical failed cell into one exact finite
Grassmann packet and proves that sibling collision histories have disjoint
tail-cell sets.  Simultaneous V34 tail moments therefore define a
projective distinct-set activity.  By tilting only with the good Lipschitz
loads, it proves an arbitrary-valence local cumulant bound with a product
of all distinct tail activities, a product of all good response scales, and
only factorial-exponential combinatorics.  This resolves the local
many-tail moment and ownership algebra.

The remaining gate is spatial: several tail descendants need the joint
pre-collision response \eqref{eq:v4v62-joint-response}.  V63 should test a
full-law forest interpolation in which every tail gradient remains inside
its own interpolated Witten edge.  If positivity of the interpolated Hessian
is uniform, this route can acquire the missing product without passing
through the false conditional \(L^\infty\) norm.  Otherwise the programme
must construct an exponential-Orlicz strong data-processing norm for the
conditional color channel.

\clearpage
\section{V63: guarded conditional expectation regularizes a tail carrier into a
Lipschitz load}
\label{sec:v4v63}

V62 proposed testing a full-law forest interpolation for the missing joint
tail response.  That test has already been performed in V37.  The third
bond derivative differentiates the Witten inverse through the unbounded
operator
\((\nabla U_e)\cdot\nabla+D^2U_e\), and a uniform convexity gap alone does
not bound the resulting forest.  Repeating that route would not add a new
estimate.

There is a more local upstream mechanism.  Before collision, a child
carrier is repeatedly replaced by a conditional expectation on the retained
skeleton.  On the guarded sector of V60, the tail is conditionally rare.
Its conditional expectation is therefore small in value.  Differentiating
that expectation in a retained direction gives an explicit derivative and
one actual conditional covariance with the boundary score.  The Witten
identity places the rare tail gradient inside this covariance and yields an
exponentially local retained derivative.  Multiplication by the smooth
guard preserves the same bound.

Consequently a guarded tail child becomes a bounded Lipschitz load of size
\(e^{-cR_Y^2}\).  Several independently regularized tail children can then
all be inserted in the ordinary V53 generator, which supplies their
separate pre-collision response factors and the complete star.  Only
branches whose first owners continue outward without yet being integrated
remain outside this regularized sector.

\subsection{The guarded conditional transform}
\label{sec:v4v63-transform}

Use the conditional geometry and the guard \(G_x\) of V60.  Let \(T_x\)
be a tail-owned carrier in the integrated color and put
\begin{equation}
 P_IT_x(\eta)
 :=
 \mathbb E_{\nu_I^\eta}T_x,
 \qquad
 U_x(\eta)
 :=
 G_x(\eta)P_IT_x(\eta).
\label{eq:v4v63-guarded-transform}
\end{equation}
Assume the V59 tail envelopes through two background marks and the weighted
boundary-score bound
\eqref{eq:v4v60-boundary-score}.  Permit an explicit retained derivative of
\(T_x\), but require its conditional \(L^2\) norm on the guarded sector to
have the same rare envelope and an exponentially summable support kernel.

\begin{lemma}[Exact retained derivative]
\label{lem:v4v63-transform-derivative}
For every retained direction \(h_y\),
\begin{equation}
 D_{\eta_y}P_IT_x[h_y]
 =
 \mathbb E_{\nu_I^\eta}D_{\eta_y}T_x[h_y]
 -
 \operatorname{Cov}_{\nu_I^\eta}
 \left(T_x,K_y^\eta[h_y]\right),
\label{eq:v4v63-transform-derivative}
\end{equation}
where \(K_y^\eta=D_{\eta_y}S_I^\eta\) is the actual boundary score.
\end{lemma}

\begin{proof}
Differentiate the numerator and normalizing denominator of the conditional
expectation.  The derivative of the density gives
\(-K_y^\eta\), and the quotient rule centers it, producing the displayed
covariance.  The explicit derivative of \(T_x\) is the first term.
\end{proof}

\begin{theorem}[Guarded tail regularization]
\label{thm:v4v63-tail-regularization}
There are \(c,\mu,C>0\), independent of volume, such that for
\(0\le j\le2\),
\begin{align}
 \sup_\eta|D_B^jU_x(\eta)|
 &\le
 C_je^{-cR_Y^2},
\label{eq:v4v63-small-value}\\
 \sup_\eta\sum_{y\in\Sigma}
 e^{\mu d(x,y)}
 \left|
 D_{\eta_y}D_B^jU_x(\eta)
 \right|
 &\le
 C_je^{-cR_Y^2}.
\label{eq:v4v63-small-Lipschitz}
\end{align}
Thus the centered retained carrier
\(U_x-\mathbb E_{\nu_\Sigma}U_x\) is a V53 Lipschitz load with individual
scale
\begin{equation}
 L_x^{\rm T}\le C e^{-cR_Y^2}.
\label{eq:v4v63-tail-Lipschitz-scale}
\end{equation}
\end{theorem}

\begin{proof}
If \(G_x(\eta)\ne0\), Lemma
\ref{lem:v4v60-owner-partition} bounds the conditional mean in the support
of \(T_x\) by \(R_Y/2\).  Lemma~\ref{lem:v4v60-guarded-tail}, with
\(p=1\), gives
\[
 |P_IT_x(\eta)|\le Ce^{-cR_Y^2},
\]
which proves \eqref{eq:v4v63-small-value}.

Differentiate \eqref{eq:v4v63-guarded-transform}:
\begin{equation}
 D_{\eta_y}U_x
 =
 (D_{\eta_y}G_x)P_IT_x
 +
 G_xD_{\eta_y}P_IT_x.
\label{eq:v4v63-product-derivative}
\end{equation}
For the second term, use
Lemma~\ref{lem:v4v63-transform-derivative}.  The explicit derivative has
the assumed rare local kernel.  In the covariance, apply the conditional
Helffer--Sjostrand identity.  The V60 guarded tail-gradient bound pays
\(e^{-cR_Y^2}\); convolving
\eqref{eq:v4v60-conditional-Witten} with the weighted score gradient
\eqref{eq:v4v60-boundary-score} gives
\begin{equation}
 \left|
 \operatorname{Cov}_{\nu_I^\eta}
 (T_x,K_y^\eta)
 \right|
 \le
 Ce^{-cR_Y^2}e^{-\mu_0d(x,y)}.
\label{eq:v4v63-tail-score-edge}
\end{equation}

For the first term in
\eqref{eq:v4v63-product-derivative}, a derivative of \(G_x\) is a sum of
terms with one derivative of \(g_{xy}\).  Such a term is supported where
\(r_y^\Sigma\le2L_{xy}\) and every other surviving guard factor is nonzero,
so the conditional-mean argument still gives
\(|P_IT_x|\le Ce^{-cR_Y^2}\).  Moreover,
\[
 |D_{\eta_y}g_{xy}|
 \le
 \frac C{L_{xy}}
 =
 \frac C{\theta R_Y}e^{-\beta d(x,y)}.
\]
Choose \(0<\mu<\min\{\mu_0,\beta\}\) and sum the lattice shells.  This proves
\eqref{eq:v4v63-small-Lipschitz} for \(j=0\).

A background derivative has the exact placements listed in V54--V56.
It either differentiates the explicit tail coefficient, creates a first or
second score, or produces two first scores.  Conditional normalized
differentiation gives cumulants of order at most four.  On every surviving
guard term, the complete product containing \(T_x\) has the rare gradient
bound; apply the fixed-order cut proof of
Theorem~\ref{thm:v4v59-tail-cumulants}.  Differentiating a guard in a
retained direction is treated by the preceding annulus estimate.  This
proves the two marked versions.  Centering changes no retained derivative,
so \eqref{eq:v4v63-tail-Lipschitz-scale} follows.
\end{proof}

\begin{firewall}[Regularization occurs after a genuine conditional
integration]
\label{fw:v4v63-real-regularization}
The small Lipschitz load in
\eqref{eq:v4v63-tail-Lipschitz-scale} is the conditional expectation
\(P_IT_x\) multiplied by its guard.  The original tail polynomial \(T_x\)
is still unbounded and is not assigned this seminorm.  A collision occurring
before the relevant conditional integration must use the V61 inserted
generator or the V62 tail moment formula.
\end{firewall}

\subsection{Propagation after regularization}
\label{sec:v4v63-propagation}

Let \(\mathcal P\) be one subsequent centered conditional-expectation
channel in the V41 color sweep.

\begin{lemma}[A regularized load remains regularized]
\label{lem:v4v63-propagation}
If \(U\) is bounded and has a weighted retained-gradient row norm \(L_U\),
then \(\mathcal PU\) is bounded by \(\|U\|_\infty\), and its centered
\(L^2\) norm contracts by the V41 factor.  Its gradient with respect to the
next retained color is bounded by
\begin{equation}
 L_{\mathcal PU}
 \le C_{\rm ch}L_U
\label{eq:v4v63-channel-Lipschitz}
\end{equation}
with a volume-independent constant and the same exponential locality after
reducing the exponent.
\end{lemma}

\begin{proof}
Conditional expectation is an \(L^\infty\) contraction.  The centered
\(L^2\) statement is Theorem~\ref{thm:v4v41-angle}.  Differentiate the
conditional expectation as in
Lemma~\ref{lem:v4v63-transform-derivative}.  The explicit derivative is
controlled by \(L_U\).  The score covariance is a Witten carrier between
the actual gradient of \(U\) and the mixed Hessian score.  Apply the
conditional weighted resolvent and convolve its kernel with the gradient
row of \(U\).  This gives \eqref{eq:v4v63-channel-Lipschitz}.
\end{proof}

Repeated color round trips are summed with the strict V41 resolvent.  Thus a
tail child that has entered the regularized state keeps its rare scale and
acquires the ordinary pre-collision response factor:
\begin{equation}
 L_i^{\rm T}
 \le
 C\tau_i e^{-m_cd(S_i,x_{\rm col})}.
\label{eq:v4v63-regularized-response}
\end{equation}

\subsection{Several regularized tails at one collision}
\label{sec:v4v63-several-tails}

\begin{theorem}[The mixed star on the fully regularized branch]
\label{thm:v4v63-regularized-star}
Consider a \(k\)-child collision with \(\ell\ge1\) regularized tail
children and \(k-\ell\) ordinary good children.  Suppose the tail children
satisfy \eqref{eq:v4v63-regularized-response}, and the good children satisfy
\eqref{eq:v4v53-child-response}.  Then
\begin{align}
 |\kappa(F_1,\ldots,F_k)|
 &\le
 k!C_*^k
 \left(\prod_{i\in{\rm T}}\tau_i\lambda_i\right)
 \left(\prod_{i\in{\rm G}}\lambda_i\right)\nonumber\\
 &\quad\times
 \sum_{r=1}^k
 \prod_{i\ne r}
 e^{-\frac12m_cd(S_i,S_r)}.
\label{eq:v4v63-regularized-star}
\end{align}
The estimate holds through two background marks.
\end{theorem}

\begin{proof}
Every child is now a centered Lipschitz load.  For a tail child, its
individual Lipschitz constant contains \(\tau_i\) by
\eqref{eq:v4v63-regularized-response}; for a good child use
\eqref{eq:v4v53-child-response}.  Apply
Theorem~\ref{thm:v4v53-Lipschitz-cumulant}.  The product contains every
child's pre-collision exponential response.  Drop the factor of a closest
child and use the V53 triangle argument to obtain the displayed star.

A mark either differentiates one regularized load or adds one actual score
load.  Theorem~\ref{thm:v4v63-tail-regularization} supplies the marked tail
Lipschitz constants, and V54--V56 supply the score constants.  The resulting
polynomial branching factors preserve the radius.
\end{proof}

\begin{corollary}[Acquisition of the V62 joint-response interface on the
regularized branch]
\label{cor:v4v63-joint-response}
For collision histories in which every tail descendant undergoes one
guarded conditional integration before meeting another tail descendant,
the sufficient interface
\eqref{eq:v4v62-joint-response} holds.  Their complete rooted history sum is
therefore covered by the V57--V58 star recursion, with one factor
\(e^{-cR_Y^2}\) for every distinct physical tail cell.
\end{corollary}

\begin{proof}
Apply Theorem~\ref{thm:v4v63-tail-regularization} to each tail descendant,
Lemma~\ref{lem:v4v63-propagation} to its pre-collision channels, and
Theorem~\ref{thm:v4v63-regularized-star} at every merger.  Distinct tail
sets combine projectively by
Lemma~\ref{lem:v4v62-projective-union}; the V57 rooted sum then applies.
\end{proof}

\subsection{The owner branch}
\label{sec:v4v63-owner-branch}

The complementary transform is
\begin{equation}
 V_{xy}(\eta)
 :=
 B_{xy}(\eta)P_IT_x(\eta).
\label{eq:v4v63-owner-transform}
\end{equation}
V60 proves that its value and retained gradient have a super-summable
skeleton \(L^p\) norm.  It is not a globally bounded Lipschitz load, because
the conditional mean may grow with the owner field.

\begin{proposition}[Exact regularized/owner dichotomy]
\label{prop:v4v63-dichotomy}
For every conditional tail transform,
\begin{equation}
 P_IT_x
 =
 U_x+\sum_yV_{xy},
\label{eq:v4v63-dichotomy}
\end{equation}
where \(U_x\) satisfies
\eqref{eq:v4v63-small-value}--%
\eqref{eq:v4v63-small-Lipschitz} and the owner family satisfies
\begin{equation}
 \sum_y
 \left(
 \|V_{xy}\|_{L^p(\nu_\Sigma)}
 +
 \|\nabla_\eta V_{xy}\|_{L^p(\nu_\Sigma)}
 \right)
 \le
 C_pe^{-c_pR_Y^2}.
\label{eq:v4v63-owner-family}
\end{equation}
\end{proposition}

\begin{proof}
Multiply \(P_IT_x\) by the exact partition
\eqref{eq:v4v60-owner-identity}.  The regularized estimate is
Theorem~\ref{thm:v4v63-tail-regularization}.  For the owner terms, use the
conditional polynomial/exponential moment envelope of \(P_IT_x\), keep
\(B_{xy}\) attached, and apply
Lemma~\ref{lem:v4v60-owner-small} and its row sum.
\end{proof}

\begin{assessment}[Only simultaneous unregularized owner chains remain]
\label{ass:v4v63-owner-chains}
V63 acquires the missing product of spatial response factors whenever each
tail descendant has reached the regularized branch.  A single
unregularized owner can still be handled by the V61 one-insertion theorem.
The unresolved geometry is a collision of two or more owner branches before
their owner cells have been conditionally integrated.  Their joint
skeleton \(L^p\) norm is rare by V62, but they are not yet separate
Lipschitz loads, so V53 does not provide one response edge per owner.
\end{assessment}

\begin{firewall}[V37 is not revived under a new name]
\label{fw:v4v63-no-global-forest}
No derivative with respect to a Higgs kinetic bond is used in V63.  The
only derivatives are retained-field responses of actual conditional
expectations, represented by one score covariance.  The V41 sweep
resolvent sums these one-use channels.  Thus the unbounded third-order
Witten insertion identified in V37 is not assumed controlled.
\end{firewall}

\subsection{V63 result ledger and V64 target}
\label{sec:v4v63-conclusion}

V63 proves that one guarded conditional integration maps an unbounded tail
carrier to a uniformly small, exponentially local Lipschitz load.  The map
is exact through two background marks, and subsequent color channels
preserve the rare scale.  Several regularized tail descendants therefore
enter the original V53 generator simultaneously and acquire a complete
star with every separate tail activity.  This proves the V62 joint-response
interface on the fully regularized branch.

V64 should analyze simultaneous owner chains.  The first test is geometric:
in the nested V38 core/collar elimination, determine whether two owner
branches can collide before at least one owner cell is integrated.  If the
elimination order forbids that event, V63 completes the typed tail
interface.  If it permits it, construct an exact stopping forest on owner
paths and use the V62 simultaneous-tail norm at their first confluence,
without applying a false \(L^p\)-to-Lipschitz conversion.

\clearpage
\section{V64: a group mean guard preserves every distinct tail-cell payment}
\label{sec:v4v64}

V63 leaves branches in which several tail descendants reach the same
retained skeleton before their owner cells are integrated.  Multiplying
their individual first-owner bounds can pay the same large boundary field
more than once.  The correction is to guard the complete distinct tail set
at once.

For a set \(S\) of \(n\) physical tail cells, consider the vector of all
conditional means on \(S\).  If its norm is at most a fixed fraction of
\(R_Y\sqrt n\), simultaneous occurrence of all \(n\) tail events forces a
centered fluctuation of that size and costs \(e^{-cR_Y^2n}\).  If the mean
vector is larger, exponential response of the mean implies that an
exponentially weighted quadratic norm of the retained boundary is at least
\(cR_Y^2n\).  A weighted version of the V31 subgaussian estimate pays the
same factor.  Thus the two branches both retain one rare activity per
distinct physical tail cell, without choosing incompatible individual
owners.

\subsection{A weighted quadratic consequence of linear subgaussianity}
\label{sec:v4v64-weighted-subgaussian}

Let \(X\) be a finite-dimensional random vector satisfying
\begin{equation}
 \mathbb E\exp\bigl(t\,a\cdot(X-\mathbb EX)\bigr)
 \le
 \exp\!\left(\frac{t^2|a|^2}{2\kappa_*}\right)
\label{eq:v4v64-linear-subgaussian}
\end{equation}
for every real \(a,t\).  This is
\eqref{eq:v4v31-linear-subgaussian} for a \(\kappa_*\)-strongly convex
law.

\begin{lemma}[Weighted quadratic exponential]
\label{lem:v4v64-weighted-quadratic}
Let \(Q\ge0\) be a deterministic matrix with
\(\|Q\|\le q_*\) and \(\operatorname{Tr}Q\le q_1n\).  If
\(0<\delta<\kappa_*/(4q_*)\), then
\begin{equation}
 \mathbb E
 \exp\!\left(
 \delta\,
 \langle X-\mathbb EX,Q(X-\mathbb EX)\rangle
 \right)
 \le
 \exp(C_{\kappa_*,q_*}\delta q_1n).
\label{eq:v4v64-weighted-quadratic}
\end{equation}
If also
\(\langle\mathbb EX,Q\mathbb EX\rangle\le q_2n\), then, after reducing
\(\delta\),
\begin{equation}
 \mathbb E e^{\delta\langle X,QX\rangle}
 \le e^{C n}.
\label{eq:v4v64-uncentered-weighted}
\end{equation}
\end{lemma}

\begin{proof}
Let \(G\) be an independent standard Gaussian vector.  The Gaussian
linearization identity gives
\[
 e^{\delta\langle z,Qz\rangle}
 =
 \mathbb E_G
 \exp\!\left(
 \sqrt{2\delta}\,G\cdot Q^{1/2}z
 \right).
\]
Apply this with \(z=X-\mathbb EX\), use Tonelli, and then use
\eqref{eq:v4v64-linear-subgaussian} conditionally on \(G\).  The result is
\[
 \mathbb E_G
 \exp\!\left(
 \frac{\delta}{\kappa_*}\langle G,QG\rangle
 \right)
 =
 \det\!\left(1-\frac{2\delta}{\kappa_*}Q\right)^{-1/2}.
\]
For \(2\delta\|Q\|/\kappa_*<1/2\), the logarithm is bounded by
\[
 C\frac{\delta}{\kappa_*}\operatorname{Tr}Q
 \le C_{\kappa_*,q_*}\delta q_1n.
\]
This proves \eqref{eq:v4v64-weighted-quadratic}.  Finally use
\[
 \langle X,QX\rangle
 \le
 2\langle X-\mathbb EX,Q(X-\mathbb EX)\rangle
 +2\langle\mathbb EX,Q\mathbb EX\rangle
\]
and the assumed mean bound.
\end{proof}

\subsection{The influence weight of a tail set}
\label{sec:v4v64-influence}

Let \(S\) be a set of \(n\ge1\) distinct integrated tail cells.  With
\(m_m\) from Lemma~\ref{lem:v4v60-mean-response}, define
\begin{equation}
 q_S(y)
 :=
 \sum_{x\in S}e^{-m_md(x,y)}.
\label{eq:v4v64-influence-weight}
\end{equation}

\begin{lemma}[Uniform Schur and trace bounds]
\label{lem:v4v64-influence-bounds}
On a bounded-degree block lattice,
\begin{equation}
 \sup_yq_S(y)\le C_m,
 \qquad
 \sum_yq_S(y)\le C_mn.
\label{eq:v4v64-influence-bounds}
\end{equation}
Moreover, if
\(m_S(\eta)=(m_x(\eta):x\in S)\), then
\begin{equation}
 \|m_S(\eta)-m_S(0)\|_{\ell^2(S)}^2
 \le
 C_m
 \sum_yq_S(y)|\eta_y|^2.
\label{eq:v4v64-vector-mean-response}
\end{equation}
\end{lemma}

\begin{proof}
The exponential lattice kernel has a finite row and column sum.  The first
bound in \eqref{eq:v4v64-influence-bounds} follows by summing over the
subset \(S\), and the second follows by interchanging the \(x,y\) sums.

For each \(x\), Lemma~\ref{lem:v4v60-mean-response} and weighted
Cauchy--Schwarz give
\[
 |m_x(\eta)-m_x(0)|^2
 \le
 C\sum_y e^{-m_md(x,y)}|\eta_y|^2.
\]
Sum over \(x\in S\) and use
\eqref{eq:v4v64-influence-weight}.
\end{proof}

The diagonal operator \(Q_S\) which multiplies the retained cell \(y\) by
\(q_S(y)\) has bounded norm and trace \(O(n)\).  The skeleton marginal is
strongly convex by V38 and has a uniformly bounded recentered mean.
Lemma~\ref{lem:v4v64-weighted-quadratic} therefore gives
\begin{equation}
 \mathbb E_{\nu_\Sigma}
 \exp\!\left(
 \delta\sum_yq_S(y)|\eta_y|^2
 \right)
 \le e^{C_\delta n}.
\label{eq:v4v64-skeleton-weighted-mgf}
\end{equation}

\subsection{The exact group guard}
\label{sec:v4v64-group-guard}

Use the same smooth cutoff \(\chi\) and define
\begin{equation}
 \rho_S(\eta)
 =
 \left(1+\|m_S(\eta)\|_{\ell^2(S)}^2\right)^{1/2},
 \qquad
 G_S(\eta)
 =
 \chi\!\left(
 \frac{4\rho_S(\eta)}{R_Y\sqrt n}
 \right),
 \qquad
 B_S=1-G_S.
\label{eq:v4v64-group-guard}
\end{equation}
This is an exact two-branch partition \(1=G_S+B_S\).

\begin{lemma}[Good group and boundary-energy alternatives]
\label{lem:v4v64-group-alternatives}
For small \(\epsilon_Y\):
\begin{enumerate}
\item on the support of \(G_S\),
\begin{equation}
 \|m_S(\eta)\|_{\ell^2(S)}
 \le\frac12R_Y\sqrt n;
\label{eq:v4v64-good-group-mean}
\end{equation}
\item on the support of \(B_S\),
\begin{equation}
 \sum_yq_S(y)|\eta_y|^2
 \ge c_0R_Y^2n.
\label{eq:v4v64-bad-group-energy}
\end{equation}
\end{enumerate}
\end{lemma}

\begin{proof}
If \(G_S\ne0\), the argument of \(\chi\) is less than \(2\), which gives
\(\rho_S<R_Y\sqrt n/2\) and proves the first assertion.

The zero-boundary mean hypothesis
\eqref{eq:v4v60-zero-boundary-mean} gives
\(\|m_S(0)\|_2\le M_0\sqrt n\le R_Y\sqrt n/8\) for small
\(\epsilon_Y\).  If \(B_S\ne0\), the argument of \(\chi\) is greater than
one, so
\(\rho_S>R_Y\sqrt n/4\).  After harmlessly increasing the lower smallness
threshold so that the added \(1\) in \(\rho_S\) is negligible,
\[
 \|m_S(\eta)-m_S(0)\|_2
 \ge cR_Y\sqrt n.
\]
Apply \eqref{eq:v4v64-vector-mean-response} and square the constants to
obtain \eqref{eq:v4v64-bad-group-energy}.
\end{proof}


\begin{lemma}[Conditional envelope in the boundary energy]
\label{lem:v4v64-conditional-envelope}
Put
\[
 E_S(\eta)=\sum_yq_S(y)|\eta_y|^2.
\]
For every sufficiently small \(\delta>0\), the finite-fibre tail packets
and their first two background marks obey
\begin{equation}
 \|D_B^jT_S\|_{L^2(\nu_I^\eta)}
 +
 \|\nabla_{\phi_I}D_B^jT_S\|_{L^2(\nu_I^\eta)}
 \le
 C_{\delta,j}^{\,n}
 \exp\!\left(\delta E_S(\eta)\right),
 \qquad 0\le j\le2.
\label{eq:v4v64-conditional-envelope}
\end{equation}
The same bound holds after one explicit retained derivative, with a changed
constant.
\end{lemma}

\begin{proof}
Strong convexity makes the conditional fluctuation about
\(m_S(\eta)\) subgaussian with constants independent of the boundary.
Each consolidated local Grassmann packet and each of its fixed marks is a
fixed-degree polynomial times a linear exponential in its local field.
Conditional Holder estimates and bounded overlap therefore give
\[
 C^n\prod_{x\in S}
 (1+|m_x(\eta)|)^p e^{c\epsilon_Y|m_x(\eta)|}
\]
for fixed \(p,c\), rather than one polynomial whose degree is hidden in
\(\|m_S(\eta)\|_2\).

By \eqref{eq:v4v64-vector-mean-response},
\[
 \|m_S(\eta)\|_2^2
 \le
 2M_0^2n+2C_mE_S(\eta).
\]
Apply Young's inequality separately at every \(x\in S\).  The product is
bounded by \(C_{\delta,j}^{\,n}\exp(\delta\|m_S\|_2^2)\), and the
preceding display then gives, after renaming \(\delta\),
\(C_{\delta,j}^{\,n}\exp(\delta E_S)\).
A field or retained derivative adds only a fixed local polynomial, a cutoff
derivative, or one boundary-score covariance; the weighted Witten estimate
has a bounded Schur row and gives the same form.
\end{proof}
\begin{theorem}[Both group branches pay every distinct tail cell]
\label{thm:v4v64-group-payment}
Let \(T_S\) be a complete product of locally consolidated V62 tail packets
with distinct cell set \(S\), including a fixed number of ordinary
fermion-tree factors and at most two background marks.  There are
\(c,C>0\) such that
\begin{align}
 \sup_{\eta:\,G_S(\eta)\ne0}
 \|T_S\|_{L^2(\nu_I^\eta)}
 &\le
 C^ne^{-cR_Y^2n},
\label{eq:v4v64-good-group-payment}\\
 \|B_S\,\mathbb E_{\nu_I^\eta}T_S\|_{L^2(\nu_\Sigma)}
 &\le
 C^ne^{-cR_Y^2n}.
\label{eq:v4v64-bad-group-payment}
\end{align}
The same estimates hold with one retained derivative, with a changed
\(C^n\).
\end{theorem}

\begin{proof}
On the good group, simultaneous occurrence of all tail events implies
\[
 \|\phi_S\|_2\ge cR_Y\sqrt n.
\]
By \eqref{eq:v4v64-good-group-mean}, the centered conditional fluctuation
has norm at least \(c'R_Y\sqrt n\).  The conditional law is
\(\kappa_{\rm p}\)-strongly convex.  Apply the Gaussian linearization
argument of Lemma~\ref{lem:v4v64-weighted-quadratic} with the identity on
the \(O(n)\)-dimensional cell space, and absorb the fixed polynomial and
linear-exponential tail envelope.  This gives
\eqref{eq:v4v64-good-group-payment}.

On the bad group, use the conditional envelope
\eqref{eq:v4v64-conditional-envelope}, choosing its exponent smaller than
the quadratic-moment budget.
Keep \(B_S\) attached.  Its support has the energy lower bound
\eqref{eq:v4v64-bad-group-energy}.  For small fixed \(\delta>0\),
\[
 \boldsymbol1_{\{B_S\ne0\}}
 \le
 e^{-\delta c_0R_Y^2n}
 \exp\!\left(
 \delta\sum_yq_S(y)|\eta_y|^2
 \right).
\]
Use Holder's inequality and
\eqref{eq:v4v64-skeleton-weighted-mgf}, with a slightly enlarged
\(\delta\), to absorb the conditional envelope.  Since \(R_Y\to\infty\),
the factor \(e^{C_\delta n}\) is absorbed by reducing the tail exponent.
This proves \eqref{eq:v4v64-bad-group-payment}.

A retained derivative either differentiates the conditional expectation or
the guard.  The former gives the explicit derivative and score covariance
of Lemma~\ref{lem:v4v63-transform-derivative}; the latter contributes
\((R_Y\sqrt n)^{-1}D_\eta m_S\) and is supported in the same transition
energy shell.  Weighted Witten response, the preceding exponential
moments, and at most polynomially many marked placements change \(C^n\)
but preserve the exponent \(cR_Y^2n\).
\end{proof}

\begin{firewall}[A shared boundary fluctuation is paid by group energy]
\label{fw:v4v64-shared-owner}
If one retained region shifts several conditional means, V64 does not
multiply several copies of the same indicator.  It measures the norm of the
entire shifted mean vector.  Bounded spatial influence converts that norm
to weighted boundary energy proportional to \(n\), and the single group
event is then paid at cost \(e^{-cR_Y^2n}\).
\end{firewall}

\subsection{A set-valued conditional transfer}
\label{sec:v4v64-set-transfer}

Define
\begin{align}
 \mathcal U_S(\eta)
 &:=
 G_S(\eta)\mathbb E_{\nu_I^\eta}T_S,
\label{eq:v4v64-good-transform}\\
 \mathcal V_S(\eta)
 &:=
 B_S(\eta)\mathbb E_{\nu_I^\eta}T_S.
\label{eq:v4v64-bad-transform}
\end{align}

\begin{theorem}[The distinct-set activity survives one elimination]
\label{thm:v4v64-set-transfer}
In the projective set norm of
\eqref{eq:v4v62-set-norm}, one conditional elimination maps a tail set
\(S\) to the exact sum
\begin{equation}
 \mathbb E_{\nu_I^\eta}T_S
 =
 \mathcal U_S+\mathcal V_S
\label{eq:v4v64-set-transfer}
\end{equation}
and both terms retain the activity
\begin{equation}
 \tau^{|S|},
 \qquad
 \tau=Ce^{-cR_Y^2}.
\label{eq:v4v64-set-activity}
\end{equation}
The good term is a bounded quasi-local Lipschitz load.  The bad term has
the same activity in skeleton \(L^2\), and its retained gradient is
quasi-local with influence weight \(q_S\).  These assertions hold through
two background marks.
\end{theorem}

\begin{proof}
Exactness is \(1=G_S+B_S\).  The activity bounds are
Theorem~\ref{thm:v4v64-group-payment}.  For the good term, differentiate
as in Theorem~\ref{thm:v4v63-tail-regularization}; the gradient of a product
of tail packets is paid by the good-group simultaneous fluctuation, while
the score covariance supplies the Witten kernel.  For the bad term, the
proof of \eqref{eq:v4v64-bad-group-payment} also applies after one retained
derivative.  Lemma~\ref{lem:v4v64-influence-bounds} gives a bounded row and
trace for \(q_S\).  Background marks have only the V54--V56 placements and
are already included in Theorem~\ref{thm:v4v64-group-payment}.
\end{proof}

\begin{corollary}[No activity loss at simultaneous owner confluence]
\label{cor:v4v64-owner-confluence}
When several unregularized owner branches meet before individual
regularization, first take the union of their disjoint physical tail sets
and apply the group guard \eqref{eq:v4v64-group-guard}.  The resulting
carrier retains one factor \(\tau\) for every distinct original failed
cell.  Same-cell occurrences were already consolidated by V62.
\end{corollary}

\begin{proof}
Sibling tail sets are disjoint by
Lemma~\ref{lem:v4v62-disjoint-children}; hence their union has additive
cardinality and its projective weight factors by
Lemma~\ref{lem:v4v62-projective-union}.  Apply
Theorem~\ref{thm:v4v64-set-transfer} to that union.
\end{proof}

\begin{assessment}[Remaining all-depth statement]
\label{ass:v4v64-all-depth}
V64 removes the loss caused by several owner branches sharing one boundary
fluctuation.  At every single elimination or owner confluence, the full
distinct tail-set activity survives, and the good output is immediately
eligible for V53.  The bad group output is small in a quasi-local skeleton
\(L^2\) norm rather than \(L^\infty\).  To close the V58 hierarchy at
arbitrary depth, the collision compiler must schedule the next elimination
of this group output before using it as an unguarded V53 load, or carry its
mixed \(L^2\)/Lipschitz state through a finite stopping forest.
\end{assessment}

\begin{firewall}[The weighted energy is not a new physical vertex]
\label{fw:v4v64-no-new-vertex}
The operator \(Q_S\) and the guard \(B_S\) occur only in the exact
partition of the conditional estimate.  They add no determinant factor,
Higgs interaction, or tail cell to the integrand.  Their role is to decide
whether the already selected physical tail set is paid by conditional
fluctuation or by retained boundary energy.
\end{firewall}

\subsection{V64 result ledger and V65 target}
\label{sec:v4v64-conclusion}

V64 proves a weighted quadratic subgaussian theorem with constants
controlled by operator norm and trace.  Exponential conditional-mean
response associates to every distinct tail set \(S\) a boundary influence
operator of bounded norm and trace \(O(|S|)\).  The exact group guard then
shows that either all physical tail cells are paid by conditional
fluctuation or the retained boundary has energy \(cR_Y^2|S|\); both branches
cost \(e^{-cR_Y^2|S|}\).  This resolves simultaneous owner confluence
without duplicate event payments and preserves two marks.

V65 must begin with the next companion audit.  It should then construct the
finite stopping forest for the remaining bad group outputs.  The desired
statement is that every path either reaches a guarded regularization leaf or
moves to a strictly larger V38 ancestor block, so path depth is logarithmic
in its support scale and the V41 sweep contraction pays propagation between
stops.  The proof must keep the same tail-set label through the path and
must not charge the group energy again at each ancestor.

\clearpage
\section{V65: companion audit and an exact first-success tail-set forest}
\label{sec:v4v65}

V64 proves that one conditional elimination preserves
\(\tau^{|S|}\), \(\tau=Ce^{-cR_Y^2}\), for every distinct physical tail set
\(S\), even when several cells share one boundary fluctuation.  The
remaining issue is repeated promotion through a block hierarchy.  One may
not pay the V64 boundary energy independently at every ancestor, and the
V41 color-sweep contraction cannot be inserted as an artificial
contraction for block promotion.

This note performs the scheduled companion audit and constructs a different
stopping forest.  At every ancestor, a background-independent weighted
boundary energy decides whether the complete tail set can be integrated
there.  The first-success weights are nonnegative and sum to one pointwise
on the original field.  Therefore summing all possible stopping levels
does not multiply the tail activity.  Their retained-field derivatives are
summable because the distance from the tail set to the exterior skeleton
grows geometrically.  Background marks never hit the guards.

Colliding packets are merged before their stopping partition is estimated.
Their physical tail sets are replaced by one disjoint union and receive one
new first-success partition at the least common ancestor.  At the root the
exterior energy is zero, so every path terminates.

\subsection{Compulsory companion audit}
\label{sec:v4v65-audit}

The exact latest companion files inspected were:
\begin{enumerate}
\item
\texttt{Renewal\_Yang\_Mills\_Mass\_Gap\_Research\_Notes\_V82.tex},
\(1\,423\,323\) bytes, SHA--256
\[
 \texttt{57F4AB279B64FBFF4B735D6EEDF2D3878E45E43AF0E9729FE62673E4457D2FE9}.
\]
YM V82 computes the scalar marginal of a complete output isotypic
projection with arbitrary multiplicity and obtains the inverse
larger-input dimension after summing the heat output.  It also proves that
this scalar product-state capacity does not imply the partial-transpose
norm required for tensor-stable bank aggregation.  The complete isotypic
output and its orientation must remain attached.
\item
\texttt{Renewal\_NS\_Stability\_Research\_Notes\_V31.tex},
\(887\,537\) bytes, SHA--256
\[
 \texttt{F1121B62238AD5491BB7CF03635EA9934942EDF573ED8FAAA9EE79E1D38A6696}.
\]
Its endpoint still proves that removing an exact flat stopping mark costs
the critical first-moment weight, that heat lifting restores that weight
only in the critical source norm, and that the signed hinge has a balanced
nullspace.  The actual recent-entry event, signed triad work, and high
parent must remain in one correlated carrier.
\end{enumerate}

The transfer is proof order.  A scalar statement that one branch costs
\(\tau^{|S|}\) is not automatically stable under the next conditional
operation, just as the YM scalar capacity is not automatically tensor
stable.  The construction below retains the full coefficient and its
tail-set label under the stopping partition.  The NS audit forbids a
separate cost at every stopping level; all levels are assembled by one
partition of unity before an absolute moment estimate is taken.

\subsection{The finite-range parent and nested buffers}
\label{sec:v4v65-parent}

Use the finite-range positive Higgs parent \(\mu_B\) of V30--V31.  The
good-modulus, phase, and tail determinant coefficients remain external
activities.  This preserves the exact conditional product of separated
cores in \eqref{eq:v4v38-conditional-product}.  Every conditional
expectation is taken in the original master probability space; an effective
marginal is not introduced as a second parent.

Let
\[
 B_0\subset B_1\subset\cdots\subset B_J
\]
be the ancestor chain of a packet carrying a distinct tail set \(S\), with
\(n=|S|\).  Its side lengths satisfy
\begin{equation}
 L_j=L_0b^j,\qquad b>1.
\label{eq:v4v65-geometric-scales}
\end{equation}
Choose the conditional interior at level \(j\) so that \(S\) lies behind a
buffer of width
\begin{equation}
 \ell_j\ge c_L L_j
\label{eq:v4v65-buffer-width}
\end{equation}
from its exterior skeleton \(\Sigma_j\).  At the root
\(\Sigma_J=\varnothing\).

\begin{firewall}[Parent representations remain separate]
\label{fw:v4v65-parent}
The V34 promoted good-modulus law may be used for the full-law V59
estimate, but its Hessian is not finite range.  V65 uses the original
finite-range strongly convex Higgs law so that V38 conditional
factorization is exact.  The V34 tail moments and V53 log-Sobolev estimates
both apply to this parent, with all determinant factors left in their
activity ancestry.
\end{firewall}

\subsection{A background-independent exterior-energy guard}
\label{sec:v4v65-energy-guard}

For \(y\in\Sigma_j\), define
\begin{equation}
 q_{j,S}(y)
 :=
 \sum_{x\in S}e^{-m_0d(x,y)},
 \qquad
 E_{j,S}(\eta)
 :=
 \sum_{y\in\Sigma_j}q_{j,S}(y)|\eta_y|^2,
\label{eq:v4v65-exterior-energy}
\end{equation}
where \(m_0>0\) is smaller than the V38 Witten exponent.  Since
\(d(S,\Sigma_j)\ge\ell_j\), exponential lattice summation gives
\begin{align}
 \|Q_{j,S}\|
 &:=
 \sup_yq_{j,S}(y)
 \le C e^{-m_1\ell_j},
\label{eq:v4v65-Qnorm}\\
 \operatorname{Tr}Q_{j,S}
 &=
 \sum_yq_{j,S}(y)
 \le Cn e^{-m_1\ell_j}
\label{eq:v4v65-Qtrace}
\end{align}
after reducing \(m_1>0\).

The V60 mean-response proof, restricted to the exterior
\(\Sigma_j\), and weighted Cauchy--Schwarz give
\begin{equation}
 \|m_{j,S}(\eta)-m_{j,S}(0)\|_2^2
 \le C E_{j,S}(\eta).
\label{eq:v4v65-mean-energy}
\end{equation}
Choose \(\zeta>0\) so small that
\(2C\zeta\le1/16\).  With the V34 cutoff \(\chi\), set
\begin{equation}
 G_{j,S}(\eta)
 :=
 \chi\!\left(
 \frac{E_{j,S}(\eta)}{\zeta R_Y^2n}
 \right),
 \qquad
 D_{j,S}=1-G_{j,S}.
\label{eq:v4v65-energy-guard}
\end{equation}
The lattice weights and cutoff are fixed independently of the background
coordinate \(B\).

\begin{lemma}[Energy guard alternatives]
\label{lem:v4v65-energy-alternatives}
For sufficiently small \(\epsilon_Y\):
\begin{enumerate}
\item on the support of \(G_{j,S}\),
\begin{equation}
 \|m_{j,S}(\eta)\|_2
 \le\frac12R_Y\sqrt n;
\label{eq:v4v65-good-mean}
\end{equation}
\item on the support of \(D_{j,S}\),
\begin{equation}
 E_{j,S}(\eta)\ge\zeta R_Y^2n;
\label{eq:v4v65-bad-energy}
\end{equation}
\item the guarded conditional tail product and the bad exterior branch
both have norm at most
\begin{equation}
 C^ne^{-cR_Y^2n}
\label{eq:v4v65-two-branch-payment}
\end{equation}
through two background marks.
\end{enumerate}
\end{lemma}

\begin{proof}
If \(G_{j,S}\ne0\), then
\(E_{j,S}<2\zeta R_Y^2n\).  Use
\eqref{eq:v4v65-mean-energy}, the choice of \(\zeta\), and
\(\|m_{j,S}(0)\|_2\le M_0\sqrt n\).  Since \(R_Y\to\infty\), this proves
\eqref{eq:v4v65-good-mean}.  If \(D_{j,S}\ne0\), the argument of \(\chi\)
is greater than one, proving \eqref{eq:v4v65-bad-energy}.

On the good branch, simultaneous tail occurrence forces a centered
conditional fluctuation of size \(cR_Y\sqrt n\).  Apply the conditional
weighted quadratic estimate of V64.  On the bad branch, use
\[
 \boldsymbol1_{\{D_{j,S}\ne0\}}
 \le
 e^{-\delta\zeta R_Y^2n}
 e^{\delta E_{j,S}}
\]
and the weighted quadratic moment
\eqref{eq:v4v64-weighted-quadratic}.  Equations
\eqref{eq:v4v65-Qnorm}--\eqref{eq:v4v65-Qtrace} make its determinant cost
at most \(e^{Cn}\).  The conditional coefficient envelope is
\eqref{eq:v4v64-conditional-envelope}.  Reduce the exponent and absorb
fixed marked polynomials into \(C^n\).
\end{proof}

\begin{lemma}[Summable retained derivatives of the guards]
\label{lem:v4v65-guard-derivatives}
For \(r=1,2\), in the corresponding operator or Hilbert--Schmidt
retained-field norm,
\begin{equation}
 \sup_\eta
 \|D_\eta^rG_{j,S}(\eta)\|
 \le
 \frac{C_r}{(R_Y\sqrt n)^r}
 e^{-c_rm_1\ell_j}.
\label{eq:v4v65-guard-derivatives}
\end{equation}
Consequently
\begin{equation}
 \sum_{j=0}^{J-1}
 \sup_\eta\|D_\eta^rG_{j,S}\|
 \le
 \frac{C_r'}{(R_Y\sqrt n)^r},
 \qquad r=1,2.
\label{eq:v4v65-guard-derivative-sum}
\end{equation}
Also \(D_BG_{j,S}=0\) for every external background mark.
\end{lemma}

\begin{proof}
On the support of a derivative of \(\chi\),
\(E_{j,S}\asymp R_Y^2n\).  Since
\[
 \nabla_\eta E_{j,S}=2Q_{j,S}\eta,
 \qquad
 \|Q_{j,S}\eta\|^2
 \le\|Q_{j,S}\|E_{j,S},
\]
the first derivative is bounded by
\(C\|Q_{j,S}\|^{1/2}/(R_Y\sqrt n)\).  The second derivative contains
\(Q_{j,S}/(R_Y^2n)\) and the square of the first term.  Insert
\eqref{eq:v4v65-Qnorm}.  The geometric growth
\eqref{eq:v4v65-geometric-scales}--%
\eqref{eq:v4v65-buffer-width} makes the resulting series summable.
Background independence is immediate from
\eqref{eq:v4v65-exterior-energy}.
\end{proof}

\subsection{The exact first-success partition}
\label{sec:v4v65-first-success}

Put
\begin{align}
 H_{0,S}
 &:=
 G_{0,S},
\label{eq:v4v65-H0}\\
 H_{j,S}
 &:=
 G_{j,S}\prod_{i=0}^{j-1}D_{i,S},
 \qquad 1\le j\le J.
\label{eq:v4v65-Hj}
\end{align}
At the root, \(E_{J,S}=0\), so \(G_{J,S}=1\).

\begin{theorem}[Pointwise first-success identity]
\label{thm:v4v65-first-success}
For every field configuration,
\begin{equation}
 H_{j,S}\ge0,
 \qquad
 \sum_{j=0}^{J}H_{j,S}=1.
\label{eq:v4v65-first-success}
\end{equation}
Moreover,
\begin{align}
 \sum_{j=0}^{J}|D_\eta H_{j,S}|
 &\le
 2\sum_{i=0}^{J-1}|D_\eta G_{i,S}|,
\label{eq:v4v65-first-derivative-partition}\\
 \sum_{j=0}^{J}|D_\eta^2 H_{j,S}|
 &\le
 C\left[
 \sum_i|D_\eta^2G_{i,S}|
 +
 \left(\sum_i|D_\eta G_{i,S}|\right)^2
 \right].
\label{eq:v4v65-second-derivative-partition}
\end{align}
Thus both derivative sums are independent of hierarchy height by
Lemma~\ref{lem:v4v65-guard-derivatives}.
\end{theorem}

\begin{proof}
The finite product telescope gives
\[
 \sum_{j=0}^{J}
 G_{j,S}\prod_{i<j}D_{i,S}
 =
 1-\prod_{i=0}^{J}D_{i,S}=1,
\]
because \(D_{J,S}=0\).  Positivity follows from \(0\le G,D\le1\).

Differentiate the product formula.  Fix \(i\).  A derivative of \(G_i\)
appears once in \(H_i\), and a derivative of
\(D_i=1-G_i\) appears in later \(H_j\).  The sum of the absolute values of
their remaining nonnegative coefficients is at most
\[
 \prod_{k<i}D_k
 \left[
 1+\sum_{j>i}G_j\prod_{i<k<j}D_k
 \right]
 \le2.
\]
Summing \(i\) proves
\eqref{eq:v4v65-first-derivative-partition}.  Two differentiations have
one second-derivative placement or two first-derivative placements.  Apply
the same coefficient telescope to each placement; this gives
\eqref{eq:v4v65-second-derivative-partition}.
\end{proof}

\begin{firewall}[No V41 contraction is used for block promotion]
\label{fw:v4v65-no-false-sweep}
The identity \eqref{eq:v4v65-first-success} is pointwise on the original
field.  It does not replace an ancestor promotion by a V41 round trip.
V41 and V44 may resolve the actual color sweeps inside one stopped
conditional calculation, at their fixed cost
\((1-q_{\rm H})^{-1}\); that cost is not raised to the number of ancestor
levels.
\end{firewall}

\subsection{Activity allocation without a level factor}
\label{sec:v4v65-activity-allocation}

Let \(T_S\) be a complete tail coefficient with distinct cell set \(S\),
including its fermion ancestry.  Before any conditional expectation, use
\begin{equation}
 T_S=\sum_{j=0}^{J}H_{j,S}T_S.
\label{eq:v4v65-tail-allocation}
\end{equation}

\begin{theorem}[Stopping-level norm has no hierarchy-height loss]
\label{thm:v4v65-no-height-loss}
For every nonnegative measurable envelope \(A\) on the master field space,
possibly containing the other factors of a moment or cumulant partition,
\begin{equation}
 \sum_j\mathbb E\!\left[H_{j,S}|T_S|A\right]
 =
 \mathbb E\!\left[|T_S|A\right].
\label{eq:v4v65-weighted-L1-partition}
\end{equation}
In particular, for the value component,
\begin{align}
 \sum_j\mathbb E\!\left[H_{j,S}|T_S|\right]
 &=
 \mathbb E|T_S|,
\label{eq:v4v65-L1-partition}\\
 \sum_j\|H_{j,S}T_S\|_2^2
 &\le
 \|T_S\|_2^2.
\label{eq:v4v65-L2-partition}
\end{align}
Through two background marks, the same identities hold with
\(D_B^rT_S\), because \(D_BH_{j,S}=0\).  Through two retained-field
derivatives, the additional guard terms are bounded by
\eqref{eq:v4v65-first-derivative-partition}--%
\eqref{eq:v4v65-second-derivative-partition} and retain
\begin{equation}
 C_r^ne^{-c_rR_Y^2n}.
\label{eq:v4v65-marked-stopping-norm}
\end{equation}
All constants are independent of \(J\) and volume.
\end{theorem}

\begin{proof}
Equation \eqref{eq:v4v65-weighted-L1-partition} follows pointwise from
\(H_j\ge0\) and \(\sum_jH_j=1\), before any absolute moment is
estimated.  Taking \(A=1\) gives
\eqref{eq:v4v65-L1-partition}.  Since \(H_j^2\le H_j\),
\[
 \sum_j\|H_jT_S\|_2^2
 \le
 \mathbb E\!\left[
 |T_S|^2\sum_jH_j
 \right]
 =
 \|T_S\|_2^2.
\]
A background derivative hits only \(T_S\).  A retained derivative of
\eqref{eq:v4v65-tail-allocation} has the exact product-rule placements.
Use Theorem~\ref{thm:v4v65-first-success} for the guard factors and the
V64 simultaneous tail envelope for \(T_S\) and its derivatives.  The
extra powers \((R_Y\sqrt n)^{-1}\) are harmless.  This proves
\eqref{eq:v4v65-marked-stopping-norm}.
\end{proof}

\begin{firewall}[Sum the positive stopping weights before taking a norm]
\label{fw:v4v65-sum-before-norm}
The proof uses
\(\sum_jH_{j,S}=1\) inside the moment.  Bounding every
\(\|H_{j,S}T_S\|_2\) separately and then summing would lose a factor
depending on the number of levels.  The stopping label is retained in a
Hilbert direct sum for covariance estimates and in the nonnegative
partition for moment estimates, just as the YM V82 output label must be
retained before tensorization.
\end{firewall}

\begin{corollary}[Compatibility with the tilted cumulant]
\label{cor:v4v65-cumulant-stopping}
Fix a moment block in the partition formula for the V62 tilted cumulant.
After all colliding tail sets in that block have been united, summing the
stopping label of the union packet introduces no constant and no factor
depending on the hierarchy height.
\end{corollary}

\begin{proof}
Write the absolute value of the remaining factors in the moment block as
the envelope \(A\) in
\eqref{eq:v4v65-weighted-L1-partition}.  The equality sums the stopping
label before the factorial partition coefficient is applied.  For
retained derivatives, use the product-rule placements and
\eqref{eq:v4v65-first-derivative-partition}--%
\eqref{eq:v4v65-second-derivative-partition}; their constants are
independent of the number of levels.
\end{proof}

\subsection{Regularization at the stopping level}
\label{sec:v4v65-regularization}

On the \(j\)-th summand, integrate the complete observable
\[
 T_{j,S}:=
 T_S\prod_{i<j}D_{i,S}
\]
in the level-\(j\) conditional law, keeping \(G_{j,S}\) attached.
The earlier failure factors are bounded observables, not changes of the
positive parent.

\begin{theorem}[Every stopped branch becomes a rare Lipschitz load]
\label{thm:v4v65-stopped-regularization}
For every \(j\), the conditional transform of
\(H_{j,S}T_S=G_{j,S}T_{j,S}\) is a centered, quasi-local Lipschitz load.
Its value and retained-gradient scale, including two background marks, are
bounded by
\begin{equation}
 C^ne^{-cR_Y^2n}
\label{eq:v4v65-stopped-Lipschitz}
\end{equation}
times the incoming fermion and already acquired response kernels.
The stopping-level family satisfies the same bound in the combined norm of
Theorem~\ref{thm:v4v65-no-height-loss}.
\end{theorem}

\begin{proof}
On the support of \(G_{j,S}\),
\eqref{eq:v4v65-good-mean} holds.  The factors \(D_{i,S}\) are bounded by
one and remain inside the observable.  The conditional simultaneous-tail
argument of V64 therefore pays every cell of \(S\).  Differentiate the
conditional expectation.  The explicit derivative and the boundary-score
covariance are exactly
\eqref{eq:v4v63-transform-derivative}; the conditional Witten inverse gives
the quasi-local response.  Derivatives of the stopping factors are
controlled by
Theorem~\ref{thm:v4v65-first-success}.  Background marks hit no guard and
have the V54--V56 placements.  Finally use the assembled stopping norm
\eqref{eq:v4v65-L1-partition}--%
\eqref{eq:v4v65-marked-stopping-norm}.
\end{proof}

\subsection{Mergers and termination}
\label{sec:v4v65-mergers}

Initialize one packet for each connected fine tail coefficient after V62
same-cell consolidation.  Before expanding its stopping levels, apply the
following deterministic rule:
\begin{enumerate}
\item if two packet interaction collars meet, replace them by one packet at
their least common ancestor and replace their disjoint tail sets by their
union;
\item recompute one energy guard and one first-success partition for the
complete union packet;
\item otherwise retain the packet's first-success partition and send every
stopped branch to the V53/V63 collision compiler.
\end{enumerate}
Ties are resolved lexicographically.

\begin{theorem}[Finite first-success forest]
\label{thm:v4v65-forest}
The merger/promotion process terminates, preserves an exact disjoint
partition of the original physical tail-cell labels, and has no terminal
bad branch.  Every final packet carries one first-success partition whose
weights sum to one and one factor \(\tau\) per distinct cell.
\end{theorem}

\begin{proof}
Use the V11 well-founded measure: the number of packets followed by the sum
of their distances to hierarchy roots.  A merger lowers the first entry;
promotion along a first-success chain lowers the second.  Tail sets are
disjoint by V62 and are replaced only by disjoint union.

At a root, \(\Sigma_J=\varnothing\), hence
\(E_{J,S}=0\), \(G_{J,S}=1\), and
\(D_{J,S}=0\).  Thus
\eqref{eq:v4v65-first-success} has no terminal failure term.  Recomputing
the union partition before taking norms discards all provisional child
partitions, exactly as V11 discards provisional child repairs after a
collision.
\end{proof}

\begin{firewall}[A merger discards provisional stopping payments]
\label{fw:v4v65-provisional}
If two packet collars meet, their separately written stopping partitions
are not multiplied.  Their product would create incompatible repeated
boundary payments.  The complete collision coefficient is first assembled,
the physical tail sets are united, and the identity
\eqref{eq:v4v65-first-success} is applied once to that union.
\end{firewall}

\subsection{The tail collision compiler}
\label{sec:v4v65-compiler}

At a stopped collision, all tail descendants are Lipschitz loads by
Theorem~\ref{thm:v4v65-stopped-regularization}; ordinary good descendants
have the V53 load.  If collision occurs inside the selected stopping core
before separate conditional transforms, use the V62 tilted multi-tail
cumulant on the already united set and then apply the one group guard.

\begin{theorem}[All-depth tail-specific interface]
\label{thm:v4v65-tail-interface}
Assume the uniform V38 block geometry, the V53--V56 local good collision
bounds, and the V64 coefficient envelopes.  Then every finite collision
history containing V34 tail cells has a hierarchical star bound with:
\begin{enumerate}
\item one factor \(Ce^{-cR_Y^2}\) per distinct physical tail cell;
\item one actual response or bounded-collar incidence edge at every merger;
\item a \(k!C^k\) coefficient at a \(k\)-child merger; and
\item volume- and hierarchy-uniform first and second background marks.
\end{enumerate}
No conditional \(L^\infty\) tail probability and no stopping-depth factor
is used.
\end{theorem}

\begin{proof}
Fine tail sets have the V34 simultaneous-tail norm after V62
consolidation.  Theorem~\ref{thm:v4v65-no-height-loss} sums all ancestor
choices without changing that norm.  Every stopped branch is a V53
Lipschitz load by
Theorem~\ref{thm:v4v65-stopped-regularization}.  A collision before
separate stopping is assembled on the union set and estimated by V62 and
the group alternative
Lemma~\ref{lem:v4v65-energy-alternatives}.

Remote approach to a collision crosses the actual V38 conditional Witten
kernel.  Packets whose collars already intersect use the bounded local
incidence kernel.  At each merger the union of child trees and the
\(k-1\) new incidences is a tree, as in V52.  The local coefficient is
factorial-exponential by V53 or V62.  The marked statement follows because
the guards are background independent and retained derivatives have the
summable bounds
\eqref{eq:v4v65-guard-derivative-sum}.  V41/V44 resolve actual internal
color sweeps only once at the stopped level.
\end{proof}

\begin{corollary}[The V58 tail correction in the hard convex parent]
\label{cor:v4v65-tail-pressure}
Under Theorem~\ref{thm:v4v65-tail-interface} and the all-stage good-field
hypotheses of Theorem~\ref{thm:v4v57-good-pressure},
\begin{equation}
 \|D_B^j(p_{\rm G+T}-p_{\rm G})\|
 \le
 C_je^{-cR_Y^2},
 \qquad j=0,1,2,
\label{eq:v4v65-tail-pressure}
\end{equation}
and normalized local ratios have the same uniform tail correction.
\end{corollary}

\begin{proof}
Theorem~\ref{thm:v4v65-tail-interface} supplies the tail-specific typed
star estimate required in V58.  Insert it in the V58 rooted fixed point and
use \(R_Y=\epsilon_Y^{-\alpha}\).
\end{proof}

\begin{assessment}[Scope after the stopping forest]
\label{ass:v4v65-scope}
Within the uniformly convex finite-range hard Higgs parent and the stated
uniform block interfaces, V65 removes the tail-specific conditional
supremum, owner-confluence, and stopping-depth gates.  The rare tail is now
controlled by V34 and V59--V65.  The all-stage good-field interface in V57
remains a renormalization-scale hypothesis: V53--V56 prove its physical
local stage, but its uniform preservation under the full Standard Model RG
map has not been proved.
\end{assessment}

\begin{firewall}[The full continuum remains open]
\label{fw:v4v65-continuum-open}
V65 is a finite-cutoff hard-sector theorem.  It does not control broken
Higgs directions, gauge/ghost terminal sectors, the residual determinant
phase, hypercharge ultraviolet behavior, reflection positivity,
Osterwalder--Schrader reconstruction, dynamical Einstein fluctuations, or
cutoff removal.
\end{firewall}

\subsection{V65 result ledger and V66 target}
\label{sec:v4v65-conclusion}

V65 records the exact YM V82 and NS V31 fingerprints.  It replaces a false
per-promotion sweep contraction by a background-independent exterior-energy
guard and an exact first-success partition.  The stopping weights sum to
one before moment estimates, and their first two retained derivatives are
summable across the geometric hierarchy.  Collision mergers discard
provisional child partitions, preserve the disjoint physical tail set, and
terminate at a root where the exterior energy vanishes.  This closes the
tail-specific all-depth interface in the hard convex parent, subject to the
existing all-stage good-carrier hypotheses.

V66 should return to the residual good-determinant phase.  The first task is
to remove its scalar phase by the V43 cocycle and prove that the centered
local factor
\(e^{i\vartheta_{\rm G}^c}-\mathbb E e^{i\vartheta_{\rm G}^c}\)
has a Lipschitz load proportional to
\(\epsilon_Y^{1-\alpha}\), rather than merely modulus one.  That estimate
must retain the exact phase branch and its first two background marks.


\clearpage
\section{V66: the residual good-determinant phase as a centered Lipschitz species}
\label{sec:v4v66}

V65 closes the rare-tail stopping interface in the hard convex parent.  The
next unresolved Yukawa species is the unit-modulus phase of the smooth good
determinant.  V35 proves a small connected phase norm through two external
background marks, but V53 cannot use pointwise smallness alone: its
all-order cumulant estimate requires a Higgs-field Lipschitz load.  This
note supplies that missing mixed norm directly from the trace words.

The proof keeps three operations separate.  A local phase factor is
centered in the positive parent before the V53 generator is used.  A
conditional mean depending on a retained skeleton remains a skeleton
function and is propagated by the V54 score identity.  Only after the V43
normalization functional is applied is a mean recorded as a scalar vacuum
factor.  With this order, the residual phase joins the same star-merger
compiler as the positive good species, while remaining outside the
positive action.

\subsection{A mixed phase norm}
\label{sec:v4v66-mixed-norm}

For a complex function \(F(\phi,B)\), put
\begin{equation}
 {\rm Lip}_\phi(F)
 :=
 \sup_{\phi,B}
 \left(\sum_y|\nabla_{\phi_y}F(\phi,B)|^2\right)^{1/2}.
\label{eq:v4v66-Lipschitz-seminorm}
\end{equation}
For the real V35 connected coefficient
\(\vartheta_{\rm G}^c(S)\), define
\begin{align}
 \Delta_S
 &:=
 \sum_{q=0}^2\rho^q
 \sup_{\|v_i\|\le1}
 \|D_B^q\vartheta_{\rm G}^c(S)[v_1,\ldots,v_q]\|_\infty,
\label{eq:v4v66-DeltaS}\\
 \mathcal L_S
 &:=
 \sum_{q=0}^2\rho^q
 \sup_{\|v_i\|\le1}
 {\rm Lip}_\phi
 \bigl(D_B^q\vartheta_{\rm G}^c(S)[v_1,\ldots,v_q]\bigr).
\label{eq:v4v66-LS}
\end{align}
The directions in \eqref{eq:v4v66-DeltaS}--\eqref{eq:v4v66-LS}
are the same admissible local background directions as in V33--V35.

\begin{theorem}[Mixed field/background connected phase norm]
\label{thm:v4v66-mixed-phase-norm}
Under the V31--V35 hard-gap and coefficient hypotheses, for every
\(0<a''<a'<a_0\) and \(0<h''<h\), there are
\(C<\infty\) and \(\epsilon_*>0\), independent of volume, such that
\begin{align}
 \sup_x\sum_{S\ni x}
 e^{h''|S|+a''\mathfrak t(S)}\Delta_S
 &\le C\epsilon_Y^{1-\alpha},
\label{eq:v4v66-phase-value-norm}\\
 \sup_x\sum_{S\ni x}
 e^{h''|S|+a''\mathfrak t(S)}\mathcal L_S
 &\le C\epsilon_Y.
\label{eq:v4v66-phase-Lipschitz-norm}
\end{align}
In particular the second right side is at most
\(C\epsilon_Y^{1-\alpha}\) for \(0<\epsilon_Y<1\).
\end{theorem}

\begin{proof}
The value estimate is Theorem~\ref{thm:v4v35-phase-norm}, with a harmless
reduction of the weights.  For the Lipschitz estimate, differentiate the
absolutely convergent cyclic-word series
\eqref{eq:v4v35-phase-series} before the V16 Mobius extraction.  In every
resulting word exactly one occurrence is replaced by
\(D_\phi V_{{\rm G},x}\).  Lemma~\ref{lem:v4v34-GT-derivatives} gives
\[
 \|D_\phi V_{{\rm G},x}\|=O(\epsilon_Y),
 \qquad
 \|CV_{{\rm G},x}\|_{a_0}=O(\epsilon_Y^{1-\alpha}).
\]
Thus a word of length \(n\) with one field mark has majorant
\[
 C\epsilon_Y\,n\zeta_Y^{\,n-1},
 \qquad
 \zeta_Y=O(\epsilon_Y^{1-\alpha})<\tfrac12.
\]
The factor \(1/n\) in the logarithm cancels the choice of the marked
cyclic occurrence.  Summation over \(n\) is therefore bounded by
\(C\epsilon_Y/(1-\zeta_Y)\).

The marked insertion is supported in one Yukawa cell.  Bounded cell overlap
turns its squared local direction norm into the Euclidean norm in
\eqref{eq:v4v66-Lipschitz-seminorm}.  The remaining covariance and vertex
kernels have the V33 exponentially weighted row and column bounds.  The
connected Mobius extraction retains a spanning tree of the visited cells;
rooted leaf summation gives the weights in
\eqref{eq:v4v66-phase-Lipschitz-norm}, after reducing \(a'\) and \(h\) to
absorb the choice of the field-marked cell.

One or two \(B\)-derivatives are placed on actual covariance or vertex
factors exactly as in the marked-minor proof of V33 and the proof of
Theorem~\ref{thm:v4v35-phase-norm}.  The field-marked factor remains in the
word, so every placement retains the leading \(O(\epsilon_Y)\).  The
marked geometric series and rooted-tree sum have the same radius.  This
proves \eqref{eq:v4v66-phase-Lipschitz-norm} through two background marks.
\end{proof}

\begin{firewall}[Pointwise size and Lipschitz load are different currencies]
\label{fw:v4v66-two-phase-currencies}
The V35 value norm pays a unary phase activity.  The V66 Lipschitz norm pays
members of a joint cumulant of order at least two.  The latter comes from
one actual \(D_\phi V_{\rm G}\) insertion and is not inferred from
\(|\vartheta_{\rm G}^c(S)|\).  Both estimates are retained because a
cluster logarithm contains unary and higher-order coefficients.
\end{firewall}

\subsection{Exact local centering and a zero-free one-factor mean}
\label{sec:v4v66-local-centering}

Fix a finite-volume positive promoted parent \(\nu_{\rm G,B}\) and put
\begin{equation}
 \theta_S:=\vartheta_{\rm G}^c(S),
 \qquad
 f_S:=e^{i\theta_S},
 \qquad
 c_S(B):=\mathbb E_{\nu_{\rm G,B}}f_S.
\label{eq:v4v66-local-phase-factor}
\end{equation}
Let \(\delta_S=\|\theta_S\|_\infty\).

\begin{lemma}[Zero-free local mean and normalized centered carrier]
\label{lem:v4v66-local-centering}
If \(\delta_S\le1/2\), then
\begin{equation}
 |c_S-1|\le\delta_S,
 \qquad
 |c_S|\ge1-\delta_S\ge\tfrac12.
\label{eq:v4v66-local-mean-disk}
\end{equation}
The exact factorization
\begin{equation}
 f_S=c_S(1+P_S),
 \qquad
 P_S:=c_S^{-1}(f_S-c_S)
\label{eq:v4v66-centered-factorization}
\end{equation}
has
\begin{align}
 \mathbb E_{\nu_{\rm G,B}}P_S&=0,
\label{eq:v4v66-P-centered}\\
 \|P_S\|_\infty&\le4\delta_S,
\label{eq:v4v66-P-value}\\
 {\rm Lip}_\phi(P_S)&\le2{\rm Lip}_\phi(\theta_S).
\label{eq:v4v66-P-Lipschitz}
\end{align}
The unnormalized centered factor
\(F_S=f_S-c_S\) satisfies the same estimates without the factors two and
four.
\end{lemma}

\begin{proof}
Because \(\theta_S\) is real,
\[
 |f_S-1|=|e^{i\theta_S}-1|\le|\theta_S|.
\]
Taking expectation proves the first inequality in
\eqref{eq:v4v66-local-mean-disk}; the second is the reverse triangle
inequality.  Equation \eqref{eq:v4v66-centered-factorization} is
algebraic, and \eqref{eq:v4v66-P-centered} follows from the definition of
\(c_S\).  Moreover,
\[
 |f_S-c_S|
 \le |f_S-1|+|c_S-1|
 \le2\delta_S,
\]
which proves \eqref{eq:v4v66-P-value}.  Since \(c_S\) is scalar in the
Higgs integration variable,
\[
 \nabla_\phi P_S
 =c_S^{-1}ie^{i\theta_S}\nabla_\phi\theta_S.
\]
Use \(|c_S|^{-1}\le2\) to obtain
\eqref{eq:v4v66-P-Lipschitz}.
\end{proof}

\begin{corollary}[Rooted norm of the raw and centered phase leaves]
\label{cor:v4v66-phase-leaves}
After reducing \(\epsilon_*\), the raw family \(z_S=f_S-1\), through two
explicit background marks, has rooted value norm at most
\(C\epsilon_Y^{1-\alpha}\) and rooted Higgs Lipschitz norm at most
\(C\epsilon_Y\).  The unmarked centered families
\(F_S=f_S-c_S\) and \(P_S\) obey the same bounds, up to a fixed constant.
Their marked bounds are proved in
Theorem~\ref{thm:v4v66-two-mark-mean}.
\end{corollary}

\begin{proof}
Differentiate \(e^{i\theta_S}\) exactly:
\begin{align*}
 D_Bf_S[v]&=if_SD_B\theta_S[v],\\
 D_B^2f_S[v,w]
 &=if_SD_B^2\theta_S[v,w]
   -f_SD_B\theta_S[v]D_B\theta_S[w].
\end{align*}
Use Theorem~\ref{thm:v4v66-mixed-phase-norm}.  The quadratic term is
bounded by the convolution algebra of the rooted norm.  The unmarked
claims for \(F_S\) and \(P_S\) follow from
\eqref{eq:v4v66-local-mean-disk}--\eqref{eq:v4v66-P-Lipschitz}.
\end{proof}

\subsection{Marks of the local mean}
\label{sec:v4v66-mean-marks}

For clarity, write the promoted positive law as
\(d\nu_B=Z_B^{-1}e^{-A_B}d\phi\).  For local background directions
\(v,w\), define centered scores
\begin{equation}
 s_v:=D_BA_B[v]-\mathbb E_{\nu_B}D_BA_B[v],
 \qquad
 s_{vw}:=D_B^2A_B[v,w]
 -\mathbb E_{\nu_B}D_B^2A_B[v,w].
\label{eq:v4v66-scores}
\end{equation}

\begin{lemma}[Exact first and second marks of a phase mean]
\label{lem:v4v66-mean-marks}
For any twice differentiable local \(f=f(\phi,B)\),
\begin{align}
 D_B\mathbb E_{\nu_B}f[v]
 &=
 \mathbb E_{\nu_B}D_Bf[v]
 -\kappa_{\nu_B}(f,s_v),
\label{eq:v4v66-first-mean-mark}\\
 D_B^2\mathbb E_{\nu_B}f[v,w]
 &=
 \mathbb E_{\nu_B}D_B^2f[v,w]
 -\kappa_{\nu_B}(D_Bf[v],s_w)
 -\kappa_{\nu_B}(D_Bf[w],s_v)
 \notag\\
 &\quad
 -\kappa_{\nu_B}(f,s_{vw})
 +\kappa_{\nu_B}(f,s_v,s_w).
\label{eq:v4v66-second-mean-mark}
\end{align}
For \(f=f_S\), every covariance or higher cumulant on the right may have
\(f_S\) replaced by \(f_S-1\).  Hence the phase leaf, rather than the
constant one, pays every term.
\end{lemma}

\begin{proof}
Normalized differentiation gives
\[
 D_B\mathbb E_{\nu_B}f[v]
 =\mathbb E_{\nu_B}D_Bf[v]
  -\mathbb E_{\nu_B}(f s_v),
\]
which is \eqref{eq:v4v66-first-mean-mark}.  Differentiate once more.  The
derivative of the centered score is
\[
 D_Bs_v[w]
 =s_{vw}+\kappa_{\nu_B}(D_BA_B[v],s_w).
\]
The scalar last term cancels the disconnected part of
\(\mathbb E(f s_vs_w)\).  Collecting the two explicit first derivatives,
the genuine second score, and the two-first-score cumulant gives
\eqref{eq:v4v66-second-mean-mark}.  Cumulants containing a constant vanish
as soon as their order is at least two.
\end{proof}

\begin{theorem}[Two-mark phase-mean bound]
\label{thm:v4v66-two-mark-mean}
Assume the V54--V56 localized score interfaces for the positive promoted
parent.  Then \(c_S-1\), its first two background marks, and the first two
marks of \(P_S\) obey the same rooted spatial norm
\begin{equation}
 C\epsilon_Y^{1-\alpha}
\label{eq:v4v66-marked-phase-mean}
\end{equation}
as the phase leaves.  Every marked term contains one phase leaf and one
actual marked star tree; no derivative of a scalar estimate is used in
place of \eqref{eq:v4v66-second-mean-mark}.
\end{theorem}

\begin{proof}
For the value, use \eqref{eq:v4v66-local-mean-disk}.  Insert the explicit
phase derivatives from Corollary~\ref{cor:v4v66-phase-leaves} into
\eqref{eq:v4v66-first-mean-mark}--%
\eqref{eq:v4v66-second-mean-mark}.  The one-score terms use the V54 marked
generator.  The last two terms use respectively the V56 genuine
second-score junction and the V55 two-first-score tree.  Since a constant
vanishes from each cumulant, \(f_S-1\) supplies the small phase value or
Lipschitz load.  Flatten each marked star with the phase leaf tree.
The V57 marked rooted series sums the resulting polynomial mark choices.
Finally differentiate \(P_S=c_S^{-1}(f_S-c_S)\); the disk
\(|c_S|\ge1/2\) bounds all inverse factors, and the exact first and second
quotient placements are controlled by the estimates just proved.
\end{proof}

\begin{firewall}[A conditional mean is not a V43 scalar]
\label{fw:v4v66-conditional-not-scalar}
If \(f_S\) is integrated in a conditional component with retained skeleton
\(\eta\), then
\(c_{S,D}(\eta)=\mathbb E_{\nu_D^\eta}f_S\) is a skeleton function.  Its
derivative is the explicit derivative minus the boundary-score covariance,
as in \eqref{eq:v4v63-transform-derivative}.  It remains a carrier until
the retained variables are integrated.  Only a number obtained after the
V43 normalization functional is applied may enter the scalar cocycle.
\end{firewall}

\subsection{Conditional recentering and collision stability}
\label{sec:v4v66-conditional-collisions}

At a conditional component \(D\), replace every incoming phase child
\(f_i\) by
\begin{equation}
 \widehat f_i
 :=f_i-\mathbb E_{\nu_D^\eta}f_i.
\label{eq:v4v66-conditional-centering}
\end{equation}
The conditional mean is sent to the unary ancestry and
\(\widehat f_i\) is sent to a collision of order at least two.

\begin{lemma}[Recentring changes no genuine collision]
\label{lem:v4v66-cumulant-centering}
For \(k\ge2\),
\begin{equation}
 \kappa_{\nu_D^\eta}(f_1,\ldots,f_k)
 =
 \kappa_{\nu_D^\eta}(\widehat f_1,\ldots,\widehat f_k).
\label{eq:v4v66-cumulant-centering}
\end{equation}
The centered children have the same Higgs Lipschitz seminorms as the
original children.  Retained-skeleton derivatives of their conditional
means have the V54 score ancestry and therefore retain an actual response
edge.
\end{lemma}

\begin{proof}
A joint cumulant of order at least two vanishes if any argument is constant;
multilinearity proves \eqref{eq:v4v66-cumulant-centering}.  Subtracting a
constant in the integrated field does not change its Higgs gradient.
Normalized conditional differentiation gives the V54 explicit-child term
and centered boundary-score covariance, so the conditional mean is not
estimated by an unmarked scalar bound.
\end{proof}

\begin{theorem}[All-depth phase-specific star interface]
\label{thm:v4v66-phase-interface}
Assume the uniform V38 block geometry and that the V54--V56 localized score
interfaces hold at every conditional stage.  Then every finite collision
history made from the V35 phase leaves has:
\begin{enumerate}
\item leaf activity
      \(\varepsilon_{\rm ph}\le C\epsilon_Y^{1-\alpha}\);
\item one actual response or bounded-collar incidence edge at every merger;
\item local \(k\)-child coefficient at most \(k!C^k\); and
\item volume-uniform first and second background-marked bounds.
\end{enumerate}
The constants are those of a complex Lipschitz load under the positive
parent; the phase is never inserted into the parent action.
\end{theorem}

\begin{proof}
Unary leaves use \eqref{eq:v4v66-phase-value-norm}.  At a merger of order
\(k\ge2\), apply Lemma~\ref{lem:v4v66-cumulant-centering} and
Theorem~\ref{thm:v4v53-Lipschitz-cumulant}; its proof already allows
complex functions.  The individual loads are supplied by
\eqref{eq:v4v66-phase-Lipschitz-norm} and contain the pre-collision
responses.  Theorem~\ref{thm:v4v53-multiway-tree} turns their product into
a star on the children.  V54 handles one mark and V55--V56 handle the two
exhaustive second-mark channels.  The conditional-mean branch is retained
by the score identity in
Lemma~\ref{lem:v4v66-cumulant-centering}.  Induction over the well-founded
collision history and V52 flattening prove the statement.
\end{proof}

\subsection{Residual phase logarithm under the all-stage interface}
\label{sec:v4v66-phase-logarithm}

For a finite family of V35 connected supports in \(\Lambda\), introduce
bookkeeping variables and write
\begin{equation}
 \mathcal Z_\Lambda^{\rm ph}(z)
 :=
 \mathbb E_{\nu_{\rm G,B}}
 \prod_{S\Subset\Lambda}
 \left[1+z_S\bigl(e^{i\theta_S}-1\bigr)\right].
\label{eq:v4v66-phase-Z}
\end{equation}
At \(z_S=1\), V35 gives exactly
\begin{equation}
 \mathcal Z_\Lambda^{\rm ph}(1)
 =
 \mathbb E_{\nu_{\rm G,B}}e^{i\vartheta_{\rm G}}.
\label{eq:v4v66-phase-Z-exact}
\end{equation}
The coefficient of distinct labels \(S_1,\ldots,S_k\) in
\(\log\mathcal Z_\Lambda^{\rm ph}(z)\) is their joint cumulant.  For
\(k=1\) it is the scalar mean; for \(k\ge2\) it is invariant under the
conditional centerings above.

\begin{theorem}[Conditional closure of the correlated-parent phase target]
\label{thm:v4v66-phase-pressure}
Under Theorem~\ref{thm:v4v66-phase-interface} and the all-stage hypotheses
of Theorem~\ref{thm:v4v57-good-pressure}, the residual phase coefficients
satisfy
\begin{equation}
 \sup_x
 \sum_{X\ni x}e^{a|X|}
 \sum_{q=0}^2\rho^q
 \|D_B^q\zeta_{\rm ph}(X)\|
 \le C_a\epsilon_Y^{1-\alpha}.
\label{eq:v4v66-phase-KP}
\end{equation}
Consequently the phase expectation in
\eqref{eq:v4v66-phase-Z-exact} is nonzero on the branch continued from
\(\epsilon_Y=0\), its phase pressure has a translation-covariant van Hove
limit, and normalized local phase ratios converge through two background
marks.
\end{theorem}

\begin{proof}
Apply the V57 rooted star-merger majorant with leaf parameter
\(\varepsilon_{\rm ph}\).  Its proof uses absolute coefficients, so it
applies unchanged to complex phase leaves.  The unordered-child factor
cancels the local \(k!\), and the marked branching factors are polynomial.
This proves \eqref{eq:v4v66-phase-KP}.  Absolute convergence defines the
logarithmic branch connected to the zero-coupling value one and hence
excludes zeros there.  Rooted translation-class summation and vanishing
boundary density give the pressure limit.  With a local insertion, all
clusters disconnected from its support cancel between numerator and
denominator; the remaining rooted sum converges.  The V43 cocycle records
the unary scalar means and final block normalizations once, including
their phase, and cancels the same constants from normalized ratios.
\end{proof}

\begin{assessment}[What V66 closes]
\label{ass:v4v66-scope}
V66 supplies the previously missing Higgs Lipschitz norm of every connected
good-determinant phase coefficient.  It proves exact centering, a zero-free
one-factor mean, and recursive complex collision bounds through two marks.
Under the same explicit all-stage block hypothesis already present in V57,
it closes the correlated-parent residual phase pressure and local-ratio
species.  V36 remains the unconditional small-hopping product-parent
closure; V66 extends the conditional arbitrary-hopping route and does not
remove its RG-scale hypothesis.
\end{assessment}

\begin{firewall}[The phase theorem is still a hard-sector statement]
\label{fw:v4v66-hard-scope}
The phase analyzed here is the anomaly-normalized phase of the V34 smooth
good determinant in a fixed hard fermion chart.  V66 does not control a
zero crossing of the hard operator, the tail determinant by a phase
estimate, gauge/ghost terminal sectors, broken Higgs directions,
hypercharge ultraviolet flow, or cutoff removal.
\end{firewall}

\subsection{V66 result ledger and V67 target}
\label{sec:v4v66-conclusion}

V66 strengthens V35 from a pointwise connected phase norm to a mixed
Higgs-Lipschitz/background-marked norm.  The exact local phase factor has a
mean in the disk \(|c_S-1|\le\delta_S\), and its normalized centered part
has Lipschitz load bounded by that of \(\theta_S\).  First and second marks
of the mean have the exact explicit, one-score, genuine-second-score, and
two-first-score placements.  Conditional means remain skeleton carriers;
only final V43 normalizations are scalars.  The complex V53 generator and
the V57 rooted sum then close the residual phase species under the stated
all-stage interface.

V67 should assemble the three now controlled Yukawa descendants---positive
good modulus, residual good phase, and stopped rare tail---in one typed
rooted recursion.  It must prove mixed mergers without multiplying three
separate convergence constants and identify precisely which hard-sector
fermion/Yukawa continuum gates remain after the complete complex
normalization is assembled.


\clearpage
\section{V67: one typed recursion for the complete hard Yukawa determinant}
\label{sec:v4v67}

V57, V65, and V66 control three different descendants of the V34 split.
They cannot be multiplied as three independent cluster estimates.  A
collision may contain phase and tail children simultaneously, and the
positive good modulus has already changed the Higgs law.  This note
assembles the exact hard-chart normalization in the correct parent and
proves one mixed rooted recursion.

The result is a conditional closure theorem for the complete hard Yukawa
determinant at finite lattice spacing and in the thermodynamic limit.  Its
condition is the uniform all-stage block interface already isolated in
V57.  The theorem does not identify that condition with cutoff removal.

\subsection{The exact no-duplication ledger}
\label{sec:v4v67-ledger}

Suppress the V16 quadratic hard-fermion normalization, which remains an
external factor fixed once.  The V34 determinant identity is
\[
 \mathcal D_{\rm Y}
 =
 \mathcal D_{\rm G}\mathcal D_{\rm T}
 =
 |\mathcal D_{\rm G}|\,
 e^{i\vartheta_{\rm G}}\,
 \mathcal D_{\rm T}.
\]
For the positive hard Higgs parent \(\mu_B\), define
\begin{align}
 Z_{\rm G,\Lambda}^{+}(B)
 &:=
 \mathbb E_{\mu_B}|\mathcal D_{\rm G}|,
\label{eq:v4v67-ZG}\\
 d\nu_{\rm G,B}
 &:=
 (Z_{\rm G,\Lambda}^{+})^{-1}
 |\mathcal D_{\rm G}|\,d\mu_B,
\label{eq:v4v67-promoted-parent}\\
 \mathcal R_{\Lambda}^{\rm Y}(B)
 &:=
 \mathbb E_{\nu_{\rm G,B}}
 \left[e^{i\vartheta_{\rm G}}\mathcal D_{\rm T}\right].
\label{eq:v4v67-relative-factor}
\end{align}
Then, exactly,
\begin{equation}
 \mathbb E_{\mu_B}\mathcal D_{\rm Y}
 =
 Z_{\rm G,\Lambda}^{+}\mathcal R_{\Lambda}^{\rm Y}.
\label{eq:v4v67-exact-factorization}
\end{equation}

\begin{lemma}[Every Yukawa factor has one owner]
\label{lem:v4v67-one-owner}
In \eqref{eq:v4v67-exact-factorization}:
\begin{enumerate}
\item every \(g_xV_x\) occurrence belongs to
      \(\mathcal D_{\rm G}\), whose modulus defines
      \(\nu_{\rm G,B}\) and whose phase defines the V35 phase leaves;
\item every \(t_xV_x\) occurrence belongs to the V62 packet
      \(\Theta_x\) inside \(\mathcal D_{\rm T}\); and
\item the covariance in every tail packet is the updated covariance
      \(C_{\rm G}=(A_0+V_{\rm G})^{-1}\).
\end{enumerate}
No modulus activity is inserted again in
\(\mathcal R_{\Lambda}^{\rm Y}\).
\end{lemma}

\begin{proof}
These statements are the three factors in
\eqref{eq:v4v34-determinant-factorization} followed by
\(\mathcal D_{\rm G}=|\mathcal D_{\rm G}|e^{i\vartheta_{\rm G}}\).
The Neumann expansion of \(C_{\rm G}\) only resolves the covariance within
the tail factor and does not create another determinant.  Normalizing the
modulus gives \eqref{eq:v4v67-promoted-parent}; division then gives
\eqref{eq:v4v67-relative-factor}.
\end{proof}

\begin{firewall}[The promoted modulus is the parent, not a relative leaf]
\label{fw:v4v67-no-double-modulus}
The V57 positive normalization controls
\(Z_{\rm G,\Lambda}^{+}\).  Inside
\(\mathcal R_{\Lambda}^{\rm Y}\), its derivatives reappear only as
conditional scores of the positive law and are controlled by V54--V56.
Adding the V57 good-modulus leaves to the relative phase/tail product would
count \(|\mathcal D_{\rm G}|\) twice.
\end{firewall}

\subsection{The two relative leaf types}
\label{sec:v4v67-leaves}

The V35 identity and V62 cell consolidation give
\begin{align}
 e^{i\vartheta_{\rm G}}
 &=
 \prod_{\varnothing\ne S\Subset\Lambda}
 (1+z_S^{\rm ph}),
 \qquad
 z_S^{\rm ph}:=e^{i\vartheta_{\rm G}^c(S)}-1,
\label{eq:v4v67-phase-product}\\
 \prod_xe^{-\mathcal V_x^{\rm T}}
 &=
 \prod_x(1+\Theta_x).
\label{eq:v4v67-tail-product}
\end{align}
The Berezin contractions of \eqref{eq:v4v67-tail-product} are always made
with \(C_{\rm G}\).  A tail coefficient retains the set of distinct
physical cells selected from the \(\Theta_x\)'s.

For a relative leaf \(I\), write
\[
 {\rm typ}(I)\in\{{\rm ph},{\rm T}\},
 \qquad
 \lambda(I)=
 \begin{cases}
  \varepsilon_{\rm ph}(I),&{\rm typ}(I)={\rm ph},\\
  \tau^{|S(I)|}a(I),&{\rm typ}(I)={\rm T},
 \end{cases}
\]
where
\begin{equation}
 \varepsilon_{\rm ph}
 \le C\epsilon_Y^{1-\alpha},
 \qquad
 \tau:=K_{\rm T}e^{-c_{\rm T}R_Y^2}.
\label{eq:v4v67-two-small-parameters}
\end{equation}
The first bound is V66.  The second is the V62/V65 projective activity.

\begin{lemma}[Exact mixed-family expansion]
\label{lem:v4v67-mixed-expansion}
At every finite cutoff,
\(\mathcal R_{\Lambda}^{\rm Y}\) is the finite sum over families of phase
supports and distinct tail cells obtained by multiplying
\eqref{eq:v4v67-phase-product} and
\eqref{eq:v4v67-tail-product}, then performing the one Berezin expectation
with covariance \(C_{\rm G}\), and finally taking the
\(\nu_{\rm G,B}\) expectation.  Every selected phase support and physical
tail cell appears exactly once.
\end{lemma}

\begin{proof}
Both products are exact finite-dimensional identities.  Distributivity
gives the selected mixed families.  Lemma~\ref{lem:v4v62-cell-consolidation}
makes the tail indices a set.  The phase indices label the distinct
connected coefficients in the exponential of
\(\sum_S i\vartheta_{\rm G}^c(S)\).  The Berezin functional is linear, so it
may be applied after the finite expansion without changing any coefficient.
No new label is created by a contraction.
\end{proof}

\subsection{A complete classification of mixed mergers}
\label{sec:v4v67-mixed-mergers}

Consider a collision of \(k\) child histories in one conditional component.
There are three mutually exclusive analytic cases:
\begin{enumerate}
\item all children are bounded Lipschitz descendants;
\item at least one tail child has not yet crossed its V65 stopping level,
      and the unregularized tail collars remain separate until stopping; or
\item unregularized tail collars meet before separate stopping and are
      first consolidated into their disjoint union.
\end{enumerate}
The second case may include stopped tail children among the other
Lipschitz loads.  The third case uses one union guard and the V62 finite
moment--cumulant polynomial.

\begin{theorem}[Typed mixed-merger bound]
\label{thm:v4v67-mixed-merger}
Assume the V53--V56 score interfaces, the V65 tail interface, and the V66
phase interface at the current conditional stage.  For children
\(I_1,\ldots,I_k\), the mixed merger has a bound
\begin{equation}
 k!C_*^k
 \left(\prod_{i=1}^k\lambda(I_i)\right)
 \sum_{r=1}^k
 \prod_{i\ne r}w(I_i,I_r),
\label{eq:v4v67-mixed-star}
\end{equation}
where \(w\) is an actual Witten-response or bounded-collar incidence
kernel with a uniform row sum.  Tail weights in the product contain exactly
one factor \(\tau\) per distinct physical tail cell.  The same statement
holds through two background marks, with the exact V54--V56 placements.
\end{theorem}

\begin{proof}
If all children are Lipschitz descendants, recenter them in the conditional
positive law.  This changes no cumulant of order \(k\ge2\).  Apply the
complex V53 generator and its star consequence
Theorem~\ref{thm:v4v53-multiway-tree}.  This covers pure phase collisions,
stopped tail collisions, and every mixture of the two.

Suppose an unregularized tail child is present.  Keep all bounded centered
children, including phase children, in the V62 good Lipschitz tilt.  That
tilt is zero-free by V53; its proof does not require the loads to be real.
Keep every tail child in the finite moment--cumulant formula
\eqref{eq:v4v62-tilted-tail-cumulant}.  The first-success forest of V65
supplies the joint response interface that remained open in V62, without a
stopping-depth factor.  Theorem~\ref{thm:v4v62-joint-response-interface}
then gives \eqref{eq:v4v67-mixed-star}.

If tail collars meet before separate stopping, V65 discards the provisional
child guards, unites their disjoint V62 cell sets, and installs one group
guard at the least common ancestor.  The V62 product norm pays
\(\tau^{|\dot\cup_iS(I_i)|}=\prod_i\tau^{|S(I_i)|}\).
The group alternative and the moment-level stopping identity
\eqref{eq:v4v65-weighted-L1-partition} allow the stop label to be summed
inside every partition moment.  Hence the same V62 tilted argument applies.

One mark uses the explicit-child or first-score alternatives of V54.  Two
marks use the exhaustive V55 identity: one genuine second score, two first
scores, one twice-marked child, or two once-marked children.  V56 supplies
the second-score junction.  V65 guards are background independent and V66
phase means have the exact score formulas
\eqref{eq:v4v66-first-mean-mark}--%
\eqref{eq:v4v66-second-mean-mark}.  Thus no mark hits an artificial norm,
and the local coefficient remains factorial-exponential.
\end{proof}

\begin{firewall}[A phase load may enter the good tilt]
\label{fw:v4v67-phase-in-tilt}
The term \emph{good} in the V62 tilt means bounded, centered, and
Higgs-Lipschitz under the positive parent.  It does not mean positive or
real valued.  The V66 phase carrier has precisely this complex Lipschitz
property.  The unbounded tail carrier still never enters the exponential
tilt.
\end{firewall}

\subsection{One multitype rooted fixed point}
\label{sec:v4v67-rooted}

Let \(R_{\rm rel}\) be the rooted absolute sum of all relative phase/tail
collision histories, retaining the V62 projective distinct-tail-set weight.
Let
\begin{equation}
 \varepsilon_{\rm rel}
 :=
 C_{\rm ph}\epsilon_Y^{1-\alpha}
 +C_{\rm T}\tau.
\label{eq:v4v67-relative-leaf}
\end{equation}
The constants include the rooted sums over elementary phase supports and
tail packets, not separate collision expansions.

\begin{theorem}[Single typed rooted recursion]
\label{thm:v4v67-typed-recursion}
Under Theorem~\ref{thm:v4v67-mixed-merger},
\begin{equation}
 R_{\rm rel}
 \le
 e^a\varepsilon_{\rm rel}
 +
 A_*\sum_{k\ge2}
 kC_*^kW^{k-1}R_{\rm rel}^k.
\label{eq:v4v67-typed-recursion}
\end{equation}
For sufficiently small \(\epsilon_Y\),
\begin{equation}
 R_{\rm rel}\le2e^a\varepsilon_{\rm rel}.
\label{eq:v4v67-relative-small}
\end{equation}
The same convergence radius controls the first two marked sums.
\end{theorem}

\begin{proof}
At a merger, the children form one unordered set even when their types
differ.  The labelled-set expansion supplies one factor \(1/k!\), which
cancels the single \(k!\) in
\eqref{eq:v4v67-mixed-star}.  Choosing a star center gives \(k\); summing
the other anchors gives \(W^{k-1}\).  Each child is summed in the same
quantity \(R_{\rm rel}\).  The phase/tail choice is already included in
the sum
\[
 R_{\rm rel}=R_{\rm ph}+R_{\rm T};
\]
expanding \(R_{\rm rel}^k\) generates every ordered type word with its
correct multinomial coefficient.  An additional factor \(2^k\) must not be
inserted.

Tail sets of siblings are disjoint and their projective weights multiply by
Lemma~\ref{lem:v4v62-projective-union}.  V65 sums stopping labels before the
absolute moment.  Therefore neither tail ownership nor stopping depth adds
a branching factor.  The scalar recursion is now exactly the V57 recursion
with \(\varepsilon\) replaced by \(\varepsilon_{\rm rel}\), so
Theorem~\ref{thm:v4v57-rooted-majorant} proves
\eqref{eq:v4v67-relative-small}.  The marked merger coefficients differ by
a fixed polynomial in \(k\); Lemma~\ref{lem:v4v57-marked-series} preserves
the radius.
\end{proof}

\begin{firewall}[Do not multiply typewise convergence constants]
\label{fw:v4v67-one-recursion}
A pure-phase sum and a pure-tail sum do not control mixed histories by
multiplication.  Equation \eqref{eq:v4v67-typed-recursion} sums the child
type at every leaf inside one fixed point.  This is why a mixed merger is
paid once and why its unordered-child factorial is cancelled once.
\end{firewall}

\subsection{The complete hard Yukawa normalization}
\label{sec:v4v67-hard-normalization}

Let \(\zeta_{\rm rel}(X,S)\) denote the connected relative coefficient with
ordinary block support \(X\) and distinct physical tail set \(S\), allowing
\(S=\varnothing\).

\begin{theorem}[Complete relative phase/tail logarithm]
\label{thm:v4v67-relative-log}
Assume the uniform all-stage interfaces in
Theorem~\ref{thm:v4v57-good-pressure},
Theorem~\ref{thm:v4v65-tail-interface}, and
Theorem~\ref{thm:v4v66-phase-interface}.  Then
\begin{equation}
 \sup_b
 \sum_{\substack{X\ni b\\S\ {\rm finite}}}
 e^{a|X|}\tau^{-|S|}
 \sum_{q=0}^2\rho^q
 \|D_B^q\zeta_{\rm rel}(X,S)\|
 \le
 C_a\left(\epsilon_Y^{1-\alpha}+\tau\right).
\label{eq:v4v67-relative-KP}
\end{equation}
Consequently \(\mathcal R_\Lambda^{\rm Y}\ne0\) on the branch continued
from \(\epsilon_Y=0\), its intensive logarithm has a
translation-covariant van Hove limit, and its normalized local ratios
converge through two background marks.
\end{theorem}

\begin{proof}
The exact mixed-family expansion is
Lemma~\ref{lem:v4v67-mixed-expansion}.  The V49 history identity groups
each connected coefficient by its first collision hierarchy.
Theorem~\ref{thm:v4v67-typed-recursion} sums all such histories and gives
\eqref{eq:v4v67-relative-KP}.  Absolute convergence defines the logarithm
on the branch connected to the empty relative gas and excludes a zero.
The rooted translation-class proof of
Theorem~\ref{thm:v4v57-good-pressure} gives the pressure limit.  For a
local insertion, clusters disconnected from its support cancel from the
ratio and the remaining rooted sum converges.  Marked convergence permits
two termwise derivatives.
\end{proof}

\begin{theorem}[Conditional hard-Yukawa thermodynamic closure]
\label{thm:v4v67-hard-Yukawa-closure}
Under the hypotheses of
Theorem~\ref{thm:v4v57-good-pressure} and
Theorem~\ref{thm:v4v67-relative-log}, the finite-volume hard-chart Yukawa
normalization is nonzero and
\begin{equation}
 p_{\rm Y}^{\rm hard}
 =
 p_{\rm G}^{+}+p_{\rm rel}^{\rm Y}
\label{eq:v4v67-hard-pressure}
\end{equation}
exists along translation-covariant van Hove sequences, after restoring the
single V16 quadratic normalization.  Hard-chart normalized local ratios
and their first two background derivatives have thermodynamic limits.  All
V43 scalar modulus and phase constants cancel exactly from those ratios.
\end{theorem}

\begin{proof}
The positive factor \(Z_{\rm G,\Lambda}^{+}\) and its pressure/local ratios
are controlled by Theorem~\ref{thm:v4v57-good-pressure}; it is strictly
positive.  The relative factor is nonzero and has the asserted limits by
Theorem~\ref{thm:v4v67-relative-log}.  Multiply the two factors in the
exact identity \eqref{eq:v4v67-exact-factorization} and add their
logarithms.  The V16 normalization is restored once.  Apply
Theorem~\ref{thm:v4v43-local-cancellation} to the common numerator and
denominator cocycle.
\end{proof}

\begin{assessment}[Acquired hard-sector statement]
\label{ass:v4v67-acquired}
V67 closes the algebraic and analytic composition of the smooth good
modulus, residual good phase, and rare tail determinant.  It includes pure
and mixed phase/tail mergers, one rare factor per distinct physical tail
cell, and two genuine background marks.  The closure is conditional on
uniform preservation of the V38 and V53--V56 interfaces at every RG scale.
At the physical local stage those interfaces are theorems; their all-scale
preservation remains the principal hard-Yukawa hypothesis.
\end{assessment}

\begin{firewall}[This is not the full Standard Model--Einstein continuum]
\label{fw:v4v67-not-continuum}
Theorem~\ref{thm:v4v67-hard-Yukawa-closure} is at fixed lattice cutoff in
one hard fermion chart and assumes an all-scale block interface.  It does
not prove passage through hard-operator zero modes, the broken Higgs
vacuum chart, nonabelian gauge/ghost terminal control, the chiral Ward and
Slavnov--Taylor identities, hypercharge ultraviolet completion,
Osterwalder--Schrader positivity, dynamical metric integration, or removal
of the ultraviolet cutoff.
\end{firewall}

\subsection{Gate ledger and V68 target}
\label{sec:v4v67-conclusion}

The complete hard Yukawa determinant now has an exact parent ledger and one
typed recursion.  The good modulus occurs once in the positive parent; the
phase and tail occur once in the relative expectation; and the tail uses
the updated covariance \(C_{\rm G}\).  Every mixed merger falls under
either the complex V53 generator or the V62 finite tail
moment--cumulant formula with the V65 stopping forest.  The single
multitype recursion closes the relative logarithm without multiplying
typewise convergence constants.

The immediate upstream bottleneck is the all-stage block hypothesis.
V68 should isolate the part that follows automatically from strong
convexity: orthogonal marginals and conditional laws preserve the
log-Sobolev curvature constant.  It should then derive the exact
Schur/Witten formula for the Hessian of a marginal action and state the
remaining quantitative contraction needed to keep its exponentially
weighted off-diagonal row sum uniform under repeated blocking.


\clearpage
\section{V68: correction of the promoted-parent mixing and a finite-range master-parent restart}
\label{sec:v4v68}

The proposed V68 target in V67 would have repeated work: V38 already proves
that strong convexity survives marginalization and gives the exact
effective Hessian formula
\eqref{eq:v4v38-effective-Hessian}.  More seriously, checking that target
against the V65 firewall reveals a parent mismatch in V67.  The global
identity
\[
 \mathbb E_{\mu_B}\mathcal D_{\rm Y}
 =
 Z_{\rm G}^{+}\,
 \mathbb E_{\nu_{\rm G,B}}
 [e^{i\vartheta_{\rm G}}\mathcal D_{\rm T}]
\]
is algebraically correct.  But V65 obtains exact conditional products by
using the finite-range parent \(\mu_B\) and leaving every determinant factor
external.  Its theorem cannot be inserted into a collision calculation
whose parent is the non-finite-range promoted law \(\nu_{\rm G,B}\).

This note corrects that error.  The good modulus, good phase, and rare tail
are kept as three external typed species over the one finite-range master
parent.  A tower/restart identity shows that every higher collision may be
recomputed directly in a conditional law of that same parent.  Thus V38
curvature and Witten constants are reused rather than propagated through
successive effective marginal actions.

\subsection{Explicit correction of V67}
\label{sec:v4v68-correction}

\begin{assessment}[Correction to the status of the V67 conclusions]
\label{corr:v4v68-v67}
Equations
\eqref{eq:v4v67-ZG}--\eqref{eq:v4v67-exact-factorization} and
Lemma~\ref{lem:v4v67-one-owner} remain exact algebraic statements.
The analytic use of the V65 interface in
Theorem~\ref{thm:v4v67-mixed-merger} has the wrong parent.  Consequently
Theorems~\ref{thm:v4v67-typed-recursion},
\ref{thm:v4v67-relative-log}, and
\ref{thm:v4v67-hard-Yukawa-closure} are withdrawn as proofs of the stated
conclusions.  They are superseded by the master-parent theorems below.
\end{assessment}

\begin{proof}
The promoted action is
\[
 S_B-\log|\mathcal D_{\rm G}|.
\]
V34 proves that its Hessian is strongly positive and exponentially local,
but \(\log|\mathcal D_{\rm G}|\) has nonzero kernels at arbitrary
separation.  Therefore separated interiors do not have the exact
conditional product
\eqref{eq:v4v38-conditional-product}.  V65 invokes that product in its
stopping forest and explicitly fixes \(\mu_B\) as parent in
Firewall~\ref{fw:v4v65-parent}.  Replacing \(\mu_B\) by
\(\nu_{\rm G,B}\) invalidates that step.  The global change-of-measure
identity itself uses no conditional factorization and remains true.
\end{proof}

\begin{firewall}[An algebraic change of measure is not an analytic transfer]
\label{fw:v4v68-change-measure}
A factor may either be promoted into the parent or retained as an activity.
After promotion, estimates requiring exact finite-range conditional
independence must be reproved for the new law.  One cannot use the
convenient global representation from one choice and a conditional theorem
proved only in the other choice.
\end{firewall}

\subsection{All three determinant descendants over the master parent}
\label{sec:v4v68-three-species}

Write
\[
 \ell_{\rm G}
 =
 \operatorname{Tr}\log(1+CV_{\rm G})
 =
 u_{\rm G}+i\vartheta_{\rm G}.
\]
Apply the V16 connected-word extraction separately to its real and
imaginary parts:
\[
 u_{\rm G}=\sum_{\varnothing\ne S}u_{\rm G}^c(S),
 \qquad
 \vartheta_{\rm G}=\sum_{\varnothing\ne S}\vartheta_{\rm G}^c(S).
\]
Define
\begin{equation}
 z_S^{\rm M}:=e^{u_{\rm G}^c(S)}-1,
 \qquad
 z_S^{\rm ph}:=e^{i\vartheta_{\rm G}^c(S)}-1.
\label{eq:v4v68-modulus-phase-leaves}
\end{equation}
Then, pointwise,
\begin{equation}
 \mathcal D_{\rm G}
 =
 \prod_S(1+z_S^{\rm M})
 \prod_S(1+z_S^{\rm ph}).
\label{eq:v4v68-DG-product}
\end{equation}
The tail determinant retains the exact V62 packets \(\Theta_x\) and the
updated covariance \(C_{\rm G}\).

\begin{lemma}[Mixed norm of the modulus leaves]
\label{lem:v4v68-modulus-leaves}
The modulus coefficients and activities in
\eqref{eq:v4v68-modulus-phase-leaves} obey, through two background marks,
the same rooted value bound
\[
 C\epsilon_Y^{1-\alpha}
\]
and the same rooted Higgs Lipschitz bound
\[
 C\epsilon_Y
\]
as the V66 phase leaves.  Their field derivatives through order three have
the weighted kernels proved in V34 and V56.
\end{lemma}

\begin{proof}
The real and imaginary parts of every trace word have the same absolute
majorant.  Replace the imaginary part by the real part in the proofs of
Theorems~\ref{thm:v4v35-phase-norm} and
\ref{thm:v4v66-mixed-phase-norm}.  Exponentiation uses
\[
 |e^u-1|\le |u|e^{|u|}
\]
and the rooted convolution algebra; the connected norm is small, so
\(e^{|u|}\) is bounded in a slightly reduced weight.  The exact first,
second, and third field derivatives are the real parts of
\eqref{eq:v4v34-good-first-response},
\eqref{eq:v4v34-good-second-response}, and
\eqref{eq:v4v56-exact-third}.  Their weighted estimates are V34 and
Theorem~\ref{thm:v4v56-weighted-third}.
\end{proof}

Let
\[
 {\rm typ}(I)\in\{{\rm M},{\rm ph},{\rm T}\}
\]
denote a modulus, phase, or tail leaf.  Its elementary activity is denoted
by \(\lambda(I)\); for a tail packet it includes
\(\tau^{|S(I)|}\), with one factor per distinct physical tail cell.

\begin{theorem}[Exact three-species master-parent representation]
\label{thm:v4v68-master-representation}
After the single V16 quadratic normalization is removed, the hard Yukawa
partition function is exactly the finite expansion
\begin{equation}
 Z_{\Lambda}^{\rm Y,hard}
 =
 \mathbb E_{\mu_B}
 \left[
 \prod_S(1+z_S^{\rm M})
 \prod_S(1+z_S^{\rm ph})
 \,\mathcal D_{\rm T}
 \right],
\label{eq:v4v68-master-representation}
\end{equation}
where \(\mathcal D_{\rm T}\) is expanded by the one-cell packets
\(\Theta_x\) and contracted once with \(C_{\rm G}\).
Every \(g_xV_x\) and \(t_xV_x\) occurrence and every determinant
normalization has exactly one ancestry.
\end{theorem}

\begin{proof}
Equation \eqref{eq:v4v68-DG-product} reconstructs the first factor
\(\mathcal D_{\rm G}\).  Lemma~\ref{lem:v4v62-cell-consolidation}
reconstructs the local tail exponential, and the Berezin expectation with
\(C_{\rm G}\) reconstructs \(\mathcal D_{\rm T}\).
The V34 determinant factorization then gives
\(\mathcal D_{\rm G}\mathcal D_{\rm T}=\mathcal D_{\rm Y}\).
The good operator occurs in the first determinant and in the resolvent
defining \(C_{\rm G}\); the latter is a covariance update inside the tail
determinant, not a second copy of \(\mathcal D_{\rm G}\).
\end{proof}

\begin{firewall}[The modulus is external exactly once]
\label{fw:v4v68-modulus-external}
Equation \eqref{eq:v4v68-master-representation} does not multiply the
promoted measure by another modulus.  It abandons the promoted measure for
the collision proof and returns to \(\mu_B\); the modulus factors then
appear once as the activities required to reconstruct
\(\mathcal D_{\rm G}\).
\end{firewall}

\subsection{The restart identities}
\label{sec:v4v68-restart}

Let \(\mathscr F_A\) be the sigma algebra generated by master fields in
\(A\).  All conditional expectations below are those of \(\mu_B\).

\begin{lemma}[Conditional product and tower restart]
\label{lem:v4v68-restart}
Let \(I_1,\ldots,I_k\) be components separated by more than the interaction
range of the master action, and let
\(\Sigma=\Lambda\setminus\dot\bigcup_iI_i\).  If \(F_i\) is measurable in
\(I_i\cup\Sigma\), then
\begin{equation}
 \prod_{i=1}^k
 \mathbb E_{\mu_B}[F_i\mid\mathscr F_\Sigma]
 =
 \mathbb E_{\mu_B}
 \left[\prod_{i=1}^kF_i\mid\mathscr F_\Sigma\right].
\label{eq:v4v68-conditional-product}
\end{equation}
If \(\mathscr F_A\subseteq\mathscr F_C\), then
\begin{equation}
 \mathbb E_{\mu_B}
 \left[
  \mathbb E_{\mu_B}[F\mid\mathscr F_C]
  \mid\mathscr F_A
 \right]
 =
 \mathbb E_{\mu_B}[F\mid\mathscr F_A].
\label{eq:v4v68-tower}
\end{equation}
Both identities remain true coefficientwise for Grassmann-valued finite
polynomials.
\end{lemma}

\begin{proof}
The conditional product law is
Theorem~\ref{thm:v4v38-conditional-product}; integrate the product against
that tensor product.  Equation \eqref{eq:v4v68-tower} is the tower property
of conditional expectation.  A finite Grassmann polynomial has finitely
many scalar coefficient functions, so both identities apply to each
coefficient.
\end{proof}

A collision history carries its underlying set of elementary determinant
leaves in addition to its current conditional transform.  When two
histories meet, let \(D\) be their least common buffered component.  Discard
the provisional child transforms and apply the conditional expectation in
\(D\) directly to the product of their underlying leaves.  If they do not
meet, \eqref{eq:v4v68-conditional-product} retains their product.

\begin{theorem}[Master-parent restart forest]
\label{thm:v4v68-restart-forest}
The preceding rule gives exactly the same finite expansion as successive
conditional elimination.  At every merger:
\begin{enumerate}
\item the integration law is a conditional law of the original
      finite-range parent \(\mu_B\);
\item its Hessian is at least \(\kappa1\), its interaction range is the
      original \(R_0\), and its V38 Witten constants are unchanged;
\item every underlying leaf is integrated once; and
\item provisional response, stopping, and normalization factors below the
      merger are discarded before the union is recomputed.
\end{enumerate}
\end{theorem}

\begin{proof}
Induct on the well-founded collision tree.  For disjoint children separated
by the retained skeleton, equation
\eqref{eq:v4v68-conditional-product} identifies the product of their
conditional transforms with the conditional transform of their product.
At a later ancestor, equation \eqref{eq:v4v68-tower} collapses the nested
conditional expectations to the direct conditional expectation of that
same product.  A collision cumulant is a finite signed sum over partitions
of its child labels.  Apply the product and tower identities separately to
every moment block in that formula.  The partition coefficients and the
elementary label sets are unchanged, so replacement by the union component
changes no cumulant coefficient.

Every union component is a subset of the original coordinate set with the
complement fixed.  Its action is the restriction of
\eqref{eq:v4v38-finite-range-action}.  A principal restriction preserves
the Hessian lower bound, interaction range, off-diagonal stencil, and
weighted commutator bound.  Lemma~\ref{lem:v4v38-conditional-Witten}
therefore applies with its original constants on an arbitrary union
component.  The elementary labels are disjoint child sets and are replaced
only by their union, proving one-use ancestry.  The discarded provisional
objects are alternative representations of the same conditional
expectation and are not additional factors.
\end{proof}

\begin{firewall}[No effective skeleton action is iterated]
\label{fw:v4v68-no-effective-iteration}
The skeleton marginal and
\eqref{eq:v4v38-effective-Hessian} remain valid objects, but V68 does not
use that marginal as the parent of the next collision.  Reapplying a
finite-range Witten theorem to successive effective Hessians would require
the unproved range propagation identified in V45.  The restart forest
instead conditions the original finite-range action on the current
exterior at every merger.
\end{firewall}

\subsection{Uniform local interfaces from the restart}
\label{sec:v4v68-uniform-interface}

\begin{theorem}[All-stage conditional analytic interface]
\label{thm:v4v68-all-stage-interface}
Assume the original V30--V31 finite-range strong-convexity and local
coefficient hypotheses and use geometrically nested V38 buffers.  In the
master-parent restart forest:
\begin{enumerate}
\item every bounded centered modulus or phase descendant is a V53 complex
      Lipschitz load;
\item every unregularized tail descendant uses the V62 finite
      moment--cumulant formula and the V65 first-success guard;
\item every stopped tail descendant is a V53 Lipschitz load;
\item every first and second retained/background mark has one of the
      exhaustive V54--V56 ancestries; and
\item all Witten-response and bounded-collar kernels have a common
      volume- and depth-independent row sum.
\end{enumerate}
Thus the all-stage interface assumed in V57 is acquired for the
fixed-cutoff master-parent hierarchy.
\end{theorem}

\begin{proof}
At every merger the conditional law has the same log-Sobolev constant and
Witten kernel bound by
Theorem~\ref{thm:v4v68-restart-forest}.  The initial modulus loads are
Lemma~\ref{lem:v4v68-modulus-leaves}; the phase loads are
Theorem~\ref{thm:v4v66-mixed-phase-norm}.  Conditional centering preserves
their field gradients.  V53 therefore applies to every merger made only
of bounded descendants.

For an unregularized tail, use the V62 tilt only in the bounded modulus,
phase, and already stopped-tail loads.  All tail variables stay in the
finite moment polynomial.  V65 is formulated in the same finite-range
master parent, and its moment-level first-success identity removes the
hierarchy-height sum.  Its root guard terminates every tail path.

Normalized differentiation of a conditional expectation has the V54
explicit-child and first-score terms.  Two derivatives have the V55
twice-child, two-child, genuine-second-score, and two-first-score terms.
For the master Higgs action, the boundary scores use only the original
finite-range stencils and their V30--V31 local derivative bounds.  Explicit
determinant-child derivatives use V56 for the real modulus, V66 for the
phase, and V65 for the tail.  No determinant is part of the conditional
action, so its derivative is never counted again as a parent score.

Geometric buffers give exponential response factors by
Theorem~\ref{thm:v4v38-buffered-influence}; intersecting collars give the
bounded incidence alternative.  The lattice exponential kernel and the
bounded-overlap collar kernel have row sums independent of the number of
ancestors.  Restart discards provisional lower-level responses at a merger,
while the retained child trees are flattened once by V52.  Induction over
the collision forest proves all five assertions.
\end{proof}

\begin{assessment}[Meaning of the all-stage statement]
\label{ass:v4v68-all-stage-meaning}
Theorem~\ref{thm:v4v68-all-stage-interface} concerns a hierarchical
conditional expansion of one fixed-cutoff master measure.  It removes the
unproved iteration of a nonlocal effective Hessian from V57.  It is not a
claim that a rescaled Wilson action remains in a uniform finite-dimensional
coupling manifold as the ultraviolet lattice spacing tends to zero.
\end{assessment}

\subsection{One correct three-type recursion}
\label{sec:v4v68-three-type-recursion}

Let
\begin{equation}
 \varepsilon_{\rm all}
 :=
 C_{\rm M}\epsilon_Y^{1-\alpha}
 +C_{\rm ph}\epsilon_Y^{1-\alpha}
 +C_{\rm T}\tau.
\label{eq:v4v68-all-leaf}
\end{equation}
Let \(R_{\rm all}\) be the rooted sum of every modulus/phase/tail collision
history in \eqref{eq:v4v68-master-representation}, retaining the projective
tail-set weight.

\begin{theorem}[Master-parent mixed merger]
\label{thm:v4v68-master-mixed-merger}
At every merger of \(k\) children in the restart forest,
\begin{equation}
 |\mathfrak M(I_1,\ldots,I_k)|
 \le
 k!C_*^k
 \left(\prod_i\lambda(I_i)\right)
 \sum_{r=1}^k\prod_{i\ne r}w(I_i,I_r),
\label{eq:v4v68-master-mixed-star}
\end{equation}
through two marks.  The type of every child may independently be
\({\rm M}\), \({\rm ph}\), or \({\rm T}\), and each distinct tail cell is
paid once.
\end{theorem}

\begin{proof}
For bounded children use the complex V53 cumulant theorem under the current
master conditional law.  For a family containing unregularized tails, use
Theorem~\ref{thm:v4v62-multi-tail-cumulant}, with every bounded modulus,
phase, or stopped-tail child in its Lipschitz tilt.  The V65 union guard
supplies the joint response when several tail collars meet.  In either
case the pre-collision loads produce the star by
Theorem~\ref{thm:v4v53-multiway-tree} or
Theorem~\ref{thm:v4v62-joint-response-interface}.  V62 disjoint-set
projectivity pays the union tail set.  The marked statement is item four
of Theorem~\ref{thm:v4v68-all-stage-interface}.
\end{proof}

\begin{theorem}[Three-type rooted closure]
\label{thm:v4v68-three-type-closure}
For sufficiently small \(\epsilon_Y\),
\begin{align}
 R_{\rm all}
 &\le
 e^a\varepsilon_{\rm all}
 +
 A_*\sum_{k\ge2}kC_*^kW^{k-1}R_{\rm all}^k,
\label{eq:v4v68-three-type-recursion}\\
 R_{\rm all}
 &\le2e^a\varepsilon_{\rm all}.
\label{eq:v4v68-three-type-small}
\end{align}
The same radius controls the first two marked sums.  No additional
\(3^k\) type factor occurs.
\end{theorem}

\begin{proof}
The unordered family of all children, without first separating their
types, contributes \(1/k!\).  It cancels the one local \(k!\) in
\eqref{eq:v4v68-master-mixed-star}.  The star center and anchor sums give
\(kW^{k-1}\).  The choice among three leaf types is already generated by
expanding
\[
 R_{\rm all}^k=(R_{\rm M}+R_{\rm ph}+R_{\rm T})^k
\]
with its multinomial coefficients.  Tail projective weights multiply and
V65 stopping weights sum inside moments.  This proves
\eqref{eq:v4v68-three-type-recursion}; apply
Theorem~\ref{thm:v4v57-rooted-majorant}.  Polynomial marked placements use
Lemma~\ref{lem:v4v57-marked-series}.
\end{proof}

\subsection{Corrected hard-Yukawa fixed-cutoff theorem}
\label{sec:v4v68-hard-theorem}

\begin{theorem}[Hard-Yukawa pressure and local ratios over the master parent]
\label{thm:v4v68-hard-Yukawa}
Under the V30--V35 hard-chart hypotheses, the V59--V66 tail and phase
interfaces, and sufficiently small
\(\epsilon_Y\), the complete finite-volume hard-Yukawa normalization
\eqref{eq:v4v68-master-representation} is nonzero on the branch continued
from zero Yukawa coupling.  Its connected coefficients satisfy a
volume-uniform KP bound through two background marks.  For
translation-covariant fixed-cutoff data:
\begin{enumerate}
\item the intensive hard-Yukawa logarithm has a van Hove limit;
\item normalized local hard-Yukawa ratios have thermodynamic limits; and
\item their first two admissible background derivatives are obtained by
      termwise differentiation.
\end{enumerate}
\end{theorem}

\begin{proof}
The exact master representation is
Theorem~\ref{thm:v4v68-master-representation}.  The collision-history
identity and Theorem~\ref{thm:v4v68-three-type-closure} give the rooted KP
bound for its connected logarithm.  Absolute convergence defines the
branch from the empty gas and excludes zeros on that branch.  Assign each
translation class to a root site.  Boundary-rooted histories have vanishing
density along a van Hove sequence, while the rooted majorant permits
dominated convergence.  With a local insertion, histories disconnected
from its support cancel in the normalized ratio.  The marked rooted sums
justify two termwise derivatives.  The V43 scalar cocycle assigns each
normalization once and cancels it from the corresponding local numerator
and denominator.
\end{proof}

\begin{corollary}[What is no longer an assumption]
\label{cor:v4v68-fixed-cutoff-interface-closed}
For the fixed-cutoff hard-Yukawa problem, the all-stage conditional
interface previously assumed in
Theorem~\ref{thm:v4v57-good-pressure} is replaced by
Theorem~\ref{thm:v4v68-all-stage-interface}.  No exponential-locality
hypothesis on an iterated effective skeleton Hessian is used.
\end{corollary}

\begin{firewall}[Fixed-cutoff hierarchy is not ultraviolet RG]
\label{fw:v4v68-not-UV}
The master-parent restart proves a convergent hierarchical organization at
one lattice cutoff.  Ultraviolet removal changes the master action,
couplings, field normalization, and hard-gap chart with the cutoff.  V68
does not supply uniform beta-function control or a continuum limit.
\end{firewall}

\subsection{V68 result ledger and V69 target}
\label{sec:v4v68-conclusion}

V68 detects and repairs the V67 parent mismatch.  The promoted-parent
factorization remains an algebraic identity, while its claimed V65-based
relative proof is withdrawn.  The corrected proof keeps the connected good
modulus, good phase, and rare tail as three external species over the
finite-range master Higgs law.  Conditional product and tower identities
allow every collision to restart in a restriction of that original law.
The curvature, finite range, and Witten constants are therefore uniform
through the entire fixed-cutoff hierarchy.  One three-type recursion then
proves the hard-Yukawa pressure and local-ratio theorem without an iterated
effective-Hessian hypothesis.

V69 should begin the ultraviolet audit of this theorem.  It must restore
the lattice spacing \(a\), canonical field and Yukawa dimensions, and the
hard covariance scale in every leaf and response norm.  The first task is
to determine whether the small parameter
\(\epsilon_Y^{1-\alpha}\) is dimensionless and uniform along a proposed
cutoff trajectory, or whether the hard-gap covariance makes it deteriorate
as \(a\downarrow0\).  A negative power count should be recorded as an
obstruction rather than hidden in fixed-cutoff constants.



\clearpage
\section{V69: ultraviolet scaling audit of the fixed-cutoff hard-Yukawa theorem}
\label{sec:v4v69}

V68 proves a fixed-cutoff hard-Yukawa cluster theorem over the finite-range
master parent.  Its constants were written in lattice units.  V3 already
proves that a four-dimensional Yukawa shell is marginal: its connected
second-order contribution has one factor per correlation cell and no
positive cutoff power.  The present note does not repeat that result.
Instead it restores the lattice spacing in the specific constants used by
V34, V38, V53, V65, and V68.

Three independent degenerations appear.  A nonzero Yukawa coupling is
dimensionless and does not shrink canonically.  A fixed physical Higgs mass
gives a strong-convexity constant of order \(a^2\) in dimensionless lattice
variables.  A hard fermion threshold fixed in physical units gives a
dimensionless hard covariance of order \(a^{-1}\).  Consequently the V68
theorem is not uniform along such a naive cutoff trajectory.  A valid
continuum use must be shellwise, with fields rescaled at every shell and
the complete marginal Higgs--Yukawa law included in the running parent as
required by V3.

\subsection{Canonical lattice variables}
\label{sec:v4v69-canonical}

In four Euclidean dimensions, let
\[
 x_{\rm phys}=ax,\qquad
 \Phi_x=aH(ax),\qquad
 \Psi_x=a^{3/2}\psi(ax),\qquad
 \overline\Psi_x=a^{3/2}\overline\psi(ax).
\]
These are dimensionless lattice fields.  The continuum Yukawa term becomes
\begin{equation}
 a^4\sum_x
 y(a)\,\overline\psi(ax)H(ax)\psi(ax)
 =
 \sum_x y(a)\,\overline\Psi_x\Phi_x\Psi_x.
\label{eq:v4v69-Yukawa-rescaling}
\end{equation}

\begin{lemma}[The Yukawa coupling has no canonical cutoff gain]
\label{lem:v4v69-Yukawa-marginal}
The coefficient in the dimensionless local lattice packet is
\begin{equation}
 \epsilon_Y(a)=\sum_\alpha|y_\alpha(a)|.
\label{eq:v4v69-epsilon-running}
\end{equation}
No factor \(a^\omega\), \(\omega>0\), is generated by canonical field
rescaling.
\end{lemma}

\begin{proof}
The powers in the left side of
\eqref{eq:v4v69-Yukawa-rescaling} are
\[
 a^4a^{-3/2}a^{-1}a^{-3/2}=a^0.
\]
Summing the finitely many internal Yukawa coefficients gives
\eqref{eq:v4v69-epsilon-running}, which is the spacing-dependent version
of the definition in V31.
\end{proof}

\begin{assessment}[Relation to the V3 shell obstruction]
\label{ass:v4v69-relation-V3}
Lemma~\ref{lem:v4v69-Yukawa-marginal} is the coefficient-level audit.
Proposition~\ref{prop:v4v3-Yukawa-scaling} is the stronger fluctuation
statement: a complete shell has connected size
\(\operatorname{tr}(Y^*Y)K^4\) per fixed physical volume.  Together they
exclude the use of canonical dimensions as a source of V68 smallness.
\end{assessment}

\subsection{The scalar curvature constant in lattice units}
\label{sec:v4v69-curvature}

For a free massive scalar component, the dimensionless lattice quadratic
action has the form
\begin{equation}
 S_{a,2}(\Phi)
 =
 \frac12\sum_{x,\mu}|\Phi_{x+\widehat\mu}-\Phi_x|^2
 +
 \frac12(am_H)^2\sum_x|\Phi_x|^2.
\label{eq:v4v69-lattice-scalar}
\end{equation}
Interactions and counterterms may change this expression, but
\eqref{eq:v4v69-lattice-scalar} is an exact test of any claimed
cutoff-uniform theorem based only on a fixed physical mass.

\begin{proposition}[Fixed physical mass does not give uniform lattice curvature]
\label{prop:v4v69-curvature-loss}
On a periodic lattice, the least Hessian eigenvalue of
\eqref{eq:v4v69-lattice-scalar} is
\begin{equation}
 \kappa_a=(am_H)^2.
\label{eq:v4v69-kappa-a}
\end{equation}
Therefore \(\inf_a\kappa_a=0\) as \(a\downarrow0\) for fixed
\(0<m_H<\infty\).
\end{proposition}

\begin{proof}
The lattice Fourier symbol is
\[
 4\sum_{\mu=1}^4\sin^2(p_\mu a/2)+(am_H)^2.
\]
Its minimum is attained at the constant mode \(p=0\), giving
\eqref{eq:v4v69-kappa-a}.
\end{proof}

A positive quartic does not repair this global lower bound by itself: at
the origin its Hessian vanishes.  In a broken chart, radial curvature does
not automatically give a uniform bound in gauge-orbit and Goldstone
directions.  Any stronger \(\kappa_a\) must therefore be an additional
running-parent theorem, not a consequence of a fixed physical mass.

\subsection{Sharp dependence of the V53 load on curvature}
\label{sec:v4v69-LSI}

The constants in V53 were allowed to depend on \(\kappa_*\).  The dependence
cannot be suppressed in a cutoff limit.

\begin{proposition}[Gaussian necessity of the normalized Lipschitz load]
\label{prop:v4v69-Gaussian-load}
Let \(X_\kappa\) be a centered real Gaussian of variance
\(\kappa^{-1}\), and let \(F=LX_\kappa\).  Then
\begin{align}
 {\rm Lip}(F)&=L,\label{eq:v4v69-Gaussian-Lip}\\
 \kappa(F,F)&=\frac{L^2}{\kappa},\label{eq:v4v69-Gaussian-cumulant}\\
 \mathbb E e^{tF}
 &=\exp\!\left(\frac{t^2L^2}{2\kappa}\right).
\label{eq:v4v69-Gaussian-generator}
\end{align}
Hence a V53 generator disk with a fixed bound at \(t=1\) requires
\begin{equation}
 \frac{L}{\sqrt\kappa}\le c
\label{eq:v4v69-normalized-load}
\end{equation}
for a fixed numerical \(c\).  No cumulant bound proportional to \(L^2\)
can have a constant uniform as \(\kappa\downarrow0\).
\end{proposition}

\begin{proof}
The first identity is immediate.  The second is the variance of
\(LX_\kappa\), and the third is the Gaussian moment-generating function.
If \(L^2/\kappa\) is unbounded, the right side of
\eqref{eq:v4v69-Gaussian-generator} is not confined to any fixed
neighborhood of one.  Equation \eqref{eq:v4v69-Gaussian-cumulant} proves
the final assertion directly.
\end{proof}

For a V66 modulus or phase leaf, the unnormalized field load is
\(L=O(\epsilon_Y(a))\).  Combining
\eqref{eq:v4v69-kappa-a} and
\eqref{eq:v4v69-normalized-load} gives the necessary fixed-mass condition
\begin{equation}
 \epsilon_Y(a)\lesssim am_H.
\label{eq:v4v69-Yukawa-curvature-condition}
\end{equation}
A canonically marginal trajectory with
\(\epsilon_Y(a)\to y_*\ne0\) violates this condition.

\begin{firewall}[Small at one cutoff is not uniform in the cutoff]
\label{fw:v4v69-small-not-uniform}
The statement \(\epsilon_Y<\epsilon_*\) in V68 uses an
\(\epsilon_*\) depending on the parent curvature and response constants.
Proposition~\ref{prop:v4v69-Gaussian-load} shows that it cannot be read as
one fixed threshold when \(\kappa_a\to0\).
\end{firewall}

\subsection{The tail exponent also contains the curvature}
\label{sec:v4v69-tail}

The subgaussian estimate behind V31 and V65 has, with constants restored,
the form
\begin{equation}
 \mathbb P(|\Phi_X-\mathbb E\Phi_X|\ge R)
 \le
 C^{|X|}\exp(-c\kappa_aR^2)
\label{eq:v4v69-tail-curvature}
\end{equation}
for a fixed-size cell.  Therefore the V34 choice
\(R_Y=\epsilon_Y^{-\alpha}\) gives the actual tail parameter
\begin{equation}
 \tau_a
 \le
 K\exp\!\left[
 -c\kappa_a\epsilon_Y(a)^{-2\alpha}
 \right].
\label{eq:v4v69-tail-parameter-a}
\end{equation}

\begin{proposition}[Failure of the fixed threshold on the naive trajectory]
\label{prop:v4v69-tail-loss}
If \(m_H\) is fixed, \(\epsilon_Y(a)\to y_*>0\), and
\(\kappa_a=(am_H)^2\), then the exponent in
\eqref{eq:v4v69-tail-parameter-a} tends to zero.  In particular V65 does
not obtain a uniformly small rare-tail activity from this threshold.
\end{proposition}

\begin{proof}
Substitution gives
\[
 \kappa_a\epsilon_Y(a)^{-2\alpha}
 =
 a^2m_H^2\epsilon_Y(a)^{-2\alpha}\longrightarrow0.
\]
Thus the displayed exponential tends to one rather than zero.
\end{proof}

One could choose
\(R_{Y,a}\gg\kappa_a^{-1/2}\), but the good determinant walk would then
contain
\(\epsilon_Y(a)R_{Y,a}\).  For fixed nonzero \(\epsilon_Y\), this grows at
least as \(a^{-1}\).  The good/tail split therefore cannot simultaneously
make the good operator perturbative and the tail rare when the full-field
curvature vanishes in lattice units.

\begin{theorem}[Good/tail incompatibility at vanishing curvature]
\label{thm:v4v69-good-tail-incompatibility}
Suppose a cutoff family uses a threshold \(R_a\) and requires both
\begin{align}
 M_{C,a}\epsilon_Y(a)R_a&\le c_0,
\label{eq:v4v69-good-requirement}\\
 \kappa_aR_a^2&\longrightarrow\infty.
\label{eq:v4v69-tail-requirement}
\end{align}
Then necessarily
\begin{equation}
 \frac{\epsilon_Y(a)M_{C,a}}{\sqrt{\kappa_a}}
 \longrightarrow0.
\label{eq:v4v69-compatibility-necessary}
\end{equation}
For fixed nonzero \(\epsilon_Y\), uniformly positive \(M_{C,a}\), and
\(\kappa_a\to0\), no such threshold exists.
\end{theorem}

\begin{proof}
Multiply and divide:
\[
 \frac{\epsilon_YM_{C,a}}{\sqrt{\kappa_a}}
 =
 \bigl(\epsilon_YM_{C,a}R_a\bigr)
 \bigl(\kappa_aR_a^2\bigr)^{-1/2}.
\]
The first factor is bounded by
\eqref{eq:v4v69-good-requirement}; the second tends to zero by
\eqref{eq:v4v69-tail-requirement}.  This proves
\eqref{eq:v4v69-compatibility-necessary}.  The final claim is immediate.
\end{proof}

\subsection{Scaling of the hard fermion covariance}
\label{sec:v4v69-fermion}

Let \(A_a^{\rm phys}\) be the dimensionful lattice Dirac operator and let
\[
 \widehat A_a:=aA_a^{\rm phys}
\]
be its dimensionless matrix in the \(\Psi\) variables.  If the hard
projector selects
\(|A_a^{\rm phys}|\ge\Lambda_{\rm h}(a)\), then
\begin{equation}
 \|\widehat C_{{\rm h},a}\|
 =
 \|(\widehat A_a|_{\rm h})^{-1}\|
 \le
 \frac1{a\Lambda_{\rm h}(a)}.
\label{eq:v4v69-hard-covariance-scale}
\end{equation}

\begin{proposition}[A uniform dimensionless hard gap is an ultraviolet split]
\label{prop:v4v69-hard-gap}
A threshold-based uniform V33/V34 hard-covariance certificate requires
\begin{equation}
 a\Lambda_{\rm h}(a)\ge\delta_{\rm h}>0.
\label{eq:v4v69-hard-threshold}
\end{equation}
Thus a certified hard threshold must be of order \(a^{-1}\).  If
\(\Lambda_{\rm h}\) is fixed in physical units, the right side of
\eqref{eq:v4v69-hard-covariance-scale} grows like \(a^{-1}\), so the V34
estimate becomes only
\begin{equation}
 \zeta_{Y,a}
 \lesssim
 \frac{\epsilon_Y(a)^{1-\alpha}}
      {a\Lambda_{\rm h}}.
\label{eq:v4v69-Neumann-loss}
\end{equation}
This estimate does not establish uniform smallness.  If the hard spectrum
contains eigenvalues comparable to \(\Lambda_{\rm h}\), then
\(M_{C,a}\asymp(a\Lambda_{\rm h})^{-1}\), and the displayed deterioration
is actual.
\end{proposition}

\begin{proof}
Equation \eqref{eq:v4v69-hard-threshold} is exactly the condition that the
right side of \eqref{eq:v4v69-hard-covariance-scale} be bounded.
The V34 walk estimate is
\(\zeta_Y\le CM_C\epsilon_Y^{1-\alpha}\); insert
\(M_C=(a\Lambda_{\rm h})^{-1}\).
\end{proof}

The hard theorem may therefore control modes at a fixed fraction of the
ultraviolet scale after shell rescaling.  It cannot place all modes above a
fixed physical mass into one uniformly hard covariance.  The remaining
fermion modes require a soft/infrared construction and the chiral Ward
ledger.

\subsection{Loss of a lattice-unit response row sum}
\label{sec:v4v69-response}

For the quadratic scalar test, the lattice correlation length is
\[
 \xi_a\asymp\kappa_a^{-1/2}\asymp(am_H)^{-1}.
\]
An exponential response majorant in lattice distance therefore has rate
\(m_a\asymp am_H\).  In four dimensions,
\begin{equation}
 \sum_{x\in\mathbb Z^4}e^{-m_a|x|}
 \asymp m_a^{-4}
 \asymp (am_H)^{-4}.
\label{eq:v4v69-row-sum-loss}
\end{equation}

\begin{lemma}[The V57 spatial constant is not cutoff uniform in microscopic units]
\label{lem:v4v69-W-loss}
If the V68 collision kernel is summed over microscopic lattice anchors
without shell rescaling, its row-sum constant \(W_a\) can grow at least on
the correlation-volume scale \(\xi_a^4\).  The fixed-point condition
\(C_*W_aR_a<1\) then has no cutoff-uniform consequence from a fixed
dimensionless leaf activity.
\end{lemma}

\begin{proof}
The lattice shell with radii \(r\) contains order \(r^3\) sites.  Comparing
the sum with
\(\int_0^\infty r^3e^{-m_ar}dr=6m_a^{-4}\) proves
\eqref{eq:v4v69-row-sum-loss}.  The V57 recursion contains \(W_a^{k-1}\);
already the quadratic term requires the product of \(W_a\) and the rooted
activity to be small.
\end{proof}

\subsection{A scale-covariant certificate}
\label{sec:v4v69-certificate}

The preceding tests identify the dimensionless combinations that a
shellwise theorem must control.  At shell \(j\), after rescaling that shell
to unit lattice spacing, let
\[
 \epsilon_{Y,j},\quad
 \kappa_j,\quad
 M_{C,j},\quad
 W_j
\]
be respectively the running Yukawa norm, parent curvature, hard covariance
norm, and response row sum.

\begin{theorem}[Necessary inputs for a uniform hard-Yukawa shell step]
\label{thm:v4v69-shell-certificate}
A cutoff-uniform use of the V34--V68 mechanism requires uniform constants
such that
\begin{align}
 M_{C,j}\epsilon_{Y,j}^{1-\alpha}
 &\le q_{\rm F}<1,
\label{eq:v4v69-shell-good}\\
 \frac{\epsilon_{Y,j}}{\sqrt{\kappa_j}}
 &\le q_{\rm LSI},
\label{eq:v4v69-shell-LSI}\\
 \kappa_j\epsilon_{Y,j}^{-2\alpha}
 &\ge q_{\rm T},
\label{eq:v4v69-shell-tail}\\
 W_j\left(
 \epsilon_{Y,j}^{1-\alpha}
 +e^{-c\kappa_j\epsilon_{Y,j}^{-2\alpha}}
 \right)
 &\le q_{\rm KP}<q_*.
\label{eq:v4v69-shell-KP}
\end{align}
Here \(q_*\) is the V57 rooted radius after the fixed local constants are
included.  These inequalities are necessary bookkeeping conditions for
the existing proof, not a construction of a running trajectory.
\end{theorem}

\begin{proof}
Equation \eqref{eq:v4v69-shell-good} is the V34 Neumann condition.
Equation \eqref{eq:v4v69-shell-LSI} is forced by
Proposition~\ref{prop:v4v69-Gaussian-load}.
Equation \eqref{eq:v4v69-shell-tail} keeps the V65 rare parameter uniformly
small.  The leaf sum in V68 is the sum of the good modulus/phase power and
the tail exponential; insertion into the V57 convergence condition gives
\eqref{eq:v4v69-shell-KP}.
\end{proof}

\begin{corollary}[The naive fixed-physical-scale trajectory fails the certificate]
\label{cor:v4v69-naive-fails}
If
\[
 \epsilon_{Y,j}\to y_*>0,\qquad
 \kappa_a\asymp a^2m_H^2,\qquad
 \Lambda_{\rm h}(a)\asymp1,\qquad
 M_{C,a}\asymp(a\Lambda_{\rm h})^{-1}
\]
in physical units, then
\eqref{eq:v4v69-shell-good},
\eqref{eq:v4v69-shell-LSI}, and
\eqref{eq:v4v69-shell-tail} all fail as \(a\downarrow0\).
\end{corollary}

\begin{proof}
Use Propositions~\ref{prop:v4v69-curvature-loss},
\ref{prop:v4v69-tail-loss}, and
\ref{prop:v4v69-hard-gap}.
\end{proof}

\begin{firewall}[The certificate is not a beta-function theorem]
\label{fw:v4v69-no-beta}
Canonical marginality says only that \(y_j\) has no forced power of the
spacing.  Whether the interacting gauge--Higgs--Yukawa beta functions
produce a bounded trajectory satisfying
\eqref{eq:v4v69-shell-good}--\eqref{eq:v4v69-shell-KP} is a separate
nonperturbative problem.  V69 neither assumes asymptotic freedom for
hypercharge nor promotes a perturbative running coupling to a constructed
measure.
\end{firewall}

\subsection{V69 result ledger and V70 target}
\label{sec:v4v69-conclusion}

V69 restores canonical lattice units in the V68 proof.  The Yukawa
coefficient is dimensionless.  A fixed physical scalar mass yields
\(\kappa_a=(am_H)^2\), forcing the normalized Lipschitz load
\(\epsilon_Y/\sqrt{\kappa_a}\) to diverge.  The rare-tail exponent becomes
\(\kappa_a\epsilon_Y^{-2\alpha}\), while a hard fermion threshold fixed in
physical units makes the dimensionless covariance grow as \(a^{-1}\).
The lattice-unit response row sum also grows like a correlation volume.
Thus the fixed-cutoff theorem is not a cutoff-uniform continuum theorem.

The viable route is the marginal-parent shell programme already isolated
in V3--V4: rescale every momentum shell to unit size, include the complete
running marginal gauge--Higgs--Yukawa packet in its normalized parent, and
apply the V68 cluster mechanism only to the shell's controlled irrelevant
and hard residual.  V70 is a compulsory companion audit.  After that audit
it should formulate an exact matching theorem between a shell marginal
projection and the hard residual certificate
\eqref{eq:v4v69-shell-good}--\eqref{eq:v4v69-shell-KP}, with every relevant
and marginal coefficient assigned to the running coupling vector rather
than charged as a small polymer.


\clearpage
\section{V70: companion audit and exact marginal-to-hard-residual matching}
\label{sec:v4v70}

V69 proves that the V68 hard-Yukawa theorem cannot be applied uniformly to
the complete canonically marginal Yukawa coupling.  V3--V4 already give the
correct upstream split: the full marginal Higgs--Yukawa packet belongs to
the running shell parent, while only the irrelevant complement may be
charged to a summable residual norm.  This compulsory checkpoint records
the current companion files and constructs the exact determinant quotient
which implements that split.

The theorem below is conditional on a uniformly controlled marginal shell
carrier.  It does not construct that carrier.  Its gain is precise: once
the complete marginal hard fermion operator is factored first, the
remaining hard determinant sees the irrelevant Yukawa coefficient
\(\varepsilon_{{\rm h},j}=O(a_j^\omega)\), not the full dimensionless
coupling \(y_j\).  The V69 certificate then holds uniformly and the
residual shell corrections are summable.

\subsection{Compulsory companion audit}
\label{sec:v4v70-audit}

The exact newest companion filenames inspected were:
\begin{enumerate}
\item
\texttt{Renewal\_Yang\_Mills\_Mass\_Gap\_Research\_Notes\_V91.tex},
\(1\,547\,662\) bytes, SHA--256
\[
 \texttt{47AEB1DB15213E86152E1E5699D85BA108617714234004495135B7D9F807D266}.
\]
At inspection time this file contained no V91 chapter; its mathematical
endpoint was V90.  V89 corrects a \(3+1\) target which divided raw edge
stocks by only one endpoint mass.  V90 then reduces failure of the literal
normalized card by one fixed finite Gram resolution, proves that any
persistent failure has a cofinal demanded-row atom, and removes diffuse,
unbalanced, and cut-crossing alternatives.  Throughout, the local tree
weight and the source-occurring bank remain attached to the same resolved
word.
\item
\texttt{Renewal\_NS\_Stability\_Research\_Notes\_V31.tex},
\(887\,537\) bytes, SHA--256
\[
 \texttt{F1121B62238AD5491BB7CF03635EA9934942EDF573ED8FAAA9EE79E1D38A6696}.
\]
Its endpoint proves that a flat dyadic mark removes stopping-family
multiplicity, but exact unmarking costs the first shell moment.  Heat
lifting pays that moment only in the critical source norm; weaker source
norms restore the missing scale factor.  The remaining opportunity is the
correlation of actual recent entry, signed triad work, and the active high
parent.
\end{enumerate}

The type-correct transfer is proof order.  The SM residual must first be
divided by the complete marginal determinant; a bound for the unfactored
full determinant cannot be converted into a residual bound, just as the YM
raw-stock atlas cannot be converted into its normalized atlas.  The
complete determinant quotient is expanded before its local jets are
projected, paralleling the assembled YM Gram square.  Finally, a uniform
bound at each shell is insufficient: the exact rescaling power of every
residual shell must remain visible until the shell sum is taken, as in the
NS marked-to-unmarked first moment.

\begin{firewall}[Companion estimates are not imported]
\label{fw:v4v70-companion-types}
No YM product-state capacity and no NS heat estimate is used as an SM
inequality.  Only the requirements of literal normalization, assembled
carriers, and scale-visible summation are transferred.
\end{firewall}

\subsection{Marginal and residual hard fermion operators}
\label{sec:v4v70-operators}

At shell \(j\), let
\[
 V_j=M_j(g_j)+R_j,
 \qquad
 \mathfrak M_jR_j=0,
\]
be the V4 decomposition.  Let \(\mathfrak P_{{\rm hY},j}R_j\) be the
fermion-bilinear hard-Yukawa part of the irrelevant remainder and denote
its multiplication operator by \(U_j(R_j,\phi)\).  All marginal Yukawa
matrices, marginal Higgs coefficients, and their prescribed counterterms
are included in the shell hard operator
\begin{equation}
 A_j^{\rm marg}(\phi,g_j).
\label{eq:v4v70-marginal-operator}
\end{equation}
Assume it is invertible on the selected hard subspace and put
\[
 C_j^{\rm marg}:=(A_j^{\rm marg})^{-1}.
\]

\begin{theorem}[Exact marginal determinant quotient]
\label{thm:v4v70-determinant-quotient}
At every finite shell cutoff,
\begin{align}
 \det(A_j^{\rm marg}+U_j)
 &=
 \det A_j^{\rm marg}\,
 \mathcal D_j^{\rm res},
\label{eq:v4v70-det-factor}\\
 \mathcal D_j^{\rm res}
 &:=
 \det(1+C_j^{\rm marg}U_j).
\label{eq:v4v70-residual-det}
\end{align}
The numerator and denominator use the same anomaly-normalized determinant
line and the same hard projector.  Hence every marginal Yukawa occurrence
belongs to \(\det A_j^{\rm marg}\), while every residual occurrence belongs
to \(\mathcal D_j^{\rm res}\).
\end{theorem}

\begin{proof}
Factor
\[
 A_j^{\rm marg}+U_j
 =
 A_j^{\rm marg}(1+C_j^{\rm marg}U_j)
\]
on the common finite-dimensional hard subspace and take determinants.
Using one fixed determinant-line branch makes the multiplicative identity
literal, including its phase.  The two summands in the operator are the
images of the complementary V4 projections, so their coefficient
ancestries are disjoint.
\end{proof}

\begin{firewall}[Project before declaring the residual coupling]
\label{fw:v4v70-project-first}
The small parameter of
\(\mathcal D_j^{\rm res}\) is the norm of
\(\mathfrak P_{{\rm hY},j}R_j\), not
\(\epsilon_{Y,j}=\sum_\alpha|y_{\alpha,j}|\) in the marginal operator.
Applying the V34 good/tail split to the unfactored full operator would
reintroduce the V69 marginal obstruction and would charge a running
coupling as an irrelevant polymer.
\end{firewall}

\subsection{The residual local norm}
\label{sec:v4v70-residual-norm}

Assume the hard-Yukawa residual projection is bounded from the V4 analytic
space to the local affine Yukawa coefficient norm:
\begin{equation}
 \varepsilon_{{\rm h},j}
 :=
 \|\mathfrak P_{{\rm hY},j}R_j\|_{{\rm loc},j}
 \le C_{\rm h}\|R_j\|_j.
\label{eq:v4v70-residual-epsilon}
\end{equation}
For the V4 backward trajectory,
\begin{equation}
 \|R_j\|_j\le C_Ra_{j+1}^{\omega},
 \qquad \omega>0.
\label{eq:v4v70-R-power}
\end{equation}
Thus
\begin{equation}
 \varepsilon_{{\rm h},j}
 \le C_0a_{j+1}^{\omega}.
\label{eq:v4v70-epsilon-power}
\end{equation}

\begin{lemma}[Why the projection bound is an additional interface]
\label{lem:v4v70-projection-interface}
The identity \(\mathfrak M_jR_j=0\) alone does not imply
\eqref{eq:v4v70-residual-epsilon}.  The latter follows when the V4 analytic
norm controls the local coefficient functional defining
\(\mathfrak P_{{\rm hY},j}\), uniformly in \(j\).
\end{lemma}

\begin{proof}
A vector can lie in the kernel of a projection while having arbitrarily
large norm, so the normalization identity gives no size estimate.
If the coefficient functional is bounded, then
\[
 \|\mathfrak P_{{\rm hY},j}R_j\|_{{\rm loc},j}
 \le\|\mathfrak P_{{\rm hY},j}\|\,\|R_j\|_j,
\]
which is \eqref{eq:v4v70-residual-epsilon}.
\end{proof}

\subsection{Uniform marginal-parent certificate}
\label{sec:v4v70-parent-certificate}

The marginal shell carrier must provide constants which do not deteriorate
with \(j\).  Assume
\begin{align}
 D_\phi^2S_j^{\rm marg}&\ge\kappa_0\,1,
\label{eq:v4v70-uniform-curvature}\\
 \|C_j^{\rm marg}\|_{a_0}&\le M_0,
\label{eq:v4v70-uniform-hard-C}\\
 \sup_jW_j&\le W_0,
\label{eq:v4v70-uniform-W}
\end{align}
with \(\kappa_0>0\).  Assume also that the positive marginal shell carrier
has a finite-range master representation, or an exact conditional-product
representation with the uniform V38 buffer and Witten constants, and that
its first three local field derivatives and first two admissible background
derivatives have the uniform V54--V56 stencil bounds.  These statements
concern the rescaled unit shell.

\begin{firewall}[The marginal phase is prior data]
\label{fw:v4v70-marginal-phase}
The complete marginal determinant, including its possibly order-one phase,
belongs to the constructed marginal shell carrier and its normalization
ledger.  V70 applies the small complex phase argument only to
\(\mathcal D_j^{\rm res}\).  It does not prove positivity or
zero-freeness of the order-one marginal chiral determinant.
\end{firewall}

Choose \(0<\alpha<1\) and set
\begin{equation}
 R_{{\rm h},j}:=\varepsilon_{{\rm h},j}^{-\alpha}.
\label{eq:v4v70-residual-threshold}
\end{equation}

\begin{theorem}[The V69 certificate holds for the irrelevant hard residual]
\label{thm:v4v70-residual-certificate}
Under
\eqref{eq:v4v70-epsilon-power}--\eqref{eq:v4v70-uniform-W}, there is
\(j_0\) such that, for \(j\ge j_0\),
\begin{align}
 M_0\varepsilon_{{\rm h},j}^{1-\alpha}
 &\longrightarrow0,
\label{eq:v4v70-good-limit}\\
 \frac{\varepsilon_{{\rm h},j}}{\sqrt{\kappa_0}}
 &\longrightarrow0,
\label{eq:v4v70-LSI-limit}\\
 \kappa_0\varepsilon_{{\rm h},j}^{-2\alpha}
 &\longrightarrow\infty,
\label{eq:v4v70-tail-limit}\\
 W_0\left[
 \varepsilon_{{\rm h},j}^{1-\alpha}
 +
 e^{-c\kappa_0\varepsilon_{{\rm h},j}^{-2\alpha}}
 \right]
 &\longrightarrow0.
\label{eq:v4v70-KP-limit}
\end{align}
Consequently the residual determinant
\(\mathcal D_j^{\rm res}\) satisfies the V34--V68 hard-sector cluster
theorem with constants uniform for all sufficiently ultraviolet shells.
\end{theorem}

\begin{proof}
Equation \eqref{eq:v4v70-epsilon-power} tends to zero geometrically because
\(a_j=q^{-j}a_0\).  The first two limits follow immediately.  The inverse
power in \eqref{eq:v4v70-tail-limit} diverges.  Its exponential tends to
zero faster than every power of \(\varepsilon_{{\rm h},j}\), proving
\eqref{eq:v4v70-KP-limit}.  These are precisely
\eqref{eq:v4v69-shell-good}--\eqref{eq:v4v69-shell-KP}, with
\(\epsilon_{Y,j}\) replaced by the projected residual norm.  The uniform
marginal-parent constants permit the V68 master-parent restart on the
rescaled shell.
\end{proof}

The finitely many shells below \(j_0\) are not part of the ultraviolet
summability argument.  They require finite-scale nonzero normalizations and
may be absorbed into the initial data only after those facts are verified.

\subsection{Assemble the residual logarithm before projecting its jets}
\label{sec:v4v70-log-before-projection}

Let \(\mathscr L_j^{\rm h}(g_j,R_j)\) be the connected logarithm generated
by \(\mathcal D_j^{\rm res}\), including its good modulus, residual phase,
tail, and all mixed collisions.  V68 constructs this as one three-type
rooted sum.  Define its contributions to the next marginal and irrelevant
coordinates by
\begin{align}
 \Delta g_{j-1}^{\rm h}
 &:=
 \mathfrak M_{j-1}\mathcal S_j
 \mathscr L_j^{\rm h}(g_j,R_j),
\label{eq:v4v70-hard-beta-piece}\\
 \Delta R_{j-1}^{\rm h}
 &:=
 \mathfrak I_{j-1}\mathcal S_j
 \mathscr L_j^{\rm h}(g_j,R_j).
\label{eq:v4v70-hard-remainder-piece}
\end{align}

\begin{theorem}[Exact matching with the V4 shell map]
\label{thm:v4v70-exact-matching}
Equations
\eqref{eq:v4v70-hard-beta-piece}--%
\eqref{eq:v4v70-hard-remainder-piece} are complementary projections of the
one residual determinant occurrence in the exact shell map
\eqref{eq:v4v4-exact-RG}.  In particular:
\begin{enumerate}
\item every local jet of engineering dimension at most four generated by
      the hard residual enters \(\Delta g_{j-1}^{\rm h}\);
\item only the projected complement enters
      \(\Delta R_{j-1}^{\rm h}\); and
\item their sum is the complete rescaled hard-residual logarithm.
\end{enumerate}
No connected word is assigned to both outputs.
\end{theorem}

\begin{proof}
Factor the marginal determinant by
Theorem~\ref{thm:v4v70-determinant-quotient} inside the one shell
integrand.  Finite-cutoff zero-freeness from
Theorem~\ref{thm:v4v70-residual-certificate} defines its connected
logarithm on the branch from \(R_j=0\).  Apply
\(\mathcal S_j\), then the complementary identity
\(\mathfrak M_{j-1}+\mathfrak I_{j-1}=1\).
This is the V4 Taylor allocation applied after the complete determinant
quotient has been assembled.  Linearity of the projections proves the
three assertions.
\end{proof}

\begin{firewall}[Do not estimate a larger unprojected target]
\label{fw:v4v70-literal-target}
A bound on \(\mathscr L_j^{\rm h}\) proves analyticity but is not by itself
the irrelevant contraction.  Its marginal jets can have their natural
nondecaying scaling and must be assigned to
\(\Delta g_{j-1}^{\rm h}\).  Only
\(\mathfrak I_{j-1}\mathcal S_j\mathscr L_j^{\rm h}\) is the literal
irrelevant target, paralleling the corrected normalization in the YM V89
audit.
\end{firewall}

\subsection{The exact shell power remains visible}
\label{sec:v4v70-shell-sum}

The V68 rooted theorem applied at shell \(j\) gives, with a fixed reduced
weight,
\begin{equation}
 \|\mathscr L_j^{\rm h}(g_j,R_j)\|_{\rm root}
 \le
 C\left[
 \varepsilon_{{\rm h},j}^{1-\alpha}
 +
 e^{-c\kappa_0\varepsilon_{{\rm h},j}^{-2\alpha}}
 \right].
\label{eq:v4v70-shell-root-bound}
\end{equation}

\begin{theorem}[Summable hard-residual shell ledger]
\label{thm:v4v70-shell-summation}
Under \eqref{eq:v4v70-epsilon-power},
\begin{align}
 \sum_{j\ge j_0}
 \varepsilon_{{\rm h},j}^{1-\alpha}
 &\le
 C a_{j_0+1}^{\omega(1-\alpha)},
\label{eq:v4v70-power-sum}\\
 \sum_{j\ge j_0}
 \exp\!\left[
 -c\kappa_0\varepsilon_{{\rm h},j}^{-2\alpha}
 \right]
 &<\infty.
\label{eq:v4v70-tail-sum}
\end{align}
The same sums control the first two marked residual logarithms.
\end{theorem}

\begin{proof}
The first summand is bounded by
\(C a_{j+1}^{\omega(1-\alpha)}\).  Since
\(a_{j+1}=q^{-(j+1)}a_0\), its sum is geometric and gives
\eqref{eq:v4v70-power-sum}.  For the second,
\[
 \varepsilon_{{\rm h},j}^{-2\alpha}
 \ge c_1q^{2\alpha\omega j}
\]
for large \(j\); the resulting double-exponential sequence is summable.
Marked V68 bounds differ only by fixed polynomial factors in the collision
valence and the same shell majorants.
\end{proof}

\begin{firewall}[A supremum over shells loses the ultraviolet payment]
\label{fw:v4v70-no-shell-sup}
Replacing
\(\varepsilon_{{\rm h},j}^{1-\alpha}\) by a uniform constant before
summing would erase the factor
\(a_j^{\omega(1-\alpha)}\).  The scale label must remain attached until
\eqref{eq:v4v70-power-sum} is taken.  This is the SM analogue of the NS V31
fact that removing a flat shell mark costs its exact first moment.
\end{firewall}

\subsection{A conditional matching theorem for the irrelevant trajectory}
\label{sec:v4v70-trajectory}

\begin{theorem}[Hard residual is compatible with V4 backward stability]
\label{thm:v4v70-backward-compatible}
Assume:
\begin{enumerate}
\item a bounded marginal trajectory \(g_j\) and normalized marginal shell
      carriers satisfying
      \eqref{eq:v4v70-uniform-curvature}--%
      \eqref{eq:v4v70-uniform-W};
\item the V4 contraction and source hypotheses on every non-hard residual
      species;
\item the projection bound
      \eqref{eq:v4v70-residual-epsilon}; and
\item the local Taylor remainder after
      \eqref{eq:v4v70-hard-remainder-piece} has scaling dimension at least
      \(4+\omega_{\rm h}\), \(\omega_{\rm h}>0\).
\end{enumerate}
Then the hard-residual contribution may be inserted into the exact V4 beta
map as \(\Delta g_{j-1}^{\rm h}\), while its complementary contribution is
an analytic irrelevant source with a positive shell power.  After reducing
the common remainder ball, the V4 backward contraction theorem and the
summable hard ledger give a unique small irrelevant trajectory subordinate
to \(g_j\).
\end{theorem}

\begin{proof}
Theorem~\ref{thm:v4v70-exact-matching} assigns the marginal hard jets to the
beta map exactly once.  The fourth hypothesis and
Proposition~\ref{prop:v4v4-scaling-certificate} give a strict scaling
factor on \(\Delta R_{j-1}^{\rm h}\).  Analyticity and the quadratic
remainder bound follow from the zero-free rooted expansion in
Theorem~\ref{thm:v4v70-residual-certificate}.  Add this source to the other
V4 irrelevant sources.  Its shell sum is finite by
Theorem~\ref{thm:v4v70-shell-summation}; the common linear contraction and
quadratic smallness then give
Theorem~\ref{thm:v4v4-backward-stability}.
\end{proof}

\begin{assessment}[Acquired and assumed data]
\label{ass:v4v70-status}
V70 acquires the exact determinant quotient, the residual small parameter,
the uniform V69 certificate for an \(O(a_j^\omega)\) affine hard-Yukawa
remainder, the complementary jet allocation, and the summable shell
ledger.  It assumes the bounded marginal trajectory and the uniform
positive marginal shell carrier.  In particular it does not prove the
order-one marginal determinant phase, the positivity/curvature of its
modulus carrier, or the hypercharge component of the exact beta map.
\end{assessment}

\begin{firewall}[The hard residual does not solve the marginal fixed point]
\label{fw:v4v70-marginal-open}
Theorem~\ref{thm:v4v70-backward-compatible} is subordinate to a solution of
\eqref{eq:v4v4-exact-beta}.  Moving \(y_j\) into
\(A_j^{\rm marg}\) repairs the residual power count; it does not construct a
nonperturbative gauge--Higgs--Yukawa trajectory or prove tightness of its
normalized shell laws.
\end{firewall}

\subsection{V70 result ledger and V71 target}
\label{sec:v4v70-conclusion}

V70 records the exact YM V91-named/V90-content and NS V31 fingerprints.  It
uses their proof-order lessons to factor the complete marginal hard
determinant before estimating the residual, retain the assembled residual
logarithm before projection, and keep the exact positive shell power until
summation.  If the V4 irrelevant trajectory has
\(\|R_j\|_j=O(a_j^\omega)\), the projected affine hard-Yukawa coefficient
has the same power.  The V69 hard, LSI, tail, and KP conditions then hold
uniformly over ultraviolet shells, and the residual corrections are
summable.

The next upstream gate is the marginal shell carrier itself.  V71 should
separate two questions which V70 leaves coupled: strong convexity of the
positive modulus carrier for an order-one running Yukawa matrix, and
zero-freeness/normalization of its marginal chiral phase.  The first attack
should derive an exact relative Hessian criterion using the marginal hard
resolvent; if it fails for order-one \(y_j\), the shell parent must retain a
different non-Gaussian joint boson--fermion representation rather than
assuming V34's small determinant promotion.



\clearpage
\section{V71: a relative Hessian and phase-load certificate for the marginal shell carrier}
\label{sec:v4v71}

V70 moves the full canonically marginal Yukawa packet into the running
shell operator and applies the V68 theorem only to the irrelevant
determinant quotient.  The remaining parent hypothesis has two parts.  The
modulus of the marginal determinant must preserve strong convexity of the
positive bosonic shell law, and its phase must have a zero-free normalized
carrier.  Neither follows merely from invertibility at one field.

This note derives quantitative sufficient conditions.  The modulus
condition is an exact trace-resolvent Hessian margin.  The phase condition
uses a fact not needed in V66: a local phase may have an arbitrary additive
constant.  After that constant is rotated away, Poincare variance makes its
mean close to one whenever its Higgs Lipschitz load is small relative to
the positive curvature.  This permits order-one phase values, but still
requires a small normalized phase gradient and a uniform hard homotopy
gap.

\subsection{The marginal hard homotopy}
\label{sec:v4v71-homotopy}

On one rescaled unit shell, let \(S_{\rm b}(\phi)\) be a positive bosonic
action with
\begin{equation}
 D_\phi^2S_{\rm b}\ge\kappa_{\rm b}\,1.
\label{eq:v4v71-base-curvature}
\end{equation}
Let \(A(\phi)\) be the complete marginal hard fermion matrix, including the
running marginal Yukawa coefficients.  Fix an anomaly-normalized reference
\(A_{\rm ref}\) and a differentiable local homotopy
\[
 A_t=A_{\rm ref}+V_t(\phi),
 \qquad 0\le t\le1,
 \qquad A_1=A.
\]

\begin{definition}[Uniform marginal hard certificate]
\label{def:v4v71-hard-certificate}
The homotopy has constants
\((M_{\rm h},J_1,J_2,a_{\rm h})\) if:
\begin{align}
 \sup_{t,\phi}\|A_t^{-1}\|_{a_{\rm h}}
 &\le M_{\rm h},
\label{eq:v4v71-resolvent-bound}\\
 \|D_\phi A_t[h]\|_{\rm loc}
 &\le J_1\|h\|_2,
\label{eq:v4v71-A1}\\
 \|D_\phi^2 A_t[h,k]\|_{\rm loc}
 &\le J_2\|h\|_2\|k\|_2,
\label{eq:v4v71-A2}
\end{align}
with the corresponding exponentially weighted local row/column bounds.
The same bounds through the derivatives required below are assumed for two
admissible background marks.
\end{definition}

The first condition is stronger than invertibility of \(A_1\): it excludes
a zero on the entire determinant-line homotopy and supplies one common
local logarithm.

\subsection{Exact modulus Hessian}
\label{sec:v4v71-modulus}

Put
\[
 C=A^{-1},
 \qquad
 \ell(\phi)=\operatorname{Tr}\log A(\phi),
 \qquad
 u=\operatorname{Re}\ell,
 \qquad
 \theta=\operatorname{Im}\ell .
\]

\begin{lemma}[First and second determinant derivatives]
\label{lem:v4v71-log-derivatives}
For field directions \(h,k\),
\begin{align}
 D_\phi\ell[h]
 &=
 \operatorname{Tr}(C A_h),
\label{eq:v4v71-log-first}\\
 D_\phi^2\ell[h,k]
 &=
 \operatorname{Tr}(C A_{hk})
 -
 \operatorname{Tr}(C A_k C A_h),
\label{eq:v4v71-log-second}
\end{align}
where \(A_h=D_\phi A[h]\) and
\(A_{hk}=D_\phi^2A[h,k]\).
\end{lemma}

\begin{proof}
Differentiate \(A^{-1}A=1\) to obtain
\(D_\phi C[h]=-CA_hC\), and differentiate
\(\operatorname{Tr}\log A\) once and twice.  Cyclicity of the finite trace
gives the displayed order.
\end{proof}

\begin{theorem}[Relative modulus curvature margin]
\label{thm:v4v71-modulus-margin}
Under Definition~\ref{def:v4v71-hard-certificate}, there is a
finite-fibre and bounded-overlap constant \(C_{\rm loc}\) such that
\begin{equation}
 |D_\phi^2u[h,h]|
 \le
 \eta_{\rm det}\|h\|_2^2,
 \qquad
 \eta_{\rm det}
 :=
 C_{\rm loc}\left(
 M_{\rm h}J_2+M_{\rm h}^2J_1^2
 \right).
\label{eq:v4v71-det-Hessian}
\end{equation}
Hence the positive modulus action
\begin{equation}
 S_+(\phi):=S_{\rm b}(\phi)-u(\phi)
\label{eq:v4v71-positive-action}
\end{equation}
satisfies
\begin{equation}
 D_\phi^2S_+\ge\kappa_+\,1,
 \qquad
 \kappa_+:=\kappa_{\rm b}-\eta_{\rm det},
\label{eq:v4v71-positive-curvature}
\end{equation}
whenever \(\eta_{\rm det}<\kappa_{\rm b}\).
The Hessian perturbation also has an exponentially weighted row bound with
the same expression and a possibly smaller exponent.
\end{theorem}

\begin{proof}
Take real parts in
\eqref{eq:v4v71-log-second}.  For the first trace, combine the weighted
row/column norm of \(C\), the local norm of \(A_{hh}\), finite internal
dimension, and bounded cell overlap.  This gives
\(C_{\rm loc}M_{\rm h}J_2\|h\|_2^2\).
For the second trace, use two resolvent row/column bounds and the two local
first-derivative insertions.  Weighted Schur convolution gives
\(C_{\rm loc}M_{\rm h}^2J_1^2\|h\|_2^2\).
This proves \eqref{eq:v4v71-det-Hessian}.  Subtract it from
\eqref{eq:v4v71-base-curvature} to obtain
\eqref{eq:v4v71-positive-curvature}.  Keeping the exponential weights in
the same two convolutions proves the final statement.
\end{proof}

\begin{corollary}[Positive marginal shell LSI and Witten decay]
\label{cor:v4v71-positive-LSI}
If \(\kappa_+>0\), the normalized law
\[
 d\nu_+=Z_+^{-1}e^{-S_+(\phi)}d\phi
\]
satisfies a log-Sobolev inequality with curvature \(\kappa_+\).
If the weighted commutator margin is chosen below \(\kappa_+\), its
one-form Witten inverse has a volume-uniform exponentially decaying kernel.
\end{corollary}

\begin{proof}
Bakry--Emery applies to
\eqref{eq:v4v71-positive-curvature}.  The off-diagonal Hessian is the base
local stencil plus the weighted determinant Hessian from
Theorem~\ref{thm:v4v71-modulus-margin}.  The V30 weighted-commutator
Neumann argument applies after choosing a smaller spatial exponent.
\end{proof}

\begin{firewall}[Order one is a numerical question, not a power count]
\label{fw:v4v71-order-one}
The condition
\(\eta_{\rm det}<\kappa_{\rm b}\) may hold for a coupling which does not
vanish with the cutoff, but only if the numerical hard gap and derivative
constants satisfy the displayed inequality.  Canonical marginality neither
proves nor disproves that inequality.
\end{firewall}

\subsection{Local logarithm densities from the hard homotopy}
\label{sec:v4v71-local-log}

Let \(\dot A_t=\partial_tA_t\).  The homotopy identity is
\begin{equation}
 \ell-\ell_{\rm ref}
 =
 \int_0^1\operatorname{Tr}(A_t^{-1}\dot A_t)\,dt.
\label{eq:v4v71-homotopy-log}
\end{equation}
With \(P_x\) the hard fermion projection over cell \(x\), define
\begin{equation}
 \ell_x(\phi)
 :=
 \int_0^1
 \operatorname{tr}
 \left(P_xA_t^{-1}\dot A_tP_x\right)\,dt.
\label{eq:v4v71-local-density}
\end{equation}
Then
\begin{equation}
 \ell-\ell_{\rm ref}=\sum_x\ell_x,
 \qquad
 \theta-\theta_{\rm ref}=\sum_x\theta_x,
 \qquad
 \theta_x:=\operatorname{Im}\ell_x.
\label{eq:v4v71-density-sum}
\end{equation}

\begin{lemma}[Quasi-local phase-gradient row]
\label{lem:v4v71-phase-gradient}
Suppose \(\dot A_t\), its field derivative, and the derivatives in
Definition~\ref{def:v4v71-hard-certificate} have uniform local weighted
bounds.  Then there are \(a_{\rm ph}>0\) and \(L_x(y)\ge0\) such that
\begin{align}
 |\nabla_{\phi_y}\theta_x|
 &\le L_x(y),
\label{eq:v4v71-local-phase-gradient}\\
 \sup_x\sum_y e^{a_{\rm ph}d(x,y)}L_x(y)
 +
 \sup_y\sum_x e^{a_{\rm ph}d(x,y)}L_x(y)
 &\le L_{\rm ph}<\infty.
\label{eq:v4v71-phase-row}
\end{align}
One may take \(L_{\rm ph}\) bounded by a fixed polynomial in
\(M_{\rm h}\) and the local first-derivative/homotopy constants.
\end{lemma}

\begin{proof}
Differentiate the integrand in
\eqref{eq:v4v71-local-density}:
\[
 D_\phi(A_t^{-1}\dot A_t)[h]
 =
 -A_t^{-1}(D_\phi A_t[h])A_t^{-1}\dot A_t
 +A_t^{-1}D_\phi\dot A_t[h].
\]
Every field derivative is local, while each resolvent has the weighted
bound \eqref{eq:v4v71-resolvent-bound}.  Weighted Schur convolution gives
exponential decay from \(x\) to the marked field cell \(y\).  Integrate in
\(t\) and take imaginary parts.  The row and column sums follow from the
Banach algebra property of the weighted kernel norm.
\end{proof}

\subsection{Phase centering by Poincare variance}
\label{sec:v4v71-phase-centering}

The next lemma permits a large additive phase.  Only its variation under
the positive law must be small.

\begin{lemma}[Rotated phase mean is close to one]
\label{lem:v4v71-rotated-mean}
Let \(\nu\) satisfy the Poincare inequality
\[
 \operatorname{Var}_\nu(F)
 \le\kappa_+^{-1}\mathbb E_\nu|\nabla F|^2.
\]
If \(\vartheta\) is real and \(L\)-Lipschitz, put
\(m=\mathbb E_\nu\vartheta\) and
\(c=\mathbb E_\nu e^{i\vartheta}\).  Then
\begin{equation}
 \left|e^{-im}c-1\right|
 \le
 \frac{L^2}{2\kappa_+}.
\label{eq:v4v71-rotated-disk}
\end{equation}
In particular, if
\(q=L^2/(2\kappa_+)<1\), then
\begin{equation}
 |c|\ge1-q>0.
\label{eq:v4v71-local-zero-free}
\end{equation}
The normalized carrier
\begin{equation}
 P:=c^{-1}e^{i\vartheta}-1
\label{eq:v4v71-normalized-phase}
\end{equation}
is centered and satisfies
\begin{equation}
 {\rm Lip}(P)\le\frac{L}{1-q}.
\label{eq:v4v71-normalized-Lip}
\end{equation}
\end{lemma}

\begin{proof}
Set \(X=\vartheta-m\).  The scalar inequality
\[
 |e^{it}-1-it|\le\frac12t^2
\]
and \(\mathbb E X=0\) give
\[
 |e^{-im}c-1|
 =
 |\mathbb E(e^{iX}-1-iX)|
 \le\frac12\mathbb E X^2
 \le\frac{L^2}{2\kappa_+}.
\]
This proves
\eqref{eq:v4v71-rotated-disk}--%
\eqref{eq:v4v71-local-zero-free}.  The definition of \(c\) gives
\(\mathbb EP=0\), and
\[
 \nabla P=c^{-1}ie^{i\vartheta}\nabla\vartheta.
\]
Use \eqref{eq:v4v71-local-zero-free}.
\end{proof}

For \(\theta_x\) in \eqref{eq:v4v71-density-sum}, set
\begin{equation}
 L_x:=
 \left(\sum_yL_x(y)^2\right)^{1/2},
 \qquad
 c_x:=\mathbb E_{\nu_+}e^{i\theta_x},
 \qquad
 P_x:=c_x^{-1}e^{i\theta_x}-1.
\label{eq:v4v71-local-centered-phases}
\end{equation}

\begin{theorem}[Zero-free local phase factors and centered loads]
\label{thm:v4v71-local-phase-factors}
If
\begin{equation}
 q_{\rm ph}:=\sup_x\frac{L_x^2}{2\kappa_+}<1,
\label{eq:v4v71-local-phase-condition}
\end{equation}
then every \(c_x\) is nonzero and
\begin{equation}
 e^{i(\theta-\theta_{\rm ref})}
 =
 \left(\prod_xc_x\right)\prod_x(1+P_x)
\label{eq:v4v71-phase-factorization}
\end{equation}
exactly.  Each \(P_x\) is centered in \(\nu_+\) and has a quasi-local
Lipschitz row bounded by
\begin{equation}
 (1-q_{\rm ph})^{-1}L_x(y).
\label{eq:v4v71-centered-row}
\end{equation}
The product of the \(c_x\)'s is a V43 scalar phase/modulus cocycle; it is not
inserted into a collision cumulant.
\end{theorem}

\begin{proof}
Apply Lemma~\ref{lem:v4v71-rotated-mean} to every \(\theta_x\).
Multiplying the exact identities
\(e^{i\theta_x}=c_x(1+P_x)\) and using
\eqref{eq:v4v71-density-sum} gives
\eqref{eq:v4v71-phase-factorization}.  The gradient formula in the lemma
and \eqref{eq:v4v71-local-phase-gradient} give
\eqref{eq:v4v71-centered-row}.  Each \(c_x\) is a number after the full
positive shell expectation and hence is eligible for the scalar cocycle.
\end{proof}

\begin{firewall}[Poincare controls variation, not the vacuum angle]
\label{fw:v4v71-angle}
The mean \(m_x=\mathbb E\theta_x\) need not be small.  It contributes to the
argument of \(c_x\) and hence to the vacuum phase pressure.  The estimate
\eqref{eq:v4v71-rotated-disk} controls only the centered fluctuation and
does not set a physical CP angle to zero.
\end{firewall}

\subsection{The centered marginal phase gas}
\label{sec:v4v71-phase-gas}

Because the unary mean of every \(P_x\) vanishes, the logarithm of
\[
 \mathbb E_{\nu_+}\prod_x(1+P_x)
\]
starts with joint cumulants of order at least two.  The V53 theorem applies
to these complex centered variables.

\begin{theorem}[Sufficient marginal phase-pressure criterion]
\label{thm:v4v71-phase-pressure}
Assume \(\kappa_+>0\),
\eqref{eq:v4v71-local-phase-condition}, the quasi-local row
\eqref{eq:v4v71-phase-row}, and the V38/V54--V56 conditional interfaces for
the positive shell carrier.  There is a numerical
\(\delta_{\rm ph}>0\) such that, if
\begin{equation}
 \frac{L_{\rm ph}}{(1-q_{\rm ph})\sqrt{\kappa_+}}
 \le\delta_{\rm ph}
\label{eq:v4v71-phase-KP-condition}
\end{equation}
after the fixed response-row constants are included, then:
\begin{enumerate}
\item the centered phase gas has an absolutely convergent logarithm and is
      zero-free on the branch from zero phase load;
\item its intensive logarithm has a translation-covariant van Hove limit;
\item the scalar product \(\prod_xc_x\) has a well-defined local logarithm
      on the disks \eqref{eq:v4v71-rotated-disk}; and
\item normalized local phase ratios and their first two admissible
      background marks converge.
\end{enumerate}
\end{theorem}

\begin{proof}
Use the normalized loads in
\eqref{eq:v4v71-centered-row}.  The Poincare/LSI normalization shows that
their dimensionless V53 sizes are
\(L_x/[(1-q_{\rm ph})\sqrt{\kappa_+}]\).
Theorem~\ref{thm:v4v53-Lipschitz-cumulant} supplies the product of all
child loads, and the V53 star theorem supplies one response tree.
The V57 unordered-child recursion converges under
\eqref{eq:v4v71-phase-KP-condition}.  This proves the first two statements
and the centered part of the fourth.

Every \(c_x\) lies in the disk centered at \(e^{im_x}\) of relative radius
\(q_{\rm ph}<1\).  Choose its logarithm by continuation from zero load.
The quasi-local phase-density and positive-parent response bounds make the
rooted first and second background variations summable; the exact
mean-derivative formulas are
\eqref{eq:v4v66-first-mean-mark}--%
\eqref{eq:v4v66-second-mean-mark}.  The V43 cocycle sums these scalar
logarithms and cancels them from local ratios.  V54--V56 give the two
marked centered collision sums.
\end{proof}

\subsection{One-mode tests}
\label{sec:v4v71-tests}

The following models distinguish a usable complex hard gap from a real
zero crossing.

\begin{proposition}[A complex one-mode certificate]
\label{prop:v4v71-complex-test}
Let
\[
 S_{\rm b}(x)=\frac\kappa2x^2,
 \qquad
 A(x)=m+iyx,
 \qquad m>0.
\]
Then \(A(x)\ne0\) for every real \(x\),
\begin{align}
 \frac{d^2}{dx^2}\log|A(x)|
 &=
 \frac{y^2(m^2-y^2x^2)}
      {(m^2+y^2x^2)^2}
 \le\frac{y^2}{m^2},
\label{eq:v4v71-toy-modulus}\\
 \left|\frac{d}{dx}\arg A(x)\right|
 &=
 \frac{|y|m}{m^2+y^2x^2}
 \le\frac{|y|}{m}.
\label{eq:v4v71-toy-phase}
\end{align}
Thus \(y^2/m^2<\kappa\) gives a positive curvature margin and a local phase
load below \(\sqrt\kappa\).  Moreover
\(\mathbb E_{\nu_+}e^{i\arg A}\) has positive real part.
\end{proposition}

\begin{proof}
Differentiate
\[
 \log|A|=\frac12\log(m^2+y^2x^2),
 \qquad
 \arg A=\arctan(yx/m).
\]
This proves the two bounds.  The positive law is even in \(x\), the sine of
the phase is odd, and
\[
 \cos(\arg A)=\frac{m}{\sqrt{m^2+y^2x^2}}>0.
\]
Its expectation is therefore strictly positive.
\end{proof}

\begin{proposition}[A real crossing destroys the certificate]
\label{prop:v4v71-real-crossing}
For
\[
 A(x)=m+yx,\qquad m>0,\quad y\ne0,
\]
the determinant vanishes at \(x=-m/y\).  No uniform resolvent bound
\eqref{eq:v4v71-resolvent-bound} holds on the real field line, and the phase
jumps by \(\pi\).  Hence neither the smooth modulus Hessian theorem nor the
phase Lipschitz theorem applies globally.
\end{proposition}

\begin{proof}
The zero is explicit.  The inverse norm diverges as the zero is approached,
and the sign of the real determinant changes across it.  A continuous real
phase representative therefore cannot exist on the whole line.
\end{proof}

\begin{assessment}[Meaning of the tests]
\label{ass:v4v71-tests}
The complex test does not model the full chiral Standard Model
determinant, and the real test is not a no-go theorem for that determinant.
They prove that order-one Yukawa size alone decides nothing: the relevant
data are the hard homotopy gap, relative modulus Hessian, and normalized
phase-gradient load.
\end{assessment}

\subsection{Marginal shell carrier certificate}
\label{sec:v4v71-certificate}

\begin{theorem}[Sufficient one-shell marginal carrier theorem]
\label{thm:v4v71-marginal-carrier}
A rescaled marginal shell supplies the positive/phase inputs assumed in V70
if:
\begin{enumerate}
\item the hard homotopy satisfies
      Definition~\ref{def:v4v71-hard-certificate};
\item \(\eta_{\rm det}<\kappa_{\rm b}\), with a strict weighted commutator
      margin;
\item its phase logarithm has the exact local density
      \eqref{eq:v4v71-local-density} and satisfies
      \eqref{eq:v4v71-phase-KP-condition};
\item the scalar determinant-line branch and its first two background
      marks are consistent between adjacent shells; and
\item the remaining marginal bosonic and gauge factors preserve the
      positive conditional-product or master representation required in
      V70.
\end{enumerate}
Under these conditions the modulus carrier is strongly log-concave, the
marginal phase normalization is zero-free with a thermodynamic pressure,
and the V70 hard residual may be attached without changing its positive
shell power.
\end{theorem}

\begin{proof}
Items one and two give
Corollary~\ref{cor:v4v71-positive-LSI}.  Item three and
Theorem~\ref{thm:v4v71-phase-pressure} give the zero-free phase carrier.
Item four identifies its scalar logarithms with the V43 shell cocycle
rather than independent branches.  Item five supplies the exact
conditional law used by the V70 residual cluster theorem.  The determinant
quotient \eqref{eq:v4v70-det-factor} then attaches the summable residual
without altering any marginal factor.
\end{proof}

\begin{firewall}[The certificate has not been verified for the full SM shell]
\label{fw:v4v71-SM-open}
V71 gives checkable inequalities but does not prove them for the complete
chiral Standard Model operator.  Gauge zero modes, Higgs-induced hard
crossings, CKM/CP phase structure, anomaly-compatible shell gluing, and the
hypercharge direction can violate or obstruct the required uniform
homotopy and phase-load bounds.
\end{firewall}

\subsection{V71 result ledger and V72 target}
\label{sec:v4v71-conclusion}

V71 derives the exact determinant Hessian and the relative curvature loss
\(C(M_{\rm h}J_2+M_{\rm h}^2J_1^2)\).  A strict margin produces a positive
log-Sobolev shell carrier and a local Witten kernel.  The homotopy trace
defines exact quasi-local phase densities.  Poincare variance then proves
that each local phase mean is zero-free after its arbitrary mean angle is
rotated away, with control by \(L^2/(2\kappa_+)\).  Centered normalized
phase factors enter V53, while their scalar means enter V43.  The resulting
certificate allows a canonically order-one marginal coupling only when its
numerical gap, curvature, and phase-gradient inequalities hold.

V72 should test the hard homotopy hypothesis on the actual chiral block
structure.  It should first separate vectorlike singular-value protection
from genuinely chiral Yukawa mixing and prove an exact finite-dimensional
criterion for
\(\inf_{t,\phi}s_{\min}(A_t(\phi))>0\).  If an affine Higgs direction can
force a zero, that direction must be assigned to the soft spectral sector
before the determinant logarithm is formed.

\clearpage
\section{V72: singular-value protection, affine Higgs pencils, and mandatory soft transfer}
\label{sec:v4v72}

V71 reduces the order-one marginal determinant problem to a uniform hard
homotopy gap, a relative modulus Hessian margin, and a normalized phase
load.  Earlier V14 and V24--V26 already construct fixed hard complements
for toron and small residual perturbations.  They do not decide whether an
unbounded affine Higgs direction can drive the chiral fermion matrix through
zero.

This note gives exact finite-dimensional tests for that question.  A common
rotated accretive cone supplies structural singular-value protection.  In
one affine Higgs direction, absence of a crossing is exactly a resolvent
condition for a generalized matrix pencil; for a nonnormal chiral block the
relevant quantity is pseudospectral, not merely the distance of eigenvalues
from the real axis.  If the test fails, the crossing singular cluster must
be transferred to the soft fermion packet before a determinant logarithm
is formed.

\subsection{Two structural hard-gap criteria}
\label{sec:v4v72-structural}

For a square complex matrix \(A\), let
\(s_{\min}(A)=\inf_{\|v\|=1}\|Av\|\).

\begin{lemma}[Norm perturbation of a singular-value gap]
\label{lem:v4v72-Weyl}
If \(s_{\min}(A_0)\ge\delta_0\), then
\begin{equation}
 s_{\min}(A_0+V)\ge\delta_0-\|V\|.
\label{eq:v4v72-Weyl}
\end{equation}
In particular \(\|V\|\le q<\delta_0\) gives a uniform hard inverse.
\end{lemma}

\begin{proof}
For every unit vector \(v\),
\[
 \|(A_0+V)v\|
 \ge\|A_0v\|-\|Vv\|
 \ge\delta_0-\|V\|.
\]
Take the infimum.
\end{proof}

This is the non-self-adjoint singular-value version of the fixed-complement
V26 estimate.  It controls a residual perturbation but cannot control an
unbounded marginal Higgs multiplication operator.

\begin{theorem}[Rotated accretivity protects a hard gap]
\label{thm:v4v72-accretive}
Suppose there is one angle \(\beta\) and \(\delta>0\) such that, for every
field and every homotopy parameter,
\begin{equation}
 \frac12\left(
 e^{-i\beta}A_t(\phi)+e^{i\beta}A_t(\phi)^*
 \right)
 \ge\delta\,1.
\label{eq:v4v72-accretive}
\end{equation}
Then
\begin{equation}
 s_{\min}(A_t(\phi))\ge\delta,
 \qquad
 \|A_t(\phi)^{-1}\|\le\delta^{-1}
\label{eq:v4v72-accretive-gap}
\end{equation}
uniformly in \(t,\phi\).
\end{theorem}

\begin{proof}
For every unit vector \(v\),
\[
 \|A_tv\|
 \ge|\langle v,A_tv\rangle|
 \ge
 \operatorname{Re}
 \langle v,e^{-i\beta}A_tv\rangle
 \ge\delta.
\]
Take the infimum.  The inverse bound follows.
\end{proof}

\begin{firewall}[Endpoint positivity is insufficient]
\label{fw:v4v72-one-angle}
The angle in \eqref{eq:v4v72-accretive} must be common to the complete
field range and determinant-line homotopy.  Choosing a different phase
rotation after inspecting each eigenvalue or each field does not define one
logarithm branch and does not prove V71's homotopy certificate.
\end{firewall}

\subsection{The exact one-direction pencil test}
\label{sec:v4v72-pencil}

Let
\begin{equation}
 A(s)=A_0+sB,\qquad s\in\mathbb R,
\label{eq:v4v72-pencil}
\end{equation}
and first suppose \(B\) is invertible.  Put \(T=B^{-1}A_0\) and define the
real-axis pseudospectral margin
\begin{equation}
 \delta_{\mathbb R}(T)
 :=
 \left[
 \sup_{s\in\mathbb R}\|(T+s1)^{-1}\|
 \right]^{-1},
\label{eq:v4v72-real-margin}
\end{equation}
with value zero if \(T+s1\) is singular for some real \(s\).

\begin{theorem}[Affine pencil criterion]
\label{thm:v4v72-pencil}
For invertible \(B\), the following are equivalent:
\begin{enumerate}
\item \(A(s)\) is invertible for every real \(s\);
\item \(\operatorname{spec}(B^{-1}A_0)\cap\mathbb R=\varnothing\);
\item \(\delta_{\mathbb R}(B^{-1}A_0)>0\).
\end{enumerate}
When they hold,
\begin{equation}
 \inf_{s\in\mathbb R}s_{\min}(A(s))
 \ge
 s_{\min}(B)\,
 \delta_{\mathbb R}(B^{-1}A_0)>0.
\label{eq:v4v72-pencil-gap}
\end{equation}
\end{theorem}

\begin{proof}
The factorization
\[
 A(s)=B(T+s1)
\]
shows that \(A(s)\) is singular exactly when \(-s\) is a real eigenvalue of
\(T\), proving the first two equivalences.  If there is no real eigenvalue,
the resolvent is continuous on the real axis and tends to zero in norm as
\(|s|\to\infty\).  Its supremum is therefore finite, proving the third
condition.  Conversely a finite supremum excludes a pole.

Finally,
\[
 s_{\min}(A(s))
 \ge s_{\min}(B)s_{\min}(T+s1)
 =
 \frac{s_{\min}(B)}{\|(T+s1)^{-1}\|}.
\]
Take the infimum over \(s\).
\end{proof}

\begin{firewall}[Eigenvalue distance is not the nonnormal margin]
\label{fw:v4v72-pseudospectral}
For a normal \(T\),
\(\delta_{\mathbb R}(T)\) is the distance of its spectrum from the real
axis.  For a nonnormal chiral block, the resolvent norm can be much larger
than the reciprocal spectral distance.  V71 requires the resolvent margin
\eqref{eq:v4v72-real-margin}; an eigenvalue plot alone is not a hard-gap
proof.
\end{firewall}

\subsection{Several real Higgs components}
\label{sec:v4v72-multipencil}

For an affine Higgs multiplet \(h\in\mathbb R^r\), write
\begin{equation}
 A(h)=A_0+\sum_{\ell=1}^rh_\ell B_\ell
 =A_0+\rho B(n),
 \qquad
 h=\rho n,\quad n\in S^{r-1}.
\label{eq:v4v72-multipencil}
\end{equation}

\begin{theorem}[Uniform directional pencil certificate]
\label{thm:v4v72-directional}
Assume:
\begin{enumerate}
\item \(B(n)\) is invertible for every \(n\in S^{r-1}\);
\item \(\inf_ns_{\min}(B(n))\ge b_0>0\); and
\item
\begin{equation}
 \inf_{n\in S^{r-1}}
 \delta_{\mathbb R}(B(n)^{-1}A_0)
 \ge d_0>0.
\label{eq:v4v72-directional-margin}
\end{equation}
\end{enumerate}
Then
\begin{equation}
 \inf_{h\in\mathbb R^r}s_{\min}(A(h))
 \ge b_0d_0.
\label{eq:v4v72-multipencil-gap}
\end{equation}
The same conclusion holds along a homotopy if the three constants are
uniform in its parameter.
\end{theorem}

\begin{proof}
For \(h=0\), use any direction \(n\); for \(h\ne0\), apply
Theorem~\ref{thm:v4v72-pencil} to the direction \(n=h/|h|\) at
\(s=\rho\).  The lower bound is uniform by the second and third
hypotheses.  Adding a homotopy parameter only adds one compact infimum.
\end{proof}

If \(B(n)\) has a kernel for some \(n\), the theorem does not declare a
zero.  It says that this uniform certificate is unavailable.  The kernel
and its coupling to the complementary block must be treated by an exact
Schur or soft-transfer analysis.

\begin{corollary}[A real generalized eigenvalue forces a crossing]
\label{cor:v4v72-real-generalized}
If, for some direction \(n\), \(B(n)\) is invertible and
\(B(n)^{-1}A_0\) has a real eigenvalue \(\lambda\), then
\[
 A(-\lambda n)
\]
is singular.  No global V71 determinant logarithm exists on the complete
real Higgs line in that direction.
\end{corollary}

\begin{proof}
This is the failed direction of
Theorem~\ref{thm:v4v72-pencil}.
\end{proof}

\subsection{Chiral block interpretation}
\label{sec:v4v72-chiral}

On a finite hard block, a Yukawa-coupled chiral matrix has schematic form
\begin{equation}
 A(h)=
 \begin{pmatrix}
 D_L&\mathsf M(h)\\
 \widetilde{\mathsf M}(h)&D_R
 \end{pmatrix},
 \qquad
 \mathsf M(h)=\sum_\ell h_\ell Y_\ell,
\label{eq:v4v72-chiral-block}
\end{equation}
with the anomaly-compatible lattice realization fixing the precise
adjoints and projectors.

\begin{proposition}[Vectorlike Hermitization criterion]
\label{prop:v4v72-hermitization}
Suppose there is a fixed unitary involution \(\Gamma=\Gamma^*=\Gamma^{-1}\)
such that
\begin{equation}
 H(h):=\Gamma A(h)=H(h)^*.
\label{eq:v4v72-hermitization}
\end{equation}
Then
\begin{equation}
 s_{\min}(A(h))
 =
 \min_{\lambda\in\operatorname{spec}H(h)}|\lambda|.
\label{eq:v4v72-hermitized-gap}
\end{equation}
A fixed V25-type soft template and the relative form bound
\begin{equation}
 \|P_{\rm h}[H(h)-H_0]P_{\rm h}\|<\delta_0
\label{eq:v4v72-hermitized-form}
\end{equation}
preserve a reference hard gap \(\delta_0\) by V26.
\end{proposition}

\begin{proof}
Multiplication by \(\Gamma\) is unitary, so \(A\) and \(H\) have the same
singular values.  Since \(H\) is self-adjoint, its singular values are the
absolute values of its eigenvalues.  Apply
Lemma~\ref{lem:v4v72-Weyl} on the fixed hard complement.
\end{proof}

A genuinely chiral determinant need not admit one such Hermitization on its
physical Weyl space.  In that case the pencil and singular-value tests,
rather than a signed eigenvalue gap, are the correct hard data.

\begin{firewall}[Hermitizing a doubled operator can erase the phase]
\label{fw:v4v72-no-fake-Hermitization}
The block matrix
\[
 \begin{pmatrix}0&A\\A^*&0\end{pmatrix}
\]
is always self-adjoint, but its determinant contains
\(|\det A|^2\) and loses the chiral phase.  It may be used to locate
singular values, not as a replacement for the anomaly-normalized
determinant in V71.
\end{firewall}

\subsection{Exact singular-cluster transfer}
\label{sec:v4v72-transfer}

For a square \(A\), let
\[
 P_{\rm s}
 =
 \boldsymbol1_{[0,\delta^2)}(A^*A),
 \qquad
 Q_{\rm s}
 =
 \boldsymbol1_{[0,\delta^2)}(AA^*),
\]
and put \(P_{\rm h}=1-P_{\rm s}\),
\(Q_{\rm h}=1-Q_{\rm s}\).

Because \(A\) is square, \(A^*A\) and \(AA^*\) have the same
singular spectrum with multiplicity, including equal nullities.  Hence
\begin{equation}
 \operatorname{rank}P_{\rm s}=\operatorname{rank}Q_{\rm s},
 \qquad
 \operatorname{rank}P_{\rm h}=\operatorname{rank}Q_{\rm h}.
\label{eq:v4v72-rank-match}
\end{equation}

\begin{lemma}[Singular subspaces intertwine]
\label{lem:v4v72-intertwine}
One has
\begin{equation}
 AP_{\rm s}=Q_{\rm s}A,
 \qquad
 AP_{\rm h}=Q_{\rm h}A.
\label{eq:v4v72-intertwine}
\end{equation}
The hard map
\[
 A_{\rm h}:=Q_{\rm h}AP_{\rm h}
\]
has
\begin{equation}
 s_{\min}(A_{\rm h})\ge\delta.
\label{eq:v4v72-transferred-gap}
\end{equation}
\end{lemma}

\begin{proof}
The polar decomposition \(A=U(A^*A)^{1/2}\) maps every positive
singular-value spectral subspace of \(A^*A\) to the corresponding subspace
of \(AA^*\).  The zero kernels obey the same intertwining.  On
\(\operatorname{Ran}P_{\rm h}\),
\[
 \|Av\|^2=\langle v,A^*Av\rangle\ge\delta^2\|v\|^2.
\]
\end{proof}

\begin{theorem}[Hard determinant only after soft transfer]
\label{thm:v4v72-transfer}
On a chart where the singular cluster below \(\delta\) has fixed rank and
the threshold \(\delta\) lies in a spectral collar, transfer the complete
equal-rank left/right pair
\((P_{\rm s},Q_{\rm s})\) to the soft fermion packet.  Then:
\begin{enumerate}
\item the hard determinant is formed only from
      \(A_{\rm h}=Q_{\rm h}AP_{\rm h}\) and has inverse norm at most
      \(\delta^{-1}\);
\item the soft--hard basis change and its determinant-line orientation are
      recorded once;
\item every soft--hard coupling generated by use of an enlarged fixed
      template is retained in one Schur/Feshbach response; and
\item a zero inside the transferred soft block is not included in the hard
      logarithm.
\end{enumerate}
\end{theorem}

\begin{proof}
Lemma~\ref{lem:v4v72-intertwine} gives the hard inverse, and
\eqref{eq:v4v72-rank-match} makes it a square map between equal-dimensional
hard spaces.  A singular-value chart supplies continuous left and right
frames for the separated cluster; their determinant-line transition is one
change of basis.  If an enlarged fixed template rather than the exact
reducing singular subspace is used, its left and right ranks must be
matched, it must contain the exact soft pair, and the operator is written in
soft/hard blocks.  The exact block \(LDU\) factorization applies whenever
the resulting square hard block is invertible.  This creates one Schur
complement and one hard determinant.  Any unmatched kernel is therefore
retained wholly as soft zero-mode/determinant-line data.  The soft Berezin
variables remain unintegrated, so a zero of their finite polynomial is not
divided by or differentiated in the hard trace logarithm.
\end{proof}

\begin{firewall}[Soft transfer does not remove the zero]
\label{fw:v4v72-soft-not-solved}
Transfer prevents a false hard resolvent estimate.  It does not prove that
the remaining soft Grassmann integral is nonzero, reflection positive, or
controlled uniformly in volume.  Those are soft-sector and anomaly/phase
problems.
\end{firewall}

\subsection{An exact first-failure rule}
\label{sec:v4v72-first-failure}

Fix a deterministic sequence of nested hard charts with thresholds
\(\delta_0>\delta_1>\cdots\).  For a field configuration \(\phi\),
let \(s_k(t,\phi)\) be the ordered singular values along the complete
determinant homotopy.  Assign index \(k\) to the first chart for which
\begin{equation}
 \inf_{t\in[0,1]}s_k(t,\phi)<2\delta_j,
\label{eq:v4v72-first-failure-rule}
\end{equation}
using a fixed half-open convention at equality.

\begin{theorem}[One owner for every potentially crossing mode]
\label{thm:v4v72-first-failure}
At finite cutoff the first-failure rule partitions all potentially
crossing singular branches.  After a branch is transferred, it never
re-enters a later hard determinant.  On each remaining hard chart the
homotopy inverse has the certified lower singular value, and no determinant
zero is charged both as a soft mode and as a hard tail event.
\end{theorem}

\begin{proof}
A finite matrix has finitely many ordered singular-value functions, and
each is continuous in the matrix entries even when singular values cross.
For every \((k,\phi)\), the ordered thresholds and the half-open equality
rule give at most one first failed index.  Delete every assigned index from
all later hard complements.  A branch retained on chart \(j\) obeys its
homotopy threshold by the negation of
\eqref{eq:v4v72-first-failure-rule}.  Since the soft and hard projectors are
complementary, a transferred index has one ancestry.  This argument is
pointwise in \(\phi\); it asserts no finite global stratification of the
noncompact field space.
\end{proof}

\begin{firewall}[No finite atlas claim on the noncompact Higgs field]
\label{fw:v4v72-noncompact}
V14 obtains a finite atlas from compact toron moduli.  The Higgs field range
is noncompact, so V72 does not claim a cutoff-uniform finite chart cover or
uniform transition constants.  The first-failure rule is an exact
finite-cutoff assignment.  A continuum theorem still needs a locally
summable soft-rank/transition ledger.
\end{firewall}

\subsection{Marginal hard-gap decision tree}
\label{sec:v4v72-decision}

\begin{theorem}[Admissible inputs to V71]
\label{thm:v4v72-V71-input}
A marginal chiral shell may use the V71 determinant certificate only on a
sector for which one of the following has been proved on the complete
homotopy:
\begin{enumerate}
\item the common rotated accretivity estimate
      \eqref{eq:v4v72-accretive};
\item a directional pencil margin such as
      \eqref{eq:v4v72-directional-margin};
\item a valid fixed Hermitization with a V26 relative hard-gap estimate; or
\item a prior singular-cluster transfer leaving the hard bound
      \eqref{eq:v4v72-transferred-gap}.
\end{enumerate}
If none holds, the sector cannot be put into
\(A_j^{\rm marg}\) under the V71 hypotheses.
\end{theorem}

\begin{proof}
The four alternatives give, respectively,
Theorem~\ref{thm:v4v72-accretive},
Theorem~\ref{thm:v4v72-directional},
Proposition~\ref{prop:v4v72-hermitization}, and
Theorem~\ref{thm:v4v72-transfer}.  Each supplies the uniform inverse used
in Definition~\ref{def:v4v71-hard-certificate}.  Without one of these or an
equivalent singular-value proof, invertibility is an unsupported
assumption.
\end{proof}

\begin{assessment}[What V72 decides]
\label{ass:v4v72-status}
V72 turns the hard homotopy gate into finite-dimensional tests.  A real
generalized eigenvalue forces an affine Higgs crossing.  A nonnormal chiral
block requires a real-axis pseudospectral margin.  Structural accretivity
or Hermitization may protect a sector.  Every unresolved crossing cluster
has an exact soft owner and is excluded from the hard logarithm.  The note
does not prove that the complete Standard Model chiral matrix satisfies one
hard alternative, nor does it control the transferred soft determinant.
\end{assessment}

\subsection{V72 result ledger and V73 target}
\label{sec:v4v72-conclusion}

V72 proves the rotated-accretivity hard gap, the exact affine pencil
criterion, its uniform multiplet version, the valid scope of vectorlike
Hermitization, and the left/right singular-cluster transfer.  These results
supply a rigorous decision tree for the V71 marginal determinant
certificate and prohibit forming a hard logarithm across a Higgs-induced
zero.

V73 should construct the soft fermion packet created by
Theorem~\ref{thm:v4v72-transfer}.  The first target is finite cutoff:
retain its Grassmann variables, Schur complement, determinant-line
orientation, and couplings to the marginal Higgs/gauge shell in one local
polynomial, and prove that its coefficient norm and transferred-rank density
are compatible with the V4 marginal projection.  Zero-freeness of its
integrated scalar normalization must remain a separate gate.
\clearpage
\section{V73: exact soft-fermion Schur packet and the locality gate}
\label{sec:v4v73}

V31--V34 construct and contract local Yukawa Grassmann activities relative
to an already invertible quadratic fermion parent.  V72 addresses a
different situation: a singular cluster has been removed from the hard
determinant because its singular value may vanish along the marginal
Higgs--gauge homotopy.  Reusing the V31 normalized Gaussian formula on that
cluster would divide by the determinant whose possible zero caused the
transfer.

This note supplies the missing algebraic bridge.  At finite cutoff, the
hard Grassmann variables can be integrated exactly and only once, leaving a
finite soft Schur polynomial with its determinant-line orientation.  A
spectral collar makes the soft projector exponentially local through a
Riesz--Combes--Thomas estimate.  These facts are enough to identify a
local marginal interface, but not to manufacture a normalized soft
Gaussian measure at a zero.  The latter remains an unintegrated Grassmann
sector until a separate zero-mode saturation or nonvanishing theorem is
proved.

\subsection{Oriented left/right block coordinates}
\label{sec:v4v73-coordinates}

Let \(E_{\rm R}\) and \(E_{\rm L}\) be equal-dimensional complex Hilbert
spaces and let \(A:E_{\rm R}\to E_{\rm L}\).  Choose orthogonal
decompositions
\[
 E_{\rm R}=E_{{\rm R},{\rm s}}\oplus E_{{\rm R},{\rm h}},
 \qquad
 E_{\rm L}=E_{{\rm L},{\rm s}}\oplus E_{{\rm L},{\rm h}},
\]
whose soft and hard dimensions match on the two sides.  In oriented
orthonormal frames write
\begin{equation}
 A=
 \begin{pmatrix}
 A_{\rm ss}&A_{\rm sh}\\
 A_{\rm hs}&A_{\rm hh}
 \end{pmatrix}.
\label{eq:v4v73-block-A}
\end{equation}
The column Grassmann variable is
\(\chi=(\chi_{\rm s},\chi_{\rm h})\), and its independent barred partner is
\(\bar\chi=(\bar\chi_{\rm s},\bar\chi_{\rm h})\).

The frame changes from the physical lattice variables to these left/right
coordinates carry one Berezin Jacobian
\(\mathcal J_{\rm LR}\).  We do not suppress it: its transition functions
are the determinant-line data of the chiral measure.  All formulas below
use the ordered hard Berezin measure for which
\[
 \int
 e^{-\bar\chi_{\rm h}D\chi_{\rm h}}\,
 d\bar\chi_{\rm h}\,d\chi_{\rm h}
 =\det D.
\]

\begin{theorem}[Exact hard Berezin integration]
\label{thm:v4v73-Berezin-Schur}
If \(A_{\rm hh}\) is invertible, define
\begin{equation}
 S_{\rm s}
 :=
 A_{\rm ss}-A_{\rm sh}A_{\rm hh}^{-1}A_{\rm hs}.
\label{eq:v4v73-soft-Schur}
\end{equation}
Then the following identity holds in the Grassmann algebra generated by the
unintegrated soft variables:
\begin{equation}
 \int
 e^{-\bar\chi A\chi}\,
 d\bar\chi_{\rm h}\,d\chi_{\rm h}
 =
 \det(A_{\rm hh})\,
 e^{-\bar\chi_{\rm s}S_{\rm s}\chi_{\rm s}}.
\label{eq:v4v73-Berezin-Schur}
\end{equation}
Consequently,
\begin{equation}
 \det A=\det(A_{\rm hh})\det(S_{\rm s})
\label{eq:v4v73-determinant-Schur}
\end{equation}
as a polynomial identity, including configurations for which
\(\det(S_{\rm s})=0\).
\end{theorem}

\begin{proof}
Set
\[
 \eta_{\rm h}
 =
 \chi_{\rm h}+A_{\rm hh}^{-1}A_{\rm hs}\chi_{\rm s},
 \qquad
 \bar\eta_{\rm h}
 =
 \bar\chi_{\rm h}+
 \bar\chi_{\rm s}A_{\rm sh}A_{\rm hh}^{-1}.
\]
A direct multiplication gives
\[
 \bar\chi A\chi
 =
 \bar\eta_{\rm h}A_{\rm hh}\eta_{\rm h}
 +
 \bar\chi_{\rm s}S_{\rm s}\chi_{\rm s}.
\]
Berezin integration is invariant under triangular translations by the
unintegrated variables.  Integrating the first term gives
\(\det A_{\rm hh}\), proving
\eqref{eq:v4v73-Berezin-Schur}.  Integrating the remaining soft variables
gives \eqref{eq:v4v73-determinant-Schur}.  Since the computation takes
place in a finite Grassmann algebra, no inverse of \(S_{\rm s}\) and no
continuity argument at its zeros is used.
\end{proof}

\begin{corollary}[Exact singular subspaces have no cross block]
\label{cor:v4v73-exact-SVD}
For the V72 projectors
\((P_{\rm s},Q_{\rm s})\), choose
\[
 E_{{\rm R},{\rm s}}=\operatorname{Ran}P_{\rm s},
 \qquad
 E_{{\rm L},{\rm s}}=\operatorname{Ran}Q_{\rm s}.
\]
Then
\[
 A_{\rm sh}=A_{\rm hs}=0,\qquad
 S_{\rm s}=Q_{\rm s}AP_{\rm s},
 \qquad
 A_{\rm hh}=Q_{\rm h}AP_{\rm h}.
\]
The hard determinant in \eqref{eq:v4v73-Berezin-Schur} is therefore exactly
the V72 hard determinant.
\end{corollary}

\begin{proof}
The intertwining identities
\eqref{eq:v4v72-intertwine} imply
\(Q_{\rm s}AP_{\rm h}=0\) and
\(Q_{\rm h}AP_{\rm s}=0\).  The remaining statements follow from the block
definitions.
\end{proof}

\begin{firewall}[The soft factor is not normalized]
\label{fw:v4v73-no-soft-normalization}
Equation \eqref{eq:v4v73-Berezin-Schur} leaves
\(e^{-\bar\chi_{\rm s}S_{\rm s}\chi_{\rm s}}\) as a polynomial.  Dividing
by \(\det S_{\rm s}\), introducing \(S_{\rm s}^{-1}\), or taking
\(\log\det S_{\rm s}\) is forbidden until a separate nonvanishing theorem
has been proved.  In particular, V73 does not reuse the normalized V31
Gaussian identity on a crossing cluster.
\end{firewall}

\subsection{Coefficient bounds for a finite soft packet}
\label{sec:v4v73-coefficients}

Let \(r=\dim E_{{\rm R},{\rm s}}=\dim E_{{\rm L},{\rm s}}\), and use the
V31 coefficient norm on the \(2r\) soft generators.

\begin{lemma}[Soft bilinear coefficient norm]
\label{lem:v4v73-bilinear-norm}
For
\(\mathcal V_{\rm s}=\bar\chi_{\rm s}S_{\rm s}\chi_{\rm s}\),
\begin{equation}
 \|\mathcal V_{\rm s}\|_{\mathfrak G}
 =
 \sum_{a,b=1}^{r}|(S_{\rm s})_{ab}|
 \le r^{3/2}\|S_{\rm s}\|.
\label{eq:v4v73-bilinear-norm}
\end{equation}
If all four blocks of \(A\) have norm at most \(M\) and
\(\|A_{\rm hh}^{-1}\|\le\delta^{-1}\), then
\begin{equation}
 \|S_{\rm s}\|
 \le M+M^2\delta^{-1}.
\label{eq:v4v73-Schur-op-bound}
\end{equation}
Consequently,
\begin{align}
 \|e^{-\mathcal V_{\rm s}}\|_{\mathfrak G}
 &\le
 \exp\!\left[
 r^{3/2}(M+M^2\delta^{-1})
 \right],
\label{eq:v4v73-soft-exp-bound}\\
 \|e^{-\mathcal V_{\rm s}}-1\|_{\mathfrak G}
 &\le
 u e^u,\qquad
 u=r^{3/2}(M+M^2\delta^{-1}).
\label{eq:v4v73-soft-activity-bound}
\end{align}
\end{lemma}

\begin{proof}
The equality is the definition of the coefficient norm.  Cauchy--Schwarz
on the \(r^2\) entries gives
\[
 \sum_{a,b}|(S_{\rm s})_{ab}|
 \le r\|S_{\rm s}\|_{\rm HS}
 \le r^{3/2}\|S_{\rm s}\|.
\]
The triangle inequality and submultiplicativity of the operator norm prove
\eqref{eq:v4v73-Schur-op-bound}.  Apply
Lemma~\ref{lem:v4v31-Grassmann-norm}.
\end{proof}

This lemma is volume-uniform for one packet only when its rank is uniformly
bounded.  Applying it once to a soft space whose rank is proportional to
the volume produces an exponential volume constant and proves no polymer
estimate.  Locality, not finite-dimensionality alone, is the next gate.

\subsection{Weighted kernel algebra}
\label{sec:v4v73-weighted}

Let the cutoff lattice be a bounded-degree metric graph with uniform polynomial volume growth.  For a block
kernel \(K=(K_{xy})\), define
\begin{equation}
 \|K\|_\mu
 :=
 \max\left\{
 \sup_x\sum_y e^{\mu d(x,y)}\|K_{xy}\|,
 \sup_y\sum_x e^{\mu d(x,y)}\|K_{xy}\|
 \right\}.
\label{eq:v4v73-weighted-norm}
\end{equation}

\begin{lemma}[Weighted Schur algebra]
\label{lem:v4v73-weighted-algebra}
For \(\mu\ge0\),
\begin{equation}
 \|K_1K_2\|_\mu
 \le\|K_1\|_\mu\|K_2\|_\mu.
\label{eq:v4v73-weighted-product}
\end{equation}
If
\(K^{(\le L)}_{xy}=\boldsymbol1_{\{d(x,y)\le L\}}K_{xy}\), then
\begin{equation}
 \|K-K^{(\le L)}\|_0
 \le e^{-\mu L}\|K\|_\mu.
\label{eq:v4v73-range-tail}
\end{equation}
\end{lemma}

\begin{proof}
The triangle inequality
\(d(x,z)\le d(x,y)+d(y,z)\) bounds the weighted row convolution by the
product of the two weighted row norms; the column estimate is identical.
For \(d(x,y)>L\), insert
\(1\le e^{-\mu L}e^{\mu d(x,y)}\) in each row and column sum.
\end{proof}

\subsection{A spectral collar localizes the singular projector}
\label{sec:v4v73-projector-locality}

Put \(H_{\rm R}=A^*A\) and \(H_{\rm L}=AA^*\).  We state the locality
criterion for either operator \(H\).  Let \(\Gamma\) be a positively
oriented contour enclosing the soft spectrum and no hard spectrum, and
assume
\begin{equation}
 \operatorname{dist}(\Gamma,\operatorname{spec}H)\ge\gamma>0.
\label{eq:v4v73-contour-gap}
\end{equation}
For a base point \(x_0\), let
\((W_{x_0,\mu}f)(x)=e^{\mu d(x,x_0)}f(x)\), with the exponent truncated on a
finite volume if necessary.

\begin{theorem}[Riesz--Combes--Thomas projector bound]
\label{thm:v4v73-Riesz-CT}
Assume uniformly in \(x_0\) that
\begin{equation}
 \left\|
 W_{x_0,\mu}HW_{x_0,\mu}^{-1}-H
 \right\|
 \le\frac{\gamma}{2}.
\label{eq:v4v73-conjugation-small}
\end{equation}
Then the Riesz projector
\begin{equation}
 P_{\rm s}
 =
 \frac{1}{2\pi i}\int_\Gamma(z-H)^{-1}\,dz
\label{eq:v4v73-Riesz-projector}
\end{equation}
obeys
\begin{equation}
 \|(P_{\rm s})_{xy}\|
 \le
 \frac{|\Gamma|}{\pi\gamma}e^{-\mu d(x,y)}.
\label{eq:v4v73-projector-decay}
\end{equation}
Finite range and a uniform kernel bound for \(H\) imply
\eqref{eq:v4v73-conjugation-small} for some \(\mu>0\) depending only on the
range, the bound, and \(\gamma\).
\end{theorem}

\begin{proof}
For \(R(z)=(z-H)^{-1}\), the spectral theorem and
\eqref{eq:v4v73-contour-gap} give \(\|R(z)\|\le\gamma^{-1}\).
The resolvent identity and
\eqref{eq:v4v73-conjugation-small} give
\[
 \left\|
 W_{x_0,\mu}R(z)W_{x_0,\mu}^{-1}
 \right\|
 \le\frac{2}{\gamma}.
\]
Choose \(x_0=y\), take the \(xy\) block, and remove the weights to obtain
\[
 \|R(z)_{xy}\|
 \le\frac{2}{\gamma}e^{-\mu d(x,y)}.
\]
Integrating around \(\Gamma\) proves
\eqref{eq:v4v73-projector-decay}.  If \(H\) has range \(R\), conjugation
multiplies a nonzero block by at most \(e^{\mu R}\).  Hence the left side of
\eqref{eq:v4v73-conjugation-small} tends to zero with \(\mu\), uniformly
under a uniform Schur-kernel bound.
\end{proof}

Apply the theorem to \(H_{\rm R}\) and \(H_{\rm L}\).  The two contours can
be chosen identically because the nonzero spectra agree; for square \(A\)
the zero multiplicities agree as well.

\begin{theorem}[Two marked projector locality]
\label{thm:v4v73-marked-projector}
Let \(H(B)\) have a common contour \(\Gamma\) satisfying
\eqref{eq:v4v73-contour-gap}--%
\eqref{eq:v4v73-conjugation-small} on a background chart.  If
\(D_BH\) and \(D_B^2H\) have uniform weighted norms at a possibly smaller
exponent \(0<\mu'<\mu\), then
\[
 P_{\rm s},\quad D_BP_{\rm s},\quad D_B^2P_{\rm s}
\]
have uniform \(\mu'\)-weighted kernel norms.
\end{theorem}

\begin{proof}
Differentiate \eqref{eq:v4v73-Riesz-projector}.  In directions \(u,v\),
\begin{align*}
 D_BR[u]&=R(D_BH[u])R,\\
 D_B^2R[u,v]
 &=R(D_B^2H[u,v])R\\
 &\quad+R(D_BH[u])R(D_BH[v])R\\
 &\quad+R(D_BH[v])R(D_BH[u])R.
\end{align*}
The pointwise resolvent estimate in the proof of
Theorem~\ref{thm:v4v73-Riesz-CT} gives a uniform weighted kernel norm after
lowering the exponent from \(\mu\) to \(\mu'\); polynomial volume growth
makes the lost exponential summable.  The weighted product algebra and the
finite length of \(\Gamma\) then bound each differentiated contour integral
uniformly.
\end{proof}

\begin{firewall}[A threshold crossing ends the chart]
\label{fw:v4v73-collar}
Theorem~\ref{thm:v4v73-Riesz-CT} requires a spectral collar around the
chosen cluster.  When a singular value reaches the contour, derivatives of
the projector need not remain bounded.  The V72 first-failure rule must
then enlarge the soft cluster and start a new chart; differentiating
through the collision is not permitted.
\end{firewall}

\subsection{Locality of the exact and enlarged soft operators}
\label{sec:v4v73-soft-locality}

\begin{corollary}[Exact soft kernel is local on a collared chart]
\label{cor:v4v73-soft-kernel}
Suppose \(A\), \(D_BA\), and \(D_B^2A\) have uniform weighted kernel bounds,
and both singular projectors satisfy
Theorem~\ref{thm:v4v73-marked-projector}.  Then
\begin{equation}
 S_{\rm s}^{\rm ex}=Q_{\rm s}AP_{\rm s}
\label{eq:v4v73-exact-soft-kernel}
\end{equation}
and its first two background derivatives have uniform weighted kernel
bounds.
\end{corollary}

\begin{proof}
Use the product rule and
Lemma~\ref{lem:v4v73-weighted-algebra} on every resulting product.
\end{proof}

For an enlarged fixed template, suppose its hard block and first two marks
obey the V16 weighted inverse identities.  Applying the same weighted
algebra to
\[
 S_{\rm s}
 =
 A_{\rm ss}-A_{\rm sh}A_{\rm hh}^{-1}A_{\rm hs}
\]
gives the identical two-marked locality conclusion.  Thus exact spectral
coordinates and fixed-template coordinates meet at the same kernel
interface, although their determinant-line transitions differ.

\begin{theorem}[Core--tail presentation of the soft bilinear]
\label{thm:v4v73-core-tail}
Under the hypotheses of Corollary~\ref{cor:v4v73-soft-kernel}, fix
\(L<\infty\) and write
\[
 S_{\rm s}=S_{\rm s}^{(\le L)}+S_{\rm s}^{(>L)}.
\]
The first term has finite range.  For \(j=0,1,2\),
\begin{equation}
 \|D_B^jS_{\rm s}^{(>L)}\|_0
 \le C_j e^{-\mu L}.
\label{eq:v4v73-soft-tail-small}
\end{equation}
The same statement holds for the enlarged-template Schur operator.
\end{theorem}

\begin{proof}
Apply \eqref{eq:v4v73-range-tail} to the kernel and to each of its first two
marks.
\end{proof}

The finite-range core is generally order one and belongs in the marginal
Grassmann parent.  Its exponentially small range tail may enter the
V31--V33 connected activity ledger after choosing \(L\) so that the
displayed constant meets the required tree norm.  This is an exact
core--tail split of the bilinear; it does not integrate the core.

\subsection{Rank density and the frame obstruction}
\label{sec:v4v73-frame}

Let \(q_{\rm f}\) be the fixed number of lattice fermion components per
site.  Every orthogonal soft projector satisfies the local trace bound
\begin{equation}
 0\le
 \operatorname{Tr}(\boldsymbol1_XP_{\rm s}\boldsymbol1_X)
 \le q_{\rm f}|X|.
\label{eq:v4v73-rank-density}
\end{equation}
Thus transferring a singular cluster does not create more than a fixed
ambient fermion fibre per site.

\begin{proof}[Proof of \eqref{eq:v4v73-rank-density}]
Since \(0\le P_{\rm s}\le1\),
\[
 0\le\boldsymbol1_XP_{\rm s}\boldsymbol1_X\le\boldsymbol1_X.
\]
Take the trace.
\end{proof}

The trace bound and projector decay still do not provide an exponentially
localized orthonormal basis.  A nontrivial family of soft subspaces can
carry topological obstruction, and a chiral theory also needs compatible
left/right determinant-line orientation.

\begin{definition}[Local soft-frame certificate]
\label{def:v4v73-frame-certificate}
A background chart has a two-marked local soft-frame certificate if it
provides left and right frames \(L_{\rm s},R_{\rm s}\) such that:
\begin{enumerate}
\item they span \(\operatorname{Ran}Q_{\rm s}\) and
      \(\operatorname{Ran}P_{\rm s}\), with equal ordered rank;
\item the frames, their duals, and their first two background marks have
      uniform exponentially weighted kernels and bounded local overlap;
\item transition matrices on chart overlaps have uniformly summable first
      two marks; and
\item their Berezin Jacobians obey the anomaly-normalized determinant-line
      cocycle fixed in V13 and are entered once in the ancestry ledger.
\end{enumerate}
\end{definition}

\begin{theorem}[Conditional interface with the V4 marginal projection]
\label{thm:v4v73-V4-interface}
On charts satisfying the V72 hard-gap certificate, the V73 spectral-collar
hypotheses, and Definition~\ref{def:v4v73-frame-certificate}, the
transferred sector has:
\begin{enumerate}
\item one nonzero hard determinant with the V71 modulus/phase treatment;
\item one unintegrated finite-fibre soft Grassmann action with a finite-range
      core and the tail estimate
      \eqref{eq:v4v73-soft-tail-small};
\item coefficientwise first and second background marks admissible for the
      V4 marginal projection and the V42 finite-fibre color sweep; and
\item one determinant-line transition on every chart overlap.
\end{enumerate}
No determinant or response term is duplicated.
\end{theorem}

\begin{proof}
Theorem~\ref{thm:v4v73-Berezin-Schur} proves the exact hard/soft ancestry.
The hard inverse is V72.  Theorems
\ref{thm:v4v73-marked-projector} and
\ref{thm:v4v73-core-tail} give the two-marked local kernel decomposition.
The frame certificate converts the basis-free kernels into local finite
Grassmann fibres with summable transitions.  V4 acts coefficientwise, and
V42 shows that a fixed finite coefficient fibre does not change the scalar
color-sweep constant.  The determinant-line Jacobian is assigned by item
4 of the certificate, so no second phase or determinant is introduced.
\end{proof}

\begin{firewall}[Projector locality is not frame triviality]
\label{fw:v4v73-projector-not-frame}
Exponential decay of \(P_{\rm s}\) controls a basis-independent kernel.  It
does not prove Definition~\ref{def:v4v73-frame-certificate}, trivialize the
chiral determinant line, or choose a gauge-covariant fermion measure.
Those are global compatibility statements and cannot be obtained by
selecting singular vectors independently at each field.
\end{firewall}

\begin{firewall}[Finite fibre does not saturate zero modes]
\label{fw:v4v73-no-saturation}
A bounded number of Grassmann generators per site makes coefficient norms
local.  It does not make the final Berezin integral nonzero.  If exact zero
modes remain, a nonzero observable requires the appropriate insertions or
an interaction that saturates every zero-mode generator.  Vacuum
zero-freeness and observable saturation must be proved separately.
\end{firewall}

\subsection{V73 result ledger and V74 target}
\label{sec:v4v73-conclusion}

V73 proves the exact oriented Berezin--Schur transfer, including determinant
ancestry at soft zeros.  It gives finite-packet coefficient bounds, proves
two-marked exponential locality of a collared singular projector, and
derives the finite-range-core/exponentially-small-tail decomposition of the
soft bilinear.  It also identifies the remaining nonlocal step precisely:
projector decay and rank density do not supply a gauge-compatible local
left/right frame or its determinant-line orientation.

V74 must therefore go upstream to the chiral measure problem.  It should
formulate the curvature of the soft determinant line from the projector,
show how the anomaly polynomial appears in its local trace, and determine
what anomaly cancellation proves: cancellation of curvature, trivial
holonomy on admissible gauge-field components, or only a local counterterm.
Only the strongest sufficient statement can discharge
Definition~\ref{def:v4v73-frame-certificate}; otherwise the residual
global phase must remain an explicit continuum gate.
\clearpage
\section{V74: local soft frames, determinant-line factorization, and anomaly redistribution}
\label{sec:v4v74}

V73 leaves two related questions.  First, an exponentially local projection
does not by itself exhibit exponentially local independent Grassmann
coordinates.  Second, the frozen first-series V138--V145 chain already
cancels the complete Standard Model local anomaly, proves trivial faithful
global holonomy, and constructs the cofinal determinant-line transport.
The transferred singular subspace must be fitted into that existing line;
it must not be assigned a second anomaly calculation or an independent
phase convention.

This note proves a projected-seed criterion that constructs local
orthonormal soft frames.  It then gives the exact determinant-line
factorization associated with the V73 Schur decomposition.  The resulting
curvature identity shows that soft transfer redistributes Berry curvature
between hard and soft factors while leaving the anomaly-normalized total
line unchanged.  Total anomaly cancellation makes the product phase
global; it does not require the two factor phases to be separately flat.

\subsection{Kato curvature of the soft determinant line}
\label{sec:v4v74-Kato}

Let \(P(b)\) be a \(C^2\), constant-rank family of orthogonal projections
over a background chart \(\mathcal B\).  Its range bundle carries the Kato
connection
\[
 \nabla^P s=P\,ds.
\]

\begin{lemma}[Range curvature and determinant curvature]
\label{lem:v4v74-Kato-curvature}
The curvature of \(\nabla^P\) is
\begin{equation}
 F^P=P(dP\wedge dP)P.
\label{eq:v4v74-range-curvature}
\end{equation}
The induced connection on \(\det\operatorname{Ran}P\) has curvature
\begin{equation}
 F^{\det P}
 =
 \operatorname{Tr}\!\left(P\,dP\wedge dP\right).
\label{eq:v4v74-det-curvature}
\end{equation}
In coordinates,
\[
 F^{\det P}_{ij}
 =
 \operatorname{Tr}\!\left(P[\partial_iP,\partial_jP]\right).
\]
\end{lemma}

\begin{proof}
Differentiate \(P^2=P\) to obtain \(P(dP)P=0\).  For a range section
\(s=Ps\),
\begin{align*}
 (\nabla_i^P\nabla_j^P-\nabla_j^P\nabla_i^P)s
 &=
 P(\partial_iP)\partial_js-P(\partial_jP)\partial_is\\
 &=
 P[(\partial_iP),(\partial_jP)]Ps,
\end{align*}
which is \eqref{eq:v4v74-range-curvature}.  The curvature of a determinant
connection is the trace of the vector-bundle curvature, giving
\eqref{eq:v4v74-det-curvature}.
\end{proof}

For the V72 left/right soft pair, define
\begin{equation}
 \mathcal L_{\rm s}
 :=
 \det(\operatorname{Ran}Q_{\rm s})
 \otimes
 \det(\operatorname{Ran}P_{\rm s})^*.
\label{eq:v4v74-soft-line}
\end{equation}

\begin{corollary}[Soft left/right Berry curvature]
\label{cor:v4v74-soft-curvature}
The Kato connection on \(\mathcal L_{\rm s}\) has curvature
\begin{equation}
 F_{\rm s}^{\rm K}
 =
 \operatorname{Tr}(Q_{\rm s}\,dQ_{\rm s}\wedge dQ_{\rm s})
 -
 \operatorname{Tr}(P_{\rm s}\,dP_{\rm s}\wedge dP_{\rm s}).
\label{eq:v4v74-soft-curvature}
\end{equation}
Its curvature coefficients, and their first background derivatives, are
local under the two-marked projector bounds of V73.  A second derivative of
the curvature itself would require a third projector mark and is not
claimed here.
\end{corollary}

\begin{proof}
Curvature is additive under tensor products and changes sign under duals,
so Lemma~\ref{lem:v4v74-Kato-curvature} gives the formula.  Each coefficient
is a trace per local diagonal of products of marked exponentially weighted
kernels.  The weighted algebra of
Lemma~\ref{lem:v4v73-weighted-algebra} makes those products summable.
\end{proof}

Let \(P_{\rm h}=1-P_{\rm s}\) and \(Q_{\rm h}=1-Q_{\rm s}\), and define
\(\mathcal L_{\rm h}\) analogously.

\begin{proposition}[Complementary split curvature cancels]
\label{prop:v4v74-complement-curvature}
For the Kato connections induced by the fixed ambient left and right
lattice spaces,
\begin{equation}
 F_{\rm s}^{\rm K}+F_{\rm h}^{\rm K}=0.
\label{eq:v4v74-split-curvature-zero}
\end{equation}
\end{proposition}

\begin{proof}
For one projection \(P\),
\begin{align*}
 F^{\det P}+F^{\det(1-P)}
 &=
 \operatorname{Tr}(dP\wedge dP).
\end{align*}
The \(ij\) coefficient on the right is
\(\operatorname{Tr}[\partial_iP,\partial_jP]=0\).  Apply this identity to
the right projector \(P_{\rm s}\) and the left projector \(Q_{\rm s}\), and
take the left-minus-right combination in
\eqref{eq:v4v74-soft-curvature}.
\end{proof}

\begin{firewall}[Split Berry curvature is not a new anomaly]
\label{fw:v4v74-not-new-anomaly}
Equation \eqref{eq:v4v74-soft-curvature} describes the connection created by
a moving hard/soft decomposition inside the already regulated finite
lattice spaces.  The physical local anomaly is the curvature of the
complete regulator determinant line and was computed for the full Standard
Model representation in the frozen first-series V138 chain.  Soft transfer
cannot change that total curvature; adding
\(F_{\rm s}^{\rm K}\) to it a second time would double count a frame
connection.
\end{firewall}

\subsection{Constructing a local frame from projected seeds}
\label{sec:v4v74-seeds}

The following criterion turns the V73 basis-free projector bound into
actual local Grassmann coordinates.  Let \(\mathcal I\) be a finite index
set with assigned lattice centers and let
\(T:\ell^2(\mathcal I)\to E\) be a local synthesis map.  Assume
\[
 \dim\ell^2(\mathcal I)=\operatorname{rank}P
\]
and set
\begin{equation}
 G:=T^*PT.
\label{eq:v4v74-seed-Gram}
\end{equation}

\begin{theorem}[Projected-seed orthonormalization]
\label{thm:v4v74-projected-seed}
Suppose, uniformly on a background chart,
\begin{equation}
 c_{\rm seed}\,1\le G\le C_{\rm seed}\,1,
 \qquad c_{\rm seed}>0,
\label{eq:v4v74-seed-gap}
\end{equation}
and \(P,T\), together with their first two background derivatives, have
uniform exponentially weighted kernels.  Then
\begin{equation}
 R:=PTG^{-1/2}
\label{eq:v4v74-frame}
\end{equation}
is an orthonormal frame of \(\operatorname{Ran}P\):
\begin{equation}
 R^*R=1,\qquad RR^*=P.
\label{eq:v4v74-frame-identities}
\end{equation}
After possibly lowering the decay exponent, \(R\), \(R^*\), and their first
two background derivatives have uniform exponentially weighted kernels.
\end{theorem}

\begin{proof}
Equation \eqref{eq:v4v74-seed-gap} makes \(G\) invertible.  Therefore
\[
 R^*R=G^{-1/2}T^*PTG^{-1/2}=1.
\]
The range of \(R\) lies in \(\operatorname{Ran}P\).  Both spaces have the
same finite dimension, so \(R\) is onto and \(RR^*=P\).

It remains to prove locality.  The weighted product lemma gives marked
exponential locality of \(G\).  Choose a contour \(\Gamma_G\) enclosing
\([c_{\rm seed},C_{\rm seed}]\) inside the slit plane on which
\(z^{-1/2}\) is holomorphic.  Exponential locality implies
\[
 \sup_{i_0}
 \|W_{i_0,\mu'}GW_{i_0,\mu'}^{-1}-G\|\longrightarrow0
 \quad\text{as }\mu'\downarrow0.
\]
Choose \(\mu'>0\) so the left side is smaller than half the distance from
\(\Gamma_G\) to the spectrum.  The same Neumann-resolvent argument as in
Theorem~\ref{thm:v4v73-Riesz-CT} gives an exponentially weighted bound for
\[
 G^{-1/2}
 =
 \frac{1}{2\pi i}
 \int_{\Gamma_G}z^{-1/2}(z-G)^{-1}\,dz.
\]
Differentiating this formula once and twice inserts, respectively,
\(DG\) and \(D^2G\), with the two ordered quadratic first-derivative terms
exactly as in Theorem~\ref{thm:v4v73-marked-projector}.  The weighted
product algebra bounds every term.  Finally use
\eqref{eq:v4v74-frame} and the product rule.
\end{proof}

\begin{corollary}[Two-sided soft-frame certificate on one chart]
\label{cor:v4v74-two-sided-frame}
If local seed maps \(T_{\rm R}\) and \(T_{\rm L}\) satisfy
Theorem~\ref{thm:v4v74-projected-seed} for
\(P_{\rm s}\) and \(Q_{\rm s}\), respectively, their frames
\(R_{\rm s}\) and \(L_{\rm s}\) satisfy items 1 and 2 of
Definition~\ref{def:v4v73-frame-certificate}.  On an overlap, the frame
transitions
\begin{equation}
 U_{\rm R}=R_{\rm s}^*\widetilde R_{\rm s},
 \qquad
 U_{\rm L}=L_{\rm s}^*\widetilde L_{\rm s}
\label{eq:v4v74-frame-transition}
\end{equation}
are unitary and inherit two-marked exponential locality.
\end{corollary}

\begin{proof}
The frame identities give the spans and equal ranks.  Locality follows from
the theorem.  The overlap matrices are products of local marked kernels;
unitarity follows because both frames are orthonormal bases of the same
range.
\end{proof}

\begin{firewall}[The seed lower bound is substantive]
\label{fw:v4v74-seed-gap}
The trace-density estimate \eqref{eq:v4v73-rank-density} and exponential
decay of \(P\) do not imply \eqref{eq:v4v74-seed-gap} for a predetermined
local seed set.  A vanishing seed Gram eigenvalue means that the chosen
local columns miss a direction of the soft range.  The chart must then
change seeds or enlarge its packet.  Uniform control of
\(c_{\rm seed}^{-1}\) is an analytic atlas requirement, not a consequence
of dimension counting.
\end{firewall}

\subsection{Exact determinant-line factorization}
\label{sec:v4v74-line-factorization}

For a map \(A:E_{\rm R}\to E_{\rm L}\), use the determinant-line convention
\[
 \mathcal L(A):=\det E_{\rm L}\otimes(\det E_{\rm R})^*.
\]
The ordered decompositions in V73 give canonical direct-sum
isomorphisms, with the soft factors placed before the hard factors.

\begin{theorem}[Hard--soft line isomorphism]
\label{thm:v4v74-line-factorization}
There is a canonical oriented isomorphism
\begin{equation}
 \Phi_{\rm sh}:
 \mathcal L(A)
 \longrightarrow
 \mathcal L_{\rm s}\otimes\mathcal L_{\rm h}.
\label{eq:v4v74-line-factorization}
\end{equation}
Under \(\Phi_{\rm sh}\), the determinant section satisfies
\begin{equation}
 \det A
 \longmapsto
 \det S_{\rm s}\otimes\det A_{\rm hh},
\label{eq:v4v74-section-factorization}
\end{equation}
where \(S_{\rm s}\) is the V73 Schur operator.  The line isomorphism remains
defined when \(\det S_{\rm s}=0\).
\end{theorem}

\begin{proof}
The exterior-power identity
\[
 \det(V_{\rm s}\oplus V_{\rm h})
 \simeq
 \det V_{\rm s}\otimes\det V_{\rm h}
\]
on the left and right spaces gives
\eqref{eq:v4v74-line-factorization}; the fixed ordering fixes its Koszul
sign.  The triangular row and column operations in the proof of
Theorem~\ref{thm:v4v73-Berezin-Schur} have determinant one.  The remaining
diagonal blocks are \(S_{\rm s}\) and \(A_{\rm hh}\), which proves
\eqref{eq:v4v74-section-factorization}.  Exterior-power isomorphisms do not
require either diagonal determinant to be nonzero.
\end{proof}

\begin{corollary}[The hard section induces the soft reference line]
\label{cor:v4v74-soft-reference}
On every chart with invertible \(A_{\rm hh}\), its determinant is a
nowhere-zero section of \(\mathcal L_{\rm h}\).  If
\(\sigma_{\rm tot}\) is a nowhere-zero reference section of the complete
anomaly-normalized line, there is a unique nowhere-zero soft reference
section \(\sigma_{\rm s}\) satisfying
\begin{equation}
 \Phi_{\rm sh}(\sigma_{\rm tot})
 =
 \sigma_{\rm s}\otimes\det A_{\rm hh}.
\label{eq:v4v74-soft-reference}
\end{equation}
The scalar coefficient of \(\det S_{\rm s}\) relative to
\(\sigma_{\rm s}\), multiplied by the hard coefficient, is exactly the
complete regulated determinant coefficient relative to
\(\sigma_{\rm tot}\).
\end{corollary}

\begin{proof}
A nonzero vector in a one-dimensional line has a unique tensor quotient.
Apply this pointwise to
\eqref{eq:v4v74-soft-reference}.  The coefficient identity is
\eqref{eq:v4v74-section-factorization}.
\end{proof}

Thus a zero of \(\det S_{\rm s}\) is a zero of a scalar coefficient in a
perfectly well-defined soft line.  It is not a failure of the line
trivialization.

\subsection{What total anomaly cancellation actually implies}
\label{sec:v4v74-anomaly}

The imported result in Section~\ref{sec:v4v1-fermion} supplies the complete
three-generation Standard Model determinant line on the faithful
\(S(U(3)\times U(2))\) gauge group, with vanishing local curvature and
trivial global gauge-anomaly character
\(\Omega_5^{\rm Spin}(BG_{\rm SM})=0\).  Its cofinal phase transport is
part of the frozen first-series determinant-line construction.

\begin{theorem}[Anomaly-free sewing of a transferred cluster]
\label{thm:v4v74-anomaly-sewing}
Let \(\sigma_{\rm tot}\) be the imported parallel total phase section on
one connected admissible background component.  On a V72--V73 chart, use
\eqref{eq:v4v74-soft-reference} to define \(\sigma_{\rm s}\).  Then:
\begin{enumerate}
\item the hard and soft phase connections add to the total connection;
\item in the total parallel frame their connection one-forms satisfy
\begin{equation}
 \alpha_{\rm s}+\alpha_{\rm h}=0;
\label{eq:v4v74-connection-cancel}
\end{equation}
\item for every loop contained in the chart,
\begin{equation}
 \operatorname{Hol}_{\rm s}(\gamma)
 \operatorname{Hol}_{\rm h}(\gamma)=1;
\label{eq:v4v74-holonomy-cancel}
\end{equation}
and
\item the complete determinant phase obtained from the product is the
      already normalized Standard Model phase.
\end{enumerate}
Neither factor is required to have zero curvature or trivial holonomy
separately.
\end{theorem}

\begin{proof}
Equip the factor lines in
\eqref{eq:v4v74-line-factorization} with the tensor-factor connections
induced by the complete regulator connection and the moving split.  Tensor
connections add.  In the parallel total section the total one-form is zero,
which gives \eqref{eq:v4v74-connection-cancel}.  Curvature and loop
holonomy are respectively the exterior derivative and exponential integral
of this additive identity, proving
\eqref{eq:v4v74-holonomy-cancel}.  The last statement follows from the
section factorization
\eqref{eq:v4v74-section-factorization}.  The imported local and global
anomaly theorems apply to the product line, exactly where they were proved.
\end{proof}

\begin{firewall}[Cancellation of the total line does not flatten each factor]
\label{fw:v4v74-factor-not-flat}
It is generally false that Standard Model anomaly cancellation gives
\(F_{\rm s}=F_{\rm h}=0\).  A moving split can have nonzero curvature
\eqref{eq:v4v74-soft-curvature}, compensated by the complementary factor.
Assigning independent parallel phases to both factors can therefore break
the already proved total sewing.
\end{firewall}

\subsection{Changing the transferred rank}
\label{sec:v4v74-rank-change}

Let two nested charts have soft pairs
\[
 (P_{\rm s},Q_{\rm s})
 \subset
 (\widetilde P_{\rm s},\widetilde Q_{\rm s}),
\]
and let \(K_{\rm R}\) and \(K_{\rm L}\) be the added equal-rank right and
left bundles.  Define the finite crossing line
\begin{equation}
 \mathcal C
 :=
 \det K_{\rm L}\otimes(\det K_{\rm R})^*.
\label{eq:v4v74-crossing-line}
\end{equation}

\begin{proposition}[Rank-change transition and cocycle]
\label{prop:v4v74-rank-change}
There are canonical oriented isomorphisms
\begin{equation}
 \widetilde{\mathcal L}_{\rm s}
 \simeq
 \mathcal L_{\rm s}\otimes\mathcal C,
 \qquad
 \mathcal L_{\rm h}
 \simeq
 \mathcal C\otimes\widetilde{\mathcal L}_{\rm h},
\label{eq:v4v74-rank-change}
\end{equation}
under which the two hard--soft factorizations of
\(\mathcal L(A)\) agree.  For three nested transfers, the transition maps
compose exactly.
\end{proposition}

\begin{proof}
Apply the determinant direct-sum identity to
\[
 \operatorname{Ran}\widetilde P_{\rm s}
 =
 \operatorname{Ran}P_{\rm s}\oplus K_{\rm R},
 \qquad
 \operatorname{Ran}\widetilde Q_{\rm s}
 =
 \operatorname{Ran}Q_{\rm s}\oplus K_{\rm L},
\]
and to the reverse decomposition of the hard spaces.  The added determinant
factors combine into \(\mathcal C\).  Associativity of exterior products,
with the fixed soft-before-hard ordering, makes the two total
factorizations equal.  The same associativity proves the three-chart
cocycle.
\end{proof}

This is the finite singular-value version of the relative spectral line
retained in the frozen V139--V142 construction.  The self-adjoint
Hermitization
\[
 \mathscr H(A)=
 \begin{pmatrix}0&A^*\\A&0\end{pmatrix}
\]
identifies the added singular pair with its symmetric finite spectral
window.  It is used here only to label the crossing line; as V72 stresses,
it does not replace the chiral determinant by \(|\det A|^2\).

\begin{theorem}[Completion of the V73 frame certificate]
\label{thm:v4v74-frame-completion}
Assume on a locally finite collared atlas:
\begin{enumerate}
\item right and left projected seeds satisfy the uniform gap
      \eqref{eq:v4v74-seed-gap};
\item overlap and rank-change transition kernels have the two-marked
      summability required in Definition~\ref{def:v4v73-frame-certificate};
      and
\item the total phase is the imported anomaly-normalized Standard Model
      section \(\sigma_{\rm tot}\).
\end{enumerate}
Then Definition~\ref{def:v4v73-frame-certificate} holds.  The soft packet
has local marked frames, exact determinant-line transitions, and one phase
ancestry compatible with the V4 marginal interface.
\end{theorem}

\begin{proof}
Theorem~\ref{thm:v4v74-projected-seed} and
Corollary~\ref{cor:v4v74-two-sided-frame} give the local frames and their
marked fixed-rank transitions.  Proposition~\ref{prop:v4v74-rank-change}
gives the rank-change maps and exact cocycle.  Theorem
\ref{thm:v4v74-anomaly-sewing} fixes their determinant phases from the
single total Standard Model section.  These are precisely the four items
of Definition~\ref{def:v4v73-frame-certificate}.
\end{proof}

\begin{assessment}[What remains after V74]
\label{ass:v4v74-status}
V74 removes a conceptual circularity: the transferred soft packet does not
need an independently anomaly-free phase.  Its phase is the tensor quotient
of the already normalized total line by the nonzero hard determinant.
The remaining frame input is analytic and explicit: a locally finite
spectral-collar atlas with uniformly positive projected-seed Gram matrices
and summable marked transitions.

Even after that certificate, the scalar vacuum coefficient can vanish when
the soft Grassmann generators are unsaturated.  No determinant-line theorem
turns such a zero into a positive normalization.
\end{assessment}

\subsection{V74 result ledger and compulsory V75 target}
\label{sec:v4v74-conclusion}

V74 proves the soft determinant-line Kato curvature, constructs exponentially
local marked frames from gapped projected seeds, factors the complete
determinant line through the exact V73 Schur packet, and proves the
rank-change cocycle.  It imports rather than repeats the established
Standard Model anomaly theorem and shows that anomaly cancellation applies
to the hard--soft product.

V75 is a compulsory companion audit.  After reading the current Yang--Mills
and Navier--Stokes notes, it should test whether their newest localization,
cofinal, or boundary-sector estimates improve the projected-seed atlas.
It should then analyze soft zero-mode saturation in the vacuum and in
physical observables.  The first exact target is a finite Grassmann
selection rule that separates an identically zero vacuum coefficient from
a saturating interaction, followed by a local-rank/tree criterion for the
saturating terms.
\clearpage
\section{V75: compulsory companion audit and exact soft zero-mode saturation}
\label{sec:v4v75}

V74 supplies a conditional local soft frame and sews its determinant line to
the already anomaly-normalized Standard Model line.  It does not show that
the final soft Berezin coefficient is nonzero.  This compulsory audit first
reads the current Yang--Mills and Navier--Stokes programmes, then uses the
relevant structural lessons to formulate the exact zero-mode selection
problem.

\subsection{Stable read-only companion snapshot}
\label{sec:v4v75-audit}

The Yang--Mills file advanced during the first audit attempt: V95 was
replaced before it could be read.  The audit was therefore restarted from
the newest unsuffixed file and accepted only after its SHA--256 digest was
unchanged across the read.  The stable snapshot was
\begin{center}
\begin{tabular}{|p{0.50\textwidth}|r|p{0.31\textwidth}|}
\hline
Companion file & Bytes & SHA--256\\
\hline
\texttt{Renewal\_Yang\_Mills\_Mass\_Gap\_Research\_Notes\_V96.tex}
& \(1\,646\,335\)
& \texttt{8B8FBD64E4001CE5E758221FBF9B236C}\\
&&\texttt{C10FDB06711386A2EA8777EF92090D54}\\
\hline
\texttt{Renewal\_NS\_Stability\_Research\_Notes\_V31.tex}
& \(887\,537\)
& \texttt{F1121B62238AD5491BB7CF03635EA993}\\
&&\texttt{4942EDF573ED8FAAA9EE79E1D38A6696}\\
\hline
\end{tabular}
\end{center}
Both reads ended at one live \verb|\end{document}|.

Yang--Mills V96 restores the complete global quartic marked-tree atlas.  Its
new mechanism is an exact first-full-connectivity coordinate, followed by a
partition of the fourteen proper four-row prefix partitions and a
coordinate-weighted tree sum.  The relevance here is combinatorial: a
saturating Grassmann term must be assigned to one first point at which its
zero-mode packets become connected.  Quoting a support-uniform local bound
and then summing coordinates would repeat the aggregation error that YM V96
repairs.

Navier--Stokes V31 proves that a flat dyadic probability mark becomes an
unweighted cutoff count only after multiplication by the physical scale;
the cost is exactly the associated first moment.  The analogous warning for
V73--V74 is that probabilities of small spectral collars or small seed Gram
eigenvalues do not control the corresponding inverse resolvent/frame loads.
Those loads require a weighted tail moment.

\subsection{Berezin integration is top-coefficient extraction}
\label{sec:v4v75-top-form}

Let \(\mathfrak G(\zeta_1,\ldots,\zeta_N)\) have the fixed ordered top
monomial
\[
 \omega=\zeta_1\cdots\zeta_N.
\]
Normalize Berezin integration by
\begin{equation}
 \int F(\zeta)\,d\zeta_N\cdots d\zeta_1
 =
 [\omega]F,
\label{eq:v4v75-top-integral}
\end{equation}
where \([\omega]F\) is the coefficient of \(\omega\).

\begin{lemma}[Exact monomial-cover rule]
\label{lem:v4v75-cover}
Let \(m_1,\ldots,m_k\) be Grassmann monomials and let
\(\operatorname{supp}m_j\subseteq[N]\) be their generator supports.  Then
\[
 [\omega]m_1\cdots m_k\ne0
\]
if and only if the supports are pairwise disjoint and their union is
\([N]\).  In that case the coefficient is the product of the scalar
coefficients times the sign of the permutation that orders the generators.
\end{lemma}

\begin{proof}
If two supports intersect, a generator is repeated and the product is zero.
If their disjoint union omits a generator, the product has degree below
\(N\).  If the union is all of \([N]\), anticommutation uniquely reorders the
product to \(\omega\), producing only the permutation sign.
\end{proof}

\begin{theorem}[Exact saturation expansion]
\label{thm:v4v75-saturation-expansion}
Let
\[
 \mathcal V=\sum_{X\in\mathcal X}\mathcal V_X
\]
be a finite sum of even positive-degree Grassmann polynomials.  Expand each
finite exponential \(e^{-\mathcal V_X}\) into monomials.  Then
\begin{equation}
 Z_{\rm soft}
 :=
 \int e^{-\mathcal V}\,d\zeta_N\cdots d\zeta_1
 =
 \sum_{\mathcal C\in\operatorname{Cov}(\zeta)}
 \operatorname{sgn}(\mathcal C)\,a(\mathcal C),
\label{eq:v4v75-cover-sum}
\end{equation}
where \(\operatorname{Cov}(\zeta)\) is exactly the finite set of selections
whose monomial supports form a disjoint cover of all \(N\) generators.
Every surviving term occurs once.
\end{theorem}

\begin{proof}
All \(\mathcal V_X\) are even, hence commute, so
\[
 e^{-\mathcal V}=\prod_Xe^{-\mathcal V_X}.
\]
Every exponential terminates.  Distribute the finite product and apply
Lemma~\ref{lem:v4v75-cover} to every resulting monomial selection.  The
choice of one term from each factor, including its scalar term, is unique,
so no surviving cover is duplicated.
\end{proof}

\begin{firewall}[A coefficient bound cannot create a cover]
\label{fw:v4v75-bound-not-cover}
Small Grassmann norm, locality, and rank density can bound the terms in
\eqref{eq:v4v75-cover-sum}; they cannot show that the cover set is nonempty.
Nonvanishing first requires at least one interaction monomial with the
charges and generator content needed to saturate every zero mode.
\end{firewall}

\subsection{Pure quadratic and zero-mode reduction}
\label{sec:v4v75-zero-reduction}

For \(r\) barred and \(r\) unbarred generators, the V73 convention gives
\begin{equation}
 \int e^{-\bar\chi S\chi}\,d\bar\chi\,d\chi=\det S.
\label{eq:v4v75-quadratic-det}
\end{equation}
Thus a singular \(S\) has zero vacuum coefficient in the purely quadratic
soft sector.

Let \(k=\dim\ker S=\dim\ker S^*\).  Choose oriented singular frames in which
\[
 S=
 \begin{pmatrix}
 0&0\\0&D
 \end{pmatrix},
 \qquad D\text{ invertible},
\]
and write the variables as zero modes \((\bar z,z)\) and nonzero modes
\((\bar\eta,\eta)\).

\begin{theorem}[Exact interacting zero-mode reduction]
\label{thm:v4v75-zero-reduction}
For an arbitrary even interaction
\(\mathcal I(\bar z,z,\bar\eta,\eta)\), define
\begin{equation}
 \mathcal F_{\rm eff}(\bar z,z)
 :=
 \frac{1}{\det D}
 \int
 e^{-\bar\eta D\eta}
 e^{-\mathcal I(\bar z,z,\bar\eta,\eta)}
 \,d\bar\eta\,d\eta.
\label{eq:v4v75-effective-zero-polynomial}
\end{equation}
Then
\begin{equation}
 Z_{\rm soft}
 =
 \det D\,
 [\,\bar z_1z_1\cdots\bar z_kz_k\,]\,
 \mathcal F_{\rm eff}.
\label{eq:v4v75-zero-mode-coefficient}
\end{equation}
In particular:
\begin{enumerate}
\item if \(k>0\) and \(\mathcal I=0\), then \(Z_{\rm soft}=0\);
\item interactions saturate the zero modes exactly when the displayed top
      coefficient is nonzero; and
\item no inverse of the zero block is used.
\end{enumerate}
\end{theorem}

\begin{proof}
Integrate the nonzero variables first.  Their Gaussian denominator is
\(\det D\ne0\), so \eqref{eq:v4v75-effective-zero-polynomial} is a
well-defined finite polynomial in the zero generators.  The remaining
Berezin integral is top-coefficient extraction
\eqref{eq:v4v75-top-integral}, proving
\eqref{eq:v4v75-zero-mode-coefficient}.  The three consequences are
immediate.
\end{proof}

\begin{firewall}[Saturation is configuration and sector dependent]
\label{fw:v4v75-sector}
The kernel of a topological chiral Dirac operator, a Higgs-induced accidental
crossing, and a boundary spectral-section crossing need not have the same
charges or interaction monomials.  A saturation proof for one cannot be
applied to the others from the equality of their dimensions alone.
\end{firewall}

\subsection{Charge selection before estimates}
\label{sec:v4v75-charges}

Let a compact torus of exact global or gauge symmetries act diagonally on
the zero generators, and let \(q(\zeta_a)\) denote their charge vectors.

\begin{proposition}[Charge obstruction to vacuum saturation]
\label{prop:v4v75-charge-obstruction}
Assume every interaction polynomial and the nonzero-mode Gaussian measure
are invariant.  If the ordered zero-mode top monomial has nonzero total
charge
\begin{equation}
 Q_0:=\sum_{a=1}^{2k}q(\zeta_a)\ne0,
\label{eq:v4v75-zero-charge}
\end{equation}
then
\[
 Z_{\rm soft}=0.
\]
More generally, an observable \(\mathcal O\) can have a nonzero soft
coefficient only if its charge cancels \(Q_0\).
\end{proposition}

\begin{proof}
The effective polynomial
\(\mathcal F_{\rm eff}\) is invariant because Berezin contraction of the
nonzero variables preserves the symmetry.  Its coefficient of a monomial
with nonzero charge must vanish.  The zero top monomial has charge \(Q_0\).
Inserting \(\mathcal O\) shifts the required charge by
\(q(\mathcal O)\), giving the second assertion.
\end{proof}

For an anomalous classical current, the fermion measure is not invariant
and the proposition must use the regulator anomaly rather than the
classical charge.  The total Standard Model gauge symmetry is anomaly free,
but baryon and lepton number have the familiar electroweak anomaly.  V76
must therefore use the actual regulated index and not impose all classical
fermion numbers as exact selection rules.

\subsection{First-connectivity ownership of saturating covers}
\label{sec:v4v75-first-connectivity}

Use the V74 local frames to assign the zero generators to finitely many
local packets \(1,\ldots,n\).  A monomial in
\(\mathcal F_{\rm eff}\) defines a hyperedge on the packets containing its
generators.  Give all candidate monomial occurrences one deterministic
order, for example by scale, cell, interaction type, and internal index.

\begin{lemma}[Unique first full connectivity]
\label{lem:v4v75-first-connectivity}
For every connected saturating cover on a fixed packet set, there is a
unique first monomial occurrence at which the cumulative hypergraph becomes
connected.  The partition of the packets immediately before that occurrence
is a unique proper partition.
\end{lemma}

\begin{proof}
The connected-component partition of a cumulative hypergraph can only
coarsen as occurrences are added.  It begins at the discrete partition and
ends at the one-block partition.  In a finite ordered list there is a unique
first index at which the latter appears, and the preceding partition is
therefore unique and proper.
\end{proof}

\begin{theorem}[No-duplication saturation decomposition]
\label{thm:v4v75-no-duplication}
Decompose every cover in
\eqref{eq:v4v75-cover-sum} into its connected packet components and apply
Lemma~\ref{lem:v4v75-first-connectivity} to each nontrivial component.
The resulting data partition the complete cover sum exactly.  Hence a
marked-tree or polymer bound may be applied separately to each first-join
family without summing one cover under several possible joining
coordinates.
\end{theorem}

\begin{proof}
Connected components of a finite hypergraph are unique.  Within each
component, the deterministic order and
Lemma~\ref{lem:v4v75-first-connectivity} give one first-join occurrence and
one preceding partition.  These operations are reversible from the
labelled cover, so the families are disjoint and exhaustive.
\end{proof}

This is the structural input imported from the successful YM V96 repair.
The Yang--Mills coefficient theorem itself is not being asserted for
Grassmann amplitudes; V75 imports only its exact ownership mechanism.

\subsection{A conditional local tree criterion}
\label{sec:v4v75-tree}

Let \(\theta_{ij}=\theta_{ji}\ge0\) be packet pair weights and put
\begin{equation}
 \rho_\theta:=\sup_i\sum_{j\ne i}\theta_{ij}.
\label{eq:v4v75-tree-row}
\end{equation}

\begin{lemma}[Rooted tree summation]
\label{lem:v4v75-tree-sum}
For a fixed root \(i\) and \(n\ge2\),
\begin{equation}
 \sum_{\substack{X\ni i\\|X|=n}}
 \ \sum_{\tau\in\mathfrak T(X)}
 \prod_{\{a,b\}\in E(\tau)}\theta_{ab}
 \le
 (e^2\rho_\theta)^{n-1}.
\label{eq:v4v75-tree-sum}
\end{equation}
\end{lemma}

\begin{proof}
Temporarily order the \(n-1\) nonroot vertices.  For each labelled tree,
root it at \(i\) and sum vertices successively away from the root; every
edge sum is at most \(\rho_\theta\).  Cayley's formula gives
\(n^{n-2}\) labelled trees.  Removing the temporary order divides by
\((n-1)!\).  The elementary bound
\(n^{n-2}/(n-1)!\le e^{2(n-1)}\) proves
\eqref{eq:v4v75-tree-sum}.  Dropping distinctness during the successive
sums only enlarges the left side.
\end{proof}

\begin{theorem}[Absolute convergence of connected saturation corrections]
\label{thm:v4v75-connected-sum}
Suppose every connected first-join family on packet set \(X\) has total
absolute amplitude at most
\begin{equation}
 K_0 K_1^{|X|-1}
 \sum_{\tau\in\mathfrak T(X)}
 \prod_{\{a,b\}\in E(\tau)}\theta_{ab}.
\label{eq:v4v75-family-majorant}
\end{equation}
If
\begin{equation}
 e^2K_1\rho_\theta<1,
\label{eq:v4v75-tree-small}
\end{equation}
then the sum of all connected corrections containing a fixed packet is
absolutely convergent, uniformly in volume.
\end{theorem}

\begin{proof}
Apply Lemma~\ref{lem:v4v75-tree-sum} at every size \(n\).  The resulting
majorant is
\[
 K_0\sum_{n\ge1}(e^2K_1\rho_\theta)^{n-1},
\]
which converges under \eqref{eq:v4v75-tree-small}.
\end{proof}

\begin{firewall}[Convergence is not nonvanishing]
\label{fw:v4v75-convergence-not-nonzero}
Theorem~\ref{thm:v4v75-connected-sum} controls an absolute series.  Complex
saturating covers can cancel.  A nonzero normalization needs a lower bound,
not only convergence.
\end{firewall}

\begin{proposition}[Dominant-cover nonvanishing certificate]
\label{prop:v4v75-dominant-cover}
Let \(z_*\) be one explicitly computed saturating contribution to
\eqref{eq:v4v75-cover-sum}, and let \(R_*\) be a rigorous upper bound on the
sum of the absolute values of every other cover.  If
\begin{equation}
 R_*<|z_*|,
\label{eq:v4v75-dominance}
\end{equation}
then
\begin{equation}
 |Z_{\rm soft}|\ge |z_*|-R_*>0.
\label{eq:v4v75-lower-bound}
\end{equation}
\end{proposition}

\begin{proof}
Write \(Z_{\rm soft}=z_*+(Z_{\rm soft}-z_*)\) and apply the reverse triangle
inequality.
\end{proof}

This is sufficient but deliberately strong.  In a topological sector the
natural distinguished contribution may be the regulated 't Hooft
zero-mode vertex, while in a Higgs crossing it may instead be a Yukawa mass
minor.  Those coefficients and their phases must be computed before
\eqref{eq:v4v75-dominance} can be tested.

\subsection{Weighted payment for collar and seed failures}
\label{sec:v4v75-weighted-failure}

Let disjoint badness annuli \(B_m\) be labelled by increasing inverse-load
scales \(K_m\ge1\), where \(K_m\) bounds any chosen combination of
\(\gamma^{-1}\), \(c_{\rm seed}^{-1}\), and the corresponding marked frame
norm.  Define the flat mark
\[
 \mathfrak m_m:=K_m^{-1}\boldsymbol1_{B_m}.
\]

\begin{lemma}[Exact mark-removal cost]
\label{lem:v4v75-mark-removal}
Pointwise,
\begin{equation}
 \sum_m\boldsymbol1_{B_m}
 =
 \sum_mK_m\mathfrak m_m.
\label{eq:v4v75-mark-removal}
\end{equation}
A bound on \(\sum_m\mathbb E\mathfrak m_m\) alone gives no uniform bound on
the unmarked failure count when \(K_m\to\infty\).  A sufficient input is the
weighted first moment
\begin{equation}
 \sum_mK_m\mathbb E\mathfrak m_m<\infty.
\label{eq:v4v75-weighted-first-moment}
\end{equation}
\end{lemma}

\begin{proof}
The identity is the definition of \(\mathfrak m_m\).  For the negative
statement, concentrate all mass at one annulus with arbitrarily large
\(K_m\); the marked expectation is smaller by \(K_m^{-1}\).  Taking
expectations in \eqref{eq:v4v75-mark-removal} proves sufficiency of
\eqref{eq:v4v75-weighted-first-moment}.
\end{proof}

This is the precise V73--V74 adaptation of the NS V31 lesson.  The required
weight is the actual analytic inverse load, not an arbitrarily chosen
dyadic radius.

\begin{theorem}[V75 soft-sector interface]
\label{thm:v4v75-interface}
A transferred soft sector can enter the SM--Einstein marginal continuum
ledger without a false inverse provided the following data are kept
separate:
\begin{enumerate}
\item the exact hard determinant and anomaly sewing of V73--V74;
\item the top-coefficient cover identity
      \eqref{eq:v4v75-cover-sum};
\item a sector-specific charge and zero-mode classification;
\item a first-connectivity tree bound such as
      \eqref{eq:v4v75-family-majorant};
\item a nonvanishing lower-bound certificate such as
      \eqref{eq:v4v75-dominance}, when normalized vacuum expectations are
      required; and
\item the weighted collar/seed failure moment
      \eqref{eq:v4v75-weighted-first-moment}.
\end{enumerate}
None of these items implies another.
\end{theorem}

\begin{proof}
Items 1--2 are exact factorization and Berezin algebra.  Proposition
\ref{prop:v4v75-charge-obstruction} shows why item 3 precedes estimates.
Theorems \ref{thm:v4v75-no-duplication} and
\ref{thm:v4v75-connected-sum} give item 4.  The firewall following that
theorem and Proposition~\ref{prop:v4v75-dominant-cover} require item 5.
Lemma~\ref{lem:v4v75-mark-removal} requires item 6 independently of the
Grassmann sum.
\end{proof}

\begin{assessment}[Direction after the V75 audit]
\label{ass:v4v75-status}
The companion audit changes the attack in two ways.  Saturating terms will
be organized by one first-full-connectivity owner before absolute values
are summed, following the repaired YM quartic construction.  Singular
projector/frame failures will be charged by their inverse-load first moment,
not by a flat bad-event probability, following NS V31.

The exact finite Grassmann question is now closed: vacuum nonvanishing is
the top zero-mode coefficient of
\(\mathcal F_{\rm eff}\).  The physical Standard Model question is not yet
closed because the zero modes and admissible saturating vertices have not
been classified sector by sector.
\end{assessment}

\subsection{V75 result ledger and V76 target}
\label{sec:v4v75-conclusion}

V75 records stable YM V96 and NS V31 snapshots, proves the exact monomial
cover and interacting zero-mode reduction, gives the charge obstruction,
assigns every saturating cover by first connectivity, and proves a
conditional tree-sum and dominant-cover nonvanishing certificate.  It also
identifies the correct weighted payment for spectral-collar and seed-frame
failures.

V76 should classify the regulated Standard Model zero modes.  It must
separate electroweak topological index sectors, QCD sectors,
Higgs/Yukawa accidental crossings, and boundary spectral-flow modes.  For
each class it should list the exact gauge and global charges of the
Berezin top form and the lowest allowed saturating operator.  In the
electroweak unit-index sector it should derive, rather than merely quote,
the three-generation twelve-doublet zero-mode count and the gauge-invariant
't Hooft selection rule.  Only after this table is exact should a
dominant-cover or paired-sector lower bound be attempted.
\clearpage
\section{V76: indexed chiral zero modes and the Standard Model saturation table}
\label{sec:v4v76}

V75 gives an exact saturation rule for a square soft Schur matrix, where
kernel and cokernel dimensions agree.  A topological Weyl operator is
different: it is a Fredholm map between chiral spaces and can have nonzero
index.  Its right kernel and left cokernel then have unequal dimensions.
Treating it as a square accidental crossing would erase precisely the
topological charge that produces the 't Hooft selection rule.

This note first extends V72--V75 to a rectangular chiral map.  It then uses
the index theorem already imported in
\eqref{eq:v4v1-curved-index} to derive the electroweak and color unit-charge
zero-mode sectors.  Topological modes, index-zero Higgs crossings, and
boundary spectral-flow lines are entered in separate rows because their
saturating operators are not interchangeable.

\subsection{Rectangular singular transfer}
\label{sec:v4v76-rectangular}

Let \(A:E_{\rm R}\to E_{\rm L}\) be a finite-dimensional complex map, with
no equality assumption on the two dimensions.  Define
\[
 P_0=\boldsymbol1_{\{0\}}(A^*A),\qquad
 Q_0=\boldsymbol1_{\{0\}}(AA^*).
\]
Thus
\[
 \operatorname{Ran}P_0=\ker A,\qquad
 \operatorname{Ran}Q_0=\ker A^*.
\]

\begin{lemma}[Positive singular pairs and index imbalance]
\label{lem:v4v76-rectangular-SVD}
Every positive singular value of \(A\) has the same multiplicity on the
right and left.  The only rank imbalance is at zero:
\begin{equation}
 \operatorname{ind}A
 =
 \dim\ker A-\dim\ker A^*.
\label{eq:v4v76-index}
\end{equation}
For every \(\delta>0\), the restrictions of \(A\) between the right and
left spectral subspaces with singular value at least \(\delta\) are square,
invertible, and have inverse norm at most \(\delta^{-1}\).
\end{lemma}

\begin{proof}
The polar decomposition maps each positive eigenspace of \(A^*A\)
unitarily onto the eigenspace of \(AA^*\) with the same eigenvalue.
Rank--nullity gives
\[
 \dim E_{\rm R}-\dim\ker A
 =
 \operatorname{rank}A
 =
 \dim E_{\rm L}-\dim\ker A^*,
\]
which is \eqref{eq:v4v76-index}.  On the spectral subspace above
\(\delta^2\), the quadratic form \(A^*A\) is at least \(\delta^2\), proving
the hard assertion.
\end{proof}

The positive soft singular cluster of V72 still transfers in matched
left/right pairs.  In addition, every unmatched zero-kernel or zero-cokernel
direction is assigned wholly to the soft determinant line.

\begin{theorem}[Rectangular hard integration]
\label{thm:v4v76-rectangular-Berezin}
Choose singular frames
\[
 E_{\rm R}=K_{\rm R}\oplus E_{\rm R}',\qquad
 E_{\rm L}=K_{\rm L}\oplus E_{\rm L}',
\]
where \(K_{\rm R}=\ker A\), \(K_{\rm L}=\ker A^*\), and
\(D=A|_{E_{\rm R}'}:E_{\rm R}'\to E_{\rm L}'\) is invertible.  Let
\(z_1,\ldots,z_{k_{\rm R}}\) be the unbarred variables on \(K_{\rm R}\)
and \(\bar w_1,\ldots,\bar w_{k_{\rm L}}\) the barred variables on
\(K_{\rm L}\).  For an arbitrary even interaction \(\mathcal I\), define
\begin{equation}
 \mathcal F_0(\bar w,z)
 :=
 \frac1{\det D}
 \int
 e^{-\bar\eta D\eta}
 e^{-\mathcal I(\bar w,z,\bar\eta,\eta)}
 \,d\bar\eta\,d\eta.
\label{eq:v4v76-rectangular-effective}
\end{equation}
Then
\begin{equation}
 Z_A
 =
 \det D\,
 [\,\bar w_1\cdots\bar w_{k_{\rm L}}
    z_1\cdots z_{k_{\rm R}}\,]\,
 \mathcal F_0,
\label{eq:v4v76-rectangular-top}
\end{equation}
up to the one fixed orientation sign.  With no interaction or external
fermion insertion, \(Z_A=0\) whenever
\(k_{\rm L}+k_{\rm R}>0\).
\end{theorem}

\begin{proof}
The singular decomposition makes the quadratic action
\(\bar\eta D\eta\); it contains none of the zero variables.  Integrate the
matched nonzero pair and use the top-coefficient rule
\eqref{eq:v4v75-top-integral} on the remaining, not necessarily equal,
sets of barred and unbarred variables.  If
\(\mathcal I=0\), the effective polynomial is one and has no positive-degree
top coefficient.
\end{proof}

\begin{corollary}[V75 is the index-zero special case]
\label{cor:v4v76-V75-special}
If \(A\) is square, then \(k_{\rm R}=k_{\rm L}\) and
Theorem~\ref{thm:v4v76-rectangular-Berezin} reduces to
Theorem~\ref{thm:v4v75-zero-reduction}.  If
\(\operatorname{ind}A\ne0\), the rectangular formula is mandatory.
\end{corollary}

\begin{firewall}[Hermitian doubling erases the index imbalance]
\label{fw:v4v76-doubling-index}
The self-adjoint matrix
\[
 \begin{pmatrix}0&A^*\\A&0\end{pmatrix}
\]
has a symmetric nonzero spectrum and total nullity
\(k_{\rm R}+k_{\rm L}\).  It can locate the zero cluster but does not retain
the signed difference
\eqref{eq:v4v76-index} unless the two chiral summands remain labelled.
Consequently it cannot replace the rectangular determinant line or its
saturation rule.
\end{firewall}

\subsection{Index input and the scope of exact counting}
\label{sec:v4v76-index-input}

For a left Weyl operator
\[
 D_R^+:
 \Gamma(S^+\otimes E_R)\longrightarrow
 \Gamma(S^-\otimes E_R)
\]
on a closed spin four-manifold, the imported curved index formula is
\begin{equation}
 \operatorname{ind}D_R^+
 =
 \int_X[\widehat A(TX)\operatorname{ch}(E_R)]_4.
\label{eq:v4v76-index-theorem}
\end{equation}
For a pure \(SU(N)\) gauge background in the normalization
\begin{equation}
 Q=\frac1{8\pi^2}\int_X\operatorname{tr}_{\rm fund}(F\wedge F)\in\mathbb Z,
\label{eq:v4v76-top-charge}
\end{equation}
the gauge part is
\begin{equation}
 \operatorname{ind}D_R^+=2T(R)Q,
 \qquad
 \operatorname{tr}_R(T^aT^b)=T(R)\delta^{ab}.
\label{eq:v4v76-index-T}
\end{equation}

\begin{lemma}[Fundamental unit-charge count]
\label{lem:v4v76-fundamental-index}
For the fundamental or antifundamental representation,
\(T=1/2\), so a unit-charge background has index of absolute value one per
left Weyl multiplet.  On a standard self-dual or anti-self-dual instanton
for which the opposite-chirality Weitzenbock operator has no kernel, this
index equals the actual zero-mode count on the selected chirality.
\end{lemma}

\begin{proof}
Substitution of \(T=1/2\) into
\eqref{eq:v4v76-index-T} gives the index.  The chiral Weitzenbock formula
writes the opposite product of Dirac operators as
\(\nabla^*\nabla\) plus the curvature term of the chirality not coupled to
the chosen duality component.  On the standard instanton with the usual
decay condition, a vector in its kernel is covariantly constant; decay
forces it to vanish.  Hence the opposite kernel is zero, and the index
difference is the kernel dimension on the selected chirality.
\end{proof}

\begin{firewall}[Index is a difference on a general Einstein background]
\label{fw:v4v76-index-not-nullity}
Without the instanton vanishing hypothesis, the index fixes
\(\dim\ker D^+-\dim\ker(D^+)^*\), not the two dimensions separately.
Metric curvature, reducibility, or an accidental paired crossing may add
equal left/right zero modes.  Those paired modes must be recorded in
addition to the topological index.
\end{firewall}

\subsection{Electroweak unit-charge sector}
\label{sec:v4v76-electroweak}

Use the left-handed Standard Model convention.  In generation \(g\),
\[
 Q_g:(\mathbf3,\mathbf2)_{1/6},
 \qquad
 L_g:(\mathbf1,\mathbf2)_{-1/2}
\]
are the only weak doublets.  The color multiplicity of \(Q_g\) makes it
three independent \(SU(2)\) fundamentals.

\begin{theorem}[Derived twelve-mode electroweak count]
\label{thm:v4v76-EW-count}
On a unit \(SU(2)_L\) instanton satisfying the vanishing hypothesis of
Lemma~\ref{lem:v4v76-fundamental-index}, one Standard Model generation has
\begin{equation}
 3\,[Q_g]+1\,[L_g]=4
\label{eq:v4v76-four-EW}
\end{equation}
left-handed zero-mode coefficients.  Three generations have
\begin{equation}
 N_{\rm EW}=3\cdot4=12.
\label{eq:v4v76-twelve-EW}
\end{equation}
The opposite topological charge gives the conjugate chirality and reverses
the global-charge selection rule.
\end{theorem}

\begin{proof}
Each color value of \(Q_g\) is one fundamental weak doublet and contributes
one zero mode by Lemma~\ref{lem:v4v76-fundamental-index}.  The lepton
doublet contributes one.  All right-handed-conjugate fields are weak
singlets and contribute zero.  Add over the three generations.
\end{proof}

This is the same twelve-doublet multiplicity that makes the ordinary
\(SU(2)\) mod-two anomaly even in the frozen anomaly calculation, but here
it is derived as an integer instanton zero-mode count.

Write the fields with color \(a,b,c\), weak indices \(i,j,k,\ell\), and
two-component spinor indices
\(\alpha,\beta,\gamma,\delta\).  One gauge- and Lorentz-invariant
four-field contraction per generation is
\begin{equation}
 \mathcal O_g
 :=
 \epsilon_{abc}\epsilon_{ij}\epsilon_{k\ell}
 \epsilon_{\alpha\beta}\epsilon_{\gamma\delta}
 Q_g^{ai\alpha}Q_g^{bj\beta}Q_g^{ck\gamma}L_g^{\ell\delta}.
\label{eq:v4v76-one-generation-vertex}
\end{equation}
The actual instanton orientation integral fixes a linear combination of
the equivalent weak/spinor pairings; the displayed term is a nonzero
representative with the required quantum numbers.

\begin{theorem}[Electroweak 't Hooft selection rule]
\label{thm:v4v76-EW-selection}
The twelve-field operator
\begin{equation}
 \mathcal O_{\rm EW}:=\prod_{g=1}^3\mathcal O_g
\label{eq:v4v76-EW-vertex}
\end{equation}
has exactly one field for every unit-instanton zero-mode coefficient and is
a singlet of
\(SU(3)_c\times SU(2)_L\times U(1)_Y\).  Its global charges are
\begin{equation}
 \Delta B=3,\qquad
 \Delta L=3,\qquad
 \Delta(B-L)=0,\qquad
 \Delta(B+L)=6.
\label{eq:v4v76-EW-charges}
\end{equation}
In the massless fixed-\(Q=1\) sector the fermion vacuum coefficient is zero,
whereas a correlation containing the conjugately oriented saturating
twelve-field insertion can have a nonzero zero-mode coefficient.
\end{theorem}

\begin{proof}
There are three \(Q_g\)'s and one \(L_g\) in
\(\mathcal O_g\), matching
\eqref{eq:v4v76-four-EW}; their product matches all twelve modes.
The color epsilon is an \(SU(3)\) singlet.  Two weak epsilons contract four
doublets to an \(SU(2)\) singlet, and the spinor epsilons make a Lorentz
scalar.  The hypercharge per generation is
\[
 3\left(\frac16\right)-\frac12=0.
\]
The baryon and lepton charges per generation are respectively
\(3(1/3)=1\) and \(1\); summing generations gives
\eqref{eq:v4v76-EW-charges}.  The rectangular top-coefficient formula
\eqref{eq:v4v76-rectangular-top} proves the vacuum statement and the
saturation possibility.  Whether the relevant observable uses
\(\mathcal O_{\rm EW}\) or its conjugate depends on the sign convention for
\(Q\) and the Berezin orientation.
\end{proof}

\begin{firewall}[Selection is not an amplitude lower bound]
\label{fw:v4v76-EW-no-lower}
Theorem~\ref{thm:v4v76-EW-selection} proves generator saturation and exact
charges.  It does not bound the instanton measure, the gauge-orientation
integral, the zero-mode wavefunction determinant, or destructive
interference with other sectors.  It therefore does not yet verify the
dominant-cover inequality \eqref{eq:v4v75-dominance}.
\end{firewall}

\subsection{Color unit-charge sector}
\label{sec:v4v76-QCD}

There are six physical Dirac quark flavours after electroweak component
splitting.  A unit \(SU(3)_c\) instanton contributes one chiral zero mode
per fundamental Dirac flavour.  In the Euclidean Dirac Berezin integral,
the independent conjugate coefficient supplies the matching barred
variable.

\begin{theorem}[QCD zero-mode determinant]
\label{thm:v4v76-QCD-vertex}
For \(N_f=6\) massless quark flavours, the unit color-instanton zero-mode
sector contains six unbarred and six barred Grassmann coefficients.  Its
lowest flavor-covariant saturating polynomial is
\begin{equation}
 \mathcal O_{\rm QCD}
 =
 \det_{1\le f,g\le6}
 \bigl(\bar q_R^{\,f}q_L^{\,g}\bigr)
\label{eq:v4v76-QCD-vertex}
\end{equation}
for one instanton orientation, with \(L\) and \(R\) exchanged for the
opposite orientation.  It preserves vector baryon number and changes the
axial quark number by \(2N_f=12\).
\end{theorem}

\begin{proof}
Lemma~\ref{lem:v4v76-fundamental-index} gives one selected-chirality mode
for each flavor.  The adjoint Grassmann expansion has one independent
coefficient on the same zero wavefunction.  Expanding the determinant in
\eqref{eq:v4v76-QCD-vertex} chooses every flavor row and column exactly
once, so Lemma~\ref{lem:v4v75-cover} shows that it saturates all twelve
coefficients.  Each bilinear has zero vector baryon charge and axial charge
two.
\end{proof}

\begin{corollary}[Mass-matrix saturation in the QCD sector]
\label{cor:v4v76-QCD-mass}
If the quadratic action contains a nonsingular flavor mass matrix \(M\),
the zero-mode integral contributes, up to the fixed orientation convention,
\begin{equation}
 \det M.
\label{eq:v4v76-QCD-mass}
\end{equation}
For Standard Model quarks before fixing a Higgs background, this is a
Yukawa--Higgs polynomial and must remain inside the common Higgs integral.
\end{corollary}

\begin{proof}
Restrict the mass bilinear
\(\bar q_R M q_L\) to the zero modes.  Its top coefficient is the determinant
by \eqref{eq:v4v75-quadratic-det}.
\end{proof}

\subsection{Index-zero Higgs crossings and boundary flow}
\label{sec:v4v76-other-sectors}

\begin{proposition}[Accidental Higgs crossing]
\label{prop:v4v76-Higgs-crossing}
A finite square Higgs--Yukawa matrix crossing with unchanged chiral index
adds equal dimensions to \(\ker A\) and \(\ker A^*\).  Its saturation is
therefore governed by the square formula
\eqref{eq:v4v75-zero-mode-coefficient}.  The lowest possible quadratic
saturator is the corresponding restricted Yukawa/mass minor; topological
't Hooft counting alone gives no information about its value.
\end{proposition}

\begin{proof}
The Fredholm index is locally constant under continuous finite-dimensional
deformations.  A square matrix has index zero, so kernel and cokernel
dimensions agree.  Apply Corollary~\ref{cor:v4v76-V75-special}; the
topological charge does not enter that algebra.
\end{proof}

\begin{proposition}[Boundary spectral flow is transition data]
\label{prop:v4v76-boundary-flow}
A crossing of a tangential boundary operator changes the chosen spectral
section by its finite crossing line.  It requires a bulk zero-mode
saturator only if the full bulk boundary-value problem acquires kernel.
Otherwise it is completely accounted for by the rank-change determinant
line \eqref{eq:v4v74-crossing-line} and its exact cocycle.
\end{proposition}

\begin{proof}
A tangential eigenvalue crossing changes the polarization, hence the finite
relative exterior-power line, whether or not the bulk operator has kernel.
The bulk Berezin integral has an unsaturated variable exactly when its
Schur boundary-value operator has kernel.  Thus tangential spectral flow
alone is transition data; the additional bulk-kernel condition is necessary
for the saturation rule.
\end{proof}

\subsection{Sector table}
\label{sec:v4v76-table}

The exact classification reached so far is
\begin{center}
\begin{tabular}{|p{0.18\textwidth}|p{0.18\textwidth}|p{0.25\textwidth}|p{0.27\textwidth}|}
\hline
Sector & Soft imbalance & Lowest identified saturator & Present status\\
\hline
\(SU(2)_L\), \(Q=1\) &
12 selected-chirality Weyl coefficients &
\(\mathcal O_{\rm EW}\), twelve fields, or an explicitly computed
Yukawa--Higgs lift &
Count and charges proved; amplitude lower bound open\\
\hline
\(SU(3)_c\), \(Q=1\) &
six quark zero pairs in Dirac variables &
\(\det(\bar q_Rq_L)\); at fixed mass, \(\det M\) &
Selection proved; common Higgs integration and uniform bound open\\
\hline
Higgs/Yukawa accidental crossing &
equal kernel/cokernel rank &
restricted soft Yukawa/mass minor or higher interaction cover &
V75 square rule applies; nonzero minor not proved\\
\hline
Boundary spectral flow &
finite relative crossing line; no automatic bulk kernel &
determinant-line transition; bulk saturator only when bulk kernel appears &
exact cocycle proved; bulk-kernel test remains\\
\hline
Curved gravitational/gauge index &
difference fixed by
\(\int[\widehat A\,\operatorname{ch}]_4\) &
sector-specific external or interaction top form &
index known; separate nullities and saturation open\\
\hline
\end{tabular}
\end{center}

\begin{firewall}[Do not multiply independent sector counts]
\label{fw:v4v76-no-count-product}
If electroweak, color, Higgs, and boundary crossings coexist, their zero
spaces can mix through Yukawa and boundary conditions.  The full kernel must
be computed before applying the cover rule.  Multiplying the separate table
entries assumes a block diagonalization that has not been proved.
\end{firewall}

\begin{theorem}[V76 typed saturation interface]
\label{thm:v4v76-interface}
For a regulated background sector, a valid soft saturation card consists
of:
\begin{enumerate}
\item the rectangular pair
      \((\ker A,\ker A^*)\) and its signed index;
\item the one determinant-line orientation inherited from V74;
\item the exact charges of the ordered top monomial;
\item a declared interaction or observable whose restricted monomial covers
      every zero generator;
\item the first-connectivity/tree majorant of V75 for all competing covers;
      and
\item a lower bound if division by the vacuum coefficient is required.
\end{enumerate}
The electroweak and QCD rows above supply items 1, 3, and the lowest
candidate for item 4 in their unit-charge instanton sectors.  They do not
supply items 5--6.
\end{theorem}

\begin{proof}
Theorem~\ref{thm:v4v76-rectangular-Berezin} proves necessity and sufficiency
of generator coverage after hard integration.  V74 supplies the line
orientation.  Proposition~\ref{prop:v4v75-charge-obstruction} gives the
charge test.  Theorems~\ref{thm:v4v75-no-duplication} and
\ref{thm:v4v75-connected-sum} control competing covers, while
Proposition~\ref{prop:v4v75-dominant-cover} explains the independent lower
bound.  The sector theorems compute the stated entries.
\end{proof}

\begin{assessment}[What V76 changes]
\label{ass:v4v76-status}
The topological sector cannot be treated as a square determinant zero.
Its unmatched chiral variables are the index.  The unit electroweak
instanton has twelve such Standard Model coefficients and the gauge-invariant
twelve-field top form carries
\(\Delta B=\Delta L=3\).  The color instanton instead has a vectorlike
flavor determinant and can be saturated by a mass determinant.  Accidental
Higgs crossings remain index-zero paired modes, and boundary spectral flow
is not automatically a bulk zero.

This classification prevents the next estimates from combining
incompatible zero-mode mechanisms.  It still gives no volume- or
cutoff-uniform lower bound for the complete complex normalization.
\end{assessment}

\subsection{V76 result ledger and V77 target}
\label{sec:v4v76-conclusion}

V76 generalizes the soft transfer and top-coefficient formula to rectangular
chiral maps.  It derives the one-mode-per-fundamental index, the twelve
electroweak zero modes, their gauge-invariant twelve-field operator and
\(B\), \(L\) selection rule, and the six-flavor QCD determinant vertex.  It
also separates index-zero Higgs crossings and boundary spectral flow.

V77 should now return to the full Higgs--Yukawa operator.  On a fixed
finite regulator and topological sector, it should compute the zero-mode
matrix of the actual Standard Model Yukawa interaction.  The first target
is the exact determinant/minor factor multiplying the electroweak and QCD
top forms, with all generation mixing, Higgs components, and CKM/PMNS
basis changes retained.  If that matrix is singular because a neutrino is
massless or a Higgs component vanishes, the corresponding external
observable sector must be kept rather than assuming a vacuum lower bound.
\clearpage
\section{V77: Yukawa index stability, electroweak external legs, and the QCD mass minor}
\label{sec:v4v77}

V76 identifies the candidate saturators sector by sector.  Its proposed
next step was to compress the actual Yukawa interaction to the zero modes.
Before computing a minor, one must ask whether a minor is allowed to remove
the index imbalance.  For the electroweak chiral complex the answer is no:
a smooth Higgs--Yukawa coupling is a zeroth-order perturbation of an elliptic
operator and cannot change its Fredholm index.  It changes the profiles and
mixes generations, but a unit weak topological sector retains an imbalance
of twelve.

Color is different.  The physical color representation is vectorlike, so
its left-minus-right index is zero.  Its unit-instanton zero modes occur in
matched Dirac pairs, and a quark mass matrix can saturate them.  This note
proves both statements and keeps the complex determinant phase that later
combines with the QCD theta angle.

\subsection{Physical-chirality index convention}
\label{sec:v4v77-index-convention}

For a gauge factor \(G\), let \(\chi_f=+1\) for a physical left-handed Weyl
field and \(\chi_f=-1\) for a physical right-handed Weyl field.  In a pure
\(G\) background of charge \(Q\), the gauge contribution to the total
physical-chirality index is
\begin{equation}
 \operatorname{Ind}_G(\mathscr D)
 =
 Q\sum_f\chi_f\,2T_G(R_f)\,m_f,
\label{eq:v4v77-physical-index}
\end{equation}
where \(m_f\) is the multiplicity from spectator gauge representations and
generation.

In an all-left-handed anomaly table, a physical right-handed field is
represented by its left-handed charge conjugate.  Its determinant line is
dualized accordingly.  Formula \eqref{eq:v4v77-physical-index} retains the
operator chirality sign explicitly; omitting it would incorrectly assign a
nonzero index to a vectorlike Dirac pair.

\begin{lemma}[Electroweak and color index sums]
\label{lem:v4v77-index-sums}
For three Standard Model generations,
\begin{align}
 \operatorname{Ind}_{SU(2)_L}(\mathscr D)
 &=12Q_2,
\label{eq:v4v77-EW-index}\\
 \operatorname{Ind}_{SU(3)_c}(\mathscr D)
 &=0.
\label{eq:v4v77-color-index}
\end{align}
\end{lemma}

\begin{proof}
Only physical left-handed weak doublets contribute to the first sum.
There are three color copies of \(Q_L\) and one \(L_L\) per generation,
each with \(2T(\mathbf2)=1\).  Thus
\[
 3\text{ generations}\times(3+1)=12.
\]
For color, one generation has two left fundamental components in \(Q_L\)
and two right fundamental fields \(u_R,d_R\).  Since
\(2T(\mathbf3)=1\),
\[
 2-(1+1)=0
\]
per generation.  Spectator weak multiplicity is the factor two in the
first term.
\end{proof}

\subsection{Zeroth-order Yukawa perturbations preserve the index}
\label{sec:v4v77-index-stability}

Let \(X\) be compact, with either no boundary or one fixed elliptic boundary
domain.  Assemble the physical Standard Model Weyl operators into
\[
 \mathscr D_0:\mathcal H^1_{\rm R}\longrightarrow\mathcal H^0_{\rm L}.
\]
A smooth Higgs field and the finite generation matrices define a
gauge-covariant bundle multiplication operator
\[
 \mathscr Y_H:\mathcal H^1_{\rm R}\longrightarrow\mathcal H^0_{\rm L}.
\]
Put
\begin{equation}
 \mathscr D_H=\mathscr D_0+\mathscr Y_H.
\label{eq:v4v77-coupled-operator}
\end{equation}

\begin{theorem}[Yukawa index stability]
\label{thm:v4v77-Yukawa-index}
Along every smooth path of Higgs fields and finite Yukawa matrices that
preserves the elliptic domain,
\begin{equation}
 \operatorname{ind}\mathscr D_H
 =
 \operatorname{ind}\mathscr D_0.
\label{eq:v4v77-Yukawa-index}
\end{equation}
Consequently, in a weak topological sector of charge \(Q_2\),
\begin{equation}
 \dim\ker\mathscr D_H-\dim\ker\mathscr D_H^*
 =
 12Q_2.
\label{eq:v4v77-coupled-EW-index}
\end{equation}
No choice of smooth Standard Model Yukawa matrices or Higgs background can
make \(\mathscr D_H\) invertible when \(Q_2\ne0\).
\end{theorem}

\begin{proof}
The principal symbol of \(\mathscr D_H\) is the principal symbol of
\(\mathscr D_0\), so ellipticity and the fixed boundary complementing
condition are unchanged.  More explicitly, multiplication by a smooth
finite matrix is bounded on \(L^2\), while the Rellich embedding of the
elliptic \(H^1\) domain into \(L^2\) is compact.  Hence
\(\mathscr Y_H:\mathcal H^1_{\rm R}\to\mathcal H^0_{\rm L}\) is a compact
perturbation of the Fredholm graph operator.  Fredholm index is invariant
under the norm-continuous path
\(\mathscr D_0+t\mathscr Y_H\), proving
\eqref{eq:v4v77-Yukawa-index}.  Combine this with
\eqref{eq:v4v77-EW-index}.  An invertible map has index zero, proving the
last assertion.
\end{proof}

\begin{remark}[Relation to the frozen mass-index obstruction]
\label{rem:v4v77-frozen-mass-index}
The frozen third-series V69 mass-matrix theorem proves the finite
rank--nullity version for opposite-chirality mirror multiplicities.  Theorem
\ref{thm:v4v77-Yukawa-index} is its elliptic background-family version:
the zeroth-order mass block can remove paired modes but not the principal
symbol index.
\end{remark}

\begin{firewall}[Noncompact asymptotics require a Callias problem]
\label{fw:v4v77-noncompact}
The compactness proof uses a compact regulated spacetime or a fixed compact
elliptic boundary problem.  On \(\mathbb R^4\), a Higgs field approaching a
nonzero value at infinity is not a compactly supported perturbation.  The
appropriate statement is then a Callias-type index theorem with declared
asymptotic data.  V77 does not infer that theorem merely from Rellich
compactness.
\end{firewall}

\subsection{Consequences for the electroweak Berezin integral}
\label{sec:v4v77-EW-Berezin}

Let
\[
 k_{\rm R}=\dim\ker\mathscr D_H,\qquad
 k_{\rm L}=\dim\ker\mathscr D_H^*.
\]
For \(Q_2=1\),
\begin{equation}
 k_{\rm R}=12+k_{\rm L}.
\label{eq:v4v77-EW-nullities}
\end{equation}
Paired accidental modes may increase both sides, but cannot alter their
difference.

\begin{theorem}[Fixed electroweak topology has no quadratic vacuum
saturation]
\label{thm:v4v77-EW-vacuum-zero}
At fixed \(Q_2\ne0\), the complete Standard Model fermion action, including
arbitrary smooth gauge-covariant Higgs--Yukawa bilinears, has nonzero
Fredholm index.  Its uninserted fermion Berezin coefficient therefore
vanishes:
\begin{equation}
 Z_{\rm F}(Q_2;A,H)=0.
\label{eq:v4v77-EW-vacuum-zero}
\end{equation}
For \(Q_2=1\), any nonzero correlation must supply at least twelve more
unbarred than barred zero-mode factors, in the orientation of
\eqref{eq:v4v77-coupled-EW-index}.
\end{theorem}

\begin{proof}
The action is quadratic in the fermions for every fixed bosonic background;
Yukawa coupling changes the matrix but not its quadratic character.
Theorem~\ref{thm:v4v77-Yukawa-index} gives unequal kernel/cokernel
dimensions.  In the rectangular reduction
\eqref{eq:v4v76-rectangular-top}, every quadratic monomial contains one
barred and one unbarred coordinate.  It cannot cover a top form with
unequal numbers of the two kinds.  Hence the vacuum coefficient is zero.
An inserted polynomial must compensate their difference, giving the
twelve-factor lower bound.
\end{proof}

\begin{corollary}[Yukawa mixing deforms the 't Hooft tensor, not its degree]
\label{cor:v4v77-EW-degree}
Generation mixing and a smooth Higgs background can rotate and deform the
twelve electroweak zero-mode wavefunctions.  They cannot change the
twelve-field degree or the charges
\eqref{eq:v4v76-EW-charges}.  In a local frame, their external top form is
the top exterior power of the unmatched kernel bundle.
\end{corollary}

\begin{proof}
The integer difference is constant by
Theorem~\ref{thm:v4v77-Yukawa-index}.  Gauge and exact global charges are
constant on a continuous representation-preserving deformation.  A basis
change on the kernel acts on its top exterior power by its determinant,
which is precisely the determinant-line transition retained in V74.
\end{proof}

The twelve-field expression
\eqref{eq:v4v76-EW-vertex} is therefore a representative at a convenient
weak basis and simple instanton background.  The physical tensor at a
general Higgs/gauge background is its Kato-transported exterior-power
section, including the one determinant-line Jacobian.

\begin{firewall}[Do not close electroweak external legs with a mass
determinant]
\label{fw:v4v77-no-EW-mass-close}
A formal product of charged-fermion masses cannot replace the twelve
unmatched electroweak factors in
Theorem~\ref{thm:v4v77-EW-vacuum-zero}.  Such a replacement would turn an
index-twelve rectangular problem into an index-zero square determinant.
Paired modes may be lifted by Yukawa matrices; the unmatched exterior power
must remain as external or higher-interaction data.
\end{firewall}

\subsection{Weak-basis covariance}
\label{sec:v4v77-basis}

Let the generation basis transformations be unitary:
\[
 Q_L\mapsto U_QQ_L,\quad
 L_L\mapsto U_LL_L,\quad
 u_R\mapsto U_uu_R,\quad
 d_R\mapsto U_dd_R,\quad
 e_R\mapsto U_ee_R.
\]
The Yukawa matrices transform in the corresponding bifundamental way, for
example
\begin{equation}
 Y_u\mapsto U_QY_uU_u^*,
 \qquad
 Y_d\mapsto U_QY_dU_d^*.
\label{eq:v4v77-Yukawa-basis}
\end{equation}

\begin{proposition}[Exterior-power covariance]
\label{prop:v4v77-basis-covariance}
Under a weak-basis transformation, every compressed zero-mode matrix
transforms by left and right multiplication, its determinant/minors
transform by the induced exterior powers, and the zero-mode top form
transforms by the inverse Berezin Jacobian.  Their complete saturated
coefficient is basis independent.  CKM or PMNS matrices may change the
component expression but cannot change the index, zero-mode degree, or
saturation criterion.
\end{proposition}

\begin{proof}
Compression is functorial:
\[
 M_0=L_0^*\mathscr Y_HR_0
 \longmapsto
 (L_0^*U_{\rm L}^*)\,
 (U_{\rm L}\mathscr Y_HU_{\rm R}^*)\,
 (U_{\rm R}R_0).
\]
Taking the \(k\)-th exterior power gives the determinant or the appropriate
minor.  Grassmann coordinate changes carry the reciprocal determinant.
The two cancel in the top-coefficient extraction.  Unitary flavor
diagonalization is a special case.
\end{proof}

\subsection{The vectorlike color zero-mode matrix}
\label{sec:v4v77-color-matrix}

Now take \(Q_3=1\), \(Q_2=0\), and a background for which the six physical
quark flavours have the standard matched color zero modes.  Normalize the
left and conjugate zero-mode wavefunctions so their scalar overlap is one.
At fixed Higgs background, let \(M_u(H)\) and \(M_d(H)\) be the \(3\times3\)
generation mass matrices obtained by restricting the Yukawa multiplication
operator to the up- and down-type zero modes.

\begin{theorem}[Exact QCD zero-mode minor]
\label{thm:v4v77-QCD-minor}
The color unit-instanton zero-mode contribution to the fermion determinant
is
\begin{equation}
 \mathcal Z_{0,{\rm QCD}}(H)
 =
 \det M_u(H)\,\det M_d(H),
\label{eq:v4v77-QCD-minor}
\end{equation}
up to the common instanton orientation and determinant-line convention.
For a constant symmetry-breaking Higgs field with
\(|H|=v/\sqrt2\),
\begin{equation}
 \mathcal Z_{0,{\rm QCD}}
 =
 \left(\frac v{\sqrt2}\right)^6
 \det Y_u\,\det Y_d.
\label{eq:v4v77-QCD-vev-minor}
\end{equation}
\end{theorem}

\begin{proof}
The physical color index vanishes by
\eqref{eq:v4v77-color-index}, so the six left and six conjugate right
coefficients form a square zero-mode mass block.  Weak charge separates it
into up and down \(3\times3\) generation blocks at a fixed neutral Higgs
orientation.  The top-coefficient determinant identity
\eqref{eq:v4v75-quadratic-det} gives the product of their determinants.
For constant \(H\), each block is
\((v/\sqrt2)Y\); taking the two \(3\times3\) determinants proves
\eqref{eq:v4v77-QCD-vev-minor}.
\end{proof}

For a nonconstant Higgs field, the entries contain the actual zero-mode
overlap:
\begin{equation}
 (M_u(H))_{fg}
 =
 \int_X
 \overline{\varphi_{R,f}(x)}\,
 Y_{u,fg}\,\widetilde H(x)\,
 \varphi_{L,g}(x)\,dx,
\label{eq:v4v77-overlap-matrix}
\end{equation}
and similarly for \(d\).  Equation
\eqref{eq:v4v77-QCD-minor}, with these compressed matrices, remains exact;
the simple factorization
\eqref{eq:v4v77-QCD-vev-minor} need not.

\begin{corollary}[Color vacuum/nonvacuum dichotomy]
\label{cor:v4v77-QCD-dichotomy}
If both compressed matrices in
\eqref{eq:v4v77-QCD-minor} are invertible, the color zero modes are
quadratically saturated.  If either is singular, the fixed color-topology
vacuum coefficient vanishes and the missing flavor factors must be supplied
by the QCD 't Hooft observable or another declared interaction.
\end{corollary}

\begin{proof}
The product in \eqref{eq:v4v77-QCD-minor} is nonzero exactly when both
determinants are nonzero.  Apply the cover rule
Theorem~\ref{thm:v4v75-saturation-expansion} to any additional
interaction.
\end{proof}

\subsection{The retained strong phase}
\label{sec:v4v77-strong-phase}

Under independent chiral flavor rephasings, the phase of
\(\det(Y_uY_d)\) changes.  The anomalous Jacobian shifts the color theta
angle by the opposite amount.

\begin{proposition}[Basis-invariant strong phase]
\label{prop:v4v77-strong-phase}
When all six quark masses are nonzero, the color unit-topology coefficient
depends on the phase combination
\begin{equation}
 \overline\theta
 =
 \theta_{\rm QCD}+\arg\det(M_uM_d)
 \pmod{2\pi},
\label{eq:v4v77-theta-bar}
\end{equation}
with the sign fixed by the common determinant-line convention.  A weak-basis
change cannot remove \(\overline\theta\); it transfers phase between the
mass minor and the anomalous measure Jacobian.
\end{proposition}

\begin{proof}
A chiral change of variables multiplies the mass determinant by the
determinant of the flavor rotations.  The regulated color index gives the
opposite phase in the fermion measure, equivalently a shift of the
coefficient of \(Q_3\).  Their sum in
\eqref{eq:v4v77-theta-bar} is unchanged.  This is the determinant-line
version of the usual anomalous change-of-variables calculation.
\end{proof}

\begin{firewall}[No positivity from nonzero quark masses]
\label{fw:v4v77-no-positive-QCD}
Nonzero
\(\det M_u\det M_d\) saturates the color zero modes but generally carries
the physical phase \eqref{eq:v4v77-theta-bar}.  It does not make the complete
SM--Einstein density positive or prove a volume-uniform lower bound after
the Higgs and gauge fields are integrated.
\end{firewall}

\subsection{Minimal neutrinos and other quadratic extensions}
\label{sec:v4v77-neutrinos}

The minimal Standard Model has no right-handed neutrino Yukawa matrix.  A
massless neutrino branch must remain in the infrared soft field; it cannot
be hidden inside an invertible marginal fermion parent.  Adding
gauge-singlet right-handed neutrinos changes the finite mass matrices and
introduces a PMNS basis, but contributes zero to
\eqref{eq:v4v77-EW-index}.  Thus it cannot change the twelve-unit weak
topological imbalance.  A Majorana mass changes the Berezin problem from a
determinant to a Pfaffian and must be supplied with its own antisymmetric
top-form rule.

\begin{theorem}[V77 zero-mode matrix decision]
\label{thm:v4v77-decision}
At fixed regulator:
\begin{enumerate}
\item in a nonzero electroweak topological sector, retain the unmatched
      exterior-power zero-mode line and external 't Hooft insertion; no
      Yukawa vacuum minor can remove the index;
\item in a pure color unit sector, use the exact compressed mass minor
      \eqref{eq:v4v77-QCD-minor};
\item in an index-zero accidental Higgs crossing, use the V75 square Schur
      minor; and
\item for a Majorana extension, replace determinant saturation by an
      explicitly oriented Pfaffian calculation.
\end{enumerate}
These four algebraic routes must not be merged.
\end{theorem}

\begin{proof}
The four items are respectively
Theorem~\ref{thm:v4v77-EW-vacuum-zero},
Theorem~\ref{thm:v4v77-QCD-minor},
Proposition~\ref{prop:v4v76-Higgs-crossing}, and the antisymmetry of a
Majorana quadratic form.
\end{proof}

\begin{assessment}[What V77 closes]
\label{ass:v4v77-status}
V77 answers the first zero-mode-matrix question more sharply than a direct
minor calculation.  Electroweak topology has an index obstruction, so the
quadratic vacuum coefficient vanishes even after Yukawa mixing; the
physical output is an external twelve-field determinant-line section.
Vectorlike color has no such imbalance, and its exact zero-mode coefficient
is the product of the compressed up- and down-generation mass
determinants, retaining the strong CP phase.

The remaining analytic problem is to control these sector coefficients
after integration over gauge, Higgs, and metric fields.  In particular,
topological-sector probabilities and inverse frame/collar loads require
the weighted ledger of V75.
\end{assessment}

\subsection{V77 result ledger and V78 target}
\label{sec:v4v77-conclusion}

V77 proves physical-chirality index sums, stability of the electroweak index
under smooth Yukawa perturbations, the resulting fixed-topology vacuum
zero, weak-basis covariance, the exact QCD compressed mass minor, and the
retained strong phase.  It also isolates the minimal-neutrino and Majorana
branches.

V78 should attach these algebraic cards to a topological-sector
decomposition of the bosonic measure.  It should define the integer
\(SU(2)\) and \(SU(3)\) charges on the admissible finite regulator, assign
each sector once, and derive a weighted sum for: the external electroweak
top form, the QCD mass minor, the hard determinant modulus/phase, and the
collar/seed inverse loads.  A dilute-instanton formula may be used only as a
test; the proof target is a sector-weighted bound under the actual
interacting gauge--Higgs law.
\clearpage
\section{V78: topological-sector ownership and the rooted charged-cluster criterion}
\label{sec:v4v78}

V77 determines the fermionic coefficient at fixed weak or color topology.
The next problem is to sum those coefficients under the interacting
gauge--Higgs law.  A direct requirement on the exponential moment of the
total topological charge is too strong in increasing volume: a gas of
locally rare instantons already makes that moment extensive.  The correct
continuum interface is a rooted connected expansion in which remote
topological components cancel between numerator and denominator.

This note gives the exact finite-regulator sector partition, derives its
fermionic selection rules, and proves a charge-marked polymer criterion
that is uniform for local observables.  It uses the overlap/geometric index
and concentration-ball results imported in
Section~\ref{sec:v4v1-fermion}; it does not reconstruct a second lattice
topological charge.

\subsection{One finite-regulator sector partition}
\label{sec:v4v78-sector-partition}

Let \(\Omega_a\) be the finite lattice field space.  Let
\(\mathcal A_a\subset\Omega_a\) be the quantitatively admissible set on which
the established overlap/geometric construction gives integer charges
\[
 Q_{2,a},Q_{3,a}:\mathcal A_a\longrightarrow\mathbb Z
\]
for \(SU(2)_L\) and \(SU(3)_c\).  Put
\begin{align}
 \Omega_a(m,n)
 &:=
 \{\,\Phi\in\mathcal A_a:
 Q_{2,a}(\Phi)=m,\ Q_{3,a}(\Phi)=n\,\},
\label{eq:v4v78-sector}\\
 \Omega_a^{\rm def}
 &:=
 \Omega_a\setminus\mathcal A_a.
\label{eq:v4v78-defect-sector}
\end{align}

\begin{lemma}[Disjoint finite-cutoff ownership]
\label{lem:v4v78-sector-partition}
At fixed lattice and volume,
\begin{equation}
 1_{\Omega_a}
 =
 1_{\Omega_a^{\rm def}}
 +
 \sum_{(m,n)\in\mathbb Z^2}1_{\Omega_a(m,n)},
\label{eq:v4v78-partition}
\end{equation}
and only finitely many terms are nonempty.  Every configuration is assigned
once.
\end{lemma}

\begin{proof}
On \(\mathcal A_a\), the ordered pair of integer indices has one value, so
its level sets are disjoint and exhaustive.  The complement is disjoint
from all of them.  A finite regulator has finite fermion matrices, hence
bounded indices, so only finitely many pairs occur.
\end{proof}

For an observable \(\mathcal O\), let \(W_a(\Phi;\mathcal O)\) denote the
complete unnormalized integrand after the hard/soft fermion ancestry has
been assembled but before normalization.  Then
\begin{equation}
 Z_a[\mathcal O]
 =
 Z_a^{\rm def}[\mathcal O]
 +
 \sum_{m,n}Z_{a;m,n}[\mathcal O]
\label{eq:v4v78-Z-sector}
\end{equation}
is an identity, not a semiclassical approximation.

\begin{firewall}[No charge on the defect set by declaration]
\label{fw:v4v78-no-defect-charge}
The rough complement
\(\Omega_a^{\rm def}\) is not assigned an arbitrary integer in order to
complete the table.  It is controlled by the existing good/defect and
concentration-ball ledger.  Only after a proved rough-index extension may
it be subdivided by charge.
\end{firewall}

\subsection{Exact fermionic sector selection}
\label{sec:v4v78-selection}

\begin{theorem}[Vacuum weak-topology selection]
\label{thm:v4v78-vacuum-selection}
For the quadratic Standard Model fermion action of V77 and no external
fermion insertion,
\begin{equation}
 Z_{a;m,n}[1]=0
 \qquad\text{whenever }m\ne0,
\label{eq:v4v78-vacuum-mzero}
\end{equation}
provided the finite regulator realizes the physical-chirality index
\[
 \operatorname{ind}\mathscr D_{H,a}=12m
\]
and the boundary domain used in that index.  No corresponding vanishing is
implied by color charge \(n\), whose physical index is zero.
\end{theorem}

\begin{proof}
For every fixed bosonic configuration in \(\Omega_a(m,n)\), V77's index
stability gives a rectangular fermion map with kernel/cokernel difference
\(12m\).  The rectangular top-coefficient formula
\eqref{eq:v4v76-rectangular-top} makes its uninserted Berezin integral zero.
Integrating the pointwise zero over the sector proves
\eqref{eq:v4v78-vacuum-mzero}.  The color index is zero by
\eqref{eq:v4v77-color-index}, so the same argument has no premise for
\(n\ne0\).
\end{proof}

Let \(B(\mathcal O)\) and \(L(\mathcal O)\) be the baryon and lepton charges
of a fermionic observable.

\begin{corollary}[Electroweak observable selection]
\label{cor:v4v78-observable-selection}
With the field and instanton orientation convention of
\eqref{eq:v4v76-EW-charges}, a necessary condition for a sector coefficient
to be nonzero is
\begin{equation}
 B(\mathcal O)-L(\mathcal O)=0,
 \qquad
 B(\mathcal O)+L(\mathcal O)=6m.
\label{eq:v4v78-BL-selection}
\end{equation}
Changing the orientation convention reverses \(m\) and the displayed
charge together.
\end{corollary}

\begin{proof}
A charge-\(m\) sector has \(12|m|\) net chiral zero-mode factors.  In the
V76 orientation, their saturating field top form carries \(B=L=3m\).
Adding and subtracting these two charge conditions gives
\eqref{eq:v4v78-BL-selection}.
\end{proof}

For color charge \(n\), let \(M_{0,n}(H)\) be the complete compressed
zero-mode mass matrix.  Its dimension is at least \(6|n|\), with equality
when there are no accidental paired modes.

\begin{proposition}[Color-sector zero-mode factor]
\label{prop:v4v78-color-factor}
On a collared chart with exactly the index-forced color zero pairs,
\begin{equation}
 \mathcal Z_{0,3}(n;H)=\det M_{0,n}(H).
\label{eq:v4v78-color-factor}
\end{equation}
If the charge-\(n\) zero modes split into \(|n|\) orthogonal unit-charge
blocks and the Higgs mass is constant, then
\begin{equation}
 \det M_{0,n}
 =
 \left[
 \left(\frac v{\sqrt2}\right)^6
 \det Y_u\det Y_d
 \right]^{|n|}
\label{eq:v4v78-dilute-color-test}
\end{equation}
up to orientation.  Without that block decomposition,
\eqref{eq:v4v78-dilute-color-test} is only a dilute test; the exact object
is \eqref{eq:v4v78-color-factor}.
\end{proposition}

\begin{proof}
The top-coefficient formula for the matched color zero pairs gives the
determinant of their complete compressed mass matrix.  Under the stated
orthogonal direct sum, that matrix is block diagonal with \(|n|\) copies of
the V77 unit-charge block, so its determinant is the displayed power.
\end{proof}

\begin{lemma}[Hadamard charge load]
\label{lem:v4v78-Hadamard}
If \(M_{0,n}\) is \(6|n|\times6|n|\) and every row has Euclidean norm at
most \(C_M\), then
\begin{equation}
 |\det M_{0,n}|
 \le C_M^{6|n|}.
\label{eq:v4v78-Hadamard}
\end{equation}
\end{lemma}

\begin{proof}
This is Hadamard's determinant inequality, applied to the row vectors.
\end{proof}

The exponential factor in \eqref{eq:v4v78-Hadamard} must be included in the
charge weight of any sector sum.  A small instanton probability without
that mass/load factor is not the required estimate.

\subsection{Why a global charge moment is the wrong target}
\label{sec:v4v78-global-moment}

\begin{proposition}[Independent-bubble obstruction]
\label{prop:v4v78-global-moment-fails}
Consider \(N\) independent cells, each carrying charge
\(\xi_j\in\{-1,0,1\}\) with probabilities
\(p/2,1-p,p/2\), where \(0<p<1\), and put
\(Q_N=\sum_j\xi_j\).  For every \(\lambda>0\),
\begin{equation}
 \mathbb E e^{\lambda|Q_N|}
 \ge
 \mathbb E e^{\lambda Q_N}
 =
 (1-p+p\cosh\lambda)^N.
\label{eq:v4v78-charge-moment-growth}
\end{equation}
Hence no volume-uniform exponential moment of total absolute charge can be
required in a phase with a positive density of local topological
fluctuations.
\end{proposition}

\begin{proof}
Since \(|Q_N|\ge Q_N\),
\(e^{\lambda|Q_N|}\ge e^{\lambda Q_N}\).  Independence gives
\[
 \mathbb E e^{\lambda Q_N}
 =
 \prod_{j=1}^N
 \left(1-p+\frac p2e^\lambda+\frac p2e^{-\lambda}\right),
\]
which is \eqref{eq:v4v78-charge-moment-growth} and grows exponentially in
\(N\).
\end{proof}

\begin{firewall}[Bogomolny action does not remove volume entropy]
\label{fw:v4v78-Bogomolny-not-enough}
The continuum bound
\[
 S_{{\rm YM},r}\ge\frac{8\pi^2}{g_r^2}|Q_r|
\]
suppresses one fixed global sector.  By itself it does not bound the number
and placement entropy of widely separated instanton--anti-instanton
components as volume grows.  A local connected sum, or an equally strong
mixing theorem, is still required.
\end{firewall}

\subsection{Charged polymer representation}
\label{sec:v4v78-charged-polymers}

Let \(\mathbb L_a\) be the block lattice.  A charged polymer is
\[
 \mathfrak p=(X,q_2,q_3),
 \qquad
 X\Subset\mathbb L_a\text{ connected},
 \qquad q_2,q_3\in\mathbb Z.
\]
Its complete activity \(z_a(\mathfrak p)\) includes, on the same ancestry:
the local gauge--Higgs action change, hard determinant modulus and phase,
the appropriate V77 zero-mode exterior power or mass minor, and all
collar/seed transition factors.  Overlapping supports are incompatible.

For \(a_0,\lambda>0\), define
\begin{equation}
 \|z_a\|_{a_0,\lambda}
 :=
 \sup_{x\in\mathbb L_a}
 \sum_{\substack{\mathfrak p=(X,q_2,q_3)\\x\in X}}
 |z_a(\mathfrak p)|
 e^{a_0|X|+\lambda(|q_2|+|q_3|)}.
\label{eq:v4v78-charged-norm}
\end{equation}
Marked versions insert the actual first or second background derivative and
the inverse collar/seed load of V75.

\begin{definition}[Rooted charged-cluster certificate]
\label{def:v4v78-RCCC}
The finite regulator satisfies RCCC if the complete complex activities have:
\begin{enumerate}
\item exact support and charge ancestry;
\item \(\|z_a\|_{a_0,\lambda}\le\varepsilon_{\rm KP}\), with the same bound
      through two background marks;
\item the V75 weighted first moment for every collar/seed failure included
      in an activity; and
\item the usual Kotecky--Preiss smallness relation between
      \(\varepsilon_{\rm KP}\) and the incompatibility neighborhood.
\end{enumerate}
\end{definition}

The charge weight must dominate any deterministic zero-mode growth, for
example the \(6\log C_M\) per unit color charge in
\eqref{eq:v4v78-Hadamard}.

\begin{theorem}[Charge marking preserves the cluster proof]
\label{thm:v4v78-charge-KP}
Assume RCCC.  Introduce charge angles
\(\vartheta=(\vartheta_2,\vartheta_3)\) by
\begin{equation}
 z_a^\vartheta(X,q_2,q_3)
 =
 e^{i(\vartheta_2q_2+\vartheta_3q_3)}
 z_a(X,q_2,q_3).
\label{eq:v4v78-charge-source}
\end{equation}
Then the polymer logarithm converges absolutely and uniformly in real
\(\vartheta\), volume, and cutoff.  It extends analytically to
\[
 |\operatorname{Im}\vartheta_2|<\lambda,
 \qquad
 |\operatorname{Im}\vartheta_3|<\lambda,
\]
after reducing the available charge weight by the imaginary parts.
The first two declared background derivatives satisfy the same rooted
cluster bounds.
\end{theorem}

\begin{proof}
For real \(\vartheta\), the multiplier in
\eqref{eq:v4v78-charge-source} has modulus one, so every ordinary and marked
KP majorant is unchanged.  For complex \(\vartheta\),
\[
 |e^{i\vartheta_rq_r}|
 \le e^{|\operatorname{Im}\vartheta_r||q_r|},
\]
which is absorbed by the exponential charge factor in
\eqref{eq:v4v78-charged-norm}.  The standard connected Ursell expansion
then converges uniformly with the reduced positive exponent.  Termwise
differentiation through two marks is justified by the marked absolute
majorant already required in RCCC.
\end{proof}

\begin{corollary}[Local observables need only rooted topology]
\label{cor:v4v78-rooted-observable}
Let \(\mathcal O\) have fixed support \(A\).  Under RCCC, the connected
response of \(\log Z_a\) to \(\mathcal O\) is a sum only over clusters
connected to \(A\), and
\begin{equation}
 \sum_{\substack{\mathcal C:\,\mathcal C\leftrightarrow A}}
 |\operatorname{wt}(\mathcal C)|
 e^{\lambda|q(\mathcal C)|}
 \le C|A|
\label{eq:v4v78-rooted-bound}
\end{equation}
uniformly in the total volume.  Remote topological clusters cancel between
the numerator and denominator.
\end{corollary}

\begin{proof}
Differentiate the convergent logarithm with respect to a source coupled only
on \(A\).  In a connected cluster, the derivative is zero unless at least
one activity meets \(A\); connectedness then roots the complete cluster at
\(A\).  The KP tree estimate sums each additional activity by its
incompatibility neighborhood, while the charge exponential is
submultiplicative because
\[
 |q(\mathcal C)|\le\sum_{\mathfrak p\in\mathcal C}|q(\mathfrak p)|.
\]
The root has at most \(C|A|\) placements.  Clusters disjoint from \(A\)
occur identically in numerator and denominator and disappear upon taking
the logarithmic derivative.
\end{proof}

\begin{firewall}[An isolated weak instanton is a zero activity in vacuum]
\label{fw:v4v78-isolated-EW-zero}
For the uninserted Standard Model vacuum, a charged polymer with net
\(q_2\ne0\) has zero fermion coefficient by
Theorem~\ref{thm:v4v78-vacuum-selection}.  A weak
instanton--anti-instanton molecule with total \(q_2=0\) can nevertheless
have a nonzero determinant through zero-mode mixing.  It must be represented
as one irreducible neutral activity or one connected fermionic contraction,
not as the product of two nonzero one-instanton vacuum activities.
\end{firewall}

\subsection{Exact charge conservation in the cluster ledger}
\label{sec:v4v78-charge-ancestry}

\begin{lemma}[Additive charge and one owner]
\label{lem:v4v78-charge-add}
If the support extraction partitions a configuration into topological
components with declared relative-boundary gluing, the charge of a cluster
is
\begin{equation}
 q_r(\mathcal C)
 =
 \sum_{\mathfrak p\in\mathcal C}q_r(\mathfrak p),
 \qquad r=2,3.
\label{eq:v4v78-charge-add}
\end{equation}
Using the V75 first-full-connectivity order assigns every mixed
fermion/topology cluster once.
\end{lemma}

\begin{proof}
Chern number is additive under gluing along boundaries on which the
trivialization and orientation are fixed; boundary transgression terms
cancel pairwise.  The support owner and its first-join coordinate are unique
by Theorem~\ref{thm:v4v75-no-duplication}.  Therefore neither charge nor
activity is repeated.
\end{proof}

The relative-boundary gluing in this lemma is an actual hypothesis.  A
naive sum of plaquette principal logarithms across cut loci need not be an
integer and cannot substitute for it.

\begin{theorem}[Finite-sector and rooted-cluster compatibility]
\label{thm:v4v78-sector-cluster}
At finite volume, summing the charged polymer expansion at fixed total
charge \((m,n)\) reproduces the admissible sector contribution
\(Z_{a;m,n}\) in \eqref{eq:v4v78-Z-sector}, provided:
\begin{enumerate}
\item the polymer charge map is the restriction of the same established
      overlap/geometric index used in
      \eqref{eq:v4v78-sector};
\item relative boundary charges obey
      \eqref{eq:v4v78-charge-add}; and
\item the activity extraction is an exact partition of the original
      unnormalized integrand.
\end{enumerate}
Thus the global sector table and the local RCCC expansion are two
organizations of one occurrence, not two factors in the measure.
\end{theorem}

\begin{proof}
Under item 3, expanding the polymer product is a coefficient-one
decomposition of every admissible configuration.  Items 1--2 say that the
sum of its polymer charge labels equals its original sector label.  Grouping
terms by that sum gives exactly the level-set integral
\(Z_{a;m,n}\).  Regrouping an absolutely convergent finite-volume sum does
not change its value.
\end{proof}

\subsection{What is proved and what is an interface}
\label{sec:v4v78-status}

\begin{theorem}[V78 topological interface]
\label{thm:v4v78-interface}
The following conclusions are exact:
\begin{enumerate}
\item the admissible finite-cutoff sector partition
      \eqref{eq:v4v78-partition};
\item the weak-topology vacuum vanishing
      \eqref{eq:v4v78-vacuum-mzero};
\item the \(B,L\) observable selection
      \eqref{eq:v4v78-BL-selection};
\item the complete color zero-mode determinant
      \eqref{eq:v4v78-color-factor}; and
\item the equivalence of global sector grouping and an exact local charged
      activity extraction.
\end{enumerate}
If RCCC is verified for those exact activities, local normalized observables
and their first two marked responses have volume-uniform rooted topological
sums.  V78 does not verify RCCC from the Wilson--Higgs action.
\end{theorem}

\begin{proof}
The exact items are Lemma~\ref{lem:v4v78-sector-partition},
Theorem~\ref{thm:v4v78-vacuum-selection},
Corollary~\ref{cor:v4v78-observable-selection},
Proposition~\ref{prop:v4v78-color-factor}, and
Theorem~\ref{thm:v4v78-sector-cluster}.  The conditional implication is
Theorem~\ref{thm:v4v78-charge-KP} and
Corollary~\ref{cor:v4v78-rooted-observable}.
\end{proof}

\begin{assessment}[Bottleneck after V78]
\label{ass:v4v78-bottleneck}
The unusable target was a uniform moment of total charge.  The replacement
RCCC is local and connected, but its verification is a genuine
gauge--Higgs--fermion estimate.  The imported smooth admissible index and
concentration-ball determinant bounds give the correct algebraic pieces;
they do not yet prove a small charged activity under the full marginal
four-dimensional law.

The most dangerous test is a neutral weak
instanton--anti-instanton molecule.  Its net index vanishes, so the
electroweak determinant need not vanish, while its quasi-zero-mode overlap
can be very small and its separation entropy large.  That molecule should
be tested before a general charged KP theorem is attempted.
\end{assessment}

\subsection{V78 result ledger and V79 target}
\label{sec:v4v78-conclusion}

V78 constructs the exact \((Q_2,Q_3)\) sector ownership, proves its
fermionic selection rules and color mass factor, falsifies a
volume-uniform total-charge moment, and replaces it by a rooted
charge-marked cluster criterion.  The charge-source theorem gives a uniform
analytic strip and shows how remote topology cancels for local observables.

V79 should analyze one separated electroweak
instanton--anti-instanton pair at fixed finite regulator.  It must compress
the fermion operator to the two twelve-dimensional quasi-zero-mode spaces,
derive the exact overlap determinant by Schur complement, and compare its
separation decay with the collective-coordinate entropy and gauge--Higgs
action.  If the resulting molecule activity is not summable, the pair must
be promoted into the marginal parent rather than hidden in RCCC.
\clearpage
\section{V79: neutral electroweak molecules and quasi-zero-mode overlap}
\label{sec:v4v79}

V78 identifies a neutral electroweak instanton--anti-instanton molecule as
the first serious test of the rooted charged-cluster criterion.  Its total
weak index is zero, so the vacuum determinant need not vanish.  At large
separation, however, its only small fermion eigenvalues arise from the
overlap of the twelve instanton modes with the twelve anti-instanton modes.
The determinant of that overlap must beat the entropy of placing the second
center.

This note derives the exact finite-regulator Schur determinant.  Under a
declared two-marked localization envelope it proves a conservative
separation decay and shows that twelve Standard Model modes are more than
enough for summability in four dimensions.  The result controls separation
only; instanton size, shape, concentration probability, and the short-pair
core remain distinct inputs.

\subsection{Two approximate zero-mode spaces}
\label{sec:v4v79-spaces}

Fix one finite regulator and one bounded shape/size chart.  Let \(x_+\) and
\(x_-\) be the centers of charge \(+1\) and \(-1\) weak lumps, and set
\[
 R=d(x_+,x_-).
\]
Let
\[
 R_+:\mathbb C^r\to E_{\rm R},
 \qquad
 L_-:\mathbb C^r\to E_{\rm L},
 \qquad r=12,
\]
be oriented orthonormal frames obtained by transporting the isolated
instanton kernel and anti-instanton cokernel.  Put
\[
 P=R_+R_+^*,\qquad Q=L_-L_-^*.
\]
Their orthogonal complements are \(P_{\rm h}=1-P\) and
\(Q_{\rm h}=1-Q\).

For the full pair operator \(A_{+-}\), write
\begin{equation}
 A_{+-}=
 \begin{pmatrix}
 T&A_{\rm sh}\\
 A_{\rm hs}&A_{\rm hh}
 \end{pmatrix},
 \qquad
 T:=L_-^*A_{+-}R_+.
\label{eq:v4v79-pair-block}
\end{equation}

\begin{theorem}[Exact molecule determinant]
\label{thm:v4v79-exact-molecule}
If \(A_{\rm hh}\) is invertible, define the quasi-zero-mode Schur matrix
\begin{equation}
 S_{+-}
 :=
 T-A_{\rm sh}A_{\rm hh}^{-1}A_{\rm hs}.
\label{eq:v4v79-molecule-Schur}
\end{equation}
Then, in the determinant-line orientation transported from V74,
\begin{equation}
 \det A_{+-}
 =
 \det A_{\rm hh}\,\det S_{+-}.
\label{eq:v4v79-molecule-det}
\end{equation}
This identity remains valid when \(S_{+-}\) is singular.
\end{theorem}

\begin{proof}
The left and right soft dimensions both equal twelve because the total pair
index is zero.  Apply the exact Berezin--Schur identity
Theorem~\ref{thm:v4v73-Berezin-Schur} to
\eqref{eq:v4v79-pair-block}.  V74 supplies the determinant-line
factorization and orientation at zeros.
\end{proof}

\begin{firewall}[The raw overlap is not always the Schur matrix]
\label{fw:v4v79-raw-overlap}
The matrix \(T\) equals the complete soft operator only when the selected
subspaces reduce \(A_{+-}\).  For transported isolated modes they generally
do not.  The response
\(A_{\rm sh}A_{\rm hh}^{-1}A_{\rm hs}\) in
\eqref{eq:v4v79-molecule-Schur} is mandatory and occurs once.
\end{firewall}

\subsection{Residual-to-Schur estimate}
\label{sec:v4v79-residual}

\begin{lemma}[Approximate kernels control every soft block]
\label{lem:v4v79-residual}
Assume
\begin{equation}
 \|A_{+-}R_+\|\le\varepsilon(R),
 \qquad
 \|A_{+-}^*L_-\|\le\varepsilon(R),
 \qquad
 \|A_{\rm hh}^{-1}\|\le\delta^{-1}.
\label{eq:v4v79-residual}
\end{equation}
Then
\begin{equation}
 \|S_{+-}\|
 \le
 \varepsilon(R)+\delta^{-1}\varepsilon(R)^2.
\label{eq:v4v79-Schur-residual}
\end{equation}
\end{lemma}

\begin{proof}
The three soft-containing blocks obey
\[
 \|T\|\le\|A_{+-}R_+\|\le\varepsilon,
 \qquad
 \|A_{\rm hs}\|
 =\|Q_{\rm h}A_{+-}R_+\|\le\varepsilon,
\]
and
\[
 \|A_{\rm sh}\|
 =\|L_-^*A_{+-}P_{\rm h}\|
 =\|P_{\rm h}A_{+-}^*L_-\|
 \le\varepsilon.
\]
Insert these estimates into
\eqref{eq:v4v79-molecule-Schur}.
\end{proof}

The same proof applies to first and second background marks, using the
product rule and the inverse identities
\[
 D A_{\rm hh}^{-1}
 =-A_{\rm hh}^{-1}(DA_{\rm hh})A_{\rm hh}^{-1}
\]
and its exact two-mark derivative.  A marked conclusion therefore requires
marked residual bounds, not only
\eqref{eq:v4v79-residual}.

\subsection{A conservative overlap convolution}
\label{sec:v4v79-convolution}

Let the four-dimensional cutoff lattice or bounded-geometry block graph
have volume growth
\[
 |B(x,\rho)|\le C_V(1+\rho)^4.
\]

\begin{lemma}[Algebraic two-center convolution]
\label{lem:v4v79-convolution}
For \(2<s<4\), there is \(C_s\) such that
\begin{equation}
 \sum_x
 \frac1{(1+d(x,x_+))^s}
 \frac1{(1+d(x,x_-))^s}
 \le
 \frac{C_s}{(1+R)^{2s-4}}.
\label{eq:v4v79-convolution}
\end{equation}
\end{lemma}

\begin{proof}
Let \(R\ge2\); bounded \(R\) is absorbed into the constant.  Split the graph
into the regions
\[
 d(x,x_+)\le R/2,\qquad
 d(x,x_-)\le R/2,
\]
and their complement.  On the first region the second factor is
\(O(R^{-s})\), while shell summation and four-dimensional volume growth give
\[
 \sum_{\rho\le R/2}(1+\rho)^{3-s}
 \le C R^{4-s}.
\]
This contributes \(CR^{4-2s}\); the second region is identical.  On the
complement, decompose by dyadic values of the smaller distance.  Both
distances are at least \(R/2\), and the shell count at scale \(2^jR\) is
\(O((2^jR)^4)\), while the product is
\(O((2^jR)^{-2s})\).  The geometric sum converges because \(2s>4\) and has
the same order \(R^{4-2s}\).
\end{proof}

\begin{theorem}[Profile localization gives residual decay]
\label{thm:v4v79-profile-residual}
Assume every column of \(R_+\) and \(L_-\), and the columns produced by its
first two declared background marks, obey the uniform envelope
\begin{equation}
 |\psi_\pm(x)|
 \le
 \frac{C_\psi}{(1+d(x,x_\pm))^s},
 \qquad 2<s<4,
\label{eq:v4v79-profile}
\end{equation}
and assume the rescaled finite-range coefficients of \(A_{+-}\) and their
two marks are uniformly bounded on the size chart.  If the isolated-mode
residual is supported by the opposite-lump perturbation with the same
envelope, then
\begin{equation}
 \|D_B^j(A_{+-}R_+)\|
 +
 \|D_B^j(A_{+-}^*L_-)\|
 \le
 C_j(1+R)^{-(2s-4)},
 \qquad 0\le j\le2.
\label{eq:v4v79-residual-decay}
\end{equation}
\end{theorem}

\begin{proof}
The isolated operator annihilates its own zero-mode frame.  The remaining
residual contains one localized opposite-lump coefficient and one isolated
profile, up to a uniformly finite number of neighboring lattice sites.
Every matrix coefficient is therefore bounded by a fixed finite sum of the
two-center convolution in
Lemma~\ref{lem:v4v79-convolution}.  First and second marks distribute among
the operator coefficient and the transported profiles; the hypothesis
gives the same envelopes for every resulting factor.  Finite internal
dimension converts the entry bounds to the displayed operator bounds.
\end{proof}

For the four-dimensional massless instanton envelope \(s=3\), this
conservative argument gives
\begin{equation}
 \varepsilon(R)\le C(1+R)^{-2}.
\label{eq:v4v79-Rminus2}
\end{equation}
The sharper semiclassical \(R^{-3}\) overlap is not needed.  If the
Yukawa--Higgs background gives exponential profile decay, the same argument
with exponential weights gives \(Ce^{-mR}\).

\begin{firewall}[The profile hypothesis is not derived from action alone]
\label{fw:v4v79-profile-open}
V79 does not infer \eqref{eq:v4v79-profile} from a Wilson--Higgs action
bound.  It is a two-marked isolated-lump spectral input on one bounded
size/shape chart.  Concentration, rough backgrounds, and degeneration of
the hard complement remain in the V75 weighted failure ledger.
\end{firewall}

\subsection{Determinant decay without an inverse soft matrix}
\label{sec:v4v79-det-decay}

\begin{lemma}[Marked determinant bound]
\label{lem:v4v79-marked-det}
Let \(S(B)\) be an \(r\times r\) matrix.  Suppose every entry of
\(D_B^jS\), \(0\le j\le2\), is bounded by \(C_jf(R)\).  Then
\begin{align}
 |\det S|
 &\le C_{r,0}f(R)^r,
\label{eq:v4v79-det-bound}\\
 |D_B\det S|
 &\le C_{r,1}f(R)^r,
\label{eq:v4v79-det-first}\\
 |D_B^2\det S|
 &\le C_{r,2}f(R)^r.
\label{eq:v4v79-det-second}
\end{align}
No inverse of \(S\) is used.
\end{lemma}

\begin{proof}
The determinant is multilinear in its \(r\) columns.  Hadamard's inequality
bounds every determinant whose columns are selected from
\(S,DS,D^2S\) by the product of their Euclidean column norms.  The first
derivative is a sum of \(r\) determinants with one marked column.  The
second derivative is a sum of \(r\) terms with one twice-marked column and
\(r(r-1)\) ordered terms with two once-marked columns.  Every term contains
exactly \(r\) columns and hence \(f(R)^r\).
\end{proof}

\begin{theorem}[Twelve-mode molecule decay]
\label{thm:v4v79-molecule-decay}
Assume the two-marked hypotheses of
Theorem~\ref{thm:v4v79-profile-residual}, the uniform hard gap in
\eqref{eq:v4v79-residual}, and \(R\) large enough that
\(\varepsilon(R)\le\delta\).  With
\(\alpha=2s-4>0\), the molecule soft determinant satisfies
\begin{equation}
 |D_B^j\det S_{+-}|
 \le
 C_j(1+R)^{-12\alpha},
 \qquad 0\le j\le2.
\label{eq:v4v79-molecule-decay}
\end{equation}
For \(s=3\),
\begin{equation}
 |D_B^j\det S_{+-}|
 \le C_j(1+R)^{-24}.
\label{eq:v4v79-molecule-R24}
\end{equation}
\end{theorem}

\begin{proof}
The marked inverse identities and
Lemma~\ref{lem:v4v79-residual} show that every entry of \(S_{+-}\) and its
first two marks is bounded by \(C(1+R)^{-\alpha}\); the quadratic Schur
response decays at least as fast after
\(\varepsilon\le\delta\).  Apply
Lemma~\ref{lem:v4v79-marked-det} with \(r=12\).
\end{proof}

The power twelve is physical multiplicity, not twelve independent uses of
the gauge action.  It arises once from the exterior power of the complete
quasi-zero-mode matrix.

\subsection{Separation entropy}
\label{sec:v4v79-entropy}

\begin{lemma}[Rooted placement sum]
\label{lem:v4v79-placement}
On a graph of polynomial volume growth degree four,
\begin{equation}
 \sup_{x_+}
 \sum_{x_-}
 (1+d(x_+,x_-))^{-\beta}
 <\infty
\label{eq:v4v79-placement}
\end{equation}
whenever \(\beta>4\).  Its tail tends to zero uniformly as the lower
separation cutoff tends to infinity.
\end{lemma}

\begin{proof}
The number of sites in a shell of radius \(R\) is at most \(C(1+R)^3\).
Thus the sum is bounded by
\(\sum_R C(1+R)^{3-\beta}\), which converges for \(\beta>4\).  The same
series proves the uniform tail statement.
\end{proof}

\begin{theorem}[Large-separation molecule is rooted-summable]
\label{thm:v4v79-rooted-summable}
Under Theorem~\ref{thm:v4v79-molecule-decay}, the position sum of the
molecule soft determinant and its first two marks is uniform in volume
provided
\begin{equation}
 12(2s-4)>4,
 \quad\text{equivalently}\quad
 s>\frac{13}{6}.
\label{eq:v4v79-summability-threshold}
\end{equation}
In particular the conservative \(s=3\) envelope is strongly summable.
For every desired tail parameter \(\epsilon>0\), a finite
\(R_{\rm mol}\) can be chosen so the rooted sum over \(R>R_{\rm mol}\) is
at most \(\epsilon\).
\end{theorem}

\begin{proof}
Use \eqref{eq:v4v79-molecule-decay} with
\(\beta=12(2s-4)\), then apply
Lemma~\ref{lem:v4v79-placement}.  The tail assertion is the final statement
of that lemma.
\end{proof}

\subsection{Core--tail molecule assignment}
\label{sec:v4v79-core-tail}

Define the short molecule core by \(R\le R_{\rm mol}\) and the separated
tail by \(R>R_{\rm mol}\).

\begin{theorem}[Molecule interface with RCCC]
\label{thm:v4v79-RCCC}
On every bounded size/shape and spectral-collar chart satisfying the V79
hypotheses:
\begin{enumerate}
\item the short neutral molecule is a generally order-one rooted
      fermion--gauge--Higgs carrier; it may be promoted to the marginal
      parent only together with a proved quasilocal support decomposition
      or a separately paid spatial truncation;
\item the separated molecule determinant and its first two marks have an
      arbitrarily small rooted position sum after increasing
      \(R_{\rm mol}\);
\item the hard determinant and soft determinant occur once through
      \eqref{eq:v4v79-molecule-det}; and
\item the molecule carries net weak charge zero but retains the internal
      \(+1,-1\) ancestry needed for zero-mode overlap.
\end{enumerate}
Thus separation entropy alone does not obstruct the RCCC tail.
\end{theorem}

\begin{proof}
Fixing \(R_{\rm mol}\) bounds the distance between the two centers, but it
does not compactly support their zero-mode or gauge tails.  Thus the core is
one rooted order-one carrier, and its promotion requires the additional
support decomposition stated in item 1; no smallness follows from
separation.  Theorem~\ref{thm:v4v79-rooted-summable} proves item 2.
Theorem~\ref{thm:v4v79-exact-molecule} proves item 3.  Additivity of the
topological charge gives item 4, while deleting the internal labels would
lose the definition of the two quasi-zero-mode frames.
\end{proof}

\begin{firewall}[Short separation is not compact support]
\label{fw:v4v79-short-not-compact}
The condition \(R\le R_{\rm mol}\) controls only the two centers.  Algebraic
instanton and zero-mode profiles still extend outside every finite block.
A compact polymer support requires an exact localization partition or a
truncation whose discarded tail is entered once in the activity ledger.
V79 does not obtain finite support by relabelling a short pair as a core.
\end{firewall}

\begin{firewall}[Fixed size chart only]
\label{fw:v4v79-size}
The rooted sum in Theorem~\ref{thm:v4v79-RCCC} fixes a bounded instanton
size and shape chart.  It does not sum small-size ultraviolet bubbles,
large diffuse lumps, gauge orientations, or chart-transition inverse
loads.  Those integrations can change the activity before the separation
sum and are the next bottleneck.
\end{firewall}

\begin{firewall}[Upper decay gives no molecule lower bound]
\label{fw:v4v79-no-lower}
The determinant estimate is an upper bound suitable for a KP tail.
Orientation or generation mixing can make \(\det S_{+-}\) vanish.  V79
does not prove that the neutral molecule supplies a nonzero vacuum
normalization or a positive effective interaction.
\end{firewall}

\begin{assessment}[What the molecule test decides]
\label{ass:v4v79-status}
The feared large-separation entropy is not the immediate obstruction under
localized Standard Model zero-mode profiles.  Even the conservative
four-dimensional convolution gives \(R^{-24}\), which is summable after
one center is rooted.  The exact Schur response preserves that conclusion
under a uniform hard gap and two marked residual bounds.

The unresolved directions are ultraviolet size integration and failure of
the spectral/profile chart on rough fields.  They cannot be repaired by a
sharper large-\(R\) asymptotic.
\end{assessment}

\subsection{V79 result ledger and compulsory V80 target}
\label{sec:v4v79-conclusion}

V79 proves the exact neutral-pair Schur determinant, controls it from
approximate-kernel residuals, derives a conservative two-center convolution,
and shows that twelve electroweak quasi-zero modes beat
four-dimensional separation entropy through two background marks.  It
splits the molecule into an order-one local core and an arbitrarily small
separated tail.

V80 is a compulsory Yang--Mills/Navier--Stokes audit.  After recording
stable companion snapshots, it should test whether the newest YM
all-order/marked-tree work supplies the missing lump-shape aggregation.
It must then analyze the instanton-size scale.  The first calculation is
the small-size power count with the actual Standard Model one-loop
coefficients and twelve zero-mode/external factors, followed by a firewall
separating that semiclassical test from a proof under the full interacting
lattice measure.
\clearpage
\section{V80: compulsory companion audit and the instanton-size exponent ledger}
\label{sec:v4v80}

V79 proves large-separation summability of one neutral electroweak molecule
on a bounded size/shape chart.  The missing sum is over the lump size
\(\rho\), especially \(\rho\downarrow0\) as the lattice cutoff is removed.
This compulsory audit first records the current Yang--Mills and
Navier--Stokes endpoints.  It then computes the Standard Model one-loop
size exponents as a falsification test and states the marked dyadic estimate
that an actual interacting proof must establish.

\subsection{Stable V80 companion audit}
\label{sec:v4v80-audit}

Both companion files were hashed before and after the read and were stable:
\begin{center}
\begin{tabular}{|p{0.50\textwidth}|r|p{0.31\textwidth}|}
\hline
Companion file & Bytes & SHA--256\\
\hline
\texttt{Renewal\_Yang\_Mills\_Mass\_Gap\_Research\_Notes\_V100.tex}
& \(1\,720\,945\)
& \texttt{A5E05A144C753426D8A461C8290E9538}\\
&&\texttt{A4CD19D8672BAB8B08B509D82FAF892F}\\
\hline
\texttt{Renewal\_NS\_Stability\_Research\_Notes\_V31.tex}
& \(887\,537\)
& \texttt{F1121B62238AD5491BB7CF03635EA993}\\
&&\texttt{4942EDF573ED8FAAA9EE79E1D38A6696}\\
\hline
\end{tabular}
\end{center}

Yang--Mills V100 proves that four source-occurring complete cofinal heat
slots supply the missing \(M^{-2}\) response in its quintic square.  With
only \(q=0,1,2,3\) such slots, it isolates the exact remaining assembled
quotient-cumulant response \(M^{-(4-q)/2}\) and proves that heat-only scalar
bounds cannot manufacture it.  The applicable lesson is that the twelve
fermion zero modes must remain one complete exterior-power determinant.
A size power may be used only when the actual zero-mode or mass factors
occur in that exact source; one cannot attach twelve scalar improvements
after splitting the determinant.

Navier--Stokes V31 again supplies the scale warning: flat dyadic marks become
unweighted shell action only after multiplication by their exact radial
first-moment weight.  Here a statement that each small-size annulus is rare
does not sum the annuli.  The full power of \(\rho\), including derivative
losses, must be positive after the annular measure is included.

\subsection{Dyadic size summation}
\label{sec:v4v80-dyadic}

Fix a macroscopic reference size \(\rho_0\) and let
\[
 I_j=(2^{-j-1}\rho_0,2^{-j}\rho_0],
 \qquad j\ge0.
\]

\begin{lemma}[Exact ultraviolet power test]
\label{lem:v4v80-power-test}
If a nonnegative size density satisfies
\begin{equation}
 f(\rho)\le C\rho^p
 \qquad(0<\rho\le\rho_0),
\label{eq:v4v80-rho-power}
\end{equation}
then
\begin{equation}
 \int_{I_j}f(\rho)\,d\rho
 \le C_p\rho_0^{p+1}2^{-j(p+1)}
\label{eq:v4v80-bin-power}
\end{equation}
when \(p>-1\).  Consequently
\[
 \int_0^{\rho_0}f(\rho)\,d\rho<\infty
 \quad\Longleftrightarrow\quad p>-1
\]
for an exact positive monomial \(f(\rho)=c\rho^p\).
At \(p=-1\) the dyadic bins have constant size and produce a logarithm.
\end{lemma}

\begin{proof}
Integrate \(\rho^p\) over \(I_j\):
\[
 \int_{2^{-j-1}\rho_0}^{2^{-j}\rho_0}\rho^p\,d\rho
 =
 \frac{\rho_0^{p+1}}{p+1}
 2^{-j(p+1)}(1-2^{-(p+1)}).
\]
This proves the estimate and the geometric-series criterion.  For
\(p=-1\), the integral of \(d\rho/\rho\) over every bin is \(\log2\).
\end{proof}

\begin{corollary}[Two-mark exponent test]
\label{cor:v4v80-mark-test}
Suppose the \(k\)-th declared background mark loses at most
\(\rho^{-\sigma_k}\), \(k=0,1,2\), from an unmarked density
\(C\rho^p\), with \(\sigma_0=0\).  A sufficient uniform ultraviolet
condition through two marks is
\begin{equation}
 p-\sigma_k>-1,
 \qquad k=0,1,2.
\label{eq:v4v80-mark-power}
\end{equation}
\end{corollary}

\begin{proof}
Apply Lemma~\ref{lem:v4v80-power-test} to each marked density
\(C_k\rho^{p-\sigma_k}\).
\end{proof}

The losses \(\sigma_k\) are actual derivatives of the complete activity.
They are not set to zero because the unmarked instanton density is
integrable.

\subsection{One-loop collective-coordinate input}
\label{sec:v4v80-one-loop}

The semiclassical one-lump size measure for a simple gauge factor has the
form
\begin{equation}
 d\rho\,
 \rho^{-5}(\mu\rho)^{b_r}
 e^{-8\pi^2/g_r(\mu)^2}
 \,\mathcal Z_{0,r}(\rho)\,
 \mathcal R_r(\rho),
\label{eq:v4v80-semiclassical-measure}
\end{equation}
where \(\mathcal R_r\) is dimensionless at the leading scaling level and
the already proved Standard Model coefficients are
\begin{equation}
 b_2=\frac{19}{6},
 \qquad
 b_3=7
\label{eq:v4v80-betas}
\end{equation}
from \eqref{eq:v4v6-SM-b}.  The factor \(\rho^{-5}\) is the
four-position/one-scale collective Jacobian after the center is retained
separately.

\begin{lemma}[Engineering factors from zero-mode saturation]
\label{lem:v4v80-zero-mode-scaling}
At the leading scale-covariant level:
\begin{enumerate}
\item a local operator containing \(N\) external Weyl fields contributes
      \(\rho^{3N/2}\);
\item \(N_f\) matched Dirac zero-mode pairs saturated by a mass matrix
      contribute
\begin{equation}
 \rho^{N_f}\det M
\label{eq:v4v80-mass-scaling}
\end{equation}
with the determinant interpreted on the \(N_f\)-dimensional flavor
space; and
\item a dimensionless gauge orientation integral changes neither power.
\end{enumerate}
\end{lemma}

\begin{proof}
A four-dimensional Weyl field has engineering dimension \(3/2\).
After the instanton center is factored as \(d^4x_0\), a local \(N\)-field
vertex must therefore carry \(\rho^{3N/2}\) to be dimensionless relative to
the collective measure.  For a normalized paired zero mode, the restricted
mass matrix has dimension one.  The natural nonzero-mode reference scale is
\(\rho^{-1}\), so each dimensionless lifted eigenvalue is \(\rho M\).
Taking the top coefficient of \(N_f\) pairs gives
\(\rho^{N_f}\det M\).  Haar measure on a compact orientation group is
dimensionless.
\end{proof}

\begin{firewall}[Engineering power is not a determinant theorem]
\label{fw:v4v80-engineering-only}
Lemma~\ref{lem:v4v80-zero-mode-scaling} determines the only possible leading
power after a semiclassical isolated-lump reduction.  It does not prove the
collective-coordinate formula, uniform control of
\(\mathcal R_r\), absence of logarithms, or validity on rough interacting
fields.  The exact finite-regulator determinant remains the V73/V77 Schur
object.
\end{firewall}

\subsection{Electroweak external sector}
\label{sec:v4v80-EW}

For a unit weak instanton, the saturating Standard Model observable has
\(N=12\) Weyl fields by
Theorem~\ref{thm:v4v76-EW-count}.  Thus
\begin{equation}
 \mathcal Z_{0,2}^{\rm ext}(\rho)
 \sim \rho^{18}\mathcal O_{\rm EW}.
\label{eq:v4v80-EW-zero-factor}
\end{equation}

\begin{proposition}[Electroweak small-size test]
\label{prop:v4v80-EW-power}
The leading one-loop size power for the twelve-field electroweak
coefficient is
\begin{equation}
 p_{\rm EW}
 =
 -5+\frac{19}{6}+18
 =
 \frac{97}{6}.
\label{eq:v4v80-EW-power}
\end{equation}
Hence its unmarked small-size integral is strongly convergent, with dyadic
gain
\begin{equation}
 p_{\rm EW}+1=\frac{103}{6}.
\label{eq:v4v80-EW-gain}
\end{equation}
Through two background marks, this test remains convergent if
\begin{equation}
 \sigma_k^{\rm EW}<\frac{103}{6},
 \qquad k=1,2.
\label{eq:v4v80-EW-mark-margin}
\end{equation}
\end{proposition}

\begin{proof}
Insert \(b_2=19/6\) and the twelve-field factor \(\rho^{18}\) into
\eqref{eq:v4v80-semiclassical-measure}.  Apply
Lemma~\ref{lem:v4v80-power-test} and
Corollary~\ref{cor:v4v80-mark-test}.
\end{proof}

The corresponding uninserted fixed-\(Q_2=1\) vacuum coefficient is exactly
zero by Theorem~\ref{thm:v4v77-EW-vacuum-zero}; it is not assigned the
bare divergent power \(-5+19/6\).

\subsection{Color mass and external sectors}
\label{sec:v4v80-QCD}

For a unit color instanton, six matched Dirac zero-mode pairs have the exact
compressed mass determinant
\(\det M_u(H)\det M_d(H)\).  The dimensionless saturation factor is
\begin{equation}
 \rho^6\det M_u(H)\det M_d(H).
\label{eq:v4v80-QCD-mass-factor}
\end{equation}

\begin{proposition}[QCD massive small-size test]
\label{prop:v4v80-QCD-mass-power}
The leading one-loop size power in the massive color vacuum sector is
\begin{equation}
 p_{\rm QCD,m}
 =
 -5+7+6=8.
\label{eq:v4v80-QCD-mass-power}
\end{equation}
It has dyadic gain \(9\), and remains integrable through two marks if
\begin{equation}
 \sigma_k^{\rm QCD,m}<9,
 \qquad k=1,2.
\label{eq:v4v80-QCD-mass-marks}
\end{equation}
\end{proposition}

\begin{proof}
Use \(b_3=7\) and
\eqref{eq:v4v80-QCD-mass-factor} in the measure
\eqref{eq:v4v80-semiclassical-measure}, then apply the marked power test.
\end{proof}

If the six quark mass pairs are not used and the twelve-field QCD
't Hooft operator is inserted instead, the leading power is
\begin{equation}
 p_{\rm QCD,ext}
 =
 -5+7+18=20,
\qquad
 p_{\rm QCD,ext}+1=21.
\label{eq:v4v80-QCD-external-power}
\end{equation}

\begin{firewall}[The Higgs is still integrated]
\label{fw:v4v80-Higgs-integrated}
The factor
\(\det M_u(H)\det M_d(H)\) is a degree-six Higgs--Yukawa polynomial before
a symmetry-breaking background is fixed.  Replacing it by physical masses
inside the microscopic SM--Einstein measure would remove the correlated
Higgs integration.  V80 uses physical masses only as a broken-phase
semiclassical test.
\end{firewall}

\subsection{Neutral-pair size test}
\label{sec:v4v80-pair}

For an instanton--anti-instanton pair, the standard scale-covariant overlap
entry has mass dimension one and, at large separation, the schematic form
\[
 T_{+-}\sim\frac{\rho_+\rho_-}{R^3}
\]
times a dimensionless orientation matrix.  V79 did not need this sharper
formula; it used the conservative separation envelope.

\begin{proposition}[Conditional pair-size margin]
\label{prop:v4v80-pair-size}
Assume on a bounded-shape pair chart that the exact V79 Schur matrix obeys,
through two marks,
\begin{equation}
 \|D_B^kS_{+-}\|
 \le
 C_k
 \frac{\rho_+\rho_-}{(R+\rho_++\rho_-)^3},
 \qquad k=0,1,2.
\label{eq:v4v80-pair-overlap-scale}
\end{equation}
Then its twelve-dimensional determinant contributes at least the explicit
size factor
\begin{equation}
 (\rho_+\rho_-)^{12}
\label{eq:v4v80-pair-size-factor}
\end{equation}
to the absolute upper bound.  Combining it with two one-loop weak
collective measures gives the separate small-size exponent
\begin{equation}
 -5+\frac{19}{6}+12
 =
 \frac{61}{6}
\label{eq:v4v80-pair-one-size}
\end{equation}
for each \(\rho_\pm\), before marked size losses.  Thus the conditional
pair-size test is integrable with gain \(67/6\) in each size variable.
\end{proposition}

\begin{proof}
Hadamard's inequality applied to the twelve columns of
\eqref{eq:v4v80-pair-overlap-scale} gives
\[
 |\det S_{+-}|
 \le
 C
 \frac{(\rho_+\rho_-)^{12}}
 {(R+\rho_++\rho_-)^{36}}.
\]
The marked determinant proof of
Lemma~\ref{lem:v4v79-marked-det} gives the same size factor for two marks
under the stated hypothesis.  Multiply by
\(\rho_\pm^{-5+19/6}\) for each collective measure and apply
Lemma~\ref{lem:v4v80-power-test}.
\end{proof}

\begin{firewall}[The pair scaling is a new analytic input]
\label{fw:v4v80-pair-scaling-open}
Equation \eqref{eq:v4v80-pair-overlap-scale} is stronger than V79's
fixed-size localization theorem and has not been proved uniformly as
\(\rho_\pm\downarrow a\).  Its favorable power count is a target for the
rescaled zero-mode and Schur-response analysis, not an established RCCC
bound.
\end{firewall}

\subsection{Actual marked size certificate}
\label{sec:v4v80-certificate}

Let \(z_{a,j}^{(k)}(X,q)\) be the complete \(k\)-marked charged activity
whose smallest lump size lies in \(I_j\), including its exact determinant
or external exterior power and every assigned frame/collar load.

\begin{definition}[Marked ultraviolet size certificate]
\label{def:v4v80-MUSC}
MUSC holds with exponent \(\eta>0\) if, for \(k=0,1,2\) and every
\(j\ge0\),
\begin{equation}
 \sup_{a}\sup_x
 \sum_{\substack{X\ni x,\ q}}
 e^{a_0|X|+\lambda|q|}
 |z_{a,j}^{(k)}(X,q)|
 \le
 C_k2^{-\eta j}.
\label{eq:v4v80-MUSC}
\end{equation}
Consequently the sum over \(j\) is bounded by
\(C_k\sum_{j\ge0}2^{-\eta j}<\infty\).
Each activity is assigned to its smallest-size annulus with a fixed
half-open tie rule.
\end{definition}

\begin{theorem}[MUSC closes the size part of RCCC]
\label{thm:v4v80-MUSC-RCCC}
If the fixed-size/shape charged activities satisfy the V78 RCCC support and
incompatibility estimates uniformly, and MUSC holds for their size
decomposition, then summing all ultraviolet size annuli preserves RCCC
through two marks.  No logarithmic number of cutoff scales appears.
\end{theorem}

\begin{proof}
The half-open smallest-size rule assigns every activity to one \(j\).
Sum its ordinary and marked charged norms using
\eqref{eq:v4v80-MUSC}.  The geometric series is independent of the terminal
lattice annulus, so its bound is uniform as \(a\downarrow0\).  Insert the
summed norm into the KP criterion of
Definition~\ref{def:v4v78-RCCC}.
\end{proof}

\begin{firewall}[One-loop power does not prove MUSC]
\label{fw:v4v80-one-loop-not-MUSC}
Propositions
\ref{prop:v4v80-EW-power}--%
\ref{prop:v4v80-pair-size}
test the ultraviolet exponent on smooth semiclassical charts.  MUSC also
requires uniform shape aggregation, rough-field ownership, exact complex
determinants, two background marks, and the interacting gauge--Higgs
density.  None follows from the sign of \(b_2\) or \(b_3\) alone.
\end{firewall}

\subsection{Audit-adjusted direction}
\label{sec:v4v80-direction}

\begin{theorem}[V80 size ledger]
\label{thm:v4v80-ledger}
The semiclassical Standard Model ultraviolet powers are:
\begin{center}
\begin{tabular}{c|c|c}
sector & density exponent \(p\) & dyadic gain \(p+1\)\\
\hline
weak \(Q_2=1\), twelve external fields
& \(97/6\) & \(103/6\)\\
color \(Q_3=1\), six mass pairs
& \(8\) & \(9\)\\
color \(Q_3=1\), twelve external fields
& \(20\) & \(21\)\\
neutral weak pair, each size, conditional
& \(61/6\) & \(67/6\)
\end{tabular}
\end{center}
Every unmarked exponent is ultraviolet integrable.  The rigorous continuum
input is nevertheless MUSC, with the actual two-mark losses and complete
assembled determinant retained.
\end{theorem}

\begin{proof}
The four rows are
Propositions~\ref{prop:v4v80-EW-power},
\ref{prop:v4v80-QCD-mass-power},
equation~\eqref{eq:v4v80-QCD-external-power}, and
Proposition~\ref{prop:v4v80-pair-size}.  The final statement is
Theorem~\ref{thm:v4v80-MUSC-RCCC} together with the one-loop firewall.
\end{proof}

\begin{assessment}[What the companion audit changes]
\label{ass:v4v80-status}
YM V100 rules out extracting independent scalar gains from the twelve
zero-mode factors after the exterior-power operator has been split.  NS V31
rules out paying infinitely many size annuli by an unweighted rarity mark.
Accordingly, MUSC is stated for the complete marked activity and includes
one geometric size gain after all losses.

The favorable one-loop powers show no elementary small-instanton divergence
for the displayed saturated Standard Model sectors.  The remaining problem
is proving a uniform rescaled determinant/profile theorem and aggregating
all shapes under the actual marginal gauge--Higgs law.
\end{assessment}

\subsection{V80 result ledger and V81 target}
\label{sec:v4v80-conclusion}

V80 records stable YM V100 and NS V31 snapshots, proves the dyadic
ultraviolet power and two-mark tests, computes the Standard Model
electroweak and color size exponents, and formulates MUSC as the exact
size-summation input to RCCC.  It keeps the neutral-pair size exponent
conditional on a declared rescaled Schur bound.

V81 should prove or falsify that rescaled bound.  On a unit ball, it should
combine the imported Uhlenbeck \(L^4\) control and fifth regularized
determinant with the V73 singular projector.  The target is a scale-uniform
two-mark estimate for the finite low spectral band and its Schur response.
If \(L^4\) control cannot bound spectral-projector marks near zero, the
collar failure must be charged by a concentration stopping rule rather than
by smooth instanton coordinates.
\clearpage
\section{V81: rescaled \(L^4\) charts, marked low bands, and the spectral-flow stop}
\label{sec:v4v81}

V80 asks whether the scale-critical Uhlenbeck estimate can prove the
two-marked low-band bound required by MUSC.  The answer has two parts.  On a
small-energy ball, the Birman--Schwinger operator is norm-small on a fixed
reference contour; the low spectral projector, its first two marks, and the
fifth regularized determinant are then controlled uniformly after
rescaling.  On a concentration core, a bounded \(L^4\) norm without
smallness does not prevent a zero crossing.  Such a core must change
spectral-section chart and retain its crossing line.

\subsection{Dimensionless unit-ball operator}
\label{sec:v4v81-rescaling}

On a ball \(B_\rho(x_0)\), use \(y=(x-x_0)/\rho\) and set
\[
 A_\rho(y)=\rho A(x_0+\rho y),\qquad
 \widehat D_{A,\rho}=\rho D_A.
\]
In four dimensions,
\begin{equation}
 \|A_\rho\|_{L^4(B_1)}
 =
 \|A\|_{L^4(B_\rho)}.
\label{eq:v4v81-L4-scale}
\end{equation}
Choose one self-adjoint hermitization \(\widehat H_0\) of the reference
Dirac boundary problem and write
\begin{equation}
 \widehat H_A=\widehat H_0+V_A,
 \qquad V_A=c(A_\rho).
\label{eq:v4v81-H}
\end{equation}
The chiral determinant line remains labelled separately; hermitization is
used only for singular bands.

Let \(\Gamma\) be a contour with
\begin{equation}
 \operatorname{dist}(\Gamma,\operatorname{spec}\widehat H_0)
 \ge\gamma_0>0
\label{eq:v4v81-reference-contour}
\end{equation}
and let \(R_0(z)=(z-\widehat H_0)^{-1}\).

\begin{lemma}[Critical Birman--Schwinger bound]
\label{lem:v4v81-BS}
Uniform unit-ball elliptic and Sobolev estimates imply
\begin{equation}
 \|V_AR_0(z)\|_{2\to2}
 \le C_{\Gamma}\|A_\rho\|_4,
 \qquad z\in\Gamma.
\label{eq:v4v81-BS}
\end{equation}
\end{lemma}

\begin{proof}
The reference resolvent maps \(L^2\) to \(H^1\) uniformly on \(\Gamma\).
The four-dimensional Sobolev embedding maps \(H^1\) to \(L^4\), and
Holder's inequality gives
\[
 \|c(A_\rho)u\|_2
 \le C\|A_\rho\|_4\|u\|_4.
\]
Apply this to \(u=R_0(z)f\).
\end{proof}

\begin{definition}[Small rescaled spectral chart]
\label{def:v4v81-small-chart}
The ball belongs to the reference chart if
\begin{equation}
 q_A:=
 \sup_{z\in\Gamma}\|V_AR_0(z)\|
 \le q<1.
\label{eq:v4v81-small-chart}
\end{equation}
\end{definition}

\subsection{Low-band stability and two marks}
\label{sec:v4v81-low-band}

\begin{theorem}[Scale-uniform low-band projector]
\label{thm:v4v81-projector}
On a small rescaled spectral chart,
\begin{align}
 R_A(z)
 &:=(z-\widehat H_A)^{-1}
 =
 R_0(z)(1-V_AR_0(z))^{-1},
\label{eq:v4v81-resolvent}\\
 \|R_A(z)\|
 &\le
 \frac{\|R_0(z)\|}{1-q},
\label{eq:v4v81-resolvent-bound}\\
 P_A
 &:=
 \frac1{2\pi i}\int_\Gamma R_A(z)\,dz
\label{eq:v4v81-projector}
\end{align}
define a projector with the same rank as the reference band.  If
\(B\mapsto A_\rho(B)\) has first two marks in \(L^4\), then
\begin{align}
 DP_A[u]
 &=
 \frac1{2\pi i}\int_\Gamma
 R_A\,DV_A[u]\,R_A\,dz,
\label{eq:v4v81-DP}\\
 D^2P_A[u,v]
 &=
 \frac1{2\pi i}\int_\Gamma
 \bigl(
 R_A D^2V_A[u,v]R_A
 +R_ADV_A[u]R_ADV_A[v]R_A\\
 &\hspace{41mm}
 +R_ADV_A[v]R_ADV_A[u]R_A
 \bigr)\,dz,
\label{eq:v4v81-D2P}
\end{align}
and
\begin{equation}
 \|D^kP_A\|
 \le C_{\Gamma,q,k}
 \mathcal P_k\bigl(
 \|DA_\rho\|_4,\|D^2A_\rho\|_4
 \bigr),
 \qquad k=0,1,2,
\label{eq:v4v81-projector-marks}
\end{equation}
with the appropriate variables omitted at lower order.  All constants are
independent of \(\rho\).
\end{theorem}

\begin{proof}
Factor
\[
 z-\widehat H_A
 =
 (1-V_AR_0(z))(z-\widehat H_0).
\]
The Neumann series gives
\eqref{eq:v4v81-resolvent}--%
\eqref{eq:v4v81-resolvent-bound}.  The contour remains in the resolvent
set along \(tV_A\), \(0\le t\le1\); Riesz rank is integer-valued and
continuous, hence constant.  Differentiate the resolvent identity.  Each
marked multiplication operator is bounded from \(H^1\) to \(L^2\) by the
proof of Lemma~\ref{lem:v4v81-BS}; the adjacent resolvents supply the
\(H^1\) regularity.  Contour length and the Neumann bound give
\eqref{eq:v4v81-projector-marks}.  Rescaling placed every norm on the unit
ball, proving scale uniformity.
\end{proof}

\begin{corollary}[Uhlenbeck small energy implies the chart]
\label{cor:v4v81-Uhlenbeck-chart}
If the rescaled Uhlenbeck gauge obeys
\[
 \|A_\rho\|_4\le C_U E(B_\rho)^{1/2},
\]
then there is \(\varepsilon_{\rm sp}>0\), depending only on the reference
contour and bounded geometry, such that
\[
 E(B_\rho)\le\varepsilon_{\rm sp}
\]
implies \eqref{eq:v4v81-small-chart}.
\end{corollary}

\begin{proof}
Choose \(\varepsilon_{\rm sp}\) so that
\(C_\Gamma C_U\varepsilon_{\rm sp}^{1/2}\le q<1\), and apply
Lemma~\ref{lem:v4v81-BS}.
\end{proof}

\subsection{Off-diagonal locality}
\label{sec:v4v81-locality}

For a one-Lipschitz function \(\varphi\), let
\(W_\mu=e^{\mu\varphi}\).  Conjugating the first-order reference operator
adds the bounded commutator
\begin{equation}
 W_\mu\widehat H_0W_\mu^{-1}-\widehat H_0
 =
 -\mu\,c(d\varphi).
\label{eq:v4v81-weight-commutator}
\end{equation}
The multiplication \(V_A\) commutes with \(W_\mu\).

\begin{theorem}[Marked off-diagonal band locality]
\label{thm:v4v81-offdiagonal}
On a small chart with fixed reserve \(1-q>0\), there is \(\mu>0\) such that
for separated sets \(X,Y\subset B_1\),
\begin{equation}
 \|1_XD^kP_A1_Y\|
 \le C_k e^{-\mu d(X,Y)},
 \qquad k=0,1,2,
\label{eq:v4v81-offdiagonal}
\end{equation}
provided the marked gauge variations are localized and have the \(L^4\)
bounds of Theorem~\ref{thm:v4v81-projector}.  With a mark supported in
\(Z\), the bound contains the corresponding minimal broken path through
\(Z\).
\end{theorem}

\begin{proof}
Choose \(\mu\) so the norm of the commutator in
\eqref{eq:v4v81-weight-commutator}, after multiplication by the reference
resolvent, uses less than half the Neumann reserve.  The same factorization
as in Theorem~\ref{thm:v4v81-projector} then bounds
\(W_\mu R_A(z)W_\mu^{-1}\) uniformly.  Choosing \(\varphi\) to increase from
\(X\) to \(Y\) gives the exponential resolvent bound.  Integrate around
\(\Gamma\).  For \eqref{eq:v4v81-DP}--%
\eqref{eq:v4v81-D2P}, insert the localized marks between weighted
resolvents and use the triangle inequality for the successive distances.
\end{proof}

This is an operator-block locality estimate.  The V74 projected-seed Gram
condition is still needed to turn it into independent local Grassmann
coordinates.

\subsection{The fifth determinant through two marks}
\label{sec:v4v81-det5}

Let
\[
 K_A=R_0(0)V_A
\]
on a reference complement where \(R_0(0)\) is defined.  The imported
four-dimensional weak-Schatten estimate gives, for every \(p>4\),
\begin{equation}
 \|K_A\|_{\mathcal S_p}
 \le C_p\|A_\rho\|_4.
\label{eq:v4v81-Sp}
\end{equation}
Use \(p=5\).

\begin{lemma}[Entire marked regularized determinant]
\label{lem:v4v81-det5-marks}
The map
\[
 K\longmapsto\det{}_5(1+K)
\]
is entire on \(\mathcal S_5\).  On every ball
\(\|K\|_{\mathcal S_5}\le M\), its first two Frechet derivatives obey
\begin{align}
 |\det{}_5(1+K)|
 &\le e^{C_5M^5},
\label{eq:v4v81-det5-zero}\\
 |D\det{}_5(1+K)[L]|
 &\le C_M\|L\|_{\mathcal S_5},
\label{eq:v4v81-det5-one}\\
 |D^2\det{}_5(1+K)[L_1,L_2]|
 &\le C_M\|L_1\|_{\mathcal S_5}\|L_2\|_{\mathcal S_5}.
\label{eq:v4v81-det5-two}
\end{align}
These bounds remain valid when \(1+K\) has kernel.
\end{lemma}

\begin{proof}
The fifth canonical product defining \(\det{}_5\) converges locally
uniformly on \(\mathcal S_5\), hence defines an entire scalar function.
The standard product estimate gives
\eqref{eq:v4v81-det5-zero}.  Apply the Banach-space Cauchy formula on a
fixed-radius complex line disk inside the ball of radius \(M+1\), first in
one direction and then on a two-dimensional polydisk.  Local boundedness
from \eqref{eq:v4v81-det5-zero} yields
\eqref{eq:v4v81-det5-one}--%
\eqref{eq:v4v81-det5-two}.  Entirety, unlike a logarithmic derivative,
does not require \(1+K\) to be invertible.
\end{proof}

\begin{theorem}[Scale-uniform small-ball fermion card]
\label{thm:v4v81-small-ball-card}
On an Uhlenbeck ball satisfying
Corollary~\ref{cor:v4v81-Uhlenbeck-chart}, the following data are uniform
under \(\rho\downarrow0\):
\begin{enumerate}
\item the rank and first two marked operator-locality bounds of the selected
      low spectral band;
\item the hard complement resolvent on the chosen contour;
\item the fifth regularized determinant and its first two marks; and
\item the exact separation of the finite low determinant line from the
      complementary regularized determinant.
\end{enumerate}
No inverse of the low-band determinant is used.
\end{theorem}

\begin{proof}
Items 1--2 are Theorems
\ref{thm:v4v81-projector} and
\ref{thm:v4v81-offdiagonal}.  Equation
\eqref{eq:v4v81-Sp}, its marked linear versions, and
Lemma~\ref{lem:v4v81-det5-marks} give item 3.  The V73 Berezin--Schur
factorization and V74 determinant-line isomorphism give item 4.
\end{proof}

\subsection{Why bounded critical energy is insufficient}
\label{sec:v4v81-obstruction}

\begin{proposition}[A norm bound does not provide a spectral collar]
\label{prop:v4v81-no-collar}
A uniform bound
\(\|A_\rho\|_4\le M<\infty\), without the small-chart condition
\eqref{eq:v4v81-small-chart}, does not imply a common spectral contour or
bounded projector marks.
\end{proposition}

\begin{proof}
Already in a two-dimensional spectral subspace, consider
\[
 H_t=\begin{pmatrix}-t&0\\0&t\end{pmatrix},
 \qquad |t|\le1.
\]
The perturbation norm is uniformly bounded, but both eigenvalues cross zero
at \(t=0\).  A spectral projector selecting one sign has a rank change or a
discontinuous choice there, and no contour around that band has a uniform
collar.  Such a matrix is the finite low-band compression of a Fredholm
family with spectral flow.  Therefore an ambient critical norm bound, which
only bounds the compression coefficients, cannot exclude the same event.
\end{proof}

\begin{firewall}[Small balls do not trivialize a topological core]
\label{fw:v4v81-topological-core}
Subdividing a concentration region into Uhlenbeck-small balls controls local
connections and high determinants.  It does not place the complete
topological zero mode in one trivial bundle chart: transition functions
between the balls carry the charge.  The low bands must be glued by the V74
crossing lines and the established overlap/geometric index.
\end{firewall}

\subsection{Spectral-flow stopping rule}
\label{sec:v4v81-stopping}

Fix a deterministic list of reference contours
\(\Gamma_0,\Gamma_1,\ldots\) and increasing soft bands.  For a rescaled ball,
assign it to the first \(j\) for which
\[
 \sup_{z\in\Gamma_j}\|V_AR_{0,j}(z)\|\le q_j<1.
\]
If no such contour exists before the declared low window is exhausted,
assign the complete finite low band to the topological core.  Equality uses
one half-open convention.

\begin{theorem}[One owner for small charts and crossing cores]
\label{thm:v4v81-stopping}
The stopping rule assigns every finite-regulator ball once.  On a successful
chart, Theorem~\ref{thm:v4v81-small-ball-card} applies.  On a failed chart,
no hard logarithm is formed across the failed contour; the newly enclosed
band is transferred to the soft crossing line.  Nested transfers obey the
exact V74 cocycle.
\end{theorem}

\begin{proof}
A finite matrix has finitely many ordered spectral values, and the
deterministic first-success rule is unique.  Neumann invertibility proves
the successful case.  In the failed case, enlarging the soft band removes
the threatened eigenvalues before a hard inverse is used.  Exterior-power
associativity is Proposition~\ref{prop:v4v74-rank-change}.
\end{proof}

\begin{theorem}[V81 consequence for MUSC]
\label{thm:v4v81-MUSC}
The small-energy part of the V80 MUSC reduces to marked probability and
support estimates for Uhlenbeck balls; its fermionic spectral constants are
scale uniform by
Theorem~\ref{thm:v4v81-small-ball-card}.  The concentration-core part still
requires a weighted bound for the stopped crossing-line activities,
including inverse contour/seed loads.  A global \(L^4\) bound does not
supply that estimate.
\end{theorem}

\begin{proof}
On small balls all fermionic factors have uniform marked bounds, so the
remaining annular gain belongs to the bosonic activity and its occurrence
probability.  Proposition~\ref{prop:v4v81-no-collar} prevents extending
those constants to failed cores.  V75's exact mark-removal lemma requires
the actual inverse loads there.
\end{proof}

\begin{assessment}[Bottleneck after V81]
\label{ass:v4v81-status}
The rescaled low-band problem is closed on genuinely small Uhlenbeck charts.
Critical \(L^4\) control is exactly the correct scale, and the fifth
determinant remains twice markable even at zeros because no logarithm is
taken.  The open problem has narrowed to concentration cores: control their
finite crossing lines, projected-seed Grams, and size-annular activities
under the interacting law.

The next step should retain the topological transition functions across a
small-ball cover and estimate the finite Cech crossing product.  Treating
the balls independently would erase the very charge that V78 assigned.
\end{assessment}

\subsection{V81 result ledger and V82 target}
\label{sec:v4v81-conclusion}

V81 proves the scale-uniform Birman--Schwinger chart, low-band projector and
two marks, operator off-diagonal locality, and twice-marked fifth
determinant.  It proves that bounded critical norm without smallness cannot
supply a collar and gives one spectral-flow stopping owner for every
failed band.

V82 should construct the finite Cech low-band complex on an Uhlenbeck
small-ball cover of one concentration core.  It should retain the transition
functions, prove that the alternating overlap map recovers the established
integer index, and bound its finite determinant/crossing line through two
marks.  If the nerve size has no action-paid bound, the core must remain one
irreducible charged polymer rather than a product of small-ball activities.
\clearpage
\section{V82: energy-paid Cech low bands and quantitative index descent}
\label{sec:v4v82}

V81 controls every genuinely small Uhlenbeck chart but leaves a topological
concentration core covered by several such charts.  The frozen programme
already reconstructs the principal bundle with exact transition cocycles
and already proves the abstract spectral-section determinant-line cocycle.
The missing assertion is quantitative: the number of overlap factors must
be paid by the core energy, and the finite local low-band descent must be
compared to the global Fredholm low band before its Euler characteristic is
called the physical index.

This note proves an energy-paid nerve bound, constructs the exact
determinant line of the finite alternating descent complex, and gives a
norm-small chain-homotopy criterion which identifies its index with the
established overlap/geometric index.  It also bounds all unitary transition
determinants through two background marks.  Verification of the comparison
defect for the actual core remains the next analytic gate.

\subsection{Stopping cover and nerve payment}
\label{sec:v4v82-cover}

Let \(K\) be one connected concentration core and let
\(\{B_i\}_{i=1}^N\) be a Vitali stopping cover.  Assume:
\begin{enumerate}
\item the shrunken balls \(\widehat B_i\) are disjoint;
\item the fixed dilates \(B_i\) have overlap multiplicity at most \(\nu\);
\item every nonresidual stopping ball satisfies
\begin{equation}
 \varepsilon_-\le
 \int_{\widehat B_i}|F|^2
 \le\varepsilon_+,
\label{eq:v4v82-stopping-energy}
\end{equation}
where \(\varepsilon_+\le\varepsilon_{\rm sp}\) is the V81 small-chart
threshold; and
\item at most \(N_{\rm res}\) residual balls are admitted, with
\(N_{\rm res}\) fixed by the boundary geometry.
\end{enumerate}
Let \(\mathcal N\) be the nerve of the dilated cover.

\begin{lemma}[Energy pays the vertices]
\label{lem:v4v82-vertices}
If \(E_K=\int_K|F|^2\), then
\begin{equation}
 N\le \varepsilon_-^{-1}E_K+N_{\rm res}.
\label{eq:v4v82-N-energy}
\end{equation}
\end{lemma}

\begin{proof}
Sum the lower bound in
\eqref{eq:v4v82-stopping-energy} over the nonresidual disjoint inner balls.
Their energies sum to at most \(E_K\).  Add the residual balls.
\end{proof}

\begin{lemma}[Energy pays the finite nerve]
\label{lem:v4v82-nerve}
Suppose every point belongs to at most \(\nu\) dilated balls and the metric
has uniformly bounded packing on the stopping family.  Then, for
\(0\le p\le\nu-1\), the number \(N_p\) of \(p\)-simplices of
\(\mathcal N\) obeys
\begin{equation}
 N_p\le C_{\nu,p}N,
 \qquad
 |\mathcal N|:=\sum_{p=0}^4N_p
 \le C_\nu(\varepsilon_-^{-1}E_K+N_{\rm res}).
\label{eq:v4v82-nerve-energy}
\end{equation}
\end{lemma}

\begin{proof}
A ball can intersect only a uniformly bounded number \(D_\nu\) of other
stopping dilates by bounded overlap and packing.  Charge every simplex to
its least-index vertex.  At that vertex there are at most
\(\binom{D_\nu}{p}\) possible \(p\)-simplices.  Sum over vertices and use
Lemma~\ref{lem:v4v82-vertices}.  No simplex can contain more than \(\nu\) vertices because its balls have a
common point.  Hence the nerve stops at degree \(\nu-1\).
\end{proof}

\begin{firewall}[Residual balls cannot proliferate freely]
\label{fw:v4v82-residual}
Without a fixed bound on \(N_{\rm res}\), arbitrarily many balls of
arbitrarily small energy can cover a core and
\eqref{eq:v4v82-nerve-energy} fails.  Such balls belong to the ordinary
small-energy background cover; they are not separately charged as
concentration-core vertices.
\end{firewall}

\subsection{Exact local low-band transitions}
\label{sec:v4v82-transitions}

On \(B_i\), choose the V81 low-band projector \(P_i\) and an oriented
orthonormal frame \(R_i\).  On a nonempty overlap \(B_{ij}\), let \(G_{ij}\)
be the exact gauge/spin transition from the established reconstructed
bundle.  When the same contour is used,
\begin{equation}
 P_i=G_{ij}P_jG_{ij}^{-1}.
\label{eq:v4v82-projector-gluing}
\end{equation}
If the contours differ, let \(K_{ij}^{+}\) and \(K_{ij}^{-}\) be the finite
bands added on the two sides.  The stabilized transition is a unitary map
\begin{equation}
 U_{ij}:
 \operatorname{Ran}P_j\oplus K_{ij}^{-}
 \longrightarrow
 \operatorname{Ran}P_i\oplus K_{ij}^{+}.
\label{eq:v4v82-stabilized-transition}
\end{equation}
Its relative line is
\begin{equation}
 \mathcal R_{ij}
 =
 \det K_{ij}^{+}\otimes(\det K_{ij}^{-})^*.
\label{eq:v4v82-relative-line}
\end{equation}

\begin{proposition}[Exact overlap cocycle]
\label{prop:v4v82-overlap-cocycle}
The stabilized maps may be ordered so that on every triple overlap
\begin{equation}
 U_{ij}U_{jk}=U_{ik}
\label{eq:v4v82-U-cocycle}
\end{equation}
after the canonical cancellation of common finite bands.  Consequently
\begin{equation}
 \mathcal R_{ij}\otimes\mathcal R_{jk}
 \simeq\mathcal R_{ik}
\label{eq:v4v82-line-cocycle}
\end{equation}
exactly.
\end{proposition}

\begin{proof}
The physical transitions satisfy
\(G_{ij}G_{jk}=G_{ik}\) exactly.  Choose one common outer spectral window on
a triple overlap.  Every local cut is obtained from it by moving a finite
orthogonal band.  Direct-sum inclusion and deletion of those bands commute
with \(G_{ij}\), giving
\eqref{eq:v4v82-U-cocycle}.  Taking top exterior powers and cancelling the
middle band gives
\eqref{eq:v4v82-line-cocycle}; this is the finite determinant-functor
cocycle already used in V74, now attached to the spatial stopping nerve.
\end{proof}

\begin{firewall}[No logarithm of the full transition]
\label{fw:v4v82-no-log}
Only near-identity analytic corrections inside a chosen trivialization may
be logarithmed.  The exact \(G_{ij}\), and hence \(U_{ij}\), can carry the
topological class of the concentration core.  Replacing all transitions by
principal logarithms would force the core into the trivial sector.
\end{firewall}

\subsection{The finite alternating descent complex}
\label{sec:v4v82-complex}

\begin{definition}[Finite low-band descent datum]
\label{def:v4v82-descent-datum}
For every nonempty ordered intersection
\(B_I=B_{i_0}\cap\cdots\cap B_{i_p}\), a finite low-band descent datum
chooses one common stabilized space \(E_I^\pm\) for the right/left chiral
rows and maps
\[
 r_{I\leftarrow J}^\pm:E_J^\pm\longrightarrow E_I^\pm
 \qquad(J\subset I)
\]
such that
\[
 r_{I\leftarrow J}^\pm r_{J\leftarrow K}^\pm
 =r_{I\leftarrow K}^\pm
 \qquad(K\subset J\subset I).
\]
These maps are additional analytic data; they are not obtained by naively
restricting eigenvectors of different ball boundary problems.
\end{definition}

Fix such a datum.  Put
\begin{equation}
 C_\pm^p
 :=
 \bigoplus_{i_0<\cdots<i_p}E_{i_0\cdots i_p}^{\pm},
\label{eq:v4v82-cochains}
\end{equation}
and define
\begin{equation}
 (\delta_\pm c)_{i_0\cdots i_{p+1}}
 =
 \sum_{a=0}^{p+1}(-1)^a
 r_{\,i_0\cdots i_{p+1}\leftarrow
 i_0\cdots\widehat i_a\cdots i_{p+1}}^\pm
 c_{i_0\cdots\widehat i_a\cdots i_{p+1}}.
\label{eq:v4v82-delta}
\end{equation}

\begin{lemma}[Cech differential]
\label{lem:v4v82-delta}
One has
\begin{equation}
 \delta_\pm^2=0.
\label{eq:v4v82-delta2}
\end{equation}
\end{lemma}

\begin{proof}
Every twice-restricted term obtained by deleting indices \(a<b\) occurs
twice, in the two deletion orders.  Functoriality makes the maps equal and
the alternating signs are opposite.
\end{proof}

Let the local compressed chiral maps \(D_I:E_I^+\to E_I^-\) intertwine the
restrictions.  They give a morphism of Cech complexes.  Its mapping cone is
the finite total low-band complex
\begin{equation}
 \mathfrak C_{\rm low}
 :=
 \operatorname{Cone}\bigl(
 D_C:(C_+^\bullet,\delta_+)
 \longrightarrow(C_-^\bullet,\delta_-)
 \bigr).
\label{eq:v4v82-low-cone}
\end{equation}

\begin{definition}[Cech low-band index and line]
\label{def:v4v82-Cech-index}
Define
\begin{align}
 \operatorname{ind}_{\rm C}\mathfrak C_{\rm low}
 &:=
 \sum_k(-1)^k\dim H^k(\mathfrak C_{\rm low}),
\label{eq:v4v82-Cech-index}\\
 \operatorname{Det}\mathfrak C_{\rm low}
 &:=
 \bigotimes_k
 \left(\det H^k(\mathfrak C_{\rm low})\right)^{(-1)^k}.
\label{eq:v4v82-Cech-line}
\end{align}
\end{definition}

\begin{theorem}[Finite determinant-functor identity]
\label{thm:v4v82-det-functor}
There is a canonical oriented isomorphism
\begin{equation}
 \bigotimes_k
 \left(\det\mathfrak C_{\rm low}^k\right)^{(-1)^k}
 \simeq
 \operatorname{Det}\mathfrak C_{\rm low}.
\label{eq:v4v82-det-functor}
\end{equation}
It is compatible with refinement and with every relative line
\eqref{eq:v4v82-line-cocycle}.
\end{theorem}

\begin{proof}
Choose complements
\[
 Z^k=B^k\oplus H^k,\qquad
 C^k=Z^k\oplus A^k,
\]
where the differential maps \(A^k\) isomorphically onto \(B^{k+1}\).
In the alternating tensor product of the \(\det C^k\), the determinant of
\(A^k\) cancels the determinant of \(B^{k+1}\).  The uncancelled factors are
exactly the alternating determinants of \(H^k\), proving
\eqref{eq:v4v82-det-functor}.  Refinement or a cut change adds an acyclic
finite cone; its paired determinant factors cancel.  Associativity is the
relative-line cocycle.
\end{proof}

\subsection{Quantitative comparison with the global Fredholm band}
\label{sec:v4v82-comparison}

Let \(\mathfrak F_{\rm low}\) be the finite global low-band Fredholm complex
defined by the established overlap/geometric regulator on \(K\).  A
partition of unity and the exact bundle transitions give chain maps
\begin{equation}
 I:\mathfrak C_{\rm low}\to\mathfrak F_{\rm low},
 \qquad
 R:\mathfrak F_{\rm low}\to\mathfrak C_{\rm low}.
\label{eq:v4v82-IR}
\end{equation}

\begin{lemma}[Norm-small homotopy correction]
\label{lem:v4v82-homotopy}
Suppose there are degree \(-1\) maps \(h_{\rm C},h_{\rm F}\) and chain-map
defects \(E_{\rm C},E_{\rm F}\) such that
\begin{align}
 1-RI
 &=
 d_{\rm C}h_{\rm C}+h_{\rm C}d_{\rm C}+E_{\rm C},
\label{eq:v4v82-C-defect}\\
 1-IR
 &=
 d_{\rm F}h_{\rm F}+h_{\rm F}d_{\rm F}+E_{\rm F},
\label{eq:v4v82-F-defect}\\
 \|E_{\rm C}\|,\|E_{\rm F}\|&<1.
\label{eq:v4v82-defect-small}
\end{align}
Then \(\mathfrak C_{\rm low}\) and \(\mathfrak F_{\rm low}\) are
chain-homotopy equivalent.  Their cohomology, Euler index, and determinant
lines are canonically isomorphic.
\end{lemma}

\begin{proof}
The Neumann series make \(1-E_{\rm C}\) and \(1-E_{\rm F}\) invertible chain
maps.  Rewrite \eqref{eq:v4v82-C-defect} as
\[
 1-E_{\rm C}
 =
 RI+d_{\rm C}h_{\rm C}+h_{\rm C}d_{\rm C}.
\]
Multiplication by \((1-E_{\rm C})^{-1}\), which commutes with the
differential, shows that the identity is chain-homotopic to an invertible
chain-map multiple of \(RI\).  Thus \(RI\) induces an isomorphism on
cohomology.  The second equation gives the same conclusion for \(IR\), so
\(I\) and \(R\) induce inverse cohomology isomorphisms.  Finite-complex
determinant functoriality gives the line isomorphism.
\end{proof}

\begin{theorem}[Conditional recovery of the physical index]
\label{thm:v4v82-index-recovery}
If the local low spaces and partition maps satisfy
\eqref{eq:v4v82-C-defect}--%
\eqref{eq:v4v82-defect-small}, then
\begin{equation}
 \boxed{
 \operatorname{ind}_{\rm C}\mathfrak C_{\rm low}
 =
 \operatorname{ind}\mathfrak F_{\rm low}
 =
 Q_{\rm geom/ov}(K),}
\label{eq:v4v82-index-recovery}
\end{equation}
where the last integer is the already established geometric/overlap index
in the chosen Standard Model representation, with the relative boundary
trivialization of the core fixed.  The Cech determinant line is
the same finite crossing line used by the global regulator.
\end{theorem}

\begin{proof}
Lemma~\ref{lem:v4v82-homotopy} gives the first equality and the determinant
line isomorphism.  The second equality is the imported overlap/geometric
local-index theorem for the global regulated band.  No new topological
charge is defined.
\end{proof}

\begin{firewall}[An arbitrary Cech complex does not compute a Dirac index]
\label{fw:v4v82-no-automatic-index}
The nerve, transition cocycle, and local band ranks alone neither construct
Definition~\ref{def:v4v82-descent-datum} nor imply
\eqref{eq:v4v82-index-recovery}.  The quantitative two-sided comparison
\eqref{eq:v4v82-C-defect}--%
\eqref{eq:v4v82-defect-small} is essential.  Without it, the Cech
cohomology can describe only the chosen descent data, not the global
Fredholm kernel and cokernel.
\end{firewall}

\subsection{Two-marked determinant transition bound}
\label{sec:v4v82-marks}

Let \(U(B)\in U(r)\) be one stabilized transition in orthonormal frames,
with \(r\le r_*\).

\begin{lemma}[Unitary line marks]
\label{lem:v4v82-unitary-marks}
If
\[
 \|DU\|\le M_1,\qquad
 \|D^2U\|\le M_2,
\]
then
\begin{align}
 |\det U|&=1,
\label{eq:v4v82-det-unitary}\\
 |D\det U|&\le rM_1,
\label{eq:v4v82-det-U-one}\\
 |D^2\det U|
 &\le rM_2+2r^2M_1^2.
\label{eq:v4v82-det-U-two}
\end{align}
\end{lemma}

\begin{proof}
Jacobi's formula is valid because \(U^{-1}=U^*\):
\[
 D\det U=\det U\,\operatorname{tr}(U^{-1}DU).
\]
The trace is bounded by \(rM_1\).  Differentiate once more; use
\(D(U^{-1})=-U^{-1}(DU)U^{-1}\), bound the resulting trace by
\(rM_1^2+rM_2\), and bound the square of the first trace by \(r^2M_1^2\).
The displayed constant is a harmless enlargement.
\end{proof}

\begin{theorem}[Energy-paid crossing product]
\label{thm:v4v82-crossing-product}
Suppose every transition associated with a nerve simplex has rank at most
\(r_*\) and the uniform mark bounds of
Lemma~\ref{lem:v4v82-unitary-marks}.  Let \(\mathcal U_K\) be their
oriented determinant-functor product.  Then
\begin{align}
 |\mathcal U_K|&=1,
\label{eq:v4v82-product-zero}\\
 |D\mathcal U_K|
 &\le C_1|\mathcal N|,
\label{eq:v4v82-product-one}\\
 |D^2\mathcal U_K|
 &\le C_2(|\mathcal N|+|\mathcal N|^2).
\label{eq:v4v82-product-two}
\end{align}
Consequently the marked transition cost is a polynomial in
\(\varepsilon_-^{-1}E_K+N_{\rm res}\).
\end{theorem}

\begin{proof}
The modulus-one factors multiply to modulus one.  The first product
derivative chooses one marked factor, giving
\eqref{eq:v4v82-product-one}.  The second derivative either marks one
factor twice or two factors once, giving respectively
\(O(|\mathcal N|)\) and \(O(|\mathcal N|^2)\) terms.  Apply
Lemma~\ref{lem:v4v82-nerve}.
\end{proof}

\begin{corollary}[Absorption by an action reserve]
\label{cor:v4v82-action-absorption}
For every \(c>0\) and fixed \(m\),
\[
 (1+E_K)^m e^{-cE_K}
 \le C_{c,m}e^{-cE_K/2}.
\]
Thus the two-marked crossing product does not destroy a strictly positive
local gauge-action reserve.
\end{corollary}

\begin{proof}
The function
\((1+E)^me^{-cE/2}\) is bounded on \([0,\infty)\).  Multiply by
\(e^{-cE/2}\).
\end{proof}

\begin{firewall}[A transition phase is not a small activity]
\label{fw:v4v82-phase-not-small}
The unmarked crossing product has modulus one, not a small factor.  Its
marks are action-absorbable only after it is attached to the complete core
activity.  It cannot be expanded as an independent KP polymer or used as a
source of decay.
\end{firewall}

\subsection{V82 concentration-core interface}
\label{sec:v4v82-interface}

\begin{theorem}[Conditional finite-core card]
\label{thm:v4v82-core-card}
Assume for one concentration core:
\begin{enumerate}
\item the stopping-cover hypotheses
      \eqref{eq:v4v82-stopping-energy};
\item the V81 small-chart estimates on every cover member and overlap;
\item exact stabilized transition cocycles;
\item the quantitative comparison defect
      \eqref{eq:v4v82-defect-small};
\item uniform rank and two-mark transition bounds; and
\item a positive action reserve after all previously assigned bosonic and
      determinant factors.
\end{enumerate}
Then the core has one finite low determinant line whose index is the
established physical index, whose transition phase is exact, and whose
first two transition marks are paid by its own action.  The core may enter
the V78 charged-polymer ledger as one irreducible activity without a
nerve-size loss.
\end{theorem}

\begin{proof}
Lemmas~\ref{lem:v4v82-vertices}--%
\ref{lem:v4v82-nerve} pay the cover.  Proposition
\ref{prop:v4v82-overlap-cocycle} and
Theorem~\ref{thm:v4v82-det-functor} give one exact line.
Theorem~\ref{thm:v4v82-index-recovery} identifies its index.
Theorem~\ref{thm:v4v82-crossing-product} and
Corollary~\ref{cor:v4v82-action-absorption} pay its two marks.
Keeping the complete core as one activity prevents overlap factors from
being counted again.
\end{proof}

\begin{assessment}[Remaining gate after V82]
\label{ass:v4v82-status}
The size of the finite Cech construction is no longer an uncontrolled
combinatorial obstruction: a stopping cover with positive inner energy and
bounded overlap has an energy-paid nerve.  Exact transition lines and their
two marks are also harmless once attached to the complete action-weighted
core.

The nontrivial analytic gate is now
\(\|E_{\rm C}\|,\|E_{\rm F}\|<1\).  It measures whether local low bands and
partition-of-unity reconstruction capture the global Fredholm band.  V81's
individual chart gaps do not by themselves bound the commutators created by
the partition gradients.
\end{assessment}

\subsection{V82 result ledger and V83 target}
\label{sec:v4v82-conclusion}

V82 proves the energy-paid stopping nerve, exact stabilized low-band
cocycle, determinant-functor line, quantitative index-recovery criterion,
and two-marked transition-product bound.  It explicitly forbids identifying
an arbitrary nerve complex with the Dirac index.

V83 should compute the comparison defects.  The first target is an IMS
formula for the rescaled Dirac graph norm with
\([D,\chi_i]=c(d\chi_i)\), followed by a bound of
\(E_{\rm C},E_{\rm F}\) in terms of local contour gaps, overlap
multiplicity, and partition-gradient scale.  If the dimensionless
commutator remains order one on every scale, the cover cannot be inverted
perturbatively and the entire concentration core must be given one global
spectral section from the outset.
\clearpage
\section{V83: the Dirac IMS tax and the buffered descent criterion}
\label{sec:v4v83}

V82 reduces Cech index recovery to norm-small two-sided chain-homotopy
defects.  This note computes their leading term.  For a first-order Dirac
operator, localization creates the exact commutator
\([D,\chi_i]=c(d\chi_i)\).  On a ball of radius \(\rho\), both this
commutator and the local hard gap scale as \(\rho^{-1}\); their ratio is
dimensionless and does not improve as \(\rho\downarrow0\).

A buffered local parametrix can still work if the transition width is large
relative to the inverse rescaled gap.  We prove that criterion, propagate it
through two background marks, and give the decision rule: a core failing
the buffered inequality must use one global core spectral section rather
than a perturbative Cech inverse.

\subsection{Exact first-order IMS identity}
\label{sec:v4v83-IMS}

Let \(D\) be a Dirac-type operator on a fixed bundle chart and let
\(\chi_1,\ldots,\chi_N\) be real Lipschitz functions satisfying
\begin{equation}
 \sum_i\chi_i^2=1.
\label{eq:v4v83-quadratic-partition}
\end{equation}

\begin{theorem}[Dirac IMS formula]
\label{thm:v4v83-IMS}
For every \(u\) in the common graph domain,
\begin{equation}
 \sum_i\|D(\chi_i u)\|_2^2
 =
 \|Du\|_2^2+
 \sum_i\|c(d\chi_i)u\|_2^2.
\label{eq:v4v83-IMS}
\end{equation}
In particular the localization tax is positive and obeys
\begin{equation}
 0\le
 \sum_i\|c(d\chi_i)u\|_2^2
 \le
 \kappa_\chi^2\|u\|_2^2,
 \qquad
 \kappa_\chi^2:=
 \left\|\sum_i|d\chi_i|^2\right\|_\infty.
\label{eq:v4v83-IMS-tax}
\end{equation}
\end{theorem}

\begin{proof}
The Leibniz rule gives
\[
 D(\chi_i u)=\chi_iDu+c(d\chi_i)u.
\]
After squaring and summing, the first terms give \(\|Du\|^2\) by
\eqref{eq:v4v83-quadratic-partition}.  The mixed term is
\[
 2\operatorname{Re}
 \left\langle Du,
 c\!\left(\sum_i\chi_i\,d\chi_i\right)u
 \right\rangle=0,
\]
because
\(\sum_i\chi_i\,d\chi_i=\frac12d\sum_i\chi_i^2=0\).
The remaining squares give \eqref{eq:v4v83-IMS}; the pointwise Clifford
bound proves \eqref{eq:v4v83-IMS-tax}.
\end{proof}

\begin{corollary}[Scale invariance of the localization ratio]
\label{cor:v4v83-scale}
Suppose a physical ball has radius \(\rho\), the partition transition width
is \(w=\ell\rho\), and the physical hard singular-value gap is at least
\(\delta_{\rm phys}=\widehat\delta/\rho\).  Then
\begin{equation}
 \frac{\kappa_\chi}{\delta_{\rm phys}}
 \le
 \frac{C_\chi}{\ell\widehat\delta},
\label{eq:v4v83-scale-ratio}
\end{equation}
which is independent of \(\rho\).
\end{corollary}

\begin{proof}
A partition changing across width \(w\) has
\(\kappa_\chi\le C_\chi/w\).  Divide by
\(\widehat\delta/\rho\) and use \(w=\ell\rho\).
\end{proof}

Thus ultraviolet refinement alone cannot prove the V82 defect small.

\subsection{Local hard parametrices}
\label{sec:v4v83-parametrix}

Let \(B_i\Subset B_i^+\) be an inner patch and its buffered enlargement.
Choose \(\chi_i\) supported in \(B_i^+\), with
\(\sum_i\chi_i^2=1\) on the core.  Let \(D_i\) be the local operator on
\(B_i^+\), let \(P_i\) contain its complete low band, and let
\[
 C_i:(1-Q_i)E_{\rm L}\longrightarrow(1-P_i)E_{\rm R}
\]
be its hard inverse.  Assume
\begin{equation}
 \|C_i\|\le\delta_i^{-1},
 \qquad
 \delta_*:=\inf_i\delta_i>0.
\label{eq:v4v83-local-gap}
\end{equation}
After exact gauge transition to one reference trivialization, define
\begin{equation}
 \mathcal Q:=\sum_i\chi_i C_i\chi_i.
\label{eq:v4v83-parametrix}
\end{equation}

Let \(\Pi_{\rm R,loc}\) and \(\Pi_{\rm L,loc}\) denote the stabilized
right and left Cech low-band projections.  The difference between local and
global low-band ancestry is represented by
operators \(M_{\rm R},M_{\rm L}\), with
\begin{equation}
 \|M_{\rm R}\|,\|M_{\rm L}\|\le\epsilon_{\rm band}.
\label{eq:v4v83-band-mismatch}
\end{equation}

\begin{theorem}[Two-sided parametrix defect]
\label{thm:v4v83-parametrix-defect}
If the local inverses and transitions are exact on their buffered patches,
then on the stabilized hard space,
\begin{align}
 D\mathcal Q
 &=1-\Pi_{\rm L,loc}+\mathcal E_{\rm L},
\label{eq:v4v83-left-parametrix}\\
 \mathcal QD
 &=1-\Pi_{\rm R,loc}+\mathcal E_{\rm R},
\label{eq:v4v83-right-parametrix}
\end{align}
where
\begin{equation}
 \|\mathcal E_{\rm L}\|+\|\mathcal E_{\rm R}\|
 \le
 C_\nu\frac{\kappa_\chi}{\delta_*}
 +2\epsilon_{\rm band}.
\label{eq:v4v83-defect-bound}
\end{equation}
Here \(C_\nu\) depends only on the overlap multiplicity and fixed fibre
dimension.
\end{theorem}

\begin{proof}
Use \(DC_i=1-Q_i\) on the local hard range and commute \(D\) through the
left cutoff:
\[
 D\mathcal Q
 =
 \sum_i\chi_i(1-Q_i)\chi_i
 +
 \sum_i[D,\chi_i]C_i\chi_i.
\]
The first sum is \(1-\Pi_{\rm L,loc}\) up to the declared left band
mismatch.
The second sum contains
\(c(d\chi_i)C_i\chi_i\).  Cauchy--Schwarz in the patch index, bounded
overlap, \eqref{eq:v4v83-IMS-tax}, and
\eqref{eq:v4v83-local-gap} bound it by
\(C_\nu\kappa_\chi/\delta_*\).  The calculation for
\(\mathcal QD\) commutes through the right cutoff and is identical.  Add
the two band mismatches.
\end{proof}

\begin{firewall}[Local inverses require compatible boundary domains]
\label{fw:v4v83-boundary-domain}
The identity \(D_iC_i=1-Q_i\) must hold on the restriction of the same
global elliptic problem, or on a local problem connected to it by an exact
Poisson/Calderon correction.  Arbitrary Dirichlet inverses on overlapping
balls do not define the Cech descent maps of V82 and can add spurious
boundary modes.
\end{firewall}

\subsection{Buffered descent}
\label{sec:v4v83-buffer}

Define
\begin{equation}
 \eta_{\rm desc}
 :=
 C_\nu\frac{\kappa_\chi}{\delta_*}
 +2\epsilon_{\rm band}.
\label{eq:v4v83-eta}
\end{equation}

\begin{theorem}[Buffered descent criterion]
\label{thm:v4v83-buffered}
If
\begin{equation}
 \eta_{\rm desc}<1,
\label{eq:v4v83-buffered-condition}
\end{equation}
then Neumann correction of
\eqref{eq:v4v83-left-parametrix}--%
\eqref{eq:v4v83-right-parametrix} supplies the two-sided chain homotopies
required in Lemma~\ref{lem:v4v82-homotopy}.  Consequently the finite Cech
low-band index and determinant line equal the global regulated index and
line.
\end{theorem}

\begin{proof}
On the respective stabilized hard complements,
\(D\mathcal Q=1+\mathcal E_{\rm L}\) and
\(\mathcal QD=1+\mathcal E_{\rm R}\), with both errors of norm below one
after distributing the bound in
\eqref{eq:v4v83-buffered-condition}.  Their inverses are the convergent
Neumann series.  Correcting the inclusion, restriction, and homotopy maps by
these inverses gives
\eqref{eq:v4v82-C-defect}--%
\eqref{eq:v4v82-defect-small}.  Apply
Theorem~\ref{thm:v4v82-index-recovery}.
\end{proof}

\begin{corollary}[Dimensionless buffer requirement]
\label{cor:v4v83-buffer-size}
Under the scale hypotheses of
Corollary~\ref{cor:v4v83-scale}, a sufficient condition is
\begin{equation}
 \frac{C_\nu C_\chi}{\ell\widehat\delta}
 +2\epsilon_{\rm band}<1.
\label{eq:v4v83-buffer-size}
\end{equation}
Making \(\rho\) smaller without increasing the relative buffer \(\ell\)
does not improve this inequality.
\end{corollary}

\begin{proof}
Insert \eqref{eq:v4v83-scale-ratio} into
\eqref{eq:v4v83-eta}.
\end{proof}

\subsection{Two background marks}
\label{sec:v4v83-marks}

The stopping cover and partition are frozen on one background chart.  Thus
the declared background derivatives do not differentiate \(\chi_i\).  Let
the local operator and inverse marks obey
\[
 DC_i=-C_i(DD_i)C_i
\]
and the corresponding exact second derivative.

\begin{theorem}[Marked corrected parametrix]
\label{thm:v4v83-marked-parametrix}
Assume
\(\eta_{\rm desc}\le1-\tau\) for some \(\tau>0\), and assume the local
operators, hard inverses, transition maps, and band mismatches have uniform
first two marks.  Then the Neumann-corrected comparison maps and homotopies
in Theorem~\ref{thm:v4v83-buffered} have uniform first two marks.  Their
bounds are polynomials in
\[
 \tau^{-1},\quad
 \delta_*^{-1},\quad
 \kappa_\chi,\quad
 \nu,
\]
and the declared local mark constants.
\end{theorem}

\begin{proof}
Differentiate each local inverse:
\begin{align*}
 DC_i[u]&=-C_i(DD_i[u])C_i,\\
 D^2C_i[u,v]
 &=
 C_i(DD_i[u])C_i(DD_i[v])C_i\\
 &\quad+
 C_i(DD_i[v])C_i(DD_i[u])C_i
 -C_i(D^2D_i[u,v])C_i.
\end{align*}
Insert these formulas in
\eqref{eq:v4v83-parametrix}.  The fixed partition contributes no marks.
For the correction
\((1+\mathcal E)^{-1}\), use
\[
 D(1+\mathcal E)^{-1}
 =-(1+\mathcal E)^{-1}(D\mathcal E)(1+\mathcal E)^{-1}
\]
and its two-mark derivative.  The reserve gives
\(\|(1+\mathcal E)^{-1}\|\le\tau^{-1}\).  Product estimates yield the
stated polynomial.
\end{proof}

\begin{firewall}[Adaptive covers are not differentiable for free]
\label{fw:v4v83-adaptive-marks}
If stopping centers or radii change with the background, differentiating
the cover produces derivatives of characteristic decisions and cutoffs.
Theorem~\ref{thm:v4v83-marked-parametrix} applies on a frozen-cover chart.
Transitions between such charts must be assigned by the V75 first-failure
rule and paid in the V82 crossing line.
\end{firewall}

\subsection{The scale-matched obstruction}
\label{sec:v4v83-obstruction}

\begin{proposition}[Unbuffered refinement gives no perturbative reserve]
\label{prop:v4v83-unbuffered}
Suppose the cutoff partition changes across a distance comparable to the
ball radius, so \(\ell\asymp1\), and suppose only a fixed rescaled gap
\(\widehat\delta>0\) is known.  Then the IMS/parametrix estimate gives only
\begin{equation}
 \eta_{\rm desc}
 \lesssim
 \frac{C_\nu}{\widehat\delta}
 +2\epsilon_{\rm band},
\label{eq:v4v83-order-one}
\end{equation}
an order-one quantity.  No refinement limit forces it below one.
\end{proposition}

\begin{proof}
This is Corollary~\ref{cor:v4v83-scale} with \(\ell\asymp1\).
All powers of \(\rho\) cancel exactly.
\end{proof}

This is a structural obstruction to the proposed local inversion, not a
proof that every core fails.  A large rescaled gap or a wide buffer can
still satisfy \eqref{eq:v4v83-buffer-size}.

\begin{theorem}[Local-or-global core decision]
\label{thm:v4v83-decision}
For each concentration core, use exactly one of the following routes:
\begin{enumerate}
\item prove the buffered inequality
      \eqref{eq:v4v83-buffered-condition} on a frozen-cover chart and use
      the V82 Cech determinant line; or
\item construct one global spectral section for the complete core operator,
      retain its finite low band and determinant line, and use the small
      balls only for high-energy determinant estimates.
\end{enumerate}
If the buffered inequality is unavailable, route 2 is mandatory; multiplying
independent local determinant lines is not justified.
\end{theorem}

\begin{proof}
Route 1 is Theorem~\ref{thm:v4v83-buffered}.  If its hypothesis fails, the
local parametrix does not establish the chain equivalence required by the
V82 index theorem.  A global spectral section directly defines the global
low band, while V81 remains applicable to high-energy pieces on each small
ball.  The alternatives assign the low determinant occurrence once.
\end{proof}

\begin{assessment}[Bottleneck after V83]
\label{ass:v4v83-status}
The Cech comparison defect is now explicit.  Its leading ratio is
partition-gradient over hard gap, plus the stabilized-band mismatch.
Because that ratio is scale invariant, ultraviolet refinement does not
close it.  A quantitative buffer may close selected cores; otherwise the
programme must go upstream to a global core spectral section.

This does not obstruct the V81 fifth-determinant estimates, which remain
local and scale uniform.  It obstructs only the claim that independently
chosen local low bands recover the global topological zero-mode line.
\end{assessment}

\subsection{V83 result ledger and V84 target}
\label{sec:v4v83-conclusion}

V83 proves the exact Dirac IMS identity, the two-sided local-parametrix
defect, the buffered Cech descent criterion, its two-mark stability, and the
scale-matched obstruction for unbuffered covers.  It gives one
local-or-global owner for every concentration-core low band.

V84 should implement route 2.  It should choose a global reference core
operator, use one finite spectral window containing every crossing along
the core interpolation, and prove a rank/action bound for that window by a
spectral counting inequality.  The target is an action-paid finite
determinant line with two marks, while local Uhlenbeck balls continue to
control only the complementary fifth determinant.
\clearpage
\section{V84: global core crossing reduction from weak-\(\mathcal S_4\) truncation}
\label{sec:v4v84}

V83 showed that a coverwise spectral construction closes only when the
Dirac IMS error is smaller than the local spectral reserve, and that
shrinking four-dimensional balls does not improve this ratio.  We therefore
work globally on each stopping core.  The large singular subspaces do not,
by themselves, make a complementary compression invertible.  The argument
below instead inverts a norm-small tail and uses an exact Sylvester
finite-rank reduction.

\subsection{Reference hard problem}

Let \(H_0\) be a closed Fredholm Dirac-type operator on a reconstructed
stopping core, with finite reference soft section \(P_0\).  Suppose its hard
restriction \(H_{0,\mathrm h}\colon\mathcal D_{\mathrm h}\to
\mathcal H_{\mathrm h}\) is bijective with bounded graph-norm inverse
\(C_0\).  For the rescaled gauge--Higgs perturbation \(V\), put
\[
 K:=VC_0,\qquad
 H_t:=H_{0,\mathrm h}+tV=(1+tK)H_{0,\mathrm h}.
\]
The required critical estimate is
\begin{equation}
 K\in\mathcal S_{4,\infty},\qquad
 \|K\|_{4,\infty}:=\sup_{n\ge1}n^{1/4}s_n(K)
 \le C_{\rm rel}\mathcal A_4 .
\label{eq:v4v84-weakS4}
\end{equation}

\begin{lemma}[Weak-Schatten count]
For every \(\alpha>0\),
\[
 N_\alpha(K):=\#\{n:s_n(K)>\alpha\}
 \le\bigl(\|K\|_{4,\infty}/\alpha\bigr)^4.
\]
\end{lemma}
\begin{proof}
If \(N_\alpha(K)=N>0\), then
\(\alpha<s_N(K)\le N^{-1/4}\|K\|_{4,\infty}\).  Rearrangement proves the
claim; \(N=0\) is immediate.
\end{proof}

\subsection{Finite large part and invertible tail}

Fix \(0<q<1\).  Choose smooth \(f_q\colon[0,\infty)\to[0,1]\) with
\(f_q=0\) on \([0,q/2]\), \(f_q>0\) on \((q/2,\infty)\), and \(f_q=1\)
on \([q,\infty)\).  For \(K=U|K|\), define
\[
 K_>:=Uf_q(|K|)|K|,\qquad K_<:=K-K_>.
\]

\begin{lemma}[Truncation bounds]
\label{lem:v4v84-trunc}
\[
 R:=\operatorname{rank}K_>
 \le\left(\frac{2\|K\|_{4,\infty}}q\right)^4,
 \qquad \|K_<\|\le q.
\]
\end{lemma}
\begin{proof}
The range of \(K_>\) is generated by singular vectors with singular value
greater than \(q/2\), which gives the rank bound.  On a singular vector of
value \(s\), the tail multiplier is \(s(1-f_q(s))\).  It vanishes for
\(s\ge q\) and is at most \(s<q\) below \(q\).
\end{proof}

Set \(T_t:=1+tK_<\).  The Neumann series gives
\begin{equation}
 T_t^{-1}=\sum_{j\ge0}(-tK_<)^j,\qquad
 \|T_t^{-1}\|\le(1-q)^{-1}.
\label{eq:v4v84-tailinv}
\end{equation}
Thus the tail has no crossing.

\subsection{Exact Sylvester block}

Let \(E=\operatorname{Ran}\mathbf1_{(q/2,\infty)}(|K|)\), let
\(B=P_E\colon\mathcal H_{\mathrm h}\to E\), and let
\(A=K_>|_E\colon E\to\mathcal H_{\mathrm h}\).  Then
\(\dim E=R\), \(A\) is injective, and \(K_>=AB\).  Put
\[
 A_t:=T_t^{-1}A,\qquad M_t:=1_E+tBA_t.
\]

\begin{theorem}[Exact finite crossing reduction]
\label{thm:v4v84-sylvester}
For \(0\le t\le1\),
\begin{equation}
 1+tK=T_t(1+tA_tB).
\label{eq:v4v84-factor}
\end{equation}
The operators \(H_t\), \(1+tA_tB\), and \(M_t\) are invertible
simultaneously.  For \(t>0\),
\[
 A_t:\ker M_t\overset{\cong}{\longrightarrow}\ker(1+tA_tB),
 \qquad
 B^*:\ker M_t^*\overset{\cong}{\longrightarrow}\ker(1+tA_tB)^*.
\]
Consequently every hard kernel and cokernel direction is represented by
the \(R\times R\) matrix \(M_t\).
\end{theorem}
\begin{proof}
Equation \eqref{eq:v4v84-factor} follows from
\(T_t(1+tT_t^{-1}K_>)=1+tK\).  If \(M_t\) is invertible, direct
multiplication proves the Woodbury identity
\[
 (1+tA_tB)^{-1}=1-tA_tM_t^{-1}B.
\]
Conversely, \(M_tc=0\) implies
\((1+tA_tB)A_tc=A_tM_tc=0\); \(A_t\) is injective.  If
\((1+tA_tB)x=0\), then \(x=-tA_tBx\), hence \(x=A_tc\), and injectivity
of \(A_t\) gives \(M_tc=0\).  This proves the first kernel isomorphism.
The same argument applied to
\((1+tA_tB)^*=1+tB^*A_t^*\) proves the second.  The factors
\(T_t\) and \(H_{0,\mathrm h}\) are invertible, completing the proof.
\end{proof}

\begin{corollary}[Determinant carrier]
\label{cor:v4v84-det}
\[
 \det\nolimits_F(1+tA_tB)=\det\nolimits_E M_t.
\]
With
\(\operatorname{Det}D=\det\ker D\otimes(\det\operatorname{coker}D)^*\),
the kernel isomorphisms induce the relative determinant-line
identification
\[
 \operatorname{Det}H_t\cong
 \operatorname{Det}H_{0,\mathrm h}\otimes\operatorname{Det}M_t,
\]
tensored with the already separated soft line of \(P_0\).
\end{corollary}
\begin{proof}
The nonzero eigenvalues of \(A_tB\) and \(BA_t\), with algebraic
multiplicity, coincide.  This is the Sylvester determinant identity.
Taking top exterior powers of the kernel and cokernel isomorphisms proves
the line statement; invertible factors have canonical trivializations.
\end{proof}

\subsection{Families, collars, and marks}

For a parameter family \(K(z)\), require a chart and \(\eta>0\) such that
\begin{equation}
 \operatorname{spec}|K(z)|\cap(q/2-\eta,q/2+\eta)=\varnothing.
\label{eq:v4v84-collar}
\end{equation}
Then \(E(z)\) is a constant-rank finite vector bundle.  Functional
calculus makes \(K_>(z),K_<(z),P_{E(z)}\), and
\(M(z,t)\in\operatorname{End}E(z)\) smooth whenever \(K(z)\) is smooth in
the marked operator topology.  In a collared chart,
\[
 \partial_aT_t^{-1}
 =-T_t^{-1}t(\partial_aK_<)T_t^{-1},
\]
and
\[
 \partial_aA_t
 =T_t^{-1}\partial_aA
  -T_t^{-1}t(\partial_aK_<)T_t^{-1}A.
\]
A second derivative follows by one more product rule; its terms contain at
most three inverse factors and cost at most \((1-q)^{-3}\).  Hence the
twice-marked crossing calculus is finite dimensional after the controlled
tail.  A frame change conjugates \(M\), so its determinant and determinant
line agree on overlaps.

\subsection{Action payment}

From Lemma~\ref{lem:v4v84-trunc} and
\eqref{eq:v4v84-weakS4},
\[
 R\le\left(\frac{2C_{\rm rel}}q\right)^4\mathcal A_4^4.
\]
On a core where an intrinsic Uhlenbeck--Higgs comparison proves
\(\mathcal A_{4,K}\le C_UE_K^{1/2}\),
\[
 R_K\le C(q,C_{\rm rel},C_U)E_K^2.
\]
For disjoint cores with \(E_K\le E_{\max}\),
\[
 \sum_KR_K\le
 C(q,C_{\rm rel},C_U)E_{\max}\sum_KE_K.
\]
Thus bounded total action pays the total crossing dimension.  Concentration
may enlarge the finite matrix, but cannot make its dimension infinite.

Since \(\mathcal S_{4,\infty}\subset\mathcal S_p\) for every \(p>4\),
the fifth regularized determinant from V81 remains defined.  It packages
the infinite-dimensional marked expansion; \(\det_E M_t\) locates every
zero.  The regularized multiplicative anomaly for
\eqref{eq:v4v84-factor} is an exponential of trace polynomials and is
nowhere zero, so it does not change the crossing divisor.

\subsection{Closed statement and remaining premise}

\begin{theorem}[Conditional global core spectral section]
\label{thm:v4v84-global}
Assume that on the exact reconstructed bundle there is a global reference
\(H_0\) with finite soft section and bounded hard inverse, that
\eqref{eq:v4v84-weakS4} holds in two marked variables, and that the collar
\eqref{eq:v4v84-collar} is available.  Then \(H_t\) has a finite-rank
global crossing bundle \(E\), with the rank bound above, and all hard
zeros, kernels, cokernels, determinant-line data, and twice-marked crossing
derivatives are represented by \(M(z,t)\).
\end{theorem}
\begin{proof}
Lemma~\ref{lem:v4v84-trunc} and \eqref{eq:v4v84-tailinv} give the finite
bundle and invertible tail.  Theorem~\ref{thm:v4v84-sylvester} and
Corollary~\ref{cor:v4v84-det} identify the crossing data.  The marked
formulas and rank estimate finish the proof.
\end{proof}

The theorem is conditional for a precise reason.  Local Uhlenbeck gauges
on a topologically nontrivial core do not automatically give one global
reference comparison, and an arbitrary Cech complex cannot be declared to
compute the Dirac index.  V85 will first perform the required YM/NS audit,
then construct the reference intrinsically on the reconstructed bundle and
test the weak-\(\mathcal S_4\) premise.

\clearpage
\section{V85: companion audit and the intrinsic reference comparison}
\label{sec:v4v85}

This is the compulsory five-version audit.  A stable snapshot of the current
Yang--Mills note V6 was read at SHA--256
\[
 \texttt{1c5292e474fc1bd1ca80a4fa8c33825a935a2c9160c05fe5c35779af5f9e311d},
\]
and the Navier--Stokes note V31 at
\[
 \texttt{f1121b62238ad5491bb7cf03635ea9934942edf573ed8faaa9ee79e1d38a6696}.
\]
Neither companion file was modified.

The YM upstream-localization theorem says that a surviving strong stock must
remain with an extracted slot, one complete prefix carrier, or a coupled
cross-prefix packet; it cannot be assigned independently to two carriers.
The NS heat-formation ladder says that a flat mark exposes, rather than
removes, the critical first-moment cost, and that a persistent interval must
not be charged again as a recent-entry interval.  The SM translation is:
\begin{enumerate}
 \item reference topological zero modes, the finite soft Schur block, and
 the hard relative perturbation require distinct owners;
 \item transition data crossing two trivializations must remain one bundle
 object until the global connection has been reconstructed;
 \item the singular cutoff of V84 does not remove the critical load:
 its rank still costs \(q^{-4}\mathcal A_4^4\); and
 \item a mode already carried by the reference soft block may not be charged
 again to the finite Sylvester crossing matrix.
\end{enumerate}

\subsection{The exact global bundle setting}
\label{sec:v4v85-bundle}

Let \(X\) be a compact four-dimensional Riemannian spin core, possibly with
smooth boundary, and let \(P\to X\) be the exact reconstructed principal
Standard Model bundle.  Its associated finite-rank fermion bundle is denoted
by \(\mathcal E\).  All statements below are invariant under bundle
isomorphism; no global trivialization of \(P\) is chosen.

Fix a smooth reference connection \(A_\circ\) on \(P\) and a smooth
reference Higgs section \(\Phi_\circ\).  Such a connection exists: choose
local connection forms and average them with a partition of unity using the
affine transformation law.  The resulting local forms obey the same
transition law and hence define a global connection.

Let \(\mathbb H_\circ\) be the doubled gauge--Higgs Dirac operator on
\(\mathcal E^+\oplus\mathcal E^-\).  If \(\partial X\ne\varnothing\), fix
once and for all an elliptic self-adjoint boundary projector for
\(\mathbb H_\circ\).  Every lower-order perturbation considered below has
the same principal symbol and Green boundary form, so the same domain is
Fredholm; a reference APS projector with a fixed Lagrangian on its kernel is
one admissible choice.

Since \(\mathbb H_\circ\) has compact resolvent, for every \(\delta_\circ>0\)
outside its absolute spectrum,
\[
 P_\circ:=\mathbf1_{[0,\delta_\circ]}(|\mathbb H_\circ|)
\]
has finite rank.  On \(\mathcal H_{\rm h}=(1-P_\circ)L^2(\mathcal E)\),
\[
 C_\circ:=
 \mathbb H_\circ^{-1}(1-P_\circ),\qquad
 \|C_\circ\|\le\delta_\circ^{-1}.
\]
The finite block \(P_\circ\) owns the reference kernel and near-zero modes.
Off-diagonal couplings to it are retained in the exact soft Schur complement
proved in V73--V74.  The V84 relative operator is the hard compression only.

\subsection{Connection differences are intrinsic tensors}

\begin{lemma}[Global zero-order difference]
\label{lem:v4v85-difference}
For any second connection \(A\) on \(P\) and Higgs section \(\Phi\),
\[
 a:=A-A_\circ\in\Omega^1(X;\operatorname{ad}P),\qquad
 \varphi:=\Phi-\Phi_\circ
\]
are global tensor fields.  On the common operator domain,
\begin{equation}
 \mathbb H_{A,\Phi}-\mathbb H_\circ
 =W(a,\varphi):=c(a)+Y(\varphi),
\label{eq:v4v85-W}
\end{equation}
where \(c\) is Clifford representation combined with the SM
representation, and \(Y\) is the finite-dimensional Yukawa map.
Pointwise,
\begin{equation}
 |W(a,\varphi)|\le
 C_{\rm rep}\bigl(|a|+|\varphi|\bigr).
\label{eq:v4v85-pointwise}
\end{equation}
\end{lemma}
\begin{proof}
In a local frame, two connection forms transform as
\[
 A_j=g_{ij}^{-1}A_ig_{ij}+g_{ij}^{-1}dg_{ij},\qquad
 (A_\circ)_j=g_{ij}^{-1}(A_\circ)_ig_{ij}+g_{ij}^{-1}dg_{ij}.
\]
Subtraction cancels the inhomogeneous term, so
\(a_j=g_{ij}^{-1}a_ig_{ij}\), the transition law of
\(\operatorname{ad}P\)-valued one-forms.  The covariant Dirac formula then
gives the Clifford multiplier \(c(a)\).  The Higgs difference is already a
section of the associated Higgs bundle, and the Yukawa term is fibrewise
linear.  Finite-dimensional operator norms give
\eqref{eq:v4v85-pointwise}.
\end{proof}

\subsection{Intrinsic weak-\(\mathcal S_4\) comparison}

We record the precise analytic input rather than hiding it inside local
gauges.

\begin{lemma}[Four-dimensional Cwikel estimate]
\label{lem:v4v85-cwikel}
Let \(D\) be a first-order elliptic Dirac-type operator with compact
resolvent on a compact four-dimensional manifold of bounded geometry, with
a fixed elliptic boundary condition when needed.  For every endomorphism
field \(w\in L^4(X;\operatorname{End}\mathcal E)\),
\begin{equation}
 \left\|M_w(1+D^2)^{-1/2}\right\|_{4,\infty}
 \le C_{X,D}\|w\|_{L^4}.
\label{eq:v4v85-cwikel}
\end{equation}
The constant is uniform on a bounded-geometry family with fixed ellipticity
and boundary constants.
\end{lemma}
\begin{proof}
Choose finitely many interior and boundary charts and a subordinate smooth
partition of unity.  A first-order elliptic parametrix writes each localized
piece of \((1+D^2)^{-1/2}\) as a classical order-\(-1\) operator plus an
order-\(-2\) remainder.  In a chart, the Birman--Solomyak weak Cwikel
inequality in dimension four gives
\[
 \|M_f\,b(x,-i\nabla)\|_{4,\infty}
 \le C\|f\|_4
\]
for every order-\(-1\) symbol \(b\) with bounded normalized symbol
seminorms.  Boundary charts have the same estimate after the standard
half-space extension; the singular Green remainder has order at most
\(-2\).  Operators of order \(-2\) lie in
\(\mathcal S_{2,\infty}\subset\mathcal S_{4,\infty}\), and multiplication
by the smooth cutoffs is bounded.  The weak Schatten quasi-triangle
inequality over the finite atlas proves \eqref{eq:v4v85-cwikel}.
Uniform symbol, extension, and atlas bounds give the final assertion.
\end{proof}

\begin{theorem}[Intrinsic global reference comparison]
\label{thm:v4v85-intrinsic}
Let \(Q_{\rm h}=1-P_\circ\), and define
\[
 K_{\rm h}:=
 Q_{\rm h}W(a,\varphi)Q_{\rm h}C_\circ.
\]
Then
\begin{equation}
 K_{\rm h}\in\mathcal S_{4,\infty},\qquad
 \|K_{\rm h}\|_{4,\infty}
 \le C_{\rm int}
 \sqrt{1+\delta_\circ^{-2}}\,
 \bigl(\|a\|_4+\|\varphi\|_4\bigr).
\label{eq:v4v85-intrinsic}
\end{equation}
The construction and estimate are invariant under simultaneous gauge
transformation of \(A,A_\circ,\Phi,\Phi_\circ\).
\end{theorem}
\begin{proof}
Let \(R_\circ=(1+\mathbb H_\circ^2)^{-1/2}\).  Spectral calculus gives
\[
 C_\circ=R_\circ G_\circ,\qquad
 G_\circ=
 \frac{\sqrt{1+\mathbb H_\circ^2}}{\mathbb H_\circ}
 \mathbf1_{(\delta_\circ,\infty)}(|\mathbb H_\circ|),
\]
with
\(\|G_\circ\|\le\sqrt{1+\delta_\circ^{-2}}\).
Lemma~\ref{lem:v4v85-cwikel}, the pointwise estimate
\eqref{eq:v4v85-pointwise}, and the two-sided ideal property yield
\[
 \|Q_{\rm h}WQ_{\rm h}R_\circ G_\circ\|_{4,\infty}
 \le C_{\rm int}\sqrt{1+\delta_\circ^{-2}}
      (\|a\|_4+\|\varphi\|_4).
\]
Gauge transformations act unitarily on every factor and preserve the
\(L^4\) norms of the tensor differences, proving invariance.
\end{proof}

\begin{corollary}[Two marked variables]
\label{cor:v4v85-marks}
If \(a(z),\varphi(z)\) have two parameter derivatives in \(L^4\), while
the sector reference and \(P_\circ\) are fixed, then for
\(|\alpha|\le2\),
\[
 \|\partial_z^\alpha K_{\rm h}(z)\|_{4,\infty}
 \le C_{\rm int}\sqrt{1+\delta_\circ^{-2}}
 \left(\|\partial_z^\alpha a(z)\|_4+
       \|\partial_z^\alpha\varphi(z)\|_4\right).
\]
Thus the marked weak-\(\mathcal S_4\) premise of V84 follows from an
intrinsic marked \(L^4\) comparison.
\end{corollary}
\begin{proof}
The reference resolvent and hard projection are fixed, so differentiation
falls only on the linear multiplier \(W(a,\varphi)\).  Apply
Theorem~\ref{thm:v4v85-intrinsic} to each derivative.
\end{proof}

\subsection{Topological index and analytic crossings have different owners}

Let \(\mathcal Q_t\) denote the chiral block before doubling.  Its principal
symbol and fixed elliptic boundary condition do not change with \(t\).
Therefore \(\mathcal Q_t\) is a norm-continuous Fredholm family and
\begin{equation}
 \operatorname{ind}\mathcal Q_t=\operatorname{ind}\mathcal Q_\circ.
\label{eq:v4v85-index}
\end{equation}
Indeed, the Fredholm set is open and its integer-valued index is locally
constant.  The reference kernel and cokernel, including the topological
imbalance, are owned by \(P_\circ\) and the finite soft Schur block.
The matrix \(M_t\) of V84 owns only additional hard crossing pairs.
Equation \eqref{eq:v4v85-index} forbids charging a topological zero mode
again to \(M_t\).  This is the exact SM form of the YM owner rule and the
NS no-double-charge rule.

\subsection{Why curvature energy alone cannot pay a fixed reference}

The intrinsic theorem does not yet prove an action-paid estimate for
\(\|A-A_\circ\|_4\).  The missing gauge choice is essential.

\begin{proposition}[Pure-gauge obstruction]
\label{prop:v4v85-puregauge}
There is no constant \(C\) such that, in every fixed trivialization,
\[
 \|A-A_\circ\|_{L^4}\le C\|F_A\|_{L^2}
\]
for all connections on a trivial bundle.
\end{proposition}
\begin{proof}
On a flat four-torus, take \(A_\circ=0\) and a nonzero integral Lie-algebra
element \(T\).  For integers \(N\), let
\(g_N(x)=\exp(Nx_1T)\), with the period normalized so that \(g_N\) is
single-valued, and put \(A_N=g_N^{-1}dg_N\).  Then \(F_{A_N}=0\), while
\[
 \|A_N\|_4=N\|T\,dx_1\|_4\longrightarrow\infty.
\]
All \(A_N\) are gauge equivalent to zero.  Hence the estimate fails before
a gauge representative is selected.
\end{proof}

This also explains why returning to a patchwise potential estimate would
be circular.  Gauge-transition derivatives are the cross-prefix packet in
the YM analogy; splitting them among charts discards their common bundle
owner.  Averaging over spectral cutoffs cannot repair the loss, just as the
NS flat mark retains its first moment.

\subsection{The exact next certificate}

\begin{definition}[Action-paid reference-gauge certificate]
\label{def:v4v85-argc}
For a bounded-action family in a fixed reconstructed topological sector,
an ARGC consists of finitely many reference connections
\(A_\circ^{(j)}\), Higgs references \(\Phi_\circ^{(j)}\), and gauge-invariant
charts \(\mathcal U_j\) such that every configuration lies in some chart
and admits a bundle automorphism \(u\) with
\begin{equation}
 \|u(A)-A_\circ^{(j)}\|_4+
 \|u(\Phi)-\Phi_\circ^{(j)}\|_4
 \le C_{\rm ARGC}\,\mathfrak E^{1/2},
\label{eq:v4v85-argc}
\end{equation}
with the same estimate for the two marked differences.  On chart overlaps,
the complete relative determinant line is transported by the bundle
isomorphism before any scalar determinant is taken.
\end{definition}

\begin{theorem}[Closure conditional only on ARGC]
\label{thm:v4v85-closure}
On every ARGC chart, the full reference-soft/hard-Sylvester construction is
global on the exact bundle.  Its hard crossing rank obeys
\[
 R_j\le
 C(q,\delta_\circ,C_{\rm int},C_{\rm ARGC})\,\mathfrak E^2.
\]
For disjoint bounded-energy cores, its total rank is action-paid as in V84.
All topological index data remains in the reference soft owner.
\end{theorem}
\begin{proof}
Equation \eqref{eq:v4v85-argc} inserted into
\eqref{eq:v4v85-intrinsic} gives the V84 weak-\(\mathcal S_4\) bound with
\(\mathcal A_4\lesssim\mathfrak E^{1/2}\).  The V84 counting theorem gives
the rank estimate.  V73--V74 sew the finite reference soft block to the hard
determinant line, while \eqref{eq:v4v85-index} fixes the topological owner.
Gauge conjugacy provides the overlap transport.
\end{proof}

V85 therefore removes the supposed need for a global trivialization and
proves the intrinsic weak-\(\mathcal S_4\) comparison.  The remaining
bottleneck is narrower: construct the finite quantitative ARGC atlas at
critical \(L^2\) curvature, including concentration cores and two marked
variables.  V86 will go upstream to the Uhlenbeck transition cocycle and
separate the small-curvature atlas, which can be proved directly, from the
large topological core that needs bubble-tree ownership.

\clearpage
\section{V86: an action-paid cocycle reference connection}
\label{sec:v4v86}

V85 reduced the global comparison problem to an action-paid reference-gauge
certificate.  We now construct a reference connection directly from the
transition cocycle of a small-energy Uhlenbeck cover.  This construction is
global on the exact reconstructed bundle and does not require the bundle to
be trivial.  It proves the unmarked \(L^4\) payment.  Its dependence on the
chosen Uhlenbeck gauges remains a separate marked-family problem.

\subsection{Good cover and conventions}

Let \(\{U_i\}_{i=1}^N\) be a finite cover of a stopping core \(X\), with
subordinate partition of unity \(\{\rho_i\}\), overlap multiplicity at most
\(\nu\), and uniformly controlled scaled geometry.  Assume
\begin{equation}
 e_i:=\int_{U_i}|F_A|^2\le\varepsilon_U
\label{eq:v4v86-small}
\end{equation}
on every enlarged chart required by the Uhlenbeck theorem.  In a local
Uhlenbeck frame, write the physical connection as \(d+A_i\), with
\begin{equation}
 d^*A_i=0,\qquad
 \|A_i\|_{L^4(U_i)}+\|\nabla A_i\|_{L^2(U_i)}
 \le C_U e_i^{1/2}.
\label{eq:v4v86-uhl}
\end{equation}
On \(U_i\cap U_j\), choose the convention
\begin{equation}
 A_j=g_{ij}^{-1}A_i g_{ij}+g_{ij}^{-1}dg_{ij},
 \qquad
 g_{ik}=g_{ij}g_{jk}.
\label{eq:v4v86-transition}
\end{equation}
The compact gauge group acts by isometries in the fixed representation.

\subsection{Partition-flat reference}

Set \(g_{ii}=1\).  In the \(i\)-th frame define
\begin{equation}
 A_{\sharp,i}:=
 -\sum_j\rho_j\,dg_{ij}g_{ij}^{-1},
\label{eq:v4v86-Asharp}
\end{equation}
where a term is understood to vanish off \(U_i\cap U_j\).

\begin{lemma}[The partition-flat forms define a connection]
\label{lem:v4v86-connection}
The forms \(\{A_{\sharp,i}\}\) obey
\[
 A_{\sharp,k}
 =g_{ik}^{-1}A_{\sharp,i}g_{ik}+g_{ik}^{-1}dg_{ik}
\]
and hence define a global connection \(A_\sharp\) on the reconstructed
bundle \(P\).
\end{lemma}
\begin{proof}
The cocycle identity gives \(g_{kj}=g_{ik}^{-1}g_{ij}\).  Therefore
\[
 d g_{kj}g_{kj}^{-1}
 =
 -g_{ik}^{-1}dg_{ik}
 +g_{ik}^{-1}(dg_{ij}g_{ij}^{-1})g_{ik}.
\]
Multiply by \(-\rho_j\), sum over \(j\), and use
\(\sum_j\rho_j=1\).  The claimed affine transformation law follows.
\end{proof}

\begin{theorem}[Exact relative-potential identity]
\label{thm:v4v86-identity}
In every chart,
\begin{equation}
 a_i:=A_i-A_{\sharp,i}
 =
 \sum_j\rho_j\,g_{ij}A_jg_{ij}^{-1}.
\label{eq:v4v86-relative}
\end{equation}
In particular \(a=A-A_\sharp\) is the global tensor predicted by V85;
all transition derivatives belong to the reference owner \(A_\sharp\), not
to the analytic relative perturbation.
\end{theorem}
\begin{proof}
Equation \eqref{eq:v4v86-transition} is equivalent to
\[
 dg_{ij}g_{ij}^{-1}=g_{ij}A_jg_{ij}^{-1}-A_i.
\]
Insert this into \eqref{eq:v4v86-Asharp}:
\[
 A_{\sharp,i}
 =-\sum_j\rho_j(g_{ij}A_jg_{ij}^{-1}-A_i)
 =A_i-\sum_j\rho_jg_{ij}A_jg_{ij}^{-1}.
\]
Rearrangement proves \eqref{eq:v4v86-relative}.
\end{proof}

The identity is stronger than a bound.  It shows exactly where the
cross-chart stock resides.  Assigning \(dg_{ij}\) separately to both
adjacent local perturbations would double count it and destroy the global
connection law.

\subsection{Action-paid \(L^4\) estimate}

\begin{theorem}[Global critical payment]
\label{thm:v4v86-payment}
There is a constant depending only on the cover geometry, representation,
and overlap multiplicity such that
\begin{align}
 \|A-A_\sharp\|_{L^4(X)}^4
 &\le C_U^4\sum_j e_j^2,
\label{eq:v4v86-L4}\\
 \|A-A_\sharp\|_{L^4(X)}^4
 &\le C_U^4\varepsilon_U\,\nu
       \int_X|F_A|^2.
\label{eq:v4v86-action}
\end{align}
\end{theorem}
\begin{proof}
Conjugation preserves the fibre norm.  By convexity of \(|\cdot|^4\) and
\(\sum_j\rho_j=1\), Theorem~\ref{thm:v4v86-identity} gives pointwise
\[
 |a_i|^4
 \le\sum_j\rho_j|A_j|^4.
\]
Integrating with a measurable choice of one chart at each point, then
enlarging to the complete chart integrals, yields
\[
 \|a\|_4^4\le\sum_j\|A_j\|_4^4
 \le C_U^4\sum_je_j^2.
\]
Since \(e_j\le\varepsilon_U\),
\(\sum_je_j^2\le\varepsilon_U\sum_je_j\).  Bounded overlap gives
\(\sum_je_j\le\nu\int_X|F_A|^2\), proving
\eqref{eq:v4v86-action}.
\end{proof}

\begin{corollary}[Action-paid Sylvester dimension]
\label{cor:v4v86-rank}
For a fixed reference operator built from \(A_\sharp\), and with a Higgs
reference satisfying the analogous \(L^4\) payment, the V84 relative
crossing rank satisfies
\[
 R_X\le C(q,\delta_\circ,C_{\rm int},C_U)
 \left(\sum_je_j^2+\|\Phi-\Phi_\sharp\|_4^4\right).
\]
In particular, if the Higgs quartic energy pays
\(\|\Phi-\Phi_\sharp\|_4^4\), then \(R_X\) is linear in total action up to
the fixed factor \(\varepsilon_U\nu\) in the gauge part.
\end{corollary}
\begin{proof}
The V85 intrinsic Cwikel estimate is linear in the \(L^4\) coefficient
load, while the V84 singular-value count takes its fourth power.  Apply
\((x+y)^4\le8(x^4+y^4)\) and
Theorem~\ref{thm:v4v86-payment}.
\end{proof}

\subsection{Low regularity and smooth approximation}

The transition equation implies
\[
 \|dg_{ij}\|_{L^4(U_i\cap U_j)}
 \le C(\|A_i\|_4+\|A_j\|_4),
\]
so \(A_\sharp\) is an \(L^4\) connection.  This is sufficient to define
the Dirac form because \(H^1(X)\hookrightarrow L^4(X)\):
multiplication by \(A_\sharp\) maps \(H^1\) boundedly to \(L^2\).

\begin{lemma}[Smooth reference approximation]
\label{lem:v4v86-smooth}
For every \(\zeta>0\), there is a smooth connection
\(A_{\sharp,\zeta}\) on the same smooth bundle such that
\[
 \|A_{\sharp,\zeta}-A_\sharp\|_4<\zeta
\]
and
\[
 \|A-A_{\sharp,\zeta}\|_4^4
 \le8\|A-A_\sharp\|_4^4+8\zeta^4.
\]
\end{lemma}
\begin{proof}
The affine space of \(L^4\) connections is a translate of
\(L^4(T^*X\otimes\operatorname{ad}P)\).  Smooth adjoint-valued one-forms
are dense in this space.  Approximate the tensor difference from any fixed
smooth connection and translate back.  The last bound is
\((x+y)^4\le8(x^4+y^4)\).
\end{proof}

This approximation proves existence of a smooth global reference for each
configuration.  It does not provide uniform bounds on derivatives of that
reference as \(\zeta\downarrow0\); no such derivative bound is used in
\eqref{eq:v4v86-action}.

\subsection{Topological ownership}

The connection \(A_\sharp\) is assembled from the complete cocycle
\(\{g_{ij}\}\), so it lives on the same bundle and carries its characteristic
class.  The Chern--Weil difference is exact:
\begin{equation}
 \operatorname{tr}(F_A\wedge F_A)
 -\operatorname{tr}(F_{A_\sharp}\wedge F_{A_\sharp})
 =
 d\,\operatorname{tr}
 \left(2a\wedge F_{A_\sharp}
       +a\wedge d_{A_\sharp}a+\frac23a\wedge a\wedge a\right)
\label{eq:v4v86-transgression}
\end{equation}
for smooth approximants, and hence distributionally in the finite-energy
limit.  On a closed core the integrated characteristic number is identical;
with boundary, the displayed transgression is the exact boundary owner.
Thus the reference soft index contains the topology, while the relative
Sylvester block contains only analytic crossings.

\begin{proof}[Proof of \eqref{eq:v4v86-transgression}]
For \(A_s=A_\sharp+sa\),
\[
 \frac d{ds}\operatorname{tr}(F_{A_s}\wedge F_{A_s})
 =2d\,\operatorname{tr}(a\wedge F_{A_s})
\]
by the Bianchi identity and invariance of the trace.  Integrating from
\(s=0\) to \(1\) and expanding
\(F_{A_s}=F_{A_\sharp}+s\,d_{A_\sharp}a+s^2a\wedge a\)
gives the formula.
\end{proof}

\subsection{Marked-family obstruction and upstream target}

For a single configuration, V86 supplies a global, topologically correct,
action-paid reference.  For a two-marked family \(A(z)\), however, local
Uhlenbeck gauges and their transitions depend on \(z\).  Differentiating
\eqref{eq:v4v86-Asharp} gives
\[
 \partial_aA_{\sharp,i}
 =-\sum_j\rho_j\,\partial_a(dg_{ij}g_{ij}^{-1}),
\]
and a second derivative contains products of first transition derivatives.
The unmarked estimate \(\|dg_{ij}\|_4\lesssim\|A_i\|_4+\|A_j\|_4\) does
not control these parameter derivatives.

\begin{proposition}[Exact marked condition]
\label{prop:v4v86-marked}
Suppose the local based Coulomb gauges can be selected so that, for
\(|\alpha|\le2\),
\[
 \sum_{i,j}
 \|\partial_z^\alpha(dg_{ij}g_{ij}^{-1})\|_{L^4(U_i\cap U_j)}
 \le C_\alpha\mathcal M_\alpha .
\]
Assume also that the chosen reference soft cluster is separated from the
hard spectrum by a uniform collar.  Then \(A_\sharp(z)\) has two
\(L^4\) marked derivatives, the reference soft projection and hard inverse
have two marked derivatives after Kato transport, and the V85 marked
weak-\(\mathcal S_4\) theorem applies.
\end{proposition}
\begin{proof}
Differentiating the finite sum \eqref{eq:v4v86-Asharp} proves the asserted
bounds for \(A_\sharp\).  Let \(\Gamma\) be a contour in the spectral
collar.  The Riesz formula
\[
 P_\sharp(z)=\frac1{2\pi i}\int_\Gamma
     (\lambda-H_\sharp(z))^{-1}\,d\lambda
\]
and its first two differentiated forms control the marked soft projection.
Kato parallel transport then identifies \(\operatorname{Ran}(1-P_\sharp(z))\)
with one fixed hard space.  On that fixed space the transported inverse
\(\widehat C_\sharp=\widehat H_\sharp^{-1}\) satisfies
\[
 \partial_a\widehat C_\sharp
 =-\widehat C_\sharp(\partial_a\widehat H_\sharp)\widehat C_\sharp.
\]
A second derivative consists of the term containing
\(\partial_a\partial_b\widehat H_\sharp\) and the two ordered products of
first derivatives, each between three hard inverses.  The spectral collar,
transition-mark bounds, and Cwikel ideal property therefore give the two
marked weak-\(\mathcal S_4\) estimates.
\end{proof}

The next step is no longer the unmarked global gauge problem:
Theorem~\ref{thm:v4v86-payment} closes it.  The upstream bottleneck is
smooth parameter selection of based Uhlenbeck gauges at critical
regularity.  V87 will fix the gauge stabilizer by a based boundary
condition, derive the linearized Coulomb equation, and identify the exact
Faddeev--Popov inverse needed for two marks.  If its gap is not
action-controlled, the construction will revert to a finite slice atlas
rather than assume a global gauge section.

\clearpage
\section{V87: frozen-centre marked charts and gauge-basic sewing}
\label{sec:v4v87}

V86 proposed differentiating a parameter-dependent Uhlenbeck gauge.  At
critical regularity this is stronger than necessary.  The determinant and
its marked insertions are local derivatives at a background configuration.
One may construct the action-paid cocycle reference at that background and
freeze it while differentiating.  Different backgrounds use different
charts; their overlap is sewn by the anomaly-normalized determinant-line
isomorphism already proved in V74.  No new anomaly calculation is made here.

\subsection{What a based Coulomb slice would provide}

On a ball, impose the based condition \(\xi|_{\partial U}=0\) on infinitesimal
gauge transformations.  In a Coulomb gauge \(d^*A=0\), the linearized gauge
fixing operator is
\[
 L_A\xi:=d^*d_A\xi.
\]

\begin{lemma}[Small-energy Faddeev--Popov coercivity]
\label{lem:v4v87-FP}
There is \(c_*>0\) such that if
\(C_{\rm ad}C_S\|A\|_4\le1/2\), then for
\(\xi\in H^1_0(U;\operatorname{ad}P)\),
\begin{equation}
 \langle L_A\xi,\xi\rangle
 \ge\frac12\|d\xi\|_2^2,
\label{eq:v4v87-coercive}
\end{equation}
and
\[
 L_A:H^1_0\longrightarrow H^{-1}
\]
is invertible with a uniform inverse bound.
\end{lemma}
\begin{proof}
Weak integration by parts gives
\[
 \langle L_A\xi,\xi\rangle
 =\|d\xi\|_2^2+\langle[A,\xi],d\xi\rangle.
\]
The representation bound and \(H^1_0\hookrightarrow L^4\) imply
\[
 |\langle[A,\xi],d\xi\rangle|
 \le C_{\rm ad}\|A\|_4\|\xi\|_4\|d\xi\|_2
 \le C_{\rm ad}C_S\|A\|_4\|d\xi\|_2^2.
\]
This proves \eqref{eq:v4v87-coercive}.  Boundedness and coercivity give
invertibility by Lax--Milgram.
\end{proof}

The lemma removes stabilizer zero modes and is enough for an \(H^1\)
implicit slice.  It does not give the endpoint estimate
\(L_A^{-1}:W^{-1,4}\to W^{1,4}\) from critical \(A\in L^4\): attempting to
absorb the perturbation requires
\(\|A\xi\|_4\le\|A\|_4\|\xi\|_\infty\), while
\(W^{1,4}\) in four dimensions does not embed in \(L^\infty\).  We therefore
do not build the two-marked theorem on an unproved endpoint inverse.

\subsection{Frozen-centre construction}

Fix one physical background \(b=(A,\Phi,g)\) on a reconstructed core.  Use
V86 once to choose its Uhlenbeck cover, cocycle, and partition-flat reference
\[
 b_\sharp=(A_\sharp,\Phi_\sharp,g_\sharp).
\]
The metric reference is chosen by the already established bounded-geometry
core chart; only the gauge--Higgs terms are varied here.  Transport all
objects by one fixed bundle isomorphism into the exact bundle carrying
\(b\), and then keep \(b_\sharp\), its reference soft projection, and its
hard resolvent fixed.

A tangent mark at \(b\) is a global tensor
\[
 h=(\dot A,\dot\Phi,\dot g).
\]
It is transformed by the one fixed background gauge, not by a
parameter-dependent gauge.  Thus no derivative of an Uhlenbeck
transformation occurs.

\begin{definition}[Frozen-centre chart]
\label{def:v4v87-frozen}
For fixed reserves \(\epsilon_{\rm rel}>0\) and
\(\epsilon_{\rm sp}>0\), let \(\mathcal U_b\) be the set of backgrounds
\(b'\) on the same reconstructed bundle for which
\[
 \|A'-A\|_4+\|\Phi'-\Phi\|_4+
 \|g'-g\|_{\mathfrak G}<\epsilon_{\rm rel}
\]
and the reference soft and singular-value contours retain collars at least
\(\epsilon_{\rm sp}\).  Here \(\|\cdot\|_{\mathfrak G}\) is the metric norm
already used in the Einstein bounded-geometry chart.  The reference
\(b_\sharp\) is frozen throughout \(\mathcal U_b\).
\end{definition}

\begin{theorem}[Two-marked weak-\(\mathcal S_4\) chart without gauge
derivatives]
\label{thm:v4v87-frozen-marks}
Let \(z\mapsto b(z)\in\mathcal U_b\) be \(C^2\), with
\[
 \partial^\alpha(A(z)-A_\sharp),\quad
 \partial^\alpha(\Phi(z)-\Phi_\sharp)
 \in L^4,\qquad 1\le|\alpha|\le2,
\]
and the corresponding two metric marks in the established Einstein
operator norm.  Then the hard relative operator
\[
 K_b(z)=Q_{\sharp,\rm h}
       \bigl(H_{b(z)}-H_{b_\sharp}\bigr)
       Q_{\sharp,\rm h}C_\sharp
\]
has two weak-\(\mathcal S_4\) derivatives, with
\begin{equation}
 \|\partial^\alpha K_b(z)\|_{4,\infty}
 \le C_b\left(
   \|\partial^\alpha A(z)\|_4+
   \|\partial^\alpha\Phi(z)\|_4+
   \|\partial^\alpha g(z)\|_{\mathfrak G}\right),
 \qquad |\alpha|\le2,
\label{eq:v4v87-mark-bound}
\end{equation}
where at order zero the right side is replaced by the relative loads.
The V84 finite crossing matrix and its first two marks are therefore
defined on \(\mathcal U_b\).
\end{theorem}
\begin{proof}
Every reference factor is fixed.  Lemma~\ref{lem:v4v85-difference} makes
the gauge and Higgs differences global zero-order multipliers, so their
derivatives are the displayed tensor marks.  The metric variation formula
for a Dirac operator is the previously proved sum of a principal first-order
metric mark and lower-order spin-connection marks; the bounded-geometry
norm controls them.  Apply the intrinsic Cwikel estimate
Theorem~\ref{thm:v4v85-intrinsic} to the zero-order pieces and the existing
Einstein relative-resolvent estimate to the metric pieces.  The fixed hard
projection and inverse produce no derivative terms.  The V84 formulas then
differentiate the tail inverse and finite Sylvester matrix twice inside the
two spectral collars.
\end{proof}

\begin{corollary}[Uniform rank at the chart centre]
\label{cor:v4v87-rank}
At \(z=b\), the gauge contribution to the crossing rank satisfies
\[
 R_b^{\rm gauge}
 \le C(q,\delta_\sharp,C_U)\sum_i e_i^2
 \le C(q,\delta_\sharp,C_U)\varepsilon_U\nu
      \int_X|F_A|^2.
\]
After shrinking \(\mathcal U_b\), the same bound holds with a fixed
multiplicative reserve for every \(b'\in\mathcal U_b\).
\end{corollary}
\begin{proof}
At the centre this is Corollary~\ref{cor:v4v86-rank}.  The relative
\(L^4\) load and spectral collars vary continuously in the chart.  Choose
the radius so that the load increases by at most the selected reserve and
apply the V84 singular-value count.
\end{proof}

\subsection{Gauge-basic derivatives}

Let \(\mathcal L_{\rm SM}\to\mathcal B\) be the complete, three-generation,
anomaly-normalized Standard Model determinant line.  V74 imports its
vanishing local anomaly curvature, trivial faithful global gauge character,
and parallel total phase section \(\sigma_{\rm tot}\).  Write the physical
determinant section as
\[
 \operatorname{Det}_{\rm SM}(b)=Z_{\rm F}(b)\,\sigma_{\rm tot}(b).
\]

\begin{theorem}[Vertical gauge marks vanish on the total coefficient]
\label{thm:v4v87-basic}
For every admissible gauge transformation \(u\),
\begin{equation}
 Z_{\rm F}(b^u)=Z_{\rm F}(b).
\label{eq:v4v87-gauge-invariance}
\end{equation}
Consequently, for a \(C^2\) two-parameter family \(b(z)\) and any
\(C^2\) family \(u(z)\),
\[
 \partial_a Z_{\rm F}(b(z)^{u(z)})
 =\partial_a Z_{\rm F}(b(z)),
 \qquad
 \partial_a\partial_c Z_{\rm F}(b(z)^{u(z)})
 =\partial_a\partial_c Z_{\rm F}(b(z)).
\]
No norm of \(\partial u\) or \(\partial^2u\) appears.
\end{theorem}
\begin{proof}
Gauge covariance identifies the determinant lines at \(b\) and \(b^u\).
Theorem~\ref{thm:v4v74-anomaly-sewing} says that the complete parallel
phase section has trivial gauge holonomy; the physical determinant section
is transported by the same line isomorphism.  Their scalar coefficient is
therefore identical, proving \eqref{eq:v4v87-gauge-invariance}.  The two
derivative identities follow by differentiating this exact equality of
scalar functions.  They remain true at determinant zeros because no
logarithm or inverse determinant is used.
\end{proof}

\begin{firewall}[Only the total line is gauge basic]
\label{fw:v4v87-total-only}
The hard determinant, the soft Schur determinant, and a single chiral
species may each acquire a frame connection or anomaly phase.  Theorem
\ref{thm:v4v87-basic} applies to their already normalized Standard Model
product.  Applying it to either factor separately would contradict the V74
curvature identity and double count the anomaly redistribution.
\end{firewall}

\subsection{Overlap of frozen-centre charts}

Let \(\mathcal U_b\cap\mathcal U_c\ne\varnothing\).  The two charts may use
different Uhlenbeck cocycles, references, soft ranks, and Sylvester ranks.
Both nevertheless factor the same physical Fredholm operator.

\begin{theorem}[Reference-change cocycle]
\label{thm:v4v87-overlap}
There is a canonical determinant-line isomorphism
\[
 \Theta_{bc}:\mathcal L_b\longrightarrow\mathcal L_c
\]
obtained by composing the two hard--soft line factorizations with the
identity on the complete physical determinant line.  It is defined at
zeros, respects the anomaly-normalized total section, and satisfies
\begin{equation}
 \Theta_{cd}\Theta_{bc}=\Theta_{bd}
\label{eq:v4v87-cocycle}
\end{equation}
on triple overlaps.  If the transferred ranks differ, the finite crossing
line of Proposition~\ref{prop:v4v74-rank-change} supplies the missing
exterior-power factor.
\end{theorem}
\begin{proof}
Theorem~\ref{thm:v4v74-line-factorization} gives an exterior-power
isomorphism from each chart factorization to the complete determinant line.
Define \(\Theta_{bc}\) by going from the \(b\)-factorization to the complete
line and back through the inverse \(c\)-factorization.  Exterior-power
isomorphisms do not require a nonzero determinant.  The V74 rank-change
proposition handles unequal finite summands.  On a triple overlap the
middle complete-line map and its inverse cancel, which proves
\eqref{eq:v4v87-cocycle}.  The anomaly-normalized section is the same
complete-line section in both constructions.
\end{proof}

\begin{corollary}[No global smooth gauge section is required]
\label{cor:v4v87-no-global-slice}
The local twice-marked coefficients and determinant-line sections from
Theorem~\ref{thm:v4v87-frozen-marks} descend to the gauge quotient by the
cocycle \eqref{eq:v4v87-cocycle}.  Their definition requires one
small-energy Uhlenbeck choice at each chart centre, but never a globally
smooth choice of Uhlenbeck gauge.
\end{corollary}
\begin{proof}
Theorem~\ref{thm:v4v87-overlap} gives descent, and
Theorem~\ref{thm:v4v87-basic} removes dependence of the total scalar
coefficient on the centre gauge.  The cocycle identity is the descent
condition.
\end{proof}

\subsection{Remaining uniformity gate}

V87 closes the two-marked gauge-selection issue at the level of local
differential estimates and determinant-line descent.  A new global claim
would still require a locally finite measurable collection of frozen-centre
charts with uniform collars and an action-paid overlap incidence.  An
arbitrary open cover of the infinite-dimensional configuration quotient
does not provide those properties.

V88 will therefore construct a countable spectral-collar atlas using
separable \(L^4\) configuration charts.  It will test whether the V84 weak
singular-value count pays both the number of rank changes and the overlap
incidence.  If uniform collars fail near crossing hypersurfaces, the notes
will retain the determinant line and stratify by finite crossing rank
instead of assuming one global scalar frame.

\clearpage
\section{V88: quantitative spectral-collar atlas}
\label{sec:v4v88}

V87 leaves a covering problem: every frozen-centre chart needs a
singular-value threshold separated from the spectrum.  Compactness gives
such a threshold pointwise, but the continuum argument needs quantitative
rank and collar bounds.  Weak-\(\mathcal S_4\) control supplies both by a
finite pigeonhole argument.

\subsection{A quantitative empty interval}

Fix
\[
 a:=\frac18,\qquad b:=\frac14.
\]
For a compact operator \(K\), write
\(\Sigma(K)=\{s_n(K):n\ge1\}\setminus\{0\}\).

\begin{theorem}[Action-paid collar selection]
\label{thm:v4v88-collar}
Suppose \(K\in\mathcal S_{4,\infty}\) and
\(\|K\|_{4,\infty}\le M\).  Put
\[
 N_M:=\left\lceil(M/a)^4\right\rceil,\qquad
 \eta_M:=\frac{b-a}{2(N_M+1)}.
\]
There is \(\tau\in[a,b]\) such that
\begin{equation}
 \operatorname{dist}(\tau,\Sigma(K))\ge\eta_M.
\label{eq:v4v88-gap}
\end{equation}
With \(q:=2\tau\), the V84 truncation has
\begin{align}
 \frac14\le q\le\frac12,\qquad
 \|(1+tK_<)^{-1}\|&\le2,
\label{eq:v4v88-reserve}\\
 \operatorname{rank}K_>
 &\le(M/a)^4.
\label{eq:v4v88-rank}
\end{align}
\end{theorem}
\begin{proof}
There are at most \((M/a)^4\le N_M\) singular values greater than \(a\).
Insert the singular values lying in \([a,b]\) between the two endpoints.
These at most \(N_M\) points divide \([a,b]\) into at most \(N_M+1\)
closed gaps whose total length is \(b-a\).  One gap has length at least
\((b-a)/(N_M+1)\).  Its midpoint \(\tau\) has distance at least
\(\eta_M\) from every point of \(\Sigma(K)\cap[a,b]\), and singular values
outside \([a,b]\) are no closer than the relevant endpoint gap.

The V84 tail norm is at most \(q\), so \(q\le1/2\) gives the inverse bound
in \eqref{eq:v4v88-reserve}.  Its finite range uses singular values greater
than \(q/2=\tau\ge a\); weak-Schatten counting gives
\eqref{eq:v4v88-rank}.
\end{proof}

\begin{lemma}[Persistence under perturbation]
\label{lem:v4v88-persist}
If \(\|K'-K\|<\eta_M/2\), then
\[
 \operatorname{dist}(\tau,\Sigma(K'))\ge\eta_M/2
\]
and the number of singular values above \(\tau\) is unchanged.
\end{lemma}
\begin{proof}
The min--max principle gives
\[
 |s_n(K')-s_n(K)|\le\|K'-K\|
\]
for every \(n\).  Hence no singular value can enter the half-sized collar.
Along the straight interpolation from \(K\) to \(K'\), none crosses
\(\tau\), so the integer-valued rank is constant.
\end{proof}

This is the spectral analogue of the stopping-family payer.  The collar
shrinks only when the weak-Schatten load creates more singular values in the
fixed window; its inverse cost is explicit:
\[
 \eta_M^{-1}=\frac{2(N_M+1)}{b-a}
 \le16\left(2+(8M)^4\right).
\]
A collar is therefore polynomially paid, not merely asserted.

\subsection{Uniform bounded-action charts}

Let \(\mathfrak C_E\) be one bounded-action, bounded-geometry sector of the
reconstructed SM--Einstein configuration space.  Assume the results of
V85--V87 give, at each centre \(c\in\mathfrak C_E\), a frozen relative
operator \(K_c\) satisfying
\begin{equation}
 \|K_c\|_{4,\infty}\le M_E
\label{eq:v4v88-uniformM}
\end{equation}
with a constant depending only on the sector bounds, not on \(c\).
Choose \(\tau_c\) by Theorem~\ref{thm:v4v88-collar} and let
\(\mathcal U_c\) be a frozen-centre neighborhood on which
\begin{equation}
 \|K_c(z)-K_c(c)\|<\eta_{M_E}/2.
\label{eq:v4v88-opball}
\end{equation}

\begin{theorem}[Uniform collared frozen atlas]
\label{thm:v4v88-atlas}
The sets \(\{\mathcal U_c\}_{c\in\mathfrak C_E}\) cover
\(\mathfrak C_E\).  Every chart has:
\begin{enumerate}
 \item tail inverse reserve at least \(1/2\);
 \item singular collar at least \(\eta_{M_E}/2\);
 \item crossing rank at most \((M_E/a)^4\); and
 \item the two-marked estimates of V87 with constants polynomial in
       \(\eta_{M_E}^{-1}\).
\end{enumerate}
\end{theorem}
\begin{proof}
Each centre belongs to its own operator-norm neighborhood.  The first three
claims are Theorem~\ref{thm:v4v88-collar} and
Lemma~\ref{lem:v4v88-persist}.  Riesz differentiation once or twice inserts
two or three resolvents around the collar, so the marked constants cost at
most fixed powers of \(\eta_{M_E}^{-1}\).  All other constants are uniform
by \eqref{eq:v4v88-uniformM}.
\end{proof}

\subsection{Countability without finite-dimensional covering claims}

The Sobolev and \(L^4\) spaces used for the gauge, Higgs, and metric
coordinates are separable.  Use based local gauge charts before taking the
gauge quotient; their disjoint union is second countable.

\begin{proposition}[Countable subatlas]
\label{prop:v4v88-countable}
The cover in Theorem~\ref{thm:v4v88-atlas} has a countable subcover
\(\{\mathcal U_n\}_{n\ge1}\).  On every overlap, the V87 reference-change
maps satisfy the determinant-line cocycle, so these countably many local
lines define one global line over the sector.
\end{proposition}
\begin{proof}
Every second-countable space is Lindelof: for each basis element whose
closure is contained in some \(\mathcal U_c\), select one such chart.  The
selected family is countable and covers.  The line descent is
Theorem~\ref{thm:v4v87-overlap}; its triple-overlap identity is exactly the
cocycle condition.
\end{proof}

There is no uniform finite-overlap theorem here.  Infinite-dimensional
Banach spaces do not have a dimension-controlled Besicovitch covering
constant.  Such a claim would import a nonexistent finite-dimensional
packing argument into configuration space.

\subsection{No overlap sum is needed}

Enumerate the subatlas and define disjoint Borel sets
\[
 S_1:=\mathcal U_1,\qquad
 S_n:=\mathcal U_n\setminus\bigcup_{j<n}\mathcal U_j.
\]
They cover the sector.

\begin{theorem}[Borel selection has zero multiplicity cost]
\label{thm:v4v88-selection}
Let \(Z_n\) be the scalar coefficient of the complete anomaly-normalized
determinant section in chart \(\mathcal U_n\).  Then
\[
 Z(c):=Z_n(c)\quad\text{for }c\in S_n
\]
is Borel and independent of the enumeration.  It equals the same global
coefficient represented by the determinant-line descent.  No sum over
overlaps and no derivative of the sets \(S_n\) occurs.
\end{theorem}
\begin{proof}
The sets \(S_n\) are Borel because the \(\mathcal U_n\) are open.
Each \(Z_n\) is continuous, hence its restriction is Borel.  On an overlap,
Theorem~\ref{thm:v4v87-overlap} identifies the physical determinant
sections and Theorem~\ref{thm:v4v87-basic} identifies their
anomaly-normalized scalar coefficients.  Thus \(Z_n=Z_m\) wherever both
are defined.  The selected value is consequently independent of the first
chart containing the point.  Since no partition of unity is inserted into
\(Z\), no overlap multiplicity or cutoff derivative is generated.
\end{proof}

\begin{firewall}[Selection is not a differentiable partition]
\label{fw:v4v88-selection}
The Borel sets \(S_n\) may have rough boundaries.  Theorem
\ref{thm:v4v88-selection} does not differentiate their indicators.
Marked derivatives are computed in any open chart containing the
background, where all chart representatives agree as derivatives of the
same global line section.  Replacing this descent by
\(\sum_n\chi_nZ_n\) would create unnecessary cutoff derivatives and an
unpaid overlap sum.
\end{firewall}

\subsection{Rank changes and crossing strata}

For a chart threshold \(\tau_n\), set
\[
 \mathfrak C_{n,R}
 :=\{c\in\mathcal U_n:
       \operatorname{rank}\mathbf1_{(\tau_n,\infty)}(|K_n(c)|)=R\}.
\]
The collar makes this rank locally constant, so
\(\mathfrak C_{n,R}\) is open relative to the collared chart.  When a path
leaves that collar, the rank-change exterior line of V74, not a scalar
division by a vanishing determinant, gives the transition.

\begin{proposition}[Finite rank change on a rectifiable fixed-reference
path]
\label{prop:v4v88-path}
Let \(K:[0,1]\to\mathcal S_{4,\infty}\) be operator-norm rectifiable, with
\(\|K(t)\|_{4,\infty}\le M\).  A greedy choice of collar centres using
Theorem~\ref{thm:v4v88-collar} needs at most
\[
 1+\left\lceil\frac{4\,\operatorname{Length}_{\|\cdot\|}K}{\eta_M}\right\rceil
\]
chart changes.
\end{proposition}
\begin{proof}
After choosing a centre, retain its chart until the operator-norm
displacement reaches \(\eta_M/4\).  Except possibly for the final segment,
each change consumes at least \(\eta_M/4\) of path length.  The stated
count follows.  Lemma~\ref{lem:v4v88-persist} preserves the collar on each
segment.
\end{proof}

Action alone does not bound the variation length of an arbitrarily
oscillatory parameter path.  This is the direct analogue of the NS
first-moment boundary: a static zeroth-order payer cannot be reused as
parameter-variation stock.

\subsection{Result and next gate}

V88 proves a countable, uniformly collared, action-ranked atlas and glues it
without overlap multiplicity.  The spectral-section obstruction isolated in
V82--V83 is now resolved at the determinant-line level: local finite
crossing blocks descend globally, including rank changes and determinant
zeros.

The next continuum gate is not another spectral cover.  The remaining
Einstein coupling in Theorem~\ref{thm:v4v87-frozen-marks} was imported as a
bounded-geometry relative estimate.  V89 will expose that estimate:
construct the Bourguignon--Gauduchon identification for two nearby metrics,
derive the exact principal and spin-connection variations, and determine
which metric norm yields weak-\(\mathcal S_4\) control without losing
diffeomorphism covariance.

\clearpage
\section{V89: metric-principal-symbol obstruction and microlocal normalization}
\label{sec:v4v89}

The metric clause in Theorem~\ref{thm:v4v87-frozen-marks} was too strong.
The gauge and Higgs variations are zero-order multipliers, but a metric
variation changes the principal symbol of the Dirac operator.  Its product
with a reference inverse is order zero and need not be compact.  This note
corrects that point, proves the obstruction, and constructs the appropriate
principal-symbol normalization.

The spinor comparison and first-variation formula used below are the
Bourguignon--Gauduchon construction (Commun.\ Math.\ Phys.\ \textbf{144}
(1992), 581--599, DOI \(10.1007/\mathrm{BF02099184}\)).

\subsection{Fixed spinor bundle and exact first variation}

Let \(g_t=g+th+O(t^2)\) be a uniformly elliptic metric path.  The positive
\(g\)-self-adjoint map \(B_t\) defined by
\[
 g_t(B_tX,B_tY)=g(X,Y)
\]
lifts canonically to an isometry \(\beta_t:\Sigma_g\to\Sigma_{g_t}\).
Transport \(D_{g_t}\) back by \(\beta_t\).  In a \(g\)-orthonormal frame
\(\{e_i\}\), and with
\[
 (\operatorname{div}h)(X):=\sum_i(\nabla_{e_i}h)(e_i,X),
\]
the first variation is
\begin{equation}
 \dot D_g[h]
 =
 -\frac12\sum_{i,j}h(e_i,e_j)c(e_i)\nabla_{e_j}
 +\frac14c\!\left(d\operatorname{tr}_g h-\operatorname{div}h\right).
\label{eq:v4v89-BG}
\end{equation}
A unitary correction for the varying volume density changes only the
zero-order term; the first-order term in \eqref{eq:v4v89-BG} is unchanged.

\begin{lemma}[Principal metric mark]
\label{lem:v4v89-principal}
The principal symbol of \(\dot D_g[h]\) is
\[
 \sigma_1(\dot D_g[h])(x,\xi)
 =
 -\frac i2\sum_{i,j}h_{ij}(x)\xi_jc(e_i).
\]
It vanishes for every \(\xi\) only when \(h=0\).
\end{lemma}
\begin{proof}
Only the covariant derivative in the first term of
\eqref{eq:v4v89-BG} contributes to the principal symbol.  Clifford
multiplication is faithful on one-forms, so vanishing for every \(\xi\)
forces the endomorphism \(h^\sharp:T_xX\to T_xX\) to vanish.
\end{proof}

\subsection{The direct weak-Schatten claim is false}

Let \(C_g=D_g^{-1}\) on a collared hard complement.

\begin{theorem}[Noncompact metric-relative operator]
\label{thm:v4v89-noncompact}
If \(h\) is nonzero on an open set, then, in general,
\[
 \dot D_g[h]\,C_g
\]
is a classical order-zero pseudodifferential operator with nonzero principal
symbol.  It is not compact on \(L^2\), and hence belongs to no Schatten or
weak-Schatten ideal \(\mathcal S_{p,\infty}\) with finite \(p\).
\end{theorem}
\begin{proof}
The hard inverse has principal symbol
\(\sigma_{-1}(C_g)=\sigma_1(D_g)^{-1}\).  By
Lemma~\ref{lem:v4v89-principal}, the product has a generally nonzero
order-zero principal symbol.  To see noncompactness directly, choose
\((x_0,\xi_0)\) where that symbol is nonzero and normalized semiclassical
wave packets \(\psi_k\) localized near \(x_0\) with frequency
\(k\xi_0\).  They converge weakly to zero, while pseudodifferential
symbol calculus gives
\[
 \|\dot D_g[h]C_g\psi_k\|_2
 \longrightarrow
 \left\|\sigma_0(\dot D_g[h]C_g)(x_0,\xi_0)v\right\|>0
\]
for a suitable spinor \(v\).  A compact operator sends every bounded weakly
null sequence to a norm-null sequence, a contradiction.
\end{proof}

\begin{corollary}[Conformal test]
\label{cor:v4v89-conformal}
For \(g_t=e^{2tf}g\) with nonzero smooth \(f\),
\[
 \dot D_g\,C_g=-M_f+\text{an operator of order }-1
\]
up to the harmless convention-dependent density conjugation.  Thus even a
smooth conformal metric mark fails the V84 compact-relative hypothesis.
\end{corollary}
\begin{proof}
Conformal covariance gives
\[
 D_{e^{2tf}g}
 =
 e^{-(n+1)tf/2}D_g e^{(n-1)tf/2}
\]
under the canonical spinor identification.  Differentiation yields
\(\dot D_g=-fD_g+\frac{n-1}{2}c(df)\).  Multiply by \(C_g\).  The first
term is \(-M_f\) on the hard space; the second has order \(-1\).
Multiplication by a nonzero function on a nonatomic \(L^2\) space is not
compact.
\end{proof}

\begin{correction}[Scope of V87--V88]
\label{corr:v4v89-scope}
Theorems~\ref{thm:v4v87-frozen-marks} and
\ref{thm:v4v88-atlas} remain valid for gauge--Higgs zero-order marks at a
fixed metric, and for metric families only after the normalization below.
Their unnormalized metric weak-\(\mathcal S_4\) sentence is withdrawn.
This correction does not affect the cocycle, collar, or rank arguments.
\end{correction}

\subsection{Principal-symbol normalization}

Fix a reference metric \(g_\circ\) and use the
Bourguignon--Gauduchon map to place \(D_g\) and \(D_\circ\) on the same
spinor bundle.  Their principal symbols are invertible for \(\xi\ne0\).
Define the homogeneous degree-zero symbol
\begin{equation}
 q_0(g,g_\circ;x,\xi)
 :=
 \sigma_1(D_\circ)(x,\xi)\,
 \sigma_1(D_g)(x,\xi)^{-1}.
\label{eq:v4v89-q0}
\end{equation}

\begin{theorem}[Microlocal principal normalization]
\label{thm:v4v89-normalize}
Assume \(g,g_\circ\) are uniformly elliptic with bounded geometry.  There
is a classical elliptic operator
\(Q_{g,\circ}\in\Psi^0\), invertible after a finite-rank smoothing
modification, whose principal symbol is \eqref{eq:v4v89-q0}, such that
\begin{equation}
 L_{g,\circ}:=Q_{g,\circ}D_g-D_\circ\in\Psi^0.
\label{eq:v4v89-L}
\end{equation}
Consequently
\begin{equation}
 K^{\rm met}_{g,\circ}:=
 L_{g,\circ}C_\circ\in\Psi^{-1}
 \subset\mathcal S_{4,\infty}.
\label{eq:v4v89-Kmetric}
\end{equation}
If \(Q_{g,\circ}\) is invertible, then
\[
 \ker D_g=\ker(Q_{g,\circ}D_g),
 \qquad
 \operatorname{coker}D_g\cong
 \operatorname{coker}(Q_{g,\circ}D_g).
\]
Thus the normalized finite crossing matrix detects exactly the physical
Dirac zeros.
\end{theorem}
\begin{proof}
Quantize the elliptic symbol \(q_0\) by a geometric order-zero
quantization and add lower symbols arbitrarily.  Equation
\eqref{eq:v4v89-q0} makes the order-one principal symbols of
\(Q_{g,\circ}D_g\) and \(D_\circ\) identical, so their difference has
order at most zero.  An elliptic order-zero operator is Fredholm; adding a
map between its finite kernel and cokernel makes it invertible whenever its
index is zero.  Here \(q_0\) is homotopic through metric principal symbols
to the identity inside a metric chart, so the index is zero.

The product of an order-zero classical operator and the hard
order-\(-1\) inverse is order \(-1\).  Weyl singular-value asymptotics in
dimension four give
\(\Psi^{-1}\subset\mathcal S_{4,\infty}\).  Finally, left multiplication
by an invertible \(Q_{g,\circ}\) preserves the kernel and identifies the
cokernel.
\end{proof}

\begin{proposition}[Quantitative metric load]
\label{prop:v4v89-load}
On a uniformly elliptic metric chart with
\[
 g-g_\circ\in L^\infty\cap W^{1,4},
\]
one may choose the normalized symbols so that
\begin{equation}
 \|K^{\rm met}_{g,\circ}\|_{4,\infty}
 \le C_{\rm geo}\left(
   \|g-g_\circ\|_\infty+\|\nabla^\circ(g-g_\circ)\|_4\right),
\label{eq:v4v89-load}
\end{equation}
locally in the metric-chart radius.  The same estimate holds through two
marks if the marked metric differences lie in
\(L^\infty\cap W^{1,4}\) and the ellipticity reserve is fixed.
\end{proposition}
\begin{proof}
The degree-zero symbol \(q_0-1\) is a smooth finite-dimensional function of
\(g-g_\circ\) and the inverse metric, hence is bounded by the
\(L^\infty\) metric difference on a fixed ellipticity set.  Composition of
\(Q_{g,\circ}\) with \(D_g\) cancels the order-one term; its remaining
order-zero coefficients consist of the spin-connection difference,
controlled by \(\nabla^\circ(g-g_\circ)\), and symbol-composition terms
with one \(x\)-derivative of \(q_0\).  Local order-\(-1\) Cwikel estimates,
a finite chart partition, and the ideal property give
\eqref{eq:v4v89-load}.  Differentiating the smooth symbol formulas twice
produces only products controlled by the fixed ellipticity reserve and the
stated marked norms.
\end{proof}

\subsection{The order-zero factor is a local gravitational owner}

The normalization does not allow one to discard \(Q_{g,\circ}\).  It
separates the fermion determinant into:
\begin{enumerate}
 \item the normalized Fredholm factor
 \(Q_{g,\circ}D_g=D_\circ+L_{g,\circ}\), whose relative perturbation is
 weak-\(\mathcal S_4\) and is handled by V84--V88; and
 \item an elliptic order-zero metric factor \(Q_{g,\circ}^{-1}\), whose
 ordinary Hilbert-space determinant does not exist.
\end{enumerate}

\begin{theorem}[Locality of the metric ultraviolet defect]
\label{thm:v4v89-localdefect}
For a smooth metric path, the cutoff dependence of the logarithmic
determinant of \(Q_{g,\circ}\), and the multiplicative anomaly in
\[
 D_g=Q_{g,\circ}^{-1}(D_\circ+L_{g,\circ}),
\]
is determined by finitely many homogeneous symbol coefficients through
degree \(-4\).  Its variation is an integral of local covariant densities of
engineering dimension at most four.  These are precisely the permitted
cosmological, Einstein--Hilbert, curvature-squared, and mixed
gauge--Higgs counterterm owners.
\end{theorem}
\begin{proof}
In four dimensions, the divergent part of a classical elliptic determinant
is the noncommutative residue of the degree-\(-4\) symbol of its logarithm,
together with the lower heat coefficients.  Symbol composition determines
each coefficient at \(x\) from finitely many jets of the metric and the
other background fields.  The residue is the cosphere integral of that
local coefficient, so its metric variation is local and diffeomorphism
covariant.  Dimensional counting restricts the scalar densities to total
dimension at most four.  The finite nonlocal remainder is carried by the
normalized Fredholm determinant rather than by a counterterm.
\end{proof}

\begin{firewall}[Local renormalization is still a gate]
\label{fw:v4v89-renorm}
Theorem~\ref{thm:v4v89-localdefect} proves locality and the finite list of
owners; it does not prove that the existing SM--Einstein bare couplings
satisfy the required cofinal renormalization conditions, reflection
positivity, or a continuum limit.  Treating \(\det Q_{g,\circ}\) as a
nowhere-zero unrenormalized scalar would hide the gravitational ultraviolet
divergence.
\end{firewall}

\subsection{Diffeomorphism covariance}

If \(g'=\phi^*g\) and \(\phi\) preserves the spin structure, its spin lift
unitarily intertwines the physical Dirac operators.  Choose the symbol
\eqref{eq:v4v89-q0} and its quantization naturally from the metric
connections.  Then the normalized Fredholm line is transported by the same
spin lift.  Any two natural quantizations with the same principal symbol
differ by \(\Psi^{-1}\); that difference belongs to the already controlled
weak-\(\mathcal S_4\) relative sector.  The remaining change of determinant
is a local symbol anomaly and belongs to the counterterm owner of
Theorem~\ref{thm:v4v89-localdefect}.

V89 therefore restores the spectral program for metric variations without
misclassifying an order-zero operator as compact.  V90 is a compulsory
companion audit.  It will then write the explicit cofinal metric
renormalization card: heat coefficients, allowed coupling owners, and the
reference-change cocycle needed to combine the normalized determinant with
the Einstein action.

\clearpage
\section{V90: companion audit and the metric heat-counterterm card}
\label{sec:v4v90}

This is the compulsory V90 companion audit.  A stable snapshot of the
current Yang--Mills V10 note was read at SHA--256
\[
 \texttt{66b4756faee43f344b0c5e07759852a9abed221a902988a4efcb88fa5c1942cf},
\]
and the unchanged Navier--Stokes V31 note at
\[
 \texttt{f1121b62238ad5491bb7cf03635ea9934942edf573ed8faaa9ee79e1d38a6696}.
\]
Neither companion file was modified.

YM V10 proves that an ownership label does not create a linear map and that
ordinary operator boundedness does not imply a heat-weighted similarity
bound.  This confirms the V89 proof order: the metric normalizer must cancel
the actual principal symbol before Schatten estimates are used; calling a
term a gravitational owner cannot substitute for a symbol calculation.
NS V31 shows that removing a cutoff mark restores its exact critical
moments and that the same interval cannot pay both persistent and
recent-entry stocks.  Here the restored moments are the quartic, quadratic,
and logarithmic heat coefficients.  They are charged to local couplings
once; the normalized Fredholm remainder is not charged again.

\subsection{Positive operator and heat convention}

Let \(\mathcal D\) be the complete Euclidean SM fermion operator after the
V89 principal-symbol normalization, and put
\[
 \mathcal P:=\mathcal D^\dagger\mathcal D.
\]
The modulus of the fermion determinant is represented by
\(\frac12\log\det\mathcal P\); its anomaly-normalized phase remains in the
V74 determinant line and is not recomputed here.

Write the positive Laplace-type operator in the convention
\begin{equation}
 \mathcal P
 =-\left(g^{\mu\nu}\nabla_\mu\nabla_\nu+\mathcal E\right),
 \qquad
 \Omega_{\mu\nu}:=[\nabla_\mu,\nabla_\nu].
\label{eq:v4v90-laplace}
\end{equation}
On a closed four-manifold,
\[
 \operatorname{Tr}(e^{-t\mathcal P})
 \sim(4\pi t)^{-2}
 \left(A_0+tA_2+t^2A_4+O(t^3)\right).
\]

\begin{theorem}[Universal four-dimensional coefficients]
\label{thm:v4v90-heat}
With the convention \eqref{eq:v4v90-laplace},
\begin{align}
 A_0&=\int_X\operatorname{tr}(I)\,dV,
\label{eq:v4v90-A0}\\
 A_2&=\int_X\operatorname{tr}
       \left(\mathcal E+\frac16RI\right)dV,
\label{eq:v4v90-A2}\\
 A_4&=\frac1{360}\int_X\operatorname{tr}\Bigl(
  60R\mathcal E+180\mathcal E^2+30\Omega_{\mu\nu}\Omega^{\mu\nu}
\nonumber\\
 &\hspace{7em}
 +(5R^2-2|{\rm Ric}|^2+2|{\rm Riem}|^2)I
 +60\Delta\mathcal E+12\Delta R\,I
 \Bigr)dV.
\label{eq:v4v90-A4}
\end{align}
For a smooth boundary with a local strongly elliptic boundary condition,
the same interior coefficients occur together with local boundary
coefficients.  For an APS condition, the interior coefficients remain the
same but its spectral boundary and eta data must be retained in the
separate gluing and phase owner.
\end{theorem}
\begin{proof}
Choose normal coordinates and a synchronous bundle frame at a point.
The heat equation recursively determines the diagonal parametrix
\[
 (4\pi t)^{-2}e^{-d(x,y)^2/4t}
 \sum_{k\ge0}t^ku_k(x,y).
\]
The transport equation gives \(u_0(x,x)=I\).  Evaluating the next two
transport equations on the diagonal yields
\(u_1=\mathcal E+RI/6\) and the displayed \(u_2\); commutators of
covariant derivatives produce \(\Omega_{\mu\nu}\Omega^{\mu\nu}\), while
normal-coordinate contractions give the three quadratic curvature
invariants.  Integrating the fibre trace gives
\eqref{eq:v4v90-A0}--\eqref{eq:v4v90-A4}.  On a boundary, the reflected
parametrix produces additional boundary integrals but does not alter the
interior recursion.
\end{proof}

\subsection{Exact cutoff moments}

Define the proper-time cutoff functional
\begin{equation}
 \Gamma_\Lambda(\mathcal P)
 :=
 -\frac12\int_{\Lambda^{-2}}^\infty
 \frac{dt}{t}\operatorname{Tr}(e^{-t\mathcal P}),
\label{eq:v4v90-propertime}
\end{equation}
with an infrared subtraction when \(\mathcal P\) has zero modes; those
finite zero modes remain in the V74/V84 determinant line.

\begin{corollary}[Quartic, quadratic, and logarithmic owners]
\label{cor:v4v90-div}
On a closed core, the ultraviolet part of \eqref{eq:v4v90-propertime} is
\begin{equation}
 \Gamma_\Lambda^{\rm div}
 =
 -\frac1{2(4\pi)^2}\left[
   \frac{\Lambda^4}{2}A_0
   +\Lambda^2A_2
   +\log\!\left(\frac{\Lambda^2}{\mu^2}\right)A_4
 \right],
\label{eq:v4v90-div}
\end{equation}
up to the sign chosen for the effective action.  No probability averaging
of the cutoff changes these three moments.
\end{corollary}
\begin{proof}
Insert the first three heat terms into
\eqref{eq:v4v90-propertime}.  The elementary integrals are
\[
 \int_{\Lambda^{-2}}t^{-3}dt=\frac12\Lambda^4+O(1),\quad
 \int_{\Lambda^{-2}}t^{-2}dt=\Lambda^2+O(1),\quad
 \int_{\Lambda^{-2}}t^{-1}dt=\log(\Lambda^2/\mu^2)+O(1).
\]
Multiplication by \(-1/[2(4\pi)^2]\) gives
\eqref{eq:v4v90-div}.  Averaging a lower limit merely averages these
explicit functions and cannot erase their moments.
\end{proof}

\begin{proposition}[Boundary cutoff moments]
\label{prop:v4v90-boundary}
For a local strongly elliptic boundary condition, write
\[
 \operatorname{Tr}(e^{-t\mathcal P})
 \sim(4\pi t)^{-2}\sum_{j=0}^4t^{j/2}B_j+O(t^{1/2}),
\]
where \(B_0,B_2,B_4\) include the preceding interior terms and every
\(B_j\) contains its boundary contribution.  Then the divergent bracket in
\eqref{eq:v4v90-div} is replaced by
\[
 \frac{\Lambda^4}{2}B_0+\frac{2\Lambda^3}{3}B_1+
 \Lambda^2B_2+2\Lambda B_3+
 \log(\Lambda^2/\mu^2)B_4.
\]
In particular, artificial core boundaries cannot be ignored before their
interface determinants and opposite-side boundary owners have been glued.
\end{proposition}
\begin{proof}
The \(j\)-th term has proper-time power \(t^{(j-6)/2}\).  Integration from
\(\Lambda^{-2}\) gives respectively
\(\frac12\Lambda^4,\frac23\Lambda^3,\Lambda^2,2\Lambda\), and
\(\log(\Lambda^2/\mu^2)\).  The last statement follows because an
uncancelled \(B_1\) or \(B_3\) is a genuine cubic or linear cutoff
divergence.
\end{proof}

\subsection{Actual SM--Einstein coupling owners}

For a Dirac operator with gauge connection and Yukawa endomorphism
\(\Psi(\Phi)\), the Lichnerowicz formula puts into \(\mathcal E\) the
scalar curvature, gauge curvature, \(\nabla\Psi\), and \(\Psi^2\), while
\(\Omega\) contains spin and gauge curvature.  Therefore
Theorem~\ref{thm:v4v90-heat} produces the following genuine local
operators:
\begin{center}
\begin{tabular}{@{}p{0.18\linewidth}p{0.72\linewidth}@{}}
\toprule
heat owner & local coupling owner\\
\midrule
\(A_0\) &
\(\int\sqrt g\), the cosmological term;\\
\(A_2\) &
\(\int\sqrt g\,R\) and the Higgs quadratic term, with representation and
Yukawa traces;\\
\(A_4\) &
gauge kinetic terms, Higgs kinetic and quartic terms,
\(R|\Phi|^2\), and the curvature-squared invariants
\(R^2,|{\rm Ric}|^2,|{\rm Riem}|^2\), modulo total derivatives and the
four-dimensional Euler relation;\\
boundary heat terms &
extrinsic-curvature and induced boundary invariants required by the chosen
boundary condition;\\
determinant phase &
the already anomaly-normalized Pontryagin, eta, and global holonomy owner.
\\
\bottomrule
\end{tabular}
\end{center}

\begin{proposition}[No fictitious counterterm owner]
\label{prop:v4v90-no-label}
A local invariant may be charged to a counterterm only if it occurs in
\eqref{eq:v4v90-A0}--\eqref{eq:v4v90-A4}, in a boundary heat coefficient,
or in the separately normalized phase.  Adding an ownership label without
one of these operator sources does not alter
\(\Gamma_\Lambda^{\rm div}\).
\end{proposition}
\begin{proof}
The divergent functional is the explicit linear combination
\eqref{eq:v4v90-div}.  Its terms are fixed by the diagonal heat recursion.
A bookkeeping label absent from \(\mathcal E\), \(\Omega\), the Riemannian
curvature, and the boundary symbol changes none of those coefficients.
This is the direct heat-kernel instance of the YM V10 role-erasure lemma.
\end{proof}

\begin{theorem}[Pure Einstein--Hilbert space is not closed]
\label{thm:v4v90-EH-notclosed}
For the Standard Model matter content on a general curved background, the
one-loop parity-even fermion determinant produces nonzero
curvature-squared terms in \(A_4\).  Consequently a bare gravitational
action containing only cosmological and Einstein--Hilbert couplings is not
closed under this renormalization step.
\end{theorem}
\begin{proof}
Even with gauge and Higgs fields set to zero, the spin connection has
curvature
\[
 \Omega_{\mu\nu}^{\rm spin}
 =\frac14R_{\mu\nu ab}\gamma^{ab}.
\]
Its square contributes a nonzero contraction of the Riemann tensor through
the \(30\Omega^2\) term of \eqref{eq:v4v90-A4}.  The identity term in the
same formula independently contains
\(5R^2-2|{\rm Ric}|^2+2|{\rm Riem}|^2\).  Taking the nonzero spinor and
generation trace leaves curvature-squared invariants which cannot be
expressed, on general metrics, as a linear combination of
\(\int\sqrt g\) and \(\int\sqrt g\,R\).  Thus corresponding couplings are
required.
\end{proof}

\subsection{One-use subtraction and finite remainder}

Let
\[
 \mathsf{CT}_\Lambda
 :=
 +\frac1{2(4\pi)^2}\left[
   \frac{\Lambda^4}{2}A_0+\Lambda^2A_2+
   \log(\Lambda^2/\mu^2)A_4\right]
\]
with the corresponding boundary terms.  Define
\[
 \Gamma_{\Lambda}^{\rm sub}
 :=\Gamma_\Lambda+\mathsf{CT}_\Lambda.
\]

\begin{lemma}[The heat stock is spent once]
\label{lem:v4v90-once}
After this subtraction, the normalized weak-\(\mathcal S_4\) Fredholm
factor of V89 may contribute only its finite regularized determinant and
finite local scheme change.  Re-inserting \(A_0,A_2,A_4\) as bounds on that
finite factor would double count the ultraviolet heat stock.
\end{lemma}
\begin{proof}
Split the proper-time integral at a fixed \(t_0\).  Subtract the three
displayed asymptotic terms on \((0,t_0)\).  The remainder is
\(O(t)\) after multiplication by \(dt/t\) and is integrable at zero; the
integral on \((t_0,\infty)\) is ultraviolet finite.  Hence every divergent
use of \(A_0,A_2,A_4\) has already been exhausted by
\(\mathsf{CT}_\Lambda\).
\end{proof}

\subsection{Reference-change cocycle}

Let \(r,s\) be two frozen metric references with normalizers
\(Q_{g,r},Q_{g,s}\).  Each gives a local determinant-line factorization and
a subtraction functional.  Denote the corresponding renormalized line
sections by \(\sigma_r^{\rm ren}\) and \(\sigma_s^{\rm ren}\).

\begin{theorem}[Local reference-change transgression]
\label{thm:v4v90-reference}
On an overlap there is a local covariant finite functional
\(\mathcal T_{rs}(g,A,\Phi)\), built from the same weight-at-most-four
invariants and boundary owners, such that
\begin{equation}
 \sigma_s^{\rm ren}
 =
 e^{\mathcal T_{rs}}\,
 \Theta_{rs}\sigma_r^{\rm ren}.
\label{eq:v4v90-reference}
\end{equation}
The transition is nowhere zero and
\begin{equation}
 \mathcal T_{rt}=\mathcal T_{rs}+\mathcal T_{st},
 \qquad
 \Theta_{rt}=\Theta_{st}\Theta_{rs},
\label{eq:v4v90-refcocycle}
\end{equation}
after one common normalization at a base reference.
\end{theorem}
\begin{proof}
Both constructions represent the same physical determinant line.  Their
ultraviolet subtractions differ by the integral of the difference of their
heat coefficients and by the order-zero multiplicative anomaly of V89.
Theorem~\ref{thm:v4v89-localdefect} makes this difference a finite local
functional of weight at most four; call it \(\mathcal T_{rs}\).
Exponentiation is nowhere zero.  The line map \(\Theta_{rs}\) is the V87
reference-change isomorphism.  On a triple overlap, subtraction differences
telescope and complete-line isomorphisms compose, proving
\eqref{eq:v4v90-refcocycle}.  A common base normalization removes the
additive constant ambiguity.
\end{proof}

The local functional \(\mathcal T_{rs}\) is a finite scheme change.  The
renormalized couplings must transform by its coefficients.  Equation
\eqref{eq:v4v90-reference} then makes the product of the fermion section and
the running local action independent of the frozen reference.

\subsection{What V90 closes and what remains}

V90 proves the explicit ultraviolet owner list, the cutoff powers, the
necessity of curvature-squared gravitational couplings, one-use
subtraction, and the reference-change cocycle.  It does not prove
uniform convergence of \(\Gamma_\Lambda^{\rm sub}\) on the
bounded-action configuration family.  Local heat asymptotics alone are not
a domination theorem for the finite proper-time remainder.

V91 will attack that next gate by splitting proper time into ultraviolet,
mesoscopic, and infrared ranges.  The ultraviolet range will use a
quantified heat remainder after \(A_4\); the mesoscopic range will use the
V89 normalized resolvent/Schatten card; and the infrared range will retain
the finite soft zero-mode sector rather than divide by a spectral gap that
may vanish.

\clearpage
\section{V91: exact low-Born/\(\det_5\) separation and the infrared soft block}
\label{sec:v4v91}

V90 proposed controlling the finite proper-time remainder uniformly after
subtracting \(A_0,A_2,A_4\).  Bounded action does not by itself supply the
higher smooth-geometry bounds needed for such a uniform heat remainder.
There is a cleaner operator decomposition.  The weak-\(\mathcal S_4\)
relative perturbation lies in \(\mathcal S_5\); its fifth regularized
determinant controls the complete high-Born tail.  The first four Born
amplitudes remain an explicit renormalized finite packet.  Their divergent
parts are the local heat owners of V90, but their finite nonlocal parts must
not be discarded.

\subsection{Weak-\(\mathcal S_4\) gives a fifth determinant}

\begin{lemma}[Critical-to-fifth embedding]
\label{lem:v4v91-S5}
If \(K\in\mathcal S_{4,\infty}\), then \(K\in\mathcal S_5\) and
\begin{equation}
 \|K\|_5^5
 \le\zeta(5/4)\|K\|_{4,\infty}^5.
\label{eq:v4v91-S5}
\end{equation}
\end{lemma}
\begin{proof}
By definition,
\(s_n(K)\le\|K\|_{4,\infty}n^{-1/4}\).  Therefore
\[
 \sum_{n\ge1}s_n(K)^5
 \le\|K\|_{4,\infty}^5\sum_{n\ge1}n^{-5/4},
\]
and the last series is \(\zeta(5/4)<\infty\).
\end{proof}

For \(K\in\mathcal S_5\), define
\begin{equation}
 \det{}_5(1+K)
 :=
 \det\!\left[
 (1+K)\exp\!\left(-K+\frac{K^2}{2}
                  -\frac{K^3}{3}+\frac{K^4}{4}\right)
 \right].
\label{eq:v4v91-det5}
\end{equation}
The operator inside the ordinary determinant differs from the identity by
a trace-class operator.  The resulting function is entire on
\(\mathcal S_5\), including at points where \(1+K\) has kernel.

\begin{theorem}[Uniform high-Born carrier]
\label{thm:v4v91-tail}
If \(\|K\|_{4,\infty}\le M\), then
\begin{equation}
 |\det{}_5(1+K)|
 \le\exp(C_5\zeta(5/4)M^5).
\label{eq:v4v91-det5-bound}
\end{equation}
On a set with two marked weak-\(\mathcal S_4\) bounds, the first two
Frechet derivatives of \(\det_5(1+K)\) are uniformly bounded multilinear
forms in the corresponding \(\mathcal S_5\) marks, even at determinant
zeros.
\end{theorem}
\begin{proof}
The canonical-product estimate is
\(|\det_5(1+K)|\le\exp(C_5\|K\|_5^5)\).
Insert Lemma~\ref{lem:v4v91-S5}.  Entire functions on Banach spaces are
locally bounded with locally bounded Frechet derivatives; the Cauchy
formula on one- and two-dimensional complex disks gives the marked bounds.
This is the V81 determinant argument with the explicit weak-to-strong
constant retained.
\end{proof}

\subsection{Why fifth order is sharp at the critical endpoint}

\begin{proposition}[Logarithmic fourth-order obstruction]
\label{prop:v4v91-sharp}
Weak-\(\mathcal S_4\) control alone does not bound
\(\operatorname{Tr}|K|^4\).  Indeed, there are compact diagonal operators
with
\[
 s_n(K)=M n^{-1/4}
\]
for which
\[
 \sum_{n\le N}s_n(K)^4
 =M^4\sum_{n\le N}\frac1n
 =M^4\log N+O(M^4),
\]
while \(\|K\|_5^5=M^5\zeta(5/4)\).
\end{proposition}
\begin{proof}
The displayed diagonal sequence defines a compact operator with weak
\(\mathcal S_4\) norm \(M\).  The two series are respectively harmonic
and \(p\)-summable with exponent \(5/4\).
\end{proof}

Thus a fourth determinant would leave an uncontrolled logarithmic trace.
The logarithm is exactly the critical \(A_4\) owner of V90.  Passing to
\(\det_5\) is forced by dimension and operator order, not by a convenient
choice of regulator.

\subsection{Exact finite-order packet}

For trace-class \(K\), the defining identity gives
\begin{equation}
 \log\det(1+K)
 =
 \log\det{}_5(1+K)
 +\sum_{j=1}^4\frac{(-1)^{j+1}}j\operatorname{Tr}K^j.
\label{eq:v4v91-traceclass}
\end{equation}
For the critical relative operators of V89, the individual traces on the
right need renormalization.

\begin{definition}[Renormalized low-Born packet]
\label{def:v4v91-lowborn}
Let \(\mathcal B_j^{\rm ren}(K)\), \(1\le j\le4\), be the renormalized
connected \(j\)-insertion amplitude obtained from
\(\operatorname{Tr}K^j\) by subtracting the local heat/symbol terms in the
V90 coupling owners, with one fixed finite normalization at the reference
background.  Put
\[
 \mathcal B_{\le4}^{\rm ren}(K)
 :=
 \sum_{j=1}^4\frac{(-1)^{j+1}}j\mathcal B_j^{\rm ren}(K).
\]
\end{definition}

\begin{firewall}[Finite low-order amplitudes are not all local]
\label{fw:v4v91-nonlocal}
The divergent parts of \(\mathcal B_j^{\rm ren}\) are local, but their
finite parts generally contain nonlocal vacuum-polarization,
stress-tensor, and vertex form factors.  V90 authorizes local counterterms;
it does not authorize deletion of those finite amplitudes.
\end{firewall}

\begin{theorem}[Renormalized determinant decomposition]
\label{thm:v4v91-decomp}
On a frozen reference chart, the renormalized complete determinant-line
section has the form
\begin{equation}
 \sigma^{\rm ren}(D_\circ+L)
 =
 \exp\!\bigl(\mathcal B_{\le4}^{\rm ren}(K)
             +\mathcal T_{\rm loc}\bigr)
 \det{}_5(1+K)\,\sigma_\circ,
\label{eq:v4v91-decomp}
\end{equation}
away from notation-dependent separation of the anomaly-normalized phase.
Here \(K=LC_\circ\), \(\mathcal T_{\rm loc}\) is the finite local
normalizer/multiplicative-anomaly functional of V89--V90, and
\(\sigma_\circ\) is the reference line section.  Equation
\eqref{eq:v4v91-decomp} extends across zeros as an equality of determinant
line sections.
\end{theorem}
\begin{proof}
For a smoothing approximation \(K_\epsilon\), use the trace-class identity
\eqref{eq:v4v91-traceclass}.  Subtract the first four local singular
symbol coefficients and impose the fixed finite normalization; by
definition their limits are \(\mathcal B_j^{\rm ren}\).
Lemma~\ref{lem:v4v91-S5} gives
\(K_\epsilon\to K\) in \(\mathcal S_5\) after choosing the standard
pseudodifferential approximation, so the entire fifth determinant
converges.  The order-zero metric normalizer changes the product by the
local functional \(\mathcal T_{\rm loc}\) from
Theorem~\ref{thm:v4v90-reference}.  The equality first holds where the
operator is invertible.  Both sides are holomorphic sections and the
\(\det_5\) factor retains the same zero divisor, so it extends across the
Fredholm zero set.
\end{proof}

The smoothing argument does not prove uniformity of the four finite
amplitudes; it identifies exactly which finite packet still needs a
continuum estimate.

\subsection{No double determinant at a crossing}

The finite Sylvester matrix \(M\) of V84 and \(\det_5(1+K)\) describe the
same zeros in different coordinates.  They are not independent factors.

\begin{proposition}[Finite matrix is a line coordinate, not a second zero]
\label{prop:v4v91-no-double}
On a collared chart, the zero order and algebraic multiplicity of
\(\det_5(1+K)\) equal those of the finite crossing determinant after
including the nowhere-zero tail and regularized multiplicative-anomaly
factors.  Multiplying \(\det_5(1+K)\) by \(\det M\) would count every hard
crossing twice.
\end{proposition}
\begin{proof}
The V84 factorization is
\[
 1+K=T(1+AB),\qquad \det_F(1+AB)=\det_E(1+BA).
\]
The tail \(T\) is invertible.  Regularized determinant multiplicativity
differs from ordinary multiplicativity by an exponential of trace
polynomials, which is nowhere zero.  Hence the only zero is the finite
Sylvester zero, with the same algebraic multiplicity.  A product of both
zero carriers squares that divisor.
\end{proof}

\subsection{Infrared range and soft zero modes}

Choose a fixed proper-time splitting
\[
 (0,\infty)=(0,t_{\rm uv})\cup
 [t_{\rm uv},t_{\rm ir}]\cup(t_{\rm ir},\infty).
\]
The ultraviolet singularities are the V90 heat owners and the finite
low-Born packet.  The mesoscopic/high-Born dependence is contained in
\(\det_5\).  For the infrared range, decompose with the collared finite soft
projection \(P_{\rm s}\).

\begin{lemma}[Infrared hard tail]
\label{lem:v4v91-IR}
If the hard spectrum of \(\mathcal P=D^\dagger D\) lies in
\([\delta_{\rm h}^2,\infty)\), then
\[
 \int_{t_{\rm ir}}^\infty\frac{dt}{t}
 \operatorname{Tr}\bigl((1-P_{\rm s})e^{-t\mathcal P}\bigr)
 \le
 \frac{e^{-\frac12t_{\rm ir}\delta_{\rm h}^2}}
      {t_{\rm ir}\delta_{\rm h}^2}
 \operatorname{Tr}\bigl((1-P_{\rm s})
 e^{-\frac12t_{\rm ir}\mathcal P}\bigr)
\]
after increasing the fixed intermediate heat time by a harmless factor.
The soft contribution is the finite sum
\[
 \sum_{\lambda_k\in{\rm soft}}
 \int_{t_{\rm ir}}^\infty e^{-t|\lambda_k|^2}\frac{dt}{t},
\]
retained as a finite determinant-line factor.
\end{lemma}
\begin{proof}
For \(x=t_{\rm ir}\lambda^2\),
\[
 \int_{t_{\rm ir}}^\infty e^{-t\lambda^2}\frac{dt}{t}
 =\int_x^\infty e^{-u}\frac{du}{u}
 \le\frac{e^{-x}}x.
\]
Since \(\lambda^2\ge\delta_{\rm h}^2\), split
\[
 e^{-t_{\rm ir}\lambda^2}
 \le e^{-\frac12t_{\rm ir}\delta_{\rm h}^2}
      e^{-\frac12t_{\rm ir}\lambda^2}
\]
and sum over the hard spectrum.  Finite rank permits direct summation on
the soft space.  No inverse soft gap is taken.
\end{proof}

At \(\lambda_k=0\), the soft integral diverges logarithmically, reflecting
the physical determinant zero.  It must be saturated by the Grassmann
insertions and selection rules of V75--V77, not bounded by an inverse
spectral gap.

\subsection{Uniform closure criterion}

\begin{theorem}[Determinant domination reduced to a finite packet]
\label{thm:v4v91-criterion}
On a bounded-action collared sector, assume:
\begin{enumerate}
 \item the combined normalized gauge--Higgs--metric relative operator has
       uniform weak-\(\mathcal S_4\) norm and two marked bounds;
 \item \(\mathcal B_j^{\rm ren}\), \(1\le j\le4\), have uniform local
       correlator bounds through two marks after the V90 subtractions; and
 \item the finite soft sector is treated by its exact zero-mode saturation.
\end{enumerate}
Then the complete renormalized fermion determinant section and its first
two marked coefficients are uniformly dominated on that sector, including
at hard determinant zeros.
\end{theorem}
\begin{proof}
Item 1 and Theorem~\ref{thm:v4v91-tail} control the entire high-Born
factor and its marks.  Item 2 controls the exponential finite packet and
the local reference transgression.  Item 3 controls the only infrared
terms not covered by Lemma~\ref{lem:v4v91-IR}.  The product rule gives the
two marked bound.  Entirety of \(\det_5\) and the determinant-line
description avoid division by a zero.
\end{proof}

V91 replaces an unproved uniform heat remainder by an exact operator
factorization.  The remaining fermionic ultraviolet gate is finite:
uniform renormalized one-, two-, three-, and four-insertion amplitudes in
the joint metric--gauge--Higgs background.  Weak-\(\mathcal S_4\) control
alone cannot prove these bounds, as
Proposition~\ref{prop:v4v91-sharp} shows.  V92 will start with the marginal
four-insertion amplitude, identify its Wodzicki-residue local part, and
test the finite nonlocal remainder against concentration scaling.

\clearpage
\section{V92: the marginal four-insertion residue and scale-matched finite part}
\label{sec:v4v92}

The four-insertion term is the critical member of the finite packet in V91.
For a smooth normalized relative operator \(K\in\Psi^{-1}\) on a closed
four-manifold, \(K^4\in\Psi^{-4}\).  Its logarithmic divergence is the
Wodzicki residue, hence local.  Its finite weighted trace is generally
nonlocal and depends on a renormalization scale.  The distinction is
essential at concentration scales.

The weighted-trace facts used here are the standard pseudodifferential
residue formulas; see, for example, the weighted trace-cochain construction
of Paycha, arXiv:\texttt{math-ph/0503033}, and the classical residue/zeta
relation.

\subsection{Residue and weighted trace}

Let \(Q\) be a positive invertible classical elliptic operator of order one.
For \(P\in\Psi^{-4}\), define initially for \(\Re s>0\)
\[
 \zeta_{P,Q}(s):=\operatorname{Tr}(PQ^{-s}).
\]

\begin{theorem}[Critical zeta pole]
\label{thm:v4v92-pole}
The function \(\zeta_{P,Q}\) extends meromorphically across \(s=0\) with at
most a simple pole,
\begin{equation}
 \zeta_{P,Q}(s)
 =
 \frac{\operatorname{res}(P)}s+
 \operatorname{tr}^{Q}(P)+O(s),
\label{eq:v4v92-laurent}
\end{equation}
where
\begin{equation}
 \operatorname{res}(P)
 =
 \frac1{(2\pi)^4}
 \int_X\int_{|\xi|=1}
 \operatorname{tr}\sigma_{-4}(P)(x,\xi)\,dS(\xi)\,dV(x).
\label{eq:v4v92-wres}
\end{equation}
The residue is independent of \(Q\), is local, and vanishes on
commutators of classical pseudodifferential operators.
\end{theorem}
\begin{proof}
In a coordinate chart, the trace of \(PQ^{-s}\) is the integral of its
complete symbol over \((x,\xi)\).  Radial integration of a homogeneous
symbol of degree \(-4-s\) gives
\[
 \int_1^\infty r^{-1-s}dr=\frac1s.
\]
All other homogeneous degrees are regular at \(s=0\), and smoothing
remainders are trace class.  Summing a partition of unity yields
\eqref{eq:v4v92-laurent}; the coefficient of the pole is precisely
\eqref{eq:v4v92-wres}.  A change of coordinates alters the degree-\(-4\)
density by a divergence, so the integral is intrinsic.  The degree-\(-4\)
symbol of a commutator is a total \((x,\xi)\)-divergence; integration over
the closed cosphere bundle gives zero.
\end{proof}

For a renormalization scale \(\mu>0\), set
\[
 \operatorname{tr}^{Q,\mu}(P)
 :=
 \operatorname*{FP}_{s=0}
 \operatorname{Tr}\!\left(P(Q/\mu)^{-s}\right).
\]

\begin{corollary}[Exact scale change]
\label{cor:v4v92-scale}
For \(\mu_1,\mu_2>0\),
\begin{equation}
 \operatorname{tr}^{Q,\mu_2}(P)
 -
 \operatorname{tr}^{Q,\mu_1}(P)
 =
 \log(\mu_2/\mu_1)\operatorname{res}(P).
\label{eq:v4v92-scale}
\end{equation}
\end{corollary}
\begin{proof}
Since \((Q/\mu)^{-s}=\mu^sQ^{-s}\), multiply
\eqref{eq:v4v92-laurent} by
\(\mu^s=1+s\log\mu+O(s^2)\) and compare finite parts.
\end{proof}

\subsection{The four-point logarithm is local}

Apply the theorem to \(P=K^4\).

\begin{theorem}[Marginal four-insertion owner]
\label{thm:v4v92-four}
For a classical normalized relative operator \(K\in\Psi^{-1}\),
\begin{equation}
 \mathcal R_4(K):=\operatorname{res}(K^4)
 =
 \frac1{(2\pi)^4}\int_{S^*X}
 \operatorname{tr}\!\left(\sigma_{-1}(K)^4\right)dS\,dV.
\label{eq:v4v92-R4}
\end{equation}
It is a local gauge- and diffeomorphism-covariant quartic functional of the
principal relative symbol.  The renormalized fourth Born amplitude changes
with scale by
\begin{equation}
 \mathcal B_4^{\rm ren}(\mu_2)
 -\mathcal B_4^{\rm ren}(\mu_1)
 =
 \log(\mu_2/\mu_1)\mathcal R_4(K),
\label{eq:v4v92-B4run}
\end{equation}
up to the fixed coefficient \(1/4\) from the logarithmic determinant
series.
\end{theorem}
\begin{proof}
The principal symbol of a product is the product of principal symbols, so
\(\sigma_{-4}(K^4)=\sigma_{-1}(K)^4\).  Equation
\eqref{eq:v4v92-R4} follows from
Theorem~\ref{thm:v4v92-pole}.  Gauge and spin lifts conjugate the symbol,
and the fibre trace is invariant; diffeomorphisms preserve the cosphere
density.  Equation \eqref{eq:v4v92-B4run} is
Corollary~\ref{cor:v4v92-scale}.
\end{proof}

For \(K=M_WC_\circ\), the symbol is
\[
 \sigma_{-1}(K)(x,\xi)
 =W(x)\sigma_1(D_\circ)(x,\xi)^{-1}.
\]
Thus \(\mathcal R_4\) is an actual local quartic contraction of the
gauge--Higgs--normalized-metric insertion.  It is the symbol-level source of
the logarithmic \(A_4\) couplings in V90, not a separate fifth owner.

\begin{corollary}[Marked residue formulas]
\label{cor:v4v92-marks}
For classical marks \(L_1,L_2\in\Psi^{-1}\),
\begin{align}
 D\mathcal R_4(K)[L_1]
 &=\operatorname{res}\!\left(
   L_1K^3+KL_1K^2+K^2L_1K+K^3L_1\right),
\label{eq:v4v92-D1}\\
 D^2\mathcal R_4(K)[L_1,L_2]
 &=\operatorname{res}\!\left(
   \sum_{\substack{\text{ordered words of length }4\\
                    \text{with two }K\text{ and }L_1,L_2}}
   \text{word}\right).
\label{eq:v4v92-D2}
\end{align}
These are local and, by the trace property of the residue, may be combined
cyclically without changing their value.
\end{corollary}
\begin{proof}
Differentiate the noncommutative polynomial \(K^4\) once and twice, retain
the order of factors, and apply linearity of the residue.
\end{proof}

\subsection{Scale-matched concentration card}

Let a core of radius \(\rho\) be rescaled to unit size by a unitary
\(U_\rho\).  Suppose the normalized critical relative operator becomes a
fixed unit-scale operator \(K_1\), while the physical weight obeys
\[
 Q_\rho=\rho^{-1}U_\rho Q_1U_\rho^{-1}.
\]

\begin{proposition}[All universal concentration growth is the residue log]
\label{prop:v4v92-concentration}
One has
\begin{equation}
 \operatorname{tr}^{Q_\rho}(K_\rho^4)
 =
 \operatorname{tr}^{Q_1}(K_1^4)
 +\log\rho\,\mathcal R_4(K_1).
\label{eq:v4v92-concentration}
\end{equation}
Equivalently, choosing the running scale \(\mu_\rho=\rho^{-1}\) makes the
dimensionless weighted trace equal to its unit-scale value.
\end{proposition}
\begin{proof}
Unitary invariance reduces the zeta function to
\[
 \operatorname{Tr}(K_1^4(\rho^{-1}Q_1)^{-s})
 =\rho^s\operatorname{Tr}(K_1^4Q_1^{-s}).
\]
Comparing finite parts at \(s=0\) gives
\eqref{eq:v4v92-concentration}.  Dividing \(Q_\rho\) by
\(\mu_\rho=\rho^{-1}\) removes the factor \(\rho^{-1}\).
\end{proof}

A fixed-scale coupling therefore acquires the physical
\(\log\rho\) running.  The log is not an uncontrolled nonlocal remainder;
it is exactly the local residue.  Whether its exponential suppresses or
enhances small cores depends on the already normalized beta-function
coefficient and cannot be decided from an absolute-value estimate alone.

\subsection{Weak Schatten control does not define the finite part}

\begin{proposition}[No finite fourth trace from weak norm alone]
\label{prop:v4v92-noFP}
There is no procedure which subtracts only
\(c\log N\) from every weak-\(\mathcal S_4\) diagonal operator and obtains
a uniformly bounded fourth partial trace from its weak norm.  For example,
for \(n\ge2\), let
\[
 s_n(K)^4
 =
 \frac{M^4}{n}\left(1+\frac1{\log(n+1)}\right).
\]
Then \(\|K\|_{4,\infty}\le2^{1/4}M\), but
\[
 \sum_{n=2}^Ns_n(K)^4-M^4\log N
 =
 M^4\log\log N+O(M^4),
\]
which diverges.
\end{proposition}
\begin{proof}
The weak bound follows from
\(n^{1/4}s_n(K)\le2^{1/4}M\).  Integral comparison gives
\[
 \sum_{n=2}^N\frac1{n\log(n+1)}
 =\log\log N+O(1),
\]
while the first summand is the harmonic logarithm.
\end{proof}

Classical symbol structure rules out this extra \(\log\log\) behavior and
selects the residue coefficient.  A rough \(L^4\) Cwikel estimate alone
does not.  Consequently the V91 finite fourth packet needs either a
uniform classical-symbol approximation or a direct renormalized
quadrilinear estimate on critical coefficient spaces.

\subsection{A sufficient finite-part topology}

\begin{theorem}[Symbol-bounded fourth packet]
\label{thm:v4v92-symbolbound}
Let \(\mathcal K\) be a family in \(\Psi^{-1}_{\rm cl}\) whose complete
symbols are bounded in sufficiently many \(S^{-1}_{1,0}\) seminorms on a
fixed bounded-geometry atlas, and whose smoothing remainders are uniformly
trace class.  Then, for fixed \(Q\) and \(\mu\),
\[
 \sup_{K\in\mathcal K}
 \left|\operatorname{tr}^{Q,\mu}(K^4)\right|<\infty.
\]
The same holds through two marked derivatives if the marked symbols obey
the corresponding seminorm bounds.
\end{theorem}
\begin{proof}
Parameter-dependent symbol composition bounds the complete symbols of
\(K^4\) in \(S^{-4}_{1,0}\).  Subtract its homogeneous degree-\(-4\)
cosphere integral in the radial symbol trace.  The remaining radial
integral is absolutely convergent and uniformly bounded by finitely many
symbol seminorms; the low-frequency compact region is bounded by the same
seminorms and the smoothing trace norm.  Differentiation once or twice
produces the ordered products in
\eqref{eq:v4v92-D1}--\eqref{eq:v4v92-D2}, to which the same argument
applies.
\end{proof}

The theorem is a sufficient smooth-background result, not an
action-only result.  Uhlenbeck \(W^{1,2}\) control gives an \(L^4\)
connection but does not uniformly bound all classical symbol seminorms.
V93 must bridge that exact gap.  It will localize the renormalized
quadrilinear form to the Uhlenbeck stopping cover and test whether subtracting
the local residue leaves a Coifman--Meyer type four-linear form bounded on
\(L^4\times L^4\times L^4\times L^4\).  Failure at the endpoint will be
recorded rather than replaced by the smooth symbol theorem.

\clearpage
\section{V93: fixed-scale endpoint failure and exact corewise RG matching}
\label{sec:v4v93}

V92 separates a local logarithmic coefficient from a nonlocal finite fourth
amplitude.  Under four-dimensional concentration the \(L^4\) insertion
norm is fixed, but the finite amplitude at a fixed subtraction scale changes
by the residue logarithm.  The correct replacement for a scale-free
endpoint estimate is evaluation of one running-coupling trajectory at each
terminal core scale.

\subsection{Sharp endpoint failure}

For a smooth compactly supported unit-scale insertion \(W\), put
\[
 W_\rho(x)=\rho^{-1}W(x/\rho).
\]
Then \(\|W_\rho\|_4=\|W\|_4\).  Let \(K_\rho\) be the corresponding
normalized order-\(-1\) relative operator.

\begin{theorem}[No fixed-scale \(L^4\) bound]
\label{thm:v4v93-noL4}
If \(\mathcal R_4(K_1)\ne0\), then at a fixed physical subtraction scale
\(\mu_0\),
\begin{equation}
 \mathcal B_4^{\rm ren}(K_\rho;\mu_0)
 =
 \mathcal B_4^{\rm ren}(K_1;1)
 +c_4\log(\rho\mu_0)\mathcal R_4(K_1)+O(1).
\label{eq:v4v93-logrho}
\end{equation}
Consequently no \(\rho\)-independent constant satisfies
\[
 |\mathcal B_4^{\rm ren}(K_\rho;\mu_0)|
 \le C\|W_\rho\|_4^4.
\]
Here \(c_4\ne0\) is the fourth coefficient in the logarithmic determinant
series.
\end{theorem}
\begin{proof}
Unitary dilation changes the physical order-one weight to
\(\rho^{-1}Q_1\).  Proposition~\ref{prop:v4v92-concentration}, multiplied
by the fixed fourth Born coefficient, gives the logarithm in
\eqref{eq:v4v93-logrho}.  Rescaled bounded-geometry errors stay bounded.
The left side is unbounded as \(\rho\downarrow0\), whereas the critical
\(L^4\) norm is constant.
\end{proof}

\begin{firewall}[A physical running log]
\label{fw:v4v93-log}
V90 already removes the ultraviolet cutoff logarithm.  The logarithm in
\eqref{eq:v4v93-logrho} compares the physical core scale with \(\mu_0\).
Subtracting it independently on each core would choose a different theory
on each core.
\end{firewall}

\subsection{Residue basis and running couplings}

Choose a linearly independent basis
\(\{\mathcal O_\alpha\}\) of the weight-four local operators present in the
V90 heat card, modulo total derivatives, geometric identities, and the
separate phase owner.  Expand
\begin{equation}
 c_4\mathcal R_4(K)
 =\sum_\alpha\beta_\alpha\mathcal O_\alpha.
\label{eq:v4v93-beta}
\end{equation}
The coefficients include all representation, generation, Yukawa, and spin
traces.

\begin{lemma}[Uniqueness of the beta vector]
\label{lem:v4v93-unique}
The coefficients in \eqref{eq:v4v93-beta} are unique.
\end{lemma}
\begin{proof}
The difference of two expansions is a vanishing linear combination of the
basis classes.  Linear independence forces every coefficient to vanish.
\end{proof}

Thus labels such as gauge, Higgs, and gravity do not replace an actual trace
calculation.  Let \(g_\alpha(\mu)\) multiply
\(\mathcal O_\alpha\) and impose
\begin{equation}
 \frac{d}{d\log\mu}g_\alpha(\mu)=-\beta_\alpha.
\label{eq:v4v93-RG}
\end{equation}

\begin{theorem}[Exact marginal RG cancellation]
\label{thm:v4v93-RG}
Wherever the flow exists,
\begin{equation}
 \mathfrak F_4(\mu)
 :=
 \mathcal B_4^{\rm ren}(K;\mu)
 +\sum_\alpha g_\alpha(\mu)\mathcal O_\alpha
\label{eq:v4v93-F}
\end{equation}
is independent of \(\mu\).
\end{theorem}
\begin{proof}
Corollary~\ref{cor:v4v92-scale} and
\eqref{eq:v4v93-beta} give
\[
 \frac d{d\log\mu}\mathcal B_4^{\rm ren}
 =\sum_\alpha\beta_\alpha\mathcal O_\alpha.
\]
The derivative of the local action is its negative by
\eqref{eq:v4v93-RG}.
\end{proof}

For coupling-dependent beta functions the same proof uses the full system
\(\dot g_\alpha=-\beta_\alpha(g)\), with the amplitude evaluated along the
flow.  Existence of that flow to arbitrarily high scale is a separate
ultraviolet-completion condition.

\subsection{Terminal-core scale matching}

Let the disjoint terminal stopping cores have radii \(\rho_K\), and set
\(\mu_K=\rho_K^{-1}\).

\begin{theorem}[Corewise evaluation of one theory]
\label{thm:v4v93-corewise}
Assume every running coupling is obtained from common data at a base scale
\(\mu_0\).  On terminal core \(K\), replacing
\[
 \mathcal B_4^{\rm ren}(\mu_0)
 +\sum_\alpha g_\alpha(\mu_0)\mathcal O_{\alpha,K}
\]
by
\[
 \mathcal B_4^{\rm ren}(\mu_K)
 +\sum_\alpha g_\alpha(\mu_K)\mathcal O_{\alpha,K}
\]
does not change the determinant-line section or local action.  It moves the
explicit concentration logarithm into the value of the same running
coupling.
\end{theorem}
\begin{proof}
Both expressions equal the scale-independent quantity
\(\mathfrak F_4\) from Theorem~\ref{thm:v4v93-RG}, restricted to the core.
The V90 reference-change cocycle transports the line section between frozen
references.
\end{proof}

\begin{firewall}[Terminal owner only]
\label{fw:v4v93-terminal}
Only disjoint terminal owners receive the local residue integral.  Ancestor
nodes localize the same physical region.  Charging every ancestor would
manufacture a logarithmic multiplicity from one density.
\end{firewall}

\subsection{Connected words across scales}

A connected four-insertion word meeting several terminal cores must be
composed before its residue is assigned.

\begin{definition}[Minimal connected owner]
\label{def:v4v93-owner}
The owner is the smallest stopping-tree node containing all four insertion
supports and all propagator segments in the ordered word.
\end{definition}

\begin{proposition}[Scale transfer of a connected word]
\label{prop:v4v93-crossscale}
Changing the weighted trace of the composed order-\(-4\) word from its
owner scale to another scale changes it by precisely its local residue
times the logarithmic scale ratio.  The remaining finite connected
amplitude is not a second local counterterm.
\end{proposition}
\begin{proof}
Apply \eqref{eq:v4v92-scale} after the four ordered factors have been
composed.  The residue is local by
Theorem~\ref{thm:v4v92-pole}; the finite part is the same weighted trace in
the two scale conventions.
\end{proof}

\subsection{The flow gate}

Corewise matching removes large logarithms from the finite coefficient but
does not prove that the RG trajectory reaches
\(\mu_K\) for arbitrarily small \(\rho_K\).  It also does not bound a
unit-scale finite amplitude on arbitrary rough \(L^4\) profiles; the V92
weak-Schatten counterexample remains.

\begin{theorem}[Conditional marginal closure]
\label{thm:v4v93-conditional}
The renormalized four-insertion packet is uniformly controlled over a
stopping family if:
\begin{enumerate}
 \item the full running-coupling trajectory remains in a compact admissible
 coupling set through every terminal scale;
 \item the scale-matched unit-core finite parts are uniformly bounded in
 the actual critical profile topology; and
 \item cross-core words use Definition~\ref{def:v4v93-owner}.
\end{enumerate}
\end{theorem}
\begin{proof}
Item 1 bounds local coupling coefficients.  Item 2 bounds the finite part
after Theorem~\ref{thm:v4v93-corewise} removes the concentration logarithm.
Item 3 and Proposition~\ref{prop:v4v93-crossscale} prevent residue
multiplicity.  Disjoint terminal residue integrals are paid by the V86
action ledger.
\end{proof}

V93 proves that a fixed-scale \(L^4\) endpoint estimate is impossible and
replaces it with exact RG covariance.  The next gate is the actual joint
SM--Einstein beta vector and its maximal ultraviolet flow interval.  V94
will compute the one-loop gauge coefficients from the three-generation
chiral representation, check their normalization and signs, and separate
asymptotically free, infrared-free, and gravitationally mixed owners.  No
ultraviolet-complete trajectory is assumed.

\clearpage
\section{V94: inherited hypercharge obstruction, Einstein two-loop growth, and the UV certificate}
\label{sec:v4v94}

V93 announced a new calculation of the one-loop gauge coefficients.  That
would duplicate Proposition~\ref{prop:v4v6-one-loop}.  The earlier result
already proves, in the convention
\[
 \mu\frac{dg_r}{d\mu}
 =-\frac{b_r}{16\pi^2}g_r^3+O(g_r^5),
\]
that
\[
 b_3=7,\qquad b_2=\frac{19}{6},\qquad b_Y=-\frac{41}{6},
\]
and Theorem~\ref{thm:v4v6-hypercharge-obstruction} gives the corresponding
conditional exact-shell no-go.  V94 imports those results and determines
their consequence for the terminal-core RG certificate of V93.  It also
adds the independent perturbative obstruction from quantized Einstein
gravity.

\subsection{Exact one-loop trajectories}

Ignore higher-order terms temporarily and set
\(z_r(\mu)=g_r(\mu)^{-2}\).

\begin{lemma}[Integrated one-loop gauge flow]
\label{lem:v4v94-flow}
For constant \(b_r\),
\begin{equation}
 z_r(\mu)
 =
 z_r(\mu_0)+\frac{b_r}{8\pi^2}\log(\mu/\mu_0).
\label{eq:v4v94-flow}
\end{equation}
Thus \(b_r>0\) gives \(g_r(\mu)\to0\) as \(\mu\to\infty\), whereas
\(b_r<0\) makes \(z_r\) vanish at a finite one-loop scale.
\end{lemma}
\begin{proof}
Differentiate \(g^{-2}\):
\[
 \frac d{d\log\mu}g^{-2}
 =-2g^{-3}\frac{dg}{d\log\mu}
 =\frac{b_r}{8\pi^2}.
\]
Integration gives \eqref{eq:v4v94-flow}.
\end{proof}

\begin{corollary}[One-loop hypercharge endpoint]
\label{cor:v4v94-Y}
For the inherited Standard Model coefficient,
\begin{equation}
 g_Y(\mu)^{-2}
 =
 g_Y(\mu_0)^{-2}
 -\frac{41}{48\pi^2}\log(\mu/\mu_0),
\label{eq:v4v94-Y}
\end{equation}
and its one-loop positive-coupling interval ends at
\begin{equation}
 \mu_Y^{\rm end}
 =
 \mu_0\exp\!\left(\frac{48\pi^2}{41g_Y(\mu_0)^2}\right).
\label{eq:v4v94-Yend}
\end{equation}
\end{corollary}
\begin{proof}
Insert \(b_Y=-41/6\) into \eqref{eq:v4v94-flow} and solve
\(g_Y(\mu)^{-2}=0\).
\end{proof}

Equation \eqref{eq:v4v94-Yend} is a one-loop endpoint, not a
nonperturbative theorem about every possible ultraviolet completion.
The earlier V6 theorem is stronger in a different direction: if the exact
shell increments remain uniformly dominated by their Gaussian-boundary
coefficient in the small-coupling wedge, positivity fails after finitely
many shells.  The only exits are to leave that wedge, change the ultraviolet
field content or gauge group, or reach a new interacting fixed point.

\begin{theorem}[Failure of the V93 all-scale premise for the unmodified
Gaussian boundary]
\label{thm:v4v94-V93fail}
Under the hypotheses of
Theorem~\ref{thm:v4v6-hypercharge-obstruction}, item \(1\) of
Theorem~\ref{thm:v4v93-conditional} cannot hold for terminal radii tending
to zero in the unmodified Standard Model hypercharge sector.
\end{theorem}
\begin{proof}
Terminal radii \(\rho_K\downarrow0\) require evaluation scales
\(\mu_K=\rho_K^{-1}\to\infty\).  The V6 exact-shell theorem says that the
positive inverse hypercharge coupling decreases by a fixed amount on every
ultraviolet shell while the dominance hypotheses hold.  It becomes
nonpositive after finitely many shells.  Therefore no positive trajectory
inside that admissible wedge reaches every \(\mu_K\).
\end{proof}

\subsection{Einstein power counting}

The V90 heat card concerns one-loop matter determinants in an external
metric.  Quantizing metric fluctuations adds a different tower.  Expand
the Einstein--Hilbert action around a smooth background.  Every graviton
propagator has momentum degree \(-2\), while every Einstein vertex has two
derivatives.

\begin{lemma}[Loop degree in Einstein gravity]
\label{lem:v4v94-power}
For a connected pure-gravity graph with \(L\) loops, \(I\) internal lines,
and \(V\) vertices, the superficial momentum degree is
\begin{equation}
 \omega=4L-2I+2V=2L+2.
\label{eq:v4v94-power}
\end{equation}
\end{lemma}
\begin{proof}
Four loop integrations contribute \(4L\), propagators contribute
\(-2I\), and vertices contribute \(2V\).  The connected identity
\(L=I-V+1\) gives the last equality.
\end{proof}

Thus \(L\)-loop local candidates have \(2L+2\) derivatives.  At two loops
the first on-shell pure-gravity candidate is cubic in curvature.

\begin{theorem}[Finite Einstein coupling space is not perturbatively closed]
\label{thm:v4v94-gravity}
Pure four-dimensional Einstein gravity has a nonzero on-shell two-loop
ultraviolet divergence proportional to a curvature-cubic invariant.
Consequently the cosmological, Einstein--Hilbert, and curvature-squared
owners of V90 do not form a closed perturbative coupling space once
graviton loops are included.
\end{theorem}
\begin{proof}
Lemma~\ref{lem:v4v94-power} permits a six-derivative invariant at two
loops.  The explicit Goroff--Sagnotti calculation proves that the
coefficient of the independent on-shell Riemann-cubic counterterm is
nonzero (Nucl.\ Phys.\ B \textbf{266} (1986), 709--736,
DOI \(10.1016/0550\text{-}3213(86)90193\text{-}8\)).
That invariant cannot be absorbed by renormalizing operators with at most
four derivatives.  Hence the finite V90 set is not closed.
\end{proof}

\begin{firewall}[Power counting and nonzero coefficients]
\label{fw:v4v94-power}
Equation \eqref{eq:v4v94-power} lists allowed counterterm degrees; it does
not prove that every allowed invariant occurs at every loop.  The
two-loop conclusion uses the separate nonzero coefficient theorem.  Claims
at higher loops require further calculations or a nonperturbative fixed
point.
\end{firewall}

\subsection{Combined obstruction for the present UV route}

\begin{theorem}[No full UV continuum from the current finite Gaussian
trajectory]
\label{thm:v4v94-combined}
Assume simultaneously:
\begin{enumerate}
 \item the ultraviolet field content and gauge group remain exactly those
 of the unmodified Standard Model;
 \item the hypercharge shells remain in the V6 Gaussian-dominance wedge;
 \item the gravitational action is restricted to the finite V90
 weight-at-most-four coupling set; and
 \item terminal core scales are unbounded above.
\end{enumerate}
Then the V93 all-scale running-coupling certificate fails.
\end{theorem}
\begin{proof}
Items \(1\), \(2\), and \(4\) contradict
Theorem~\ref{thm:v4v94-V93fail}.  Independently, item \(3\) contradicts
Theorem~\ref{thm:v4v94-gravity} when graviton loops are included.  Either
obstruction is sufficient.
\end{proof}

This theorem does not say that no SM-compatible quantum gravity exists.  It
says that the present proof cannot reach the full ultraviolet continuum
while keeping both the unmodified hypercharge Gaussian boundary and a
finite Einstein coupling list.

\subsection{A precise ultraviolet-completion certificate}

\begin{definition}[SM--Einstein ultraviolet-completion certificate]
\label{def:v4v94-UVCC}
A UVCC consists of:
\begin{enumerate}
 \item an ultraviolet field and symmetry system whose low-energy matching
 gives the faithful anomaly-free Standard Model determinant line and the
 Einstein sector;
 \item a gauge- and diffeomorphism-invariant regulator with reflection
 positivity, or a proved replacement sufficient for
 Osterwalder--Schrader reconstruction;
 \item a coupling space containing every local operator generated by the
 exact shell map, together with a trajectory defined for all
 \(t=\log(\mu/\mu_0)\ge0\);
 \item positivity of the physical kinetic forms and a relatively compact
 normalized trajectory, possibly approaching an interacting fixed point;
 \item uniform Ward/BRST-complete shell remainder estimates, including the
 large-field sector and two marked insertions;
 \item compatible anomaly, determinant-line, soft-zero, and topological
 transports across every matching scale; and
 \item a cofinal measure limit whose Schwinger functions satisfy symmetry,
 locality, clustering, and reflection positivity.
\end{enumerate}
\end{definition}

\begin{theorem}[What a UVCC would discharge]
\label{thm:v4v94-UVCC}
A UVCC supplies the all-scale coupling premise of V93 and a common owner
for every local residue generated in V90--V93.  Combined with the already
isolated determinant, stopping-core, and zero-mode estimates, it reduces
the remaining continuum proof to the stated uniform shell remainder and
cofinal measure limits.
\end{theorem}
\begin{proof}
Items \(1\), \(3\), and \(4\) provide a trajectory at every terminal scale,
so Theorem~\ref{thm:v4v93-corewise} applies.  Items \(2\), \(5\), and
\(7\) provide the constructive domination and reconstruction gates.
Item \(6\) supplies the determinant-line and topological compatibility.
No residue is assigned outside the coupling space in item \(3\).
\end{proof}

\begin{firewall}[A certificate is not its construction]
\label{fw:v4v94-UVCC}
Definition~\ref{def:v4v94-UVCC} does not prove asymptotic safety, a grand
unified embedding, a string completion, or any other ultraviolet
mechanism.  It records the exact properties that a proposed mechanism must
prove before it can be inserted into this programme.
\end{firewall}

\subsection{Two honest continuation branches}

A finite ultraviolet cutoff below the first failure scale gives an
effective SM--Einstein theory.  V84--V93 can still organize its finite
determinant, local counterterms, topological sectors, and marked
correlators, but cutoff independence beyond that scale is not claimed.

A full continuum requires changing the current ultraviolet route.  The
mathematically live options are a nonabelian hypercharge embedding, a
non-Gaussian joint gauge--gravity fixed point, or additional ultraviolet
degrees of freedom with anomaly-free matching.  The programme will not
choose among them by name.  It will test candidate routes against
Definition~\ref{def:v4v94-UVCC}.

V95 is a compulsory YM/NS audit.  It will then go upstream to the common
nonperturbative requirement rather than repeat perturbative coefficients:
formulate an exact Wilsonian fixed-point compactness theorem whose
hypotheses are strong enough to imply the UVCC trajectory and weak enough
to be testable by the existing renewal shell maps.

\clearpage
\section{V95: companion audit and a Wilsonian inverse-limit theorem}
\label{sec:v4v95}

This is the compulsory V95 audit.  The newest completed Yang--Mills note
is V14 at SHA--256
\[
 \texttt{2985c1e0b699ec6ac91bd16a861e166afeaa49436674da77b98d36e5fd3987a1}.
\]
The current V15 file was byte-identical to V14 at audit time.  The
Navier--Stokes note V31 remains at
\[
 \texttt{f1121b62238ad5491bb7cf03635ea9934942edf573ed8faaa9ee79e1d38a6696}.
\]
No companion file was modified.

YM V14 obtains arbitrary polynomial absorption because a cofinal scalar
heat multiplier occurs in the literal source word.  A stable label alone
does not give contraction.  NS V31 supplies the complementary warning:
compactness of a marked family does not pay the first moment needed to
remove shell marks.  The inverse-limit theorem below gives a coherent
trajectory, but not response summability.

\subsection{Inverse orientation of shell maps}

Let \(K_j\) be the admissible normalized data at ultraviolet depth \(j\).
Integrating the shell between \(j+1\) and \(j\) defines
\[
 \mathcal R_j:K_{j+1}\longrightarrow K_j.
\]
A full ultraviolet ancestor of fixed low-energy data \(x_0\in K_0\) is
\begin{equation}
 x_j\in K_j,\qquad
 \mathcal R_j(x_{j+1})=x_j,\qquad j\ge0.
\label{eq:v4v95-compatible}
\end{equation}
Forward iteration from \(x_0\) has the wrong physical orientation and does
not construct ultraviolet bare data.

\begin{theorem}[Finite-depth compatibility implies an all-scale
trajectory]
\label{thm:v4v95-inverse}
Suppose every \(K_j\) is nonempty compact Hausdorff, every
\(\mathcal R_j\) is continuous, and for every \(N\) there are
\(x_1^{(N)},\ldots,x_N^{(N)}\) satisfying the first \(N\) equations in
\eqref{eq:v4v95-compatible} with the same \(x_0\).  Then an infinite
compatible trajectory exists, and the set of all such trajectories is
compact.
\end{theorem}
\begin{proof}
The product \(\mathcal K=\prod_{j\ge1}K_j\) is compact.  Let
\(C_N\subset\mathcal K\) impose the first \(N\) compatibility equations.
Continuity makes \(C_N\) closed, finite-depth compatibility makes it
nonempty, and \(C_{N+1}\subset C_N\).  Compactness gives
\(\bigcap_NC_N\ne\varnothing\).  This intersection is exactly the
trajectory set and is closed in \(\mathcal K\).
\end{proof}

\begin{corollary}[Surjective compact shell maps]
\label{cor:v4v95-surjective}
If all shell maps are continuous and surjective, every \(x_0\in K_0\) has
an all-scale ultraviolet trajectory.
\end{corollary}
\begin{proof}
Lift \(x_0\) successively to every finite depth and apply
Theorem~\ref{thm:v4v95-inverse}.
\end{proof}

Surjectivity is stronger than needed.  It suffices that \(x_0\) belong to
every finite composed image
\(\mathcal R_0\cdots\mathcal R_{N-1}(K_N)\).

\subsection{A genuine compact-image criterion}

Closed bounded sets in an infinite-dimensional interaction space are not
compact.  Let \(X_{\rm s}\hookrightarrow X_{\rm w}\) be a compact embedding
of separable Banach spaces.

\begin{theorem}[Strong-to-weak shell compactness]
\label{thm:v4v95-compact}
Suppose every shell image is contained in one uniform strong ball
\(B_{X_{\rm s}}(C)\), and the admissibility constraints defining \(K_j\)
are closed in \(X_{\rm w}\).  Then the \(X_{\rm w}\)-closure of each shell
image is compact.  Nested compact closures of nonempty finite-depth images
supply the compact sets required by
Theorem~\ref{thm:v4v95-inverse}.
\end{theorem}
\begin{proof}
The compact embedding makes the weak-space closure of
\(B_{X_{\rm s}}(C)\) compact.  Intersecting it with a closed admissibility
set preserves compactness.  Continuous finite compositions preserve the
nested image relations.
\end{proof}

This requires actual smoothing or localization compactness.  A marginal
coupling bounded in the same norm supplies no compact gain.

\subsection{Literal stable factors}

Split a shell coordinate as \(x_j=(u_j,v_j)\), where \(u_j\) is
relevant/marginal and the literal shell formula gives
\begin{equation}
 v_j=L_jv_{j+1}+r_j,\qquad \|L_j\|\le\theta<1.
\label{eq:v4v95-stable}
\end{equation}

\begin{lemma}[Exact stable ancestry]
\label{lem:v4v95-stable}
If
\[
 \sum_{k=j}^\infty\theta^{k-j}\|r_k\|<\infty
\]
and \(\theta^{N-j}\|v_N\|\to0\), then
\begin{equation}
 v_j=\sum_{k=j}^\infty
 L_jL_{j+1}\cdots L_{k-1}r_k.
\label{eq:v4v95-series}
\end{equation}
This is the unique solution with that ultraviolet boundary condition.
\end{lemma}
\begin{proof}
Iterating \eqref{eq:v4v95-stable} to depth \(N\) gives the partial sum in
\eqref{eq:v4v95-series} plus
\(L_j\cdots L_{N-1}v_N\).  The boundary term tends to zero and the series
converges absolutely.  The difference of two solutions obeys the
homogeneous equation and vanishes by the same argument.
\end{proof}

\begin{firewall}[A label is not \(L_j\)]
\label{fw:v4v95-literal}
The lemma applies only when the exact shell word exhibits \(L_j\) and proves
its norm bound.  YM V14 works because the scalar heat is present in the
complete word.  Engineering irrelevance alone does not prove
\eqref{eq:v4v95-stable}.  Compactness of \(\{r_j\}\) also does not imply
the weighted first-moment sum; this is the NS V31 boundary.
\end{firewall}

\subsection{Closed constraints and measures}

\begin{proposition}[Closed constraints survive]
\label{prop:v4v95-closed}
If finite-cutoff Ward identities, normalization, reality, and reflection
positivity are closed equalities or inequalities in \(X_{\rm w}\), every
inverse-limit trajectory satisfies them at each level.
\end{proposition}
\begin{proof}
Each coordinate is a limit of admissible finite-depth coordinates in the
closed set \(K_j\).
\end{proof}

\begin{theorem}[Cofinal measure consequence]
\label{thm:v4v95-measure}
Suppose finite-cutoff probability measures are tight on a Polish
distribution space, their required cylinder moments are uniformly
integrable, and their coarse marginals obey the exact shell identities.
Then a cofinal subsequence converges weakly to a probability measure with
those marginals.  If reflection positivity holds at each cutoff, it holds
for the limiting cylinder functionals.
\end{theorem}
\begin{proof}
Prokhorov compactness supplies a weakly convergent subsequence.  Exact
marginals pass to bounded cylinder functions, and uniform integrability
extends this to the stated moments.  Reflection positivity is
nonnegativity of expectations of quadratic reflected cylinder
functionals, so it passes through the same convergence.
\end{proof}

\subsection{Why compactness cannot evade V94}

\begin{theorem}[Hypercharge excludes the naive compact inverse system]
\label{thm:v4v95-hypercharge}
Under the V6 Gaussian-dominance hypotheses, compact admissible sets with
positive hypercharge kinetic coefficient cannot contain finite-depth
ancestors of the fixed \(x_0\) for every depth.
\end{theorem}
\begin{proof}
Otherwise Theorem~\ref{thm:v4v95-inverse} would give a positive all-scale
trajectory, contradicting Theorem~\ref{thm:v4v94-V93fail}.
\end{proof}

Likewise, the finite weight-at-most-four Einstein coupling list is not
invariant under the two-loop shell map by
Theorem~\ref{thm:v4v94-gravity}.  Compactness diagnoses the missing
ultraviolet mechanism; it does not create one.

\subsection{Renewal inverse-limit certificate}

\begin{definition}[RILC]
\label{def:v4v95-RILC}
A proposed ultraviolet completion satisfies the renewal inverse-limit
certificate if it supplies:
\begin{enumerate}
 \item compact admissible sets containing every generated operator;
 \item continuous Ward-complete shell maps;
 \item finite-depth ancestors of the matched SM--Einstein datum at every
 depth;
 \item literal stable factors and summable source moments;
 \item uniform two-marked estimates and tight cutoff measures; and
 \item closed anomaly, positivity, and reflection constraints.
\end{enumerate}
\end{definition}

\begin{theorem}[RILC implication]
\label{thm:v4v95-RILC}
The RILC gives an all-scale coherent trajectory, compactness of its solution
set, a cofinal measure subsequence, and inheritance of the closed Ward and
reflection-positive constraints.  It discharges the trajectory and measure
compactness parts of Definition~\ref{def:v4v94-UVCC}.
\end{theorem}
\begin{proof}
Apply Theorems~\ref{thm:v4v95-inverse},
\ref{thm:v4v95-compact}, and \ref{thm:v4v95-measure}, together with
Proposition~\ref{prop:v4v95-closed}.  Lemma~\ref{lem:v4v95-stable}
constructs the stable coordinates.
\end{proof}

The RILC remains a certificate rather than a construction.  V96 will make
its first analytic hypothesis quantitative: define strong and weak
weighted kernel norms for one renewal shell, separate smoothing irrelevant
kernels from marginal local coefficients, and prove compactness of the
former by a weighted Rellich argument.  Hypercharge and gravitational
marginal directions remain in the finite-depth compatibility gate.

\clearpage
\section{V96: weighted Rellich compactness for the irrelevant kernel sector}
\label{sec:v4v96}

V95 reduces an all-scale trajectory to compact shell images and finite-depth
compatibility.  This note constructs the compactness mechanism for
quasilocal irrelevant interactions.  Local marginal distributions are
removed first, translations are anchored at one insertion, and the strong
norm has both one extra derivative and a stronger exponential/arity weight.
The resulting inclusion is compact by a quantitative Rellich argument.

\subsection{Rooted connected kernels}

For an \(n\)-point connected kernel, fix \(x_1=0\) and write
\[
 y=(y_2,\ldots,y_n),\qquad y_i=x_i-x_1
 \in\mathbb R^{4(n-1)}.
\]
Let \(\ell_n(y)\) be the minimum total length of a spanning tree on
\(\{0,y_2,\ldots,y_n\}\).  Then
\begin{equation}
 \max_{2\le i\le n}|y_i|\le\ell_n(y).
\label{eq:v4v96-coercive}
\end{equation}
Thus \(\{\ell_n\le R\}\) is bounded.  Bundle and species indices take
values in fixed finite-dimensional fibres and are included in the norm.

Fix an integer \(m\ge0\), weights
\[
 0<\mu_{\rm w}<\mu_{\rm s},
 \qquad 0<a_{\rm w}<a_{\rm s},
\]
and define
\begin{align}
 \|F\|_{\mathfrak X_{\rm s}}
 &:=
 \sum_{n\ge2}a_{\rm s}^n
 \sum_{|\alpha|\le m+1}
 \|e^{\mu_{\rm s}\ell_n}\partial^\alpha F_n\|_{L^2},
\label{eq:v4v96-strong}\\
 \|F\|_{\mathfrak X_{\rm w}}
 &:=
 \sum_{n\ge2}a_{\rm w}^n
 \sum_{|\alpha|\le m}
 \|e^{\mu_{\rm w}\ell_n}\partial^\alpha F_n\|_{L^2}.
\label{eq:v4v96-weak}
\end{align}
The zero-, one-, and two-marked sectors are a finite direct sum of spaces
of the same form, with marked species labels retained.

\begin{lemma}[Uniform spatial tail]
\label{lem:v4v96-tail}
For every \(R>0\),
\begin{equation}
 \sum_{n\ge2}a_{\rm w}^n
 \sum_{|\alpha|\le m}
 \|1_{\{\ell_n>R\}}e^{\mu_{\rm w}\ell_n}
   \partial^\alpha F_n\|_2
 \le
 e^{-(\mu_{\rm s}-\mu_{\rm w})R}
 \|F\|_{\mathfrak X_{\rm s}}.
\label{eq:v4v96-tail}
\end{equation}
\end{lemma}
\begin{proof}
On \(\{\ell_n>R\}\),
\[
 e^{\mu_{\rm w}\ell_n}
 \le e^{-(\mu_{\rm s}-\mu_{\rm w})R}
    e^{\mu_{\rm s}\ell_n}.
\]
Also \(a_{\rm w}^n\le a_{\rm s}^n\).  Sum the strong-norm terms with
\(|\alpha|\le m\).
\end{proof}

\begin{lemma}[Uniform arity tail]
\label{lem:v4v96-arity}
For every \(N\ge2\),
\begin{equation}
 \sum_{n>N}a_{\rm w}^n
 \sum_{|\alpha|\le m}
 \|e^{\mu_{\rm w}\ell_n}\partial^\alpha F_n\|_2
 \le
 \left(\frac{a_{\rm w}}{a_{\rm s}}\right)^{N+1}
 \|F\|_{\mathfrak X_{\rm s}}.
\label{eq:v4v96-arity}
\end{equation}
\end{lemma}
\begin{proof}
For \(n>N\),
\(a_{\rm w}^n\le(a_{\rm w}/a_{\rm s})^{N+1}a_{\rm s}^n\), and the weaker
spatial and derivative weights are bounded by the strong ones.
\end{proof}

\subsection{Compact embedding}

\begin{theorem}[Weighted rooted Rellich theorem]
\label{thm:v4v96-Rellich}
The inclusion
\begin{equation}
 \mathfrak X_{\rm s}\hookrightarrow\mathfrak X_{\rm w}
\label{eq:v4v96-embedding}
\end{equation}
is compact.  The same is true for the finite direct sum through two marked
sectors.
\end{theorem}
\begin{proof}
Let \(\{F^{(k)}\}\) be bounded in \(\mathfrak X_{\rm s}\).  Choose
\(\varepsilon>0\).  Lemma~\ref{lem:v4v96-arity} gives \(N\) for which the
weak arity tail is below \(\varepsilon\), uniformly in \(k\).  Then
Lemma~\ref{lem:v4v96-tail} gives \(R\) for which the part with
\(\ell_n>R\) is below \(\varepsilon\), simultaneously for
\(2\le n\le N\).

For each of the finitely many remaining \(n\), equation
\eqref{eq:v4v96-coercive} places \(\{\ell_n\le R\}\) inside a bounded
domain of \(\mathbb R^{4(n-1)}\).  Multiplication by a fixed cutoff equal
to one on that set maps the strong bound to a uniform \(H^{m+1}\) bound on
a slightly larger bounded domain.  Rellich--Kondrachov gives a subsequence
converging in \(H^m\).  A finite diagonal extraction handles all
\(2\le n\le N\) and all finitely many species and marked sectors.

The convergent interior pieces plus the two uniform tails make that
subsequence Cauchy in \(\mathfrak X_{\rm w}\).  Completeness gives a weak
norm limit, proving compactness.
\end{proof}

\begin{corollary}[Quantitative finite-rank approximation]
\label{cor:v4v96-finite-rank}
For every strong ball \(B_{\mathfrak X_{\rm s}}(C)\) and
\(\varepsilon>0\), there is a finite-rank map
\(\Pi_\varepsilon:\mathfrak X_{\rm s}\to\mathfrak X_{\rm w}\) such that
\[
 \sup_{\|F\|_{\mathfrak X_{\rm s}}\le C}
 \|F-\Pi_\varepsilon F\|_{\mathfrak X_{\rm w}}<\varepsilon.
\]
\end{corollary}
\begin{proof}
Choose \(N,R\) by the two tail lemmas.  On the finitely many bounded
domains, approximate the compact Sobolev embeddings by finite-rank maps.
Extend these maps by zero to the discarded arities and spatial tails.
\end{proof}

\subsection{Why rooting and norm gaps are substantive}

\begin{proposition}[Translation obstruction without a root]
\label{prop:v4v96-translation}
On an unrooted translation-invariant kernel space, a nonzero compactly
supported kernel and its translates form a bounded strong-norm sequence
with no strongly convergent weak-norm subsequence.
\end{proposition}
\begin{proof}
Translation preserves unweighted Sobolev norms.  Choose translates with
pairwise disjoint supports.  The weak norm of the difference of two
translates is the square root of the sum of their squared norms, bounded
away from zero.  Hence no subsequence is Cauchy.
\end{proof}

\begin{proposition}[A single critical norm is not compact]
\label{prop:v4v96-single}
The identity from a Sobolev space to itself on any infinite-dimensional
closed subspace is not compact.  Therefore boundedness in one critical
kernel norm cannot replace the derivative, exponential, or arity gaps in
Theorem~\ref{thm:v4v96-Rellich}.
\end{proposition}
\begin{proof}
An orthonormal sequence in the closed subspace is bounded and has pairwise
distance \(\sqrt2\), so its image has no convergent subsequence.
\end{proof}

Rooting removes only the global translation orbit.  It does not identify
physical insertion points or discard their connected relative geometry.

\subsection{Local terms are separate coordinates}

A local operator has a distributional kernel supported on the total
diagonal and need not belong to the weighted \(L^2\) spaces above.  Let
\(\Pi_{\rm loc}\) be the prescribed Taylor/symbol extraction onto the local
operators retained by the ultraviolet candidate, and write
\[
 \mathcal I=m+F,\qquad
 m:=\Pi_{\rm loc}\mathcal I,\qquad
 F:=(1-\Pi_{\rm loc})\mathcal I.
\]
The coordinate \(m\) contains relevant and marginal couplings, including
the hypercharge and gravitational tower; \(F\) is the rooted quasilocal
remainder.

\begin{firewall}[Delta kernels are not Rellich compact]
\label{fw:v4v96-local}
Theorem~\ref{thm:v4v96-Rellich} does not prove compactness of a sequence of
local delta kernels.  Those terms must lie in an independently compact
admissible coupling set or be controlled by a different ultraviolet fixed
point.  Hiding them inside \(F\) would falsely classify the V94
obstructions as smoothing directions.
\end{firewall}

\subsection{A literal smoothing mechanism}

Let \(C\) be a rescaled shell kernel satisfying, for some
\(\mu_{\rm s}>0\),
\begin{equation}
 \sum_{|\gamma|\le1}
 \int_{\mathbb R^4}e^{\mu_{\rm s}|x|}
 |\partial^\gamma C(x)|\,dx\le C_{\rm sh}.
\label{eq:v4v96-shellkernel}
\end{equation}

\begin{lemma}[One shell convolution gains one derivative]
\label{lem:v4v96-smoothing}
For every \(f\) with finite weighted \(H^m\) norm,
\begin{equation}
 \|C*f\|_{H^{m+1}_{\mu_{\rm s}}}
 \le C'_{\rm sh}\|f\|_{H^m_{\mu_{\rm s}}},
\label{eq:v4v96-smoothing}
\end{equation}
where a slightly smaller exponential weight may be used at no cost.
The same estimate holds when the convolution acts in one relative
coordinate of a rooted connected kernel.
\end{lemma}
\begin{proof}
For derivatives of order at most \(m\), commute derivatives with
convolution.  For order \(m+1\), place one derivative on \(C\).
The inequality
\(e^{\mu|x|}\le e^{\mu|x-y|}e^{\mu|y|}\) and weighted Young convolution
give every term in \eqref{eq:v4v96-smoothing}.  The relative-coordinate
version is identical with the other variables treated as parameters.
\end{proof}

A nonlinear contraction of several kernels preserves the weighted tree
length because joining rooted trees gives a spanning tree for the output.
Weighted Young and Cauchy--Schwarz estimates therefore preserve
quasilocality.  A derivative gain, however, follows only when every
surviving connected remainder word contains an actual shell convolution
after local extraction.

\begin{theorem}[Compact irrelevant shell image]
\label{thm:v4v96-shellcompact}
Suppose a Ward-complete shell map has the exact split
\[
 \mathcal R(m,F)=
 \bigl(\mathcal R_{\rm loc}(m,F),
       \mathcal R_{\rm irr}(m,F)\bigr),
\]
where:
\begin{enumerate}
 \item the local image lies in a compact admissible set \(K_{\rm loc}\);
 \item the irrelevant image, through two marks, is uniformly bounded in
       \(\mathfrak X_{\rm s}\); and
 \item its source formula contains the literal smoothing/negative-scaling
       factors used to prove that strong bound.
\end{enumerate}
Then the shell image is relatively compact in
\(K_{\rm loc}\times\mathfrak X_{\rm w}\).
\end{theorem}
\begin{proof}
The first factor is compact by assumption.  The second is relatively
compact by Theorem~\ref{thm:v4v96-Rellich}.  A product of two relatively
compact sets has compact closure.  Item 3 records the analytic source of
item 2 and prevents an engineering-label shortcut.
\end{proof}

\begin{corollary}[Input to the V95 inverse limit]
\label{cor:v4v96-V95}
If every finite-depth shell map satisfies
Theorem~\ref{thm:v4v96-shellcompact} with uniform constants and has a
nonempty matched ancestor, the compact closures of its finite-depth images
supply the compact sets in Theorem~\ref{thm:v4v95-inverse}.
\end{corollary}
\begin{proof}
Apply Theorem~\ref{thm:v4v96-shellcompact} at every depth and take the
nested closures described in
Theorem~\ref{thm:v4v95-compact}.
\end{proof}

V96 proves the functional-analytic compactness of the irrelevant kernel
sector.  It does not prove the uniform strong bound for the complete
SM--Einstein shell map.  V97 will inspect the actual local-extraction
remainder: a Taylor formula will show which source words acquire an external
momentum factor, and the shell heat/resolvent kernel will be tested for the
literal one-derivative gain required by
Theorem~\ref{thm:v4v96-shellcompact}.

\clearpage
\section{V97: covariant local extraction and graphwise shell smoothing}
\label{sec:v4v97}

V96 proves compactness once the irrelevant image is uniformly bounded in a
strong weighted kernel norm.  V97 identifies the literal source of that
bound at finite graph order.  A covariant Taylor projector removes diagonal
local distributions.  Every remaining connected graph must contain a
heat-shell tree joining its external supports.  Weighted convolution along
that tree gives the derivative and localization gains required by V96.

\subsection{Exact Taylor remainder}

In a rescaled normal-coordinate chart impose momentum conservation by
\[
 p_n=-\sum_{i=1}^{n-1}p_i,\qquad p=(p_1,\ldots,p_{n-1}).
\]
Let \(A(p)\) be a smooth amputated shell amplitude near zero and let
\(T^dA\) be its Taylor polynomial of total degree \(d\).

\begin{lemma}[Integral Taylor remainder]
\label{lem:v4v97-Taylor}
For every \(d\ge0\),
\begin{equation}
 (1-T^d)A(p)
 =
 (d+1)\sum_{|\alpha|=d+1}\frac{p^\alpha}{\alpha!}
 \int_0^1(1-t)^d(\partial^\alpha A)(tp)\,dt.
\label{eq:v4v97-Taylor}
\end{equation}
Consequently
\[
 |(1-T^d)A(p)|
 \le C_d|p|^{d+1}
 \max_{\substack{|\alpha|=d+1\\|q|\le|p|}}
 |\partial^\alpha A(q)|.
\]
\end{lemma}
\begin{proof}
Apply the one-variable Taylor formula with integral remainder to
\(f(t)=A(tp)\).  The chain rule gives
\[
 f^{(d+1)}(t)=
 \sum_{|\alpha|=d+1}\frac{(d+1)!}{\alpha!}
 p^\alpha\partial^\alpha A(tp).
\]
Substitution and the triangle inequality prove both statements.
\end{proof}

The polynomial \(T^dA\) is a finite sum of derivatives of the total
diagonal delta distribution in position space and belongs to the local
coordinate of V96.  The exact remainder is the nonlocal kernel.

On a curved core, use symmetrized covariant derivatives in geodesic normal
coordinates and parallel transport all fibre indices to the root.  Two such
choices differ by curvature polynomials of degree at most \(d\), hence by a
permitted local-coordinate change transported by the V90 cocycle.

\subsection{Shell scale and heat kernels}

Use a proper-time shell covariance
\[
 C_j=\int_{\tau_{j+1}}^{\tau_j}e^{-tP}B(t)\,dt.
\]
After rescaling to unit shell size, bounded geometry gives, for fixed \(M\),
\begin{equation}
 \sum_{|\gamma|\le M}\int e^{\mu|x-y|}
 |\nabla_x^\gamma C_j(x,y)|\,dy\le C_{M,\mu},
\label{eq:v4v97-heatkernel}
\end{equation}
uniformly in \(j\), with \(\mu\) below the Gaussian reserve.

\begin{lemma}[Restoring the physical scale]
\label{lem:v4v97-scale}
If the shell momentum is \(K_j\) and an amplitude has superficial degree
\(\omega\), then
\[
 |\partial_p^\alpha A_j(p)|
 \le C_\alpha K_j^{\omega-|\alpha|}.
\]
Therefore
\begin{equation}
 |(1-T^d)A_j(p)|
 \le C_dK_j^{\omega-d-1}|p|^{d+1}.
\label{eq:v4v97-negative}
\end{equation}
\end{lemma}
\begin{proof}
Write \(A_j(p)=K_j^\omega\widehat A_j(p/K_j)\).
Each momentum derivative gives \(K_j^{-1}\), and the Taylor lemma completes
the estimate.
\end{proof}

For \(d=\lfloor\omega\rfloor\), the exponent in
\eqref{eq:v4v97-negative} is negative.  This is the literal negative
scaling factor created by local extraction.

\subsection{Connected shell trees}

Let \(G\) be a fixed connected contraction graph.  Its vertices carry
rooted kernels and its shell edges satisfy
\eqref{eq:v4v97-heatkernel}.  Choose a spanning tree
\(\mathcal T\subset G\).

\begin{lemma}[Tree-length allocation]
\label{lem:v4v97-tree}
For every internal-vertex assignment,
\begin{equation}
 \ell_{\rm out}
 \le\sum_{e\in\mathcal T}|x_{e,+}-x_{e,-}|
   +\sum_{v\in G}\ell_v.
\label{eq:v4v97-tree}
\end{equation}
\end{lemma}
\begin{proof}
Join minimizing internal trees at the vertices with the graph spanning
tree.  The union connects all external points.  Deleting cycles produces an
external spanning tree no longer than the union.
\end{proof}

\begin{theorem}[Graphwise strong kernel bound]
\label{thm:v4v97-graph}
Fix \(m\) and \(G\).  Suppose:
\begin{enumerate}
 \item every input nonlocal vertex has finite
       \(\mathfrak X_{\rm w}\) norm;
 \item every external component is joined to the root by a literal
       heat-shell edge satisfying \eqref{eq:v4v97-heatkernel} through
       order \(m+1\); and
 \item every superficially nonnegative subgraph is replaced by its exact
       Taylor remainder after assigning its local jet to the local
       coordinate.
\end{enumerate}
Then the renormalized nonlocal output \(F_G\) lies in
\(\mathfrak X_{\rm s}\) and
\begin{equation}
 \|F_G\|_{\mathfrak X_{\rm s}}
 \le C_G\prod_{v\in G}\|F_v\|_{\mathfrak X_{\rm w}},
\label{eq:v4v97-graph}
\end{equation}
including versions with one or two marked vertices.
\end{theorem}
\begin{proof}
Use \eqref{eq:v4v97-tree} to distribute the output exponential weight among
spanning-tree edges and input vertices.  Integrate leaves toward the root.
Weighted Young convolution uses \eqref{eq:v4v97-heatkernel}; non-tree edges
are bounded by their weighted \(L^1\) or Hilbert--Schmidt norms.

For derivatives through order \(m\), use the input weak bounds and the
product rule.  For the additional derivative, follow the path from the
differentiated external coordinate to the root and place one derivative on
the first heat-shell kernel.  Its marked derivatives are bounded by
\eqref{eq:v4v97-heatkernel}.  Taylor local terms were removed, while their
remainders carry \eqref{eq:v4v97-negative}.  One or two marks merely select
the corresponding factors in the same tree integration.
\end{proof}

\begin{firewall}[An internal loop need not smooth externally]
\label{fw:v4v97-loop}
A tadpole shell loop at one vertex may alter a local coefficient without
joining two external supports.  It belongs to the Taylor-extracted owner.
Theorem~\ref{thm:v4v97-graph} requires a shell edge on every
external-to-root path.
\end{firewall}

\subsection{Subdivergences and forest order}

For each divergent proper subgraph \(\gamma\), let \(T_\gamma\) be its
covariant Taylor projector through its superficial degree.

\begin{definition}[Covariant forest subtraction]
\label{def:v4v97-forest}
Define
\[
 R_G=(1-T_G)\sum_{\mathcal F}
 \prod_{\gamma\in\mathcal F}(-T_\gamma)A_G,
\]
where the sum is over forests of disjoint or nested divergent proper
subgraphs, with inner projectors applied first.
\end{definition}

\begin{proposition}[Every forest term has one owner]
\label{prop:v4v97-forest}
The forest formula decomposes \(A_G\) into local jets assigned to declared
ultraviolet operators and one nonlocal remainder satisfying
\eqref{eq:v4v97-graph}.  No divergent subgraph is subtracted twice.
\end{proposition}
\begin{proof}
Use \(1=T_\gamma+(1-T_\gamma)\) from inner to outer subgraphs.  Nested
application assigns each inner jet before the quotient graph is expanded.
The alternating forest sum is inclusion--exclusion for overlapping choices.
After every nonnegative subgraph and the overall graph are treated, all
remainders have negative Taylor degree and the heat-tree argument applies.
\end{proof}

\subsection{The graph sum}

Let \(w_j(G)\) contain symmetry factors and local couplings.

\begin{theorem}[Sufficient all-graph condition]
\label{thm:v4v97-sum}
If, uniformly in shell and cutoff,
\begin{equation}
 \sum_{G\ni{\rm root}}e^{\theta|G|}|w_j(G)|C_G
 \prod_{v\in G}\|F_v\|_{\mathfrak X_{\rm w}}
 \le C_{\rm irr},
\label{eq:v4v97-sum}
\end{equation}
and the same holds with up to two marked choices, then the complete
irrelevant shell image is bounded in \(\mathfrak X_{\rm s}\) and relatively
compact in \(\mathfrak X_{\rm w}\).
\end{theorem}
\begin{proof}
Absolute convergence permits summing
\eqref{eq:v4v97-graph} and its marked versions.  The strong norm is bounded
by \(C_{\rm irr}\).  Apply Theorem~\ref{thm:v4v96-Rellich}.
\end{proof}

\begin{firewall}[Finite graph order is not the shell map]
\label{fw:v4v97-finite}
The graphwise theorem proves literal smoothing at every fixed graph.  It
does not prove \eqref{eq:v4v97-sum}, restore the large-field sector, or
show that the ultraviolet coupling space contains all gravitational forest
counterterms.
\end{firewall}

V97 closes the graphwise part of the V96 strong bound and isolates one
summation inequality.  V98 will attack it with a Ward-complete polymer norm,
keeping the small-field analytic sum and large-field restoration in one
exact decomposition and importing the earlier V6 criterion without
assuming its hypotheses for the joint SM--Einstein shell.

\clearpage
\section{V98: Ward-complete polymer summation and the conformal-factor barrier}
\label{sec:v4v98}

V97 reduces compactness of the irrelevant shell image to an all-graph sum.
The small-field part can be closed by a Banach-valued polymer criterion,
provided Ward-related graphs are assembled before norms are taken.  The
large-field restoration is more serious: the Yang--Mills rare-event lemma
presupposes a probability measure.  The naive Euclidean
Einstein--Hilbert weight has no such measure because its conformal mode is
unbounded below.

\subsection{Ward-complete activities}

Let \(\Lambda\) be the rescaled shell cell set.  A polymer is a finite
connected \(X\subset\Lambda\), and two polymers are incompatible when they
overlap.  For each \(X\), assemble all gauge, ghost, Goldstone, fermion,
Higgs, and metric graphs required by the finite-cutoff Ward/BRST identity
into one activity
\[
 w(X)\in\mathfrak X_{\rm s}.
\]
Taylor forest subtractions are performed inside this sum.  Gauge-variant
longitudinal pieces are not normed separately.

For \(a>0\), assume
\begin{equation}
 \sup_{x\in\Lambda}
 \sum_{X\ni x}e^{a|X|}\|w(X)\|_{\mathfrak X_{\rm s}}
 \le a.
\label{eq:v4v98-KP}
\end{equation}

\begin{lemma}[Local condition implies the KP condition]
\label{lem:v4v98-KP}
For every polymer \(X_0\),
\begin{equation}
 \sum_{X:X\cap X_0\ne\varnothing}
 e^{a|X|}\|w(X)\|_{\mathfrak X_{\rm s}}
 \le a|X_0|.
\label{eq:v4v98-KPX}
\end{equation}
\end{lemma}
\begin{proof}
Bound the left side by summing first over a chosen intersection point:
\[
 \sum_{x\in X_0}\sum_{X\ni x}
 e^{a|X|}\|w(X)\|_{\mathfrak X_{\rm s}}
 \le\sum_{x\in X_0}a=a|X_0|.
\]
Multiple counting only enlarges the bound.
\end{proof}

\begin{theorem}[Banach-valued connected polymer sum]
\label{thm:v4v98-polymer}
Under \eqref{eq:v4v98-KP}, the logarithm of the hard-core polymer
partition function converges absolutely in the rooted
\(\mathfrak X_{\rm s}\) norm.  More precisely, with
\(\phi^T(X_1,\ldots,X_n)\) the Ursell coefficient,
\begin{equation}
 \sum_{n\ge1}\frac1{(n-1)!}
 \sum_{\substack{X_1\ni x\\X_2,\ldots,X_n}}
 |\phi^T(X_1,\ldots,X_n)|
 \prod_{i=1}^n\|w(X_i)\|_{\mathfrak X_{\rm s}}
 e^{a\sum_i|X_i|}
 \le a.
\label{eq:v4v98-rooted}
\end{equation}
The resulting irrelevant shell image is uniformly bounded in
\(\mathfrak X_{\rm s}\).
\end{theorem}
\begin{proof}
The tree-graph inequality bounds the absolute Ursell coefficient by the
sum over rooted spanning trees of incompatibility edges.  Sum leaves
toward the root.  For a leaf attached to \(X_i\),
Lemma~\ref{lem:v4v98-KP} pays the sum over all incompatible leaf polymers
by \(a|X_i|\).  The exponential \(e^{a|X_i|}\) supplies exactly the
standard combinatorial reserve for assigning its children.  Induction over
the rooted tree gives the Kotecky--Preiss bound
\eqref{eq:v4v98-rooted}.  Submultiplicativity of the weighted kernel norm
allows the same leaf elimination for Banach-valued activities.
\end{proof}

\begin{corollary}[Two marked sums]
\label{cor:v4v98-marks}
Suppose \eqref{eq:v4v98-KP} holds with \(a\) replaced by a strictly larger
reserve \(a+\delta\), uniformly for the unmarked activities and their first
two normalized marks.  Then the connected logarithm has first two marked
derivatives satisfying \eqref{eq:v4v98-rooted} with weight \(a\).
\end{corollary}
\begin{proof}
One derivative chooses one activity in a cluster; two derivatives choose
one twice-marked activity or an ordered pair of once-marked activities.
The resulting factors \(n\) and \(n^2\) are bounded by
\(C_\delta e^{\delta n}\), which is absorbed by the reserve.
Termwise differentiation follows from absolute convergence.
\end{proof}

Combined with Theorems~\ref{thm:v4v97-graph} and
\ref{thm:v4v96-Rellich}, this proves compactness of the small-field
irrelevant image whenever the Ward-complete activities satisfy the one
local estimate \eqref{eq:v4v98-KP}.

\begin{firewall}[Ward assembly precedes absolute values]
\label{fw:v4v98-Ward}
The condition is imposed on the complete Ward activity \(w(X)\), not on
each gauge, ghost, Goldstone, or metric graph separately.  Taking absolute
values first can destroy the longitudinal cancellations and create a
false mass or conformal owner.
\end{firewall}

\subsection{Exact small/large decomposition}

Let \(\chi_{\rm sf}+\chi_{\rm lf}=1\) be a gauge- and
diffeomorphism-invariant finite-cutoff partition.  For every observable
\(\mathcal O\),
\[
 \mathbb E\mathcal O
 =
 \mathbb E(\chi_{\rm sf}\mathcal O)
 +\mathbb E(\chi_{\rm lf}\mathcal O).
\]
If the underlying weight is a probability measure, the V6 rare-event
argument gives, for \(p>1\),
\[
 |\mathbb E(\chi_{\rm lf}\mathcal O)|
 \le\|\mathcal O\|_{L^p}
    \mathbb P(\chi_{\rm lf}\ne0)^{1-1/p}.
\]
Thus a stretched-exponential large-field probability pays every fixed
polynomial insertion moment.  This statement is conditional on existence
of the probability measure.

\subsection{Euclidean conformal factor}

Use the Euclidean sign
\begin{equation}
 S_{\rm EH}(g)
 =
 -\frac1{16\pi G}\int_X\sqrt g\,(R(g)-2\Lambda).
\label{eq:v4v98-EH}
\end{equation}
Let \(g=e^{2\varphi}\bar g\) on a closed four-manifold.  The conformal
transformation law is
\[
 R(g)=e^{-2\varphi}
 \left(\bar R-6\bar\Delta\varphi-6|\bar\nabla\varphi|^2\right).
\]

\begin{lemma}[Exact conformal Einstein action]
\label{lem:v4v98-conformal}
On a closed core,
\begin{equation}
 \int_X\sqrt g\,R(g)
 =
 \int_X\sqrt{\bar g}\,e^{2\varphi}
 \left(\bar R+6|\bar\nabla\varphi|^2\right).
\label{eq:v4v98-conformal}
\end{equation}
\end{lemma}
\begin{proof}
Insert \(\sqrt g=e^{4\varphi}\sqrt{\bar g}\).  Integration by parts gives
\[
 -6\int e^{2\varphi}\bar\Delta\varphi
 =12\int e^{2\varphi}|\bar\nabla\varphi|^2.
\]
Combining this with the original negative gradient term leaves the
coefficient \(+6\).
\end{proof}

\begin{theorem}[Conformal-factor unboundedness]
\label{thm:v4v98-unbounded}
For \(G>0\), the action \eqref{eq:v4v98-EH} is unbounded below on every
metric neighborhood that permits arbitrarily high-frequency conformal
variations of fixed nonzero amplitude.  The cosmological term does not
repair this.
\end{theorem}
\begin{proof}
Choose a coordinate ball and a smooth cutoff \(\chi\), and set
\(\varphi_N=\varepsilon\chi(x)\sin(Nx_1)\) with fixed small
\(\varepsilon>0\).  Uniform bounds on \(\varphi_N\) keep
\(e^{2\varphi_N}\) above and below by positive constants, while
\[
 \int e^{2\varphi_N}|\nabla\varphi_N|^2
 \ge c\varepsilon^2N^2-O(N).
\]
By Lemma~\ref{lem:v4v98-conformal}, the curvature integral grows like
\(+N^2\), so its negative coefficient in \(S_{\rm EH}\) drives the action
to \(-\infty\).  The volume term
\(\int e^{4\varphi_N}\sqrt{\bar g}\) stays uniformly bounded because the
amplitude is fixed.
\end{proof}

\begin{corollary}[No action-derived Einstein rare-event bound]
\label{cor:v4v98-no-probability}
The factor \(e^{-S_{\rm EH}(g)}\) supplies no large-field suppression on
the real conformal contour.  Consequently the Yang--Mills V6 rare-event
restoration cannot be imported into the full metric sector from the
Einstein--Hilbert action alone.
\end{corollary}
\begin{proof}
Along the sequence in Theorem~\ref{thm:v4v98-unbounded},
\(e^{-S_{\rm EH}}\to\infty\) exponentially in \(N^2\), rather than
suppressing the sequence.  Thus an action-derived probability tail bound
fails; a claim about a normalized metric functional requires additional
measure or contour data.
\end{proof}

\subsection{What a gravity integration prescription must prove}

A proposed metric prescription may rotate the conformal contour, impose
causal/Lorentzian variables, add higher-derivative stabilizers, or use a
nonperturbative fixed-point measure.  The programme does not choose one by
name.  It requires:
\begin{enumerate}
 \item a well-defined normalized finite-cutoff functional;
 \item exact gauge and diffeomorphism Ward identities on its contour or
 variable space;
 \item a positive/reflection-positive physical observable sector;
 \item a large-field estimate strong enough for two marked activities; and
 \item compatibility with the V90 counterterms and V94 gravitational
 tower.
\end{enumerate}

\begin{theorem}[Conditional completion of the all-graph compactness step]
\label{thm:v4v98-conditional}
If a metric integration prescription satisfies the five properties above,
the Ward-complete small-field activities obey
\eqref{eq:v4v98-KP}, and its large-field sector has the V6 insertion moments
and stretched-exponential rarity, then the complete irrelevant shell image
is relatively compact through two marks.
\end{theorem}
\begin{proof}
Theorem~\ref{thm:v4v98-polymer} and
Corollary~\ref{cor:v4v98-marks} control the small-field image.
The V6 rare-event estimate restores the large-field part without breaking
the Ward packet.  Their sum is uniformly strong-bounded, so
Theorem~\ref{thm:v4v96-Rellich} gives relative compactness.
\end{proof}

V98 proves the polymer summation mechanism and identifies a prior
gravitational obstruction: before large-field restoration, the programme
needs a metric integration prescription that overcomes the real-contour
conformal instability and still supports reflection positivity.  V99 will
go upstream to formulate contour-independent Schwinger-function
requirements and test which parts of the inverse-limit theorem need a
positive microscopic metric measure, rather than assuming that a formal
contour rotation preserves Osterwalder--Schrader positivity.

\clearpage
\section{V99: reflection positivity at the Schwinger-functional level}
\label{sec:v4v99}

V98 proves that the real Euclidean Einstein action does not provide a
large-field probability bound.  This does not logically exclude every
complex contour or non-measure construction, but reflection positivity must
then be proved from the resulting Schwinger functional.  It cannot be
inferred from formal contour rotation.

Let \(\mathcal A_+\) be the algebra of gauge-invariant cylinder
observables supported at positive Euclidean time, and let
\(\Theta:\mathcal A_+\to\mathcal A_-\) be an antilinear time reflection.
A normalized Schwinger functional \(S\) is reflection positive when
\begin{equation}
 \sum_{i,j=1}^m\overline{c_i}c_j
 S\bigl((\Theta F_i)F_j\bigr)\ge0
\label{eq:v4v99-OS}
\end{equation}
for every finite family \(F_i\in\mathcal A_+\).

\begin{theorem}[Closedness of reflection positivity]
\label{thm:v4v99-closed}
Let \(S_N\) be normalized reflection-positive functionals.  If
\[
 S_N((\Theta F)G)\longrightarrow S((\Theta F)G)
\]
for every \(F,G\in\mathcal A_+\), then \(S\) is reflection positive.
No positive microscopic measure is required for this implication.
\end{theorem}
\begin{proof}
For a fixed finite family, form the matrices
\[
 M_N(i,j)=S_N((\Theta F_i)F_j).
\]
Each \(M_N\) is positive semidefinite.  Entrywise convergence gives a
matrix \(M\), and
\(c^*Mc=\lim_Nc^*M_Nc\ge0\) for every \(c\).  This is
\eqref{eq:v4v99-OS}.
\end{proof}

\begin{corollary}[Cofinal use]
\label{cor:v4v99-cofinal}
The V95 cofinal measure theorem may be replaced by convergence of the
reflected cylinder matrix elements themselves.  Positivity of every
finite matrix and their convergence suffice for the limiting OS form.
\end{corollary}
\begin{proof}
Apply Theorem~\ref{thm:v4v99-closed} to the cofinal sequence.
\end{proof}

\subsection{A transfer-matrix certificate}

\begin{theorem}[Positive factorization]
\label{thm:v4v99-factor}
Suppose there are Hilbert spaces \(\mathcal H_N\) and linear maps
\(T_N:\mathcal A_+\to\mathcal H_N\) such that
\begin{equation}
 S_N((\Theta F)G)=\langle T_NF,T_NG\rangle_{\mathcal H_N}.
\label{eq:v4v99-factor}
\end{equation}
Then \(S_N\) is reflection positive.  If the inner products converge on a
dense cylinder algebra, the limit is reflection positive and its OS Hilbert
space is the completion of the quotient by the nullspace.
\end{theorem}
\begin{proof}
The OS quadratic form is
\[
 \sum_{i,j}\overline{c_i}c_j
 \langle T_NF_i,T_NF_j\rangle
 =\left\|\sum_i c_iT_NF_i\right\|^2\ge0.
\]
Pass to the limit by Theorem~\ref{thm:v4v99-closed}.  Quotienting a
positive semidefinite form by its nullspace makes it definite; completion
gives the Hilbert space.
\end{proof}

This factorization may arise from a positive transfer operator even when
the chosen field coordinates do not carry a positive density.  It is an
operator identity that must be constructed.

\subsection{Contour deformation}

Let
\[
 S_\Gamma(F)=Z_\Gamma^{-1}\int_\Gamma F(z)e^{-I(z)}\,\Omega(z)
\]
be a finite-dimensional regulated contour integral.

\begin{lemma}[When contour deformation preserves correlators]
\label{lem:v4v99-contour}
Suppose the differential form
\(F(z)e^{-I(z)}\Omega(z)\) is holomorphic in a region joining
\(\Gamma_0\) and \(\Gamma_1\), the contours are homologous there, and all
boundary-at-infinity integrals vanish, for \(F=1\) and every cylinder
observable under consideration.  Then
\[
 S_{\Gamma_0}(F)=S_{\Gamma_1}(F).
\]
\end{lemma}
\begin{proof}
Stokes' theorem applied to the interpolating chain gives equality of the
unnormalized integrals.  The case \(F=1\) gives equality of normalizations.
\end{proof}

\begin{corollary}[Positivity transported by exact equality]
\label{cor:v4v99-contour}
If \(\Gamma_0\) has a positive factorization
\eqref{eq:v4v99-factor} and the hypotheses of
Lemma~\ref{lem:v4v99-contour} hold for every product
\((\Theta F)G\), then \(\Gamma_1\) defines the same OS form and is
reflection positive.
\end{corollary}
\begin{proof}
All entries of every reflected Gram matrix are equal on the two contours.
\end{proof}

\begin{firewall}[Formal conformal rotation]
\label{fw:v4v99-formal}
Replacing a real conformal variable by an imaginary one does not establish
the hypotheses of Lemma~\ref{lem:v4v99-contour}.  Singularities,
nonvanishing end caps, a changed integration cycle, or failure of
\(\Theta\)-compatibility can alter the reflected correlators.  Positivity
must be proved by exact contour equality or by a direct factorization.
\end{firewall}

\subsection{Gauge quotient and determinant line}

The observables in \(\mathcal A_+\) are gauge invariant.  The complete
fermion determinant is the V74 anomaly-normalized line section, sewn across
the V87 charts.  A local hard or soft determinant factor need not be
reflection real by itself.

\begin{proposition}[Total-line reality condition]
\label{prop:v4v99-line}
Suppose the complete finite-cutoff integrand obeys an antilinear line
isomorphism
\[
 \Theta^*\overline{\mathcal L_{\rm SM}}\cong\mathcal L_{\rm SM}
\]
which fixes the normalized total section and is compatible with the
bosonic contour.  Then reflected products have chart-independent scalar
coefficients.  Hard, soft, and finite crossing factors must be combined
before the OS matrix is formed.
\end{proposition}
\begin{proof}
V74 and V87 identify every chart factorization with the same complete
line.  Applying the antilinear reflection to that line makes the product
\((\Theta F)G\) a scalar.  The overlap cocycle cancels between the two
descriptions.  Separate factors retain compensating Berry and anomaly
connections and therefore need not be scalar-real.
\end{proof}

\subsection{An explicit OS certificate for the gravity prescription}

\begin{definition}[Metric OS certificate]
\label{def:v4v99-MOSC}
A finite-cutoff metric integration prescription satisfies MOSC if:
\begin{enumerate}
 \item it defines normalized reflected cylinder matrix elements;
 \item the contour/variable space and determinant line admit the compatible
 antilinear reflection in Proposition~\ref{prop:v4v99-line};
 \item every finite reflected matrix is positive semidefinite, proved by a
 transfer factorization or exact deformation from one;
 \item the matrix elements have cutoff-uniform domination through the two
 marked observables needed for reconstruction; and
 \item they converge on a dense cylinder algebra.
\end{enumerate}
\end{definition}

\begin{theorem}[MOSC supplies the missing positivity gate]
\label{thm:v4v99-MOSC}
MOSC gives a limiting reflection-positive Schwinger functional and its OS
Hilbert space.  Combined with Euclidean covariance, symmetry, clustering,
and regularity, it supplies the positivity component of
Osterwalder--Schrader reconstruction without assuming a positive
microscopic metric density.
\end{theorem}
\begin{proof}
Items 1--3 give finite-cutoff positivity.  Items 4--5 give convergence of
all reflected Gram entries.  Apply Theorems~\ref{thm:v4v99-closed} and
\ref{thm:v4v99-factor}.  The remaining listed properties are the other OS
axioms and are not consequences of positivity alone.
\end{proof}

V99 closes the logical question raised by the conformal-factor barrier:
microscopic positivity of the metric density is sufficient but not
necessary; positivity of the Schwinger functional is necessary and can be
proved by MOSC.  No present contour has been shown to satisfy MOSC.

V100 will consolidate the fourth-series results into a conditional
continuum assembly theorem and an exact obstruction ledger.  The file will
then be frozen as \(\_f4\), and a new compact V1 will restate the major
proved results, current context, and next acquisition target.

\clearpage
\section{V100: conditional continuum assembly and fourth-series obstruction ledger}
\label{sec:v4v100}

This note closes the fourth series at the level actually proved.  It
assembles the determinant, renormalization, compactness, and positivity
certificates into one conditional continuum theorem.  It does not assert
that their unresolved hypotheses hold.

\begin{theorem}[Conditional SM--Einstein continuum assembly]
\label{thm:v4v100-assembly}
Let a cofinal family of regulated SM--Einstein functionals satisfy:
\begin{enumerate}
 \item the UVCC and RILC of V94--V95, including finite-depth ancestors,
 compact shell images, and an all-scale matched trajectory;
 \item the V90 local heat-counterterm closure, the V91 renormalized
 low-Born/\(\det_5\) decomposition, and uniform bounds on its finite
 one-through-four insertion packet;
 \item the V74 anomaly-normalized determinant-line sewing, the V84--V88
 finite crossing atlas, and exact V75--V77 soft-zero saturation;
 \item the V96--V98 rooted irrelevant-kernel compactness and
 Ward-complete small/large-field polymer estimates through two marks;
 \item MOSC from V99; and
 \item cutoff-uniform Euclidean covariance, permutation symmetry,
 locality, clustering, and the regularity bounds required by the
 Osterwalder--Schrader theorem.
\end{enumerate}
Then a cofinal subsequence of Schwinger functionals converges on the dense
cylinder algebra to a gauge-invariant, anomaly-free, reflection-positive
SM--Einstein Schwinger functional.  Its OS reconstruction yields a Hilbert
space, vacuum, positive Hamiltonian, and local real-time fields with the
matched low-energy Standard Model and gravitational observables.
\end{theorem}
\begin{proof}
Items 1 and 4 give compactness of the shell trajectory and irrelevant
interactions by Theorems~\ref{thm:v4v95-inverse} and
\ref{thm:v4v96-shellcompact}.  Item 2 gives uniform fermion determinant
domination by Theorem~\ref{thm:v4v91-criterion}; item 3 extends it through
topological sectors and determinant zeros without dividing by a gap.
Tightness and uniform integrability yield a cofinal functional limit as in
Theorem~\ref{thm:v4v95-measure}.  Item 5 and
Theorem~\ref{thm:v4v99-closed} give reflection positivity.  Gauge
invariance and anomaly-free line transport pass through the exact Ward and
cocycle identities.  Item 6 supplies the remaining OS axioms, so the
standard reconstruction theorem gives the stated real-time theory.
\end{proof}

\begin{corollary}[What remains after fermion spectral descent]
\label{cor:v4v100-spectral}
The local-to-global fermion spectral problem is no longer an independent
gate: under the intrinsic \(L^4\) comparison, its hard crossings are finite,
action-ranked, collared, and globally sewn.  The remaining obstacles lie in
the ultraviolet shell trajectory, finite marginal amplitudes, metric
integration, and constructive measure estimates.
\end{corollary}
\begin{proof}
This is the combined content of V84--V93.  None of those steps proves the
bosonic shell or MOSC hypotheses.
\end{proof}

\subsection{Exact open gates}

\begin{enumerate}
 \item \textbf{Hypercharge.}  In the V6 Gaussian-dominance wedge the
 positive inverse \(U(1)_Y\) coupling cannot have ancestors at arbitrarily
 high scales.
 \item \textbf{Quantum gravity.}  The Einstein coupling list is not closed:
 the nonzero two-loop curvature-cubic counterterm lies beyond V90.
 \item \textbf{Metric integration.}  The real Euclidean conformal action is
 unbounded below, and no alternative prescription has yet proved MOSC.
 \item \textbf{Marginal finite packet.}  Weak-\(\mathcal S_4\) control
 defines \(\det_5\) but does not bound the renormalized first four Born
 amplitudes on rough critical profiles.
 \item \textbf{All-graph shell bound.}  V97 proves graphwise smoothing and
 V98 proves the polymer theorem, but the joint Ward-complete activities and
 gravity large-field restoration have not been shown to satisfy its local
 hypothesis.
 \item \textbf{Measure and reconstruction.}  Tightness, uniform
 integrability, clustering, and all OS axioms remain certificates rather
 than established estimates.
\end{enumerate}

\begin{theorem}[Present-route no-go boundary]
\label{thm:v4v100-no-go}
The conditional assembly theorem cannot be applied while simultaneously
retaining the unmodified hypercharge Gaussian boundary, a finite
Einstein coupling space, and the real Euclidean Einstein contour.
\end{theorem}
\begin{proof}
The first choice violates
Theorem~\ref{thm:v4v94-V93fail}; the second violates
Theorem~\ref{thm:v4v94-gravity}; the third supplies no action-derived
large-field tail by Corollary~\ref{cor:v4v98-no-probability} and has no
proved MOSC replacement.
\end{proof}

The next series must therefore test an actual ultraviolet and metric
integration mechanism against UVCC, RILC, and MOSC.  Repeating spectral
covers, one-loop gauge coefficients, or formal conformal rotation would not
advance the proof.

\subsection{Fourth-series result ledger}

V84--V100 establish:
\begin{enumerate}
 \item exact weak-\(\mathcal S_4\) finite crossing reduction by a Sylvester
 matrix, with action-paid rank;
 \item an intrinsic global reference connection from the reconstructed
 transition cocycle;
 \item frozen-centre marked charts and quantitative spectral collars;
 \item the order-zero metric-symbol obstruction and its microlocal
 normalization;
 \item the complete four-dimensional heat owner list, boundary cutoff
 moments, and reference-change cocycle;
 \item the exact low-Born/\(\det_5\) separation and Wodzicki-residue scale
 law;
 \item fixed-scale endpoint failure and corewise RG matching;
 \item the inherited hypercharge and two-loop gravity obstructions;
 \item the Wilsonian inverse-limit and weighted Rellich compactness
 theorems;
 \item graphwise covariant forest smoothing and Banach-valued polymer
 summation; and
 \item the conformal-factor barrier and Schwinger-functional OS criterion.
\end{enumerate}

The full SM--Einstein continuum has not been constructed.  The series has
replaced several false shortcuts by exact finite-dimensional,
pseudodifferential, RG, compactness, and positivity statements and has
localized the remaining programme to the six gates above.

\end{document}
